\documentclass[12pt,a4paper,oneside,openany]{book}
\usepackage[UTF8]{ctex}
\usepackage{amsmath,amssymb,amsthm,mathtools}
\usepackage{esint}
\usepackage{geometry}
\geometry{left=2.2cm,right=2.2cm,top=2.2cm,bottom=2.2cm}
\usepackage{graphicx}
\usepackage{fancyhdr}
\usepackage{enumitem}
\usepackage{needspace}
\usepackage{etoolbox}
\usepackage[colorlinks=true,linkcolor=black,bookmarksnumbered=true]{hyperref}

% 常用数学算子补充
\DeclareMathOperator{\sgn}{sgn}
\DeclareMathOperator{\arccot}{arccot}
\DeclareMathOperator{\arsinh}{arsinh}
\DeclareMathOperator{\arcosh}{arcosh}
\DeclareMathOperator{\artanh}{artanh}

% 行距与断行微调（适配行内 \displaystyle 公式的舒适行距）
\linespread{1.35}
\setlength{\lineskiplimit}{3pt}
\setlength{\lineskip}{4pt}
\setlength{\parskip}{3pt plus 1pt minus 1pt}
\setlength{\emergencystretch}{3em}
\setlength{\headheight}{15pt}

% 页眉页脚设置
\pagestyle{fancy}
\fancyhf{}
\fancyhead[L]{\small\leftmark}
\fancyhead[R]{\small\rightmark}
\fancyfoot[C]{\small\thepage}
\renewcommand{\headrulewidth}{0.4pt}
\renewcommand{\footrulewidth}{0pt}

% 部、章与节标题（进入目录、书签与页眉）
\newcommand{\probpart}[1]{%
  \cleardoublepage
  \thispagestyle{empty}
  \vspace*{5cm}
  \begin{center}
    {\Huge\bfseries #1\par}
  \end{center}
  \addcontentsline{toc}{part}{#1}%
  \clearpage
}
\newcommand{\probchapter}[1]{%
  \chapter*{#1}%
  \addcontentsline{toc}{chapter}{#1}%
  \markboth{#1}{}%
}
\newcommand{\probsec}[1]{%
  \par\bigskip
  \needspace{8cm}%
  \section*{#1}%
  \addcontentsline{toc}{section}{#1}%
  \markright{#1}%
}

% ============================================================
% 做题本题目间距控制（可按需在此统一调整 S / M / L 三档留白）
%   S: 简短计算/不等式/定义域/简单极限（约 3~4 题/页）
%   M: 标准证明/中等计算/双小题（约 2~3 题/页）
%   L: 多小题/长证明/复杂作图（约 1~2 题/页）
% ============================================================
\newlength{\probspaceS}
\newlength{\probspaceM}
\newlength{\probspaceL}
\setlength{\probspaceS}{3.8cm}
\setlength{\probspaceM}{6.2cm}
\setlength{\probspaceL}{9.5cm}

\newenvironment{prob}[2][M]{%
  \def\probtier{#1}%
  \par
  \ifstrequal{#1}{S}{\needspace{5.2cm}}{%
  \ifstrequal{#1}{L}{\needspace{11.0cm}}{%
    \needspace{7.6cm}}}%
  \noindent\textbf{#2.}\enspace\ignorespaces
}{%
  \par
  \ifstrequal{\probtier}{S}{\vspace{\probspaceS}}{%
  \ifstrequal{\probtier}{L}{\vspace{\probspaceL}}{%
    \vspace{\probspaceM}}}%
  \pagebreak[1]%
}

% 共用题干（与后续首题尽量保持同页）
\newcommand{\probgroup}[1]{%
  \par\medskip
  \needspace{6.5cm}%
  \noindent\textbf{#1}\par\smallskip
}

% 题目提示
\newcommand{\probhint}[1]{%
  \par\smallskip
  \noindent{\small\textit{提示：}#1}%
}

% 题目插图
\newcommand{\probfig}[2]{%
  \begin{center}
    \includegraphics[width=0.38\linewidth,height=0.26\textheight,keepaspectratio]{figures/#1}\\[3pt]
    {\small #2}
  \end{center}
}

\begin{document}

% 封面页
\begin{titlepage}
  \centering
  \vspace*{4.5cm}
  {\Huge\bfseries 吉米多维奇数学分析习题集\par}
  \vspace{1.2cm}
  {\LARGE 做题本\par}
  \vspace{2.5cm}
  {\large 全八章（共 4562 题）\par}
  \vfill
  {\normalsize 根据 2010 年俄文版中译本整理 \quad 仅含习题与必要提示、题图\par}
  \vspace*{1.5cm}
\end{titlepage}

\frontmatter
\tableofcontents
\clearpage

\mainmatter
\everymath{\displaystyle}
\probpart{第一部分\quad 一元函数}
\probchapter{第一章\quad 分析引论}

\probsec{§1. 实数}

\probgroup{利用数学归纳法求证下列等式对任何正整数 $n$ 皆成立：}

\begin{prob}[M]{1}
$1 + 2 + \cdots + n = \dfrac{n(n+1)}{2}$.
\end{prob}

\begin{prob}[M]{2}
$1^2 + 2^2 + \cdots + n^2 = \dfrac{n(n+1)(2n+1)}{6}$.
\end{prob}

\begin{prob}[M]{3}
$1^3 + 2^3 + \cdots + n^3 = (1 + 2 + \cdots + n)^2$.
\end{prob}

\begin{prob}[M]{4}
$1 + 2 + 2^2 + \cdots + 2^{n-1} = 2^n - 1$.
\end{prob}

\begin{prob}[M]{5}
设 $a^{[n]} = a(a-h)\cdots[a-(n-1)h]$ 及 $a^{[0]} = 1$，求证
\[
(a+b)^{[n]} = \sum_{m=0}^{n} \mathrm{C}_n^m a^{[n-m]} b^{[m]},
\]
其中 $\mathrm{C}_n^m$ 是由 $n$ 个元素中选取 $m$ 个的组合数，由此推出牛顿二项式公式.
\end{prob}

\begin{prob}[M]{6}
证明伯努利不等式：
\[
(1+x_1)(1+x_2)\cdots(1+x_n) \geqslant 1 + x_1 + x_2 + \cdots + x_n,
\]
式中 $x_1, x_2, \cdots, x_n$ 是符号相同且大于 $-1$ 的数.
\end{prob}

\begin{prob}[M]{7}
证明：若 $x > -1$，则不等式
\[
(1+x)^n \geqslant 1 + nx \quad (n > 1)
\]
为真，且仅当 $x = 0$ 时，等式成立.
\end{prob}

\begin{prob}[M]{8}
证明不等式：
\[
n! < \left(\frac{n+1}{2}\right)^n \quad (n > 1).
\]
\probhint{利用不等式
\[
\left(\frac{n+2}{n+1}\right)^{n+1} = \left(1 + \frac{1}{n+1}\right)^{n+1} > 2 \quad (n = 1, 2, \cdots).
\]}
\end{prob}

\begin{prob}[M]{9}
证明如下不等式：
\begin{enumerate}[label=(\alph*), itemsep=2pt, topsep=3pt]
  \item $2! \cdot 4! \cdots (2n)! > [(n+1)!]^n \ (n > 1)$;
  \item $\dfrac{1}{2} \cdot \dfrac{3}{4} \cdots \dfrac{2n-1}{2n} < \dfrac{1}{\sqrt{2n+1}}$.
\end{enumerate}
\end{prob}

\begin{prob}[L]{10}
证明如下不等式：
\begin{enumerate}[label=(\alph*), itemsep=2pt, topsep=3pt]
  \item $1 + \dfrac{1}{\sqrt{2}} + \dfrac{1}{\sqrt{3}} + \cdots + \dfrac{1}{\sqrt{n}} > \sqrt{n} \ (n \geqslant 2)$;
  \item $n^{n+1} > (n+1)^n \ (n \geqslant 3)$;
  \item $\left|\sin\left(\displaystyle\sum_{k=1}^{n} x_k\right)\right| \leqslant \displaystyle\sum_{k=1}^{n} \sin x_k \ (0 \leqslant x_k \leqslant \pi, \, k = 1, 2, \cdots, n)$;
  \item $(2n)! < 2^{2n}(n!)^2$.
\end{enumerate}
\end{prob}

\begin{prob}[M]{11}
设 $c$ 为正整数，而不为整数的平方，且 $A/B$ 为确定实数 $\sqrt{c}$ 的分割，其中 $B$ 类包含所有满足 $b^2 > c$ 的正有理数 $b$，而 $A$ 类包含所有其余的有理数. 求证：在 $A$ 类中无最大数，而在 $B$ 类中无最小数.
\end{prob}

\begin{prob}[M]{12}
确定数 $\sqrt[3]{2}$ 的分割 $A/B$ 用下面的方法来建立：$A$ 类包含所有满足 $a^3 < 2$ 的有理数 $a$；$B$ 类包含所有其余的有理数. 证明：在 $A$ 类中无最大数，而在 $B$ 类中无最小数.
\end{prob}

\begin{prob}[M]{13}
作出适当的分割，然后证明等式：

(a) $\sqrt{2} + \sqrt{8} = \sqrt{18}$; \qquad (b) $\sqrt{2}\sqrt{3} = \sqrt{6}$.
\end{prob}

\begin{prob}[M]{14}
建立确定数 $2^{\sqrt{2}}$ 的分割.
\end{prob}

\begin{prob}[M]{15}
求证：任何非空且下方有界的数集有下确界，而且任何非空且上方有界的数集有上确界.
\end{prob}

\begin{prob}[M]{16}
证明：一切有理真分数 $\dfrac{m}{n}$（式中 $m$ 及 $n$ 为正整数，且 $0 < m < n$）的集合无最小及最大的元素. 求这个集合的上确界及下确界.
\end{prob}

\begin{prob}[M]{17}
求满足不等式
\[
r^2 < 2
\]
的有理数 $r$ 所成集合的下确界和上确界.
\end{prob}

\begin{prob}[M]{18}
设 $\{-x\}$ 为数 $x \in \{x\}$ 的相反数的集合，证明：

(a) $\inf\{-x\} = -\sup\{x\}$; \qquad (b) $\sup\{-x\} = -\inf\{x\}$.
\end{prob}

\begin{prob}[M]{19}
设 $\{x+y\}$ 为所有 $x+y$ 这些和的集合，其中 $x \in \{x\}$ 及 $y \in \{y\}$，证明等式：

(a) $\inf\{x+y\} = \inf\{x\} + \inf\{y\}$; \qquad (b) $\sup\{x+y\} = \sup\{x\} + \sup\{y\}$.
\end{prob}

\begin{prob}[M]{20}
设 $\{xy\}$ 为所有 $xy$ 乘积的集合，其中 $x \in \{x\}$ 及 $y \in \{y\}$，且 $x \geqslant 0$ 及 $y \geqslant 0$. 证明等式：

(a) $\inf\{xy\} = \inf\{x\}\inf\{y\}$; \qquad (b) $\sup\{xy\} = \sup\{x\}\sup\{y\}$.
\end{prob}

\begin{prob}[M]{21}
求证不等式：

(a) $|x - y| \geqslant ||x| - |y||$; \qquad (b) $|x + x_1 + \cdots + x_n| \geqslant |x| - (|x_1| + \cdots + |x_n|)$.
\end{prob}

\probgroup{解不等式：}

\begin{prob}[S]{22}
$|x + 1| < 0.01$.
\end{prob}

\begin{prob}[S]{23}
$|x - 2| \geqslant 10$.
\end{prob}

\begin{prob}[S]{24}
$|x| > |x + 1|$.
\end{prob}

\begin{prob}[S]{25}
$|2x - 1| < |x - 1|$.
\end{prob}

\begin{prob}[S]{26}
$|x + 2| + |x - 2| \leqslant 12$.
\end{prob}

\begin{prob}[S]{27}
$|x + 2| - |x| > 1$.
\end{prob}

\begin{prob}[S]{28}
$||x + 1| - |x - 1|| < 1$.
\end{prob}

\begin{prob}[S]{29}
$|x(1 - x)| < 0.05$.
\end{prob}

\begin{prob}[M]{30}
证明恒等式：
\[
\left(\frac{x + |x|}{2}\right)^2 + \left(\frac{x - |x|}{2}\right)^2 = x^2.
\]
\end{prob}

\begin{prob}[S]{31}
当测量长度 $10\,\mathrm{cm}$ 时，绝对误差为 $0.5\,\mathrm{mm}$；当测量距离 $500\,\mathrm{km}$ 时，绝对误差等于 $200\,\mathrm{m}$. 哪种测量较为精确？
\end{prob}

\begin{prob}[S]{32}
设数 $x = 2.3752$ 的相对误差为 $1\%$，试确定此数包含多少位精确数字？
\end{prob}

\begin{prob}[S]{33}
数 $x = 12.125$ 包含 $3$ 位精确数字，试求此数的相对误差.
\end{prob}

\begin{prob}[M]{34}
矩形的边长等于：
\[
x = 2.50\,\mathrm{cm} \pm 0.01\,\mathrm{cm}, \quad y = 4.00\,\mathrm{cm} \pm 0.02\,\mathrm{cm}.
\]
这个矩形的面积 $S$ 界于什么范围内？当其边长取平均值时，矩形面积的绝对误差 $\Delta$ 和相对误差 $\delta$ 是多少？
\end{prob}

\begin{prob}[M]{35}
物体的质量 $P = 12.59\,\mathrm{g} \pm 0.01\,\mathrm{g}$，其体积 $v = 3.2\,\mathrm{cm}^3 \pm 0.2\,\mathrm{cm}^3$. 若对物体的质量和体积都取其平均值，试求物体的密度，并估计密度的绝对误差和相对误差.
\end{prob}

\begin{prob}[M]{36}
圆半径
\[
r = 7.2\,\mathrm{m} \pm 0.1\,\mathrm{m}.
\]
若取 $\pi = 3.14$，则求出的圆面积的最小相对误差是多少？
\end{prob}

\begin{prob}[M]{37}
已测得长方体各边长为
\[
x = 24.7\,\mathrm{m} \pm 0.2\,\mathrm{m}, \quad y = 6.5\,\mathrm{m} \pm 0.1\,\mathrm{m}, \quad z = 1.2\,\mathrm{m} \pm 0.1\,\mathrm{m}.
\]
此长方体的体积 $V$ 界于什么范围内？若测量的各结果都取其平均值，则所求出的体积可能有的绝对误差和相对误差是多少？
\end{prob}

\begin{prob}[M]{38}
正方形的边长 $x$ 满足 $2\,\mathrm{m} < x < 3\,\mathrm{m}$. 应当以多小的绝对误差来测量边长，在计算此正方形的面积时才有可能精确到 $0.001\,\mathrm{m}^2$？
\end{prob}

\begin{prob}[M]{39}
假设矩形每边的长皆不超过 $10\,\mathrm{m}$，为了使根据测量所计算出来的面积与原面积之差不超过 $0.01\,\mathrm{cm}^2$，问测量矩形的边 $x$ 与 $y$ 时，许可的绝对误差 $\Delta$ 的值多大？
\end{prob}

\begin{prob}[M]{40}
设 $\delta(x)$ 及 $\delta(y)$ 为数 $x$ 和 $y$ 的相对误差，$\delta(xy)$ 为数 $xy$ 的相对误差. 求证
\[
\delta(xy) \leqslant \delta(x) + \delta(y) + \delta(x)\delta(y).
\]
为数列 $x_n$ 的（有限或无穷）极限存在的充分必要条件.
\end{prob}

\probsec{§2. 数列理论}

\begin{prob}[M]{41}
设
\[
x_n = \frac{n}{n+1} \quad (n=1,2,\cdots).
\]
证明：
\[
\lim_{n\to\infty} x_n = 1,
\]
即对于任一个给定的 $\varepsilon > 0$，求出数 $N = N(\varepsilon)$，使得
\[
\text{在 } n > N \text{ 时，} \quad |x_n - 1| < \varepsilon.
\]
填下表：
\begin{center}
\begin{tabular}{c|c|c|c|c|c}
\hline
$\quad \varepsilon \quad$ & $\quad 0.1 \quad$ & $\quad 0.01 \quad$ & $\quad 0.001 \quad$ & $\quad 0.0001 \quad$ & $\quad \cdots \quad$ \\
\hline
$N$ & & & & & \\
\hline
\end{tabular}
\end{center}
\end{prob}

\begin{prob}[L]{42}
设
\begin{enumerate}[label=(\alph*), itemsep=2pt, topsep=3pt]
\item $x_n = \dfrac{(-1)^{n+1}}{n}$；
\item $x_n = \dfrac{2n}{n^3+1}$；
\item $x_n = \dfrac{1}{n!}$；
\item $x_n = (-1)^n \cdot 0.999^n$，
\end{enumerate}
对于任何的 $\varepsilon > 0$，求出数 $N = N(\varepsilon)$，使得
\[
\text{当 } n > N \text{ 时，} \quad |x_n| < \varepsilon;
\]
从而证明：$x_n\ (n=1,2,\cdots)$ 为无穷小量（就是说，它的极限值为 $0$）.

对应着上面四种情形，填下表：
\begin{center}
\begin{tabular}{c|c|c|c|c}
\hline
$\quad \varepsilon \quad$ & $\quad 0.1 \quad$ & $\quad 0.01 \quad$ & $\quad 0.001 \quad$ & $\quad \cdots \quad$ \\
\hline
$N$ & & & & \\
\hline
\end{tabular}
\end{center}
\end{prob}

\begin{prob}[L]{43}
证明：数列
\begin{enumerate}[label=(\alph*), itemsep=2pt, topsep=3pt]
\item $x_n = (-1)^n n$，
\item $x_n = 2^{\sqrt{n}}$，
\item $x_n = \lg(\lg n)\ (n \geqslant 2)$
\end{enumerate}
当 $n \to \infty$ 时有无穷极限（即成为无穷大），即：对任意的 $E > 0$，求出数 $N = N(E)$，使得
\[
\text{当 } n > N \text{ 时，} \quad |x_n| > E.
\]
对应着上面的每一种情形，填下表：
\begin{center}
\begin{tabular}{c|c|c|c|c|c}
\hline
$\quad E \quad$ & $\quad 10 \quad$ & $\quad 100 \quad$ & $\quad 1000 \quad$ & $\quad 10000 \quad$ & $\quad \cdots \quad$ \\
\hline
$N$ & & & & & \\
\hline
\end{tabular}
\end{center}
\end{prob}

\begin{prob}[M]{44}
求证：
\[
x_n = n^{(-1)^n} \quad (n=1,2,\cdots)
\]
无界，但当 $n \to \infty$ 时，它并不成为无穷大.
\end{prob}

\begin{prob}[L]{45}
用不等式表述下列命题：
\begin{enumerate}[label=(\alph*), itemsep=2pt, topsep=3pt]
\item $\lim\limits_{n\to\infty} x_n = \infty$；
\item $\lim\limits_{n\to\infty} x_n = -\infty$；
\item $\lim\limits_{n\to\infty} x_n = +\infty$.
\end{enumerate}
\end{prob}

\probgroup{设 $n$ 遍历正整数列，求下列各式之值：}

\begin{prob}[S]{46}
$\lim\limits_{n\to\infty} \dfrac{10\,000n}{n^2+1}$.
\end{prob}

\begin{prob}[S]{47}
$\lim\limits_{n\to\infty} (\sqrt{n+1} - \sqrt{n})$.
\end{prob}

\begin{prob}[S]{48}
$\lim\limits_{n\to\infty} \dfrac{\sqrt[3]{n^2}\sin n!}{n+1}$.
\end{prob}

\begin{prob}[S]{49}
$\lim\limits_{n\to\infty} \dfrac{(-2)^n + 3^n}{(-2)^{n+1} + 3^{n+1}}$.
\end{prob}

\begin{prob}[S]{50}
$\lim\limits_{n\to\infty} \dfrac{1 + a + a^2 + \cdots + a^n}{1 + b + b^2 + \cdots + b^n}\ (|a|<1,\ |b|<1)$.
\end{prob}

\begin{prob}[S]{51}
$\lim\limits_{n\to\infty} \left(\dfrac{1}{n^2} + \dfrac{2}{n^2} + \cdots + \dfrac{n-1}{n^2}\right)$.
\end{prob}

\begin{prob}[S]{52}
$\lim\limits_{n\to\infty} \left|\dfrac{1}{n} - \dfrac{2}{n} + \dfrac{3}{n} - \cdots + \dfrac{(-1)^{n-1}n}{n}\right|$.
\end{prob}

\begin{prob}[S]{53}
$\lim\limits_{n\to\infty} \left[\dfrac{1^2}{n^3} + \dfrac{2^2}{n^3} + \cdots + \dfrac{(n-1)^2}{n^3}\right]$.
\end{prob}

\begin{prob}[S]{54}
$\lim\limits_{n\to\infty} \left[\dfrac{1^2}{n^3} + \dfrac{3^2}{n^3} + \cdots + \dfrac{(2n-1)^2}{n^3}\right]$.
\end{prob}

\begin{prob}[S]{55}
$\lim\limits_{n\to\infty} \left(\dfrac{1}{2} + \dfrac{3}{2^2} + \dfrac{5}{2^3} + \cdots + \dfrac{2n-1}{2^n}\right)$.
\end{prob}

\begin{prob}[S]{56}
$\lim\limits_{n\to\infty} \left[\dfrac{1}{1\cdot 2} + \dfrac{1}{2\cdot 3} + \cdots + \dfrac{1}{n(n+1)}\right]$.
\end{prob}

\begin{prob}[S]{57}
$\lim\limits_{n\to\infty} \left(\sqrt{2}\sqrt[4]{2}\sqrt[8]{2}\cdots\sqrt[2^n]{2}\right)$.
\end{prob}

\probgroup{证明下列等式：}

\begin{prob}[M]{58}
$\lim\limits_{n\to\infty} \dfrac{n}{2^n} = 0$.
\end{prob}

\begin{prob}[M]{59}
$\lim\limits_{n\to\infty} \dfrac{2^n}{n!} = 0$.
\end{prob}

\begin{prob}[M]{60}
$\lim\limits_{n\to\infty} \dfrac{n^k}{a^n} = 0\ (a > 1)$.
\end{prob}

\begin{prob}[M]{61}
$\lim\limits_{n\to\infty} \dfrac{a^n}{n!} = 0$.
\end{prob}

\begin{prob}[M]{62}
$\lim\limits_{n\to\infty} nq^n = 0$，若 $|q| < 1$.
\end{prob}

\begin{prob}[M]{63}
$\lim\limits_{n\to\infty} \sqrt[n]{a} = 1\ (a > 0)$.
\end{prob}

\begin{prob}[M]{64}
$\lim\limits_{n\to\infty} \dfrac{\log_a n}{n} = 0\ (a > 1)$.
\end{prob}

\begin{prob}[M]{65}
$\lim\limits_{n\to\infty} \sqrt[n]{n} = 1$.
\end{prob}

\begin{prob}[M]{66}
$\lim\limits_{n\to\infty} \dfrac{1}{\sqrt[n]{n!}} = 0$.
\end{prob}

\begin{prob}[L]{67}
当 $n$ 充分大时，下列各组表达式中哪个更大？
\begin{enumerate}[label=(\alph*), itemsep=2pt, topsep=3pt]
\item $100n + 200,\ 0.01n^2$；
\item $2^n,\ n^{1000}$；
\item $1000^n,\ n!$.
\end{enumerate}
\end{prob}

\begin{prob}[M]{68}
证明：
\[
\lim_{n\to\infty} \left(\frac{1}{2}\cdot\frac{3}{4}\cdots\frac{2n-1}{2n}\right) = 0.
\]
\probhint{参阅习题 9 (b).}
\end{prob}

\begin{prob}[M]{69}
证明：数列
\[
x_n = \left(1 + \frac{1}{n}\right)^n \quad (n=1,2,\cdots)
\]
是单调增的，且上方有界，而数列
\[
y_n = \left(1 + \frac{1}{n}\right)^{n+1} \quad (n=1,2,\cdots)
\]
是单调减的，且上方有界. 由此推出这些数列有公共的极限：
\[
\lim_{n\to\infty} \left(1 + \frac{1}{n}\right)^n = \lim_{n\to\infty} \left(1 + \frac{1}{n}\right)^{n+1} = \mathrm{e}.
\]
\probhint{先作出比 $\dfrac{x_{n+1}}{x_n},\ \dfrac{y_n}{y_{n-1}}$，并利用习题 7 的不等式.}
\end{prob}

\begin{prob}[M]{70}
证明：
\[
0 < \mathrm{e} - \left(1 + \frac{1}{n}\right)^n < \frac{3}{n} \quad (n=1,2,\cdots).
\]
当指数 $n$ 是什么样的数值时，表达式 $\left(1 + \dfrac{1}{n}\right)^n$ 与数 $\mathrm{e}$ 之差小于 $0.001$？
\end{prob}

\begin{prob}[M]{71}
设 $p_n\ (n=1,2,\cdots)$ 为趋于 $+\infty$ 的任意数列，而 $q_n\ (n=1,2,\cdots)$ 为趋于 $-\infty$ 的任意数列 ($p_n, q_n \notin [-1,0]$). 求证：
\[
\lim_{n\to\infty} \left(1 + \frac{1}{p_n}\right)^{p_n} = \lim_{n\to\infty} \left(1 + \frac{1}{q_n}\right)^{q_n} = \mathrm{e}.
\]
\end{prob}

\begin{prob}[L]{72}
已知
\[
\lim_{n\to\infty} \left(1 + \frac{1}{n}\right)^n = \mathrm{e},
\]
求证：
\[
\lim_{n\to\infty} \left(1 + 1 + \frac{1}{2!} + \frac{1}{3!} + \cdots + \frac{1}{n!}\right) = \mathrm{e}.
\]
由此推出公式
\[
\mathrm{e} = 1 + 1 + \frac{1}{2!} + \frac{1}{3!} + \cdots + \frac{1}{n!} + \frac{\theta_n}{n!n}, \tag{$*$}
\]
其中 $0 < \theta_n < 1$，并计算数 $\mathrm{e}$，精确到 $10^{-5}$.
\end{prob}

\begin{prob}[M]{73}
证明：数 $\mathrm{e}$ 为无理数.
\end{prob}

\begin{prob}[M]{74}
证明不等式：
\[
\left(\frac{n}{\mathrm{e}}\right)^n < n! < \mathrm{e}\left(\frac{n}{2}\right)^n.
\]
\end{prob}

\begin{prob}[M]{75}
证明不等式：
\begin{enumerate}[label=(\alph*), itemsep=2pt, topsep=3pt]
\item $\dfrac{1}{n+1} < \ln\left(1 + \dfrac{1}{n}\right) < \dfrac{1}{n}$，其中 $n$ 为任意的正整数.
\item $1 + a < \mathrm{e}^a$，其中 $a$ 为异于零的实数.
\end{enumerate}
\end{prob}

\begin{prob}[M]{76}
求证：
\[
\lim_{n\to\infty} n\left(a^{\frac{1}{n}} - 1\right) = \ln a \quad (a > 0),
\]
式中 $\ln a$ 是取 $\mathrm{e} = 2.718\cdots$ 作底时，数 $a$ 的对数.
\end{prob}

\probgroup{利用关于单调有界数列极限存在的定理，证明以下各数列的收敛性：}

\begin{prob}[M]{77}
$x_n = p_0 + \dfrac{p_1}{10} + \cdots + \dfrac{p_n}{10^n}\ (n=1,2,\cdots)$，其中 $p_i\ (i=0,1,2,\cdots)$ 是非负整数，并且从 $p_1$ 起不大于 $9$.
\end{prob}

\begin{prob}[M]{78}
$x_n = \dfrac{10}{1}\cdot\dfrac{11}{3}\cdots\dfrac{n+9}{2n-1}$.
\end{prob}

\begin{prob}[M]{79}
$x_n = \left(1 - \dfrac{1}{2}\right)\left(1 - \dfrac{1}{4}\right)\cdots\left(1 - \dfrac{1}{2^n}\right)$.
\end{prob}

\begin{prob}[M]{80}
$x_n = \left(1 + \dfrac{1}{2}\right)\left(1 + \dfrac{1}{4}\right)\cdots\left(1 + \dfrac{1}{2^n}\right)$.
\end{prob}

\begin{prob}[M]{81}
$x_1 = \sqrt{2},\ x_2 = \sqrt{2 + \sqrt{2}},\ \cdots,\ x_n = \underbrace{\sqrt{2 + \sqrt{2 + \cdots + \sqrt{2}}}}_{n\text{ 重根号}},\ \cdots$.
\end{prob}

\probgroup{利用柯西准则，证明以下各数列的收敛性：}

\begin{prob}[M]{82}
$x_n = a_0 + a_1 q + \cdots + a_n q^n$，其中 $|a_k| < M\ (k=0,1,2,\cdots)$ 且 $|q| < 1$.
\end{prob}

\begin{prob}[M]{83}
$x_n = \dfrac{\sin 1}{2} + \dfrac{\sin 2}{2^2} + \cdots + \dfrac{\sin n}{2^n}$.
\end{prob}

\begin{prob}[M]{84}
$x_n = \dfrac{\cos 1!}{1\cdot 2} + \dfrac{\cos 2!}{2\cdot 3} + \cdots + \dfrac{\cos n!}{n(n+1)}$.
\end{prob}

\begin{prob}[M]{85}
$x_n = 1 + \dfrac{1}{2^2} + \dfrac{1}{3^2} + \cdots + \dfrac{1}{n^2}$.
\probhint{利用不等式 $\dfrac{1}{n^2} < \dfrac{1}{n-1} - \dfrac{1}{n}\ (n=2,3,\cdots)$.}
\end{prob}

\begin{prob}[L]{86}
对于数列 $x_n\ (n=1,2,\cdots)$，若存在数 $C$，使得
\[
|x_2 - x_1| + |x_3 - x_2| + \cdots + |x_n - x_{n-1}| < C \quad (n=2,3,\cdots),
\]
则称数列 $x_n$ 有有界变差.

证明：凡有有界变差的数列是收敛的.

举出一个收敛数列而无有界变差的例子.
\end{prob}

\begin{prob}[S]{87}
试述“某数列不满足柯西准则”的含义.
\end{prob}

\begin{prob}[M]{88}
利用柯西准则，证明数列：
\[
x_n = 1 + \frac{1}{2} + \frac{1}{3} + \cdots + \frac{1}{n}
\]
的发散性.
\end{prob}

\begin{prob}[M]{89}
证明：若数列 $x_n\ (n=1,2,\cdots)$ 收敛，则它的任何子数列 $x_{p_n}$ 也收敛且有同一极限：$\lim\limits_{n\to\infty} x_{p_n} = \lim\limits_{n\to\infty} x_n$.
\end{prob}

\begin{prob}[M]{90}
证明：若单调数列的某一子数列收敛，则此单调数列本身是收敛的.
\end{prob}

\begin{prob}[M]{91}
证明：若 $\lim\limits_{n\to\infty} x_n = a$，则 $\lim\limits_{n\to\infty} |x_n| = |a|$.
\end{prob}

\begin{prob}[S]{92}
设 $x_n \to a$，则极限
\[
\lim_{n\to\infty} \frac{x_{n+1}}{x_n}
\]
是什么？
\end{prob}

\begin{prob}[M]{93}
证明：收敛的数列是有界的.
\end{prob}

\begin{prob}[M]{94}
证明：收敛的数列或达到其上确界，或达到其下确界，或两者都达到. 举出这三类数列的例子.
\end{prob}

\begin{prob}[M]{95}
证明：趋于 $+\infty$ 的数列 $x_n\ (n=1,2,\cdots)$ 必定达到其下确界.
\end{prob}

\probgroup{求数列 $x_n\ (n=1,2,\cdots)$ 的最大项，设：}

\begin{prob}[S]{96}
$x_n = \dfrac{n^2}{2^n}$.
\end{prob}

\begin{prob}[S]{97}
$x_n = \dfrac{\sqrt{n}}{100+n}$.
\end{prob}

\begin{prob}[S]{98}
$x_n = \dfrac{1000^n}{n!}$.
\end{prob}

\probgroup{求数列 $x_n\ (n=1,2,\cdots)$ 的最小项，设：}

\begin{prob}[S]{99}
$x_n = n^2 - 9n - 100$.
\end{prob}

\begin{prob}[S]{100}
$x_n = n + \dfrac{100}{n}$.
\end{prob}

\probgroup{求数列 $x_n\ (n=1,2,\cdots)$ 的 $\inf x_n,\ \sup x_n,\ \varliminf\limits_{n\to\infty} x_n$ 及 $\varlimsup\limits_{n\to\infty} x_n$，设：}

\begin{prob}[M]{101}
(a) $x_n = 1 - \dfrac{1}{n}$； \qquad (b) $x_n = (-1)^{n-1}\left(2 + \dfrac{3}{n}\right)$.
\end{prob}

\begin{prob}[S]{102}
$x_n = \dfrac{(-1)^n}{n} + \dfrac{1+(-1)^n}{2}$.
\end{prob}

\begin{prob}[S]{103}
$x_n = 1 + \dfrac{n}{n+1}\cos\dfrac{n\pi}{2}$.
\end{prob}

\begin{prob}[S]{104}
$x_n = 1 + 2\cdot(-1)^{n+1} + 3\cdot(-1)^{\frac{n(n-1)}{2}}$.
\end{prob}

\begin{prob}[S]{105}
$x_n = \dfrac{n-1}{n+1}\cos\dfrac{2n\pi}{3}$.
\end{prob}

\begin{prob}[S]{106}
$x_n = (-1)^n n$.
\end{prob}

\begin{prob}[S]{107}
$x_n = -n[2 + (-1)^n]$.
\end{prob}

\begin{prob}[S]{108}
$x_n = n^{(-1)^n}$.
\end{prob}

\begin{prob}[S]{109}
$x_n = 1 + n\sin\dfrac{n\pi}{2}$.
\end{prob}

\begin{prob}[S]{110}
$x_n = \dfrac{1}{n - 10.2}$.
\end{prob}

\probgroup{求 $\varliminf\limits_{n\to\infty} x_n$ 及 $\varlimsup\limits_{n\to\infty} x_n$，设：}

\begin{prob}[S]{111}
$x_n = \dfrac{n^2}{1+n^2}\cos\dfrac{2n\pi}{3}$.
\end{prob}

\begin{prob}[S]{112}
$x_n = \left(1 + \dfrac{1}{n}\right)^n \cdot (-1)^n + \sin\dfrac{n\pi}{4}$.
\end{prob}

\begin{prob}[S]{113}
$x_n = \dfrac{n}{n+1}\sin^2\dfrac{n\pi}{4}$.
\end{prob}

\begin{prob}[S]{114}
$x_n = \sqrt[n]{1 + 2^{n(-1)^n}}$.
\end{prob}

\begin{prob}[S]{115}
$x_n = \cos^n\dfrac{2n\pi}{3}$.
\end{prob}

\probgroup{求以下各数列的子列极限：}

\begin{prob}[S]{116}
$\dfrac{1}{2},\ \dfrac{1}{2},\ \dfrac{1}{4},\ \dfrac{3}{4},\ \dfrac{1}{8},\ \dfrac{7}{8},\ \cdots,\ \dfrac{1}{2^n},\ \dfrac{2^n-1}{2^n},\ \cdots$.
\end{prob}

\begin{prob}[S]{117}
$1,\ \dfrac{1}{2},\ 1+\dfrac{1}{2},\ \dfrac{1}{3},\ 1+\dfrac{1}{3},\ \dfrac{1}{2}+\dfrac{1}{3},\ \dfrac{1}{4},\ 1+\dfrac{1}{4},\ \dfrac{1}{2}+\dfrac{1}{4},\ \dfrac{1}{3}+\dfrac{1}{4},\ \dfrac{1}{5},\ \cdots,\ \dfrac{1}{n},\ 1+\dfrac{1}{n},\ \dfrac{1}{2}+\dfrac{1}{n},\ \cdots,\ \dfrac{1}{n-1}+\dfrac{1}{n},\ \dfrac{1}{n+1},\ \cdots$.
\end{prob}

\begin{prob}[S]{118}
$\dfrac{1}{2},\ \dfrac{1}{3},\ \dfrac{2}{3},\ \dfrac{1}{4},\ \dfrac{2}{4},\ \dfrac{3}{4},\ \dfrac{1}{5},\ \dfrac{2}{5},\ \dfrac{3}{5},\ \dfrac{4}{5},\ \cdots$.
\end{prob}

\begin{prob}[S]{119}
$x_n = 3\left(1 - \dfrac{1}{n}\right) + 2(-1)^n$.
\end{prob}

\begin{prob}[S]{120}
$x_n = \dfrac{1}{2}[(a+b) + (-1)^n(a-b)]$.
\end{prob}

\begin{prob}[M]{121}
试举出以已知数
\[
a_1, a_2, \cdots, a_p
\]
作为子列极限的数列的例子.
\end{prob}

\begin{prob}[M]{122}
试举出数列的例子，对此数列而言，已知数列
\[
a_1, a_2, \cdots, a_n, \cdots
\]
的所有各项皆为其子列极限. 所举数列还必有怎样的子列极限?
\end{prob}

\begin{prob}[L]{123}
举出数列的例子:
\begin{enumerate}[label=(\alph*), itemsep=2pt, topsep=3pt]
  \item 没有有限的子列极限;
  \item 有唯一有限的子列极限, 但不收敛;
  \item 有无穷多的子列极限;
  \item 以每一实数作为子列极限.
\end{enumerate}
\end{prob}

\begin{prob}[M]{124}
证明: 数列 $x_n$ 和 $y_n = x_n \sqrt[n]{n}$ ($n = 1, 2, \cdots$) 有相同的子列极限.
\end{prob}

\begin{prob}[M]{125}
证明: 从有界的数列 $x_n$ ($n = 1, 2, \cdots$) 中, 永远可选出收敛的子数列 $x_{p_n}$ ($n = 1, 2, \cdots$).
\end{prob}

\begin{prob}[M]{126}
证明: 若数列 $x_n$ ($n = 1, 2, \cdots$) 无界, 则存在子数列 $x_{p_n}$, 使得
\[
\lim_{n\to\infty} x_{p_n} = \infty.
\]
\end{prob}

\begin{prob}[M]{127}
设数列 $x_n$ ($n = 1, 2, \cdots$) 收敛, 而数列 $y_n$ ($n = 1, 2, \cdots$) 发散, 则能否断定关于数列 (a) $x_n + y_n$, (b) $x_n y_n$ 的收敛性? 举出适当的例子.
\end{prob}

\begin{prob}[M]{128}
设数列 $x_n$ 和 $y_n$ 发散 ($n = 1, 2, \cdots$). 可否断定数列 (a) $x_n + y_n$, (b) $x_n y_n$ 也发散呢?
\end{prob}

\begin{prob}[M]{129}
设 $\lim\limits_{n\to\infty} x_n = 0$, $y_n$ ($n = 1, 2, \cdots$) 为任意数列. 能否断定
\[
\lim_{n\to\infty} x_n y_n = 0?
\]
举出适当的例子.
\end{prob}

\begin{prob}[M]{130}
设 $\lim\limits_{n\to\infty} x_n y_n = 0$, 是否由此可得出: 或 $\lim\limits_{n\to\infty} x_n = 0$, 或 $\lim\limits_{n\to\infty} y_n = 0$? 考虑例子:
\[
x_n = \frac{1 + (-1)^n}{2}, \quad y_n = \frac{1 - (-1)^n}{2} \quad (n = 1, 2, \cdots).
\]
\end{prob}

\begin{prob}[L]{131}
证明\footnote{在习题 131---133 中, 假设极限之和或积皆有意义. —— 译注}:
\begin{enumerate}[label=(\alph*), itemsep=2pt, topsep=3pt]
  \item $\varliminf\limits_{n\to\infty} x_n + \varliminf\limits_{n\to\infty} y_n \leqslant \varliminf\limits_{n\to\infty} (x_n + y_n) \leqslant \varliminf\limits_{n\to\infty} x_n + \varlimsup\limits_{n\to\infty} y_n$;
  \item $\varliminf\limits_{n\to\infty} x_n + \varlimsup\limits_{n\to\infty} y_n \leqslant \varlimsup\limits_{n\to\infty} (x_n + y_n) \leqslant \varlimsup\limits_{n\to\infty} x_n + \varlimsup\limits_{n\to\infty} y_n$.
\end{enumerate}
举出在上面关系式中严格不等式成立的例子.
\end{prob}

\begin{prob}[L]{132}
设 $x_n \geqslant 0$ 且 $y_n \geqslant 0$ ($n = 1, 2, \cdots$), 证明:
\begin{enumerate}[label=(\alph*), itemsep=2pt, topsep=3pt]
  \item $\varliminf\limits_{n\to\infty} x_n \cdot \varliminf\limits_{n\to\infty} y_n \leqslant \varliminf\limits_{n\to\infty} (x_n y_n) \leqslant \varliminf\limits_{n\to\infty} x_n \cdot \varlimsup\limits_{n\to\infty} y_n$;
  \item $\varliminf\limits_{n\to\infty} x_n \cdot \varlimsup\limits_{n\to\infty} y_n \leqslant \varlimsup\limits_{n\to\infty} (x_n y_n) \leqslant \varlimsup\limits_{n\to\infty} x_n \cdot \varlimsup\limits_{n\to\infty} y_n$.
\end{enumerate}
举出在这些关系式中严格不等式成立的例子.
\end{prob}

\begin{prob}[M]{133}
证明: 若 $\lim\limits_{n\to\infty} x_n$ 存在, 则对于任何数列 $y_n$ ($n = 1, 2, \cdots$) 有:
\begin{enumerate}[label=(\alph*), itemsep=2pt, topsep=3pt]
  \item $\varlimsup\limits_{n\to\infty} (x_n + y_n) = \lim\limits_{n\to\infty} x_n + \varlimsup\limits_{n\to\infty} y_n$;
  \item $\varlimsup\limits_{n\to\infty} (x_n y_n) = \lim\limits_{n\to\infty} x_n \cdot \varlimsup\limits_{n\to\infty} y_n \quad (x_n \geqslant 0)$.
\end{enumerate}
\end{prob}

\begin{prob}[M]{134}
证明: 若对于某数列 $x_n$ ($n = 1, 2, \cdots$), 无论数列 $y_n$ ($n = 1, 2, \cdots$) 如何选取, 以下两个等式中都至少有一个成立:
\begin{enumerate}[label=(\alph*), itemsep=2pt, topsep=3pt]
  \item $\varlimsup\limits_{n\to\infty} (x_n + y_n) = \varlimsup\limits_{n\to\infty} x_n + \varlimsup\limits_{n\to\infty} y_n$;
  \item $\varlimsup\limits_{n\to\infty} (x_n y_n) = \varlimsup\limits_{n\to\infty} x_n \cdot \varlimsup\limits_{n\to\infty} y_n \quad (x_n \geqslant 0)$,
\end{enumerate}
则数列 $x_n$ 收敛或发散于 $+\infty$\footnote{此处的原文是: 则序列 $x_n$ 收敛. —— 译注}.
\end{prob}

\begin{prob}[M]{135}
证明: 若 $x_n > 0$ ($n = 1, 2, \cdots$) 及
\[
\varlimsup_{n\to\infty} x_n \cdot \varlimsup_{n\to\infty} \frac{1}{x_n} = 1,
\]
则数列 $x_n$ 是收敛的.
\end{prob}

\begin{prob}[M]{136}
证明: 若数列 $x_n$ ($n = 1, 2, \cdots$) 有界, 且
\[
\lim_{n\to\infty} (x_{n+1} - x_n) = 0,
\]
则此数列的子列极限充满于下极限和上极限
\[
l = \varliminf_{n\to\infty} x_n \quad \text{和} \quad L = \varlimsup_{n\to\infty} x_n
\]
之间, 即区间 $[l, L]$ 中的任意一个数都是该数列的子列极限.
\end{prob}

\begin{prob}[M]{137}
设数列 $x_1, x_2, \cdots, x_n, \cdots$ 满足条件
\[
0 \leqslant x_{m+n} \leqslant x_m + x_n \quad (m, n = 1, 2, \cdots),
\]
证明: $\lim\limits_{n\to\infty} \dfrac{x_n}{n}$ 存在.
\end{prob}

\begin{prob}[M]{138}
证明: 若数列 $x_n$ 收敛, 则算术平均值数列
\[
\xi_n = \frac{1}{n}(x_1 + x_2 + \cdots + x_n) \quad (n = 1, 2, \cdots)
\]
也收敛, 且
\[
\lim_{n\to\infty} \frac{x_1 + x_2 + \cdots + x_n}{n} = \lim_{n\to\infty} x_n.
\]
逆命题不成立, 举例说明.
\end{prob}

\begin{prob}[M]{139}
证明: 若 $\lim\limits_{n\to\infty} x_n = +\infty$, 则
\[
\lim_{n\to\infty} \frac{x_1 + x_2 + \cdots + x_n}{n} = +\infty.
\]
\end{prob}

\begin{prob}[M]{140}
证明: 若数列 $x_n$ ($n = 1, 2, \cdots$) 收敛且 $x_n > 0$, 则
\[
\lim_{n\to\infty} \sqrt[n]{x_1 x_2 \cdots x_n} = \lim_{n\to\infty} x_n.
\]
\end{prob}

\begin{prob}[M]{141}
证明: 若 $x_n > 0$ ($n = 1, 2, \cdots$) 且 $\lim\limits_{n\to\infty} \dfrac{x_{n+1}}{x_n}$ 存在, 则
\[
\lim_{n\to\infty} \sqrt[n]{x_n} = \lim_{n\to\infty} \frac{x_{n+1}}{x_n}.
\]
\end{prob}

\begin{prob}[M]{142}
证明:
\[
\lim_{n\to\infty} \frac{n}{\sqrt[n]{n!}} = \mathrm{e}.
\]
\end{prob}

\begin{prob}[L]{143}
证明施托尔茨定理: 若
\begin{enumerate}[label=(\alph*), itemsep=2pt, topsep=3pt]
  \item $y_{n+1} > y_n$ ($n = 1, 2, \cdots$);
  \item $\lim\limits_{n\to\infty} y_n = +\infty$;
  \item $\lim\limits_{n\to\infty} \dfrac{x_{n+1} - x_n}{y_{n+1} - y_n}$ 存在,
\end{enumerate}
则
\[
\lim_{n\to\infty} \frac{x_n}{y_n} = \lim_{n\to\infty} \frac{x_{n+1} - x_n}{y_{n+1} - y_n}.
\]
\end{prob}

\begin{prob}[M]{144}
求:
\begin{enumerate}[label=(\alph*), itemsep=2pt, topsep=3pt]
  \item $\lim\limits_{n\to+\infty} \dfrac{n^2}{a^n}$ ($a > 1$);
  \item $\lim\limits_{n\to+\infty} \dfrac{\lg n}{n}$.
\end{enumerate}
\end{prob}

\begin{prob}[L]{145}
证明: 若 $p$ 为正整数, 则
\begin{enumerate}[label=(\alph*), itemsep=2pt, topsep=3pt]
  \item $\lim\limits_{n\to\infty} \dfrac{1^p + 2^p + \cdots + n^p}{n^{p+1}} = \dfrac{1}{p+1}$,
  \item $\lim\limits_{n\to\infty} \left(\dfrac{1^p + 2^p + \cdots + n^p}{n^p} - \dfrac{n}{p+1}\right) = \dfrac{1}{2}$,
  \item $\lim\limits_{n\to\infty} \dfrac{1^p + 3^p + \cdots + (2n-1)^p}{n^{p+1}} = \dfrac{2^p}{p+1}$.
\end{enumerate}
\end{prob}

\begin{prob}[M]{146}
证明: 数列
\[
x_n = 1 + \frac{1}{2} + \frac{1}{3} + \cdots + \frac{1}{n} - \ln n \quad (n = 1, 2, \cdots)
\]
收敛.

因此有公式
\[
1 + \frac{1}{2} + \frac{1}{3} + \cdots + \frac{1}{n} = C + \ln n + \varepsilon_n,
\]
式中 $C = 0.577216\cdots$ 称为欧拉常数, 且当 $n \to \infty$ 时 $\varepsilon_n \to 0$.
\end{prob}

\begin{prob}[M]{147}
求
\[
\lim_{n\to\infty} \left(\frac{1}{n+1} + \frac{1}{n+2} + \cdots + \frac{1}{2n}\right).
\]
\end{prob}

\begin{prob}[M]{148}
数列 $x_n$ ($n = 1, 2, \cdots$) 由下列各式
\[
x_1 = a, \quad x_2 = b, \quad x_n = \frac{x_{n-1} + x_{n-2}}{2} \quad (n = 3, 4, \cdots)
\]
所确定, 求 $\lim\limits_{n\to\infty} x_n$.
\end{prob}

\begin{prob}[M]{149}
设 $x_n$ ($n = 1, 2, \cdots$) 为由以下各式
\[
x_0 > 0, \quad x_{n+1} = \frac{1}{2}\left(x_n + \frac{1}{x_n}\right) \quad (n = 0, 1, 2, \cdots)
\]
所确定的数列. 求证
\[
\lim_{n\to\infty} x_n = 1.
\]
\end{prob}

\begin{prob}[M]{150}
证明: 由下列各式
\[
x_1 = a, \quad y_1 = b, \quad x_{n+1} = \sqrt{x_n y_n}, \quad y_{n+1} = \frac{x_n + y_n}{2},
\]
确定的数列 $x_n$ 和 $y_n$ ($n = 1, 2, \cdots$) 有共同的极限
\[
\mu(a, b) = \lim_{n\to\infty} x_n = \lim_{n\to\infty} y_n
\]
(数 $a$ 和 $b$ 的算术几何平均值).
\end{prob}

\probsec{§3. 函数的概念}

\probgroup{求下列函数的存在域:}

\begin{prob}[S]{151}
$y = \dfrac{x^2}{1 + x}$.
\end{prob}

\begin{prob}[S]{152}
$y = \sqrt{3x - x^3}$.
\end{prob}

\begin{prob}[S]{153}
$y = (x - 2)\sqrt{\dfrac{1 + x}{1 - x}}$.
\end{prob}

\begin{prob}[M]{154}
(a) $y = \log(x^2 - 4)$; \qquad (b) $y = \log(x + 2) + \log(x - 2)$.
\end{prob}

\begin{prob}[S]{155}
$y = \sqrt{\sin(\sqrt{x})}$.
\end{prob}

\begin{prob}[S]{156}
$y = \sqrt{\cos x^2}$.
\end{prob}

\begin{prob}[S]{157}
$y = \lg\left(\sin\dfrac{\pi}{x}\right)$.
\end{prob}

\begin{prob}[S]{158}
$y = \dfrac{\sqrt{x}}{\sin \pi x}$.
\end{prob}

\begin{prob}[S]{159}
$y = \arcsin\dfrac{2x}{1 + x}$.
\end{prob}

\begin{prob}[S]{160}
$y = \arccos(2\sin x)$.
\end{prob}

\begin{prob}[S]{161}
$y = \lg[\cos(\lg x)]$.
\end{prob}

\begin{prob}[S]{162}
$y = (x + |x|)\sqrt{x\sin^2 \pi x}$.
\end{prob}

\begin{prob}[S]{163}
$y = \cot \pi x + \arccos(2^x)$.
\end{prob}

\begin{prob}[S]{164}
$y = \arcsin(1 - x) + \lg(\lg x)$.
\end{prob}

\begin{prob}[L]{165}
\begin{enumerate}[label=(\alph*), itemsep=2pt, topsep=3pt]
  \item $y = (2x)!$;
  \item $y = \log_2 \log_3 \log_4 x$;
  \item $y = \sqrt[4]{\lg\tan x}$;
  \item $y = \sqrt{\sin 2x} + \sqrt{\sin 3x}$ ($0 \leqslant x \leqslant 2\pi$).
\end{enumerate}
\end{prob}

\probgroup{求下列函数的存在域和值域:}

\begin{prob}[S]{166}
$y = \sqrt{2 + x - x^2}$.
\end{prob}

\begin{prob}[S]{167}
$y = \lg(1 - 2\cos x)$.
\end{prob}

\begin{prob}[S]{168}
$y = \arccos\dfrac{2x}{1 + x^2}$.
\end{prob}

\begin{prob}[S]{169}
$y = \arcsin\left(\lg\dfrac{x}{10}\right)$.
\end{prob}

\begin{prob}[S]{170}
$y = (-1)^x$.
\end{prob}

\begin{prob}[L]{171}
在底 $AC = b$ 和高 $BD = h$ 的三角形 $ABC$ 中 (图 1) 内接一个高 $NM = x$ 的矩形 $KLMN$. 把矩形 $KLMN$ 的周长 $P$ 及其面积 $S$ 表示为 $x$ 的函数.

作函数 $P = P(x)$ 及 $S = S(x)$ 的图像.
\probfig{fig_171.png}{图 1}
\end{prob}

\begin{prob}[L]{172}
在三角形 $ABC$ 中, 边 $AB = 6\,\mathrm{cm}$, $AC = 8\,\mathrm{cm}$, 角 $BAC = x$. 把边 $BC = a$ 和面积 $ABC = S$ 表示为变量 $x$ 的函数. 作函数 $a = a(x)$ 及 $S = S(x)$ 的图像.
\end{prob}

\begin{prob}[L]{173}
在等腰梯形 $ABCD$ 中 (图 2), 底 $AD = a$, $BC = b$ ($a > b$), 高 $HB = h$, 引直线 $MN \parallel BH$, $MN$ 与顶点 $A$ 相距 $AM = x$. 把图形 $ABNMA$ 的面积 $S$ 表示为变量 $x$ 的函数. 作函数 $S = S(x)$ 的图像.
\probfig{fig_173.png}{图 2}
\end{prob}

\begin{prob}[L]{174}
有 $2\,\mathrm{g}$ 质量均匀分布在 $Ox$ 轴上的闭区间 $0 \leqslant x \leqslant 1$ 上, 另有两个质量为 $1\,\mathrm{g}$ 的质点分别位于点 $x = 2$ 和 $x = 3$.

设 $m(x)$ 是区间 $(-\infty, x)$ 内的质量的值, 求函数 $m = m(x)$ ($-\infty < x < +\infty$) 的解析表达式, 并作这个函数的图像.
\end{prob}

\begin{prob}[M]{175}
函数 $y = \operatorname{sgn} x$ 被定义为:
\[
\operatorname{sgn} x = \begin{cases}
-1, & x < 0, \\
0, & x = 0, \\
1, & x > 0.
\end{cases}
\]
作这个函数的图像. 证明:
\[
|x| = x \operatorname{sgn} x.
\]
\end{prob}

\begin{prob}[M]{176}
函数 $y = [x]$ (数 $x$ 的整数部分) 用下法定义: 若 $x = n + r$, 式中 $n$ 为整数且 $0 \leqslant r < 1$, 则
\[
[x] = n.
\]
作这个函数的图像.
\end{prob}

\begin{prob}[M]{177}
设
\[
y = \pi(x) \quad (x \geqslant 0)
\]
表示不超过数 $x$ 的素数的数目, 作这个函数在 $0 \leqslant x \leqslant 20$ 时的图像.
\end{prob}

\probgroup{下列函数 $y = f(x)$ 把集合 $E_x$ 映射成怎样的集合 $E_y$?}

\begin{prob}[S]{178}
$y = x^2$, $E_x = \{-1 \leqslant x \leqslant 2\}$.
\end{prob}

\begin{prob}[S]{179}
$y = \lg x$, $E_x = \{10 < x < 1000\}$.
\end{prob}

\begin{prob}[S]{180}
$y = \dfrac{1}{\pi}\operatorname{arccot} x,\ E_x = \{-\infty < x < +\infty\}$.
\end{prob}

\begin{prob}[S]{181}
$y = \cot\dfrac{\pi x}{4},\ E_x = \{0 < |x| \leqslant 1\}$.
\end{prob}

\begin{prob}[S]{182}
$y = |x|,\ E_x = \{1 \leqslant |x| \leqslant 2\}$.
\end{prob}

\probgroup{设变量 $x$ 遍历区间 $0 < x < 1$，试确定变量 $y$ 所遍历的集合：}

\begin{prob}[S]{183}
$y = a + (b - a)x$.
\end{prob}

\begin{prob}[S]{184}
$y = \dfrac{1}{1 - x}$.
\end{prob}

\begin{prob}[S]{185}
$y = \dfrac{x}{2x - 1}$.
\end{prob}

\begin{prob}[S]{186}
$y = \sqrt{x - x^2}$.
\end{prob}

\begin{prob}[S]{187}
$y = \cot\pi x$.
\end{prob}

\begin{prob}[S]{188}
$y = x + [2x]$.
\end{prob}

\begin{prob}[S]{189}
设 $f(x) = x^4 - 6x^3 + 11x^2 - 6x$，求 $f(0),\, f(1),\, f(2),\, f(3),\, f(4)$。
\end{prob}

\begin{prob}[S]{190}
设 $f(x) = \lg x^2$，求 $f(-1),\, f(-0.001),\, f(100)$。
\end{prob}

\begin{prob}[S]{191}
设 $f(x) = 1 + [x]$，求 $f(0.9),\, f(0.99),\, f(0.999),\, f(1)$。
\end{prob}

\begin{prob}[M]{192}
设
\[
f(x) = \begin{cases}
1 + x, & -\infty < x \leqslant 0, \\
2^x, & 0 < x < +\infty,
\end{cases}
\]
求 $f(-2),\, f(-1),\, f(0),\, f(1),\, f(2)$。
\end{prob}

\begin{prob}[M]{193}
设
\[
f(x) = \frac{1 - x}{1 + x},
\]
求 $f(0),\, f(-x),\, f(x + 1),\, f(x) + 1,\, f\!\left(\dfrac{1}{x}\right),\, \dfrac{1}{f(x)}$。
\end{prob}

\begin{prob}[L]{194}
设
\begin{enumerate}[label=(\alph*), itemsep=2pt, topsep=3pt]
\item $f(x) = x - x^3$;
\item $f(x) = \sin\dfrac{\pi}{x}$;
\item $f(x) = (x + |x|)(1 - x)$.
\end{enumerate}
求满足以下各式的 $x$ 值：(1) $f(x) = 0$；(2) $f(x) > 0$；(3) $f(x) < 0$。
\end{prob}

\begin{prob}[L]{195}
设
\begin{enumerate}[label=(\alph*), itemsep=2pt, topsep=3pt]
\item $f(x) = ax + b$;
\item $f(x) = x^2$;
\item $f(x) = a^x$,
\end{enumerate}
求 $\varphi(x) = \dfrac{f(x + h) - f(x)}{h}$。
\end{prob}

\begin{prob}[M]{196}
设 $f(x) = ax^2 + bx + c$，证明：
\[
f(x + 3) - 3f(x + 2) + 3f(x + 1) - f(x) \equiv 0.
\]
\end{prob}

\begin{prob}[M]{197}
若 $f(0) = -2,\, f(3) = 5$，求线性函数
\[
f(x) = ax + b.
\]
$f(1)$ 及 $f(2)$ 等于什么（线性插值法）？
\end{prob}

\begin{prob}[M]{198}
若 $f(-2) = 0,\, f(0) = 1,\, f(1) = 5$，求二次有理函数：
\[
f(x) = ax^2 + bx + c.
\]
$f(-1)$ 及 $f(0.5)$ 等于什么（二次插值法）？
\end{prob}

\begin{prob}[M]{199}
设 $f(-1) = 0,\, f(0) = 2,\, f(1) = -3,\, f(2) = 5$。求三次有理函数：
\[
f(x) = ax^3 + bx^2 + cx + d.
\]
\end{prob}

\begin{prob}[M]{200}
设 $f(0) = 15,\, f(2) = 30,\, f(4) = 90$。求形为
\[
f(x) = a + bc^x
\]
的函数。
\end{prob}

\begin{prob}[M]{201}
证明：对于线性函数
\[
f(x) = ax + b,
\]
若自变量的值 $x = x_n\ (n = 1, 2, \cdots)$ 组成等差数列，则对应的函数值 $y_n = f(x_n)\ (n = 1, 2, \cdots)$ 也组成等差数列。
\end{prob}

\begin{prob}[M]{202}
证明：对于指数函数
\[
f(x) = a^x \quad (a > 0),
\]
若自变量的值 $x = x_n\ (n = 1, 2, \cdots)$ 组成等差数列，则对应的函数值 $y_n = f(x_n)\ (n = 1, 2, \cdots)$ 组成等比数列。
\end{prob}

\begin{prob}[L]{203}
设当 $0 < u < 1$ 时函数 $f(u)$ 有定义。求下列函数的定义域：
\begin{enumerate}[label=(\alph*), itemsep=2pt, topsep=3pt]
\item $f(\sin x)$;
\item $f(\ln x)$;
\item $f\!\left(\dfrac{[x]}{x}\right)$.
\end{enumerate}
\end{prob}

\begin{prob}[M]{204}
设
\[
f(x) = \frac{1}{2}(a^x + a^{-x}) \quad (a > 0).
\]
证明：
\[
f(x + y) + f(x - y) = 2f(x)f(y).
\]
\end{prob}

\begin{prob}[L]{205}
设 $f(x) + f(y) = f(z)$。求出 $z$，若
\begin{enumerate}[label=(\alph*), itemsep=2pt, topsep=3pt]
\item $f(x) = ax$;
\item $f(x) = \dfrac{1}{x}$;
\item $f(x) = \arctan x\ (|x| < 1)$;
\item $f(x) = \log\dfrac{1 + x}{1 - x}$.
\end{enumerate}
\end{prob}

\probgroup{求 $\varphi[\varphi(x)],\, \psi[\psi(x)],\, \varphi[\psi(x)]$ 及 $\psi[\varphi(x)]$，设}

\begin{prob}[M]{206}
$\varphi(x) = x^2$ 及 $\psi(x) = 2^x$。
\end{prob}

\begin{prob}[M]{207}
$\varphi(x) = \operatorname{sgn} x$ 及 $\psi(x) = \dfrac{1}{x}$。
\end{prob}

\begin{prob}[M]{208}
\[
\varphi(x) = \begin{cases}
0, & x \leqslant 0, \\
x, & x > 0.
\end{cases}
\quad \text{及} \quad
\psi(x) = \begin{cases}
0, & x \leqslant 0, \\
-x^2, & x > 0.
\end{cases}
\]
\end{prob}

\begin{prob}[M]{209}
设
\[
f(x) = \frac{1}{1 - x},
\]
求 $f[f(x)],\, f\{f[f(x)]\}$。
\end{prob}

\begin{prob}[M]{210}
设
\[
f_n(x) = \underbrace{f(f(\cdots f(x)))}_{n\text{ 次}},
\]
若
\[
f(x) = \frac{x}{\sqrt{1 + x^2}},
\]
求 $f_n(x)$。
\end{prob}

\begin{prob}[S]{211}
设 $f(x + 1) = x^2 - 3x + 2$，求 $f(x)$。
\end{prob}

\begin{prob}[M]{212}
设
\[
f\!\left(x + \frac{1}{x}\right) = x^2 + \frac{1}{x^2} \quad (|x| \geqslant 2),
\]
求 $f(x)$。
\end{prob}

\begin{prob}[M]{213.1}
设
\[
f\!\left(\frac{1}{x}\right) = x + \sqrt{1 + x^2} \quad (x > 0),
\]
求 $f(x)$。
\end{prob}

\begin{prob}[M]{213.2}
设
\[
f\!\left(\frac{x}{x + 1}\right) = x^2,
\]
求 $f(x)$。
\end{prob}

\probgroup{证明：下列各函数在所给区间内是单调增函数：}

\begin{prob}[S]{214}
$f(x) = x^2\ (0 \leqslant x < +\infty)$.
\end{prob}

\begin{prob}[S]{215}
$f(x) = \sin x\ \left(-\dfrac{\pi}{2} \leqslant x \leqslant \dfrac{\pi}{2}\right)$.
\end{prob}

\begin{prob}[S]{216}
$f(x) = \tan x\ \left(-\dfrac{\pi}{2} < x < \dfrac{\pi}{2}\right)$.
\end{prob}

\begin{prob}[M]{217}
$f(x) = 2x + \sin x\ (-\infty < x < +\infty)$.
\end{prob}

\probgroup{证明：下列各函数在所给区间内是单调减函数：}

\begin{prob}[S]{218}
$f(x) = x^2\ (-\infty < x \leqslant 0)$.
\end{prob}

\begin{prob}[S]{219}
$f(x) = \cos x\ (0 \leqslant x \leqslant \pi)$.
\end{prob}

\begin{prob}[S]{220}
$f(x) = \cot x\ (0 < x < \pi)$.
\end{prob}

\begin{prob}[L]{221}
研究下列函数的单调性：
\begin{enumerate}[label=(\alph*), itemsep=2pt, topsep=3pt]
\item $f(x) = ax + b$;
\item $f(x) = ax^2 + bx + c$;
\item $f(x) = x^3$;
\item $f(x) = \dfrac{ax + b}{cx + d}$;
\item $f(x) = a^x\ (a > 0)$.
\end{enumerate}
\end{prob}

\begin{prob}[S]{222}
不等式能否逐项取对数？
\end{prob}

\begin{prob}[M]{223}
设 $\varphi(x),\, \psi(x)$ 及 $f(x)$ 为单调增函数。证明：若
\[
\varphi(x) \leqslant f(x) \leqslant \psi(x),
\]
则
\[
\varphi[\varphi(x)] \leqslant f[f(x)] \leqslant \psi[\psi(x)].
\]
\end{prob}

\probgroup{求反函数 $x = \varphi(y)$ 和它的存在域，若：}

\begin{prob}[S]{224}
$y = 2x + 3\ (-\infty < x < +\infty)$.
\end{prob}

\begin{prob}[M]{225}
$y = x^2$；(a) $(-\infty < x \leqslant 0)$；\qquad (b) $(0 \leqslant x < +\infty)$。
\end{prob}

\begin{prob}[S]{226}
$y = \dfrac{1 - x}{1 + x}\ (x \neq -1)$.
\end{prob}

\begin{prob}[M]{227}
$y = \sqrt{1 - x^2}$；(a) $(-1 \leqslant x \leqslant 0)$；\qquad (b) $(0 \leqslant x \leqslant 1)$。
\end{prob}

\begin{prob}[M]{228}
$y = \sinh x$，式中 $\sinh x = \dfrac{1}{2}(\mathrm{e}^x - \mathrm{e}^{-x})\ (-\infty < x < +\infty)$。
\end{prob}

\begin{prob}[M]{229}
$y = \tanh x$，其中 $\tanh x = \dfrac{\mathrm{e}^x - \mathrm{e}^{-x}}{\mathrm{e}^x + \mathrm{e}^{-x}}\ (-\infty < x < +\infty)$。
\end{prob}

\begin{prob}[M]{230}
\[
y = \begin{cases}
x, & -\infty < x < 1, \\
x^2, & 1 \leqslant x \leqslant 4, \\
2^x, & 4 < x < +\infty.
\end{cases}
\]
\end{prob}

\begin{prob}[L]{231}
设函数 $f(x)$ 定义于对称区间 $(-l, l)$，若
\[
f(-x) \equiv f(x),
\]
则称 $f(x)$ 为偶函数；若
\[
f(-x) \equiv -f(x),
\]
则称 $f(x)$ 为奇函数。

确定下列函数中哪些是偶函数，哪些是奇函数：
\begin{enumerate}[label=(\alph*), itemsep=2pt, topsep=3pt]
\item $f(x) = 3x - x^3$;
\item $f(x) = \sqrt[3]{(1 - x)^2} + \sqrt[3]{(1 + x)^2}$;
\item $f(x) = a^x + a^{-x}\ (a > 0)$;
\item $f(x) = \ln\dfrac{1 - x}{1 + x}$;
\item $f(x) = \ln\!\left(x + \sqrt{1 + x^2}\right)$.
\end{enumerate}
\end{prob}

\begin{prob}[M]{232}
证明：定义于对称区间 $(-l, l)$ 的任何函数 $f(x)$ 可以表示为偶函数与奇函数之和的形式。
\end{prob}

\begin{prob}[L]{233}
若存在数 $T > 0$（函数的周期——在广义的意义上！），使定义于集合 $E$ 的函数 $f(x)$ 满足等式
\[
f(x \pm T) = f(x) \quad (x \in E),
\]
则函数 $f(x)$ 称为周期函数。

说明下列函数中哪些是周期函数，并求它们的最小周期：
\begin{enumerate}[label=(\alph*), itemsep=2pt, topsep=3pt]
\item $f(x) = A\cos\lambda x + B\sin\lambda x$;
\item $f(x) = \sin x + \dfrac{1}{2}\sin 2x + \dfrac{1}{3}\sin 3x$;
\item $f(x) = 2\tan\dfrac{x}{2} - 3\tan\dfrac{x}{3}$;
\item $f(x) = \sin^2 x$;
\item $f(x) = \sin x^2$;
\item $f(x) = \sqrt{\tan x}$;
\item $f(x) = \tan\sqrt{x}$;
\item $f(x) = \sin x + \sin(x\sqrt{2})$.
\end{enumerate}
\end{prob}

\begin{prob}[M]{234}
证明：对于狄利克雷函数
\[
\chi(x) = \begin{cases}
1, & x\text{ 为有理数}, \\
0, & x\text{ 为无理数},
\end{cases}
\]
任何有理数皆为其周期。
\end{prob}

\begin{prob}[M]{235.1}
证明：定义于共同的集合且周期可公度的两个周期函数之和及其乘积也是周期函数。
\end{prob}

\begin{prob}[M]{235.2}
若
\[
f(x + T) \equiv -f(x) \quad (T > 0),
\]
则函数 $f(x)$ 称为反周期函数。证明：此函数是以 $2T$ 为周期的函数。
\end{prob}

\begin{prob}[M]{236}
证明：若函数 $f(x)\ (-\infty < x < +\infty)$ 满足等式 $f(x + T) = kf(x)$，式中 $k$ 和 $T$ 为正的常数，则 $f(x) = a^x\varphi(x)$，式中 $a$ 为常数，而 $\varphi(x)$ 为以 $T$ 为周期的函数。
\end{prob}

\probsec{§4. 函数的图像表示法}

\begin{prob}[M]{237}
作出线性齐次函数 $y = ax$ 当 $a = 0,\, \dfrac{1}{2},\, 1,\, 2,\, -1$ 时的图像。
\end{prob}

\begin{prob}[M]{238}
作出线性函数 $y = x + b$ 当 $b = 0,\, 1,\, 2,\, -1$ 时的图像。
\end{prob}

\begin{prob}[L]{239}
作出线性函数的图像：
\begin{enumerate}[label=(\alph*), itemsep=2pt, topsep=3pt]
\item $y = 2x + 3$;
\item $y = 2 - 0.1x$;
\item $y = -\dfrac{x}{2} - 1$.
\end{enumerate}
\end{prob}

\begin{prob}[M]{240}
铁的线膨胀系数 $\alpha = 1.2 \times 10^{-5}\,\mathrm{K}^{-1}$。在适当的尺度下作出函数
\[
l = f(T) \quad (233\,\mathrm{K} \leqslant T \leqslant 373\,\mathrm{K})
\]
的图像，其中 $T$ 表示温度，$l$ 表示当温度为 $T$ 时细长铁棒的长，已知当 $T = 273\,\mathrm{K}$ 时，$l = 100\,\mathrm{cm}$。
\end{prob}

\begin{prob}[M]{241}
二质点沿数轴运动，第一个质点在初始时刻 $t = 0$ 时位于坐标原点左方 $20\,\mathrm{m}$ 处，其速度为 $v_1 = 10\,\mathrm{m/s}$；第二个质点在 $t = 0$ 时位于原点 $O$ 右方 $30\,\mathrm{m}$ 处，其速度为 $v_2 = -20\,\mathrm{m/s}$。作出此二点运动方程的图像并求它们相遇的时刻和位置。
\end{prob}

\begin{prob}[L]{242}
作出二次有理函数的图像（抛物线）：
\begin{enumerate}[label=(\alph*), itemsep=2pt, topsep=3pt]
\item $y = ax^2$，当 $a = 1,\, \dfrac{1}{2},\, 2,\, -1$;
\item $y = (x - x_0)^2$，当 $x_0 = 0,\, 1,\, 2,\, -1$;
\end{enumerate}
(c) $y = x^2 + c$，当 $c = 0, 1, 2, -1$.
\end{prob}

\begin{prob}[L]{243}
把二次三项式 $y = ax^2 + bx + c$ 化为 $y = y_0 + a(x - x_0)^2$ 的形式，再作出它的图像。研究例子：
\begin{enumerate}[label=(\alph*), itemsep=2pt, topsep=3pt]
  \item $y = 8x - 2x^2$;
  \item $y = x^2 - 3x + 2$;
  \item $y = -x^2 + 2x - 1$;
  \item $y = \dfrac{1}{2}x^2 + x + 1$.
\end{enumerate}
\end{prob}

\begin{prob}[M]{244}
质点以初速度 $v_0 = 600\,\mathrm{m/s}$ 沿与水平面成角 $\alpha = 45^\circ$ 的方向射出。作出运动轨迹的图像并求最大上升高度及射程（取 $g \approx 10\,\mathrm{m/s^2}$，不计空气阻力）。
\end{prob}

\probgroup{作出下列高于二次的有理函数的图像：}

\begin{prob}[S]{245}
$y = x^3 + 1$.
\end{prob}

\begin{prob}[S]{246}
$y = (1 - x^2)(2 + x)$.
\end{prob}

\begin{prob}[S]{247}
$y = x^2 - x^4$.
\end{prob}

\begin{prob}[S]{248}
$y = x(a - x)^2(a + x)^3 \ (a > 0)$.
\end{prob}

\probgroup{作出下列分式线性函数的图像（双曲线）：}

\begin{prob}[S]{249}
$y = \dfrac{1}{x}$.
\end{prob}

\begin{prob}[S]{250}
$y = \dfrac{1 - x}{1 + x}$.
\end{prob}

\begin{prob}[M]{251}
把分式线性函数
\[
y = \frac{ax + b}{cx + d} \quad (ad - bc \neq 0,\ c \neq 0)
\]
化为以下形式：
\[
y = y_0 + \frac{m}{x - x_0},
\]
再作出它的图像。

研究例子
\[
y = \frac{3x + 2}{2x - 3}.
\]
\end{prob}

\begin{prob}[M]{252}
设气体当压强 $p_0 = 1\,\mathrm{Pa}$ 时占有体积 $V_0 = 12\,\mathrm{m^3}$。若气体的温度保持不变，作出气体体积 $V$ 对压强 $p$ 的依赖关系的图像（玻意耳--马略特定律）。
\end{prob}

\probgroup{作下列分式有理函数的图像：}

\begin{prob}[S]{253}
$y = x + \dfrac{1}{x}$.
\end{prob}

\begin{prob}[S]{254}
$y = x^2 + \dfrac{1}{x}$（牛顿三叉线）。
\end{prob}

\begin{prob}[S]{255}
$y = x + \dfrac{1}{x^2}$.
\end{prob}

\begin{prob}[S]{256}
$y = \dfrac{1}{1 + x^2}$（箕舌线）。
\end{prob}

\begin{prob}[S]{257}
$y = \dfrac{2x}{1 + x^2}$（牛顿蛇形线）。
\end{prob}

\begin{prob}[S]{258}
$y = \dfrac{1}{1 - x^2}$.
\end{prob}

\begin{prob}[S]{259}
$y = \dfrac{x}{1 - x^2}$.
\end{prob}

\begin{prob}[S]{260}
$y = \dfrac{1}{1 + x} - \dfrac{2}{x} + \dfrac{1}{1 - x}$.
\end{prob}

\begin{prob}[S]{261}
$y = \dfrac{1}{1 + x} - \dfrac{2}{x^2} + \dfrac{1}{1 - x}$.
\end{prob}

\begin{prob}[S]{262}
$y = \dfrac{(x + 1)(x - 2)}{(x - 1)(x + 2)}$.
\end{prob}

\begin{prob}[M]{263}
把函数
\[
y = \frac{ax^2 + bx + c}{a_1 x + b_1} \quad (a_1 \neq 0)
\]
化为以下形式：
\[
y = kx + m + \frac{n}{x - x_0},
\]
然后作出它的草图。

研究例子
\[
y = \frac{x^2 - 4x + 3}{x + 1}.
\]
\end{prob}

\begin{prob}[M]{264}
一质点与引力中心相距 $x$，质点所受引力的大小为 $F$，并且当 $x = 1\,\mathrm{m}$ 时 $F = 10\,\mathrm{N}$。作出引力 $F$ 的图像（牛顿定律）。
\end{prob}

\begin{prob}[M]{265}
根据范德瓦耳斯定律，当温度不变时，真实气体的体积 $V$ 与压强 $p$ 以关系式
\[
\left(p + \frac{a}{V^2}\right)(V - b) = c
\]
相联系。若 $a = 2, b = 0.1$ 及 $c = 10$，作出函数 $p = p(V)$ 的图像。
\end{prob}

\probgroup{作下列无理函数的图像：}

\begin{prob}[S]{266}
$y = \pm\sqrt{-x - 2}$（抛物线）。
\end{prob}

\begin{prob}[S]{267}
$y = \pm x\sqrt{x}$（半三次抛物线）。
\end{prob}

\begin{prob}[S]{268}
$y = \pm\dfrac{1}{2}\sqrt{100 - x^2}$（椭圆）。
\end{prob}

\begin{prob}[S]{269}
$y = \pm\sqrt{x^2 - 1}$（双曲线）。
\end{prob}

\begin{prob}[S]{270}
$y = \pm\sqrt{\dfrac{1 - x}{1 + x}}$.
\end{prob}

\begin{prob}[S]{271}
$y = \pm x\sqrt{100 - x^2}$.
\end{prob}

\begin{prob}[S]{272}
$y = \pm x\sqrt{\dfrac{x}{10 - x}}$（蔓叶线）。
\end{prob}

\begin{prob}[S]{273}
$y = \pm\sqrt{(x^2 - 1)(9 - x^2)}$.
\end{prob}

\begin{prob}[M]{274}
作幂函数 $y = x^n$ 当 (a) $n = 1, 3, 5$; (b) $n = 2, 4, 6$ 时的图像。
\end{prob}

\begin{prob}[M]{275}
作幂函数 $y = x^n$ 当 (a) $n = -1, -3$; (b) $n = -2, -4$ 时的图像。
\end{prob}

\begin{prob}[M]{276}
作根式 $y = \sqrt[m]{x}$ 当 (a) $m = 2, 4$; (b) $m = 3, 5$ 时的图像。
\end{prob}

\begin{prob}[L]{277}
设：
\begin{enumerate}[label=(\alph*), itemsep=2pt, topsep=3pt]
  \item $m = 2,\ k = 1$;
  \item $m = 2,\ k = 3$;
  \item $m = 3,\ k = 1$;
  \item $m = 3,\ k = 2$;
  \item $m = 3,\ k = 4$;
  \item $m = 4,\ k = 2$;
  \item $m = 4,\ k = 3$.
\end{enumerate}
作根式 $y = \sqrt[m]{x^k}$ 的图像。
\end{prob}

\begin{prob}[M]{278}
作指数函数 $y = a^x$ 当 $a = \dfrac{1}{2}, 1, 2, \mathrm{e}, 10$ 时的图像。
\end{prob}

\begin{prob}[L]{279}
作复合指数函数 $y = \mathrm{e}^{y_1}$ 的图像，设：
\begin{enumerate}[label=(\alph*), itemsep=2pt, topsep=3pt]
  \item $y_1 = x^2$;
  \item $y_1 = -x^2$;
  \item $y_1 = \dfrac{1}{x}$;
  \item $y_1 = \dfrac{1}{x^2}$;
  \item $y_1 = -\dfrac{1}{x^2}$;
  \item $y_1 = \dfrac{2x}{1 - x^2}$.
\end{enumerate}
\end{prob}

\begin{prob}[M]{280}
作对数函数 $y = \log_a x$ 当 $a = \dfrac{1}{2}, 2, \mathrm{e}, 10$ 时的图像。
\end{prob}

\begin{prob}[M]{281}
作下列函数的图像：

(a) $y = \ln(-x)$; \qquad (b) $y = -\ln(x)$.
\end{prob}

\begin{prob}[L]{282}
设
\begin{enumerate}[label=(\alph*), itemsep=2pt, topsep=3pt]
  \item $y_1 = 1 + x^2$;
  \item $y_1 = (x - 1)(x - 2)^2(x - 3)^3$;
  \item $y_1 = \dfrac{1 - x}{1 + x}$;
  \item $y_1 = \dfrac{1}{x^2}$;
  \item $y_1 = 1 + \mathrm{e}^x$.
\end{enumerate}
作出对数复合函数 $y = \ln y_1$ 的图像。
\end{prob}

\begin{prob}[S]{283}
作函数 $y = \log_x 2$ 的图像。
\end{prob}

\begin{prob}[S]{284}
作函数 $y = A\sin x$ 当 $A = 1, 10, -2$ 时的图像。
\end{prob}

\begin{prob}[M]{285}
作函数 $y = \sin(x - x_0)$ 在 $x_0 = 0, \dfrac{\pi}{4}, \dfrac{\pi}{2}, \dfrac{3\pi}{4}, \pi$ 时的图像。
\end{prob}

\begin{prob}[M]{286}
作函数 $y = \sin nx$ 的图像。设 $n = 1, 2, 3, \dfrac{1}{2}, \dfrac{1}{3}$。
\end{prob}

\begin{prob}[M]{287}
把函数 $y = a\cos x + b\sin x$ 化为以下形式 $y = A\sin(x - x_0)$，再作它的图像。研究例子：$y = 6\cos x + 8\sin x$。
\end{prob}

\probgroup{作下列三角函数的图像：}

\begin{prob}[S]{288}
$y = \cos x$.
\end{prob}

\begin{prob}[S]{289}
$y = \tan x$.
\end{prob}

\begin{prob}[S]{290}
$y = \cot x$.
\end{prob}

\begin{prob}[S]{291}
$y = \sec x$.
\end{prob}

\begin{prob}[S]{292}
$y = \csc x$.
\end{prob}

\begin{prob}[S]{293}
$y = \sin^2 x$.
\end{prob}

\begin{prob}[S]{294}
$y = \sin^3 x$.
\end{prob}

\begin{prob}[S]{295}
$y = \cot^2 x$.
\end{prob}

\begin{prob}[S]{296}
$y = \sin x \cdot \sin 3x$.
\end{prob}

\begin{prob}[S]{297}
$y = \pm\sqrt{\cos x}$.
\end{prob}

\probgroup{作下列函数的图像：}

\begin{prob}[S]{298}
$y = \sin x^2$.
\end{prob}

\begin{prob}[S]{299}
$y = \sin\dfrac{1}{x}$.
\end{prob}

\begin{prob}[M]{300}
(a) $y = x\cos\dfrac{\pi}{x}$; \qquad (b) $y = \sin x \cdot \sin\dfrac{1}{x}$.
\end{prob}

\begin{prob}[M]{301}
(a) $y = \tan\dfrac{\pi}{x}$; \qquad (b) $y = \sec\dfrac{1}{x}$.
\end{prob}

\begin{prob}[S]{302}
$y = x\left(2 + \sin\dfrac{1}{x}\right)$.
\end{prob}

\begin{prob}[S]{303}
$y = \pm\sqrt{1 - x^2}\sin\dfrac{\pi}{x}$.
\end{prob}

\begin{prob}[S]{304}
$y = \dfrac{\sin x}{x}$.
\end{prob}

\begin{prob}[S]{305}
$y = \mathrm{e}^x\cos x$.
\end{prob}

\begin{prob}[S]{306}
$y = \pm 2^{-x}\sqrt{\sin\pi x}$.
\end{prob}

\begin{prob}[S]{307}
$y = \dfrac{\cos x}{1 + x^2}$.
\end{prob}

\begin{prob}[S]{308}
$y = \ln(\cos x)$.
\end{prob}

\begin{prob}[S]{309}
$y = \cos(\ln x)$.
\end{prob}

\begin{prob}[S]{310}
$y = \mathrm{e}^{\frac{1}{\sin x}}$.
\end{prob}

\probgroup{作下列反三角函数的图像：}

\begin{prob}[S]{311}
$y = \arcsin x$.
\end{prob}

\begin{prob}[S]{312}
$y = \arccos x$.
\end{prob}

\begin{prob}[S]{313}
$y = \arctan x$.
\end{prob}

\begin{prob}[S]{314}
$y = \operatorname{arccot} x$.
\end{prob}

\begin{prob}[S]{315}
$y = \arcsin\dfrac{1}{x}$.
\end{prob}

\begin{prob}[S]{316}
$y = \arccos\dfrac{1}{x}$.
\end{prob}

\begin{prob}[S]{317}
$y = \arctan\dfrac{1}{x}$.
\end{prob}

\begin{prob}[S]{318}
$y = \arcsin(\sin x)$.
\end{prob}

\begin{prob}[S]{319}
$y = \arcsin(\cos x)$.
\end{prob}

\begin{prob}[S]{320}
$y = \arccos(\cos x)$.
\end{prob}

\begin{prob}[S]{321}
$y = \arctan(\tan x)$.
\end{prob}

\begin{prob}[S]{322}
$y = \arcsin(2\sin x)$.
\end{prob}

\begin{prob}[L]{323}
设：
\begin{enumerate}[label=(\alph*), itemsep=2pt, topsep=3pt]
  \item $y_1 = 1 - \dfrac{x}{2}$;
  \item $y_1 = \dfrac{2x}{1 + x^2}$;
  \item $y_1 = \dfrac{1 - x}{1 + x}$;
  \item $y_1 = \mathrm{e}^x$.
\end{enumerate}
作函数 $y = \arcsin y_1$ 的图像。
\end{prob}

\begin{prob}[L]{324.1}
设：
\begin{enumerate}[label=(\alph*), itemsep=2pt, topsep=3pt]
  \item $y_1 = x^2$;
  \item $y_1 = \dfrac{1}{x^2}$;
  \item $y_1 = \ln x$;
  \item $y_1 = \dfrac{1}{\sin x}$.
\end{enumerate}
作函数 $y = \arctan y_1$ 的图像。
\end{prob}

\begin{prob}[L]{324.2}
作下列函数的图像：
\begin{enumerate}[label=(\alph*), itemsep=2pt, topsep=3pt]
  \item $y = x^3 - 3x + 2$;
  \item $y = \dfrac{x^3}{(1 - x)(1 + x)^2}$;
  \item $y = \dfrac{x^2}{|x| - 1}$;
  \item $y = \sqrt{x(1 - x^2)}$;
  \item $y = 3\sin\left(\dfrac{x}{2} + \dfrac{\pi}{4}\right)$;
  \item $y = \cot\dfrac{\pi x}{1 + x^2}$;
  \item $y = \dfrac{1}{1 - 2^{\frac{x}{1 - x}}}$;
  \item $y = \lg(x^2 - 3x + 2)$;
  \item $y = \arcsin\left(\dfrac{3}{2} - \sin x\right)$;
  \item $y = \arctan\left(\dfrac{1}{x - 1} + \dfrac{1}{x - 2} + \dfrac{1}{x - 3}\right)$;
  \item $y = \log_{\cos x} \sin x$;
  \item $y = (\sin x)^{\cot x}$.
\end{enumerate}
\end{prob}

\begin{prob}[L]{325}
已知函数 $y = f(x)$ 的图像，作下列各函数的图像：
\begin{enumerate}[label=(\alph*), itemsep=2pt, topsep=3pt]
  \item $y = -f(x)$;
  \item $y = f(-x)$;
  \item $y = -f(-x)$;
  \item $y = f(x - x_0)$;
  \item $y = y_0 + f(x - x_0)$;
  \item $y = f(2x)$;
  \item $y = f(kx + b) \ (k \neq 0)$.
\end{enumerate}
\end{prob}

\begin{prob}[M]{326}
设
\[
f(x) = \begin{cases}
1 - |x|, & |x| \leqslant 1, \\
0, & |x| > 1.
\end{cases}
\]
作函数
\[
y = \frac{1}{2}(f(x - t) + f(x + t))
\]
当 $t = 0, t = 1$ 及 $t = 2$ 时的图像。
\end{prob}

\begin{prob}[L]{327}
作函数的图像：
\begin{enumerate}[label=(\alph*), itemsep=2pt, topsep=3pt]
  \item $y = 2 + \sqrt{1 - x}$;
  \item $y = 1 - \mathrm{e}^{-x}$;
  \item $y = \ln(1 + x)$;
  \item $y = -\arcsin(1 + x)$;
  \item $y = 3 + 2\cos 3x$.
\end{enumerate}
\end{prob}

\begin{prob}[L]{328}
已知函数 $y = f(x)$ 的图像，作下列函数的图像：
\begin{enumerate}[label=(\alph*), itemsep=2pt, topsep=3pt]
  \item $y = |f(x)|$;
  \item $y = \dfrac{1}{2}(|f(x)| + f(x))$;
  \item $y = \dfrac{1}{2}(|f(x)| - f(x))$;
  \item $y = f^2(x)$;
  \item $y = \sqrt{f(x)}$;
  \item $y = \ln f(x)$;
  \item $y = f(f(x))$;
  \item $y = \operatorname{sgn} f(x)$;
  \item $y = [f(x)]$.
\end{enumerate}
\end{prob}

\begin{prob}[L]{329.1}
设
\[
f(x) = (x - a)(b - x) \quad (a < b),
\]
作下列函数的图像：
\begin{enumerate}[label=(\alph*), itemsep=2pt, topsep=3pt]
  \item $y = f(x)$;
  \item $y = f^2(x)$;
  \item $y = \dfrac{1}{f(x)}$;
  \item $y = \sqrt{f(x)}$;
  \item $y = \mathrm{e}^{f(x)}$;
  \item $y = \lg f(x)$;
  \item $y = \operatorname{arccot} f(x)$.
\end{enumerate}
\end{prob}

\begin{prob}[L]{329.2}
若 (1) $f(x) = x^2$; (2) $f(x) = x^3$，作如下函数的图像：
\begin{enumerate}[label=(\alph*), itemsep=2pt, topsep=3pt]
  \item $y = \arcsin(\sin f(x))$;
  \item $y = \arcsin(\cos f(x))$;
  \item $y = \arccos(\sin f(x))$;
  \item $y = \arccos(\cos f(x))$;
  \item $y = \arctan(\tan f(x))$.
\end{enumerate}
\end{prob}

\begin{prob}[L]{330}
已知函数 $y = f(x)$ 和 $y = g(x)$ 的图像，作下列函数的图像：
\begin{enumerate}[label=(\alph*), itemsep=2pt, topsep=3pt]
  \item $y = f(x) + g(x)$;
  \item $y = f(x)g(x)$;
  \item $y = f(g(x))$.
\end{enumerate}
\end{prob}

\probgroup{利用图像的加法，作下列函数的图像：}

\begin{prob}[S]{331}
$y = 1 + x + \mathrm{e}^x$.
\end{prob}

\begin{prob}[S]{332}
$y = (x + 1)^{-2} + (x - 1)^{-2}$.
\end{prob}

\begin{prob}[S]{333}
$y = x + \sin x$.
\end{prob}

\begin{prob}[S]{334}
$y = x + \arctan x$.
\end{prob}

\begin{prob}[S]{335}
$y = \cos x + \dfrac{1}{2}\cos 2x + \dfrac{1}{3}\cos 3x$.
\end{prob}

\begin{prob}[S]{336}
$y = \sin x - \dfrac{1}{3}\sin 3x + \dfrac{1}{5}\sin 5x$.
\end{prob}

\begin{prob}[S]{337}
$y = \sin^4 x + \cos^4 x$.
\end{prob}

\begin{prob}[S]{338}
$y = |1 - x| + |1 + x|$.
\end{prob}

\begin{prob}[S]{339}
$y = |1 - x| - |1 + x|$.
\end{prob}

\begin{prob}[L]{340}
作双曲函数的图像：
\begin{enumerate}[label=(\alph*), itemsep=2pt, topsep=3pt]
  \item $y = \cosh x$，式中 $\cosh x = \dfrac{1}{2}(\mathrm{e}^x + \mathrm{e}^{-x})$;
  \item $y = \sinh x$，式中 $\sinh x = \dfrac{1}{2}(\mathrm{e}^x - \mathrm{e}^{-x})$;
  \item $y = \tanh x$，式中 $\tanh x = \dfrac{\sinh x}{\cosh x}$.
\end{enumerate}
\end{prob}

\probgroup{利用图像的乘法，作下列函数的图像：}

\begin{prob}[S]{341}
$y = x\sin x$.
\end{prob}

\begin{prob}[S]{342}
$y = x\cos x$.
\end{prob}

\begin{prob}[S]{343}
$y = x^2 \sin^2 x$.
\end{prob}

\begin{prob}[S]{344}
$y = \dfrac{\sin x}{1 + x^2}$.
\end{prob}

\begin{prob}[S]{345}
$y = \mathrm{e}^{-x^2} \cos 2x$.
\end{prob}

\begin{prob}[S]{346}
$y = x \operatorname{sgn}(\sin x)$.
\end{prob}

\begin{prob}[S]{347}
$y = [x] |\sin \pi x|$.
\end{prob}

\begin{prob}[S]{348}
$y = \cos x \cdot \operatorname{sgn}(\sin x)$.
\end{prob}

\begin{prob}[L]{349}
设
\[
f(x) = \begin{cases}
1 - |x|, & |x| \leqslant 1, \\
0, & |x| > 1.
\end{cases}
\]
作函数 $y = f(x) f(a - x)$ 当 (a) $a = 0$; (b) $a = 1$; (c) $a = 2$ 时的图像.
\end{prob}

\begin{prob}[M]{350}
作函数 $y = x + \sqrt{x} \operatorname{sgn}(\sin \pi x)$ 的图像.
\end{prob}

\probgroup{作函数 $y = \dfrac{1}{f(x)}$ 的图像，设：}

\begin{prob}[S]{351}
$f(x) = x^2(1 - x^2)$.
\end{prob}

\begin{prob}[S]{352}
$f(x) = x(1 - x)^2$.
\end{prob}

\begin{prob}[S]{353}
$f(x) = \sin^2 x$.
\end{prob}

\begin{prob}[S]{354}
$f(x) = \ln x$.
\end{prob}

\begin{prob}[S]{355}
$f(x) = \mathrm{e}^x \sin x$.
\end{prob}

\begin{prob}[M]{356}
设
\[
f(u) = \begin{cases}
-1, & -\infty < u < -1, \\
u, & -1 \leqslant u \leqslant 1, \\
1, & 1 < u < +\infty.
\end{cases}
\]
作复合函数 $y = f(u)$ 的图像，其中 $u = 2\sin x$.
\end{prob}

\begin{prob}[L]{357}
设
\[
\varphi(x) = \frac{1}{2}(x + |x|) \quad \text{和} \quad \psi(x) = \begin{cases}
x, & x < 0, \\
x^2, & x \geqslant 0.
\end{cases}
\]
作下列函数的图像：
\begin{enumerate}[label=(\alph*), itemsep=2pt, topsep=3pt]
  \item $y = \varphi(\varphi(x))$;
  \item $y = \varphi(\psi(x))$;
  \item $y = \psi(\varphi(x))$;
  \item $y = \psi(\psi(x))$.
\end{enumerate}
\end{prob}

\begin{prob}[L]{358}
设
\[
\varphi(x) = \begin{cases}
1, & |x| \leqslant 1, \\
0, & |x| > 1
\end{cases} \quad \text{及} \quad \psi(x) = \begin{cases}
2 - x^2, & |x| \leqslant 2, \\
2, & |x| > 2.
\end{cases}
\]
作函数
\begin{enumerate}[label=(\alph*), itemsep=2pt, topsep=3pt]
  \item $y = \varphi(\varphi(x))$;
  \item $y = \varphi(\psi(x))$;
  \item $y = \psi(\varphi(x))$;
  \item $y = \psi(\psi(x))$.
\end{enumerate}
的图像.
\end{prob}

\begin{prob}[L]{359}
把定义于正数区域 $x > 0$ 的函数 $f(x)$ 延拓到负数区域 $x < 0$，使所得的函数为：(1) 偶函数；(2) 奇函数，并分别作出所得函数的图像：
\begin{enumerate}[label=(\alph*), itemsep=2pt, topsep=3pt]
  \item $f(x) = 1 - x$;
  \item $f(x) = 2x - x^2$;
  \item $f(x) = \sqrt{x}$;
  \item $f(x) = \sin x$;
  \item $f(x) = \mathrm{e}^x$;
  \item $f(x) = \ln x$.
\end{enumerate}
\end{prob}

\begin{prob}[L]{360}
确定下列函数的图像在竖直方向上的对称轴：
\begin{enumerate}[label=(\alph*), itemsep=2pt, topsep=3pt]
  \item $y = ax^2 + bx + c$;
  \item $y = \dfrac{1}{x^2} + \dfrac{1}{(1 - x)^2}$;
  \item $y = \sqrt{a + x} + \sqrt{b - x}$ ($0 < a < b$);
  \item $y = a + b\cos x$.
\end{enumerate}
\end{prob}

\begin{prob}[L]{361}
确定下列函数的图像的对称中心：
\begin{enumerate}[label=(\alph*), itemsep=2pt, topsep=3pt]
  \item $y = ax + b$;
  \item $y = \dfrac{ax + b}{cx + d}$;
  \item $y = ax^3 + bx^2 + cx + d$;
  \item $y = \dfrac{1}{x - 1} + \dfrac{1}{x - 2} + \dfrac{1}{x - 3}$;
  \item $y = 1 + \sqrt[3]{x - 2}$.
\end{enumerate}
\end{prob}

\begin{prob}[L]{362}
作周期函数的图像：
\begin{enumerate}[label=(\alph*), itemsep=2pt, topsep=3pt]
  \item $y = |\sin x|$;
  \item $y = \operatorname{sgn}\cos x$;
  \item $y = f(x)$，其中 $f(x) = A\dfrac{x}{l}\left(2 - \dfrac{x}{l}\right)$，假设 $0 \leqslant x \leqslant 2l$ 和 $f(x + 2l) \equiv f(x)$;
  \item $y = [x] - 2\left[\dfrac{x}{2}\right]$;
  \item $y = (x)$，此处 $(x)$ 为从数 $x$ 到最近的整数的距离.
\end{enumerate}
\end{prob}

\begin{prob}[M]{363}
证明：若函数 $y = f(x)$ ($-\infty < x < +\infty$) 的图像关于竖直方向上的两条直线 $x = a$ 和 $x = b$ ($b > a$) 对称，则函数 $f(x)$ 为周期函数.
\end{prob}

\begin{prob}[M]{364}
证明：若函数 $y = f(x)$ ($-\infty < x < +\infty$) 的图像关于两点 $A(a, y_0)$ 和 $B(b, y_1)$ ($b > a$) 对称，则函数 $f(x)$ 是线性函数与周期函数的和. 特别是，若 $y_0 = y_1$，则函数 $f(x)$ 是周期函数.
\end{prob}

\begin{prob}[M]{365}
证明：若函数 $y = f(x)$ ($-\infty < x < +\infty$) 的图像关于点 $A(a, y_0)$ 和直线 $x = b$ ($b \neq a$) 对称，则函数 $f(x)$ 为周期函数.
\end{prob}

\begin{prob}[M]{366}
设 $f(x + 1) = 2f(x)$，且当 $0 \leqslant x \leqslant 1$ 时 $f(x) = x(1 - x)$，作函数 $y = f(x)$ ($-\infty < x < +\infty$) 的图像.
\end{prob}

\begin{prob}[M]{367}
设 $f(x + \pi) = f(x) + \sin x$，且当 $0 \leqslant x \leqslant \pi$ 时 $f(x) = 0$. 作函数 $y = f(x)$ ($-\infty < x < +\infty$) 的图像.
\end{prob}

\begin{prob}[L]{368}
作函数 $y = y(x)$ 的图像，设：
\begin{enumerate}[label=(\alph*), itemsep=2pt, topsep=3pt]
  \item $x = y - y^3$;
  \item $x = \dfrac{1 - y}{1 + y^2}$;
  \item $x = y - \ln y$;
  \item $x^2 = \sin y$.
\end{enumerate}
\end{prob}

\begin{prob}[L]{369}
作出下列用参数形式表示的函数的图像：
\begin{enumerate}[label=(\alph*), itemsep=2pt, topsep=3pt]
  \item $x = 1 - t$, $y = 1 - t^2$;
  \item $x = t + \dfrac{1}{t}$, $y = t + \dfrac{1}{t^2}$;
  \item $x = 10\cos t$, $y = \sin t$ (椭圆);
  \item $x = \cosh t$, $y = \sinh t$ (双曲线);
  \item $x = 5\cos^2 t$, $y = 3\sin^2 t$;
  \item $x = 2(t - \sin t)$, $y = 2(1 - \cos t)$ (摆线);
  \item $x = \sqrt[t+1]{t}$, $y = \sqrt[t]{t + 1}$ ($t > 0$).
\end{enumerate}
\end{prob}

\begin{prob}[L]{370.1}
作下列隐函数的图像：
\begin{enumerate}[label=(\alph*), itemsep=2pt, topsep=3pt]
  \item $x^2 - xy + y^2 = 1$ (椭圆);
  \item $x^3 + y^3 - 3xy = 0$ (笛卡儿叶形线);
  \item $\sqrt{x} + \sqrt{y} = 1$ (抛物线);
  \item $x^{\frac{2}{3}} + y^{\frac{2}{3}} = 4$ (星形线);
  \item $\sin x = \sin y$;
  \item $\cos(\pi x^2) = \cos(\pi y)$;
  \item $x^y = y^x$ ($x > 0, y > 0$);
  \item $x - |x| = y - |y|$.
\end{enumerate}
\end{prob}

\begin{prob}[L]{370.2}
作下列隐函数的图像：
\begin{enumerate}[label=(\alph*), itemsep=2pt, topsep=3pt]
  \item $\min(x, y) = 1$;
  \item $\max(x, y) = 1$;
  \item $\max(|x|, |y|) = 1$;
  \item $\min(x^2, y) = 1$.
\end{enumerate}
\end{prob}

\begin{prob}[L]{371.1}
在极坐标系 $(r, \varphi)$ 中作出下列函数 $r = r(\varphi)$ 的图像：
\begin{enumerate}[label=(\alph*), itemsep=2pt, topsep=3pt]
  \item $r = \varphi$ (阿基米德螺线);
  \item $r = \dfrac{\pi}{\varphi}$ (双曲螺线);
  \item $r = \dfrac{\varphi}{\varphi + 1}$ ($0 \leqslant \varphi < +\infty$);
  \item $r = 2^{\frac{\varphi}{2\pi}}$ (对数螺线);
  \item $r = 2(1 + \cos\varphi)$ (心脏线);
  \item $r = 10\sin 3\varphi$ (三瓣玫瑰线);
  \item $r^2 = 36\cos 2\varphi$ (伯努利双纽线);
  \item $\varphi = \dfrac{r}{r - 1}$ ($r > 1$);
  \item $\varphi = 2\pi\sin r$.
\end{enumerate}
\end{prob}

\begin{prob}[L]{371.2}
在极坐标系 $(r, \varphi)$ 中作出下列函数的图像：
\begin{enumerate}[label=(\alph*), itemsep=2pt, topsep=3pt]
  \item $\varphi = 4r - r^2$;
  \item $\varphi = \dfrac{12r}{1 + r^2}$;
  \item $r^2 + \varphi^2 = 100$.
\end{enumerate}
\end{prob}

\begin{prob}[M]{371.3}
在极坐标系 $(r, \varphi)$ 中作出以参数方程表示的函数的图像 ($t \geqslant 0$ 为参数)：
\begin{enumerate}[label=(\alph*), itemsep=2pt, topsep=3pt]
  \item $\begin{cases} \varphi = t\cos^2 t, \\ r = t\sin^2 t; \end{cases}$
  \item $\begin{cases} \varphi = 1 - 2^{-t}\sin\dfrac{\pi t}{2}, \\ r = 1 - 2^{-t}\cos\dfrac{\pi t}{2}. \end{cases}$
\end{enumerate}
\end{prob}

\begin{prob}[M]{372}
利用函数 $y = x^3 - 3x + 1$ 的图像近似地求解方程 $x^3 - 3x + 1 = 0$.
\end{prob}

\probgroup{用图解法解下列方程：}

\begin{prob}[S]{373}
$x^3 - 4x - 1 = 0$.
\end{prob}

\begin{prob}[S]{374}
$x^4 - 4x + 1 = 0$.
\end{prob}

\begin{prob}[S]{375}
$x = 2^{-x}$.
\end{prob}

\begin{prob}[S]{376}
$\lg x = 0.1x$.
\end{prob}

\begin{prob}[S]{377}
$10^x = x^2$.
\end{prob}

\begin{prob}[S]{378}
$\tan x = x$ ($0 \leqslant x \leqslant 2\pi$).
\end{prob}

\probgroup{用图解法解下列方程组：}

\begin{prob}[S]{379}
$\begin{cases} x + y^2 = 1, \\ 16x^2 + y = 4. \end{cases}$
\end{prob}

\begin{prob}[S]{380}
$\begin{cases} x^2 + y^2 = 100, \\ y = 10(x^2 - x - 2). \end{cases}$
\end{prob}

\probsec{§5. 函数的极限}

\begin{prob}[M]{381}
证明：函数
\[
f(x) = \begin{cases}
n, & x = \dfrac{m}{n} \text{（$m$ 和 $n$ 为互素的整数，且 $n > 0$）}, \\
0, & x \text{ 为无理数},
\end{cases}
\]
在每一点 $x$ 为有限的，但并非有界的（即在这点的任何邻域中是无界的）.
\end{prob}

\begin{prob}[M]{382}
若函数 $f(x)$ 在：(a) 开区间，(b) 闭区间内的每一点有定义且局部有界，则此函数在该开区间或闭区间内是否有界的？举出适当的例子.
\end{prob}

\begin{prob}[M]{383}
证明：函数
\[
f(x) = \frac{1 + x^2}{1 + x^4}
\]
在区间 $-\infty < x < +\infty$ 中是有界的.
\end{prob}

\begin{prob}[M]{384}
证明：函数
\[
f(x) = \frac{1}{x} \cos \frac{1}{x}
\]
在点 $x = 0$ 的任何邻域内是无界的，但在 $x \to 0$ 时不成为无穷大.
\end{prob}

\begin{prob}[M]{385}
研究函数 $f(x) = \ln x \cdot \sin^2 \dfrac{\pi}{x}$ 在区间 $0 < x < \varepsilon$ 内的有界性.
\end{prob}

\begin{prob}[M]{386}
证明：函数
\[
f(x) = \frac{x}{1 + x}
\]
在区域 $0 \leqslant x < +\infty$ 内有下确界 $m = 0$ 和上确界 $M = 1$.
\end{prob}

\begin{prob}[S]{387}
函数 $f(x)$ 在闭区间 $[a, b]$ 上有定义且单调上升，则在此闭区间内函数的下确界和上确界等于什么？
\end{prob}

\probgroup{求下列函数的下确界和上确界：}

\begin{prob}[S]{388}
$f(x) = x^2$ 在 $[-2, 5]$ 内.
\end{prob}

\begin{prob}[S]{389}
$f(x) = \dfrac{1}{1 + x^2}$ 在 $(-\infty, +\infty)$ 内.
\end{prob}

\begin{prob}[S]{390}
$f(x) = \dfrac{2x}{1 + x^2}$ 在 $(0, +\infty)$ 内.
\end{prob}

\begin{prob}[S]{391}
$f(x) = x + \dfrac{1}{x}$ 在 $(0, +\infty)$ 内.
\end{prob}

\begin{prob}[S]{392}
$f(x) = \sin x$ 在 $(0, +\infty)$ 内.
\end{prob}

\begin{prob}[S]{393}
$f(x) = \sin x + \cos x$ 在 $[0, 2\pi]$ 内.
\end{prob}

\begin{prob}[S]{394}
$f(x) = 2^x$ 在 $(-1, 2)$ 内.
\end{prob}

\begin{prob}[M]{395}
$f(x) = [x]$：(a) 在 $(0, 2)$ 内，(b) 在 $[0, 2]$ 内.
\end{prob}

\begin{prob}[S]{396}
$f(x) = x - [x]$ 在 $[0, 1]$ 内.
\end{prob}

\begin{prob}[L]{397}
求函数 $f(x) = x^2$ 在下列区间内的振幅：
\begin{enumerate}[label=(\alph*), itemsep=2pt, topsep=3pt]
  \item $(1, 3)$;
  \item $(1.9, 2.1)$;
  \item $(1.99, 2.01)$;
  \item $(1.999, 2.001)$.
\end{enumerate}
\end{prob}

\begin{prob}[L]{398}
求函数 $f(x) = \arctan\dfrac{1}{x}$ 在下列区间内的振幅：
\begin{enumerate}[label=(\alph*), itemsep=2pt, topsep=3pt]
  \item $(-1, 1)$;
  \item $(-0.1, 0.1)$;
  \item $(-0.01, 0.01)$;
  \item $(-0.001, 0.001)$.
\end{enumerate}
\end{prob}

\begin{prob}[M]{399}
设 $m[f]$ 和 $M[f]$ 分别为函数 $f(x)$ 在区间 $(a, b)$ 内的下确界和上确界.
证明：若 $f_1(x)$ 和 $f_2(x)$ 为定义于 $(a, b)$ 的函数，则
\[
m[f_1 + f_2] \geqslant m[f_1] + m[f_2],
\]
\[
M[f_1 + f_2] \leqslant M[f_1] + M[f_2].
\]
举出函数 $f_1(x)$ 和 $f_2(x)$ 的例子，使上述两个关系式为：(a) 等式，(b) 不等式.
\end{prob}

\begin{prob}[M]{400}
设函数 $f(x)$ 定义于区域 $[a, +\infty)$，并且在每一个闭区间 $[a, b]$ 上是有界的. 令
\[
m(x) = \inf_{a \leqslant \xi \leqslant x} f(\xi), \quad M(x) = \sup_{a \leqslant \xi \leqslant x} f(\xi).
\]
作函数 $y = m(x)$ 和 $y = M(x)$ 的图像，设
(a) $f(x) = \sin x$; \qquad (b) $f(x) = \cos x$.
\end{prob}

\begin{prob}[M]{401}
利用“$\varepsilon$-$\delta$”语言，证明：
\[
\lim_{x \to 2} x^2 = 4.
\]
填下表：
\begin{center}
\begin{tabular}{c|c|c|c|c|c}
\hline
$\varepsilon$ & $0.1$ & $0.01$ & $0.001$ & $0.0001$ & $\cdots$ \\
\hline
$\delta$ & & & & & \\
\hline
\end{tabular}
\end{center}
\end{prob}

\begin{prob}[M]{402}
以“$\varepsilon$-$\delta$”语言，证明：
\[
\lim_{x \to 1} \frac{1}{(1 - x)^2} = +\infty.
\]
填下表：
\begin{center}
\begin{tabular}{c|c|c|c|c|c}
\hline
$E$ & $10$ & $100$ & $1000$ & $10000$ & $\cdots$ \\
\hline
$\delta$ & & & & & \\
\hline
\end{tabular}
\end{center}
\end{prob}

\begin{prob}[L]{403}
利用不等式表示下列各式：
\begin{enumerate}[label=(\alph*), itemsep=2pt, topsep=3pt]
  \item $\lim\limits_{x \to a} f(x) = b$;
  \item $\lim\limits_{x \to a-0} f(x) = b$;
  \item $\lim\limits_{x \to a+0} f(x) = b$.
\end{enumerate}
举出适当的例子.
\end{prob}

\probgroup{利用不等式表示下列结论，并举出适当的例子：}

\begin{prob}[L]{404}
\begin{enumerate}[label=(\alph*), itemsep=2pt, topsep=3pt]
  \item $\lim\limits_{x \to \infty} f(x) = b$;
  \item $\lim\limits_{x \to -\infty} f(x) = b$;
  \item $\lim\limits_{x \to +\infty} f(x) = b$.
\end{enumerate}
\end{prob}

\begin{prob}[L]{405}
\begin{enumerate}[label=(\alph*), itemsep=2pt, topsep=3pt]
  \item $\lim\limits_{x \to a} f(x) = \infty$;
  \item $\lim\limits_{x \to a} f(x) = -\infty$;
  \item $\lim\limits_{x \to a} f(x) = +\infty$;
  \item $\lim\limits_{x \to a-0} f(x) = \infty$;
  \item $\lim\limits_{x \to a-0} f(x) = -\infty$;
  \item $\lim\limits_{x \to a-0} f(x) = +\infty$;
  \item $\lim\limits_{x \to a+0} f(x) = \infty$;
  \item $\lim\limits_{x \to a+0} f(x) = -\infty$;
  \item $\lim\limits_{x \to a+0} f(x) = +\infty$.
\end{enumerate}
\end{prob}

\begin{prob}[L]{406}
\begin{enumerate}[label=(\alph*), itemsep=2pt, topsep=3pt]
  \item $\lim\limits_{x \to \infty} f(x) = \infty$;
  \item $\lim\limits_{x \to \infty} f(x) = -\infty$;
  \item $\lim\limits_{x \to \infty} f(x) = +\infty$;
  \item $\lim\limits_{x \to -\infty} f(x) = \infty$;
  \item $\lim\limits_{x \to -\infty} f(x) = -\infty$;
  \item $\lim\limits_{x \to -\infty} f(x) = +\infty$;
  \item $\lim\limits_{x \to +\infty} f(x) = \infty$;
  \item $\lim\limits_{x \to +\infty} f(x) = -\infty$;
  \item $\lim\limits_{x \to +\infty} f(x) = +\infty$.
\end{enumerate}
\end{prob}

\begin{prob}[L]{407}
命 $y = f(x)$. 利用不等式表示下列结论：
\begin{enumerate}[label=(\alph*), itemsep=2pt, topsep=3pt]
  \item 当 $x \to a$ 时，$y \to b - 0$;
  \item 当 $x \to a - 0$ 时，$y \to b - 0$;
  \item 当 $x \to a + 0$ 时，$y \to b - 0$;
  \item 当 $x \to a$ 时，$y \to b + 0$;
  \item 当 $x \to a - 0$ 时，$y \to b + 0$;
  \item 当 $x \to a + 0$ 时，$y \to b + 0$;
  \item 当 $x \to \infty$ 时，$y \to b - 0$;
  \item 当 $x \to -\infty$ 时，$y \to b - 0$;
  \item 当 $x \to +\infty$ 时，$y \to b - 0$;
  \item 当 $x \to \infty$ 时，$y \to b + 0$;
  \item 当 $x \to -\infty$ 时，$y \to b + 0$;
  \item 当 $x \to +\infty$ 时，$y \to b + 0$.
\end{enumerate}
举出适当的例子.
\end{prob}

\begin{prob}[M]{408}
设
\[
P(x) = a_0 x^n + a_1 x^{n-1} + \cdots + a_n,
\]
式中 $a_i\ (i = 0, 1, \cdots, n;\ n \geqslant 1,\ a_0 \neq 0)$ 为实数. 证明：
\[
\lim_{x \to \infty} |P(x)| = +\infty.
\]
\end{prob}

\begin{prob}[M]{409}
设
\[
R(x) = \frac{a_0 x^n + a_1 x^{n-1} + \cdots + a_n}{b_0 x^m + b_1 x^{m-1} + \cdots + b_m},
\]
式中 $a_0 \neq 0, b_0 \neq 0$，证明：
\[
\lim_{x \to \infty} |R(x)| = \begin{cases}
\infty, & n > m, \\
\dfrac{a_0}{b_0}, & n = m, \\
0, & n < m.
\end{cases}
\]
\end{prob}

\begin{prob}[M]{410}
设 $R(x) = \dfrac{P(x)}{Q(x)}$，式中 $P(x)$ 和 $Q(x)$ 为 $x$ 的多项式，且 $P(a) = Q(a) = 0$. 表达式
\[
\lim_{x \to a} \frac{P(x)}{Q(x)}
\]
可能取何值？
\end{prob}

\probgroup{求下列各式之值：}

\begin{prob}[L]{411}
\begin{enumerate}[label=(\alph*), itemsep=2pt, topsep=3pt]
  \item $\lim\limits_{x \to 0} \dfrac{x^2 - 1}{2x^2 - x - 1}$;
  \item $\lim\limits_{x \to 1} \dfrac{x^2 - 1}{2x^2 - x - 1}$;
  \item $\lim\limits_{x \to \infty} \dfrac{x^2 - 1}{2x^2 - x - 1}$.
\end{enumerate}
\end{prob}

\begin{prob}[S]{412}
$\lim\limits_{x \to 0} \dfrac{(1 + x)(1 + 2x)(1 + 3x) - 1}{x}$.
\end{prob}

\begin{prob}[S]{413}
$\lim\limits_{x \to 0} \dfrac{(1 + x)^5 - (1 + 5x)}{x^2 + x^5}$.
\end{prob}

\begin{prob}[S]{414}
$\lim\limits_{x \to 0} \dfrac{(1 + mx)^n - (1 + nx)^m}{x^2}$（$m$ 与 $n$ 为正整数）.
\end{prob}

\begin{prob}[S]{415}
$\lim\limits_{x \to \infty} \dfrac{(x - 1)(x - 2)(x - 3)(x - 4)(x - 5)}{(5x - 1)^5}$.
\end{prob}

\begin{prob}[S]{416}
$\lim\limits_{x \to \infty} \dfrac{(2x - 3)^{20}(3x + 2)^{30}}{(2x + 1)^{50}}$.
\end{prob}

\begin{prob}[S]{417}
$\lim\limits_{x \to \infty} \dfrac{(x + 1)(x^2 + 1)\cdots(x^n + 1)}{[(nx)^n + 1]^{\frac{n+1}{2}}}$.
\end{prob}

\begin{prob}[S]{418}
$\lim\limits_{x \to 3} \dfrac{x^2 - 5x + 6}{x^2 - 8x + 15}$.
\end{prob}

\begin{prob}[S]{419}
$\lim\limits_{x \to 1} \dfrac{x^3 - 3x + 2}{x^4 - 4x + 3}$.
\end{prob}

\begin{prob}[S]{420}
$\lim\limits_{x \to 1} \dfrac{x^4 - 3x + 2}{x^5 - 4x + 3}$.
\end{prob}

\begin{prob}[S]{421}
$\lim\limits_{x \to 2} \dfrac{x^3 - 2x^2 - 4x + 8}{x^4 - 8x^2 + 16}$.
\end{prob}

\begin{prob}[S]{422}
$\lim\limits_{x \to -1} \dfrac{x^3 - 2x - 1}{x^5 - 2x - 1}$.
\end{prob}

\begin{prob}[S]{423}
$\lim\limits_{x \to 2} \dfrac{(x^2 - x - 2)^{20}}{(x^3 - 12x + 16)^{10}}$.
\end{prob}

\begin{prob}[M]{424}
(a) $\lim\limits_{x \to 1} \dfrac{x + x^2 + \cdots + x^n - n}{x - 1}$; \qquad (b) $\lim\limits_{x \to 1} \dfrac{x^{100} - 2x + 1}{x^{50} - 2x + 1}$.
\end{prob}

\begin{prob}[S]{425}
$\lim\limits_{x \to 1} \dfrac{x^m - 1}{x^n - 1}$（$m$ 和 $n$ 为正整数）.
\end{prob}

\begin{prob}[S]{426}
$\lim\limits_{x \to a} \dfrac{(x^n - a^n) - na^{n-1}(x - a)}{(x - a)^2}$（$n$ 为正整数）.
\end{prob}

\begin{prob}[S]{427}
$\lim\limits_{x \to 1} \dfrac{x^{n+1} - (n + 1)x + n}{(x - 1)^2}$（$n$ 为正整数）.
\end{prob}

\begin{prob}[S]{428}
$\lim\limits_{x \to 1} \left(\dfrac{m}{1 - x^m} - \dfrac{n}{1 - x^n}\right)$（$m$ 和 $n$ 为正整数）.
\end{prob}

\begin{prob}[S]{429}
$\lim\limits_{n \to \infty} \dfrac{1}{n} \left[\left(x + \dfrac{a}{n}\right) + \left(x + \dfrac{2a}{n}\right) + \cdots + \left(x + \dfrac{(n - 1)a}{n}\right)\right]$.
\end{prob}

\begin{prob}[M]{430}
$\lim\limits_{n \to \infty} \dfrac{1}{n} \left[\left(x + \dfrac{a}{n}\right)^2 + \left(x + \dfrac{2a}{n}\right)^2 + \cdots + \left(x + \dfrac{(n - 1)a}{n}\right)^2\right]$.
\probhint{参阅题 2.}
\end{prob}

\begin{prob}[S]{431}
$\lim\limits_{n \to \infty} \dfrac{1^2 + 3^2 + \cdots + (2n - 1)^2}{2^2 + 4^2 + \cdots + (2n)^2}$.
\end{prob}

\begin{prob}[M]{432}
$\lim\limits_{n \to \infty} \left(\dfrac{1^3 + 2^3 + \cdots + n^3}{n^3} - \dfrac{\pi}{4}\right)$.
\probhint{参阅题 3.}
\end{prob}

\begin{prob}[S]{433}
$\lim\limits_{n \to \infty} \dfrac{1^3 + 4^3 + 7^3 + \cdots + (3n - 2)^3}{[1 + 4 + 7 + \cdots + (3n - 2)]^2}$.
\end{prob}

\begin{prob}[M]{434}
把由抛物线 $y = b\left(\dfrac{x}{a}\right)^2$，$Ox$ 轴及直线 $x = a$ 所围成的曲边三角形 $OAM$（图 3）的面积，当作以 $\dfrac{a}{n}$ 为底的各内接矩形面积之和当 $n \to \infty$ 的极限值，求此面积.
\probfig{fig_434.png}{图 3}
\end{prob}

\probgroup{求极限：}

\begin{prob}[S]{435}
$\lim\limits_{x \to +\infty} \dfrac{\sqrt{x + \sqrt{x + \sqrt{x}}}}{\sqrt{x + 1}}$.
\end{prob}

\begin{prob}[S]{436}
$\lim\limits_{x \to +\infty} \dfrac{\sqrt{x} + \sqrt[3]{x} + \sqrt[4]{x}}{\sqrt{2x + 1}}$.
\end{prob}

\begin{prob}[S]{437}
$\lim\limits_{x \to 4} \dfrac{\sqrt{1 + 2x} - 3}{\sqrt{x} - 2}$.
\end{prob}

\begin{prob}[S]{438}
$\lim\limits_{x \to -8} \dfrac{\sqrt{1 - x} - 3}{2 + \sqrt[3]{x}}$.
\end{prob}

\begin{prob}[S]{439}
$\lim\limits_{x \to a} \dfrac{\sqrt{x} - \sqrt{a} + \sqrt{x - a}}{\sqrt{x^2 - a^2}}$.
\end{prob}

\begin{prob}[S]{440}
$\lim\limits_{x \to 3} \dfrac{\sqrt{x + 13} - 2\sqrt{x + 1}}{x^2 - 9}$.
\end{prob}

\begin{prob}[S]{441}
$\lim\limits_{x \to -2} \dfrac{\sqrt[3]{x - 6} + 2}{x^3 + 8}$.
\end{prob}

\begin{prob}[S]{442}
$\lim\limits_{x \to 16} \dfrac{\sqrt[4]{x} - 2}{\sqrt{x} - 4}$.
\end{prob}

\begin{prob}[S]{443}
$\lim\limits_{x \to 8} \dfrac{\sqrt{9 + 2x} - 5}{\sqrt[3]{x} - 2}$.
\end{prob}

\begin{prob}[S]{444}
$\lim\limits_{x \to 0} \dfrac{\sqrt[n]{1 + x} - 1}{x}$（$n$ 为整数）.
\end{prob}

\begin{prob}[S]{445}
$\lim\limits_{x \to 0} \dfrac{\sqrt{1 - 2x - x^2} - (1 + x)}{x}$.
\end{prob}

\begin{prob}[S]{446}
$\lim\limits_{x \to 0} \dfrac{\sqrt[3]{8 + 3x - x^2} - 2}{x + x^2}$.
\end{prob}

\begin{prob}[S]{447}
$\lim\limits_{x \to 0} \dfrac{\sqrt[3]{27 + x} - \sqrt[3]{27 - x}}{x + 2\sqrt[3]{x^4}}$.
\end{prob}

\begin{prob}[S]{448}
$\lim\limits_{x \to 0} \dfrac{\sqrt{1 + x} - \sqrt{1 - x}}{\sqrt[3]{1 + x} - \sqrt[3]{1 - x}}$.
\end{prob}

\begin{prob}[S]{449}
$\lim\limits_{x \to 7} \dfrac{\sqrt{x + 2} - \sqrt[3]{x + 20}}{\sqrt[4]{x + 9} - 2}$.
\end{prob}

\begin{prob}[S]{450}
$\lim\limits_{x \to 0} \dfrac{\sqrt[3]{1 + \dfrac{x}{3}} - \sqrt[4]{1 + \dfrac{x}{4}}}{1 - \sqrt{1 - \dfrac{x}{2}}}$.
\end{prob}

\begin{prob}[S]{451}
$\lim\limits_{x \to 0} \dfrac{x^2}{\sqrt[5]{1 + 5x} - (1 + x)}$.
\end{prob}

\begin{prob}[S]{452}
$\lim\limits_{x \to 0} \dfrac{\sqrt[m]{1 + \alpha x} - \sqrt[n]{1 + \beta x}}{x}$（$m$ 及 $n$ 为整数）.
\end{prob}

\begin{prob}[S]{453}
$\lim\limits_{x \to 0} \dfrac{\sqrt[m]{1 + \alpha x}\sqrt[n]{1 + \beta x} - 1}{x}$（$m$ 及 $n$ 为整数）.
\end{prob}

\begin{prob}[M]{454}
设 $P(x) = a_1 x + a_2 x^2 + \cdots + a_n x^n$ 又 $m$ 为整数，求证：
\[
\lim_{x \to 0} \frac{\sqrt[m]{1 + p(x)} - 1}{x} = \frac{a_1}{m}.
\]
\end{prob}

\probgroup{求下列的极限：}

\begin{prob}[S]{455.1}
$\lim\limits_{x \to 1} \dfrac{\sqrt[m]{x} - 1}{\sqrt[n]{x} - 1}$（$m$ 及 $n$ 为整数）.
\end{prob}

\begin{prob}[S]{455.2}
$\lim\limits_{x \to 1} \left(\dfrac{3}{1 - \sqrt{x}} - \dfrac{2}{1 - \sqrt[3]{x}}\right)$.
\end{prob}

\begin{prob}[S]{456}
$\lim\limits_{x \to 1} \dfrac{(1 - \sqrt{x})(1 - \sqrt[3]{x})\cdots(1 - \sqrt[n]{x})}{(1 - x)^{n-1}}$.
\end{prob}

\begin{prob}[S]{457}
$\lim\limits_{x \to +\infty} \left[\sqrt{(x + a)(x + b)} - x\right]$.
\end{prob}

\begin{prob}[S]{458}
$\lim\limits_{x \to +\infty} \left(\sqrt{x + \sqrt{x + \sqrt{x}}} - \sqrt{x}\right)$.
\end{prob}

\begin{prob}[S]{459}
$\lim\limits_{x \to +\infty} x\left(\sqrt{x^2 + 2x} - 2\sqrt{x^2 + x} + x\right)$.
\end{prob}

\begin{prob}[M]{460}
$\lim\limits_{x \to +0} \left(\sqrt{\dfrac{1}{x} + \sqrt{\dfrac{1}{x} + \sqrt{\dfrac{1}{x}}}} - \sqrt{\dfrac{1}{x} - \sqrt{\dfrac{1}{x} + \sqrt{\dfrac{1}{x}}}}\right)$.
\end{prob}

\begin{prob}[S]{461}
$\lim\limits_{x \to \infty} \left(\sqrt[3]{x^3 + x^2 + 1} - \sqrt[3]{x^3 - x^2 + 1}\right)$.
\end{prob}

\begin{prob}[S]{462}
$\lim\limits_{x \to +\infty} \left(\sqrt[3]{x^3 + 3x^2} - \sqrt{x^2 - 2x}\right)$.
\end{prob}

\begin{prob}[S]{463}
$\lim\limits_{x \to \infty} x^{\frac{1}{3}} \left[(x + 1)^{\frac{2}{3}} - (x - 1)^{\frac{2}{3}}\right]$.
\end{prob}

\begin{prob}[S]{464}
$\lim\limits_{x \to +\infty} x^{\frac{3}{2}} \left(\sqrt{x + 2} - 2\sqrt{x + 1} + \sqrt{x}\right)$.
\end{prob}

\begin{prob}[S]{465}
$\lim\limits_{x \to +\infty} \left[\sqrt[n]{(x + a_1)\cdots(x + a_n)} - x\right]$.
\end{prob}

\begin{prob}[S]{466}
$\lim\limits_{x \to +\infty} \dfrac{\left(x - \sqrt{x^2 - 1}\right)^n + \left(x + \sqrt{x^2 - 1}\right)^n}{x^n}$（$n$ 为正整数）.
\end{prob}

\begin{prob}[S]{467}
$\lim\limits_{x \to 0} \dfrac{\left(\sqrt{1 + x^2} + x\right)^n - \left(\sqrt{1 + x^2} - x\right)^n}{x}$（$n$ 为正整数）.
\end{prob}

\begin{prob}[M]{468}
设二次方程 $ax^2 + bx + c = 0$ 的系数 $a$ 趋于零，系数 $b$ 与 $c$ 为常数，且 $b \neq 0$，试研究此二次方程之二根 $x_1$ 及 $x_2$ 的性质.
\end{prob}

\begin{prob}[M]{469}
从条件
\[
\lim_{x \to \infty} \left(\frac{x^2 + 1}{x + 1} - ax - b\right) = 0
\]
求常数 $a$ 和 $b$.
\end{prob}

\begin{prob}[M]{470}
从条件
\[
\lim_{x \to -\infty} \left(\sqrt{x^2 - x + 1} - a_1 x - b_1\right) = 0,
\]
\[
\lim_{x \to +\infty} \left(\sqrt{x^2 - x + 1} - a_2 x - b_2\right) = 0
\]
求常数 $a_i$ 和 $b_i\ (i = 1, 2)$.
\end{prob}

\probgroup{求下列极限：}

\begin{prob}[S]{471}
$\lim\limits_{x \to 0} \dfrac{\sin 5x}{x}$.
\end{prob}

\begin{prob}[S]{472}
$\lim\limits_{x \to \infty} \dfrac{\sin x}{x}$.
\end{prob}

\begin{prob}[S]{473}
$\lim\limits_{x \to \pi} \dfrac{\sin mx}{\sin nx}$（$m$ 及 $n$ 为整数）.
\end{prob}

\begin{prob}[L]{474}
\begin{enumerate}[label=(\alph*), itemsep=2pt, topsep=3pt]
  \item $\lim\limits_{x \to 0} \dfrac{1 - \cos x}{x^2}$;
  \item $\lim\limits_{x \to 0} \dfrac{\tan x}{x}$;
  \item $\lim\limits_{x \to 0} x \cot 3x$.
\end{enumerate}
\end{prob}

\begin{prob}[S]{475}
$\lim\limits_{x \to 0} \dfrac{\tan x - \sin x}{\sin^3 x}$.
\end{prob}

\begin{prob}[S]{476}
$\lim\limits_{x \to 0} \dfrac{\sin 5x - \sin 3x}{\sin x}$.
\end{prob}

\begin{prob}[S]{477}
$\lim\limits_{x \to 0} \dfrac{\cos x - \cos 3x}{x^2}$.
\end{prob}

\begin{prob}[S]{478}
$\lim\limits_{x \to 0} \dfrac{1 + \sin x - \cos x}{1 + \sin px - \cos px}$.
\end{prob}

\begin{prob}[S]{479}
$\lim\limits_{x \to \frac{\pi}{4}} \tan 2x \tan\left(\dfrac{\pi}{4} - x\right)$.
\end{prob}

\begin{prob}[S]{480}
$\lim\limits_{x \to 1} (1 - x) \tan \dfrac{\pi x}{2}$.
\end{prob}

\begin{prob}[L]{481}
证明等式：
\begin{enumerate}[label=(\alph*), itemsep=2pt, topsep=3pt]
  \item $\lim\limits_{x \to a} \sin x = \sin a$;
  \item $\lim\limits_{x \to a} \cos x = \cos a$;
  \item $\lim\limits_{x \to a} \tan x = \tan a\ \left(a \neq \dfrac{2n - 1}{2}\pi;\ n = 0, \pm 1, \pm 2, \cdots\right)$.
\end{enumerate}
\end{prob}

\probgroup{求下列极限：}

\begin{prob}[S]{482}
$\lim\limits_{x \to a} \dfrac{\sin x - \sin a}{x - a}$.
\end{prob}

\begin{prob}[S]{483}
$\lim\limits_{x \to a} \dfrac{\cos x - \cos a}{x - a}$.
\end{prob}

\begin{prob}[S]{484}
$\lim\limits_{x \to a} \dfrac{\tan x - \tan a}{x - a}$.
\end{prob}

\begin{prob}[S]{485}
$\lim\limits_{x \to a} \dfrac{\cot x - \cot a}{x - a}$.
\end{prob}

\begin{prob}[S]{486}
$\lim\limits_{x \to a} \dfrac{\sec x - \sec a}{x - a}$.
\end{prob}

\begin{prob}[S]{487}
$\lim\limits_{x \to a} \dfrac{\csc x - \csc a}{x - a}$.
\end{prob}

\begin{prob}[S]{488}
$\lim\limits_{x \to 0} \dfrac{\sin(a + 2x) - 2\sin(a + x) + \sin a}{x^2}$.
\end{prob}

\begin{prob}[S]{489}
$\displaystyle \lim_{x\to 0} \frac{\cos(a+2x) - 2\cos(a+x) + \cos a}{x^2}.$
\end{prob}

\begin{prob}[S]{490}
$\displaystyle \lim_{x\to 0} \frac{\tan(a+2x) - 2\tan(a+x) + \tan a}{x^2}.$
\end{prob}

\begin{prob}[S]{491}
$\displaystyle \lim_{x\to 0} \frac{\cot(a+2x) - 2\cot(a+x) + \cot a}{x^2}.$
\end{prob}

\begin{prob}[S]{492}
$\displaystyle \lim_{x\to 0} \frac{\sin(a+x)\sin(a+2x) - \sin^2 a}{x}.$
\end{prob}

\begin{prob}[S]{493}
$\displaystyle \lim_{x\to \frac{\pi}{6}} \frac{2\sin^2 x + \sin x - 1}{2\sin^2 x - 3\sin x + 1}.$
\end{prob}

\begin{prob}[S]{494}
$\displaystyle \lim_{x\to 0} \frac{1 - \cos x \cos 2x \cos 3x}{1 - \cos x}.$
\end{prob}

\begin{prob}[S]{495}
$\displaystyle \lim_{x\to \frac{\pi}{3}} \frac{\sin\left(x - \dfrac{\pi}{3}\right)}{1 - 2\cos x}.$
\end{prob}

\begin{prob}[S]{496}
$\displaystyle \lim_{x\to \frac{\pi}{3}} \frac{\tan^3 x - 3\tan x}{\cos\left(x + \dfrac{\pi}{6}\right)}.$
\end{prob}

\begin{prob}[S]{497}
$\displaystyle \lim_{x\to 0} \frac{\tan(a+x)\tan(a-x) - \tan^2 a}{x^2}.$
\end{prob}

\begin{prob}[S]{498}
$\displaystyle \lim_{x\to \frac{\pi}{4}} \frac{1 - \cot^3 x}{2 - \cot x - \cot^3 x}.$
\end{prob}

\begin{prob}[S]{499}
$\displaystyle \lim_{x\to 0} \frac{\sqrt{1+\tan x} - \sqrt{1+\sin x}}{x^3}.$
\end{prob}

\begin{prob}[S]{500}
$\displaystyle \lim_{x\to 0} \frac{x^2}{\sqrt{1+x\sin x} - \sqrt{\cos x}}.$
\end{prob}

\begin{prob}[S]{501}
$\displaystyle \lim_{x\to 0} \frac{\sqrt{\cos x} - \sqrt[3]{\cos x}}{\sin^2 x}.$
\end{prob}

\begin{prob}[S]{502}
$\displaystyle \lim_{x\to 0} \frac{\sqrt{1-\cos x^2}}{1-\cos x}.$
\end{prob}

\begin{prob}[S]{503}
$\displaystyle \lim_{x\to 0} \frac{1-\sqrt{\cos x}}{1-\cos(\sqrt{x})}.$
\end{prob}

\begin{prob}[S]{504}
$\displaystyle \lim_{x\to 0} \frac{1 - \cos x\sqrt{\cos 2x}\sqrt[3]{\cos 3x}}{x^2}.$
\end{prob}

\begin{prob}[S]{505}
$\displaystyle \lim_{x\to +\infty} (\sin\sqrt{x+1} - \sin\sqrt{x}).$
\end{prob}

\begin{prob}[L]{506}
\begin{enumerate}[label=(\alph*), itemsep=2pt, topsep=3pt]
  \item $\displaystyle \lim_{x\to 0} \left(\frac{1+x}{2+x}\right)^{\frac{1-\sqrt{x}}{1-x}}$;
  \item $\displaystyle \lim_{x\to 1} \left(\frac{1+x}{2+x}\right)^{\frac{1-\sqrt{x}}{1-x}}$;
  \item $\displaystyle \lim_{x\to +\infty} \left(\frac{1+x}{2+x}\right)^{\frac{1-\sqrt{x}}{1-x}}.$
\end{enumerate}
\end{prob}

\begin{prob}[S]{507}
$\displaystyle \lim_{x\to \infty} \left(\frac{x+2}{2x-1}\right)^{x^2}.$
\end{prob}

\begin{prob}[S]{508}
$\displaystyle \lim_{x\to \infty} \left(\frac{3x^2 - x + 1}{2x^2 + x + 1}\right)^{\frac{x^3}{1-x}}.$
\end{prob}

\begin{prob}[S]{509}
$\displaystyle \lim_{n\to \infty} \left(\sin^n \frac{2n\pi}{3n+1}\right).$
\end{prob}

\begin{prob}[S]{510}
$\displaystyle \lim_{x\to \frac{\pi}{4}+0} \left[\tan\left(\frac{\pi}{8}+x\right)\right]^{\tan 2x}.$
\end{prob}

\begin{prob}[S]{511}
$\displaystyle \lim_{x\to \infty} \left(\frac{x^2-1}{x^2+1}\right)^{\frac{x-1}{x+1}}.$
\end{prob}

\begin{prob}[S]{512}
$\displaystyle \lim_{x\to \infty} \left(\frac{x^2+1}{x^2-2}\right)^{x^2}.$
\end{prob}

\begin{prob}[S]{513}
$\displaystyle \lim_{x\to \infty} \left(\frac{x^2+2x-1}{2x^2-3x-2}\right)^{\frac{1}{x}}.$
\end{prob}

\begin{prob}[S]{514}
$\displaystyle \lim_{x\to 0} \sqrt[x]{1-2x}.$
\end{prob}

\begin{prob}[S]{515}
$\displaystyle \lim_{x\to \infty} \left(\frac{x+a}{x-a}\right)^x.$
\end{prob}

\begin{prob}[S]{516}
$\displaystyle \lim_{x\to +\infty} \left(\frac{a_1 x + b_1}{a_2 x + b_2}\right)^x \ (a_1 > 0, a_2 > 0).$
\end{prob}

\begin{prob}[S]{517}
$\displaystyle \lim_{x\to 0} (1+x^2)^{\cot^2 x}.$
\end{prob}

\begin{prob}[S]{518}
$\displaystyle \lim_{x\to 1} (1+\sin\pi x)^{\cot\pi x}.$
\end{prob}

\begin{prob}[M]{519}
(a) $\displaystyle \lim_{x\to 0} \left(\frac{1+\tan x}{1+\sin x}\right)^{\frac{1}{\sin x}}$; \qquad (b) $\displaystyle \lim_{x\to 0} \left(\frac{1+\tan x}{1+\sin x}\right)^{\frac{1}{\sin^3 x}}.$
\end{prob}

\begin{prob}[S]{520}
$\displaystyle \lim_{x\to a} \left(\frac{\sin x}{\sin a}\right)^{\frac{1}{x-a}}.$
\end{prob}

\begin{prob}[S]{521}
$\displaystyle \lim_{x\to 0} \left(\frac{\cos x}{\cos 2x}\right)^{\frac{1}{x^2}}.$
\end{prob}

\begin{prob}[S]{522}
$\displaystyle \lim_{x\to \frac{\pi}{4}} (\tan x)^{\tan 2x}.$
\end{prob}

\begin{prob}[S]{523}
$\displaystyle \lim_{x\to \frac{\pi}{2}} (\sin x)^{\tan x}.$
\end{prob}

\begin{prob}[S]{524}
$\displaystyle \lim_{x\to 0} \left[\tan\left(\frac{\pi}{4}-x\right)\right]^{\cot x}.$
\end{prob}

\begin{prob}[S]{525}
$\displaystyle \lim_{x\to \infty} \left(\sin\frac{1}{x} + \cos\frac{1}{x}\right)^x.$
\end{prob}

\begin{prob}[S]{526}
$\displaystyle \lim_{x\to 0} \sqrt[x]{\cos\sqrt{x}}.$
\end{prob}

\begin{prob}[S]{527}
$\displaystyle \lim_{n\to \infty} \left(\frac{n+x}{n-1}\right)^n.$
\end{prob}

\begin{prob}[S]{528}
$\displaystyle \lim_{n\to \infty} \cos^n \frac{x}{\sqrt{n}}.$
\end{prob}

\begin{prob}[S]{529}
$\displaystyle \lim_{x\to 0} \frac{\ln(1+x)}{x}.$
\end{prob}

\begin{prob}[S]{530}
$\displaystyle \lim_{x\to +\infty} x[\ln(x+1) - \ln x].$
\end{prob}

\begin{prob}[S]{531}
$\displaystyle \lim_{x\to a} \frac{\ln x - \ln a}{x-a} \ (a > 0).$
\end{prob}

\begin{prob}[S]{532}
$\displaystyle \lim_{x\to +\infty} [\sin\ln(x+1) - \sin\ln x].$
\end{prob}

\begin{prob}[S]{533}
$\displaystyle \lim_{x\to +\infty} \frac{\ln(x^2-x+1)}{\ln(x^{10}+x+1)}.$
\end{prob}

\begin{prob}[S]{534}
$\displaystyle \lim_{x\to \infty} \left(\lg\frac{100+x^2}{1+100x^2}\right).$
\end{prob}

\begin{prob}[S]{535}
$\displaystyle \lim_{x\to +\infty} \frac{\ln(2+\mathrm{e}^{3x})}{\ln(3+\mathrm{e}^{2x})}.$
\end{prob}

\begin{prob}[S]{536}
$\displaystyle \lim_{x\to +\infty} \frac{\ln(1+\sqrt{x}+\sqrt[3]{x})}{\ln(1+\sqrt[3]{x}+\sqrt[4]{x})}.$
\end{prob}

\begin{prob}[S]{537}
$\displaystyle \lim_{h\to 0} \frac{\log(x+h) + \log(x-h) - 2\log x}{h^2} \ (x > 0).$
\end{prob}

\begin{prob}[S]{538}
$\displaystyle \lim_{x\to 0} \frac{\ln\tan\left(\dfrac{\pi}{4}+ax\right)}{\sin bx}.$
\end{prob}

\begin{prob}[S]{539}
$\displaystyle \lim_{x\to 0} \frac{\ln\cos ax}{\ln\cos bx}.$
\end{prob}

\begin{prob}[M]{540}
(a) $\displaystyle \lim_{x\to 0} \left(\ln\frac{nx+\sqrt{1-n^2 x^2}}{x+\sqrt{1-x^2}}\right)$; \qquad (b) $\displaystyle \lim_{x\to 0} \frac{\ln(nx+\sqrt{1-n^2 x^2})}{\ln(x+\sqrt{1-x^2})}.$
\end{prob}

\begin{prob}[S]{541}
$\displaystyle \lim_{x\to 0} \frac{a^x - 1}{x} \ (a > 0).$
\end{prob}

\begin{prob}[S]{542}
$\displaystyle \lim_{x\to a} \frac{a^x - x^a}{x-a} \ (a > 0).$
\end{prob}

\begin{prob}[S]{543}
$\displaystyle \lim_{x\to a} \frac{x^x - a^a}{x-a} \ (a > 0).$
\end{prob}

\begin{prob}[S]{544}
$\displaystyle \lim_{x\to 0} (x+\mathrm{e}^x)^{\frac{1}{x}}.$
\end{prob}

\begin{prob}[L]{545}
\begin{enumerate}[label=(\alph*), itemsep=2pt, topsep=3pt]
  \item $\displaystyle \lim_{x\to 0} \left(\frac{1+x\cdot 2^x}{1+x\cdot 3^x}\right)^{\frac{1}{x^2}}$;
  \item $\displaystyle \lim_{x\to 0} \left(\frac{1+\sin x\cos\alpha x}{1+\sin x\cos\beta x}\right)^{\cot^3 x}$;
  \item $\displaystyle \lim_{x\to 1} \frac{\sin(\pi x^\alpha)}{\sin(\pi x^\beta)}$;
  \item $\displaystyle \lim_{x\to 1} \frac{\sin^2(\pi\cdot 2^x)}{\ln(\cos(\pi\cdot 2^x))}.$
\end{enumerate}
\end{prob}

\begin{prob}[S]{546}
$\displaystyle \lim_{n\to \infty} \tan^n \left(\frac{\pi}{4}+\frac{1}{n}\right).$
\end{prob}

\begin{prob}[S]{547}
$\displaystyle \lim_{x\to 0} \frac{\mathrm{e}^{\alpha x} - \mathrm{e}^{\beta x}}{\sin\alpha x - \sin\beta x}.$
\end{prob}

\begin{prob}[S]{548}
$\displaystyle \lim_{x\to a} \frac{x^\alpha - a^\alpha}{x^\beta - a^\beta} \ (a > 0).$
\end{prob}

\begin{prob}[S]{549}
$\displaystyle \lim_{x\to b} \frac{a^x - a^b}{x-b} \ (a > 0).$
\end{prob}

\begin{prob}[S]{550}
$\displaystyle \lim_{h\to 0} \frac{a^{x+h} + a^{x-h} - 2a^x}{h^2} \ (a > 0).$
\end{prob}

\begin{prob}[S]{551}
$\displaystyle \lim_{x\to \infty} \frac{(x+a)^{x+a}(x+b)^{x+b}}{(x+a+b)^{2x+a+b}}.$
\end{prob}

\begin{prob}[S]{552}
$\displaystyle \lim_{n\to \infty} n(\sqrt[n]{x} - 1) \ (x > 0).$
\end{prob}

\begin{prob}[S]{553}
$\displaystyle \lim_{n\to \infty} n^2(\sqrt[n]{x} - \sqrt[n+1]{x}) \ (x > 0).$
\end{prob}

\begin{prob}[S]{554}
$\displaystyle \lim_{n\to \infty} \left(\frac{a - 1 + \sqrt[n]{b}}{a}\right)^n \ (a > 0, b > 0).$
\end{prob}

\begin{prob}[S]{555}
$\displaystyle \lim_{n\to \infty} \left(\frac{\sqrt[n]{a} + \sqrt[n]{b}}{2}\right)^n \ (a > 0, b > 0).$
\end{prob}

\begin{prob}[S]{556}
$\displaystyle \lim_{x\to 0} \left(\frac{a^x + b^x + c^x}{3}\right)^{\frac{1}{x}} \ (a > 0, b > 0, c > 0).$
\end{prob}

\begin{prob}[S]{557}
$\displaystyle \lim_{x\to 0} \left(\frac{a^{x+1} + b^{x+1} + c^{x+1}}{a+b+c}\right)^{\frac{1}{x}} \ (a > 0, b > 0, c > 0).$
\end{prob}

\begin{prob}[S]{558}
$\displaystyle \lim_{x\to 0} \left(\frac{a^{x^2} + b^{x^2}}{a^x + b^x}\right)^{\frac{1}{x}} \ (a > 0, b > 0).$
\end{prob}

\begin{prob}[S]{559}
$\displaystyle \lim_{x\to 0} \frac{a^{x^2} - b^{x^2}}{(a^x - b^x)^2} \ (a > 0, b > 0).$
\end{prob}

\begin{prob}[S]{560}
$\displaystyle \lim_{x\to a} \frac{a^{a^x} - a^{x^a}}{a^x - x^a} \ (a > 0).$
\end{prob}

\begin{prob}[M]{561}
(a) $\displaystyle \lim_{x\to -\infty} \frac{\ln(1+3^x)}{\ln(1+2^x)}$; \qquad (b) $\displaystyle \lim_{x\to +\infty} \frac{\ln(1+3^x)}{\ln(1+2^x)}.$
\end{prob}

\begin{prob}[S]{562}
$\displaystyle \lim_{x\to +\infty} \ln(1+2^x)\ln\left(1+\frac{3}{x}\right).$
\end{prob}

\begin{prob}[S]{563}
$\displaystyle \lim_{x\to 1} (1-x)\log_x 2.$
\end{prob}

\begin{prob}[M]{564}
证明：
\[
\lim_{x\to +\infty} \frac{x^n}{a^x} = 0 \quad (a > 1, n > 0).
\]
\end{prob}

\begin{prob}[M]{565}
证明：
\[
\lim_{x\to +\infty} \frac{\log_a x}{x^\varepsilon} = 0 \quad (a > 1, \varepsilon > 0).
\]
\end{prob}

\probgroup{求下列极限：}

\begin{prob}[M]{566}
(a) $\displaystyle \lim_{x\to 0} \frac{\ln(x^2+\mathrm{e}^x)}{\ln(x^4+\mathrm{e}^{2x})}$; \qquad (b) $\displaystyle \lim_{x\to +\infty} \frac{\ln(x^2+\mathrm{e}^x)}{\ln(x^4+\mathrm{e}^{2x})}.$
\end{prob}

\begin{prob}[S]{567}
$\displaystyle \lim_{x\to 0} \frac{\ln(1+x\mathrm{e}^x)}{\ln(x+\sqrt{1+x^2})}.$
\end{prob}

\begin{prob}[S]{568}
$\displaystyle \lim_{x\to +\infty} [(x+2)\ln(x+2) - 2(x+1)\ln(x+1) + x\ln x].$
\end{prob}

\begin{prob}[S]{569}
$\displaystyle \lim_{x\to +0} \left[\ln(x\ln a)\cdot \ln\left(\frac{\ln ax}{\ln\frac{x}{a}}\right)\right] \ (a > 1).$
\end{prob}

\begin{prob}[S]{570}
$\displaystyle \lim_{x\to +\infty} \left(\ln\frac{x+\sqrt{x^2+1}}{x+\sqrt{x^2-1}}\cdot \ln^{-2}\frac{x+1}{x-1}\right).$
\end{prob}

\begin{prob}[S]{571}
$\displaystyle \lim_{x\to 0} \frac{\sqrt{1+x\sin x} - 1}{\mathrm{e}^{x^2} - 1}.$
\end{prob}

\begin{prob}[S]{572}
$\displaystyle \lim_{x\to 0} \frac{\cos(x\mathrm{e}^x) - \cos(x\mathrm{e}^{-x})}{x^3}.$
\end{prob}

\begin{prob}[S]{573}
$\displaystyle \lim_{x\to 0} \left(2\mathrm{e}^{\frac{x}{x+1}} - 1\right)^{\frac{x^2+1}{x}}.$
\end{prob}

\begin{prob}[S]{574}
$\displaystyle \lim_{x\to 1} (2-x)^{\sec\frac{\pi x}{2}}.$
\end{prob}

\begin{prob}[S]{575}
$\displaystyle \lim_{x\to \frac{\pi}{2}} \frac{1 - \sin^{\alpha+\beta} x}{\sqrt{(1-\sin^\alpha x)(1-\sin^\beta x)}} \ (\alpha > 0, \beta > 0).$
\end{prob}

\begin{prob}[L]{576}
\begin{enumerate}[label=(\alph*), itemsep=2pt, topsep=3pt]
  \item $\displaystyle \lim_{x\to 0} \frac{\sinh x}{x}$;
  \item $\displaystyle \lim_{x\to 0} \frac{\cosh x - 1}{x^2}$;
  \item $\displaystyle \lim_{x\to 0} \frac{\tanh x}{x}$ （参考习题 340）。
\end{enumerate}
\end{prob}

\begin{prob}[M]{577}
\begin{enumerate}[label=(\alph*), itemsep=2pt, topsep=3pt]
  \item $\displaystyle \lim_{x\to 0} \frac{\sinh^2 x}{\ln(\cosh 3x)}$ （参考习题 340）；
  \item $\displaystyle \lim_{x\to +\infty} \frac{\sinh\sqrt{x^2+x} - \sinh\sqrt{x^2-x}}{\cosh x}.$
\end{enumerate}
\end{prob}

\begin{prob}[L]{578}
\begin{enumerate}[label=(\alph*), itemsep=2pt, topsep=3pt]
  \item $\displaystyle \lim_{x\to a} \frac{\sinh x - \sinh a}{x-a}$;
  \item $\displaystyle \lim_{x\to a} \frac{\cosh x - \cosh a}{x-a}$;
  \item $\displaystyle \lim_{x\to 0} \frac{\ln\cosh x}{\ln\cos x}.$
\end{enumerate}
\end{prob}

\begin{prob}[M]{579}
(a) $\displaystyle \lim_{x\to +\infty} (x - \ln\cosh x)$; \qquad (b) $\displaystyle \lim_{x\to 0} \frac{\mathrm{e}^{\sin 2x} - \mathrm{e}^{\sin x}}{\tanh x}.$
\end{prob}

\begin{prob}[S]{580}
$\displaystyle \lim_{n\to \infty} \left(\frac{\cosh\dfrac{\pi}{n}}{\cos\dfrac{\pi}{n}}\right)^{n^2}.$
\end{prob}

\begin{prob}[S]{581}
$\displaystyle \lim_{x\to \infty} \arcsin\frac{1-x}{1+x}.$
\end{prob}

\begin{prob}[S]{582}
$\displaystyle \lim_{x\to +\infty} \arccos(\sqrt{x^2+x} - x).$
\end{prob}

\begin{prob}[S]{583}
$\displaystyle \lim_{x\to 2} \arctan\frac{x-4}{(x-2)^2}.$
\end{prob}

\begin{prob}[S]{584}
$\displaystyle \lim_{x\to -\infty} \operatorname{arccot}\frac{x}{\sqrt{1+x^2}}.$
\end{prob}

\begin{prob}[S]{585}
$\displaystyle \lim_{h\to 0} \frac{\arctan(x+h) - \arctan x}{h}.$
\end{prob}

\begin{prob}[S]{586}
$\displaystyle \lim_{x\to 0} \frac{\ln\dfrac{1+x}{1-x}}{\arctan(1+x) - \arctan(1-x)}.$
\end{prob}

\begin{prob}[S]{587}
$\displaystyle \lim_{n\to \infty} \left[n\arctan\frac{1}{n(x^2+1)+x}\cdot \tan^n \left(\frac{\pi}{4}+\frac{x}{2n}\right)\right].$
\end{prob}

\begin{prob}[S]{588}
$\displaystyle \lim_{x\to \infty} x\left(\frac{\pi}{4} - \arctan\frac{x}{x+1}\right).$
\end{prob}

\begin{prob}[S]{589}
$\displaystyle \lim_{x\to +\infty} x\left(\frac{\pi}{2} - \arcsin\frac{x}{\sqrt{x^2+1}}\right).$
\end{prob}

\begin{prob}[S]{590}
$\displaystyle \lim_{n\to \infty} \left[1 + \frac{(-1)^n}{n}\right]^{\csc(\pi\sqrt{1+n^2})}.$
\end{prob}

\begin{prob}[M]{591}
(a) $\displaystyle \lim_{x\to 0} \frac{1}{x^{100}}\mathrm{e}^{-\frac{1}{x^2}}$; \qquad (b) $\displaystyle \lim_{x\to +0} x\ln x.$
\end{prob}

\begin{prob}[M]{592}
(a) $\displaystyle \lim_{x\to -\infty} (\sqrt{x^2+x} - x)$; \qquad (b) $\displaystyle \lim_{x\to +\infty} (\sqrt{x^2+x} - x).$
\end{prob}

\begin{prob}[M]{593}
(a) $\displaystyle \lim_{x\to -\infty} (\sqrt{1+x+x^2} - \sqrt{1-x+x^2})$; \qquad (b) $\displaystyle \lim_{x\to +\infty} (\sqrt{1+x+x^2} - \sqrt{1-x+x^2}).$
\end{prob}

\begin{prob}[M]{594}
设
\[
f(x) = \ln\frac{x+\sqrt{x^2+a^2}}{x+\sqrt{x^2+b^2}},
\]
求
\[
h = \lim_{x\to +\infty} f(x) - \lim_{x\to -\infty} f(x).
\]
\end{prob}

\begin{prob}[M]{595}
(a) $\displaystyle \lim_{x\to 1-0} \arctan\frac{1}{1-x}$; \qquad (b) $\displaystyle \lim_{x\to 1+0} \arctan\frac{1}{1-x}.$
\end{prob}

\begin{prob}[M]{596}
(a) $\displaystyle \lim_{x\to -0} \frac{1}{1+\mathrm{e}^{\frac{1}{x}}}$; \qquad (b) $\displaystyle \lim_{x\to +0} \frac{1}{1+\mathrm{e}^{\frac{1}{x}}}.$
\end{prob}

\begin{prob}[M]{597}
(a) $\displaystyle \lim_{x\to -\infty} \frac{\ln(1+\mathrm{e}^x)}{x}$; \qquad (b) $\displaystyle \lim_{x\to -\infty} \frac{\ln(1+\mathrm{e}^x)}{x}.$
\end{prob}

\begin{prob}[M]{598}
证明：
\begin{enumerate}[label=(\alph*), itemsep=2pt, topsep=3pt]
  \item 当 $x \to -\infty$ 时，$\dfrac{2x}{1+x} \to 2 + 0$；
  \item 当 $x \to +\infty$ 时，$\dfrac{2x}{1+x} \to 2 - 0$。
\end{enumerate}
\end{prob}

\begin{prob}[M]{599}
证明：
\begin{enumerate}[label=(\alph*), itemsep=2pt, topsep=3pt]
  \item 当 $x \to -0$ 时，$2^x \to 1 - 0$；
  \item 当 $x \to +0$ 时，$2^x \to 1 + 0$。
\end{enumerate}
\end{prob}

\begin{prob}[S]{600}
设 $f(x) = x + [x^2]$，求 $f(1),\, f(1-0),\, f(1+0)$。
\end{prob}

\begin{prob}[S]{601}
设 $f(x) = \operatorname{sgn}(\sin\pi x)$，求 $f(n),\, f(n-0),\, f(n+0) \ (n = 0, \pm 1, \cdots)$。
\end{prob}

\probgroup{求：}

\begin{prob}[S]{602}
$\displaystyle \lim_{x\to 0} x\sqrt{\cos\frac{1}{x}}.$
\end{prob}

\begin{prob}[S]{603}
$\displaystyle \lim_{x\to 0} x\left[\frac{1}{x}\right].$
\end{prob}

\begin{prob}[S]{604}
$\displaystyle \lim_{n\to \infty} \sin(\pi\sqrt{n^2+1}).$
\end{prob}

\begin{prob}[S]{605}
$\displaystyle \lim_{n\to \infty} \sin^2(\pi\sqrt{n^2+n}).$
\end{prob}

\begin{prob}[S]{606}
$\displaystyle \lim_{n\to \infty} \underbrace{\sin\sin\cdots\sin}_{n\text{ 次}} x.$
\end{prob}

\begin{prob}[L]{607}
设 $\lim_{x\to a} \varphi(x) = A$ 及 $\lim_{x\to A} \psi(x) = B$，则由此是否可推出
\[
\lim_{x\to a} \psi(\varphi(x)) = B?
\]
研究例子：
\[
\varphi(x) = \begin{cases}
\dfrac{1}{q}, & x = \dfrac{p}{q} \ (p, q \text{ 是互素的整数}), \\[8pt]
0, & x \text{ 为无理数},
\end{cases}
\qquad
\psi(x) = \begin{cases}
1, & x \neq 0, \\
0, & x = 0,
\end{cases}
\]
并且 $x \to 0$。
\end{prob}

\begin{prob}[L]{608}
证明柯西定理：若函数 $f(x)$ 定义于区间 $(a, +\infty)$，且在每一个有限的区间 $(a, b)$ 内是有界的，则
\begin{enumerate}[label=(\alph*), itemsep=2pt, topsep=3pt]
  \item $\displaystyle \lim_{x\to +\infty} \frac{f(x)}{x} = \lim_{x\to +\infty} [f(x+1) - f(x)]$；
  \item $\displaystyle \lim_{x\to +\infty} [f(x)]^{\frac{1}{x}} = \lim_{x\to +\infty} \frac{f(x+1)}{f(x)} \ (f(x) \geqslant C > 0)$，
\end{enumerate}
假定等式右端的极限都存在。
\end{prob}

\begin{prob}[M]{609}
证明：若 (1) 函数 $f(x)$ 定义于区域 $x > a$；(2) $f(x)$ 在每一个有限的区域 $a < x < b$ 内是有界的；(3) $\lim_{x\to +\infty} [f(x+1) - f(x)] = \pm\infty$，则
\[
\lim_{x\to +\infty} \frac{f(x)}{x} = \pm\infty.
\]
\end{prob}

\begin{prob}[M]{610}
证明：若 (1) 函数 $f(x)$ 定义于区域 $x > a$；(2) $f(x)$ 在每一个有限的区域 $a < x < b$ 内是有界的；(3) 存在着有限的或无穷的极限
\[
\lim_{x \to +\infty} \frac{f(x+1) - f(x)}{x^n} = l,
\]
则
\[
\lim_{x \to +\infty} \frac{f(x)}{x^{n+1}} = \frac{l}{n+1}.
\]
\end{prob}

\begin{prob}[M]{611}
证明：
\begin{enumerate}[label=(\alph*), itemsep=2pt, topsep=3pt]
\item $\lim_{n \to \infty} \left(1 + \frac{x}{n}\right)^n = \mathrm{e}^x$;
\item $\lim_{n \to \infty} \left(1 + x + \frac{x^2}{2!} + \cdots + \frac{x^n}{n!}\right) = \mathrm{e}^x$.
\end{enumerate}
\end{prob}

\begin{prob}[M]{612}
证明：
\[
\lim_{n \to \infty} n\sin(2\pi \mathrm{e} n!) = 2\pi.
\]
\probhint{利用 72 题的公式 $(*)$.}
\end{prob}

\probgroup{作下列函数的图像：}

\begin{prob}[M]{613}
(a) $y = 1 - x^{100}$; \qquad (b) $y = \lim_{n \to \infty} (1 - x^{2n})$ $(-1 \leqslant x \leqslant 1)$.
\end{prob}

\begin{prob}[M]{614}
(a) $y = \frac{x^{100}}{1 + x^{100}}$ $(x \geqslant 0)$; \qquad (b) $y = \lim_{n \to \infty} \frac{x^n}{1 + x^n}$ $(x \geqslant 0)$.
\end{prob}

\begin{prob}[M]{615}
$y = \lim_{n \to \infty} \frac{x^n - x^{-n}}{x^n + x^{-n}}$ $(x \neq 0)$.
\end{prob}

\begin{prob}[M]{616}
$y = \lim_{n \to \infty} \sqrt{x^2 + \frac{1}{n^2}}$.
\end{prob}

\begin{prob}[M]{617}
$y = \lim_{n \to \infty} \sqrt[n]{1 + x^n}$ $(x \geqslant 0)$.
\end{prob}

\begin{prob}[M]{618}
$y = \lim_{n \to \infty} \sqrt[n]{1 + x^n + \left(\frac{x^2}{2}\right)^n}$ $(x \geqslant 0)$.
\end{prob}

\begin{prob}[M]{619}
$y = \lim_{n \to \infty} \frac{x^{n+2}}{\sqrt{2^{2n} + x^{2n}}}$ $(x \geqslant 0)$.
\end{prob}

\begin{prob}[M]{620}
(a) $y = \sin^{1000} x$; \qquad (b) $y = \lim_{n \to \infty} \sin^{2n} x$.
\end{prob}

\begin{prob}[M]{621}
$y = \lim_{n \to \infty} \frac{\ln(2^n + x^n)}{n}$ $(x \geqslant 0)$.
\end{prob}

\begin{prob}[M]{622}
$y = \lim_{n \to \infty} (x - 1)\arctan x^n$.
\end{prob}

\begin{prob}[M]{623}
$y = \lim_{n \to \infty} \sqrt[n]{1 + \mathrm{e}^{n(x+1)}}$.
\end{prob}

\begin{prob}[M]{624}
(a) $y = \lim_{t \to +\infty} \frac{x + \mathrm{e}^{tx}}{1 + \mathrm{e}^{tx}}$; \qquad (b) $y = \lim_{t \to x} \frac{1}{t - x}\ln\frac{t}{x}$ $(x > 0)$.
\end{prob}

\begin{prob}[L]{625}
\begin{enumerate}[label=(\alph*), itemsep=2pt, topsep=3pt]
\item $y = \lim_{n \to \infty} \frac{x\tan^{2n}\frac{\pi x}{4} + \sqrt{x}}{\tan^{2n}\frac{\pi x}{4} + 1}$ $(x \geqslant 0)$;
\item $y = \lim_{n \to \infty} x\operatorname{sgn}|\sin^2(n!\pi x)|$;
\item 作曲线 $\lim_{n \to \infty} \sqrt[n]{|x|^n + |y|^n} = 1$.
\end{enumerate}
\end{prob}

\begin{prob}[M]{626}
若有
\[
\lim_{x \to \infty} [f(x) - (kx + b)] = 0,
\]
则直线 $y = kx + b$ 称为曲线 $y = f(x)$ 的（斜）渐近线。利用此方程推出渐近线存在的充分必要条件。
\end{prob}

\begin{prob}[L]{627}
求下列曲线的渐近线并作其图像：
\begin{enumerate}[label=(\alph*), itemsep=2pt, topsep=3pt]
\item $y = \frac{x^3}{x^2 + x - 2}$;
\item $y = \sqrt{x^2 + x}$;
\item $y = \sqrt[3]{x^2 - x^3}$;
\item $y = \frac{x\mathrm{e}^x}{\mathrm{e}^x - 1}$;
\item $y = \ln(1 + \mathrm{e}^x)$;
\item $y = x + \arccos\frac{1}{x}$.
\end{enumerate}
\end{prob}

\probgroup{求下列极限：}

\begin{prob}[M]{628}
$\lim_{n \to \infty} \left[\frac{x^{n+1}}{(n+1)!} + \frac{x^{n+2}}{(n+2)!} + \cdots + \frac{x^{2n}}{(2n)!}\right]$.
\end{prob}

\begin{prob}[M]{629}
$\lim_{n \to \infty} [(1+x)(1+x^2)(1+x^4)\cdots(1+x^{2n})]$，若 $|x| < 1$.
\end{prob}

\begin{prob}[M]{630}
$\lim_{n \to \infty} \left(\cos\frac{x}{2}\cos\frac{x}{4}\cdots\cos\frac{x}{2^n}\right)$.
\end{prob}

\begin{prob}[L]{631}
设
\[
\lim_{x \to 0} \frac{\varphi(x)}{\psi(x)} = 1,
\]
其中 $\psi(x) > 0$，再设 $n \to \infty$ 时 $\alpha_{mn} \rightrightarrows 0$ $(m = 1, 2, \cdots, n)$，换言之，对于任意 $\varepsilon > 0$，存在正整数 $N(\varepsilon)$，当 $m = 1, 2, \cdots, n$ 且 $n > N(\varepsilon)$ 时 $0 < |\alpha_{mn}| < \varepsilon$。

证明：
\[
\lim_{n \to \infty} [\varphi(\alpha_{1n}) + \varphi(\alpha_{2n}) + \cdots + \varphi(\alpha_{nn})] = \lim_{n \to \infty} [\psi(\alpha_{1n}) + \psi(\alpha_{2n}) + \cdots + \psi(\alpha_{nn})], \tag{1}
\]
此处假定等式 (1) 右端的极限存在。
\end{prob}

\probgroup{利用上述定理，求：}

\begin{prob}[M]{632}
$\lim_{n \to \infty} \sum_{k=1}^{n} \left(\sqrt[3]{1 + \frac{k}{n^2}} - 1\right)$.
\end{prob}

\begin{prob}[M]{633}
$\lim_{n \to \infty} \sum_{k=1}^{n} \left(\sin\frac{ka}{n^2}\right)$.
\end{prob}

\begin{prob}[M]{634}
$\lim_{n \to \infty} \sum_{k=1}^{n} \left(a^{\frac{k}{n^2}} - 1\right)$ $(a > 0)$.
\end{prob}

\begin{prob}[M]{635}
$\lim_{n \to \infty} \prod_{k=1}^{n} \left(1 + \frac{k}{n^2}\right)$.
\end{prob}

\begin{prob}[M]{636}
$\lim_{n \to \infty} \prod_{k=1}^{n} \cos\frac{ka}{n\sqrt{n}}$.
\end{prob}

\begin{prob}[M]{637.1}
数列 $x_n$ 由以下的等式所给定：
\[
x_1 = \sqrt{a}, \quad x_2 = \sqrt{a + \sqrt{a}}, \quad x_3 = \sqrt{a + \sqrt{a + \sqrt{a}}}, \quad \cdots \quad (a > 0).
\]
求 $\lim_{n \to \infty} x_n$.
\end{prob}

\begin{prob}[M]{637.2}
数列 $x_n$ 按下述方式给定：
\[
x_1 = 0, \quad x_2 = 1, \quad x_n = \frac{1}{2}(x_{n-1} + x_{n-2}) \quad (n = 3, 4, \cdots).
\]
求 $\lim_{n \to \infty} x_n$.
\end{prob}

\begin{prob}[M]{637.3}
数列 $y_n$ 借助于数列 $x_n$，由下述关系式定义：
\[
y_0 = x_0, \quad y_n = x_n - \alpha x_{n-1} \quad (n = 1, 2, \cdots),
\]
其中 $|\alpha| < 1$。若 $\lim_{n \to \infty} y_n = b$，求 $\lim_{n \to \infty} x_n$.
\end{prob}

\begin{prob}[M]{637.4}
数列 $x_n$ 按下述方式给定：
\[
x_0 = 1, \quad x_n = \frac{1}{1 + x_{n-1}} \quad (n = 1, 2, \cdots).
\]
求 $\lim_{n \to \infty} x_n$.
\probhint{考虑 $x_n$ 与方程 $x = \frac{1}{1+x}$ 的根之差.}
\end{prob}

\begin{prob}[M]{638}
函数序列
\[
y_n = y_n(x) \quad (0 \leqslant x \leqslant 1)
\]
用以下的方法来确定：
\begin{enumerate}[label=(\alph*), itemsep=2pt, topsep=3pt]
\item $y_1 = \frac{x}{2}, \; y_n = \frac{x}{2} - \frac{y_{n-1}^2}{2}$ $(n = 2, 3, \cdots)$;
\item $y_1 = \frac{x}{2}, \; y_n = \frac{x}{2} + \frac{y_{n-1}^2}{2}$ $(n = 2, 3, \cdots)$.
\end{enumerate}
求 $\lim_{n \to \infty} y_n$.
\end{prob}

\begin{prob}[M]{639.1}
设 $x > 0$ 且 $y_n = y_{n-1}(2 - x y_{n-1})$ $(n = 1, 2, \cdots)$。证明：若 $y_i > 0$ $(i = 0, 1)$，则序列 $y_n$ 收敛且
\[
\lim_{n \to \infty} y_n = \frac{1}{x}.
\]
\probhint{研究差 $\frac{1}{x} - y_n$.}
\end{prob}

\begin{prob}[M]{639.2}
为了求 $y = \sqrt{x}$，其中 $x > 0$，应用如下的步骤：$y_0 > 0$ 是任意的，
\[
y_n = \frac{1}{2}\left(y_{n-1} + \frac{x}{y_{n-1}}\right) \quad (n = 1, 2, \cdots).
\]
证明：
\[
\lim_{n \to \infty} y_n = \sqrt{x}.
\]
\probhint{利用公式
\[
\frac{y_n - \sqrt{x}}{y_n + \sqrt{x}} = \left(\frac{y_{n-1} - \sqrt{x}}{y_{n-1} + \sqrt{x}}\right)^2 \quad (n \geqslant 1).
\]}
\end{prob}

\begin{prob}[L]{640}
为了求开普勒方程
\[
x - \varepsilon \sin x = m \quad (0 < \varepsilon < 1) \tag{1}
\]
的近似解，假设
\[
x_0 = m, \quad x_1 = m + \varepsilon \sin x_0, \quad \cdots, \quad x_n = m + \varepsilon \sin x_{n-1}, \quad \cdots
\]
（逐步逼近法）。

证明：存在 $\xi = \lim_{n \to \infty} x_n$，且数 $\xi$ 为方程 (1) 的唯一的根。
\end{prob}

\begin{prob}[L]{641}
若 $\omega_h(f)$ 为函数 $f(x)$ 在区间 $|x - \xi| \leqslant h$ $(h > 0)$ 上的振幅，则数
\[
\omega_0(f) = \lim_{h \to 0} \omega_h(f)
\]
称为函数 $f(x)$ 在 $\xi$ 点的振幅。

若 $f(x) = 0$，而在 $x \neq 0$ 时 $f(x)$ 为下列函数，求 $f(x)$ 在点 $x = 0$ 的振幅：
\begin{enumerate}[label=(\alph*), itemsep=2pt, topsep=3pt]
\item $f(x) = \sin\frac{1}{x}$;
\item $f(x) = \frac{1}{x^2}\cos^2\frac{1}{x}$;
\item $f(x) = x\left(2 + \sin\frac{1}{x}\right)$;
\item $f(x) = \frac{1}{\pi}\arctan\frac{1}{x}$;
\item $f(x) = \frac{|\sin x|}{x}$;
\item $f(x) = \frac{1}{1 + \mathrm{e}^{\frac{1}{x}}}$;
\item $f(x) = (1 + |x|)^{\frac{1}{x}}$.
\end{enumerate}
\end{prob}

\begin{prob}[M]{642}
命
\[
f(x) = \sin\frac{1}{x},
\]
证明：对于满足条件 $-1 \leqslant \alpha \leqslant 1$ 的任何数 $\alpha$，可以选出数列 $x_n \to 0$ $(n = 1, 2, \cdots)$，使得
\[
\lim_{n \to \infty} f(x_n) = \alpha.
\]
\end{prob}

\begin{prob}[L]{643}
设
\begin{enumerate}[label=(\alph*), itemsep=2pt, topsep=3pt]
\item $f(x) = \sin^2\frac{1}{x} + \frac{2}{\pi}\arctan\frac{1}{x}$;
\item $f(x) = (2 - x^2)\cos\frac{1}{x}$;
\item $f(x) = \left(1 + \cos^2\frac{1}{x}\right)^{\sec^2\frac{1}{x}}$.
\end{enumerate}
求
\[
l = \varliminf_{x \to 0} f(x) \quad \text{和} \quad L = \varlimsup_{x \to 0} f(x).
\]
\end{prob}

\begin{prob}[L]{644}
设
\begin{enumerate}[label=(\alph*), itemsep=2pt, topsep=3pt]
\item $f(x) = \sin x$;
\item $f(x) = x^2\cos^2 x$;
\item $f(x) = 2^{\sin x^2}$;
\item $f(x) = \frac{x}{1 + x^2\sin^2 x}$ $(x \geqslant 0)$.
\end{enumerate}
求
\[
l = \varliminf_{x \to \infty} f(x) \quad \text{和} \quad L = \varlimsup_{x \to \infty} f(x).
\]
\end{prob}

\probsec{§6. 符号 $O$}

\begin{prob}[L]{645}
把圆心角 $AOB = x$（图 4）当作 $1$ 阶无穷小量，求下列各无穷小量的阶：
\probfig{fig_645.png}{图 4}
\begin{enumerate}[label=(\alph*), itemsep=2pt, topsep=3pt]
\item 弦 $AB$;
\item 拱高 $CD$;
\item 扇形 $AOB$ 的面积;
\item 三角形 $ABC$ 的面积;
\item 梯形 $ABB_1A_1$ 的面积;
\item 弓形 $ABC$ 的面积.
\end{enumerate}
\end{prob}

\begin{prob}[L]{646}
设 $o(f(x))$ 是当 $x \to a$ 时比函数 $f(x)$ 低阶增长的任意函数，并且 $O(f(x))$ 是当 $x \to a$ 时与函数 $f(x)$ 同阶增长的任意函数，其中 $f(x) > 0$。

证明：
\begin{enumerate}[label=(\alph*), itemsep=2pt, topsep=3pt]
\item $o(o(f(x))) = o(f(x))$;
\item $O(o(f(x))) = o(f(x))$;
\end{enumerate}
\begin{enumerate}[label=(\alph*), start=3, itemsep=2pt, topsep=3pt]
  \item $o(O(f(x))) = o(f(x))$;
  \item $O(O(f(x))) = O(f(x))$;
  \item $O(f(x)) + o(f(x)) = O(f(x))$.
\end{enumerate}
\end{prob}

\begin{prob}[L]{647}
设 $x \to 0$ 和 $n > 0$. 证明:
\begin{enumerate}[label=(\alph*), itemsep=2pt, topsep=3pt]
  \item $C O(x^n) = O(x^n)$ ($C \neq 0$ 为常数);
  \item $O(x^n) + O(x^m) = O(x^n)$ ($n < m$);
  \item $O(x^n) O(x^m) = O(x^{n+m})$.
\end{enumerate}
\end{prob}

\begin{prob}[L]{648}
设 $x \to +\infty, n > 0$. 证明:
\begin{enumerate}[label=(\alph*), itemsep=2pt, topsep=3pt]
  \item $C O(x^n) = O(x^n)$ ($C$ 为常数);
  \item $O(x^n) + O(x^m) = O(x^n)$ ($n > m$);
  \item $O(x^n) O(x^m) = O(x^{n+m})$.
\end{enumerate}
\end{prob}

\begin{prob}[L]{649}
证明符号 $\sim$ 具有如下性质:
\begin{enumerate}[label=(\arabic*), itemsep=2pt, topsep=3pt]
  \item 自反性: $\varphi(x) \sim \varphi(x)$;
  \item 对称性: 若 $\varphi(x) \sim \psi(x)$, 则 $\psi(x) \sim \varphi(x)$;
  \item 传递性: 若 $\varphi(x) \sim \psi(x)$ 及 $\psi(x) \sim \chi(x)$, 则 $\varphi(x) \sim \chi(x)$.
\end{enumerate}
\end{prob}

\begin{prob}[L]{650}
设 $x \to 0$, 证明下列等式:
\begin{enumerate}[label=(\alph*), itemsep=2pt, topsep=3pt]
  \item $2x - x^2 = O(x)$;
  \item $x\sin\sqrt{x} = O(x^{\frac{3}{2}})$;
  \item $x\sin\dfrac{1}{x} = O(|x|)$;
  \item $\ln x = o\left(\dfrac{1}{x^\varepsilon}\right)$ ($\varepsilon > 0$);
  \item $\sqrt{x + \sqrt{x + \sqrt{x}}} \sim \sqrt[8]{x}$;
  \item $\arctan\dfrac{1}{x} = O(1)$;
  \item $(1 + x)^n = 1 + nx + o(x)$.
\end{enumerate}
\end{prob}

\begin{prob}[L]{651}
设 $x \to +\infty$. 证明下列等式:
\begin{enumerate}[label=(\alph*), itemsep=2pt, topsep=3pt]
  \item $2x^3 - 3x^2 + 1 = O(x^3)$;
  \item $\dfrac{x + 1}{x^2 + 1} = O\left(\dfrac{1}{x}\right)$;
  \item $x + x^2\sin x = O(x^2)$;
  \item $\dfrac{\arctan x}{1 + x^2} = O\left(\dfrac{1}{x^2}\right)$;
  \item $\ln x = o(x^\varepsilon)$ ($\varepsilon > 0$);
  \item $x^p \mathrm{e}^{-x} = o\left(\dfrac{1}{x^2}\right)$;
  \item $\sqrt{x + \sqrt{x + \sqrt{x}}} \sim \sqrt{x}$;
  \item $x^2 + x\ln^{100} x \sim x^2$.
\end{enumerate}
\end{prob}

\begin{prob}[L]{652.1}
证明: 当 $x$ 充分大时, 下面的不等式成立:
\begin{enumerate}[label=(\alph*), itemsep=2pt, topsep=3pt]
  \item $x^2 + 10x + 100 < 0.001x^3$;
  \item $\ln^{1000} x < \sqrt{x}$;
  \item $x^{10}\mathrm{e}^x < \mathrm{e}^{2x}$.
\end{enumerate}
\end{prob}

\begin{prob}[M]{652.2}
证明渐近公式: 当 $x \to +\infty$ 时,
\[
  \sqrt{x^2 + px + q} = x + \frac{p}{2} + O\left(\frac{1}{x}\right).
\]
\end{prob}

\begin{prob}[L]{653}
设 $x \to 0$. 分出下列函数的形如 $Cx^n$ ($C$ 为常数) 的主部, 并求其对于无穷小变量 $x$ 的阶:
\begin{enumerate}[label=(\alph*), itemsep=2pt, topsep=3pt]
  \item $f(x) = 2x - 3x^3 + x^5$;
  \item $f(x) = \sqrt{1 + x} - \sqrt{1 - x}$;
  \item $f(x) = \sqrt{1 - 2x} - \sqrt[3]{1 - 3x}$;
  \item $f(x) = \tan x - \sin x$.
\end{enumerate}
\end{prob}

\begin{prob}[M]{654}
设 $x \to 0$. 证明: 无穷小量
\[
  \text{(a) } f(x) = \frac{1}{\ln x}; \qquad \text{(b) } f(x) = \mathrm{e}^{-\frac{1}{x^2}}
\]
无论对任何的 $n$, 也不能与无穷小量 $x^n (x > 0)$ 相比较. 即, 对于任何的 $n$, 等式
\[
  \lim_{x \to 0} \frac{f(x)}{x^n} = k
\]
都不成立, 式中 $k$ 为异于零的有限量.
\end{prob}

\begin{prob}[L]{655}
设 $x \to 1$. 分出下列函数的形如 $C(x - 1)^n$ 的主部, 并求其对于无穷小量 $x - 1$ 的阶:
\begin{enumerate}[label=(\alph*), itemsep=2pt, topsep=3pt]
  \item $f(x) = x^3 - 3x + 2$;
  \item $f(x) = \sqrt[3]{1 - \sqrt{x}}$;
  \item $f(x) = \ln x$;
  \item $f(x) = \mathrm{e}^x - \mathrm{e}$;
  \item $f(x) = x^x - 1$.
\end{enumerate}
\end{prob}

\begin{prob}[L]{656}
设 $x \to +\infty$. 分出下列函数的形如 $Cx^n$ 的主部, 并求其对于无穷大量 $x$ 的阶:
\begin{enumerate}[label=(\alph*), itemsep=2pt, topsep=3pt]
  \item $f(x) = x^2 + 100x + 10\,000$;
  \item $f(x) = \dfrac{2x^5}{x^3 - 3x + 1}$;
  \item $f(x) = \sqrt[3]{x^2 - x} + \sqrt{x}$;
  \item $f(x) = \sqrt{1 + \sqrt{1 + \sqrt{x}}}$.
\end{enumerate}
\end{prob}

\begin{prob}[L]{657}
设 $x \to +\infty$. 分出下列函数的形如 $C\left(\dfrac{1}{x}\right)^n$ 的主部, 并求其对于无穷小量 $\dfrac{1}{x}$ 的阶:
\begin{enumerate}[label=(\alph*), itemsep=2pt, topsep=3pt]
  \item $f(x) = \dfrac{x + 1}{x^4 + 1}$;
  \item $f(x) = \sqrt{x + 1} - \sqrt{x}$;
  \item $f(x) = \sqrt{x + 2} - 2\sqrt{x + 1} + \sqrt{x}$;
  \item $f(x) = \dfrac{1}{x}\sin\dfrac{1}{x}$.
\end{enumerate}
\end{prob}

\begin{prob}[L]{658}
设 $x \to 1$. 分出下列函数的形如 $C\left(\dfrac{1}{x - 1}\right)^n$ 的主部, 并求其对于无穷大量 $\dfrac{1}{x - 1}$ 的阶:
\begin{enumerate}[label=(\alph*), itemsep=2pt, topsep=3pt]
  \item $f(x) = \dfrac{x^2}{x^2 - 1}$;
  \item $f(x) = \sqrt{\dfrac{1 + x}{1 - x}}$;
  \item $f(x) = \dfrac{x}{\sqrt[3]{1 - x^3}}$;
  \item $f(x) = \dfrac{1}{\sin\pi x}$;
  \item $f(x) = \dfrac{\ln x}{(1 - x)^2}$.
\end{enumerate}
\end{prob}

\begin{prob}[M]{659}
设 $x \to +\infty$, $f_n(x) = x^n$ ($n = 1, 2, 3, \cdots$). 证明:
\begin{enumerate}[label=(\arabic*), itemsep=2pt, topsep=3pt]
  \item 每一个函数 $f_n(x)$ 都比其前面的一个函数 $f_{n-1}(x)$ 增加得快;
  \item 函数 $\mathrm{e}^x$ 比函数 $f_n(x)$ ($n = 1, 2, \cdots$) 中的每一个都增加得快.
\end{enumerate}
\end{prob}

\begin{prob}[M]{660}
设 $x \to +\infty$, $f_n(x) = \sqrt[n]{x}$ ($n = 1, 2, 3, \cdots$). 证明:
\begin{enumerate}[label=(\arabic*), itemsep=2pt, topsep=3pt]
  \item 每一个函数 $f_n(x)$ 都比其前面的一个函数 $f_{n-1}(x)$ 增加得慢;
  \item 函数 $f(x) = \ln x$ 比函数 $f_n(x)$ ($n = 1, 2, \cdots$) 中的每一个都增加得慢.
\end{enumerate}
\end{prob}

\begin{prob}[M]{661}
证明: 对于任意的函数序列
\[
  f_1(x), f_2(x), \cdots, f_n(x), \cdots \quad (x_0 < x < +\infty),
\]
可举出一函数 $f(x)$, 当 $x \to +\infty$ 时它比函数 $f_n(x)$ ($n = 1, 2, \cdots$) 中的每一个都增加得快.
\end{prob}

\probsec{§7. 函数的连续性}

\begin{prob}[M]{662}
已给连续函数 $y = f(x)$ 的图像, 对于给定点 $a$ 与给定数 $\varepsilon > 0$, 用几何方法表示出这样的数 $\delta > 0$, 使当 $|x - a| < \delta$ 时, $|f(x) - f(a)| < \varepsilon$.
\end{prob}

\begin{prob}[L]{663}
需要制造一块边长 $x_0 = 10\,\mathrm{cm}$ 的正方形金属板, 其面积 $y = x^2$ 与设计值 $y_0 = 100\,\mathrm{cm}^2$ 之差不可超过: (a) $\pm 1\,\mathrm{cm}^2$; (b) $\pm 0.1\,\mathrm{cm}^2$; (c) $\pm 0.01\,\mathrm{cm}^2$; (d) $\pm \varepsilon\,\mathrm{cm}^2$, 问其边长 $x$ 可以在什么范围内变化?
\end{prob}

\begin{prob}[L]{664}
立方体的边介于 $2\,\mathrm{m}$ 和 $3\,\mathrm{m}$ 之间, 在测量此立方体边长 $x$ 时容许怎样的绝对误差 $\Delta$, 方可使计算立方体体积时的绝对误差不超过: (a) $\varepsilon = 0.1\,\mathrm{m}^3$; (b) $\varepsilon = 0.01\,\mathrm{m}^3$; (c) $\varepsilon = 0.001\,\mathrm{m}^3$?
\end{prob}

\begin{prob}[M]{665}
问在 $x_0 = 100$ 的尽可能多大的邻域内, 函数 $y = \sqrt{x}$ 的图像的纵坐标与 $y_0 = 10$ 之差小于 $\varepsilon = 10^{-n}$ ($n \geqslant 0$)? 求当 $n = 0, 1, 2, 3$ 时这个邻域的大小.
\end{prob}

\begin{prob}[M]{666}
利用 “$\varepsilon$--$\delta$” 语言, 证明: 函数 $f(x) = x^2$ 当 $x = 5$ 时连续. 填下表:
\begin{center}
\begin{tabular}{c|c|c|c|c|c}
\hline
$\varepsilon$ & \quad $1$ \quad & \quad $0.1$ \quad & \quad $0.01$ \quad & \quad $0.001$ \quad & \quad $\cdots$ \quad \\
\hline
$\delta$ & & & & & \\
\hline
\end{tabular}
\end{center}
\end{prob}

\begin{prob}[M]{667}
设 $f(x) = \dfrac{1}{x}$, $\varepsilon = 0.001$. 对于数值 $x_0 = 0.1; 0.01; 0.001; \cdots$ 求出最大的正数 $\delta = \delta(\varepsilon, x_0)$, 使得可从不等式 $|x - x_0| < \delta$ 推出不等式 $|f(x) - f(x_0)| < \varepsilon$.

可否对于已知的 $\varepsilon = 0.001$ 选出这样的 $\delta > 0$, 使它对区间 $(0, 1)$ 中的一切 $x_0$ 值都适用? 换言之, 可否找到这样的 $\delta > 0$, 使得只要 $|x - x_0| < \delta$, 就有 $|f(x) - f(x_0)| < \varepsilon$, 而不论 $x_0 \in (0, 1)$ 为何值?
\end{prob}

\begin{prob}[M]{668}
用 “$\varepsilon$--$\delta$” 语言以肯定的方式表述以下论断: 函数 $f(x)$ 在点 $x_0$ 有定义, 而在这一点不连续.
\end{prob}

\begin{prob}[M]{669}
设对于某些数 $\varepsilon > 0$ 可找到对应的数 $\delta = \delta(\varepsilon, x_0) > 0$, 使得只要 $|x - x_0| < \delta$, 则 $|f(x) - f(x_0)| < \varepsilon$.

若 (a) 诸数 $\varepsilon$ 组成一有限的集合; (b) 数 $\varepsilon$ 组成分数 $\varepsilon = \dfrac{1}{2^n}$ ($n = 1, 2, \cdots$) 的无穷集合, 则可否断定函数 $f(x)$ 在点 $x_0$ 连续?
\end{prob}

\begin{prob}[M]{670}
设已知函数
\[
  f(x) = x + 0.001[x].
\]
证明: 对于每一个 $\varepsilon > 0.001$, 可以选出 $\delta = \delta(\varepsilon, x) > 0$, 使得只要 $|x' - x| < \delta$, 则 $|f(x') - f(x)| < \varepsilon$. 然而, 当 $0 < \varepsilon \leqslant 0.001$ 时, 对于所有 $x$ 值都无法选出这样的 $\delta$.

这个函数的连续性在哪些点遭到破坏?
\end{prob}

\begin{prob}[M]{671}
设对于每一个充分小的数 $\delta > 0$, 都存在 $\varepsilon = \varepsilon(\delta, x_0) > 0$, 使得只要 $|x - x_0| < \delta$ 则不等式 $|f(x) - f(x_0)| < \varepsilon$ 成立. 由此是否可知函数 $f(x)$ 在 $x = x_0$ 连续? 这些不等式描述了函数 $f(x)$ 的什么性质?
\end{prob}

\begin{prob}[M]{672}
设对于每一个数 $\varepsilon > 0$, 都存在数 $\delta = \delta(\varepsilon, x_0) > 0$, 使得若 $|f(x) - f(x_0)| < \varepsilon$, 则 $|x - x_0| < \delta$. 由此是否可知函数 $f(x)$ 在 $x = x_0$ 连续? 这些不等式描述了函数的什么性质?
\end{prob}

\begin{prob}[M]{673}
设对于每一个数 $\delta > 0$, 都存在数 $\varepsilon = \varepsilon(\delta, x_0) > 0$, 使得若 $|f(x) - f(x_0)| < \varepsilon$, 则 $|x - x_0| < \delta$. 由此是否可知函数 $f(x)$ 在 $x = x_0$ 连续? 这些不等式描述了函数的什么性质?

研究例子:
\[
  f(x) = \begin{cases}
    \arctan x, & x\text{ 为有理数}, \\
    \pi - \arctan x, & x\text{ 为无理数}.
  \end{cases}
\]
\end{prob}

\begin{prob}[L]{674}
利用 “$\varepsilon$--$\delta$” 语言证明下列函数的连续性:
\begin{enumerate}[label=(\alph*), itemsep=2pt, topsep=3pt]
  \item $ax + b$;
  \item $x^2$;
  \item $x^3$;
  \item $\sqrt{x}$;
  \item $\sqrt[3]{x}$;
  \item $\sin x$;
  \item $\cos x$;
  \item $\arctan x$.
\end{enumerate}
\end{prob}

\probgroup{研究下列函数的连续性并绘出其图像:}

\begin{prob}[M]{675}
$f(x) = |x|$.
\end{prob}

\begin{prob}[M]{676}
$f(x) = \begin{cases} \dfrac{x^2 - 4}{x - 2}, & x \neq 2, \\ A, & x = 2. \end{cases}$
\end{prob}

\begin{prob}[M]{677}
$f(x) = \begin{cases} \dfrac{1}{(1 + x)^2}, & x \neq -1, \\ \text{任意值}, & x = -1. \end{cases}$
\end{prob}

\begin{prob}[M]{678}
(a) $f_1(x) = \begin{cases} \left|\dfrac{\sin x}{x}\right|, & x \neq 0, \\ 1, & x = 0; \end{cases}$ \qquad (b) $f_2(x) = \begin{cases} \dfrac{\sin x}{|x|}, & x \neq 0, \\ 1, & x = 0. \end{cases}$
\end{prob}

\begin{prob}[M]{679}
$f(x) = \begin{cases} \sin\dfrac{1}{x}, & x \neq 0, \\ \text{任意值}, & x = 0. \end{cases}$
\end{prob}

\begin{prob}[M]{680}
$f(x) = \begin{cases} x\sin\dfrac{1}{x}, & x \neq 0, \\ 0, & x = 0. \end{cases}$
\end{prob}

\begin{prob}[M]{681}
$f(x) = \begin{cases} \mathrm{e}^{-\frac{1}{x^2}}, & x \neq 0, \\ 0, & x = 0. \end{cases}$
\end{prob}

\begin{prob}[M]{682}
$f(x) = \begin{cases} \dfrac{1}{1 + \mathrm{e}^{\frac{1}{x - 1}}}, & x \neq 1, \\ \text{任意值}, & x = 1. \end{cases}$
\end{prob}

\begin{prob}[M]{683}
$f(x) = \begin{cases} x\ln x^2, & x \neq 0, \\ a, & x = 0. \end{cases}$
\end{prob}

\begin{prob}[M]{684}
$f(x) = \operatorname{sgn} x$.
\end{prob}

\begin{prob}[M]{685}
$f(x) = [x]$.
\end{prob}

\begin{prob}[M]{686}
$f(x) = \sqrt{x} - [\sqrt{x}]$.
\end{prob}

\probgroup{求出下列函数的间断点并研究这些点的性质:}

\begin{prob}[S]{687}
$y = \dfrac{x}{(1 + x)^2}$.
\end{prob}

\begin{prob}[S]{688}
$y = \dfrac{1 + x}{1 + x^3}$.
\end{prob}

\begin{prob}[S]{689}
$y = \dfrac{x^2 - 1}{x^3 - 3x + 2}$.
\end{prob}

\begin{prob}[S]{690}
$y = \dfrac{\dfrac{1}{x} - \dfrac{1}{x + 1}}{\dfrac{1}{x - 1} - \dfrac{1}{x}}$.
\end{prob}

\begin{prob}[S]{691}
$y = \dfrac{x}{\sin x}$.
\end{prob}

\begin{prob}[S]{692}
$y = \sqrt{\dfrac{1 - \cos \pi x}{4 - x^2}}$.
\end{prob}

\begin{prob}[S]{693}
$y = \cos^2 \dfrac{1}{x}$.
\end{prob}

\begin{prob}[S]{694}
$y = \operatorname{sgn}\left(\sin \dfrac{\pi}{x}\right)$.
\end{prob}

\begin{prob}[S]{695}
$y = \dfrac{\cos \dfrac{\pi}{x}}{\cos \dfrac{\pi}{x}}$.
\end{prob}

\begin{prob}[S]{696}
$y = \arctan \dfrac{1}{x}$.
\end{prob}

\begin{prob}[S]{697}
$y = \sqrt{x} \arctan \dfrac{1}{x}$.
\end{prob}

\begin{prob}[S]{698}
$y = \mathrm{e}^{x + \frac{1}{x}}$.
\end{prob}

\begin{prob}[S]{699}
$y = \dfrac{1}{\ln x}$.
\end{prob}

\begin{prob}[S]{700}
$y = \dfrac{1}{1 - \mathrm{e}^{\frac{x}{1 - x}}}$.
\end{prob}

\probgroup{研究下列函数的连续性并绘出其大略图像:}

\begin{prob}[S]{701}
$y = \operatorname{sgn}(\sin x)$.
\end{prob}

\begin{prob}[S]{702}
$y = x - [x]$.
\end{prob}

\begin{prob}[S]{703}
$y = x[x]$.
\end{prob}

\begin{prob}[S]{704}
$y = [x] \sin \pi x$.
\end{prob}

\begin{prob}[S]{705}
$y = x^2 - [x^2]$.
\end{prob}

\begin{prob}[S]{706}
$y = \left[\dfrac{1}{x}\right]$.
\end{prob}

\begin{prob}[S]{707}
$y = x\left[\dfrac{1}{x}\right]$.
\end{prob}

\begin{prob}[S]{708}
$y = \operatorname{sgn}\left(\cos \dfrac{1}{x}\right)$.
\end{prob}

\begin{prob}[S]{709}
$y = \left[\dfrac{1}{x^2}\right] \operatorname{sgn}\left(\sin \dfrac{\pi}{x}\right)$.
\end{prob}

\begin{prob}[S]{710}
$y = \cot \dfrac{\pi}{x}$.
\end{prob}

\begin{prob}[S]{711}
$y = \sec^2 \dfrac{1}{x}$.
\end{prob}

\begin{prob}[S]{712}
$y = (-1)^{[x^2]}$.
\end{prob}

\begin{prob}[S]{713}
$y = \arctan\left(\dfrac{1}{x} + \dfrac{1}{x - 1} + \dfrac{1}{x - 2}\right)$.
\end{prob}

\begin{prob}[S]{714}
$y = \dfrac{1}{x^2 \sin^2 x}$.
\end{prob}

\begin{prob}[S]{715}
$y = \dfrac{1}{\sin(x^2)}$.
\end{prob}

\begin{prob}[S]{716}
$y = \ln \dfrac{x^2}{(x + 1)(x - 3)}$.
\end{prob}

\begin{prob}[S]{717}
$y = \mathrm{e}^{-\frac{1}{x}}$.
\end{prob}

\begin{prob}[S]{718}
$y = 1 - \mathrm{e}^{-\frac{1}{x^2}}$.
\end{prob}

\begin{prob}[S]{719}
$y = \tanh \dfrac{2x}{1 - x^2}$.
\end{prob}

\probgroup{研究下列函数的连续性并作出其图像:}

\begin{prob}[M]{720}
$y = \lim\limits_{n \to \infty} \dfrac{1}{1 + x^n} \ (x \geqslant 0)$.
\end{prob}

\begin{prob}[M]{721}
$y = \lim\limits_{n \to \infty} \dfrac{n^x - n^{-x}}{n^x + n^{-x}}$.
\end{prob}

\begin{prob}[M]{722}
$y = \lim\limits_{n \to \infty} \sqrt[n]{1 + x^{2n}}$.
\end{prob}

\begin{prob}[M]{723}
$y = \lim\limits_{n \to \infty} \cos^{2n} x$.
\end{prob}

\begin{prob}[M]{724}
$y = \lim\limits_{n \to \infty} \dfrac{x}{1 + (2\sin x)^{2n}}$.
\end{prob}

\begin{prob}[M]{725}
$y = \lim\limits_{n \to \infty} [x \arctan(n \cot x)]$.
\end{prob}

\begin{prob}[M]{726}
$y = \lim\limits_{n \to \infty} \dfrac{x + x^2 \mathrm{e}^{nx}}{1 + \mathrm{e}^{nx}}$.
\end{prob}

\begin{prob}[M]{727}
$y = \lim\limits_{t \to +\infty} \dfrac{\ln(1 + \mathrm{e}^{xt})}{\ln(1 + \mathrm{e}^t)}$.
\end{prob}

\begin{prob}[M]{728}
$y = \lim\limits_{t \to +\infty} (1 + x)\tanh tx$.
\end{prob}

\begin{prob}[M]{729}
函数
\[
f(x) = \begin{cases}
2x, & 0 \leqslant x \leqslant 1, \\
2 - x, & 1 < x \leqslant 2
\end{cases}
\]
是否连续?
\end{prob}

\begin{prob}[M]{730}
设
\[
f(x) = \begin{cases}
\mathrm{e}^x, & x < 0, \\
a + x, & x \geqslant 0.
\end{cases}
\]
怎样选择数 $a$, 可使函数 $f(x)$ 连续?
\end{prob}

\begin{prob}[L]{731}
研究下列函数的连续性并说明间断点的性质:
\begin{enumerate}[label=(\alph*), itemsep=2pt, topsep=3pt]
\item $f(x) = \begin{cases} x^2, & 0 \leqslant x \leqslant 1, \\ 2 - x, & 1 < x \leqslant 2; \end{cases}$
\item $f(x) = \begin{cases} x, & |x| \leqslant 1, \\ 1, & |x| > 1; \end{cases}$
\item $f(x) = \begin{cases} \cos\dfrac{\pi x}{2}, & |x| \leqslant 1, \\ |x - 1|, & |x| > 1; \end{cases}$
\item $f(x) = \begin{cases} \cot^2 \pi x, & x \text{ 非整数}, \\ 0, & x \text{ 为整数}; \end{cases}$
\item $f(x) = \begin{cases} \sin \pi x, & x \text{ 为有理数}, \\ 0, & x \text{ 为无理数}. \end{cases}$
\end{enumerate}
\end{prob}

\begin{prob}[M]{732}
函数 $d = d(x)$ 是数轴 $Ox$ 上的点 $x$ 与由线段 $0 \leqslant x \leqslant 1$ 及 $2 \leqslant x \leqslant 3$ 所构成点集间的最短距离. 求函数 $d$ 的解析表达式, 作出其图像并研究其连续性.
\end{prob}

\begin{prob}[L]{733}
图像 $E$ 由高为 $1$ 底为 $1$ 的等腰三角形及底为 $1$ 高为 $2$ 及 $3$ 的二矩形构成 (图 5). 函数 $S = S(y) \ (0 \leqslant y < +\infty)$ 是图像 $E$ 介于平行线 $Y = 0$ 及 $Y = y$ 之间的那一部分面积, 而函数 $b = b(y) \ (0 \leqslant y < +\infty)$ 是用平行线 $Y = y$ 去截图像所得截痕之长. 求函数 $S$ 及 $b$ 的解析表达式, 作出它们的图像并研究其连续性.
\probfig{fig_733.png}{图 5}
\end{prob}

\begin{prob}[M]{734}
证明: 狄利克雷函数
\[
\chi(x) = \lim_{m \to \infty} \left\{\lim_{n \to \infty} \cos^n(\pi m! x)\right\}
\]
对于每个 $x$ 值都是不连续的.
\end{prob}

\begin{prob}[M]{735}
设有函数
\[
f(x) = x\chi(x),
\]
式中 $\chi(x)$ 为狄利克雷函数 (参阅上题), 研究函数 $f(x)$ 的连续性并作出它的草图.
\end{prob}

\begin{prob}[L]{736}
证明: 黎曼函数
\[
f(x) = \begin{cases}
\dfrac{1}{n}, & x = \dfrac{m}{n}, \text{ 其中 } m \text{ 和 } n \text{ 为互素的数}, \\
0, & x \text{ 为无理数},
\end{cases}
\]
当 $x$ 取任一个有理值时是不连续的, 而当 $x$ 取任一个无理值时是连续的. 作出这个函数的草图.
\end{prob}

\begin{prob}[M]{737}
研究函数
\[
f(x) = \begin{cases}
\dfrac{nx}{n + 1}, & x \text{ 为既约的有理分数 } \dfrac{m}{n} \ (n \geqslant 1), \\
|x|, & x \text{ 为无理数}
\end{cases}
\]
的连续性并作出此函数的草图.
\end{prob}

\begin{prob}[S]{738}
函数 $f(x) = \dfrac{1 - \cos x}{x^2}$ 当自变量 $x$ 取一切非零值时都有定义. 为了使此函数在 $x = 0$ 时连续, 应当以什么数值作为函数在点 $x = 0$ 的值?
\end{prob}

\begin{prob}[M]{739}
证明: 无论怎样选取数 $f(1)$, 函数 $f(x) = \dfrac{1}{1 - x}$ 在 $x = 1$ 时都是不连续的.
\end{prob}

\begin{prob}[L]{740}
下列函数 $f(x)$ 当 $x = 0$ 时没有意义. 定义 $f(0)$ 的值, 使得 $f(x)$ 在 $x = 0$ 连续:
\begin{enumerate}[label=(\alph*), itemsep=2pt, topsep=3pt]
\item $f(x) = \dfrac{\sqrt{1 + x} - 1}{\sqrt[3]{1 + x} - 1}$;
\item $f(x) = \dfrac{\tan 2x}{x}$;
\item $f(x) = \sin x \sin \dfrac{1}{x}$;
\item $f(x) = (1 + x)^{\frac{1}{x}}$;
\item $f(x) = \dfrac{1}{x^2} \mathrm{e}^{-\frac{1}{x^2}}$;
\item $f(x) = x^x \ (x > 0)$;
\item $f(x) = x \ln^2 x$.
\end{enumerate}
\end{prob}

\begin{prob}[M]{741}
设: (a) 函数 $f(x)$ 当 $x = x_0$ 时是连续的, 而函数 $g(x)$ 当 $x = x_0$ 时是不连续的; (b) 当 $x = x_0$ 时函数 $f(x)$ 和 $g(x)$ 二者都是不连续的, 则此二函数的和 $f(x) + g(x)$ 在已知点 $x_0$ 是否必为不连续的? 举出适当的例子.
\end{prob}

\begin{prob}[M]{742}
设: (a) 函数 $f(x)$ 于 $x_0$ 连续, 而函数 $g(x)$ 于 $x_0$ 不连续; (b) 当 $x = x_0$ 时, 函数 $f(x)$ 和 $g(x)$ 二者都不连续, 则此二函数的乘积 $f(x)g(x)$ 在已知点 $x_0$ 是否必不连续? 举出适当的例子.
\end{prob}

\begin{prob}[M]{743}
可否断定不连续函数的平方仍为不连续函数? 举出处处不连续但其平方连续的函数的例子.
\end{prob}

\begin{prob}[L]{744}
研究函数 $f[g(x)]$ 和 $g[f(x)]$ 的连续性:
\begin{enumerate}[label=(\alph*), itemsep=2pt, topsep=3pt]
\item $f(x) = \operatorname{sgn} x$ 及 $g(x) = 1 + x^2$;
\item $f(x) = \operatorname{sgn} x$ 及 $g(x) = x(1 - x^2)$;
\item $f(x) = \operatorname{sgn} x$ 及 $g(x) = 1 + x - [x]$.
\end{enumerate}
\end{prob}

\begin{prob}[M]{745}
设
\[
f(u) = \begin{cases}
u, & 0 < u \leqslant 1, \\
2 - u, & 1 < u < 2,
\end{cases}
\]
\[
\varphi(x) = \begin{cases}
x, & x \text{ 为有理数}, \\
2 - x, & x \text{ 为无理数},
\end{cases}
\]
其中 $0 < x < 1$. 研究复合函数 $y = f(u)$ 的连续性, 其中 $u = \varphi(x)$.
\end{prob}

\begin{prob}[M]{746}
证明: 若 $f(x)$ 是连续函数, 则下列函数也是连续的:
\[
F(x) = |f(x)|.
\]
\end{prob}

\begin{prob}[M]{747}
证明: 若函数 $f(x)$ 是连续的, 则函数
\[
f_c(x) = \begin{cases}
-c, & f(x) < -c, \\
f(x), & |f(x)| \leqslant c, \\
c, & f(x) > c
\end{cases}
\]
(式中 $c$ 为任意的正数) 也是连续函数.
\end{prob}

\begin{prob}[M]{748}
证明: 若函数 $f(x)$ 在闭区间 $[a, b]$ 上连续, 则函数
\[
m(x) = \inf_{a \leqslant \xi \leqslant x} \{f(\xi)\} \quad \text{和} \quad M(x) = \sup_{a \leqslant \xi \leqslant x} \{f(\xi)\}
\]
在 $[a, b]$ 上也是连续的.
\end{prob}

\begin{prob}[M]{749}
证明: 若函数 $f(x)$ 和 $g(x)$ 是连续的, 则函数
\[
\varphi(x) = \min[f(x), g(x)] \quad \text{和} \quad \psi(x) = \max[f(x), g(x)]
\]
也是连续的.
\end{prob}

\begin{prob}[M]{750}
设函数 $f(x)$ 在闭区间 $[a, b]$ 上有定义并有界. 证明: 函数
\[
m(x) = \inf_{a \leqslant \xi < x} \{f(\xi)\} \quad \text{和} \quad M(x) = \sup_{a \leqslant \xi < x} \{f(\xi)\}
\]
在闭区间 $[a, b]$ 上是左连续的.
\end{prob}

\begin{prob}[M]{751}
证明: 若函数 $f(x)$ 在区间 $a \leqslant x < +\infty$ 上连续, 且有有限的
\[
\lim_{x \to +\infty} f(x),
\]
则此函数在已知区间上是有界的.
\end{prob}

\begin{prob}[M]{752}
设函数 $f(x)$ 在区间 $(x_0, +\infty)$ 上连续并有界. 证明: 对于任何数 $T$, 可求得数列 $x_n \to +\infty$, 使
\[
\lim_{n \to \infty} [f(x_n + T) - f(x_n)] = 0.
\]
\end{prob}

\begin{prob}[M]{753}
设 $\varphi(x)$ 及 $\psi(x)$ 为定义于 $-\infty < x < +\infty$ 的连续周期函数, 且
\[
\lim_{x \to +\infty} [\varphi(x) - \psi(x)] = 0.
\]
证明: $\varphi(x) \equiv \psi(x)$.
\end{prob}

\begin{prob}[M]{754}
证明: 单调有界函数的一切间断点皆为第一类间断点.
\end{prob}

\begin{prob}[M]{755}
证明: 若函数 $f(x)$ 具有下列性质: (1) 在闭区间 $[a, b]$ 上有定义且单调; (2) 取介于 $f(a)$ 和 $f(b)$ 之间所有的数作为其函数值, 则此函数在 $[a, b]$ 上连续.
\end{prob}

\begin{prob}[M]{756}
证明: 函数
\[
f(x) = \begin{cases}
\sin \dfrac{1}{x - a}, & x \neq a, \\
0, & x = a
\end{cases}
\]
在任意闭区间 $[a, b]$ 上取介于 $f(a)$ 和 $f(b)$ 之间的一切中间值, 但在 $[a, b]$ 上并不连续.
\end{prob}

\begin{prob}[M]{757}
证明: 若函数 $f(x)$ 在区间 $(a, b)$ 内连续, 且 $x_1, x_2, \cdots, x_n$ 为此区间中的任意值, 则在它们之间可找到一个数值 $\xi$, 使得
\[
f(\xi) = \frac{1}{n}[f(x_1) + f(x_2) + \cdots + f(x_n)].
\]
\end{prob}

\begin{prob}[M]{758}
设函数 $f(x)$ 在区间 $(a, b)$ 上连续, 且
\[
l = \varliminf_{x \to a} f(x), \quad L = \varlimsup_{x \to a} f(x).
\]
证明: 对于满足 $l \leqslant \lambda \leqslant L$ 的任意的数 $\lambda$, 存在数列 $x_n \to a \ (n = 1, 2, \cdots)$, 使得
\[
\lim_{n \to \infty} f(x_n) = \lambda.
\]
\end{prob}

\probsec{§8. 反函数. 用参数形式表示的函数}

\begin{prob}[M]{759}
求分式线性函数
\[
y = \frac{ax + b}{cx + d} \quad (ad - bc \neq 0)
\]
的反函数. 在怎样的情形下, 反函数与已知函数相同?
\end{prob}

\begin{prob}[S]{760}
设 $y = x + [x]$, 求反函数 $x = x(y)$.
\end{prob}

\begin{prob}[M]{761}
证明: 存在唯一的连续函数 $y = y(x) \ (-\infty < x < +\infty)$ 满足开普勒方程
\[
y - \varepsilon \sin y = x \quad (0 \leqslant \varepsilon < 1).
\]
\end{prob}

\begin{prob}[M]{762}
证明: 对于每一个实数 $k \ (-\infty < k < +\infty)$ 方程
\[
\cot x = kx
\]
在区间 $0 < x < \pi$ 中有唯一连续的根 $x = x(k)$.
\end{prob}

\begin{prob}[M]{763}
非单调的函数 $y = f(x) \ (-\infty < x < +\infty)$ 可否有单值的反函数? 研究例子:
\[
y = \begin{cases}
x, & x \text{ 为有理数}, \\
-x, & x \text{ 为无理数}.
\end{cases}
\]
\end{prob}

\begin{prob}[S]{764}
在什么情况下, 函数 $y = f(x)$ 和反函数 $x = f^{-1}(y)$ 是同一个函数?
\end{prob}

\begin{prob}[M]{765}
证明: 不连续函数
\[
y = (1 + x^2)\operatorname{sgn} x
\]
的反函数是连续函数.
\end{prob}

\begin{prob}[M]{766}
证明: 若函数 $f(x)$ 在闭区间 $[a, b]$ 上有定义并且是严格单调的, 且
\[
\lim_{n \to \infty} f(x_n) = f(a) \quad (a \leqslant x_n \leqslant b),
\]
则
\[
\lim_{n \to \infty} x_n = a.
\]
\end{prob}

\probgroup{求下列函数的反函数的连续单值分支：}

\begin{prob}[S]{767}
$y = x^2.$
\end{prob}

\begin{prob}[S]{768}
$y = 2x - x^2.$
\end{prob}

\begin{prob}[S]{769}
$y = \frac{2x}{1 + x^2}.$
\end{prob}

\begin{prob}[S]{770}
$y = \sin x.$
\end{prob}

\begin{prob}[S]{771}
$y = \cos x.$
\end{prob}

\begin{prob}[S]{772}
$y = \tan x.$
\end{prob}

\begin{prob}[M]{773}
证明：连续函数 $y = 1 + \sin x$ 在区间 $0 < x < 2\pi$ 内的值域是闭区间.
\end{prob}

\begin{prob}[M]{774}
证明：等式
\[
\arcsin x + \arccos x = \frac{\pi}{2}.
\]
\end{prob}

\begin{prob}[M]{775}
证明：等式
\[
\arctan x + \arctan \frac{1}{x} = \frac{\pi}{2}\operatorname{sgn} x \quad (x \neq 0).
\]
\end{prob}

\begin{prob}[L]{776}
证明反正切加法定理：
\[
\arctan x + \arctan y = \arctan \frac{x + y}{1 - xy} + \varepsilon \pi,
\]
式中 $\varepsilon = \varepsilon(x, y)$ 为取 $0, 1, -1$ 这三个值的函数.

当给定 $x$ 的值时，对于怎样的 $y$ 值，函数 $\varepsilon$ 可能不连续？在 $Oxy$ 平面上作出函数 $\varepsilon$ 连续的相应区域，并求此函数在所得区域内的值.
\end{prob}

\begin{prob}[M]{777}
证明反正弦加法定理：
\[
\arcsin x + \arcsin y = (-1)^{\varepsilon} \arcsin \left( x\sqrt{1 - y^2} + y\sqrt{1 - x^2} \right) + \varepsilon \pi \quad (|x| \leqslant 1, |y| \leqslant 1),
\]
式中
\[
\varepsilon = \begin{cases}
0, & xy \leqslant 0 \text{ 或 } x^2 + y^2 \leqslant 1, \\
\operatorname{sgn} x, & xy > 0 \text{ 和 } x^2 + y^2 > 1.
\end{cases}
\]
\end{prob}

\begin{prob}[M]{778}
证明反余弦相加的定理：
\[
\arccos x + \arccos y = (-1)^{\varepsilon} \arccos \left( xy - \sqrt{1 - x^2}\sqrt{1 - y^2} \right) + 2\pi \varepsilon \quad (|x| \leqslant 1, |y| \leqslant 1),
\]
式中
\[
\varepsilon = \begin{cases}
0, & x + y \geqslant 0, \\
1, & x + y < 0.
\end{cases}
\]
\end{prob}

\begin{prob}[M]{779}
作函数的图像：
\begin{enumerate}[label=(\alph*), itemsep=2pt, topsep=3pt]
    \item $y = \arcsin x - \arcsin \sqrt{1 - x^2}$;
    \item $y = \arcsin \left(2x\sqrt{1 - x^2}\right) - 2\arcsin x.$
\end{enumerate}
\end{prob}

\begin{prob}[M]{780}
函数 $y = y(x)$ 由下面的方程给出：
\[
x = \arctan t, \quad y = \operatorname{arccot} t \quad (-\infty < t < +\infty),
\]
求此函数. 此函数的定义域是什么？
\end{prob}

\begin{prob}[M]{781}
设
\[
x = \cosh t, \quad y = \sinh t \quad (-\infty < t < +\infty).
\]
在参量 $t$ 的哪些变化域中可以把变量 $y$ 看作变量 $x$ 的单值函数？求 $y$ 在不同变化域上的表达式.
\end{prob}

\begin{prob}[M]{782}
设 $y$ 对 $x$ 的函数关系由方程组
\[
x = \varphi(t), \quad y = \psi(t) \quad (\alpha < t < \beta)
\]
给出，问 $y = y(x)$ 是单值函数的充分必要条件是什么？

研究例子：
\[
x = \sin^2 t, \quad y = \cos^2 t.
\]
\end{prob}

\begin{prob}[M]{783}
在怎样的条件下，两个方程组
\[
x = \varphi(t), \quad y = \psi(t) \quad (a < t < b)
\]
和
\[
x = \varphi(\chi(\tau)), \quad y = \psi(\chi(\tau)) \quad (\alpha < \tau < \beta)
\]
定义出同一个函数 $y = y(x)$？
\end{prob}

\begin{prob}[M]{784}
设函数 $\varphi(x)$ 和 $\psi(x)$ 在区间 $(a, b)$ 内有定义并且连续，且
\[
A = \inf_{a < x < b} \varphi(x), \quad B = \sup_{a < x < b} \varphi(x).
\]
在何种情况下存在单值函数 $f(x)$，它在区间 $(A, B)$ 上有定义，并且当 $a < x < b$ 时 $\psi(x) = f[\varphi(x)]$？
\end{prob}

\probsec{§9. 函数的一致连续性}

\begin{prob}[L]{785}
某工厂车间制造正方形平板，其边长 $x$ 在 $1\,\mathrm{cm}$ 至 $10\,\mathrm{cm}$ 之间取值. 当边长的允许误差 $\delta$ 为何值时，不论制造多大的平板（边长在上述范围内），其面积 $y$ 与设计值之差都小于 $\varepsilon$？考虑下列情形：
\begin{enumerate}[label=(\alph*), itemsep=2pt, topsep=3pt]
    \item $\varepsilon = 1\,\mathrm{cm}^2$;
    \item $\varepsilon = 0.01\,\mathrm{cm}^2$;
    \item $\varepsilon = 0.0001\,\mathrm{cm}^2.$
\end{enumerate}
\end{prob}

\begin{prob}[L]{786}
圆柱形套筒的宽为 $\varepsilon$，长为 $\delta$. 将套筒套在曲线 $y = \sqrt[3]{x}$ 上并使它沿此曲线滑动，但套筒的轴须保持平行于 $Ox$ 轴. $\delta$ 应取何值，才可以使此筒顺利地经过此曲线上由不等式 $-10 \leqslant x \leqslant 10$ 所限定的部分？考虑下列情形：
\begin{enumerate}[label=(\alph*), itemsep=2pt, topsep=3pt]
    \item $\varepsilon = 1$;
    \item $\varepsilon = 0.1$;
    \item $\varepsilon = 0.001$;
    \item $\varepsilon$ 为任意小数.
\end{enumerate}
\end{prob}

\begin{prob}[M]{787}
以“$\varepsilon\text{-}\delta$”语言用肯定的方式表达下述论断：函数 $f(x)$ 在某集合（开区间，闭区间，等等）上连续，但在此集合上并不一致连续.
\end{prob}

\begin{prob}[M]{788}
证明：函数 $f(x) = \frac{1}{x}$ 在区间 $(0, 1)$ 上是连续的，但在此区间上并非一致连续的.
\end{prob}

\begin{prob}[M]{789}
证明：函数 $f(x) = \sin \frac{\pi}{x}$ 在区间 $(0, 1)$ 上是连续的并且有界，但在此区间上并非一致连续的.
\end{prob}

\begin{prob}[M]{790}
证明：函数 $f(x) = \sin x^2$ 在无穷区间 $-\infty < x < +\infty$ 上连续并且有界，但在此区间上并非一致连续的.
\end{prob}

\begin{prob}[M]{791}
证明：若函数 $f(x)$ 在区域 $a \leqslant x < +\infty$ 上有定义并且连续，而且有限的
\[
\lim_{x \to +\infty} f(x)
\]
存在，则 $f(x)$ 在此区域上是一致连续的.
\end{prob}

\begin{prob}[M]{792}
证明：无界函数 $f(x) = x + \sin x$ 于全轴 $-\infty < x < +\infty$ 上一致连续.
\end{prob}

\begin{prob}[M]{793}
函数 $f(x) = x^2$ 在下列区间中是否为一致连续的：(a) $(-l, l)$，这里 $l$ 为任意大的正数；\qquad (b) 在区间 $(-\infty, +\infty)$ 上？
\end{prob}

\probgroup{研究下列函数在给定区域上的一致连续性：}

\begin{prob}[S]{794}
$f(x) = \frac{x}{4 - x^2} \quad (-1 \leqslant x \leqslant 1).$
\end{prob}

\begin{prob}[S]{795}
$f(x) = \ln x \quad (0 < x < 1).$
\end{prob}

\begin{prob}[S]{796}
$f(x) = \frac{\sin x}{x} \quad (0 < x < \pi).$
\end{prob}

\begin{prob}[S]{797}
$f(x) = \mathrm{e}^x \cos \frac{1}{x} \quad (0 < x < 1).$
\end{prob}

\begin{prob}[S]{798}
$f(x) = \arctan x \quad (-\infty < x < +\infty).$
\end{prob}

\begin{prob}[S]{799}
$f(x) = \sqrt{x} \quad (1 \leqslant x < +\infty).$
\end{prob}

\begin{prob}[S]{800}
$f(x) = x \sin x \quad (0 \leqslant x < +\infty).$
\end{prob}

\begin{prob}[M]{801.1}
证明：函数 $f(x) = \frac{|\sin x|}{x}$ 在每个区间
\[
J_1 = (-1 < x < 0) \quad \text{及} \quad J_2 = (0 < x < 1)
\]
上是一致连续的，但在它们的和
\[
J_1 + J_2 = \{0 < |x| < 1\}
\]
上并非一致连续的.
\end{prob}

\begin{prob}[M]{801.2}
证明：若函数 $f(x)$ 在每一个闭区间 $[a, c]$ 与 $[c, b]$ 上是一致连续的，则这个函数在合起来的闭区间 $[a, b]$ 上也是一致连续的.
\end{prob}

\begin{prob}[L]{802}
对于 $\varepsilon > 0$，求使下列函数 $f(x)$ 在所给区间上满足一致连续条件的 $\delta = \delta(\varepsilon)$（任意的！）：
\begin{enumerate}[label=(\alph*), itemsep=2pt, topsep=3pt]
    \item $f(x) = 5x - 3 \quad (-\infty < x < +\infty)$;
    \item $f(x) = x^2 - 2x - 1 \quad (-2 \leqslant x \leqslant 5)$;
    \item $f(x) = \frac{1}{x} \quad (0.1 \leqslant x \leqslant 1)$;
    \item $f(x) = \sqrt{x} \quad (0 \leqslant x < +\infty)$;
    \item $f(x) = 2\sin x - \cos x \quad (-\infty < x < +\infty)$;
    \item $f(x) = \begin{cases} x \sin \frac{1}{x}, & x \neq 0, \\ 0, & x = 0 \end{cases} \quad (0 \leqslant x \leqslant \pi).$
\end{enumerate}
\end{prob}

\begin{prob}[M]{803}
将闭区间 $[1, 10]$ 等分为多少段，就足以让函数 $f(x) = x^2$ 在每一段上的振幅小于 $0.0001$？
\end{prob}

\begin{prob}[M]{804}
证明：有限个在区间 $(a, b)$ 上一致连续的函数的和与积在此区间上仍是一致连续的.
\end{prob}

\begin{prob}[M]{805}
证明：若单调有界的函数 $f(x)$ 在有限或无穷的区间 $(a, b)$ 上是连续的，则此函数在区间 $(a, b)$ 上是一致连续的.
\end{prob}

\begin{prob}[M]{806.1}
证明：若函数 $f(x)$ 在有限区间 $(a, b)$ 上一致连续，则存在极限
\[
A = \lim_{x \to a+0} f(x) \quad \text{及} \quad B = \lim_{x \to b-0} f(x).
\]
这个定理对无穷区间 $(a, b)$ 是否成立？
\end{prob}

\begin{prob}[M]{806.2}
证明：在有限区间 $(a, b)$ 上有定义且连续的函数 $f(x)$，可用连续的方法延拓到闭区间 $[a, b]$ 上，其充分必要条件是函数 $f(x)$ 在区间 $(a, b)$ 上是一致连续的.
\end{prob}

\begin{prob}[M]{807}
函数
\[
\omega_f(\delta) = \sup |f(x_1) - f(x_2)|
\]
（式中 $x_1$ 和 $x_2$ 为 $(a, b)$ 中受条件 $|x_1 - x_2| \leqslant \delta$ 限制的任意两点）称为函数 $f(x)$ 在区间 $(a, b)$ 上的连续模.

证明：函数 $f(x)$ 在区间 $(a, b)$ 上一致连续的充分必要条件是
\[
\lim_{\delta \to +0} \omega_f(\delta) = 0.
\]
\end{prob}

\begin{prob}[L]{808}
设
\begin{enumerate}[label=(\alph*), itemsep=2pt, topsep=3pt]
    \item $f(x) = x^3 \quad (0 \leqslant x \leqslant 1)$;
    \item $f(x) = \sqrt{x} \quad (0 \leqslant x \leqslant a) \text{ 及 } (a < x < +\infty)$;
    \item $f(x) = \sin x + \cos x \quad (0 \leqslant x \leqslant 2\pi).$
\end{enumerate}
对 $f(x)$ 的连续模 $\omega_f(\delta)$（参阅前题）作如下形式的估计：
\[
\omega_f(\delta) \leqslant C\delta^{\alpha},
\]
式中 $C$ 和 $\alpha$ 为常数.
\end{prob}

\probsec{§10. 函数方程}

\begin{prob}[M]{809}
证明：对于所有实数 $x$ 和 $y$ 都满足方程
\[
f(x + y) = f(x) + f(y) \tag{1}
\]
的唯一的连续函数 $f(x)$ $(-\infty < x < +\infty)$ 是线性齐次函数
\[
f(x) = ax,
\]
式中 $a = f(1)$ 是任意的常数.
\end{prob}

\begin{prob}[M]{810}
证明：满足方程 $(1)$ 的单调函数 $f(x)$ 是线性齐次的.
\end{prob}

\begin{prob}[M]{811}
证明：满足方程 $(1)$ 且在任意小区间 $(-\varepsilon, \varepsilon)$ 上有界的函数 $f(x)$，是线性齐次函数.
\end{prob}

\begin{prob}[M]{812}
证明：对于所有 $x$ 和 $y$ 都满足方程
\[
f(x + y) = f(x) f(y) \tag{2}
\]
的唯一不恒等于零的连续函数 $f(x)$ $(-\infty < x < +\infty)$ 是指数函数：
\[
f(x) = a^x,
\]
式中 $a = f(1)$ 为正的常数.
\end{prob}

\begin{prob}[M]{813}
证明：在区间 $(0, \varepsilon)$ 中有界并满足方程 $(2)$ 的不恒等于零的函数 $f(x)$ 是指数函数.
\end{prob}

\begin{prob}[M]{814}
证明：对于所有正数 $x$ 和 $y$ 都满足方程
\[
f(xy) = f(x) + f(y)
\]
的唯一不恒等于零的连续函数 $f(x)$ $(0 < x < +\infty)$ 是对数函数：
\[
f(x) = \log_a x,
\]
式中 $a$ 为正的常数.
\end{prob}

\begin{prob}[M]{815}
证明：对于所有正数 $x$ 和 $y$ 都满足方程
\[
f(xy) = f(x) f(y), \tag{3}
\]
的唯一不恒等于零的连续函数 $f(x)$ $(0 < x < +\infty)$ 是幂函数：
\[
f(x) = x^a,
\]
式中 $a$ 为常数.
\end{prob}

\begin{prob}[M]{816}
求对于所有实数 $x$ 和 $y$ 都满足方程 $(3)$ 的所有的连续函数 $f(x)$ $(-\infty < x < +\infty)$.
\end{prob}

\begin{prob}[S]{817}
证明：不连续函数 $f(x) = \operatorname{sgn} x$ 满足方程 $(3)$.
\end{prob}

\begin{prob}[M]{818}
求对于所有实数 $x$ 和 $y$ 都满足方程
\[
f(x + y) + f(x - y) = 2f(x) f(y)
\]
的所有的连续函数 $f(x)$ $(-\infty < x < +\infty)$.
\end{prob}

\begin{prob}[M]{819}
求对于所有实数 $x$ 和 $y$ 都满足方程组
\[
f(x + y) = f(x) f(y) - g(x) g(y), \quad g(x + y) = f(x) g(y) + f(y) g(x),
\]
以及规范条件
\[
f(0) = 1 \quad \text{和} \quad g(0) = 0
\]
的所有的有界连续函数 $f(x)$ 和 $g(x)$ $(-\infty < x < +\infty)$.
\probhint{研究函数 $F(x) = f^2(x) + g^2(x)$.}
\end{prob}

\begin{prob}[M]{820}
设
\[
\Delta f(x) = f(x + \Delta x) - f(x), \quad \Delta^2 f(x) = \Delta\{\Delta f(x)\}
\]
分别为函数 $f(x)$ 的一阶、二阶有限差分.

证明：若函数 $f(x)$ $(-\infty < x < +\infty)$ 是连续的且
\[
\Delta^2 f(x) \equiv 0,
\]
则此函数是线性函数，即
\[
f(x) = ax + b,
\]
式中 $a$ 和 $b$ 为常数.
\end{prob}

\probchapter{第二章\quad 一元函数微分学}

\probsec{§1. 显函数的导数}

\begin{prob}[S]{821}
若 $x$ 由 $1$ 变到 $1000$, 求自变量 $x$ 的增量 $\Delta x$ 和函数 $y = \lg x$ 的相应增量 $\Delta y$.
\end{prob}

\begin{prob}[S]{822}
若 $x$ 由 $0.01$ 变到 $0.001$, 求自变量 $x$ 的增量 $\Delta x$ 和函数 $y = \frac{1}{x^2}$ 的相应增量 $\Delta y$.
\end{prob}

\begin{prob}[L]{823}
设:
\begin{enumerate}[label=(\alph*), itemsep=2pt, topsep=3pt]
\item $y = ax + b$;
\item $y = ax^2 + bx + c$;
\item $y = a^x$.
\end{enumerate}
若变量 $x$ 的增量为 $\Delta x$, 求增量 $\Delta y$.
\end{prob}

\begin{prob}[M]{824}
证明:
\begin{enumerate}[label=(\alph*), itemsep=2pt, topsep=3pt]
\item $\Delta[f(x) + g(x)] = \Delta f(x) + \Delta g(x)$;
\item $\Delta[f(x)g(x)] = g(x + \Delta x)\Delta f(x) + f(x)\Delta g(x)$.
\end{enumerate}
\end{prob}

\begin{prob}[L]{825}
过曲线 $y = x^2$ 上的两点 $A(2, 4)$ 和 $A'(2 + \Delta x, 4 + \Delta y)$ 引割线 $AA'$, 求此割线的斜率, 设:
\begin{enumerate}[label=(\alph*), itemsep=2pt, topsep=3pt]
\item $\Delta x = 1$;
\item $\Delta x = 0.1$;
\item $\Delta x = 0.01$;
\item $\Delta x$ 为任意小量.
\end{enumerate}
在该曲线上 $A$ 点的切线的斜率等于什么?
\end{prob}

\begin{prob}[L]{826}
利用函数 $y = x^3$ 把 $Ox$ 轴上的线段 $1 \leqslant x \leqslant 1 + h$ 映射到 $Oy$ 轴上, 求其平均伸长系数. 设:
\begin{enumerate}[label=(\alph*), itemsep=2pt, topsep=3pt]
\item $h = 0.1$;
\item $h = 0.01$;
\item $h = 0.001$,
\end{enumerate}
计算此系数的值. 在 $x = 1$, 此映射的伸长系数等于什么?
\end{prob}

\begin{prob}[L]{827}
动点沿 $Ox$ 轴运动的规律由下式给出:
\[
x = 10t + 5t^2,
\]
式中 $t$ 为以 $\mathrm{s}$ (秒) 计的时间, $x$ 为以 $\mathrm{m}$ (米) 计的距离. 求在 $20 \leqslant t \leqslant 20 + \Delta t$ 时间内运动的平均速度. 设:
\begin{enumerate}[label=(\alph*), itemsep=2pt, topsep=3pt]
\item $\Delta t = 1$;
\item $\Delta t = 0.1$;
\item $\Delta t = 0.01$,
\end{enumerate}
计算此速度的值. 当 $t = 20$ 时运动速度等于什么?
\end{prob}

\begin{prob}[L]{828}
根据导数的定义, 直接求下列函数的导数:
\begin{enumerate}[label=(\alph*), itemsep=2pt, topsep=3pt]
\item $x^2$;
\item $x^3$;
\item $\frac{1}{x}$;
\item $\sqrt{x}$;
\item $\sqrt[3]{x}$;
\item $\tan x$;
\item $\cot x$;
\item $\arcsin x$;
\item $\arccos x$;
\item $\arctan x$.
\end{enumerate}
\end{prob}

\begin{prob}[M]{829}
设 $f(x) = (x - 1)(x - 2)^2(x - 3)^3$, 求 $f'(1)$, $f'(2)$ 和 $f'(3)$.
\end{prob}

\begin{prob}[S]{830}
设 $f(x) = x^2\sin(x - 2)$, 求 $f'(2)$.
\end{prob}

\begin{prob}[M]{831}
设 $f(x) = x + (x - 1)\arcsin\sqrt{\frac{x}{x + 1}}$, 求 $f'(1)$.
\end{prob}

\begin{prob}[M]{832}
设函数 $f(x)$ 在 $a$ 点可微, 求
\[
\lim_{x \to a} \frac{f(x) - f(a)}{x - a}.
\]
\end{prob}

\begin{prob}[L]{833}
证明: 若函数 $f(x)$ 可微, $n$ 为正整数, 则
\[
\lim_{n \to \infty} n\left[ f\left(x + \frac{1}{n}\right) - f(x) \right] = f'(x). \tag{1}
\]
反之, 若对于函数 $f(x)$ 有极限 $(1)$ 存在, 则可否断定此函数有导数? 研究狄利克雷函数的例子 (参阅习题 734).
\end{prob}

\probgroup{利用导数表, 求下列函数的导数:}

\begin{prob}[S]{834}
$y = 2 + x - x^2$. 问 $y'(0)$, $y'\left(\frac{1}{2}\right)$, $y'(1)$, $y'(-10)$ 等于什么?
\end{prob}

\begin{prob}[L]{835}
$y = \frac{x^3}{3} + \frac{x^2}{2} - 2x$. 当 $x$ 为何值时:
\begin{enumerate}[label=(\alph*), itemsep=2pt, topsep=3pt]
\item $y'(x) = 0$;
\item $y'(x) = -2$;
\item $y'(x) = 10$?
\end{enumerate}
\end{prob}

\begin{prob}[S]{836}
$y = a^5 + 5a^3x^2 - x^5$.
\end{prob}

\begin{prob}[S]{837}
$y = \frac{ax + b}{a + b}$.
\end{prob}

\begin{prob}[S]{838}
$y = (x - a)(x - b)$.
\end{prob}

\begin{prob}[S]{839}
$y = (x + 1)(x + 2)^2(x + 3)^3$.
\end{prob}

\begin{prob}[S]{840}
$y = (x\sin\alpha + \cos\alpha)(x\cos\alpha - \sin\alpha)$.
\end{prob}

\begin{prob}[S]{841}
$y = (1 + nx^m)(1 + mx^n)$.
\end{prob}

\begin{prob}[M]{842}
(a) $y = (1 - x)(1 - x^2)^2(1 - x^3)^3$; \qquad (b) $y = (5 + 2x)^{10}(3 - 4x)^{20}$.
\end{prob}

\begin{prob}[S]{843}
$y = \frac{1}{x} + \frac{2}{x^2} + \frac{3}{x^3}$.
\end{prob}

\begin{prob}[M]{844}
证明: 公式 $\left(\frac{ax + b}{cx + d}\right)' = \frac{\begin{vmatrix} a & b \\ c & d \end{vmatrix}}{(cx + d)^2}$.
\end{prob}

\probgroup{求下列函数的导数:}

\begin{prob}[S]{845}
$y = \frac{2x}{1 - x^2}$.
\end{prob}

\begin{prob}[S]{846}
$y = \frac{1 + x - x^2}{1 - x + x^2}$.
\end{prob}

\begin{prob}[S]{847}
$y = \frac{x}{(1 - x)^2(1 + x)^3}$.
\end{prob}

\begin{prob}[S]{848}
$y = \frac{(2 - x^2)(3 - x^3)}{(1 - x)^2}$.
\end{prob}

\begin{prob}[S]{849}
$y = \frac{(1 - x)^p}{(1 + x)^q}$.
\end{prob}

\begin{prob}[S]{850}
$y = \frac{x^p(1 - x)^q}{1 + x}$.
\end{prob}

\begin{prob}[S]{851}
$y = x + \sqrt{x} + \sqrt[3]{x}$.
\end{prob}

\begin{prob}[S]{852}
$y = \frac{1}{x} + \frac{1}{\sqrt{x}} + \frac{1}{\sqrt[3]{x}}$.
\end{prob}

\begin{prob}[S]{853}
$y = \sqrt[3]{x^2} - \frac{2}{\sqrt{x}}$.
\end{prob}

\begin{prob}[S]{854}
$y = x\sqrt{1 + x^2}$.
\end{prob}

\begin{prob}[S]{855}
$y = (1 + x)\sqrt{2 + x^2}\sqrt[3]{3 + x^3}$.
\end{prob}

\begin{prob}[S]{856}
$y = \sqrt[m+n]{(1 - x)^m(1 + x)^n}$.
\end{prob}

\begin{prob}[S]{857}
$y = \frac{x}{\sqrt{a^2 - x^2}}$.
\end{prob}

\begin{prob}[S]{858}
$y = \sqrt[3]{\frac{1 + x^3}{1 - x^3}}$.
\end{prob}

\begin{prob}[S]{859}
$y = \frac{1}{\sqrt{1 + x^2}(x + \sqrt{1 + x^2})}$.
\end{prob}

\begin{prob}[S]{860}
$y = \sqrt{x + \sqrt{x + \sqrt{x}}}$.
\end{prob}

\begin{prob}[S]{861}
$y = \sqrt[3]{1 + \sqrt[3]{1 + \sqrt[3]{x}}}$.
\end{prob}

\begin{prob}[S]{862}
$y = \cos 2x - 2\sin x$.
\end{prob}

\begin{prob}[S]{863}
$y = (2 - x^2)\cos x + 2x\sin x$.
\end{prob}

\begin{prob}[S]{864}
$y = \sin(\cos^2 x)\cdot\cos(\sin^2 x)$.
\end{prob}

\begin{prob}[S]{865}
$y = \sin^n x \cos nx$.
\end{prob}

\begin{prob}[S]{866}
$y = \sin[\sin(\sin x)]$.
\end{prob}

\begin{prob}[S]{867}
$y = \frac{\sin^2 x}{\sin x^2}$.
\end{prob}

\begin{prob}[S]{868}
$y = \frac{\cos x}{2\sin^2 x}$.
\end{prob}

\begin{prob}[S]{869}
$y = \frac{1}{\cos^n x}$.
\end{prob}

\begin{prob}[S]{870}
$y = \frac{\sin x - x\cos x}{\cos x + x\sin x}$.
\end{prob}

\begin{prob}[S]{871}
$y = \tan\frac{x}{2} - \cot\frac{x}{2}$.
\end{prob}

\begin{prob}[S]{872}
$y = \tan x - \frac{1}{3}\tan^3 x + \frac{1}{5}\tan^5 x$.
\end{prob}

\begin{prob}[S]{873}
$y = 4\sqrt[3]{\cot^2 x} + \sqrt[3]{\cot^8 x}$.
\end{prob}

\begin{prob}[S]{874}
$y = \sec^2\frac{x}{a} + \csc^2\frac{x}{a}$.
\end{prob}

\begin{prob}[S]{875}
$y = \sin[\cos^2(\tan^3 x)]$.
\end{prob}

\begin{prob}[S]{876}
$y = \mathrm{e}^{-x^2}$.
\end{prob}

\begin{prob}[S]{877}
$y = 2^{\tan\frac{1}{x}}$.
\end{prob}

\begin{prob}[S]{878}
$y = \mathrm{e}^x(x^2 - 2x + 2)$.
\end{prob}

\begin{prob}[S]{879}
$y = \left[\frac{1 - x^2}{2}\sin x - \frac{(1 + x)^2}{2}\cos x\right]\mathrm{e}^{-x}$.
\end{prob}

\begin{prob}[S]{880}
$y = \mathrm{e}^x\left(1 + \cot\frac{x}{2}\right)$.
\end{prob}

\begin{prob}[S]{881}
$y = \frac{\ln 3 \cdot \sin x + \cos x}{3^x}$.
\end{prob}

\begin{prob}[S]{882}
$y = \mathrm{e}^{ax}\frac{a\sin bx - b\cos bx}{\sqrt{a^2 + b^2}}$.
\end{prob}

\begin{prob}[S]{883}
$y = \mathrm{e}^x + \mathrm{e}^{\mathrm{e}^x} + \mathrm{e}^{\mathrm{e}^{\mathrm{e}^x}}$.
\end{prob}

\begin{prob}[S]{884}
$y = \left(\frac{a}{b}\right)^x \left(\frac{b}{x}\right)^a \left(\frac{x}{a}\right)^b$ $(a > 0, b > 0)$.
\end{prob}

\begin{prob}[S]{885}
$y = x^{a^a} + a^{x^a} + a^{a^x}$ $(a > 0)$.
\end{prob}

\begin{prob}[S]{886}
$y = \lg^3 x^2$.
\end{prob}

\begin{prob}[S]{887}
$y = \ln(\ln(\ln x))$.
\end{prob}

\begin{prob}[S]{888}
$y = \ln(\ln^2(\ln^3 x))$.
\end{prob}

\begin{prob}[S]{889}
$y = \frac{1}{2}\ln(1 + x) - \frac{1}{4}\ln(1 + x^2) - \frac{1}{2(1 + x)}$.
\end{prob}

\begin{prob}[S]{890}
$y = \frac{1}{4}\ln\frac{x^2 - 1}{x^2 + 1}$.
\end{prob}

\begin{prob}[S]{891}
$y = \frac{1}{4(1 + x^4)} + \frac{1}{4}\ln\frac{x^4}{1 + x^4}$.
\end{prob}

\begin{prob}[S]{892}
$y = \frac{1}{2\sqrt{6}}\ln\frac{x\sqrt{3} - \sqrt{2}}{x\sqrt{3} + \sqrt{2}}$.
\end{prob}

\begin{prob}[S]{893}
$y = \frac{1}{1 - k}\ln\frac{1 + x}{1 - x} + \frac{\sqrt{k}}{1 - k}\ln\frac{1 + x\sqrt{k}}{1 - x\sqrt{k}}$ $(0 < k < 1)$.
\end{prob}

\begin{prob}[S]{894}
$y = \sqrt{x + 1} - \ln(1 + \sqrt{x + 1})$.
\end{prob}

\begin{prob}[S]{895}
$y = \ln(x + \sqrt{x^2 + 1})$.
\end{prob}

\begin{prob}[S]{896}
$y = x\ln(x + \sqrt{1 + x^2}) - \sqrt{1 + x^2}$.
\end{prob}

\begin{prob}[S]{897}
$y = x\ln^2(x + \sqrt{1 + x^2}) - 2\sqrt{1 + x^2}\ln(x + \sqrt{1 + x^2}) + 2x$.
\end{prob}

\begin{prob}[S]{898}
$y = \frac{x}{2}\sqrt{x^2 + a^2} + \frac{a^2}{2}\ln(x + \sqrt{x^2 + a^2})$.
\end{prob}

\begin{prob}[S]{899}
$y = \frac{1}{2\sqrt{ab}}\ln\frac{\sqrt{a} + x\sqrt{b}}{\sqrt{a} - x\sqrt{b}}$ $(a > 0, b > 0)$.
\end{prob}

\begin{prob}[S]{900}
$y = \frac{2 + 3x^2}{x^4}\sqrt{1 - x^2} + 3\ln\frac{1 + \sqrt{1 - x^2}}{x}$.
\end{prob}

\begin{prob}[S]{901}
$y = \ln\tan\frac{x}{2}$.
\end{prob}

\begin{prob}[S]{902}
$y = \ln\tan\left(\frac{x}{2} + \frac{\pi}{4}\right)$.
\end{prob}

\begin{prob}[S]{903}
$y = \frac{1}{2}\cot^2 x + \ln\sin x$.
\end{prob}

\begin{prob}[S]{904}
$y = \ln\sqrt{\frac{1 - \sin x}{1 + \sin x}}$.
\end{prob}

\begin{prob}[S]{905}
$y = -\frac{\cos x}{2\sin^2 x} + \ln\sqrt{\frac{1 + \cos x}{\sin x}}$.
\end{prob}

\begin{prob}[S]{906}
$y = \ln\frac{b + a\cos x + \sqrt{b^2 - a^2}\sin x}{a + b\cos x}$ $(0 \leqslant |a| < |b|)$.
\end{prob}

\begin{prob}[S]{907}
$y = \frac{1}{x}(\ln^3 x + 3\ln^2 x + 6\ln x + 6)$.
\end{prob}

\begin{prob}[S]{908}
$y = \frac{1}{4x^4}\ln\frac{1}{x} - \frac{1}{16x^4}$.
\end{prob}

\begin{prob}[S]{909}
$y = \frac{3}{2}\left(1 - \sqrt[3]{1 + x^2}\right)^2 + 3\ln\left(1 + \sqrt[3]{1 + x^2}\right)$.
\end{prob}

\begin{prob}[S]{910}
$y = \ln\left[\frac{1}{x} + \ln\left(\frac{1}{x} + \ln\frac{1}{x}\right)\right]$.
\end{prob}

\begin{prob}[S]{911}
$y = x[\sin(\ln x) - \cos(\ln x)]$.
\end{prob}

\begin{prob}[S]{912}
$y = \ln\tan\frac{x}{2} - \cos x \cdot \ln\tan x$.
\end{prob}

\begin{prob}[S]{913}
$y = \arcsin\frac{x}{2}$.
\end{prob}

\begin{prob}[S]{914}
$y = \arccos\frac{1 - x}{\sqrt{2}}$.
\end{prob}

\begin{prob}[S]{915}
$y = \arctan\frac{x^2}{a}$.
\end{prob}

\begin{prob}[S]{916}
$y = \frac{1}{\sqrt{2}}\operatorname{arccot}\frac{\sqrt{2}}{x}$.
\end{prob}

\begin{prob}[S]{917}
$y = \sqrt{x} - \arctan\sqrt{x}$.
\end{prob}

\begin{prob}[S]{918}
$y = x + \sqrt{1 - x^2}\cdot\arccos x$.
\end{prob}

\begin{prob}[S]{919}
$y = x\arcsin\sqrt{\frac{x}{1 + x}} + \arctan\sqrt{x} - \sqrt{x}$.
\end{prob}

\begin{prob}[S]{920}
$y = \arccos\frac{1}{x}$.
\end{prob}

\begin{prob}[S]{921}
$y = \arcsin(\sin x)$.
\end{prob}

\begin{prob}[S]{922}
$y = \arccos(\cos^2 x)$.
\end{prob}

\begin{prob}[S]{923}
$y = \arcsin(\sin x - \cos x)$.
\end{prob}

\begin{prob}[S]{924}
$y = \arccos\sqrt{1 - x^2}$.
\end{prob}

\begin{prob}[S]{925}
$y = \arctan\frac{1 + x}{1 - x}$.
\end{prob}

\begin{prob}[S]{926}
$y = \operatorname{arccot}\left(\frac{\sin x + \cos x}{\sin x - \cos x}\right)$.
\end{prob}

\begin{prob}[S]{927}
$y = \frac{2}{\sqrt{a^2 - b^2}} \arctan \left( \sqrt{\frac{a - b}{a + b}} \tan \frac{x}{2} \right) \ (a > b \geqslant 0)$.
\end{prob}

\begin{prob}[S]{928}
$y = \arcsin \frac{1 - x^2}{1 + x^2}$.
\end{prob}

\begin{prob}[S]{929}
$y = \frac{1}{\arccos^2(x^2)}$.
\end{prob}

\begin{prob}[S]{930}
$y = \arctan x + \frac{1}{3} \arctan(x^3)$.
\end{prob}

\begin{prob}[S]{931}
$y = \ln(1 + \sin^2 x) - 2 \sin x \cdot \arctan(\sin x)$.
\end{prob}

\begin{prob}[S]{932}
$y = \ln \left( \arccos \frac{1}{\sqrt{x}} \right)$.
\end{prob}

\begin{prob}[S]{933}
$y = \ln \frac{x + a}{\sqrt{x^2 + b^2}} + \frac{a}{b} \arctan \frac{x}{b}$.
\end{prob}

\begin{prob}[S]{934}
$y = \frac{x}{2} \sqrt{a^2 - x^2} + \frac{a^2}{2} \arcsin \frac{x}{a} \ (a > 0)$.
\end{prob}

\begin{prob}[S]{935}
$y = \frac{1}{6} \ln \frac{(x + 1)^2}{x^2 - x + 1} + \frac{1}{\sqrt{3}} \arctan \frac{2x - 1}{\sqrt{3}}$.
\end{prob}

\begin{prob}[S]{936}
$y = \frac{1}{4\sqrt{2}} \ln \frac{x^2 + x\sqrt{2} + 1}{x^2 - x\sqrt{2} + 1} - \frac{1}{2\sqrt{2}} \arctan \frac{x\sqrt{2}}{x^2 - 1}$.
\end{prob}

\begin{prob}[S]{937}
$y = x(\arcsin x)^2 + 2\sqrt{1 - x^2} \arcsin x - 2x$.
\end{prob}

\begin{prob}[S]{938}
$y = \frac{\arccos x}{x} + \frac{1}{2} \ln \frac{1 - \sqrt{1 - x^2}}{1 + \sqrt{1 - x^2}}$.
\end{prob}

\begin{prob}[S]{939}
$y = \arctan \sqrt{x^2 - 1} - \frac{\ln x}{\sqrt{x^2 - 1}}$.
\end{prob}

\begin{prob}[S]{940}
$y = \frac{\arcsin x}{\sqrt{1 - x^2}} + \frac{1}{2} \ln \frac{1 - x}{1 + x}$.
\end{prob}

\begin{prob}[S]{941}
$y = \frac{1}{12} \ln \frac{x^4 - x^2 + 1}{(x^2 + 1)^2} - \frac{1}{2\sqrt{3}} \arctan \frac{\sqrt{3}}{2x^2 - 1}$.
\end{prob}

\begin{prob}[S]{942}
$y = \frac{x^6}{1 + x^{12}} - \operatorname{arccot} x^6$.
\end{prob}

\begin{prob}[S]{943}
$y = \ln \frac{1 - \sqrt[3]{x}}{\sqrt{1 + \sqrt[3]{x} + \sqrt[3]{x^2}}} + \sqrt{3} \arctan \frac{1 + 2\sqrt[3]{x}}{\sqrt{3}}$.
\end{prob}

\begin{prob}[S]{944}
$y = \arctan \frac{x}{1 + \sqrt{1 - x^2}}$.
\end{prob}

\begin{prob}[S]{945}
$y = \operatorname{arccot} \frac{a - 2x}{2\sqrt{ax - x^2}} \ (a > 0)$.
\end{prob}

\begin{prob}[S]{946}
$y = \frac{3 - x}{2} \sqrt{1 - 2x - x^2} + 2\arcsin \frac{1 + x}{\sqrt{2}}$.
\end{prob}

\begin{prob}[S]{947}
$y = \frac{1}{4} \ln \frac{\sqrt[4]{1 + x^4} + x}{\sqrt[4]{1 + x^4} - x} - \frac{1}{2} \arctan \frac{\sqrt[4]{1 + x^4}}{x}$.
\end{prob}

\begin{prob}[S]{948}
$y = \arctan(\tan^2 x)$.
\end{prob}

\begin{prob}[S]{949}
$y = \sqrt{1 - x^2} \cdot \ln \sqrt{\frac{1 - x}{1 + x}} + \frac{1}{2} \ln \frac{1 - \sqrt{1 - x^2}}{1 + \sqrt{1 - x^2}} + \sqrt{1 - x^2} + \arcsin x$.
\end{prob}

\begin{prob}[S]{950}
$y = x\arctan x - \frac{1}{2} \ln(1 + x)^2 - \frac{1}{2} (\arctan x)^2$.
\end{prob}

\begin{prob}[S]{951}
$y = \ln(\mathrm{e}^x + \sqrt{1 + \mathrm{e}^{2x}})$.
\end{prob}

\begin{prob}[S]{952}
$y = \arctan(x + \sqrt{1 + x^2})$.
\end{prob}

\begin{prob}[S]{953}
$y = \arcsin \left( \frac{\sin a \sin x}{1 - \cos a \cos x} \right)$.
\end{prob}

\begin{prob}[S]{954}
$y = \frac{1}{4\sqrt{3}} \ln \frac{\sqrt{x^2 + 2} - x\sqrt{3}}{\sqrt{x^2 + 2} + x\sqrt{3}} + \frac{1}{2} \arctan \frac{\sqrt{x^2 + 2}}{x}$.
\end{prob}

\begin{prob}[S]{955}
$y = \frac{1}{2\sqrt{2}} \arctan \frac{x\sqrt{2}}{\sqrt{1 + x^4}} - \frac{1}{4\sqrt{2}} \ln \frac{\sqrt{1 + x^4} - x\sqrt{2}}{\sqrt{1 + x^4} + x\sqrt{2}}$.
\end{prob}

\begin{prob}[S]{956}
$y = \frac{x\sqrt{1 - x^2}}{1 + x^2} - \frac{3}{\sqrt{2}} \operatorname{arccot} \frac{x\sqrt{2}}{\sqrt{1 - x^2}}$.
\end{prob}

\begin{prob}[S]{957}
$y = \arccos(\sin x^2 - \cos x^2)$.
\end{prob}

\begin{prob}[S]{958}
$y = \arcsin(\sin x^2) + \arccos(\cos x^2)$.
\end{prob}

\begin{prob}[S]{959}
$y = \mathrm{e}^{m\arcsin x}[\cos(m\arcsin x) + \sin(m\arcsin x)]$.
\end{prob}

\begin{prob}[L]{960}
\begin{enumerate}[label=(\alph*), itemsep=2pt, topsep=3pt]
\item $y = \arctan \mathrm{e}^x - \ln \sqrt{\frac{\mathrm{e}^{2x}}{\mathrm{e}^{2x} + 1}}$;
\item $y = \sqrt{1 + \sqrt[3]{1 + \sqrt[4]{1 + x^4}}}$;
\item $y = \operatorname{arccot} \frac{1}{\sqrt{\cot \frac{1}{x^2}}}$;
\item $y = \ln^2(\sec 2^{\sqrt[3]{x}})$.
\end{enumerate}
\end{prob}

\begin{prob}[S]{961}
$y = x + x^x + x^{x^x} \ (x > 0)$.
\end{prob}

\begin{prob}[S]{962}
$y = x^{x^a} + x^{a^x} + a^{x^x} \ (a > 0, x > 0)$.
\end{prob}

\begin{prob}[S]{963}
$y = \sqrt[x]{x} \ (x > 0)$.
\end{prob}

\begin{prob}[S]{964}
$y = (\sin x)^{\cos x} + (\cos x)^{\sin x}$.
\end{prob}

\begin{prob}[S]{965.1}
$y = \frac{(\ln x)^x}{x^{\ln x}}$.
\end{prob}

\begin{prob}[S]{965.2}
$y = \left[ \frac{\arcsin(\sin^2 x)}{\arccos(\cos^2 x)} \right]^{\arctan^2 x}$.
\end{prob}

\begin{prob}[S]{966}
$y = \log_x \mathrm{e}$.
\end{prob}

\begin{prob}[S]{967}
$y = \ln(\cosh x) + \frac{1}{2\cosh^2 x}$.
\end{prob}

\begin{prob}[S]{968}
$y = \frac{\cosh x}{\sinh^2 x} - \ln\left(\coth \frac{x}{2}\right)$.
\end{prob}

\begin{prob}[S]{969}
$y = \arctan(\tanh x)$.
\end{prob}

\begin{prob}[S]{970}
$y = \arccos\left(\frac{1}{\cosh x}\right)$.
\end{prob}

\begin{prob}[S]{971}
$y = \frac{b}{a} x + \frac{2\sqrt{a^2 - b^2}}{a} \arctan \left( \sqrt{\frac{a - b}{a + b}} \tanh \frac{x}{2} \right) \ (0 \leqslant |b| < a)$.
\end{prob}

\begin{prob}[M]{972}
引入中间变量 $u = \cos^2 x$，求函数
\[
y = \ln(\cos^2 x + \sqrt{1 + \cos^4 x})
\]
的导数。
\end{prob}

\probgroup{利用 972 题所示的方法，求下列函数的导数：}

\begin{prob}[S]{973}
$y = (\arccos x)^2 \left[ \ln^2(\arccos x) - \ln(\arccos x) + \frac{1}{2} \right]$.
\end{prob}

\begin{prob}[S]{974}
$y = \frac{1}{2} \arctan(\sqrt[4]{1 + x^4}) + \frac{1}{4} \ln \frac{\sqrt[4]{1 + x^4} + 1}{\sqrt[4]{1 + x^4} - 1}$.
\end{prob}

\begin{prob}[S]{975}
$y = \frac{\mathrm{e}^{-x^2} \arcsin(\mathrm{e}^{-x^2})}{\sqrt{1 - \mathrm{e}^{-2x^2}}} + \frac{1}{2} \ln(1 - \mathrm{e}^{-2x^2})$.
\end{prob}

\begin{prob}[S]{976}
$y = \frac{a^x}{1 + a^{2x}} - \frac{1 - a^{2x}}{1 + a^{2x}} \operatorname{arccot} a^{-x}$.
\end{prob}

\begin{prob}[L]{977}
求下列函数的导数并作出函数及其导数的图像：
\begin{enumerate}[label=(\alph*), itemsep=2pt, topsep=3pt]
\item $y = |x|$;
\item $y = x|x|$;
\item $y = \ln |x|$.
\end{enumerate}
\end{prob}

\begin{prob}[L]{978}
求下列函数的导数：
\begin{enumerate}[label=(\alph*), itemsep=2pt, topsep=3pt]
\item $y = |(x - 1)^2 (x + 1)^3|$;
\item $y = |\sin^3 x|$;
\item $y = \arccos \frac{1}{|x|}$;
\item $y = [x] \sin^2 \pi x$.
\end{enumerate}
\end{prob}

\probgroup{求下列函数的导数并作出函数及其导数的图像：}

\begin{prob}[M]{979}
\[
y = \begin{cases}
1 - x, & -\infty < x < 1, \\
(1 - x)(2 - x), & 1 \leqslant x \leqslant 2, \\
-(2 - x), & 2 < x < +\infty.
\end{cases}
\]
\end{prob}

\begin{prob}[M]{980}
\[
y = \begin{cases}
(x - a)^2 (x - b)^2, & a \leqslant x \leqslant b, \\
0, & \text{在线段 } [a, b] \text{ 之外}.
\end{cases}
\]
\end{prob}

\begin{prob}[M]{981}
\[
y = \begin{cases}
x, & x < 0, \\
\ln(1 + x), & x \geqslant 0.
\end{cases}
\]
\end{prob}

\begin{prob}[M]{982}
\[
y = \begin{cases}
\arctan x, & |x| \leqslant 1, \\
\dfrac{\pi}{4} \operatorname{sgn} x + \dfrac{x - 1}{2}, & |x| > 1.
\end{cases}
\]
\end{prob}

\begin{prob}[M]{983}
\[
y = \begin{cases}
x^2 \mathrm{e}^{-x^2}, & |x| \leqslant 1, \\
\dfrac{1}{\mathrm{e}}, & |x| > 1.
\end{cases}
\]
\end{prob}

\begin{prob}[L]{984}
所给函数 $y = f(x)$ 的对数的导数称为此函数的对数导数：
\[
\frac{y'}{y} = \frac{\mathrm{d}}{\mathrm{d}x} \ln f(x) = \frac{f'(x)}{f(x)} \quad (f(x) > 0).
\]
求下列函数 $y$ 的对数导数：
\begin{enumerate}[label=(\alph*), itemsep=2pt, topsep=3pt]
\item $y = x \sqrt{\frac{1 - x}{1 + x}}$;
\item $y = \frac{x^2}{1 - x} \sqrt[3]{\frac{3 - x}{(3 + x)^2}}$;
\item $y = (x - a_1)^{\alpha_1} (x - a_2)^{\alpha_2} \cdots (x - a_n)^{\alpha_n}$;
\item $y = \left(x + \sqrt{1 + x^2}\right)^n$.
\end{enumerate}
\end{prob}

\begin{prob}[L]{985}
设 $\varphi(x)$ 及 $\psi(x)$ 为 $x$ 的可微函数. 求下列函数 $y$ 的导数：
\begin{enumerate}[label=(\alph*), itemsep=2pt, topsep=3pt]
\item $y = \sqrt{\varphi^2(x) + \psi^2(x)}$;
\item $y = \arctan \frac{\varphi(x)}{\psi(x)}$;
\item $y = \sqrt[\varphi(x)]{\psi(x)} \ (\varphi(x) \neq 0, \ \psi(x) > 0)$;
\item $y = \log_{\varphi(x)} \psi(x) \ (\varphi(x) > 0, \ \psi(x) > 0)$.
\end{enumerate}
\end{prob}

\begin{prob}[L]{986.1}
求 $y'$，设
\begin{enumerate}[label=(\alph*), itemsep=2pt, topsep=3pt]
\item $y = f(x^2)$;
\item $y = f(\sin^2 x) + f(\cos^2 x)$;
\item $y = f(\mathrm{e}^x) \cdot \mathrm{e}^{f(x)}$;
\item $y = f\{f[f(x)]\}$，
\end{enumerate}
其中 $f(u)$ 为可微函数.
\end{prob}

\begin{prob}[M]{986.2}
设 $f(x) = x(x - 1)(x - 2)\cdots(x - 1000)$，求 $f'(0)$.
\end{prob}

\begin{prob}[L]{987}
证明：$n$ 阶行列式微分法：
\[
\begin{vmatrix}
f_{11}(x) & f_{12}(x) & \cdots & f_{1n}(x) \\
\vdots & \vdots & & \vdots \\
f_{k1}(x) & f_{k2}(x) & \cdots & f_{kn}(x) \\
\vdots & \vdots & & \vdots \\
f_{n1}(x) & f_{n2}(x) & \cdots & f_{nn}(x)
\end{vmatrix}'
= \sum_{k=1}^{n}
\begin{vmatrix}
f_{11}(x) & f_{12}(x) & \cdots & f_{1n}(x) \\
\vdots & \vdots & & \vdots \\
f'_{k1}(x) & f'_{k2}(x) & \cdots & f'_{kn}(x) \\
\vdots & \vdots & & \vdots \\
f_{n1}(x) & f_{n2}(x) & \cdots & f_{nn}(x)
\end{vmatrix}.
\]
\end{prob}

\begin{prob}[M]{988}
设
\[
F(x) = \begin{vmatrix}
x - 1 & 1 & 2 \\
-3 & x & 3 \\
-2 & -3 & x + 1
\end{vmatrix},
\]
求 $F'(x)$.
\end{prob}

\begin{prob}[M]{989}
设
\[
F(x) = \begin{vmatrix}
x & x^2 & x^3 \\
1 & 2x & 3x^2 \\
0 & 2 & 6x
\end{vmatrix},
\]
求 $F'(x)$.
\end{prob}

\begin{prob}[M]{990}
已知函数的图像，近似地作出其导数的图像.
\end{prob}

\begin{prob}[M]{991}
证明：函数
\[
f(x) = \begin{cases}
x^2 \sin \dfrac{1}{x}, & x \neq 0, \\
0, & x = 0
\end{cases}
\]
有不连续的导数.
\end{prob}

\begin{prob}[L]{992}
在什么条件下，函数
\[
f(x) = \begin{cases}
x^n \sin \dfrac{1}{x}, & x \neq 0, \\
0, & x = 0
\end{cases}
\]
\begin{enumerate}[label=(\alph*), itemsep=2pt, topsep=3pt]
\item 在 $x = 0$ 处是连续的；
\item 在 $x = 0$ 处可微；
\item 在 $x = 0$ 处有连续的导数？
\end{enumerate}
\end{prob}

\begin{prob}[M]{993}
在什么条件下，函数
\[
f(x) = \begin{cases}
|x|^n \sin \dfrac{1}{|x|^m}, & x \neq 0, \\
0, & x = 0
\end{cases}
\quad (m > 0)
\]
(a) 于坐标原点的邻域上有有界的导数； \qquad (b) 在此邻域上有无穷导数？
\end{prob}

\begin{prob}[M]{994}
设
\[
f(x) = (x - a)\varphi(x),
\]
其中函数 $\varphi(x)$ 在 $x = a$ 处是连续的，求 $f'(a)$.
\end{prob}

\begin{prob}[M]{995}
设
\[
f(x) = |x - a|\varphi(x),
\]
其中 $\varphi(x)$ 为连续函数且 $\varphi(a) \neq 0$，证明：此函数在 $a$ 点没有导数. 单侧导数 $f'_-(a)$ 及 $f'_+(a)$ 等于什么？
\end{prob}

\begin{prob}[M]{996}
举出在已知点 $a_1, a_2, \cdots, a_n$ 没有导数的连续函数的例子.
\end{prob}

\begin{prob}[M]{997}
证明：函数
\[
f(x) = \begin{cases}
x^2 \left|\cos \dfrac{\pi}{x}\right|, & x \neq 0, \\
0, & x = 0
\end{cases}
\]
在点 $x = 0$ 的任何邻域上都有不可微点，但在该点是可微的. 作出此函数的略图.
\end{prob}

\begin{prob}[M]{998}
证明：函数
\[
f(x) = \begin{cases}
x^2, & x \text{ 为有理数}, \\
0, & x \text{ 为无理数}
\end{cases}
\]
仅在 $x = 0$ 时有导数.
\end{prob}

\begin{prob}[L]{999}
研究下列函数的可微性：
\begin{enumerate}[label=(\alph*), itemsep=2pt, topsep=3pt]
\item $y = |(x - 1)(x - 2)^2(x - 3)^3|$;
\item $y = |\cos x|$;
\item $y = |\pi^2 - x^2|\sin^2 x$;
\item $y = \arcsin(\cos x)$;
\item $y = \begin{cases} \dfrac{x - 1}{4}(x + 1)^2, & |x| \leqslant 1, \\ |x| - 1, & |x| > 1. \end{cases}$
\end{enumerate}
\end{prob}

\probgroup{求函数 $f(x)$ 的左导数 $f'_-(x)$ 和右导数 $f'_+(x)$，设：}

\begin{prob}[S]{1000}
$f(x) = |x|$.
\end{prob}

\begin{prob}[S]{1001}
$f(x) = [x] \sin \pi x$.
\end{prob}

\begin{prob}[M]{1002}
$f(x) = \begin{cases}
x \left|\cos \dfrac{\pi}{x}\right|, & x \neq 0, \\
0, & x = 0.
\end{cases}$
\end{prob}

\begin{prob}[S]{1003}
$f(x) = \sqrt{\sin x^2}$.
\end{prob}

\begin{prob}[M]{1004}
$f(x) = \begin{cases}
\dfrac{x}{1 + \mathrm{e}^{\frac{1}{x}}}, & x \neq 0, \\
0, & x = 0.
\end{cases}$
\end{prob}

\begin{prob}[S]{1005}
$f(x) = \sqrt{1 - \mathrm{e}^{-x^2}}$.
\end{prob}

\begin{prob}[S]{1006}
$f(x) = |\ln|x|| \ (x \neq 0)$.
\end{prob}

\begin{prob}[S]{1007}
$f(x) = \arcsin \frac{2x}{1 + x^2}$.
\end{prob}

\begin{prob}[M]{1008}
$f(x) = \begin{cases}
(x - 2)\arctan \dfrac{1}{x - 2}, & x \neq 2, \\
0, & x = 2.
\end{cases}$
\end{prob}

\begin{prob}[M]{1009.1}
证明：函数
\[
f(x) = \begin{cases}
x \sin \dfrac{1}{x}, & x \neq 0, \\
0, & x = 0
\end{cases}
\]
在点 $x = 0$ 连续，但在此点既无左导数，又无右导数.
\end{prob}

\begin{prob}[M]{1009.2}
设 $x_0$ 为函数 $f(x)$ 的第一类间断点. 表达式
\[
f'_-(x_0) = \lim_{h \to -0} \frac{f(x_0 + h) - f(x_0 - 0)}{h},
\]
\[
f'_+(x_0) = \lim_{h \to +0} \frac{f(x_0 + h) - f(x_0 + 0)}{h}
\]
称为函数 $f(x)$ 在点 $x_0$ 的广义单侧导数（左导数和右导数）.
求函数 $f(x)$ 在间断点 $x_0$ 的 $f'_-(x_0)$ 和 $f'_+(x_0)$：
\begin{enumerate}[label=(\alph*), itemsep=2pt, topsep=3pt]
  \item $f(x) = \dfrac{\sqrt{x^2 + x^3}}{x}$；
  \item $f(x) = \arctan \dfrac{1+x}{1-x}$；
  \item $f(x) = \dfrac{1}{1 + \mathrm{e}^{\frac{1}{x}}}$．
\end{enumerate}
\end{prob}

\begin{prob}[M]{1010}
设
\[
f(x) = \begin{cases}
x^2, & x \leqslant x_0, \\
ax + b, & x > x_0.
\end{cases}
\]
为了使函数 $f(x)$ 在点 $x = x_0$ 处连续而且可微，应当如何选取系数 $a$ 和 $b$？
\end{prob}

\begin{prob}[M]{1011}
设
\[
F(x) = \begin{cases}
f(x), & x \leqslant x_0, \\
ax + b, & x > x_0,
\end{cases}
\]
其中函数 $f(x)$ 在点 $x = x_0$ 处为左可微的．应当如何选取系数 $a$ 和 $b$，使函数 $F(x)$ 在点 $x_0$ 处连续而且可微？
\end{prob}

\begin{prob}[M]{1012}
适当地选定参数 $A$ 和 $c$，用立方抛物线
\[
y = A(x - a)(x - b)(x - c)
\]
在闭区间 $a \leqslant x \leqslant b$ 上把两条射线：
\[
y = k_1(x - a) \quad (-\infty < x < a) \quad \text{及} \quad y = k_2(x - b) \quad (b < x < +\infty)
\]
光滑地连接起来．
\end{prob}

\begin{prob}[M]{1013}
用抛物线 $y = a + bx^2$ ($|x| \leqslant c$)（其中 $a$ 与 $b$ 为未知的参数）去补充曲线 $y = \dfrac{m^2}{|x|}$ ($|x| > c$)，以便得到一条光滑曲线．
\end{prob}

\begin{prob}[M]{1014}
若 (a) 函数 $f(x)$ 在点 $x_0$ 有导数，而函数 $g(x)$ 在此点没有导数；(b) 函数 $f(x)$ 和 $g(x)$ 二者在点 $x_0$ 都没有导数，可否断定它们的和
\[
F(x) = f(x) + g(x)
\]
在点 $x = x_0$ 没有导数？
\end{prob}

\begin{prob}[M]{1015}
若 (a) 函数 $f(x)$ 在点 $x_0$ 有导数，而函数 $g(x)$ 在此点没有导数；(b) 函数 $f(x)$ 和 $g(x)$ 二者在点 $x_0$ 都没有导数，可否断定它们的积
\[
F(x) = f(x)g(x)
\]
在点 $x = x_0$ 没有导数？

研究两个例子：
(a) $f(x) = x$, $g(x) = |x|$, $x_0 = 0$； \qquad (b) $f(x) = |x|$, $g(x) = |x|$, $x_0 = 0$．
\end{prob}

\begin{prob}[L]{1016}
若 (a) 函数 $f(x)$ 于点 $x = g(x_0)$ 有导数，而函数 $g(x)$ 于点 $x = x_0$ 没有导数；(b) 函数 $f(x)$ 于点 $x = g(x_0)$ 没有导数，而函数 $g(x)$ 于点 $x = x_0$ 有导数；(c) 函数 $f(x)$ 于点 $x = g(x_0)$ 没有导数，函数 $g(x)$ 于点 $x = x_0$ 没有导数，则函数
\[
F(x) = f(g(x))
\]
于已知点 $x = x_0$ 的可微性怎样？

研究三个例子：
\begin{enumerate}[label=(\alph*), itemsep=2pt, topsep=3pt]
  \item $f(x) = x^2$, $g(x) = |x|$, $x_0 = 0$；
  \item $f(x) = |x|$, $g(x) = x^2$, $x_0 = 0$；
  \item $f(x) = 2x + |x|$, $g(x) = \dfrac{2}{3}x - \dfrac{1}{3}|x|$, $x_0 = 0$．
\end{enumerate}
\end{prob}

\begin{prob}[M]{1017}
函数
\[
y = x + \sqrt[3]{\sin x}
\]
的图像在哪些点处有竖直切线？作出此图像．
\end{prob}

\begin{prob}[M]{1018}
函数 $f(x)$ 在其间断点可否有：(a) 有限的导数；(b) 无穷的导数？研究例子：
\[
f(x) = \operatorname{sgn} x.
\]
\end{prob}

\begin{prob}[M]{1019}
若函数 $f(x)$ 在有限的区间 $(a, b)$ 上可微，且
\[
\lim_{x \to a} f(x) = \infty,
\]
则是否必有
\[
\text{(1) } \lim_{x \to a} f'(x) = \infty; \qquad \text{(2) } \overline{\lim}_{x \to a} |f'(x)| = +\infty\text{？}
\]
研究例子：$f(x) = \dfrac{1}{x} + \cos \dfrac{1}{x}$ 当 $x \to 0$ 时．
\end{prob}

\begin{prob}[M]{1020}
若函数 $f(x)$ 在有限的区间 $(a, b)$ 上可微，且
\[
\lim_{x \to a} f'(x) = \infty,
\]
则是否必有
\[
\lim_{x \to a} f(x) = \infty\text{？}
\]
研究例子：$f(x) = \sqrt[3]{x}$ 当 $x \to 0$ 时．
\end{prob}

\begin{prob}[M]{1021}
设函数 $f(x)$ 在区间 $(x_0, +\infty)$ 上可微，且 $\lim_{x \to +\infty} f(x)$ 存在．由此能否推出 $\lim_{x \to +\infty} f'(x)$ 存在？研究例子：
\[
f(x) = \frac{\sin(x^2)}{x}.
\]
\end{prob}

\begin{prob}[M]{1022}
设有界函数 $f(x)$ 在区间 $(x_0, +\infty)$ 上可微，且 $\lim_{x \to +\infty} f'(x)$ 存在．由此可否推出有限的或无穷的 $\lim_{x \to +\infty} f(x)$ 存在？研究例子：
\[
f(x) = \cos(\ln x).
\]
\end{prob}

\begin{prob}[S]{1023}
对函数之间的不等式可否逐项微分？
\end{prob}

\begin{prob}[L]{1024}
导出求和公式：
\begin{enumerate}[label=(\alph*), itemsep=2pt, topsep=3pt]
  \item $P_n = 1 + 2x + 3x^2 + \cdots + nx^{n-1}$，\\
  $Q_n = 1^2 + 2^2 x + 3^2 x^2 + \cdots + n^2 x^{n-1}$；
  \probhint{研究 $(x + x^2 + \cdots + x^n)'$．}
  \item $S_n = \sin x + \sin 2x + \cdots + \sin nx$，\\
  $T_n = \cos x + 2\cos 2x + \cdots + n\cos nx$．
\end{enumerate}
\end{prob}

\begin{prob}[M]{1025}
导出求和公式：
\[
S_n = \cosh x + 2\cosh 2x + \cdots + n\cosh nx.
\]
\probhint{$S_n = (\sinh x + \sinh 2x + \cdots + \sinh nx)'$．}
\end{prob}

\begin{prob}[M]{1026}
利用恒等式
\[
\cos \frac{x}{2} \cos \frac{x}{4} \cdots \cos \frac{x}{2^n} = \frac{\sin x}{2^n \sin \frac{x}{2^n}}
\]
推出求和公式：
\[
S_n = \frac{1}{2} \tan \frac{x}{2} + \frac{1}{4} \tan \frac{x}{4} + \cdots + \frac{1}{2^n} \tan \frac{x}{2^n}.
\]
\end{prob}

\begin{prob}[M]{1027}
证明：可微偶函数的导数为奇函数，而可微奇函数的导数为偶函数．给出这个事实的几何解释．
\end{prob}

\begin{prob}[M]{1028}
证明：可微周期函数的导数仍为具有相同周期的周期函数．
\end{prob}

\begin{prob}[S]{1029}
若圆半径以 $2\,\mathrm{cm/s}$ 的速度匀速增加，则当圆半径 $R = 10\,\mathrm{cm}$ 时，圆面积增加的速度如何？
\end{prob}

\begin{prob}[M]{1030}
若矩形的一边以 $1\,\mathrm{m/s}$ 的速度缩短，另一边以 $2\,\mathrm{m/s}$ 的速度伸长，则当这两条边达到 $x = 20\,\mathrm{m}$ 和 $y = 15\,\mathrm{m}$ 时，矩形的面积和对角线变化的速度如何？
\end{prob}

\begin{prob}[M]{1031}
二轮船 $A$ 和 $B$ 从同一码头同时出发，$A$ 船往北，$B$ 船往东．若 $A$ 船的速度为 $30\,\mathrm{km/h}$，$B$ 船的速度为 $40\,\mathrm{km/h}$，问二船间的距离增加的速度如何？
\end{prob}

\begin{prob}[L]{1032}
设
\[
f(x) = \begin{cases}
x, & 0 \leqslant x \leqslant 2, \\
2x - 2, & 2 < x < +\infty,
\end{cases}
\]
又设 $S(x)$ 表示由曲线 $y = f(x)$，轴 $Ox$ 及通过点 $x$ ($x \geqslant 0$) 且垂直于 $Ox$ 的直线三者围成的面积．写出函数 $S(x)$ 的解析表达式，求出导数 $S'(x)$，并作出函数 $y = S'(x)$ 的图像．
\end{prob}

\begin{prob}[L]{1033}
函数 $S(x)$ 是由圆弧 $y = \sqrt{a^2 - x^2}$、轴 $Ox$ 及分别通过点 $O$ 和 $x$ ($|x| \leqslant a$) 且垂直于轴 $Ox$ 的两条直线四者围成的面积．写出函数 $S(x)$ 的解析表达式，求出导数 $S'(x)$，并作此导数 $y = S'(x)$ 的图像．
\end{prob}

\probsec{§2. 反函数的导数. 用参数形式给出的函数的导数. 隐函数的导数}

\begin{prob}[M]{1034}
证明：由方程 $y^3 + 3y = x$ 确定的单值函数 $y = y(x)$ 存在，并求它的导数 $y'_x$．
\end{prob}

\begin{prob}[M]{1035}
证明：由方程 $y - \varepsilon \sin y = x$ ($0 \leqslant \varepsilon < 1$) 确定的单值函数 $y = y(x)$ 存在，并求它的导数 $y'_x$．
\end{prob}

\begin{prob}[L]{1036}
设：
\begin{enumerate}[label=(\alph*), itemsep=2pt, topsep=3pt]
  \item $y = x + \ln x$ ($x > 0$)；
  \item $y = x + \mathrm{e}^x$；
  \item $y = \sinh x$；
  \item $y = \tanh x$．
\end{enumerate}
求反函数 $x = x(y)$ 的存在域，并求它们的导数．
\end{prob}

\begin{prob}[L]{1037}
设：
\begin{enumerate}[label=(\alph*), itemsep=2pt, topsep=3pt]
  \item $y = 2x^2 - x^4$；
  \item $y = \dfrac{x^2}{1 + x^2}$；
  \item $y = 2\mathrm{e}^{-x} - \mathrm{e}^{-2x}$．
\end{enumerate}
选出反函数 $x = x(y)$ 的单值连续的各支，求它们的导数并作其图像．
\end{prob}

\begin{prob}[L]{1038}
作出函数 $y = y(x)$ 的略图，并求其导数 $y'_x(x)$．设：$x = -1 + 2t - t^2$，$y = 2 - 3t + t^3$．当 $x = 0$ 及 $x = -1$ 时，$y'_x(x)$ 等于什么？在何点 $M(x, y)$ 的导数 $y'_x(x) = 0$？
\end{prob}

\probgroup{求导数 $y'_x$（参数是正数），设：}

\begin{prob}[S]{1039}
$x = \sqrt[3]{1 - \sqrt{t}}$，$y = \sqrt{1 - \sqrt[3]{t}}$．
\end{prob}

\begin{prob}[S]{1040}
$x = \sin^2 t$，$y = \cos^2 t$．
\end{prob}

\begin{prob}[S]{1041}
$x = a\cos t$，$y = b\sin t$．
\end{prob}

\begin{prob}[S]{1042}
$x = a\cosh t$，$y = b\sinh t$．
\end{prob}

\begin{prob}[S]{1043}
$x = a\cos^3 t$，$y = a\sin^3 t$．
\end{prob}

\begin{prob}[S]{1044}
$x = a(t - \sin t)$，$y = a(1 - \cos t)$．
\end{prob}

\begin{prob}[S]{1045}
$x = \mathrm{e}^{2t} \cos^2 t$，$y = \mathrm{e}^{2t} \sin^2 t$．
\end{prob}

\begin{prob}[S]{1046}
$x = \arcsin \dfrac{t}{\sqrt{1 + t^2}}$，$y = \arccos \dfrac{1}{\sqrt{1 + t^2}}$．
\end{prob}

\begin{prob}[M]{1047}
证明：由方程组
\[
\begin{cases}
x = 2t + |t|, \\
y = 5t^2 + 4t|t|
\end{cases}
\]
所确定的函数 $y = y(x)$ 当 $t = 0$ 时可微，但它的导数在此点不能用普通的公式求得．
\end{prob}

\probgroup{求下列隐函数的导数 $y'_x$：}

\begin{prob}[M]{1048}
$x^2 + 2xy - y^2 = 2x$．

当 $x = 2$ 与 $y = 4$ 时及当 $x = 2$ 与 $y = 0$ 时，$y'$ 等于什么？
\end{prob}

\begin{prob}[S]{1049}
$y^2 = 2px$（抛物线）．
\end{prob}

\begin{prob}[S]{1050}
$\dfrac{x^2}{a^2} + \dfrac{y^2}{b^2} = 1$（椭圆）．
\end{prob}

\begin{prob}[S]{1051}
$\sqrt{x} + \sqrt{y} = \sqrt{a}$（抛物线）．
\end{prob}

\begin{prob}[S]{1052}
$x^{\frac{2}{3}} + y^{\frac{2}{3}} = a^{\frac{2}{3}}$（星形线）．
\end{prob}

\begin{prob}[S]{1053}
$\arctan \dfrac{y}{x} = \ln \sqrt{x^2 + y^2}$（对数螺线）．
\end{prob}

\begin{prob}[L]{1054}
设 $y'_x$，设：
\begin{enumerate}[label=(\alph*), itemsep=2pt, topsep=3pt]
  \item $r = a\varphi$（阿基米德螺线）；
  \item $r = a(1 + \cos\varphi)$（心脏线）；
  \item $r = a\mathrm{e}^{m\varphi}$（对数螺线），
\end{enumerate}
其中 $r = \sqrt{x^2 + y^2}$ 及 $\varphi = \arctan \dfrac{y}{x}$ 是极坐标．
\end{prob}

\probsec{§3. 导数的几何意义}

\begin{prob}[L]{1055}
写出曲线
\[
y = (x + 1)\sqrt[3]{3 - x}
\]
在下列各点处的切线和法线方程：
\begin{enumerate}[label=(\alph*), itemsep=2pt, topsep=3pt]
  \item $A(-1, 0)$；
  \item $B(2, 3)$；
  \item $C(3, 0)$．
\end{enumerate}
\end{prob}

\begin{prob}[M]{1056}
在曲线
\[
y = 2 + x - x^2
\]
上的哪些点，其切线 (a) 平行于 $Ox$ 轴；(b) 平行于第一象限角的平分线？
\end{prob}

\begin{prob}[M]{1057}
证明：抛物线
\[
y = a(x - x_1)(x - x_2) \quad (a \neq 0, x_1 < x_2)
\]
与 $Ox$ 轴相交所成的两角 $\alpha$ 及 $\beta$ $\left(0 < \alpha < \dfrac{\pi}{2}, 0 < \beta < \dfrac{\pi}{2}\right)$ 彼此相等．
\end{prob}

\begin{prob}[M]{1058}
在曲线
\[
y = 2\sin x \quad (-\pi \leqslant x \leqslant \pi)
\]
上求出“曲线的坡度”（即 $|y'|$）大于 $1$ 的部分.
\end{prob}

\begin{prob}[M]{1059}
函数
\[
y = x \quad \text{及} \quad y_1 = x + 0.01\sin 1000\pi x
\]
二者相差不大于 $0.01$. 这些函数的导数的差的最大值如何？作出相应的图像.
\end{prob}

\begin{prob}[S]{1060}
曲线 $y = \ln x$ 与 $Ox$ 轴相交的角如何？
\end{prob}

\begin{prob}[S]{1061}
曲线 $y = x^2$ 与 $x = y^2$ 相交的角如何？
\end{prob}

\begin{prob}[S]{1062}
曲线 $y = \sin x$ 与 $y = \cos x$ 相交的角如何？
\end{prob}

\begin{prob}[M]{1063.1}
如何选择参数 $n$，可使曲线
\[
y = \arctan nx \quad (n > 0)
\]
与 $Ox$ 轴相交所成的角大于 $89^\circ$？
\end{prob}

\begin{prob}[M]{1063.2}
证明：曲线 $y = |x|^\alpha$ 在不同情形下与下列坐标轴相切：
(a) 当 $0 < \alpha < 1$ 时与 $Oy$ 轴相切；\qquad (b) 当 $1 < \alpha < +\infty$ 时与 $Ox$ 轴相切.
\end{prob}

\begin{prob}[M]{1063.3}
证明：对于函数
\[
y = \begin{cases}
|x|^\alpha, & x \neq 0, \, \alpha \neq 0, \\
1, & x = 0
\end{cases}
\]
的图像，经过点 $A(0, 1)$ 的割线的极限位置是 $Oy$ 轴.
\end{prob}

\begin{prob}[M]{1064}
求出曲线：(a) $y = \sqrt{1 - \mathrm{e}^{-a^2 x^2}}$ 于点 $x = 0$ 处；(b) $y = \arcsin \dfrac{2x}{1 + x^2}$ 于点 $x = 1$ 处的左切线与右切线间的夹角.
\end{prob}

\begin{prob}[M]{1065}
证明：对数螺线 $r = a\mathrm{e}^{m\varphi}$（$a$ 及 $m$ 为常数）的切线与切点的径向量所成的角度为一常量.
\end{prob}

\begin{prob}[M]{1066}
求曲线 $y = ax^n$ 的次切线长，由此给出作此曲线的切线的方法.
\end{prob}

\begin{prob}[M]{1067}
证明：对于抛物线 $y^2 = 2px$，(a) 次切线长等于切点的横坐标（绝对值）的两倍；(b) 次法线为一常量. 给出作抛物线的切线的方法.
\end{prob}

\begin{prob}[M]{1068}
证明：指数曲线 $y = a^x$ ($a > 0$) 有定长的次切线. 给出作指数曲线的切线的方法.
\end{prob}

\begin{prob}[M]{1069}
求悬链线
\[
y = a\cosh \frac{x}{a}
\]
上任何点 $M(x_0, y_0)$ 处的法线长.
\end{prob}

\begin{prob}[M]{1070}
证明：星形线
\[
x^{\frac{2}{3}} + y^{\frac{2}{3}} = a^{\frac{2}{3}} \quad (a > 0)
\]
的切线介于坐标轴间的部分的长为一常量.
\end{prob}

\begin{prob}[S]{1071}
若抛物线 $y = ax^2 + bx + c$ 与 $Ox$ 轴相切，则系数 $a, b, c$ 间的关系如何？
\end{prob}

\begin{prob}[M]{1072}
在什么条件下立方抛物线
\[
y = x^3 + px + q
\]
与 $Ox$ 轴相切？
\end{prob}

\begin{prob}[S]{1073}
当参数 $a$ 为何值时，抛物线 $y = ax^2$ 与曲线 $y = \ln x$ 相切？
\end{prob}

\begin{prob}[M]{1074}
证明：曲线
\[
y = f(x) \quad (f(x) > 0)
\]
和
\[
y = f(x)\sin ax
\]
于公共点彼此相切，其中 $f(x)$ 为可微函数.
\end{prob}

\begin{prob}[M]{1075}
证明：双曲线族 $x^2 - y^2 = a$ 及 $xy = b$ 形成一正交网，即这两族曲线成直角相交.
\end{prob}

\begin{prob}[M]{1076}
证明：抛物线族
\[
y^2 = 4a(a - x) \quad (a > 0)
\]
及
\[
y^2 = 4b(b + x) \quad (b > 0)
\]
形成正交网.
\end{prob}

\begin{prob}[M]{1077}
写出曲线 $x = 2t - t^2$ 及 $y = 3t - t^3$ 在下列各点的切线和法线的方程：
(a) $t = 0$；\qquad (b) $t = 1$.
\end{prob}

\begin{prob}[L]{1078}
写出曲线
\[
x = \frac{2t + t^2}{1 + t^3}, \quad y = \frac{2t - t^2}{1 + t^3}
\]
在下列各点的切线与法线方程：(a) $t = 0$，(b) $t = 1$，(c) $t = \infty$.
\end{prob}

\begin{prob}[M]{1079}
写出摆线（旋轮线）
\[
x = a(t - \sin t), \quad y = a(1 - \cos t)
\]
上任意一点 $t = t_0$ 处的切线方程. 给出摆线的切线的作法.
\end{prob}

\begin{prob}[M]{1080}
证明：曳物线
\[
x = a\left(\ln \tan \frac{t}{2} + \cos t\right), \quad y = a\sin t \quad (a > 0, \, 0 < t < \pi)
\]
有长度不变的切线.
\end{prob}

\probgroup{写出下列曲线在指定点的切线与法线方程：}

\begin{prob}[S]{1081}
$\dfrac{x^2}{100} + \dfrac{y^2}{64} = 1, \; M(6, 6.4)$.
\end{prob}

\begin{prob}[S]{1082}
$xy + \ln y = 1, \; M(1, 1)$.
\end{prob}

\probsec{§4. 函数的微分}

\begin{prob}[L]{1083}
设
(a) $\Delta x = 1$；\qquad (b) $\Delta x = 0.1$；\qquad (c) $\Delta x = 0.01$.
对于函数
\[
f(x) = x^3 - 2x + 1,
\]
求：(1) $\Delta f(1)$，(2) $\mathrm{d}f(1)$，并进行对比.
\end{prob}

\begin{prob}[L]{1084}
设运动方程由公式
\[
x = 5t^2
\]
给出，其中 $t$ 以 $\mathrm{s}$（秒）来度量，$x$ 以 $\mathrm{m}$（米）来度量. 若 (a) $\Delta t = 1\,\mathrm{s}$，(b) $\Delta t = 0.1\,\mathrm{s}$，(c) $\Delta t = 0.001\,\mathrm{s}$，对于 $t = 2\,\mathrm{s}$ 的时刻，求路程的增量 $\Delta x$ 及微分 $\mathrm{d}x$，并作比较.
\end{prob}

\probgroup{求下列函数 $y$ 的微分：}

\begin{prob}[S]{1085}
$y = \dfrac{1}{x}$.
\end{prob}

\begin{prob}[S]{1086}
$y = \dfrac{1}{a}\arctan \dfrac{x}{a} \; (a \neq 0)$.
\end{prob}

\begin{prob}[S]{1087}
$y = \dfrac{1}{2a}\ln\left|\dfrac{x - a}{x + a}\right|$.
\end{prob}

\begin{prob}[S]{1088}
$y = \ln\left|x + \sqrt{x^2 + a}\right|$.
\end{prob}

\begin{prob}[S]{1089}
$y = \arcsin \dfrac{x}{a} \; (a \neq 0)$.
\end{prob}

\begin{prob}[L]{1090}
求：
\begin{enumerate}[label=(\alph*), itemsep=2pt, topsep=3pt]
\item $\mathrm{d}(x\mathrm{e}^x)$；
\item $\mathrm{d}(\sin x - x\cos x)$；
\item $\mathrm{d}\left(\dfrac{1}{x^3}\right)$；
\item $\mathrm{d}\left(\dfrac{\ln x}{\sqrt{x}}\right)$；
\item $\mathrm{d}(\sqrt{a^2 + x^2})$；
\item $\mathrm{d}\left(\dfrac{x}{\sqrt{1 - x^2}}\right)$；
\item $\mathrm{d}\ln(1 - x^2)$；
\item $\mathrm{d}\left(\arccos \dfrac{1}{|x|}\right)$；
\item $\mathrm{d}\left[\dfrac{\sin x}{2\cos^2 x} + \dfrac{1}{2}\ln\left|\tan\left(\dfrac{x}{2} + \dfrac{\pi}{4}\right)\right|\right]$.
\end{enumerate}
\end{prob}

\probgroup{设 $u, v, w$ 为 $x$ 的可微函数，求下列函数 $y$ 的微分：}

\begin{prob}[S]{1091}
$y = uvw$.
\end{prob}

\begin{prob}[S]{1092}
$y = \dfrac{u}{v^2}$.
\end{prob}

\begin{prob}[S]{1093}
$y = \dfrac{1}{\sqrt{u^2 + v^2}}$.
\end{prob}

\begin{prob}[S]{1094}
$y = \arctan \dfrac{u}{v}$.
\end{prob}

\begin{prob}[S]{1095}
$y = \ln\sqrt{u^2 + v^2}$.
\end{prob}

\begin{prob}[L]{1096}
求：
\begin{enumerate}[label=(\alph*), itemsep=2pt, topsep=3pt]
\item $\dfrac{\mathrm{d}}{\mathrm{d}x^3}(x^3 - 2x^6 - x^9)$；
\item $\dfrac{\mathrm{d}}{\mathrm{d}x^2}\left(\dfrac{\sin x}{x}\right)$；
\item $\dfrac{\mathrm{d}(\sin x)}{\mathrm{d}(\cos x)}$；
\item $\dfrac{\mathrm{d}(\tan x)}{\mathrm{d}(\cot x)}$；
\item $\dfrac{\mathrm{d}(\arcsin x)}{\mathrm{d}(\arccos x)}$.
\end{enumerate}
\end{prob}

\begin{prob}[M]{1097}
有半径为 $R = 100\,\mathrm{cm}$ 及圆心角 $\alpha = 60^\circ$ 的扇形. 若 (a) 其半径 $R$ 增加 $1\,\mathrm{cm}$；(b) 角 $\alpha$ 减少 $30'$，则扇形面积的变化如何？求出精确的和近似的解.
\end{prob}

\begin{prob}[M]{1098}
摆的振动周期按下式确定：
\[
T = 2\pi\sqrt{\frac{l}{g}},
\]
其中 $l$ 为摆长，$g = 9.8\,\mathrm{m/s^2}$ 为重力加速度.
为了使周期 $T$ 增大 $0.05\,\mathrm{s}$，对摆长 $l = 20\,\mathrm{cm}$ 需要作多少修改？
\end{prob}

\probgroup{利用函数之微分代替其增量，求下列各式之近似值：}

\begin{prob}[S]{1099}
$\sqrt[3]{1.02}$.
\end{prob}

\begin{prob}[S]{1100}
$\sin 29^\circ$.
\end{prob}

\begin{prob}[S]{1101}
$\cos 151^\circ$.
\end{prob}

\begin{prob}[S]{1102}
$\arctan 1.05$.
\end{prob}

\begin{prob}[S]{1103}
$\lg 11$.
\end{prob}

\begin{prob}[L]{1104.1}
证明近似公式：
\[
\sqrt{a^2 + x} \approx a + \frac{x}{2a} \quad (a > 0),
\]
其中 $|x| \ll a$（正数 $A$ 和 $B$ 间的关系式 $A \ll B$ 表示 $A$ 远小于 $B$）. 利用此公式近似地计算：
(a) $\sqrt{5}$；\qquad (b) $\sqrt{34}$；\qquad (c) $\sqrt{120}$，
并与平方根表中的数值作比较.
\end{prob}

\begin{prob}[M]{1104.2}
证明公式：
\[
\sqrt{a^2 + x} = a + \frac{x}{2a} - r \quad (a > 0, \, x > 0),
\]
其中 $0 < r < \dfrac{x^2}{8a^3}$.
\end{prob}

\begin{prob}[L]{1105}
证明近似公式：
\[
\sqrt[n]{a^n + x} \approx a + \frac{x}{na^{n-1}} \quad (a > 0),
\]
其中 $|x| \ll a$. 利用此公式近似地计算：
(a) $\sqrt[3]{9}$；\qquad (b) $\sqrt[4]{80}$；\qquad (c) $\sqrt[7]{100}$；\qquad (d) $\sqrt[10]{1000}$.
\end{prob}

\begin{prob}[M]{1106}
正方形的边 $x = 2.4\,\mathrm{m} \pm 0.05\,\mathrm{m}$. 由此计算所得正方形的面积的相对误差和绝对误差如何？
\end{prob}

\begin{prob}[M]{1107}
为了在 $1\%$ 的精度下计算出球的体积，问度量球半径 $R$ 时所允许的相对误差如何？
\end{prob}

\begin{prob}[M]{1108}
为了确定重力加速度，可以利用摆的振动公式 $g = \dfrac{4\pi^2 l}{T^2}$，其中 $l$ 为摆长，$T$ 为振动周期. 当测量 (a) 长度 $l$，(b) 周期 $T$ 时，相对误差 $\delta$ 如何影响 $g$ 的值？
\end{prob}

\begin{prob}[M]{1109}
求数 $x$ ($x > 0$) 的常用对数（以 $10$ 为底）的绝对误差，设此数的相对误差等于 $\delta$.
\end{prob}

\begin{prob}[M]{1110}
证明：根据正切对数表求得的角度比根据具有同样多位小数的正弦对数表求得的角度更为精确.
\end{prob}

\probsec{§5. 高阶的导数和微分}

\probgroup{求 $y''$，设：}

\begin{prob}[S]{1111}
$y = x\sqrt{1 + x^2}$.
\end{prob}

\begin{prob}[S]{1112}
$y = \dfrac{x}{\sqrt{1 - x^2}}$.
\end{prob}

\begin{prob}[S]{1113}
$y = \mathrm{e}^{-x^2}$.
\end{prob}

\begin{prob}[S]{1114}
$y = \tan x$.
\end{prob}

\begin{prob}[S]{1115}
$y = (1 + x^2)\arctan x$.
\end{prob}

\begin{prob}[S]{1116}
$y = \frac{\arcsin x}{\sqrt{1 - x^2}}$.
\end{prob}

\begin{prob}[S]{1117}
$y = x\ln x$.
\end{prob}

\begin{prob}[S]{1118}
$y = \ln f(x)$.
\end{prob}

\begin{prob}[S]{1119}
$y = x[\sin(\ln x) + \cos(\ln x)]$.
\end{prob}

\begin{prob}[M]{1120}
设 $y = \mathrm{e}^{\sin x}\cos(\sin x)$, 求 $y(0)$, $y'(0)$ 及 $y''(0)$.
\end{prob}

\probgroup{设 $u = \varphi(x)$ 及 $v = \psi(x)$ 为二阶可微函数. 求 $y''$, 设:}

\begin{prob}[S]{1121}
$y = u^2$.
\end{prob}

\begin{prob}[S]{1122}
$y = \ln\frac{u}{v}$.
\end{prob}

\begin{prob}[S]{1123}
$y = \sqrt{u^2 + v^2}$.
\end{prob}

\begin{prob}[S]{1124}
$y = u^v$ ($u > 0$).
\end{prob}

\probgroup{设 $f(x)$ 为三阶可微函数. 求 $y''$ 及 $y'''$, 设:}

\begin{prob}[S]{1125}
$y = f(x^2)$.
\end{prob}

\begin{prob}[S]{1126}
$y = f\left(\frac{1}{x}\right)$.
\end{prob}

\begin{prob}[S]{1127}
$y = f(\mathrm{e}^x)$.
\end{prob}

\begin{prob}[S]{1128}
$y = f(\ln x)$.
\end{prob}

\begin{prob}[S]{1129}
$y = f(\varphi(x))$, 其中 $\varphi(x)$ 是充分多次可微函数.
\end{prob}

\begin{prob}[M]{1130}
对于以下两种情形: (a) $x$ 为自变量, (b) $x$ 为中间变量, 求函数 $y = \mathrm{e}^x$ 的 $\mathrm{d}^2 y$.
\end{prob}

\probgroup{若 $x$ 为自变量, 求 $\mathrm{d}^2 y$, 设:}

\begin{prob}[S]{1131}
$y = \sqrt{1 + x^2}$.
\end{prob}

\begin{prob}[S]{1132}
$y = \frac{\ln x}{x}$.
\end{prob}

\begin{prob}[S]{1133}
$y = x^x$.
\end{prob}

\probgroup{令 $u$ 及 $v$ 为变量 $x$ 的二阶可微函数, 求 $\mathrm{d}^2 y$, 设:}

\begin{prob}[S]{1134}
$y = uv$.
\end{prob}

\begin{prob}[S]{1135}
$y = \frac{u}{v}$.
\end{prob}

\begin{prob}[S]{1136}
$y = u^m v^n$ ($m$ 及 $n$ 为常数).
\end{prob}

\begin{prob}[S]{1137}
$y = a^u$ ($a > 0$).
\end{prob}

\begin{prob}[S]{1138}
$y = \ln\sqrt{u^2 + v^2}$.
\end{prob}

\begin{prob}[S]{1139}
$y = \arctan\frac{u}{v}$.
\end{prob}

\probgroup{对以参数形式给出的函数 $y = y(x)$ 的导数 $y'_x$, $y''_{x^2}$, $y'''_{x^3}$, 设:}

\begin{prob}[S]{1140}
$x = 2t - t^2$, $y = 3t - t^3$.
\end{prob}

\begin{prob}[S]{1141}
$x = a\cos t$, $y = a\sin t$.
\end{prob}

\begin{prob}[S]{1142}
$x = a(t - \sin t)$, $y = a(1 - \cos t)$.
\end{prob}

\begin{prob}[S]{1143}
$x = \mathrm{e}^t \cos t$, $y = \mathrm{e}^t \sin t$.
\end{prob}

\begin{prob}[S]{1144}
$x = f'(t)$, $y = t f'(t) - f(t)$.
\end{prob}

\begin{prob}[M]{1145}
设函数 $y = f(x)$ 充分多次可微. 求反函数 $x = f^{-1}(y)$ 的导数 $x'$, $x''$, $x'''$, $x^{(4)}$ (设这些导数都存在).
\end{prob}

\probgroup{求由下列隐函数给出的 $y = y(x)$ 的 $y'_x$, $y''_{x^2}$ 及 $y'''_{x^3}$:}

\begin{prob}[M]{1146}
$x^2 + y^2 = 25$. 在点 $M(3, 4)$ 的 $y'$, $y''$ 及 $y'''$ 等于什么?
\end{prob}

\begin{prob}[S]{1147}
$y^2 = 2px$.
\end{prob}

\begin{prob}[S]{1148}
$x^2 - xy + y^2 = 1$.
\end{prob}

\probgroup{求 $y'_x$ 及 $y''_{x^2}$, 设:}

\begin{prob}[S]{1149}
$y^2 + 2\ln y = x^4$.
\end{prob}

\begin{prob}[S]{1150}
$\sqrt{x^2 + y^2} = a\mathrm{e}^{\arctan\frac{y}{x}}$ ($a > 0$).
\end{prob}

\begin{prob}[M]{1151}
设函数 $f(x)$ 当 $x \leqslant x_0$ 时有定义且二阶可微. 应当如何选择系数 $a$, $b$ 及 $c$, 使函数
\[
F(x) = \begin{cases}
f(x), & x \leqslant x_0, \\
a(x - x_0)^2 + b(x - x_0) + c, & x > x_0
\end{cases}
\]
是二阶可微的函数?
\end{prob}

\begin{prob}[M]{1152}
质点作直线运动的规律为
\[
s = 10 + 20t - 5t^2.
\]
求其运动速度和加速度. 在 $t = 2$ 的时刻, 速度与加速度等于什么?
\end{prob}

\begin{prob}[M]{1153}
质点 $M(x, y)$ 沿圆周 $x^2 + y^2 = a^2$ 匀速运动, 运动一周的时间为 $T$. 若质点在时刻 $t = 0$ 位于点 $M_0(a, 0)$, 求质点的速度和加速度在 $Ox$ 轴上的投影 $v$ 和 $w$.
\end{prob}

\begin{prob}[L]{1154}
在重力场中, 质点 $M(x, y)$ 在竖直平面 $Oxy$ 内以初速度 $v_0$ 沿与水平面成 $\alpha$ 角的方向抛出. 建立运动方程 (忽略空气阻力), 并确定速度 $v$ 和加速度 $w$ 的大小及运动轨迹. 质点的最大上升高度和射程等于多少?
\end{prob}

\begin{prob}[M]{1155}
质点的运动方程为
\[
x = 4\sin\omega t - 3\cos\omega t, \quad y = 3\sin\omega t + 4\cos\omega t \quad (\omega \text{ 为常数}).
\]
求运动轨迹、速度与加速度的大小.
\end{prob}

\probgroup{求下列指定阶的导数:}

\begin{prob}[S]{1156}
$y = x(2x - 1)^2(x + 3)^3$, 求 $y^{(6)}$ 及 $y^{(7)}$.
\end{prob}

\begin{prob}[S]{1157}
$y = \frac{a}{x^m}$, 求 $y'''$.
\end{prob}

\begin{prob}[S]{1158}
$y = \sqrt{x}$, 求 $y^{(10)}$.
\end{prob}

\begin{prob}[S]{1159}
$y = \frac{x^2}{1 - x}$, 求 $y^{(8)}$.
\end{prob}

\begin{prob}[S]{1160}
$y = \frac{1 + x}{\sqrt{1 - x}}$, 求 $y^{(100)}$.
\end{prob}

\begin{prob}[S]{1161}
$y = x^2\mathrm{e}^{2x}$, 求 $y^{(20)}$.
\end{prob}

\begin{prob}[S]{1162}
$y = \frac{\mathrm{e}^x}{x}$, 求 $y^{(10)}$.
\end{prob}

\begin{prob}[S]{1163}
$y = x\ln x$, 求 $y^{(5)}$.
\end{prob}

\begin{prob}[S]{1164}
$y = \frac{\ln x}{x}$, 求 $y^{(5)}$.
\end{prob}

\begin{prob}[S]{1165}
$y = x^2\sin 2x$, 求 $y^{(50)}$.
\end{prob}

\begin{prob}[S]{1166}
$y = \frac{\cos 3x}{\sqrt[3]{1 - 3x}}$, 求 $y'''$.
\end{prob}

\begin{prob}[S]{1167}
$y = \sin x\sin 2x\sin 3x$, 求 $y^{(10)}$.
\end{prob}

\begin{prob}[S]{1168}
$y = x\sinh x$, 求 $y^{(100)}$.
\end{prob}

\begin{prob}[S]{1169}
$y = \mathrm{e}^x\cos x$, 求 $y^{(4)}$.
\end{prob}

\begin{prob}[S]{1170}
$y = \sin^2 x\ln x$, 求 $y^{(6)}$.
\end{prob}

\probgroup{在下列各例中, 视 $x$ 为自变量, 求指定阶的微分.}

\begin{prob}[S]{1171}
$y = x^5$, 求 $\mathrm{d}^5 y$.
\end{prob}

\begin{prob}[S]{1172}
$y = \frac{1}{\sqrt{x}}$, 求 $\mathrm{d}^3 y$.
\end{prob}

\begin{prob}[S]{1173}
$y = x\cos 2x$, 求 $\mathrm{d}^{10} y$.
\end{prob}

\begin{prob}[S]{1174}
$y = \mathrm{e}^x\ln x$, 求 $\mathrm{d}^4 y$.
\end{prob}

\begin{prob}[S]{1175}
$y = \cos x\cdot\cosh x$, 求 $\mathrm{d}^6 y$.
\end{prob}

\probgroup{设 $u$ 为 $x$ 的充分多次可微函数, 在下列各例中求指定阶的微分.}

\begin{prob}[S]{1176}
$y = u^2$, 求 $\mathrm{d}^{10} y$.
\end{prob}

\begin{prob}[S]{1177}
$y = \mathrm{e}^u$, 求 $\mathrm{d}^4 y$.
\end{prob}

\begin{prob}[S]{1178}
$y = \ln u$, 求 $\mathrm{d}^3 y$.
\end{prob}

\begin{prob}[M]{1179}
视 $x$ 为某个自变量的函数, 由函数 $y = f(x)$ 求 $\mathrm{d}^2 y$, $\mathrm{d}^3 y$ 及 $\mathrm{d}^4 y$.
\end{prob}

\begin{prob}[M]{1180}
以变量 $x$ 和 $y$ 的逐次微分来表示函数 $y = f(x)$ 的导数 $y''$ 及 $y'''$, 但不假定 $x$ 为自变量.
\end{prob}

\begin{prob}[M]{1181}
证明: 函数
\[
y = C_1\cos x + C_2\sin x,
\]
其中 $C_1$ 及 $C_2$ 为任意的常数, 满足方程
\[
y'' + y = 0.
\]
\end{prob}

\begin{prob}[M]{1182}
证明: 函数
\[
y = C_1\cosh x + C_2\sinh x,
\]
其中 $C_1$ 及 $C_2$ 为任意的常数, 满足方程
\[
y'' - y = 0.
\]
\end{prob}

\begin{prob}[M]{1183}
证明: 函数
\[
y = C_1\mathrm{e}^{\lambda_1 x} + C_2\mathrm{e}^{\lambda_2 x},
\]
其中 $C_1$ 及 $C_2$ 为任意的常数, $\lambda_1$ 及 $\lambda_2$ 为常数, 满足方程
\[
y'' - (\lambda_1 + \lambda_2)y' + \lambda_1\lambda_2 y = 0.
\]
\end{prob}

\begin{prob}[M]{1184}
证明: 函数
\[
y = x^n[C_1\cos(\ln x) + C_2\sin(\ln x)],
\]
其中 $C_1$ 及 $C_2$ 为任意的常数, $n$ 为常数, 满足方程
\[
x^2 y'' + (1 - 2n)xy' + (1 + n^2)y = 0.
\]
\end{prob}

\begin{prob}[L]{1185}
证明: 函数
\[
y = \mathrm{e}^{\frac{x}{\sqrt{2}}}\left(C_1\cos\frac{x}{\sqrt{2}} + C_2\sin\frac{x}{\sqrt{2}}\right) + \mathrm{e}^{-\frac{x}{\sqrt{2}}}\left(C_3\cos\frac{x}{\sqrt{2}} + C_4\sin\frac{x}{\sqrt{2}}\right),
\]
其中 $C_1, C_2, C_3$ 及 $C_4$ 为任意常数, 满足方程
\[
y^{(4)} + y = 0.
\]
\end{prob}

\begin{prob}[M]{1186}
证明: 若函数 $f(x)$ 有 $n$ 阶导数, 则
\[
[f(ax + b)]^{(n)} = a^n f^{(n)}(ax + b).
\]
\end{prob}

\begin{prob}[M]{1187}
若
\[
P(x) = a_0 x^n + a_1 x^{n-1} + \cdots + a_n,
\]
求 $P^{(n)}(x)$.
\end{prob}

\probgroup{求 $y^{(n)}$, 设:}

\begin{prob}[S]{1188}
$y = \frac{ax + b}{cx + d}$.
\end{prob}

\begin{prob}[S]{1189}
$y = \frac{1}{x(1 - x)}$.
\end{prob}

\begin{prob}[M]{1190}
$y = \frac{1}{x^2 - 3x + 2}$.
\probhint{分解函数为最简分式.}
\end{prob}

\begin{prob}[S]{1191}
$y = \frac{1}{\sqrt{1 - 2x}}$.
\end{prob}

\begin{prob}[S]{1192}
$y = \frac{x}{\sqrt[3]{1 + x}}$.
\end{prob}

\begin{prob}[S]{1193}
$y = \sin^2 x$.
\end{prob}

\begin{prob}[S]{1194}
$y = \cos^2 x$.
\end{prob}

\begin{prob}[S]{1195}
$y = \sin^3 x$.
\end{prob}

\begin{prob}[S]{1196}
$y = \cos^3 x$.
\end{prob}

\begin{prob}[S]{1197}
$y = \sin ax\sin bx$.
\end{prob}

\begin{prob}[S]{1198}
$y = \cos ax\cos bx$.
\end{prob}

\begin{prob}[S]{1199}
$y = \sin ax\cos bx$.
\end{prob}

\begin{prob}[S]{1200}
$y = \sin^2 ax\cos bx$.
\end{prob}

\begin{prob}[S]{1201}
$y = \sin^4 x + \cos^4 x$.
\end{prob}

\begin{prob}[S]{1202}
$y = x\cos ax$.
\end{prob}

\begin{prob}[S]{1203}
$y = x^2\sin ax$.
\end{prob}

\begin{prob}[S]{1204}
$y = (x^2 + 2x + 2)\mathrm{e}^{-x}$.
\end{prob}

\begin{prob}[S]{1205}
$y = \frac{\mathrm{e}^x}{x}$.
\end{prob}

\begin{prob}[S]{1206}
$y = \mathrm{e}^x\cos x$.
\end{prob}

\begin{prob}[S]{1207}
$y = \mathrm{e}^x\sin x$.
\end{prob}

\begin{prob}[S]{1208}
$y = \ln\frac{a + bx}{a - bx}$.
\end{prob}

\begin{prob}[S]{1209}
$y = \mathrm{e}^{ax} P(x)$, 其中 $P(x)$ 为多项式.
\end{prob}

\begin{prob}[S]{1210}
$y = x\sinh x$.
\end{prob}

\probgroup{求 $\mathrm{d}^n y$, 设:}

\begin{prob}[S]{1211}
$y = x^n\mathrm{e}^x$.
\end{prob}

\begin{prob}[S]{1212}
$y = \frac{\ln x}{x}$.
\end{prob}

\begin{prob}[L]{1213}
证明等式:
\begin{enumerate}[label=(\arabic*), itemsep=2pt, topsep=3pt]
  \item $[\mathrm{e}^{ax}\sin(bx + c)]^{(n)} = \mathrm{e}^{ax}(a^2 + b^2)^{\frac{n}{2}}\sin(bx + c + n\varphi)$,
  \item $[\mathrm{e}^{ax}\cos(bx + c)]^{(n)} = \mathrm{e}^{ax}(a^2 + b^2)^{\frac{n}{2}}\cos(bx + c + n\varphi)$,
\end{enumerate}
其中 $\sin\varphi = \frac{b}{\sqrt{a^2 + b^2}}$, $\cos\varphi = \frac{a}{\sqrt{a^2 + b^2}}$.
\end{prob}

\begin{prob}[M]{1214}
求 $y^{(n)}$, 设:
(a) $y = \cosh ax\cos bx$; \qquad (b) $y = \cosh ax\sin bx$.
\end{prob}

\begin{prob}[L]{1215}
将函数
\[
f(x) = \sin^{2p} x,
\]
其中 $p$ 为正整数, 化为三角多项式:
\[
f(x) = \sum_{k=0}^{p} A_k\cos 2kx,
\]
以求 $f^{(n)}(x)$.
\probhint{令 $\sin x = \frac{1}{2\mathrm{i}}(t - \bar{t})$, 其中 $t = \cos x + \mathrm{i}\sin x$ 及 $\bar{t} = \cos x - \mathrm{i}\sin x$, 并且利用棣莫弗公式来计算.}
\end{prob}

\begin{prob}[L]{1216}
设:
\begin{enumerate}[label=(\alph*), itemsep=2pt, topsep=3pt]
  \item $f(x) = \sin^{2p+1} x$;
  \item $f(x) = \cos^{2p} x$;
  \item $f(x) = \cos^{2p+1} x$,
\end{enumerate}
其中 $p$ 为正整数, 求 $f^{(n)}(x)$ (参阅前题).
\end{prob}

\probgroup{若
\[
f(x) = f_1(x) + \mathrm{i} f_2(x),
\]
其中 $\mathrm{i}$ 是虚数单位, $f_1(x)$ 和 $f_2(x)$ 为实变量 $x$ 的实函数, 则按定义有:
\[
f'(x) = f'_1(x) + \mathrm{i} f'_2(x).
\]}

\begin{prob}[L]{1217}
利用恒等式
\[
\frac{1}{x^2 + 1} = \frac{1}{2\mathrm{i}}\left(\frac{1}{x - \mathrm{i}} - \frac{1}{x + \mathrm{i}}\right)
\]
证明:
\[
\left(\frac{1}{x^2 + 1}\right)^{(n)} = \frac{(-1)^n n!}{(1 + x^2)^{\frac{n+1}{2}}}\sin[(n + 1)\operatorname{arccot} x].
\]
\probhint{利用棣莫弗公式.}
\end{prob}

\begin{prob}[M]{1218}
求函数 $f(x) = \arctan x$ 的 $n$ 阶导数.
\end{prob}

\probgroup{求 $f^{(n)}(0)$, 设:}

\begin{prob}[M]{1219}
(a) $f(x) = \frac{1}{(1 - 2x)(1 + x)}$; \qquad (b) $f(x) = \frac{x}{\sqrt{1 - x}}$.
\end{prob}

\begin{prob}[L]{1220}
\begin{enumerate}[label=(\alph*), itemsep=2pt, topsep=3pt]
  \item $f(x) = x^2\mathrm{e}^{ax}$;
  \item $f(x) = \arctan x$;
  \item $f(x) = \arcsin x$.
\end{enumerate}
\end{prob}

\begin{prob}[M]{1221}
(a) $f(x) = \cos(m\arcsin x)$; \qquad (b) $f(x) = \sin(m\arcsin x)$.
\end{prob}

\begin{prob}[M]{1222}
(a) $f(x) = (\arctan x)^2$; \qquad (b) $f(x) = (\arcsin x)^2$.
\end{prob}

\begin{prob}[M]{1223}
设
\[
f(x) = (x - a)^n\varphi(x),
\]
其中函数 $\varphi(x)$ 在点 $a$ 的邻域内有 $n - 1$ 阶的连续导数, 求 $f^{(n)}(a)$.
\end{prob}

\begin{prob}[M]{1224}
证明: 函数
\[
f(x) = \begin{cases}
x^{2n}\sin\dfrac{1}{x}, & x \neq 0, \\
0, & x = 0
\end{cases}
\]
($n$ 为正整数), 于点 $x = 0$ 有一直到 $n$ 阶的导数, 而无 $n + 1$ 阶导数.
\end{prob}

\begin{prob}[M]{1225}
证明: 函数
\[
f(x) = \begin{cases}
\mathrm{e}^{-1/x^2}, & x \neq 0, \\
0, & x = 0
\end{cases}
\]
在 $x = 0$ 处是无穷次可微的. 作出此函数的图像.
\end{prob}

\begin{prob}[M]{1226}
证明: 切比雪夫多项式
\[
T_m(x) = \frac{1}{2^{m-1}}\cos(m\arccos x) \quad (m = 1, 2, \cdots)
\]
满足方程
\[
(1 - x^2)T''_m(x) - xT'_m(x) + m^2 T_m(x) = 0.
\]
\end{prob}

\begin{prob}[L]{1227}
证明: 勒让德多项式
\[
P_m(x) = \frac{1}{2^m m!}[(x^2 - 1)^m]^{(m)} \quad (m = 0, 1, 2, \cdots)
\]
满足方程
\[
(1 - x^2)P''_m(x) - 2xP'_m(x) + m(m + 1)P_m(x) = 0.
\]
\probhint{把等式 $(x^2 - 1)u' = 2mxu$ 作 $m + 1$ 次微分, 其中 $u = (x^2 - 1)^m$.}
\end{prob}

\begin{prob}[L]{1228}
切比雪夫-拉盖尔多项式定义如下:
\[
L_m(x) = \mathrm{e}^x(x^m\mathrm{e}^{-x})^{(m)} \quad (m = 0, 1, 2, \cdots).
\]
求多项式 $L_m(x)$ 的显式表达式.
证明: $L_m(x)$ 满足方程
\[
xL''_m(x) + (1 - x)L'_m(x) + mL(x) = 0.
\]
\probhint{利用等式 $xu' + (x - m)u = 0$，其中 $u = x^m \mathrm{e}^{-x}$.}
\end{prob}

\begin{prob}[M]{1229}
设 $y = f(u)$ 及 $u = \varphi(x)$，其中 $f(u)$ 及 $\varphi(x)$ 为 $n$ 阶可微函数. 证明：
\[
\frac{\mathrm{d}^n y}{\mathrm{d}x^n} = \sum_{k=1}^{n} A_k(x) f^{(k)}(u),
\]
其中系数 $A_k(x)$ $(k = 0, 1, \cdots, n)$ 与函数 $f(u)$ 无关.
\end{prob}

\begin{prob}[M]{1230}
证明：对于复合函数 $y = f(x^2)$ 的 $n$ 阶导数，成立公式
\[
\frac{\mathrm{d}^n y}{\mathrm{d}x^n} = (2x)^n f^{(n)}(x^2) + \frac{n(n-1)}{1!}(2x)^{n-2} f^{(n-1)}(x^2) + \frac{n(n-1)(n-2)(n-3)}{2!}(2x)^{n-4} f^{(n-2)}(x^2) + \cdots.
\]
\end{prob}

\begin{prob}[L]{1231}
切比雪夫--埃尔米特多项式定义如下：
\[
H_m(x) = (-1)^m \mathrm{e}^{x^2} (\mathrm{e}^{-x^2})^{(m)} \quad (m = 0, 1, 2, \cdots).
\]
求多项式 $H_m(x)$ 的显式表达式.

证明：$H_m(x)$ 满足方程
\[
H_m''(x) - 2x H_m'(x) + 2m H_m(x) = 0.
\]
\probhint{利用等式 $u' + 2xu = 0$，其中 $u = \mathrm{e}^{-x^2}$.}
\end{prob}

\begin{prob}[M]{1232.1}
证明等式：
\[
\left(x^{n-1} \mathrm{e}^{\frac{1}{x}}\right)^{(n)} = \frac{(-1)^n}{x^{n+1}} \mathrm{e}^{\frac{1}{x}}.
\]
\probhint{运用数学归纳法.}
\end{prob}

\begin{prob}[M]{1232.2}
证明公式：
\begin{enumerate}[label=(\alph*), itemsep=2pt, topsep=3pt]
  \item $\dfrac{\mathrm{d}^n}{\mathrm{d}x^n}(x^n \ln x) = n!\left(\ln x + \sum\limits_{k=1}^{n} \dfrac{1}{k}\right) \quad (x > 0)$;
  \item $\dfrac{\mathrm{d}^{2n}}{\mathrm{d}x^{2n}}\left(\dfrac{\sin x}{x}\right) = \dfrac{(2n)!}{x^{2n+1}}[C_n(x)\sin x - S_n(x)\cos x]$，其中
  \[
  C_n(x) = 1 - \frac{x^2}{2!} + \cdots + (-1)^n \frac{x^{2n}}{(2n)!}, \quad S_n(x) = x - \frac{x^3}{3!} + \cdots + (-1)^{n-1} \frac{x^{2n-1}}{(2n-1)!}.
  \]
\end{enumerate}
\end{prob}

\begin{prob}[M]{1233}
设 $\dfrac{\mathrm{d}}{\mathrm{d}x} = D$ 表示微分算子，
\[
f(D) = \sum_{k=0}^{n} p_k(x) D^k
\]
为微分符号的多项式，其中 $p_k(x)$ $(k = 0, 1, \cdots, n)$ 为 $x$ 的某连续函数. 证明：
\[
f(D)\{\mathrm{e}^{\lambda x} u(x)\} = \mathrm{e}^{\lambda x} f(D + \lambda) u(x),
\]
其中 $\lambda$ 为常数.
\end{prob}

\begin{prob}[M]{1234}
证明：若在方程
\[
\sum_{k=0}^{n} a_k x^k y_x^{(k)} = 0
\]
中令
\[
x = \mathrm{e}^t,
\]
其中 $t$ 为自变量，则此方程化为
\[
\sum_{k=0}^{n} a_k D(D - 1)\cdots(D - k + 1) y = 0,
\]
其中 $D = \dfrac{\mathrm{d}}{\mathrm{d}t}$.
\end{prob}

\probsec{§6. 罗尔定理、拉格朗日定理和柯西定理}

\begin{prob}[S]{1235}
检验罗尔定理对于函数
\[
f(x) = (x - 1)(x - 2)(x - 3)
\]
的正确性.
\end{prob}

\begin{prob}[M]{1236}
函数 $f(x) = 1 - \sqrt[3]{x^2}$ 当 $x_1 = -1$ 及 $x_2 = 1$ 时为零，但是当 $-1 \leqslant x \leqslant 1$ 时，$f'(x) \neq 0$. 说明与罗尔定理表面上的矛盾.
\end{prob}

\begin{prob}[M]{1237}
设函数 $f(x)$ 在有限或无穷的区间 $(a, b)$ 中的任意一点有有限的导数 $f'(x)$，且
\[
\lim_{x \to a+0} f(x) = \lim_{x \to b-0} f(x).
\]
证明：
\[
f'(c) = 0,
\]
其中 $c$ 为区间 $(a, b)$ 中的某点.
\end{prob}

\begin{prob}[M]{1238}
设函数 $f(x)$：(1) 在闭区间 $[x_0, x_n]$ 上有定义且有 $n - 1$ 阶的连续导数 $f^{(n-1)}(x)$；(2) 在区间 $(x_0, x_n)$ 内有 $n$ 阶导数 $f^{(n)}(x)$；(3) 下列等式成立：
\[
f(x_0) = f(x_1) = \cdots = f(x_n) \quad (x_0 < x_1 < \cdots < x_n).
\]
证明：在区间 $(x_0, x_n)$ 内最少存在一点 $\xi$，使
\[
f^{(n)}(\xi) = 0.
\]
\end{prob}

\begin{prob}[M]{1239}
设函数 $f(x)$：(1) 在闭区间 $[a, b]$ 上有定义且有 $p + q$ 阶的连续导数 $f^{(p+q)}(x)$；(2) 在区间 $(a, b)$ 内有 $p + q + 1$ 阶的导数 $f^{(p+q+1)}(x)$；(3) 下列等式成立：
\[
f(a) = f'(a) = \cdots = f^{(p)}(a) = 0,
\]
\[
f(b) = f'(b) = \cdots = f^{(q)}(b) = 0.
\]
证明：在此种情形下
\[
f^{(p+q+1)}(c) = 0,
\]
其中 $c$ 为区间 $(a, b)$ 内的某点.
\end{prob}

\begin{prob}[M]{1240}
证明：若具实系数 $a_k$ $(k = 0, 1, \cdots, n)$ 的多项式
\[
P_n(x) = a_0 x^n + a_1 x^{n-1} + \cdots + a_n \quad (a_0 \neq 0)
\]
之一切根为实数，则其逐次的导数 $P_n'(x), P_n''(x), \cdots, P_n^{(n-1)}(x)$ 也仅有实根.
\end{prob}

\begin{prob}[M]{1241}
证明：勒让德多项式
\[
P_n(x) = \frac{1}{2^n n!} \frac{\mathrm{d}^n}{\mathrm{d}x^n}\{(x^2 - 1)^n\}
\]
的一切根都是实数且包含于区间 $(-1, 1)$ 中.
\end{prob}

\begin{prob}[M]{1242}
证明：切比雪夫--拉盖尔多项式
\[
L_n(x) = \mathrm{e}^x \frac{\mathrm{d}^n}{\mathrm{d}x^n}(x^n \mathrm{e}^{-x})
\]
所有的根都是正数.
\end{prob}

\begin{prob}[M]{1243}
证明：切比雪夫--埃尔米特多项式
\[
H_n(x) = (-1)^n \mathrm{e}^{x^2} \frac{\mathrm{d}^n}{\mathrm{d}x^n}(\mathrm{e}^{-x^2})
\]
所有的根都是实数.
\end{prob}

\begin{prob}[S]{1244}
在曲线 $y = x^3$ 上某点的切线，平行于连接点 $A(-1, -1)$ 及点 $B(2, 8)$ 所成的弦，求出此点.
\end{prob}

\begin{prob}[S]{1245}
若 $ab < 0$，有限增量公式对于函数
\[
f(x) = \frac{1}{x}
\]
在闭区间 $[a, b]$ 上是否正确？
\end{prob}

\begin{prob}[L]{1246.1}
设
\begin{enumerate}[label=(\alph*), itemsep=2pt, topsep=3pt]
  \item $f(x) = ax^2 + bx + c$ $(a \neq 0)$;
  \item $f(x) = x^3$;
  \item $f(x) = \dfrac{1}{x}$;
  \item $f(x) = \mathrm{e}^x$,
\end{enumerate}
求满足
\[
f(x + \Delta x) - f(x) = \Delta x f'(x + \theta \Delta x) \quad (0 < \theta < 1)
\]
的函数 $\theta = \theta(x, \Delta x)$.
\end{prob}

\begin{prob}[M]{1246.2}
设 $f(x) \in C^{(1)}(-\infty, +\infty)$，且对于任何 $x$ 和 $h$ 成立恒等式
\[
f(x + h) - f(x) \equiv h f'(x).
\]
证明：
\[
f(x) = ax + b,
\]
其中 $a$ 和 $b$ 为常数.
\end{prob}

\begin{prob}[M]{1246.3}
设 $f(x) \in C^{(2)}(-\infty, +\infty)$，且对于任何 $x$ 和 $h$ 成立恒等式
\[
f(x + h) - f(x) \equiv h f'\left(x + \frac{h}{2}\right).
\]
证明：
\[
f(x) = ax^2 + bx + c,
\]
其中 $a, b$ 和 $c$ 为常数.
\end{prob}

\begin{prob}[M]{1247}
证明：若 $x \geqslant 0$，则
\[
\sqrt{x + 1} - \sqrt{x} = \frac{1}{2\sqrt{x + \theta(x)}},
\]
其中 $\dfrac{1}{4} \leqslant \theta(x) \leqslant \dfrac{1}{2}$，并且 $\lim\limits_{x \to 0} \theta(x) = \dfrac{1}{4}$，$\lim\limits_{x \to +\infty} \theta(x) = \dfrac{1}{2}$.
\end{prob}

\begin{prob}[M]{1248}
设
\[
f(x) = \begin{cases}
\dfrac{3 - x^2}{2}, & 0 \leqslant x \leqslant 1, \\[4pt]
\dfrac{1}{x}, & 1 < x < +\infty,
\end{cases}
\]
在闭区间 $[0, 2]$ 上对于函数 $f(x)$ 求有限增量公式中的中间值 $c$.
\end{prob}

\begin{prob}[M]{1249}
设
\[
f(x) - f(0) = x f'[\xi(x)],
\]
其中 $0 < \xi(x) < x$. 证明：若
\[
f(x) = \begin{cases}
x \sin(\ln x), & x > 0, \\
0, & x = 0,
\end{cases}
\]
则函数 $\xi = \xi(x)$ 在任意小的区间 $(0, \varepsilon)$ 内 $(\varepsilon > 0)$ 是不连续的.
\end{prob}

\begin{prob}[M]{1250}
设函数 $f(x)$ 在区间 $(a, b)$ 内有连续的导数 $f'(x)$. 对于区间 $(a, b)$ 内任何一点 $\xi$，可否从此区间中指出另外的两点 $x_1$ 及 $x_2$，使得
\[
\frac{f(x_2) - f(x_1)}{x_2 - x_1} = f'(\xi) \quad (x_1 < \xi < x_2)?
\]
研究一个例子：
\[
f(x) = x^3 \quad (-1 \leqslant x \leqslant 1), \quad \xi = 0.
\]
\end{prob}

\begin{prob}[L]{1251}
证明下列不等式：
\begin{enumerate}[label=(\alph*), itemsep=2pt, topsep=3pt]
  \item $|\sin x - \sin y| \leqslant |x - y|$;
  \item $p y^{p-1}(x - y) \leqslant x^p - y^p \leqslant p x^{p-1}(x - y) \quad (0 < y < x, p > 1)$;
  \item $|\arctan a - \arctan b| \leqslant |a - b|$;
  \item $\dfrac{a - b}{a} < \ln \dfrac{a}{b} < \dfrac{a - b}{b} \quad (0 < b < a)$.
\end{enumerate}
\end{prob}

\begin{prob}[S]{1252}
说明在闭区间 $[-1, 1]$ 上柯西定理对于函数 $f(x) = x^2$ 及 $g(x) = x^3$ 何以不真？
\end{prob}

\begin{prob}[M]{1253}
设函数 $f(x)$ 在闭区间 $[x_1, x_2]$ 上可微，并且 $x_1 x_2 > 0$. 证明：
\[
\frac{1}{x_1 - x_2}
\begin{vmatrix}
x_1 & x_2 \\
f(x_1) & f(x_2)
\end{vmatrix}
= f(\xi) - \xi f'(\xi),
\]
其中 $x_1 < \xi < x_2$.
\end{prob}

\begin{prob}[M]{1254}
证明：若函数 $f(x)$ 在有限的区间 $(a, b)$ 内可微但无界，则其导数 $f'(x)$ 在区间 $(a, b)$ 内也无界. 逆命题不真（举出例子）.
\end{prob}

\begin{prob}[M]{1255}
证明：若函数 $f(x)$ 在有限或无穷的区间 $(a, b)$ 内有有界的导数 $f'(x)$，则 $f(x)$ 在 $(a, b)$ 中一致连续.
\end{prob}

\begin{prob}[M]{1256}
证明：若函数 $f(x)$ 在无穷的区间 $(x_0, +\infty)$ 内可微，且
\[
\lim_{x \to +\infty} f'(x) = 0,
\]
则
\[
\lim_{x \to +\infty} \frac{f(x)}{x} = 0,
\]
即当 $x \to +\infty$ 时，$f(x) = o(x)$.
\end{prob}

\begin{prob}[M]{1257}
证明：若函数 $f(x)$ 在无穷的区间 $(x_0, +\infty)$ 内可微且
\[
\text{当 } x \to +\infty \text{ 时，} f(x) = o(x),
\]
则
\[
\lim_{x \to +\infty} |f'(x)| = 0.
\]
\end{prob}

\begin{prob}[M]{1258.1}
证明：若函数 $f(x)$：(1) 在闭区间 $[x_0, X]$ 上有定义并且是连续的；(2) 在区间 $(x_0, X)$ 内有有限的导数 $f'(x)$；(3) 存在有限或无穷的极限
\[
\lim_{x \to x_0 + 0} f'(x) = f'(x_0 + 0),
\]
则相应地存在有限或无穷的单侧导数 $f'_+(x_0)$ 且
\[
f'_+(x_0) = f'(x_0 + 0).
\]
\end{prob}

\begin{prob}[M]{1258.2}
证明：对于函数
\[
f(x) = \begin{cases}
\arctan \dfrac{1 + x}{1 - x}, & x \neq 1, \\
0, & x = 1
\end{cases}
\]
存在有限的极限
\[
\lim_{x \to 1} f'(x),
\]
但是函数 $f(x)$ 没有单侧导数 $f'_-(1)$ 及 $f'_+(1)$. 给出这个事实的几何解释. 然而，在这个点存在广义单侧导数（见习题 1009.2）.
\end{prob}

\begin{prob}[S]{1259}
证明：若当 $a < x < b$ 时 $f'(x) = 0$，则
\[
\text{当 } a < x < b \text{ 时，} f(x) = \text{常数}.
\]
\end{prob}

\begin{prob}[S]{1260}
证明：导数为常数
\[
f'(x) = k
\]
的唯一函数 $f(x)$ $(-\infty < x < +\infty)$ 是线性函数
\[
f(x) = kx + b.
\]
\end{prob}

\begin{prob}[S]{1261.1}
若 $f^{(n)}(x) = 0$，则函数 $f(x)$ 有什么性质？
\end{prob}

\begin{prob}[M]{1261.2}
设 $f(x) \in C^{(\infty)}(-\infty, +\infty)$，并且对每个 $x$ 都存在正整数 $n_x$ $(n_x \leqslant n)$，使得
\[
f^{(n_x)}(x) = 0.
\]
证明：函数 $f(x)$ 是多项式.
\end{prob}

\begin{prob}[M]{1262}
证明：满足方程
\[
y' = \lambda y \quad (\lambda = \text{常数})
\]
的唯一函数 $y = y(x)$ ($-\infty < x < +\infty$) 是指数函数 $y = C\mathrm{e}^{\lambda x}$，其中 $C$ 为任意常数.
\probhint{研究 $(y\mathrm{e}^{-\lambda x})'$.}
\end{prob}

\begin{prob}[M]{1263}
检验函数
\[
f(x) = \arctan \frac{1 + x}{1 - x}, \quad g(x) = \arctan x
\]
在区间：(1) $x < 1$ 及 (2) $x > 1$ 内有相同的导数. 推出这些函数间的关系.
\end{prob}

\begin{prob}[M]{1264}
证明下列恒等式：
\begin{enumerate}[label=(\alph*), itemsep=2pt, topsep=3pt]
\item $2\arctan x + \arcsin \dfrac{2x}{1 + x^2} = \pi \operatorname{sgn} x$，当 $|x| \geqslant 1$；
\item $3\arccos x - \arccos(3x - 4x^3) = \pi$，当 $|x| \leqslant \dfrac{1}{2}$.
\end{enumerate}
\end{prob}

\begin{prob}[M]{1265}
证明：若函数 $f(x)$：(1) 在闭区间 $[a, b]$ 上是连续的；(2) 于此区间内有有限的导数 $f'(x)$；(3) 不是线性函数，则在区间 $(a, b)$ 内至少能找到一点 $c$，使得
\[
|f'(c)| > \left|\frac{f(b) - f(a)}{b - a}\right|.
\]
给出这个事实的几何解释.
\end{prob}

\begin{prob}[M]{1266}
证明：若函数 $f(x)$：(1) 在闭区间 $[a, b]$ 上有二阶导数 $f''(x)$；(2) $f'(a) = f'(b) = 0$，则在区间 $(a, b)$ 内至少存在一点 $c$，使得
\[
|f''(c)| \geqslant \frac{4}{(b - a)^2}|f(b) - f(a)|.
\]
\end{prob}

\begin{prob}[M]{1267}
汽车从某点开始行驶，于 $t$ 秒内走完了路程，所经过的距离为 $s$. 证明：汽车运动的加速度在某瞬时不小于 $\dfrac{4s}{t^2}$.
\end{prob}

\probsec{§7. 增函数与减函数. 不等式}

\probgroup{求下列函数的严格单调 (增或减) 区间：}

\begin{prob}[S]{1268}
$y = 2 + x - x^2$.
\end{prob}

\begin{prob}[S]{1269}
$y = 3x - x^3$.
\end{prob}

\begin{prob}[S]{1270}
$y = \dfrac{2x}{1 + x^2}$.
\end{prob}

\begin{prob}[S]{1271}
$y = \dfrac{\sqrt{x}}{x + 100} \quad (x \geqslant 0)$.
\end{prob}

\begin{prob}[S]{1272}
$y = x + \sin x$.
\end{prob}

\begin{prob}[S]{1273}
$y = x + |\sin 2x|$.
\end{prob}

\begin{prob}[S]{1274}
$y = \cos \dfrac{\pi}{x}$.
\end{prob}

\begin{prob}[S]{1275}
$y = \dfrac{x^2}{2^x}$.
\end{prob}

\begin{prob}[S]{1276}
$y = x^n \mathrm{e}^{-x} \quad (n > 0, x \geqslant 0)$.
\end{prob}

\begin{prob}[S]{1277}
$y = x^2 - \ln x^2$.
\end{prob}

\begin{prob}[M]{1278}
若
\[
f(x) = \begin{cases}
x\left(\sqrt{\dfrac{3}{2}} + \sin\ln x\right), & x > 0, \\
0, & x = 0.
\end{cases}
\]
\end{prob}

\begin{prob}[M]{1279}
证明：圆的内接正 $n$ 边形的周长 $p_n$ 当边的数目 $n$ 增加时增加，而此圆的外切正 $n$ 边形的周长 $P_n$ 此时则减小. 利用这点来证明，当 $n \to \infty$ 时，$p_n$ 及 $P_n$ 有相同的极限.
\end{prob}

\begin{prob}[M]{1280}
证明：函数
\[
\left(1 + \frac{1}{x}\right)^x
\]
在区间 $(-\infty, -1)$ 及 $(0, +\infty)$ 内递增.
\end{prob}

\begin{prob}[M]{1281}
证明：有理函数
\[
P(x) = a_0 + a_1 x + \cdots + a_n x^n \quad (n \geqslant 1, a_n \neq 0)
\]
是区间 $(-\infty, -x_0)$ 及 $(x_0, +\infty)$ 上的严格单调函数，其中 $x_0$ 为充分大的正数.
\end{prob}

\begin{prob}[M]{1282}
证明：有理函数
\[
R(x) = \frac{a_0 + a_1 x + \cdots + a_n x^n}{b_0 + b_1 x + \cdots + b_m x^m} \quad (a_n b_m \neq 0)
\]
是区间 $(-\infty, -x_0)$ 及 $(x_0, +\infty)$ 上的严格单调函数，其中 $x_0$ 为充分大的正数.
\end{prob}

\begin{prob}[M]{1283}
单调函数的导数是否也必为单调的？研究一个例子：
\[
f(x) = x + \sin x.
\]
\end{prob}

\begin{prob}[M]{1284}
证明：若 $\varphi(x)$ 为可微的单调增函数，且当 $x \geqslant x_0$ 时，$|f'(x)| \leqslant \varphi'(x)$，则当 $x \geqslant x_0$ 时，$|f(x) - f(x_0)| \leqslant \varphi(x) - \varphi(x_0)$.
给出这个事实的几何解释.
\end{prob}

\begin{prob}[M]{1285}
设函数 $f(x)$ 在区间 $a \leqslant x < +\infty$ 内连续，而且当 $x > a$ 时，$f'(x) > k > 0$，其中 $k$ 为常数. 证明：若 $f(a) < 0$，则在区间 $\left(a, a - \dfrac{f(a)}{k}\right)$ 内方程 $f(x) = 0$ 有且仅有一个实根.
\end{prob}

\begin{prob}[M]{1286}
若在某邻域 $|x - x_0| < \delta$ 内，函数增量 $\Delta f(x_0) = f(x) - f(x_0)$ 的符号与自变量增量 $\Delta x_0 = x - x_0$ 的符号相同，则称函数 $f(x)$ 为在 $x_0$ 点的增函数.

证明：若函数 $f(x)$ ($a < x < b$) 在有限或无穷的区间 $(a, b)$ 内的每一点皆为增函数，则它在此区间内为增函数.
\end{prob}

\begin{prob}[M]{1287}
证明：函数
\[
f(x) = \begin{cases}
x + x^2 \sin\dfrac{2}{x}, & x \neq 0, \\
0, & x = 0
\end{cases}
\]
在点 $x = 0$ 是增函数，但在包含此点的任何区间 $(-\varepsilon, \varepsilon)$ 中并非增函数，其中 $\varepsilon > 0$ 为任意小的数. 作出此函数的略图.
\end{prob}

\begin{prob}[M]{1288}
证明定理：设 (1) 函数 $\varphi(x)$ 及 $\psi(x)$ 为 $n$ 阶可微函数；(2) $\varphi^{(k)}(x_0) = \psi^{(k)}(x_0)$ ($k = 0, 1, \cdots, n - 1$)；(3) 当 $x > x_0$ 时，$\varphi^{(n)}(x) > \psi^{(n)}(x)$，则当 $x > x_0$ 时有不等式
\[
\varphi(x) > \psi(x).
\]
\end{prob}

\begin{prob}[L]{1289}
证明下列不等式：
\begin{enumerate}[label=(\alph*), itemsep=2pt, topsep=3pt]
\item $\mathrm{e}^x > 1 + x \quad (x \neq 0)$；
\item $x - \dfrac{x^2}{2} < \ln(1 + x) < x \quad (x > 0)$；
\item $x - \dfrac{x^3}{6} < \sin x < x \quad (x > 0)$；
\item $\tan x > x + \dfrac{x^3}{3} \quad \left(0 < x < \dfrac{\pi}{2}\right)$；
\item $(x^\alpha + y^\alpha)^{\frac{1}{\alpha}} > (x^\beta + y^\beta)^{\frac{1}{\beta}} \quad (x > 0, y > 0, 0 < \alpha < \beta)$.
\end{enumerate}
给出不等式 (a)---(d) 的几何解释.
\end{prob}

\begin{prob}[M]{1290}
证明：不等式
\[
\frac{2}{\pi} x < \sin x < x
\]
当 $0 < x < \dfrac{\pi}{2}$ 时成立.
\end{prob}

\begin{prob}[M]{1291}
证明：当 $x > 0$ 时有不等式
\[
\left(1 + \frac{1}{x}\right)^x < \mathrm{e} < \left(1 + \frac{1}{x}\right)^{x+1}.
\]
\end{prob}

\begin{prob}[M]{1292}
等差数列与等比数列的项的数目相同且有相同的首项与末项，并且一切项都是正的. 证明：等差数列各项之和大于等于等比数列各项之和.
\end{prob}

\begin{prob}[M]{1293}
用不等式
\[
\sum_{k=1}^{n} (a_k x + b_k)^2 \geqslant 0,
\]
其中 $x, a_k, b_k$ ($k = 1, \cdots, n$) 为实数，来证明柯西--布尼亚科夫斯基不等式：
\[
\left(\sum_{k=1}^{n} a_k b_k\right)^2 \leqslant \sum_{k=1}^{n} a_k^2 \cdot \sum_{k=1}^{n} b_k^2.
\]
\end{prob}

\begin{prob}[M]{1294}
证明：正数的算术平均值的平方不大于这些数的平方平均值，即
\[
\frac{1}{n}\sum_{k=1}^{n} x_k \leqslant \sqrt{\frac{1}{n}\sum_{k=1}^{n} x_k^2}.
\]
\end{prob}

\begin{prob}[M]{1295}
证明：正数的几何平均值不大于这些数的算术平均值，即
\[
\sqrt[n]{x_1 x_2 \cdots x_n} \leqslant \frac{1}{n}(x_1 + x_2 + \cdots + x_n).
\]
\probhint{用数学归纳法.}
\end{prob}

\begin{prob}[L]{1296}
设 $a$ 及 $b$ 为两个正数，则由等式
\[
\Delta_s(a, b) = \begin{cases}
\left(\dfrac{a^s + b^s}{2}\right)^{\frac{1}{s}}, & s \neq 0, \\
\lim\limits_{t \to 0} \Delta_t(a, b), & s = 0
\end{cases}
\]
所定义的函数称为正数 $a$ 与 $b$ 之 $s$ 阶平均值.

例如，当 $s = -1$ 时得调和平均值，当 $s = 0$ 时得几何平均值 (试证明！)；当 $s = 1$ 时得算术平均值；当 $s = 2$ 时得平方平均值.

证明：(1) $\min(a, b) \leqslant \Delta_s(a, b) \leqslant \max(a, b)$；
(2) 当 $a \neq b$ 时，函数 $\Delta_s(a, b)$ 是变量 $s$ 的增函数；
(3) $\lim\limits_{s \to -\infty} \Delta_s(a, b) = \min(a, b)$，$\lim\limits_{s \to +\infty} \Delta_s(a, b) = \max(a, b)$.
\probhint{考察 $\dfrac{\mathrm{d}}{\mathrm{d}s}[\ln \Delta_s(a, b)]$.}
\end{prob}

\begin{prob}[L]{1297}
证明下列不等式：
\begin{enumerate}[label=(\alph*), itemsep=2pt, topsep=3pt]
\item $x^\alpha - 1 > \alpha(x - 1)$，若 $\alpha \geqslant 2, x > 1$；
\item $\sqrt[n]{x} - \sqrt[n]{a} < \sqrt[n]{x - a}$，若 $n > 1, x > a > 0$；
\item $1 + 2\ln x \leqslant x^2$，若 $x > 0$.
\end{enumerate}
\end{prob}

\probsec{§8. 凹凸性. 拐点}

\begin{prob}[M]{1298}
研究曲线
\[
y = 1 + \sqrt[3]{x}
\]
于 $A(-1, 0)$，$B(1, 2)$ 及 $C(0, 1)$ 诸点的凹凸性.
\end{prob}

\probgroup{求下列函数的图像的凹或凸的区间及拐点：}

\begin{prob}[S]{1299}
$y = 3x^2 - x^3$.
\end{prob}

\begin{prob}[S]{1300}
$y = \dfrac{a^2}{a^2 + x^2} \quad (a > 0)$.
\end{prob}

\begin{prob}[S]{1301}
$y = x + x^{\frac{5}{3}}$.
\end{prob}

\begin{prob}[S]{1302}
$y = \sqrt{1 + x^2}$.
\end{prob}

\begin{prob}[S]{1303}
$y = x + \sin x$.
\end{prob}

\begin{prob}[S]{1304}
$y = \mathrm{e}^{-x^2}$.
\end{prob}

\begin{prob}[S]{1305}
$y = \ln(1 + x^2)$.
\end{prob}

\begin{prob}[S]{1306}
$y = x \sin(\ln x) \quad (x > 0)$.
\end{prob}

\begin{prob}[S]{1307}
$y = x^x \quad (x > 0)$.
\end{prob}

\begin{prob}[M]{1308}
证明：曲线
\[
y = \frac{x + 1}{x^2 + 1}
\]
在位于同一直线上的三个拐点. 作出这个函数的图像.
\end{prob}

\begin{prob}[M]{1309}
如何选择参量 $h$，可使“概率曲线”
\[
y = \frac{h}{\sqrt{\pi}} \mathrm{e}^{-h^2 x^2} \quad (h > 0)
\]
有拐点 $x = \pm \sigma$？
\end{prob}

\begin{prob}[M]{1310}
研究摆线 (旋轮线)
\[
x = a(t - \sin t), \quad y = a(1 - \cos t) \quad (a > 0)
\]
的凹凸性.
\end{prob}

\begin{prob}[M]{1311}
设函数 $f(x)$ 在区间 $a \leqslant x < +\infty$ 中二阶可微，并且：
(1) $f(a) = A > 0$； (2) $f'(a) < 0$； (3) 当 $x > a$ 时，$f''(x) \leqslant 0$.
证明：在区间 $(a, +\infty)$ 内方程 $f(x) = 0$ 有且仅有一个实根.
\end{prob}

\begin{prob}[M]{1312}
若对于区间 $(a, b)$ 内的任意两点 $x_1$ 与 $x_2$ 及任意两数 $\lambda_1$ 与 $\lambda_2$ ($\lambda_1 > 0, \lambda_2 > 0, \lambda_1 + \lambda_2 = 1$) 有不等式：
\[
f(\lambda_1 x_1 + \lambda_2 x_2) < \lambda_1 f(x_1) + \lambda_2 f(x_2)
\]
(或对应地，相反的不等式 $f(\lambda_1 x_1 + \lambda_2 x_2) > \lambda_1 f(x_1) + \lambda_2 f(x_2)$)，则称函数 $f(x)$ 在区间 $(a, b)$ 上是凸 (凹) 的.

证明：(1) 若当 $a < x < b$ 时，$f''(x) > 0$，则函数 $f(x)$ 在 $(a, b)$ 上是凸的；(2) 若当 $a < x < b$ 时，$f''(x) < 0$，则 $f(x)$ 在 $(a, b)$ 上是凹的.
\end{prob}

\begin{prob}[M]{1313}
证明：函数
\[
x^n \ (n > 1), \quad \mathrm{e}^x, \quad x \ln x
\]
在区间 $(0, +\infty)$ 上是凸的，而函数
\[
x^n \ (0 < n < 1), \quad \ln x
\]
在区间 $(0, +\infty)$ 上是凹的.
\end{prob}

\begin{prob}[L]{1314.1}
证明下列不等式，并解释其几何意义：
\begin{enumerate}[label=(\alph*), itemsep=2pt, topsep=3pt]
\item $\dfrac{1}{2}(x^n + y^n) > \left(\dfrac{x + y}{2}\right)^n \ (x > 0, y > 0, x \neq y, n > 1)$；
\item $\dfrac{\mathrm{e}^x + \mathrm{e}^y}{2} > \mathrm{e}^{\frac{x+y}{2}} \ (x \neq y)$；
\item $x \ln x + y \ln y > (x + y)\ln\dfrac{x + y}{2} \ (x > 0, y > 0)$.
\end{enumerate}
\end{prob}

\begin{prob}[M]{1314.2}
设当 $a \leqslant x \leqslant b$ 时 $f''(x) \geqslant 0$，证明：对于任何 $x_1, x_2 \in [a, b]$，都有
\[
f\left(\frac{x_1 + x_2}{2}\right) \leqslant \frac{1}{2}[f(x_1) + f(x_2)].
\]
\end{prob}

\begin{prob}[M]{1315}
证明：有界的凸函数处处连续，并有左导数和右导数.
\end{prob}

\begin{prob}[M]{1316}
设函数 $f(x)$ 在区间 $(a, b)$ 内二阶可微，且 $f''(\xi) \neq 0$，其中 $a < \xi < b$. 证明：在区间 $(a, b)$ 中可找出两个值 $x_1$ 与 $x_2$，满足
\[
\frac{f(x_2) - f(x_1)}{x_2 - x_1} = f'(\xi).
\]
\end{prob}

\begin{prob}[M]{1317}
证明：若函数 $f(x)$ 在无穷的区间 $(x_0, +\infty)$ 内二阶可微，且
\[
\lim_{x \to x_0 + 0} f(x) = 0, \quad \lim_{x \to +\infty} f(x) = 0,
\]
则在区间 $(x_0, +\infty)$ 内最少有一点 $\xi$，满足 $f''(\xi) = 0$.
\end{prob}

\probsec{§9. 不定式的求值法}

\probgroup{求出下列各式之值：}

\begin{prob}[S]{1318}
$\lim_{x \to 0} \frac{\sin ax}{\sin bx}$.
\end{prob}

\begin{prob}[S]{1319}
$\lim_{x \to 0} \frac{\cosh x - \cos x}{x^2}$.
\end{prob}

\begin{prob}[S]{1320}
$\lim_{x \to 0} \frac{\tan x - x}{x - \sin x}$.
\end{prob}

\begin{prob}[S]{1321}
$\lim_{x \to 0} \frac{3\tan 4x - 12\tan x}{3\sin 4x - 12\sin x}$.
\end{prob}

\begin{prob}[S]{1322}
$\lim_{x \to \frac{\pi}{2}} \frac{\tan 3x}{\tan x}$.
\end{prob}

\begin{prob}[S]{1323}
$\lim_{x \to 0} \frac{x\cot x - 1}{x^2}$.
\end{prob}

\begin{prob}[S]{1324}
$\lim_{x \to \frac{\pi}{4}} \frac{\sqrt[3]{\tan x} - 1}{2\sin^2 x - 1}$.
\end{prob}

\begin{prob}[S]{1325}
$\lim_{x \to 0} \frac{x(\mathrm{e}^x + 1) - 2(\mathrm{e}^x - 1)}{x^3}$.
\end{prob}

\begin{prob}[S]{1326}
$\lim_{x \to 0} \frac{1 - \cos x^2}{x^2 \sin x^2}$.
\end{prob}

\begin{prob}[S]{1327}
$\lim_{x \to 0} \frac{\arcsin 2x - 2\arcsin x}{x^3}$.
\end{prob}

\begin{prob}[S]{1328}
$\lim_{x \to 0} \frac{1}{x\sqrt{x}} \left( \sqrt{a}\arctan\sqrt{\frac{x}{a}} - \sqrt{b}\arctan\sqrt{\frac{x}{b}} \right)$.
\end{prob}

\begin{prob}[S]{1329}
$\lim_{x \to 0} \frac{a^x - a^{\sin x}}{x^3} \quad (a > 0)$.
\end{prob}

\begin{prob}[S]{1330}
$\lim_{x \to 1} \left( \frac{x^x - x}{\ln x - x + 1} \right)$.
\end{prob}

\begin{prob}[S]{1331}
$\lim_{x \to 0} \frac{\ln(\sin ax)}{\ln(\sin bx)}$.
\end{prob}

\begin{prob}[S]{1332}
$\lim_{x \to 0} \frac{\ln(\cos ax)}{\ln(\cos bx)}$.
\end{prob}

\begin{prob}[S]{1333}
$\lim_{x \to 0} \frac{\cos(\sin x) - \cos x}{x^4}$.
\end{prob}

\begin{prob}[S]{1334}
$\lim_{x \to 0} \frac{1}{x} \left( \frac{1}{\tanh x} - \frac{1}{\tan x} \right)$.
\end{prob}

\begin{prob}[S]{1335}
$\lim_{x \to 0} \frac{\operatorname{arsinh}(\sinh x) - \operatorname{arsinh}(\sin x)}{\sinh x - \sin x}$，其中 $\operatorname{arsinh} x = \ln(x + \sqrt{1 + x^2})$。
\end{prob}

\begin{prob}[S]{1336}
$\lim_{x \to +\infty} \frac{\ln x}{x^\varepsilon} \quad (\varepsilon > 0)$.
\end{prob}

\begin{prob}[S]{1337}
$\lim_{x \to +\infty} \frac{x^n}{\mathrm{e}^{ax}} \quad (a > 0, n > 0)$.
\end{prob}

\begin{prob}[S]{1338}
$\lim_{x \to 0} \frac{\mathrm{e}^{-\frac{1}{x^2}}}{x^{1000}}$.
\end{prob}

\begin{prob}[S]{1339}
$\lim_{x \to +\infty} x^2 \mathrm{e}^{-0.01x}$.
\end{prob}

\begin{prob}[S]{1340}
$\lim_{x \to 1-0} \ln x \cdot \ln(1 - x)$.
\end{prob}

\begin{prob}[S]{1341}
$\lim_{x \to +0} x^\varepsilon \ln x \quad (\varepsilon > 0)$.
\end{prob}

\begin{prob}[S]{1342}
$\lim_{x \to +0} x^x$.
\end{prob}

\begin{prob}[S]{1343}
$\lim_{x \to 0} x^{x^x - 1}$.
\end{prob}

\begin{prob}[S]{1344}
$\lim_{x \to 0} \left( x^{x^x} - 1 \right)$.
\end{prob}

\begin{prob}[S]{1345}
$\lim_{x \to +0} x^{\frac{k}{1 + \ln x}}$.
\end{prob}

\begin{prob}[S]{1346}
$\lim_{x \to 1} x^{\frac{1}{1 - x}}$.
\end{prob}

\begin{prob}[S]{1347}
$\lim_{x \to 1} (2 - x)^{\tan\frac{\pi x}{2}}$.
\end{prob}

\begin{prob}[S]{1348}
$\lim_{x \to \frac{\pi}{4}} (\tan x)^{\tan 2x}$.
\end{prob}

\begin{prob}[S]{1349}
$\lim_{x \to 0} (\cot x)^{\sin x}$.
\end{prob}

\begin{prob}[S]{1350}
$\lim_{x \to +0} \left( \ln\frac{1}{x} \right)^x$.
\end{prob}

\begin{prob}[S]{1351}
$\lim_{x \to \infty} \left( \tan\frac{\pi x}{2x + 1} \right)^{\frac{1}{x}}$.
\end{prob}

\begin{prob}[S]{1352}
$\lim_{x \to a} \left( \frac{\tan x}{\tan a} \right)^{\cot(x - a)}$.
\end{prob}

\begin{prob}[S]{1353}
$\lim_{x \to 0} \left( \frac{a^x - x\ln a}{b^x - x\ln b} \right)^{\frac{1}{x^2}}$.
\end{prob}

\begin{prob}[S]{1354}
$\lim_{x \to 0} \left( \frac{1}{x} - \frac{1}{\mathrm{e}^x - 1} \right)$.
\end{prob}

\begin{prob}[S]{1355}
$\lim_{x \to 1} \left( \frac{1}{\ln x} - \frac{1}{x - 1} \right)$.
\end{prob}

\begin{prob}[S]{1356}
$\lim_{x \to 0} \left( \cot x - \frac{1}{x} \right)$.
\end{prob}

\begin{prob}[S]{1357}
$\lim_{x \to 0} \left[ \frac{1}{\ln(x + \sqrt{1 + x^2})} - \frac{1}{\ln(1 + x)} \right]$.
\end{prob}

\begin{prob}[S]{1358}
$\lim_{x \to a} \frac{a^x - x^a}{x - a} \quad (a > 0)$.
\end{prob}

\begin{prob}[S]{1359}
$\lim_{x \to 0} \frac{(1 + x)^{\frac{1}{x}} - \mathrm{e}}{x}$.
\end{prob}

\begin{prob}[S]{1360}
$\lim_{x \to 0} \frac{(a + x)^x - a^x}{x^2} \quad (a > 0)$.
\end{prob}

\begin{prob}[S]{1361}
$\lim_{x \to +\infty} \left( \frac{2}{\pi}\arctan x \right)^x$.
\end{prob}

\begin{prob}[S]{1362}
$\lim_{x \to +\infty} (\tanh x)^x$.
\end{prob}

\begin{prob}[L]{1363}
\begin{enumerate}[label=(\alph*), itemsep=2pt, topsep=3pt]
\item $\lim_{x \to 0} \left( \frac{\arcsin x}{x} \right)^{\frac{1}{x^2}}$;
\item $\lim_{x \to 0} \left( \frac{\sin x}{x} \right)^{\frac{1}{x^2}}$;
\item $\lim_{x \to 0} \left( \frac{\tan x}{x} \right)^{\frac{1}{x^2}}$;
\item $\lim_{x \to 0} \left( \frac{\arctan x}{x} \right)^{\frac{1}{x^2}}$;
\item $\lim_{x \to 0} \left( \frac{\operatorname{arsinh} x}{x} \right)^{\frac{1}{x^2}}$，其中 $\operatorname{arsinh} x = \ln(x + \sqrt{1 + x^2})$。
\end{enumerate}
\end{prob}

\begin{prob}[S]{1364}
$\lim_{x \to 0} \left[ \frac{(1 + x)^{\frac{1}{x}}}{\mathrm{e}} \right]^{\frac{1}{x}}$.
\end{prob}

\begin{prob}[S]{1365}
$\lim_{x \to 0} \left( \frac{2}{\pi}\arccos x \right)^{\frac{1}{x}}$.
\end{prob}

\begin{prob}[S]{1366}
$\lim_{x \to 0} \left( \frac{\cos x}{\cosh x} \right)^{\frac{1}{x^2}}$.
\end{prob}

\begin{prob}[S]{1367}
$\lim_{x \to 0} \frac{\ln\cosh x}{\sqrt[m]{\cosh x} - \sqrt[n]{\cosh x}}$.
\end{prob}

\begin{prob}[M]{1368}
(a) $\lim_{x \to 0} \left( \frac{1 + \mathrm{e}^x}{2} \right)^{\coth x}$; \qquad (b) $\lim_{x \to +\infty} \frac{x^{\ln x}}{(\ln x)^x}$.
\end{prob}

\begin{prob}[S]{1369}
$\lim_{x \to +\infty} \left[ \sqrt[3]{x^3 + x^2 + x + 1} - \sqrt{x^2 + x + 1} \cdot \frac{\ln(\mathrm{e}^x + x)}{x} \right]$.
\end{prob}

\begin{prob}[S]{1370}
$\lim_{x \to +\infty} \left[ (x + a)^{1 + \frac{1}{x}} - x^{1 + \frac{1}{x + a}} \right]$.
\end{prob}

\begin{prob}[M]{1371}
若当 $x \to 0$ 时，曲线 $y = f(x)$ 通过坐标原点 $(0, 0)$ [$\lim_{x \to 0} f(x) = 0$]，且在此有倾角 $\alpha$，求 $\lim_{x \to 0} \frac{y}{x}$。
\end{prob}

\begin{prob}[M]{1372}
若当 $x \to +0$ 时，连续曲线 $y = f(x)$ 通过坐标原点 $(0, 0)$ [$\lim_{x \to +0} f(x) = 0$]，并且当 $0 < x < \varepsilon$ 时，此曲线完全位于两直线 $y = -kx$ 及 $y = kx$ ($k \neq \infty$) 所组成的锐角之内，证明：$\lim_{x \to +0} x^{f(x)} = 1$。
\end{prob}

\begin{prob}[M]{1373.1}
证明：若函数 $f(x)$ 的二阶导数 $f''(x)$ 存在，则
\[
f''(x) = \lim_{h \to 0} \frac{f(x + h) + f(x - h) - 2f(x)}{h^2}.
\]
\end{prob}

\begin{prob}[M]{1373.2}
研究函数
\[
f(x) = \begin{cases} \dfrac{1}{x} - \dfrac{1}{\mathrm{e}^x - 1}, & x \neq 0, \\ \dfrac{1}{2}, & x = 0 \end{cases}
\]
在点 $x = 0$ 的可微性。
\end{prob}

\begin{prob}[M]{1373.3}
求曲线
\[
y = \frac{x^{1 + x}}{(1 + x)^x} \quad (x > 0)
\]
的渐近线。
\end{prob}

\begin{prob}[L]{1374}
研究运用洛必达法则于下列各例的可能性：
\begin{enumerate}[label=(\alph*), itemsep=2pt, topsep=3pt]
\item $\lim_{x \to 0} \frac{x^2 \sin\frac{1}{x}}{\sin x}$;
\item $\lim_{x \to \infty} \frac{x - \sin x}{x + \sin x}$;
\item $\lim_{x \to +\infty} \frac{\mathrm{e}^{-2x}(\cos x + 2\sin x) + \mathrm{e}^{-x^2}\sin^2 x}{\mathrm{e}^{-x}(\cos x + \sin x)}$;
\item $\lim_{x \to \infty} \frac{1 + x + \sin x\cos x}{(x + \sin x\cos x)\mathrm{e}^{\sin x}}$.
\end{enumerate}
\end{prob}

\begin{prob}[M]{1375}
设有一弓形，其弦长为 $b$，拱高为 $h$，半径为 $R$，又有内接于此弓形的等腰三角形。若当 $R$ 不变时弓形的弧长趋于零，求弓形面积与内接等腰三角形面积之比的极限。利用所得结果推出弓形面积的近似公式：
\[
S \approx \frac{2}{3} bh.
\]
\end{prob}

\probsec{§10. 泰勒公式}

\begin{prob}[M]{1376}
将多项式
\[
P(x) = 1 + 3x + 5x^2 - 2x^3
\]
改写为二项式 $x + 1$ 的非负整数次幂多项式。
\end{prob}

\probgroup{写出下列函数按变量 $x$ 的非负整数次幂的展开式，到含有所指阶数的项为止：}

\begin{prob}[S]{1377}
$\frac{1 + x + x^2}{1 - x + x^2}$ 到含 $x^4$ 的项。$f^{(4)}(0)$ 等于什么？
\end{prob}

\begin{prob}[S]{1378}
$\frac{(1 + x)^{100}}{(1 - 2x)^{40}(1 + 2x)^{60}}$ 到含 $x^2$ 的项。
\end{prob}

\begin{prob}[S]{1379}
$\sqrt[m]{a^m + x} \quad (a > 0)$ 到含 $x^2$ 的项。
\end{prob}

\begin{prob}[S]{1380}
$\sqrt{1 - 2x + x^3} - \sqrt[3]{1 - 3x + x^2}$ 到含 $x^3$ 的项。
\end{prob}

\begin{prob}[S]{1381}
$\mathrm{e}^{2x - x^2}$ 到含 $x^5$ 的项。
\end{prob}

\begin{prob}[S]{1382}
$\frac{x}{\mathrm{e}^x - 1}$ 到含 $x^4$ 的项。
\end{prob}

\begin{prob}[S]{1383}
$\sqrt[3]{\sin x^3}$ 到含 $x^{13}$ 的项。
\end{prob}

\begin{prob}[S]{1384}
$\ln\cos x$ 到含 $x^6$ 的项。
\end{prob}

\begin{prob}[S]{1385}
$\sin(\sin x)$ 到含 $x^3$ 的项。
\end{prob}

\begin{prob}[S]{1386}
$\tan x$ 到含 $x^5$ 的项。
\end{prob}

\begin{prob}[S]{1387}
$\ln\frac{\sin x}{x}$ 到含 $x^6$ 的项。
\end{prob}

\begin{prob}[M]{1388}
求函数 $f(x) = \sqrt{x}$ 按照 $x - 1$ 的非负整数次幂展开式的前三项。
\end{prob}

\begin{prob}[M]{1389}
将函数 $f(x) = x^x - 1$ 按照 $x - 1$ 的非负整数次幂展开至含有 $(x - 1)^3$ 的项。
\end{prob}

\begin{prob}[M]{1390}
在点 $x = 0$ 的邻域中，用抛物线（二次多项式）近似地代替函数
\[
y = a\cosh\frac{x}{a} \quad (a > 0).
\]
\end{prob}

\begin{prob}[M]{1391}
按公式 $\frac{1}{x}$ 的非负整数次幂展开函数 $f(x) = \sqrt{1 + x^2} - x$ $(x > 0)$ 到含 $\frac{1}{x^3}$ 的项。
\end{prob}

\begin{prob}[M]{1392}
求函数 $f(h) = \ln(x + h)$ $(x > 0)$ 按增量 $h$ 的正整数次幂的展开式，到含 $h^n$ 的项（$n$ 为正整数）。
\end{prob}

\begin{prob}[M]{1393.1}
设
\[
f(x + h) = f(x) + hf'(x) + \cdots + \frac{h^n}{n!} f^{(n)}(x + \theta h) \quad (0 < \theta < 1),
\]
且 $f^{(n+1)}(x) \neq 0$。证明：
\[
\lim_{h \to 0} \theta = \frac{1}{n + 1}.
\]
\end{prob}

\begin{prob}[M]{1393.2}
设当 $x \to 0$ 时有 $f(x) = 1 + kx + o(x)$。证明：
\[
\lim_{x \to 0} [f(x)]^{\frac{1}{x}} = \mathrm{e}^k.
\]
\end{prob}

\begin{prob}[M]{1393.3}
设 $f(x) \in C^{(2)}[0, 1]$，$f(0) = f(1) = 0$，且当 $x \in (0, 1)$ 时 $|f''(x)| \leqslant A$。证明：当 $0 \leqslant x \leqslant 1$ 时
\[
|f'(x)| \leqslant \frac{A}{2}.
\]
\end{prob}

\begin{prob}[M]{1393.4}
设 $f(x)$ $(-\infty < x < +\infty)$ 为二阶可微函数，且
\[
M_k = \sup_{-\infty < x < +\infty} |f^{(k)}(x)| < +\infty \quad (k = 0, 1, 2).
\]
证明不等式:
\[
M_1^2 \leqslant 2M_0 M_2.
\]
\end{prob}

\begin{prob}[L]{1394}
估计下列近似公式的绝对误差：
\begin{enumerate}[label=(\alph*), itemsep=2pt, topsep=3pt]
    \item $\mathrm{e}^x \approx 1 + x + \dfrac{x^2}{2!} + \cdots + \dfrac{x^n}{n!}$，当 $0 \leqslant x \leqslant 1$；
    \item $\sin x \approx x - \dfrac{x^3}{6}$，当 $|x| \leqslant \dfrac{1}{2}$；
    \item $\tan x \approx x + \dfrac{x^3}{3}$，当 $|x| \leqslant 0.1$；
    \item $\sqrt{1 + x} \approx 1 + \dfrac{x}{2} - \dfrac{x^2}{8}$，当 $0 \leqslant x \leqslant 1$.
\end{enumerate}
\end{prob}

\begin{prob}[M]{1395.1}
近似公式
\[
\cos x = 1 - \frac{x^2}{2}
\]
对于怎样的 $x$ 精确到 $0.0001$?
\end{prob}

\begin{prob}[M]{1395.2}
证明公式：
\[
\sqrt[n]{a^n + x} = a + \frac{x}{na^{n-1}} - r \quad (n \geqslant 2, \, a > 0, \, x > 0),
\]
其中 $0 < r < \dfrac{n - 1}{2n^2} \cdot \dfrac{x^2}{a^{2n - 1}}$.
\end{prob}

\begin{prob}[L]{1396}
利用泰勒公式近似地计算并估计误差：
\begin{enumerate}[label=(\alph*), itemsep=2pt, topsep=3pt]
    \item $\sqrt[3]{30}$；
    \item $\sqrt[5]{250}$；
    \item $\sqrt[12]{4000}$；
    \item $\sqrt{\mathrm{e}}$；
    \item $\sin 18^\circ$；
    \item $\ln 1.2$；
    \item $\arctan 0.8$；
    \item $\arcsin 0.45$；
    \item $1.1^{1.2}$.
\end{enumerate}
\end{prob}

\begin{prob}[L]{1397}
计算：
\begin{enumerate}[label=(\alph*), itemsep=2pt, topsep=3pt]
    \item $\mathrm{e}$，精确到 $10^{-9}$；
    \item $\sin 1^\circ$，精确到 $10^{-8}$；
    \item $\cos 9^\circ$，精确到 $10^{-5}$；
    \item $\sqrt{5}$，精确到 $10^{-4}$；
    \item $\lg 11$，精确到 $10^{-5}$.
\end{enumerate}
\end{prob}

\probgroup{利用展开式 $\mathrm{I}\text{---}\mathrm{V}$，求下列极限：}

\begin{prob}[S]{1398}
$\lim\limits_{x \to 0} \dfrac{\cos x - \mathrm{e}^{-\frac{x^2}{2}}}{x^4}$.
\end{prob}

\begin{prob}[S]{1399}
$\lim\limits_{x \to 0} \dfrac{\mathrm{e}^x \sin x - x(1 + x)}{x^3}$.
\end{prob}

\begin{prob}[S]{1400}
$\lim\limits_{x \to +\infty} x^{\frac{3}{2}} (\sqrt{x + 1} + \sqrt{x - 1} - 2\sqrt{x})$.
\end{prob}

\begin{prob}[S]{1401}
$\lim\limits_{x \to +\infty} (\sqrt[6]{x^6 + x^5} - \sqrt[6]{x^6 - x^5})$.
\end{prob}

\begin{prob}[S]{1402}
$\lim\limits_{x \to +\infty} \left[\left(x^3 - x^2 + \dfrac{x}{2}\right)\mathrm{e}^{\frac{1}{x}} - \sqrt{x^6 + 1}\right]$.
\end{prob}

\begin{prob}[S]{1403}
$\lim\limits_{x \to 0} \dfrac{a^x + a^{-x} - 2}{x^2} \ (a > 0)$.
\end{prob}

\begin{prob}[S]{1404}
$\lim\limits_{x \to \infty} \left[x - x^2 \ln\left(1 + \dfrac{1}{x}\right)\right]$.
\end{prob}

\begin{prob}[S]{1405}
$\lim\limits_{x \to 0} \left(\dfrac{1}{x} - \dfrac{1}{\sin x}\right)$.
\end{prob}

\begin{prob}[L]{1406}
\begin{enumerate}[label=(\alph*), itemsep=2pt, topsep=3pt]
    \item $\lim\limits_{x \to 0} \dfrac{1}{x} \left(\dfrac{1}{x} - \cot x\right)$；
    \item $\lim\limits_{x \to 0} \dfrac{\sin(\sin x) - x\sqrt[3]{1 - x^2}}{x^5}$；
    \item $\lim\limits_{x \to 0} \dfrac{1 - (\cos x)^{\sin x}}{x^3}$；
    \item $\lim\limits_{x \to 0} \dfrac{\sinh(\tan x) - x}{x^3}$.
\end{enumerate}
\end{prob}

\probgroup{当 $x \to 0$ 时，求出无穷小量 $y$ 的形如 $Cx^n$（$C$ 为常数）的主项，设：}

\begin{prob}[S]{1407}
$y = \tan(\sin x) - \sin(\tan x)$.
\end{prob}

\begin{prob}[S]{1408}
$y = (1 + x)^x - 1$.
\end{prob}

\begin{prob}[S]{1409}
$y = 1 - \dfrac{(1 + x)^{\frac{1}{x}}}{\mathrm{e}}$.
\end{prob}

\begin{prob}[M]{1410.1}
当选择怎样的系数 $a$ 与 $b$ 时，量
\[
x - (a + b\cos x)\sin x
\]
对于 $x$ 为 $5$ 阶无穷小?
\end{prob}

\begin{prob}[M]{1410.2}
选择系数 $A$ 和 $B$，使当 $x \to 0$ 时成立渐近公式
\[
\cot x = \frac{1 + Ax^2}{x + Bx^3} + O(x^5).
\]
\end{prob}

\begin{prob}[M]{1410.3}
系数 $A, B, C, D$ 为何值时，在 $x \to 0$ 时有渐近公式
\[
\mathrm{e}^x = \frac{1 + Ax + Bx^2}{1 + Cx + Dx^2} + O(x^5)\text{?}
\]
\end{prob}

\begin{prob}[L]{1411}
假设 $|x|$ 为小量，推出下列各式的简单的近似公式：
\begin{enumerate}[label=(\alph*), itemsep=2pt, topsep=3pt]
    \item $\dfrac{1}{R^2} - \dfrac{1}{(R + x)^2} \ (R > 0)$；
    \item $\sqrt[3]{\dfrac{1 + x}{1 - x}} - \sqrt[3]{\dfrac{1 - x}{1 + x}}$；
    \item $\dfrac{A}{x}\left[1 - \left(1 + \dfrac{x}{100}\right)^{-n}\right]$；
    \item $\dfrac{\ln 2}{\ln\left(1 + \dfrac{x}{100}\right)}$.
\end{enumerate}
\end{prob}

\begin{prob}[M]{1412}
假设 $x$ 的绝对值为小量，推出形如
\[
x = \alpha \sin x + \beta \tan x
\]
且精确到 $x^5$ 项的近似公式. 应用此公式近似地求小角度弧长.
\end{prob}

\begin{prob}[M]{1413}
估计下面的切比雪夫法则的相对误差：圆弧长近似地等于以此弧的弦为底、以其拱高的 $\sqrt{\dfrac{4}{3}}$ 为高的等腰三角形的两腰之和.
\end{prob}

\probsec{§11. 函数的极值. 函数的最大值和最小值}

\probgroup{研究下列函数的极值：}

\begin{prob}[S]{1414}
$y = 2 + x - x^2$.
\end{prob}

\begin{prob}[S]{1415}
$y = (x - 1)^3$.
\end{prob}

\begin{prob}[S]{1416}
$y = (x - 1)^4$.
\end{prob}

\begin{prob}[S]{1417}
$y = x^m (1 - x)^n$（$m$ 及 $n$ 为正整数）.
\end{prob}

\begin{prob}[S]{1418}
$y = \cos x + \cosh x$.
\end{prob}

\begin{prob}[S]{1419}
$y = (x + 1)^{10} \mathrm{e}^{-x}$.
\end{prob}

\begin{prob}[S]{1420}
$y = \left(1 + x + \dfrac{x^2}{2!} + \cdots + \dfrac{x^n}{n!}\right)\mathrm{e}^{-x}$（$n$ 为正整数）
\end{prob}

\begin{prob}[S]{1421}
$y = |x|$.
\end{prob}

\begin{prob}[S]{1422}
$y = x^{\frac{1}{3}} (1 - x)^{\frac{2}{3}}$.
\end{prob}

\begin{prob}[M]{1423}
设
\[
f(x) = (x - x_0)^n \varphi(x) \quad (n \text{ 为正整数}),
\]
其中函数 $\varphi(x)$ 当 $x = x_0$ 时连续，且 $\varphi(x_0) \neq 0$. 研究此函数在点 $x = x_0$ 的极值.
\end{prob}

\begin{prob}[M]{1424}
设
\[
f(x) = \frac{P(x)}{Q(x)}, \quad f'(x) = \frac{P_1(x)}{Q^2(x)},
\]
$x_0$ 为函数 $f(x)$ 的临界点，即
\[
P_1(x_0) = 0, \quad Q(x_0) \neq 0.
\]
证明：
\[
\operatorname{sgn} f''(x_0) = \operatorname{sgn} P_1'(x_0).
\]
\end{prob}

\begin{prob}[M]{1425}
可否断定：若函数 $f(x)$ 在点 $x_0$ 有极大值，则在此点的某充分小邻域内，函数 $f(x)$ 在点 $x_0$ 的左侧递增，而在其右侧递减? 研究一个例子：
\[
f(x) = \begin{cases}
2 - x^2 \left(2 + \sin\dfrac{1}{x}\right), & x \neq 0, \\
2, & x = 0.
\end{cases}
\]
\end{prob}

\begin{prob}[L]{1426}
证明：函数
\[
f(x) = \begin{cases}
\mathrm{e}^{-\frac{1}{x^2}}, & x \neq 0, \\
0, & x = 0
\end{cases}
\]
在点 $x = 0$ 有极小值，而函数
\[
g(x) = \begin{cases}
x\mathrm{e}^{-\frac{1}{x^2}}, & x \neq 0, \\
0, & x = 0
\end{cases}
\]
在点 $x = 0$ 没有极值，尽管
\[
f^{(n)}(0) = 0, \quad g^{(n)}(0) = 0 \quad (n = 1, 2, \cdots).
\]
作出这些函数的图像.
\end{prob}

\begin{prob}[L]{1427}
研究下列函数的极值并作出其图像：
\begin{enumerate}[label=(\alph*), itemsep=2pt, topsep=3pt]
    \item $f(x) = \begin{cases} \mathrm{e}^{-\frac{1}{|x|}}\left(\sqrt{2} + \sin\dfrac{1}{x}\right), & x \neq 0, \\ 0, & x = 0\text{；} \end{cases}$
    \item $f(x) = \begin{cases} \mathrm{e}^{-\frac{1}{|x|}}\left(\sqrt{2} + \cos\dfrac{1}{x}\right), & x \neq 0, \\ 0, & x = 0. \end{cases}$
\end{enumerate}
\end{prob}

\begin{prob}[M]{1428}
研究函数
\[
f(x) = \begin{cases}
|x|\left(2 + \cos\dfrac{1}{x}\right), & x \neq 0, \\
0, & x = 0
\end{cases}
\]
在点 $x = 0$ 处的极值，并作出此函数的图像.
\end{prob}

\probgroup{求下列函数的极值：}

\begin{prob}[S]{1429}
$y = x^3 - 6x^2 + 9x - 4$.
\end{prob}

\begin{prob}[S]{1430}
$y = 2x^2 - x^4$.
\end{prob}

\begin{prob}[S]{1431}
$y = x(x - 1)^2 (x - 2)^3$.
\end{prob}

\begin{prob}[S]{1432}
$y = x + \dfrac{1}{x}$.
\end{prob}

\begin{prob}[S]{1433}
$y = \dfrac{2x}{1 + x^2}$.
\end{prob}

\begin{prob}[S]{1434}
$y = \dfrac{x^2 - 3x + 2}{x^2 + 2x + 1}$.
\end{prob}

\begin{prob}[S]{1435}
$y = \sqrt{2x - x^2}$.
\end{prob}

\begin{prob}[S]{1436}
$y = x\sqrt[3]{x - 1}$.
\end{prob}

\begin{prob}[S]{1437}
$y = x\mathrm{e}^{-x}$.
\end{prob}

\begin{prob}[S]{1438}
$y = \sqrt{x} \ln x$.
\end{prob}

\begin{prob}[S]{1439}
$y = \dfrac{\ln^2 x}{x}$.
\end{prob}

\begin{prob}[S]{1440}
$y = \cos x + \dfrac{1}{2} \cos 2x$.
\end{prob}

\begin{prob}[S]{1441}
$y = \dfrac{10}{1 + \sin^2 x}$.
\end{prob}

\begin{prob}[S]{1442}
$y = \arctan x - \dfrac{1}{2} \ln(1 + x^2)$.
\end{prob}

\begin{prob}[S]{1443}
$y = \mathrm{e}^x \sin x$.
\end{prob}

\begin{prob}[S]{1444}
$y = |x|\mathrm{e}^{-|x - 1|}$.
\end{prob}

\probgroup{求下列函数在所给闭区间上的最大值和最小值：}

\begin{prob}[S]{1445}
$f(x) = 2^x, \ [-1, 5]$.
\end{prob}

\begin{prob}[S]{1446}
$f(x) = x^2 - 4x + 6, \ [-3, 10]$.
\end{prob}

\begin{prob}[S]{1447}
$f(x) = |x^2 - 3x + 2|, \ [-10, 10]$.
\end{prob}

\begin{prob}[S]{1448}
$f(x) = x + \dfrac{1}{x}, \ [0.01, 100]$.
\end{prob}

\begin{prob}[S]{1449}
$f(x) = \sqrt{5 - 4x}, \ [-1, 1]$.
\end{prob}

\probgroup{求下列函数在所给区间上的下确界 ($\inf$) 与上确界 ($\sup$)：}

\begin{prob}[S]{1450}
$f(x) = x\mathrm{e}^{-0.01x}, \ (0, +\infty)$.
\end{prob}

\begin{prob}[S]{1451}
$f(x) = \left(1 + x + \dfrac{x^2}{2!} + \cdots + \dfrac{x^n}{n!}\right)\mathrm{e}^{-x}, \ (0, +\infty)$.
\end{prob}

\begin{prob}[S]{1452}
$f(x) = \dfrac{1 + x^2}{1 + x^4}, \ (0, +\infty)$.
\end{prob}

\begin{prob}[S]{1453}
$f(x) = \mathrm{e}^{-x^2} \cos x^2, \ (-\infty, +\infty)$.
\end{prob}

\begin{prob}[M]{1454.1}
求函数 $f(\xi) = \dfrac{1 + \xi}{3 + \xi^2}$ 在区间 $x < \xi < +\infty$ 内的下确界与上确界. 作出下列函数的图像：
\[
M(x) = \sup_{x < \xi < +\infty} f(\xi), \quad m(x) = \inf_{x < \xi < +\infty} f(\xi).
\]
\end{prob}

\begin{prob}[M]{1454.2}
设
\[
M_k = \sup_x |f^{(k)}(x)| \quad (k = 0, 1, 2, \cdots).
\]
若 $f(x) = \mathrm{e}^{-x^2}$，求 $M_0, M_1$ 和 $M_2$.
\end{prob}

\begin{prob}[L]{1455}
求以下各数列的最大项：
\begin{enumerate}[label=(\alph*), itemsep=2pt, topsep=3pt]
    \item $\dfrac{n^{10}}{2^n} \ (n = 1, 2, \cdots)$；
    \item $\dfrac{\sqrt{n}}{n + 10\,000} \ (n = 1, 2, \cdots)$；
    \item $\sqrt[n]{n} \ (n = 1, 2, \cdots)$.
\end{enumerate}
\end{prob}

\begin{prob}[L]{1456.1}
证明下列不等式：
\begin{enumerate}[label=(\alph*), itemsep=2pt, topsep=3pt]
    \item $|3x - x^3| \leqslant 2 \ (|x| \leqslant 2)$；
    \item $\dfrac{1}{2^{p - 1}} \leqslant x^p + (1 - x)^p \leqslant 1 \ (0 \leqslant x \leqslant 1, \, p > 1)$；
    \item $x^m (a - x)^n \leqslant \dfrac{m^m n^n}{(m + n)^{m + n}} a^{m + n} \ (0 \leqslant x \leqslant a, \, m > 0, \, n > 0)$；
    \item $\dfrac{x + a}{2^{\frac{n - 1}{n}}} \leqslant \sqrt[n]{x^n + a^n} \leqslant x + a \ (x > 0, \, a > 0, \, n > 1)$；
    \item $|a\sin x + b\cos x| \leqslant \sqrt{a^2 + b^2}$.
\end{enumerate}
\end{prob}

\begin{prob}[M]{1456.2}
证明不等式：
\[
\frac{2}{3} \leqslant \frac{x^2 + 1}{x^2 + x + 1} \leqslant 2 \quad (-\infty < x < +\infty).
\]
\end{prob}

\begin{prob}[M]{1457}
求多项式
\[
P(x) = x(x - 1)^2 (x + 2)
\]
在闭区间 $[-2, 1]$ 上``与零的偏差''，就是求
\[
E_P = \sup_{-2 \leqslant x \leqslant 1} |P(x)|.
\]
\end{prob}

\begin{prob}[M]{1458}
应当选择怎样的系数 $q$，使多项式
\[
P(x) = x^2 + q
\]
在闭区间 $[-1, 1]$ 上与零的偏差最小，即
\[
E_P = \sup_{-1 \leqslant x \leqslant 1} |P(x)| = \min ?
\]
\end{prob}

\begin{prob}[M]{1459}
数
\[
\Delta = \sup_{a \leqslant x \leqslant b} |f(x) - g(x)|
\]
称为函数 $f(x)$ 与 $g(x)$ 在闭区间 $[a, b]$ 上的绝对偏差.

求函数 $f(x) = x^2$ 和 $g(x) = x^3$ 在闭区间 $[0, 1]$ 上的绝对偏差.
\end{prob}

\begin{prob}[M]{1460}
在闭区间 $[x_1, x_2]$ 上用线性函数
\[
g(x) = (x_1 + x_2)x + b
\]
近似地代替函数
\[
f(x) = x^2,
\]
使函数 $f(x)$ 与 $g(x)$ 的绝对偏差（参阅上题）最小，并求此最小的绝对偏差.
\end{prob}

\begin{prob}[M]{1461}
求函数
\[
f(x) = \max\{2|x|, |1+x|\}
\]
的极小值.
\end{prob}

\probgroup{确定下列各方程实根的数目，并划分出这些根所在的区间：}

\begin{prob}[S]{1462}
$x^3 - 6x^2 + 9x - 10 = 0.$
\end{prob}

\begin{prob}[S]{1463}
$x^3 - 3x^2 - 9x + h = 0.$
\end{prob}

\begin{prob}[S]{1464}
$3x^4 - 4x^3 - 6x^2 + 12x - 20 = 0.$
\end{prob}

\begin{prob}[S]{1465}
$x^5 - 5x = a.$
\end{prob}

\begin{prob}[S]{1466}
$\ln x = kx.$
\end{prob}

\begin{prob}[S]{1467}
$\mathrm{e}^x = ax^2.$
\end{prob}

\begin{prob}[S]{1468}
当 $0 \leqslant x \leqslant \pi$ 时，$\sin^3 x \cdot \cos x = a.$
\end{prob}

\begin{prob}[S]{1469}
$\cosh x = kx.$
\end{prob}

\begin{prob}[M]{1470}
在什么条件下方程
\[
x^3 + px + q = 0
\]
有：(a) 一个实根；(b) 三个实根？在平面 $(p, q)$ 上画出相应的区域.
\end{prob}

\probsec{§12. 依据函数的特征点作函数图像}

\probgroup{作出下列函数的图像：}

\begin{prob}[S]{1471}
$y = 3x - x^3.$
\end{prob}

\begin{prob}[S]{1472}
$y = 1 + x^2 - \frac{x^4}{2}.$
\end{prob}

\begin{prob}[S]{1473}
$y = (x+1)(x-2)^2.$
\end{prob}

\begin{prob}[S]{1474*}
$y = \frac{2 - x^2}{1 + x^4}.$
\end{prob}

\begin{prob}[S]{1475*}
$y = \frac{x^2 - 1}{x^2 - 5x + 6}.$
\end{prob}

\begin{prob}[S]{1476*}
$y = \frac{x}{(1+x)(1-x)^2}.$
\end{prob}

\begin{prob}[S]{1477}
$y = \frac{x^4}{(1+x)^3}.$
\end{prob}

\begin{prob}[S]{1478}
$y = \left(\frac{1+x}{1-x}\right)^4.$
\end{prob}

\begin{prob}[S]{1479}
$y = \frac{x^2(x-1)}{(x+1)^2}.$
\end{prob}

\begin{prob}[S]{1480}
$y = \frac{x}{(1-x^2)^2}.$
\end{prob}

\begin{prob}[S]{1481}
$y = \frac{(x+1)^3}{(x-1)^2}.$
\end{prob}

\begin{prob}[S]{1482*}
$y = \frac{x^4 + 8}{x^3 + 1}.$
\end{prob}

\begin{prob}[S]{1483}
$y = \frac{1}{1+x} - \frac{10}{3x^2} + \frac{1}{1-x}.$
\end{prob}

\begin{prob}[S]{1484}
$y = (x-3)\sqrt{x}.$
\end{prob}

\begin{prob}[M]{1485}
(a) $y = \pm\sqrt{8x^2 - x^4}$; \qquad (b) $y = \frac{x-2}{\sqrt{x^2+1}}.$
\end{prob}

\begin{prob}[S]{1486}
$y = \pm\sqrt{(x-1)(x-2)(x-3)}.$
\end{prob}

\begin{prob}[S]{1487*}
$y = \sqrt[3]{x^3 - x^2 - x + 1}.$
\end{prob}

\begin{prob}[S]{1488}
$y = \sqrt[3]{x^2} - \sqrt[3]{x^2 + 1}.$
\end{prob}

\begin{prob}[S]{1489}
$y = (x+2)^{\frac{2}{3}} - (x-2)^{\frac{2}{3}}.$
\end{prob}

\begin{prob}[S]{1490}
$y = (x+1)^{\frac{2}{3}} + (x-1)^{\frac{2}{3}}.$
\end{prob}

\begin{prob}[S]{1491}
$y = \frac{x}{\sqrt[3]{x^2 - 1}}.$
\end{prob}

\begin{prob}[S]{1492}
$y = \frac{x^2\sqrt{x^2 - 1}}{2x^2 - 1}.$
\end{prob}

\begin{prob}[S]{1493}
$y = \frac{|1+x|^{\frac{3}{2}}}{\sqrt{x}}.$
\end{prob}

\begin{prob}[S]{1494}
$y = 1 - x + \sqrt{\frac{x^3}{3+x}}.$
\end{prob}

\begin{prob}[S]{1495}
$y = \sqrt[3]{\frac{x^2}{x+1}}.$
\end{prob}

\begin{prob}[S]{1496*}
$y = \sqrt{\frac{x^4 + 3}{x^2 + 1}}.$
\end{prob}

\begin{prob}[S]{1497}
$y = \sin x + \cos^2 x.$
\end{prob}

\begin{prob}[S]{1498}
$y = (7 + 2\cos x)\sin x.$
\end{prob}

\begin{prob}[S]{1499}
$y = \sin x + \frac{1}{3}\sin 3x.$
\end{prob}

\begin{prob}[S]{1500}
$y = \cos x - \frac{1}{2}\cos 2x.$
\end{prob}

\begin{prob}[S]{1501}
$y = \sin^4 x + \cos^4 x.$
\end{prob}

\begin{prob}[S]{1502}
$y = \sin x \cdot \sin 3x.$
\end{prob}

\begin{prob}[S]{1503}
$y = \frac{\sin x}{\sin\left(x + \dfrac{\pi}{4}\right)}.$
\end{prob}

\begin{prob}[M]{1504}
(a) $y = \frac{\cos x}{\cos 2x}$; \qquad (b) $y = \frac{\sin x}{2 + \cos x}.$
\end{prob}

\begin{prob}[S]{1505}
$y = 2x - \tan x.$
\end{prob}

\begin{prob}[S]{1506}
$y = \mathrm{e}^{2x - x^2}.$
\end{prob}

\begin{prob}[S]{1507}
$y = (1 + x^2)\mathrm{e}^{-x^2}.$
\end{prob}

\begin{prob}[S]{1508}
$y = x + \mathrm{e}^{-x}.$
\end{prob}

\begin{prob}[M]{1509}
(a) $y = x^{\frac{2}{3}}\mathrm{e}^{-x}$; \qquad (b) $y = \mathrm{e}^{-2x}\sin^2 x.$
\end{prob}

\begin{prob}[S]{1510}
$y = \frac{\mathrm{e}^x}{1+x}.$
\end{prob}

\begin{prob}[S]{1511}
$y = \sqrt{1 - \mathrm{e}^{-x^2}}.$
\end{prob}

\begin{prob}[S]{1512}
$y = \frac{\ln x}{\sqrt{x}}.$
\end{prob}

\begin{prob}[S]{1513}
$y = \ln(x + \sqrt{x^2 + 1}).$
\end{prob}

\begin{prob}[S]{1514}
$y = \sqrt{x^2 + 1}\cdot\ln(x + \sqrt{x^2 + 1}).$
\end{prob}

\begin{prob}[S]{1515}
$y = \frac{\arcsin x}{\sqrt{1 - x^2}}.$
\end{prob}

\begin{prob}[S]{1516}
$y = x + \arctan x.$
\end{prob}

\begin{prob}[S]{1517}
$y = \frac{x}{2} + \operatorname{arccot} x.$
\end{prob}

\begin{prob}[S]{1518}
$y = x\arctan x.$
\end{prob}

\begin{prob}[S]{1519}
$y = \arcsin\frac{2x}{1 + x^2}.$
\end{prob}

\begin{prob}[S]{1520}
$y = \arccos\frac{1 - x^2}{1 + x^2}.$
\end{prob}

\begin{prob}[S]{1521}
$y = (x+2)\mathrm{e}^{\frac{1}{x}}.$
\end{prob}

\begin{prob}[S]{1522}
$y = 2^{\sqrt{x^2+1} - \sqrt{x^2-1}}.$
\end{prob}

\begin{prob}[S]{1523*}
$y = \ln\frac{x^2 - 3x + 2}{x^2 + 1}.$
\end{prob}

\begin{prob}[S]{1524}
$y = a\arcsin\frac{x}{a} - \sqrt{a^2 - x^2}\ (a > 0).$
\end{prob}

\begin{prob}[S]{1525}
$y = \arccos\frac{1 - x}{1 - 2x}.$
\end{prob}

\begin{prob}[S]{1526}
$y = x^x.$
\end{prob}

\begin{prob}[S]{1527}
$y = x^{\frac{1}{x}}.$
\end{prob}

\begin{prob}[S]{1528}
$y = (1 + x)^{\frac{1}{x}}.$
\end{prob}

\begin{prob}[S]{1529*}
$y = x\left(1 + \frac{1}{x}\right)^x\ (x > 0).$
\end{prob}

\begin{prob}[S]{1530*}
$y = \frac{\mathrm{e}^{\frac{1}{1 - x^2}}}{1 + x^2}$ (不研究凹凸性).
\end{prob}

\probgroup{作出下列参数方程所表示的曲线：}

\begin{prob}[S]{1531}
$x = \frac{(t+1)^2}{4},\ y = \frac{(t-1)^2}{4}.$
\end{prob}

\begin{prob}[S]{1532}
$x = 2t - t^2,\ y = 3t - t^3.$
\end{prob}

\begin{prob}[S]{1533*}
$x = \frac{t^2}{t - 1},\ y = \frac{t}{t^2 - 1}.$
\end{prob}

\begin{prob}[S]{1534}
$x = \frac{t^2}{1 - t^2},\ y = \frac{1}{1 + t^2}.$
\end{prob}

\begin{prob}[S]{1535}
$x = t + \mathrm{e}^{-t},\ y = 2t + \mathrm{e}^{-2t}.$
\end{prob}

\begin{prob}[S]{1536}
$x = a\cos 2t,\ y = a\cos 3t\ (a > 0).$
\end{prob}

\begin{prob}[S]{1537}
$x = \cos^4 t,\ y = \sin^4 t.$
\end{prob}

\begin{prob}[S]{1538}
$x = t\ln t,\ y = \frac{\ln t}{t}.$
\end{prob}

\begin{prob}[S]{1539}
$x = \frac{a}{\cos^3 t},\ y = a\tan^3 t\ (a > 0).$
\end{prob}

\begin{prob}[S]{1540}
$x = a(\sinh t - t),\ y = a(\cosh t - 1)\ (a > 0).$
\end{prob}

\probgroup{把下列曲线方程化为参数方程，然后作出这些曲线的图像：}

\begin{prob}[M]{1541}
$x^3 + y^3 - 3axy = 0\ (a > 0).$
\probhint{令 $y = tx$.}
\end{prob}

\begin{prob}[S]{1542}
$x^2 + y^2 = x^4 + y^4.$
\end{prob}

\begin{prob}[S]{1543}
$x^2 y^2 = x^3 - y^3.$
\end{prob}

\begin{prob}[S]{1544}
$x^y = y^x\ (x > 0, y > 0).$
\end{prob}

\begin{prob}[S]{1545}
作出曲线 $\cosh^2 x - \cosh^2 y = 1$ 的图像.
\end{prob}

\probgroup{作出下列用极坐标 $(\varphi, r)\ (r \geqslant 0)$ 表示的函数的图像：}

\begin{prob}[S]{1546}
$r = a + b\cos\varphi\ (0 < a \leqslant b).$
\end{prob}

\begin{prob}[S]{1547}
$r = a\sin 3\varphi\ (a > 0).$
\end{prob}

\begin{prob}[S]{1548}
$r = \frac{a}{\sqrt{\cos 3\varphi}}\ (a > 0).$
\end{prob}

\begin{prob}[S]{1549*}
$r = a\frac{\tanh\varphi}{\varphi - 1}$，其中 $\varphi > 1\ (a > 0).$
\end{prob}

\begin{prob}[S]{1550*}
$\varphi = \arccos\frac{r - 1}{r^2}.$
\end{prob}

\probgroup{作出下列曲线族的图像（$a$ 为参变量）：}

\begin{prob}[S]{1551}
$y = x^2 - 2x + a.$
\end{prob}

\begin{prob}[S]{1552}
$y = x + \frac{a^2}{x}.$
\end{prob}

\begin{prob}[S]{1553}
$y = x \pm \sqrt{a(1 - x^2)}.$
\end{prob}

\begin{prob}[S]{1554}
$y = \frac{x}{2} + \mathrm{e}^{-ax}.$
\end{prob}

\begin{prob}[S]{1555}
$y = x\mathrm{e}^{-\frac{x}{a}}.$
\end{prob}

\probsec{§13. 函数的极大值与极小值问题}

\begin{prob}[M]{1556}
证明：若函数 $f(x)$ 不为负，则函数
\[
F(x) = Cf^2(x)\ (C > 0)
\]
与函数 $f(x)$ 有相同的极值点.
\end{prob}

\begin{prob}[M]{1557}
证明：若当 $-\infty < x < +\infty$ 时函数 $\varphi(x)$ 严格单调递增，则函数
\[
f(x) \quad \text{和} \quad \varphi(f(x))
\]
有相同的极值点.
\end{prob}

\begin{prob}[M]{1558}
若二正数之和等于常数 $a$，求此二数的 $m$ 次幂与 $n$ 次幂 $(m > 0, n > 0)$ 之积的最大值.
\end{prob}

\begin{prob}[M]{1559}
若二正数之积等于常数 $a$，求此二数的 $m$ 次幂与 $n$ 次幂 $(m > 0, n > 0)$ 之和的最小值.
\end{prob}

\begin{prob}[M]{1560}
当对数之底取何值时存在这样的数，它本身和它的对数相等？
\end{prob}

\begin{prob}[M]{1561}
从面积为给定值 $S$ 的一切矩形中，求其周长为最小者.
\end{prob}

\begin{prob}[M]{1562}
若直角三角形的一直角边与斜边之和为常数，求具有最大面积的直角三角形.
\end{prob}

\begin{prob}[M]{1563}
要使容积为给定值 $V$ 的圆柱形闭合容器具有最小的表面积，其尺寸如何？
\end{prob}

\begin{prob}[M]{1564}
在不超过半圆的给定弓形内作出具有最大面积的内接矩形.
\end{prob}

\begin{prob}[M]{1565}
在椭圆
\[
\frac{x^2}{a^2} + \frac{y^2}{b^2} = 1
\]
内作出具有最大面积而边平行于椭圆轴的内接矩形.
\end{prob}

\begin{prob}[M]{1566}
在底边为 $b$ 及高为 $h$ 的三角形内作出具有最大周长的内接矩形. 研究此问题有解的可能性.
\end{prob}

\begin{prob}[M]{1567}
从直径为 $d$ 的原木（圆形横截面）切出横截面为矩形的梁，此矩形的底等于 $b$，高等于 $h$. 若梁的强度\footnote{此处的强度指梁的截面模量. —— 译注}与 $bh^2$ 成正比，问梁的尺寸如何其强度最大？
\end{prob}

\begin{prob}[M]{1568}
在半径为 $R$ 的半球中作出具有最大体积且底为正方形的内接长方体.
\end{prob}

\begin{prob}[M]{1569}
在半径为 $R$ 的球内作出具有最大体积的内接圆柱体.
\end{prob}

\begin{prob}[M]{1570}
在半径为 $R$ 的球内作出具有最大表面积的内接圆柱体.
\end{prob}

\begin{prob}[M]{1571}
在已知球外作出具有最小体积的外切圆锥体.
\end{prob}

\begin{prob}[M]{1572}
求母线为给定值 $l$ 的圆锥之最大体积.
\end{prob}

\begin{prob}[M]{1573}
在顶角为 $2\alpha$ 底半径为 $R$ 的直圆锥体中作出具有最大表积的内接圆柱体.
\end{prob}

\begin{prob}[M]{1574}
求从点 $M(p, p)$ 到抛物线 $y^2 = 2px$ 的最短距离.
\end{prob}

\begin{prob}[M]{1575}
求从点 $A(2, 0)$ 到圆 $x^2 + y^2 = 1$ 的最短与最长距离.
\end{prob}

\begin{prob}[M]{1576}
求椭圆 $\dfrac{x^2}{a^2} + \dfrac{y^2}{b^2} = 1\ (0 < b < a)$ 的经过顶点 $B(0, -b)$ 的最大弦.
\end{prob}

\begin{prob}[M]{1577}
过椭圆 $\dfrac{x^2}{a^2} + \dfrac{y^2}{b^2} = 1$ 上的点 $M(x, y)$ 引切线，使由此切线与坐标轴构成的三角形具有最小的面积.
\end{prob}

\begin{prob}[M]{1578}
直径相同的圆柱体与半球体拼接在一起构成一物体，其体积为 $V$. 要使此物体具有最小的表面积，其尺寸如何？
\end{prob}

\begin{prob}[M]{1579}
设明渠的横截面为等腰梯形，渠中流水的横截面面积为 $S$，水面的高为 $h$. 问水渠侧边的倾角 $\varphi$ 如何，才使其横截面边界被水浸湿的部分具有最小的长度？
\end{prob}

\begin{prob}[M]{1580}
设封闭曲线所围面积为 $S$，则该曲线的周长与面积同为 $S$ 的圆的周长之比称为封闭曲线的弯曲度.

设等腰梯形 $ABCD\ (AD // BC)$ 的底边 $AD = 2a$，锐角 $BAD = \alpha$，问等腰梯形的形状如何，才有最小的弯曲度？
\end{prob}

\begin{prob}[M]{1581}
从半径为 $R$ 的圆中应切去怎样的扇形，才能使余下的部分可卷成一漏斗，其容积为最大？
\end{prob}

\begin{prob}[M]{1582}
从南至北的铁路经过 $B$ 城，工厂 $A$ 距此铁路的最短距离为 $a$，距偏北方的 $B$ 城所在纬线的距离为 $b$. 为了从 $A$ 至 $B$ 运输货物最经济，从工厂修建一条专线铁路. 若每吨货物沿专线运输的价格是 $p\ \text{元/km}$，而沿常规铁路为 $q\ \text{元/km}\ (p > q)$，则支线应向常规铁路取怎样的角度 $\varphi$？
\end{prob}

\begin{prob}[M]{1583}
二船各以恒定的速度 $u$ 和 $v$ 沿直线前进，二者前进方向所成的角为 $\theta$. 若在某时刻它们到其路线交点的距离分别为 $a$ 和 $b$，求二船的最小距离.
\end{prob}

\begin{prob}[M]{1584}
在 $A$ 与 $B$ 二点处各有一光源，其发光强度分别为 $S_1$ 与 $S_2$. 在线段 $AB = a$ 上求出最小照度的点 $M$.
\end{prob}

\begin{prob}[M]{1585}
发光点位于半径为 $R$ 与 $r\ (R > r)$ 的二互不相交之球的连心线上，并在此二球的外面. 此发光点的位置如何，才可使二球表面上被照明部分之和为最大？
\end{prob}

\begin{prob}[M]{1586}
设圆桌面的半径为 $a$. 应当在正对桌面中央多高的地方安置电灯，才可使桌面边缘的照度为最大？
\probhint{照度可用下列公式表示：
\[
I = I_0\frac{\sin\varphi}{r^2},
\]
式中 $\varphi$ 为光线对桌面的倾角，$r$ 为光源与被照面的距离，$I_0$ 为光源的发光强度.}
\end{prob}

\begin{prob}[M]{1587}
向宽为 $a$ 的河修建一宽为 $b$ 的运河，二者成直角相交，问能驶进这运河的船，其最大的长度如何？
\end{prob}

\begin{prob}[M]{1588}
船航行一昼夜的耗费由两部分组成：固定部分等于 $a$ 元，变动部分与速度的立方成正比增加. 在怎样的速度 $v$ 时，航行最为经济？
\end{prob}

\begin{prob}[M]{1589}
重量为 $P$ 的物体位于粗糙的水平面上，需用力把物体从原位置移动. 若物体摩擦因子等于 $k$，问作用力对水平面的倾斜程度如何，才使所需的力为最小？
\end{prob}

\begin{prob}[M]{1590}
有一茶杯，其形状是半径为 $a$ 的半球，在茶杯中放一长为 $l > 2a$ 的杆，求杆的平衡位置.
\end{prob}

\probsec{§14. 曲线的相切. 曲率圆. 渐屈线}

\begin{prob}[M]{1591}
选择直线
\[
y = kx + b
\]
的参数 $k$ 与 $b$，使它与曲线
\[
y = x^3 - 3x^2 + 2
\]
二阶或更高阶相切.
\end{prob}

\begin{prob}[M]{1592}
应该怎样选择系数 $a, b$ 和 $c$，才能使抛物线
\[
y = ax^2 + bx + c
\]
在点 $x = x_0$ 与曲线 $y = \mathrm{e}^x$ 二阶相切？
\end{prob}

\begin{prob}[L]{1593}
下列曲线与 $Ox$ 轴在点 $x = 0$ 相切的阶如何：
\begin{enumerate}[label=(\alph*), itemsep=2pt, topsep=3pt]
\item $y = 1 - \cos x$;
\item $y = \tan x - \sin x$;
\item $y = \mathrm{e}^x - \left(1 + x + \frac{x^2}{2}\right)$.
\end{enumerate}
\end{prob}

\begin{prob}[M]{1594}
证明：曲线
\[
y = \begin{cases} \mathrm{e}^{-\frac{1}{x^2}}, & x \neq 0, \\ 0, & x = 0 \end{cases}
\]
在点 $x = 0$ 与 $Ox$ 轴相切的阶为无穷大.
\end{prob}

\begin{prob}[M]{1595}
求双曲线
\[
xy = 1
\]
在下列各点的曲率半径和曲率中心：
(a) $M(1, 1)$; \qquad (b) $N(100, 0.01)$.
\end{prob}

\probgroup{求下列曲线的曲率半径：}

\begin{prob}[S]{1596}
抛物线 $y^2 = 2px$.
\end{prob}

\begin{prob}[S]{1597}
椭圆 $\frac{x^2}{a^2} + \frac{y^2}{b^2} = 1 \ (a \geqslant b > 0)$.
\end{prob}

\begin{prob}[S]{1598}
双曲线 $\frac{x^2}{a^2} - \frac{y^2}{b^2} = 1$.
\end{prob}

\begin{prob}[S]{1599}
星形线 $x^{\frac{2}{3}} + y^{\frac{2}{3}} = a^{\frac{2}{3}}$.
\end{prob}

\begin{prob}[S]{1600}
椭圆 $x = a\cos t,\ y = b\sin t$.
\end{prob}

\begin{prob}[S]{1601}
摆线 $x = a(t - \sin t),\ y = a(1 - \cos t)$.
\end{prob}

\begin{prob}[S]{1602}
圆的渐伸线 $x = a(\cos t + t\sin t),\ y = a(\sin t - t\cos t)$.
\end{prob}

\begin{prob}[M]{1603}
证明：二次曲线
\[
y^2 = 2px - qx^2
\]
的曲率半径与法线段的立方成正比.
\end{prob}

\begin{prob}[S]{1604}
写出以极坐标表示的曲线的曲率半径公式.
\end{prob}

\probgroup{求下列极坐标方程所表示的曲线的曲率半径：}

\begin{prob}[S]{1605}
阿基米德螺线 $r = a\varphi$.
\end{prob}

\begin{prob}[S]{1606}
对数螺线 $r = a\mathrm{e}^{m\varphi}$.
\end{prob}

\begin{prob}[S]{1607}
心脏线 $r = a(1 + \cos\varphi)$.
\end{prob}

\begin{prob}[S]{1608}
双纽线 $r^2 = a^2 \cos 2\varphi$.
\end{prob}

\begin{prob}[M]{1609}
在曲线 $y = \ln x$ 上求曲率最大的点.
\end{prob}

\begin{prob}[M]{1610}
三次抛物线 $y = \frac{kx^3}{6} \ (0 \leqslant x < +\infty, k > 0)$ 的最大曲率等于 $\frac{1}{1000}$，求达到此最大曲率的点 $x$.
\end{prob}

\probgroup{求下列各曲线的渐屈线方程：}

\begin{prob}[S]{1611}
抛物线 $y^2 = 2px$.
\end{prob}

\begin{prob}[S]{1612}
椭圆 $\frac{x^2}{a^2} + \frac{y^2}{b^2} = 1$.
\end{prob}

\begin{prob}[S]{1613}
星形线 $x^{\frac{2}{3}} + y^{\frac{2}{3}} = a^{\frac{2}{3}}$.
\end{prob}

\begin{prob}[S]{1614}
曳物线 $x = a\ln\frac{a + \sqrt{a^2 - y^2}}{y} - \sqrt{a^2 - y^2}$.
\end{prob}

\begin{prob}[S]{1615}
对数螺线 $r = a\mathrm{e}^{m\varphi}$.
\end{prob}

\begin{prob}[M]{1616}
证明：摆线
\[
x = a(t - \sin t),\quad y = a(1 - \cos t)
\]
的渐屈线仍为一摆线，仅其位置与所给摆线不同而已.
\end{prob}

\probsec{§15. 方程的近似解法}

\probgroup{利用比例法求下列方程的根（精确到 $0.001$）：}

\begin{prob}[S]{1617}
$x^3 - 6x + 2 = 0$.
\end{prob}

\begin{prob}[S]{1618}
$x^4 - x - 1 = 0$.
\end{prob}

\begin{prob}[S]{1619}
$x - 0.1\sin x = 2$.
\end{prob}

\begin{prob}[S]{1620}
$\cos x = x^2$.
\end{prob}

\probgroup{利用牛顿法求下列方程的根（精确到所指定的精度）：}

\begin{prob}[S]{1621}
$x^2 + \frac{1}{x^2} = 10x$（精确到 $10^{-3}$）.
\end{prob}

\begin{prob}[S]{1622}
$x\lg x = 1$（精确到 $10^{-4}$）.
\end{prob}

\begin{prob}[S]{1623}
$\cos x \cdot \cosh x = 1$（精确到 $10^{-3}$）（二正根）.
\end{prob}

\begin{prob}[S]{1624}
$x + \mathrm{e}^x = 0$（精确到 $10^{-5}$）.
\end{prob}

\begin{prob}[S]{1625}
$x\tanh x = 1$（精确到 $10^{-6}$）.
\end{prob}

\begin{prob}[M]{1626}
求方程
\[
\tan x = x
\]
最小的三个正根（精确到 $0.001$）.
\end{prob}

\begin{prob}[M]{1627}
求方程
\[
\cot x = \frac{1}{x} - \frac{x}{2}
\]
的二正根（精确到 $10^{-3}$）.
\end{prob}

\probchapter{第三章\quad 不定积分}

\probsec{§1. 最简单的不定积分}

\probgroup{利用最简积分表，求下列积分：}

\begin{prob}[S]{1628}
$\displaystyle\int (3 - x^2)^3 \,\mathrm{d}x.$
\end{prob}

\begin{prob}[S]{1629}
$\displaystyle\int x^2 (5 - x)^4 \,\mathrm{d}x.$
\end{prob}

\begin{prob}[S]{1630}
$\displaystyle\int (1 - x)(1 - 2x)(1 - 3x) \,\mathrm{d}x.$
\end{prob}

\begin{prob}[S]{1631}
$\displaystyle\int \left(\frac{1 - x}{x}\right)^2 \,\mathrm{d}x.$
\end{prob}

\begin{prob}[S]{1632}
$\displaystyle\int \left(\frac{a}{x} + \frac{a^2}{x^2} + \frac{a^3}{x^3}\right) \,\mathrm{d}x.$
\end{prob}

\begin{prob}[S]{1633}
$\displaystyle\int \frac{x + 1}{\sqrt{x}} \,\mathrm{d}x.$
\end{prob}

\begin{prob}[S]{1634}
$\displaystyle\int \frac{\sqrt{x} - 2\sqrt[3]{x^2} + 1}{\sqrt[4]{x}} \,\mathrm{d}x.$
\end{prob}

\begin{prob}[S]{1635}
$\displaystyle\int \frac{(1 - x)^3}{x\sqrt[3]{x}} \,\mathrm{d}x.$
\end{prob}

\begin{prob}[S]{1636}
$\displaystyle\int \left(1 - \frac{1}{x^2}\right)\sqrt{x\sqrt{x}} \,\mathrm{d}x.$
\end{prob}

\begin{prob}[S]{1637}
$\displaystyle\int \frac{(\sqrt{2x} - \sqrt[3]{3x})^2}{x} \,\mathrm{d}x.$
\end{prob}

\begin{prob}[S]{1638}
$\displaystyle\int \frac{\sqrt{x^4 + x^{-4} + 2}}{x^3} \,\mathrm{d}x.$
\end{prob}

\begin{prob}[S]{1639}
$\displaystyle\int \frac{x^2 \,\mathrm{d}x}{1 + x^2}.$
\end{prob}

\begin{prob}[S]{1640}
$\displaystyle\int \frac{x^2 \,\mathrm{d}x}{1 - x^2}.$
\end{prob}

\begin{prob}[S]{1641}
$\displaystyle\int \frac{x^2 + 3}{x^2 - 1} \,\mathrm{d}x.$
\end{prob}

\begin{prob}[S]{1642}
$\displaystyle\int \frac{\sqrt{1 + x^2} + \sqrt{1 - x^2}}{\sqrt{1 - x^4}} \,\mathrm{d}x.$
\end{prob}

\begin{prob}[S]{1643}
$\displaystyle\int \frac{\sqrt{x^2 + 1} - \sqrt{x^2 - 1}}{\sqrt{x^4 - 1}} \,\mathrm{d}x.$
\end{prob}

\begin{prob}[S]{1644}
$\displaystyle\int (2^x + 3^x)^2 \,\mathrm{d}x.$
\end{prob}

\begin{prob}[S]{1645}
$\displaystyle\int \frac{2^{x+1} - 5^{x-1}}{10^x} \,\mathrm{d}x.$
\end{prob}

\begin{prob}[S]{1646}
$\displaystyle\int \frac{\mathrm{e}^{3x} + 1}{\mathrm{e}^x + 1} \,\mathrm{d}x.$
\end{prob}

\begin{prob}[S]{1647}
$\displaystyle\int (1 + \sin x + \cos x) \,\mathrm{d}x.$
\end{prob}

\begin{prob}[S]{1648}
$\displaystyle\int \sqrt{1 - \sin 2x} \,\mathrm{d}x \quad (0 \leqslant x \leqslant \pi).$
\end{prob}

\begin{prob}[S]{1649}
$\displaystyle\int \cot^2 x \,\mathrm{d}x.$
\end{prob}

\begin{prob}[S]{1650}
$\displaystyle\int \tan^2 x \,\mathrm{d}x.$
\end{prob}

\begin{prob}[S]{1651}
$\displaystyle\int (a\sinh x + b\cosh x) \,\mathrm{d}x.$
\end{prob}

\begin{prob}[S]{1652}
$\displaystyle\int \tanh^2 x \,\mathrm{d}x.$
\end{prob}

\begin{prob}[S]{1653}
$\displaystyle\int \coth^2 x \,\mathrm{d}x.$
\end{prob}

\begin{prob}[M]{1654}
证明：若
\[
\int f(x)\,\mathrm{d}x = F(x) + C,
\]
则
\[
\int f(ax + b)\,\mathrm{d}x = \frac{1}{a}F(ax + b) + C \quad (a \neq 0).
\]
\end{prob}

\probgroup{求下列积分：}

\begin{prob}[S]{1655}
$\displaystyle\int \frac{\mathrm{d}x}{x + a}.$
\end{prob}

\begin{prob}[S]{1656}
$\displaystyle\int (2x - 3)^{10} \,\mathrm{d}x.$
\end{prob}

\begin{prob}[S]{1657}
$\displaystyle\int \sqrt[3]{1 - 3x} \,\mathrm{d}x.$
\end{prob}

\begin{prob}[S]{1658}
$\displaystyle\int \frac{\mathrm{d}x}{\sqrt{2 - 5x}}.$
\end{prob}

\begin{prob}[S]{1659}
$\displaystyle\int \frac{\mathrm{d}x}{(5x - 2)^{\frac{5}{2}}}.$
\end{prob}

\begin{prob}[S]{1660}
$\displaystyle\int \frac{\sqrt[5]{1 - 2x + x^2}}{1 - x} \,\mathrm{d}x.$
\end{prob}

\begin{prob}[S]{1661}
$\displaystyle\int \frac{\mathrm{d}x}{2 + 3x^2}.$
\end{prob}

\begin{prob}[S]{1662}
$\displaystyle\int \frac{\mathrm{d}x}{2 - 3x^2}.$
\end{prob}

\begin{prob}[S]{1663}
$\displaystyle\int \frac{\mathrm{d}x}{\sqrt{2 - 3x^2}}.$
\end{prob}

\begin{prob}[S]{1664}
$\displaystyle\int \frac{\mathrm{d}x}{\sqrt{3x^2 - 2}}.$
\end{prob}

\begin{prob}[S]{1665}
$\displaystyle\int (\mathrm{e}^{-x} + \mathrm{e}^{-2x}) \,\mathrm{d}x.$
\end{prob}

\begin{prob}[S]{1666}
$\displaystyle\int (\sin 5x - \sin 5\alpha) \,\mathrm{d}x.$
\end{prob}

\begin{prob}[S]{1667}
$\displaystyle\int \frac{\mathrm{d}x}{\sin^2\left(2x + \frac{\pi}{4}\right)}.$
\end{prob}

\begin{prob}[S]{1668}
$\displaystyle\int \frac{\mathrm{d}x}{1 + \cos x}.$
\end{prob}

\begin{prob}[S]{1669}
$\displaystyle\int \frac{\mathrm{d}x}{1 - \cos x}.$
\end{prob}

\begin{prob}[S]{1670}
$\displaystyle\int \frac{\mathrm{d}x}{1 + \sin x}.$
\end{prob}

\begin{prob}[S]{1671}
$\displaystyle\int [\sinh(2x + 1) + \cosh(2x - 1)] \,\mathrm{d}x.$
\end{prob}

\begin{prob}[S]{1672}
$\displaystyle\int \frac{\mathrm{d}x}{\cosh^2 \frac{x}{2}}.$
\end{prob}

\begin{prob}[S]{1673}
$\displaystyle\int \frac{\mathrm{d}x}{\sinh^2 \frac{x}{2}}.$
\end{prob}

\probgroup{用适当地变换被积函数的方法求下列积分：}

\begin{prob}[S]{1674}
$\displaystyle\int \frac{x\,\mathrm{d}x}{\sqrt{1 - x^2}}.$
\end{prob}

\begin{prob}[S]{1675}
$\displaystyle\int x^2\sqrt[3]{1 + x^3}\,\mathrm{d}x.$
\end{prob}

\begin{prob}[S]{1676}
$\displaystyle\int \frac{x\,\mathrm{d}x}{3 - 2x^2}.$
\end{prob}

\begin{prob}[S]{1677}
$\displaystyle\int \frac{x\,\mathrm{d}x}{(1 + x^2)^2}.$
\end{prob}

\begin{prob}[S]{1678}
$\displaystyle\int \frac{x\,\mathrm{d}x}{4 + x^2}.$
\end{prob}

\begin{prob}[S]{1679}
$\displaystyle\int \frac{x^3\,\mathrm{d}x}{x^8 - 2}.$
\end{prob}

\begin{prob}[S]{1680}
$\displaystyle\int \frac{\mathrm{d}x}{(1 + x)\sqrt{x}}.$
\probhint{$\frac{\mathrm{d}x}{\sqrt{x}} = 2\,\mathrm{d}(\sqrt{x})$.}
\end{prob}

\begin{prob}[S]{1681}
$\displaystyle\int \sin\frac{1}{x} \cdot \frac{\mathrm{d}x}{x^2}.$
\end{prob}

\begin{prob}[S]{1682}
$\displaystyle\int \frac{\mathrm{d}x}{x\sqrt{x^2 + 1}}.$
\end{prob}

\begin{prob}[S]{1683}
$\displaystyle\int \frac{\mathrm{d}x}{x\sqrt{x^2 - 1}}.$
\end{prob}

\begin{prob}[S]{1684}
$\displaystyle\int \frac{\mathrm{d}x}{(x^2 + 1)^{\frac{3}{2}}}.$
\end{prob}

\begin{prob}[S]{1685}
$\displaystyle\int \frac{x\,\mathrm{d}x}{(x^2 - 1)^{\frac{3}{2}}}.$
\end{prob}

\begin{prob}[S]{1686}
$\displaystyle\int \frac{x^2\,\mathrm{d}x}{(8x^3 + 27)^{\frac{2}{3}}}.$
\end{prob}

\begin{prob}[S]{1687}
$\displaystyle\int \frac{\mathrm{d}x}{\sqrt{x(1 + x)}}.$
\end{prob}

\begin{prob}[S]{1688}
$\displaystyle\int \frac{\mathrm{d}x}{\sqrt{x(1 - x)}}.$
\end{prob}

\begin{prob}[S]{1689}
$\displaystyle\int x\mathrm{e}^{-x^2}\,\mathrm{d}x.$
\end{prob}

\begin{prob}[S]{1690}
$\displaystyle\int \frac{\mathrm{e}^x\,\mathrm{d}x}{2 + \mathrm{e}^x}.$
\end{prob}

\begin{prob}[S]{1691}
$\displaystyle\int \frac{\mathrm{d}x}{\mathrm{e}^x + \mathrm{e}^{-x}}.$
\end{prob}

\begin{prob}[S]{1692}
$\displaystyle\int \frac{\mathrm{d}x}{\sqrt{1 + \mathrm{e}^{2x}}}.$
\end{prob}

\begin{prob}[S]{1693}
$\displaystyle\int \frac{\ln^2 x}{x}\,\mathrm{d}x.$
\end{prob}

\begin{prob}[S]{1694}
$\displaystyle\int \frac{\mathrm{d}x}{x\ln x\ln(\ln x)}.$
\end{prob}

\begin{prob}[S]{1695}
$\displaystyle\int \sin^5 x\cos x\,\mathrm{d}x.$
\end{prob}

\begin{prob}[S]{1696}
$\displaystyle\int \frac{\sin x}{\sqrt{\cos^3 x}}\,\mathrm{d}x.$
\end{prob}

\begin{prob}[S]{1697}
$\displaystyle\int \tan x\,\mathrm{d}x.$
\end{prob}

\begin{prob}[S]{1698}
$\displaystyle\int \cot x\,\mathrm{d}x.$
\end{prob}

\begin{prob}[S]{1699}
$\displaystyle\int \frac{\sin x + \cos x}{\sqrt[3]{\sin x - \cos x}}\,\mathrm{d}x.$
\end{prob}

\begin{prob}[L]{1700}
\begin{enumerate}[label=(\alph*), itemsep=2pt, topsep=3pt]
\item $\displaystyle\int \frac{\sin x\cos x}{\sqrt{a^2\sin^2 x + b^2\cos^2 x}}\,\mathrm{d}x;$
\item $\displaystyle\int \frac{\sin x}{\sqrt{\cos 2x}}\,\mathrm{d}x;$
\item $\displaystyle\int \frac{\cos x}{\sqrt{\cos 2x}}\,\mathrm{d}x;$
\item $\displaystyle\int \frac{\sinh x}{\sqrt{\cosh 2x}}\,\mathrm{d}x.$
\end{enumerate}
\end{prob}

\begin{prob}[S]{1701}
$\displaystyle\int \frac{\mathrm{d}x}{\sin^2 x\sqrt[4]{\cot x}}.$
\end{prob}

\begin{prob}[S]{1702}
$\displaystyle\int \frac{\mathrm{d}x}{\sin^2 x + 2\cos^2 x}.$
\end{prob}

\begin{prob}[S]{1703}
$\displaystyle\int \frac{\mathrm{d}x}{\sin x}.$
\end{prob}

\begin{prob}[S]{1704}
$\displaystyle\int \frac{\mathrm{d}x}{\cos x}.$
\end{prob}

\begin{prob}[S]{1705}
$\displaystyle\int \frac{\mathrm{d}x}{\sinh x}.$
\end{prob}

\begin{prob}[S]{1706}
$\displaystyle\int \frac{\mathrm{d}x}{\cosh x}.$
\end{prob}

\begin{prob}[S]{1707}
$\displaystyle\int \frac{\sinh x\cosh x}{\sqrt{\sinh^4 x + \cosh^4 x}}\,\mathrm{d}x.$
\end{prob}

\begin{prob}[S]{1708}
$\displaystyle\int \frac{\mathrm{d}x}{\cosh^2 x\sqrt[3]{\tanh^2 x}}.$
\end{prob}

\begin{prob}[S]{1709}
$\displaystyle\int \frac{\arctan x}{1 + x^2}\,\mathrm{d}x.$
\end{prob}

\begin{prob}[S]{1710}
$\displaystyle\int \frac{\mathrm{d}x}{(\arcsin x)^2\sqrt{1 - x^2}}.$
\end{prob}

\begin{prob}[S]{1711}
$\displaystyle\int \sqrt{\frac{\ln(x + \sqrt{1 + x^2})}{1 + x^2}}\,\mathrm{d}x.$
\end{prob}

\begin{prob}[S]{1712}
$\displaystyle\int \frac{x^2 + 1}{x^4 + 1}\,\mathrm{d}x.$
\probhint{$\left(1 + \frac{1}{x^2}\right)\,\mathrm{d}x = \mathrm{d}\left(x - \frac{1}{x}\right)$.}
\end{prob}

\begin{prob}[S]{1713}
$\displaystyle\int \frac{x^2 - 1}{x^4 + 1}\,\mathrm{d}x.$
\end{prob}

\begin{prob}[S]{1714}
$\displaystyle\int \frac{x^4\,\mathrm{d}x}{(x^5 + 1)^4}.$
\end{prob}

\begin{prob}[S]{1715}
$\displaystyle\int \frac{x^{\frac{n}{2}}\,\mathrm{d}x}{\sqrt{1 + x^{n+2}}}.$
\end{prob}

\begin{prob}[S]{1716}
$\displaystyle\int \frac{1}{1 - x^2}\ln\frac{1 + x}{1 - x}\,\mathrm{d}x.$
\end{prob}

\begin{prob}[S]{1717}
$\displaystyle\int \frac{\cos x\,\mathrm{d}x}{\sqrt{2 + \cos 2x}}.$
\end{prob}

\begin{prob}[S]{1718}
$\displaystyle\int \frac{\sin x\cos x}{\sin^4 x + \cos^4 x}\,\mathrm{d}x.$
\end{prob}

\begin{prob}[S]{1719}
$\displaystyle\int \frac{2^x \cdot 3^x}{9^x - 4^x}\,\mathrm{d}x.$
\end{prob}

\begin{prob}[S]{1720}
$\displaystyle\int \frac{x\,\mathrm{d}x}{\sqrt{1 + x^2 + \sqrt{(1 + x^2)^3}}}.$
\end{prob}

\probgroup{用分项积分法计算下列积分：}

\begin{prob}[M]{1721}
(a) $\displaystyle\int x^2(2 - 3x^2)^2\,\mathrm{d}x;$ \qquad (b) $\displaystyle\int x(1 - x)^{10}\,\mathrm{d}x.$
\end{prob}

\begin{prob}[S]{1722}
$\displaystyle\int \frac{1 + x}{1 - x}\,\mathrm{d}x.$
\end{prob}

\begin{prob}[S]{1723}
$\displaystyle\int \frac{x^2}{1 + x}\,\mathrm{d}x.$
\end{prob}

\begin{prob}[S]{1724}
$\displaystyle\int \frac{x^3}{3 + x}\,\mathrm{d}x.$
\end{prob}

\begin{prob}[S]{1725}
$\displaystyle\int \frac{(1 + x)^2}{1 + x^2}\,\mathrm{d}x.$
\end{prob}

\begin{prob}[S]{1726}
$\displaystyle\int \frac{(2 - x)^2}{2 - x^2}\,\mathrm{d}x.$
\end{prob}

\begin{prob}[S]{1727}
$\displaystyle\int \frac{x^2}{(1 - x)^{100}}\,\mathrm{d}x.$
\end{prob}

\begin{prob}[S]{1728}
$\displaystyle\int \frac{x^5}{x + 1}\,\mathrm{d}x.$
\end{prob}

\begin{prob}[S]{1729}
$\displaystyle\int \frac{\mathrm{d}x}{\sqrt{x + 1} + \sqrt{x - 1}}.$
\end{prob}

\begin{prob}[S]{1730}
$\displaystyle\int x\sqrt{2 - 5x}\,\mathrm{d}x.$
\probhint{$x \equiv -\frac{1}{5}(2 - 5x) + \frac{2}{5}$.}
\end{prob}

\begin{prob}[S]{1731}
$\displaystyle\int \frac{x\,\mathrm{d}x}{\sqrt[3]{1 - 3x}}.$
\end{prob}

\begin{prob}[S]{1732}
$\displaystyle\int x^3\sqrt[3]{1 + x^2}\,\mathrm{d}x.$
\end{prob}

\begin{prob}[S]{1733}
$\displaystyle\int \frac{\mathrm{d}x}{(x - 1)(x + 3)}.$
\probhint{$1 \equiv \frac{1}{4}[(x + 3) - (x - 1)]$.}
\end{prob}

\begin{prob}[S]{1734}
$\displaystyle\int \frac{\mathrm{d}x}{x^2 + x - 2}.$
\end{prob}

\begin{prob}[S]{1735}
$\displaystyle\int \frac{\mathrm{d}x}{(x^2 + 1)(x^2 + 2)}.$
\end{prob}

\begin{prob}[S]{1736}
$\displaystyle\int \frac{\mathrm{d}x}{(x^2 - 2)(x^2 + 3)}.$
\end{prob}

\begin{prob}[S]{1737}
$\displaystyle\int \frac{x\,\mathrm{d}x}{(x + 2)(x + 3)}.$
\end{prob}

\begin{prob}[S]{1738}
$\int \frac{x\,\mathrm{d}x}{x^4 + 3x^2 + 2}$.
\end{prob}

\begin{prob}[S]{1739}
$\int \frac{\mathrm{d}x}{(x + a)^2(x + b)^2} \quad (a \neq b)$.
\end{prob}

\begin{prob}[S]{1740}
$\int \frac{\mathrm{d}x}{(x^2 + a^2)(x^2 + b^2)} \quad (a^2 \neq b^2)$.
\end{prob}

\begin{prob}[S]{1741}
$\int \sin^2 x\,\mathrm{d}x$.
\end{prob}

\begin{prob}[S]{1742}
$\int \cos^2 x\,\mathrm{d}x$.
\end{prob}

\begin{prob}[S]{1743}
$\int \sin x \sin(x + \alpha)\,\mathrm{d}x$.
\end{prob}

\begin{prob}[S]{1744}
$\int \sin 3x \cdot \sin 5x\,\mathrm{d}x$.
\end{prob}

\begin{prob}[S]{1745}
$\int \cos\frac{x}{2} \cdot \cos\frac{x}{3}\,\mathrm{d}x$.
\end{prob}

\begin{prob}[S]{1746}
$\int \sin\left(2x - \frac{\pi}{6}\right)\cos\left(3x + \frac{\pi}{4}\right)\,\mathrm{d}x$.
\end{prob}

\begin{prob}[S]{1747}
$\int \sin^3 x\,\mathrm{d}x$.
\end{prob}

\begin{prob}[S]{1748}
$\int \cos^3 x\,\mathrm{d}x$.
\end{prob}

\begin{prob}[S]{1749}
$\int \sin^4 x\,\mathrm{d}x$.
\end{prob}

\begin{prob}[S]{1750}
$\int \cos^4 x\,\mathrm{d}x$.
\end{prob}

\begin{prob}[S]{1751}
$\int \cot^2 x\,\mathrm{d}x$.
\end{prob}

\begin{prob}[S]{1752}
$\int \tan^3 x\,\mathrm{d}x$.
\end{prob}

\begin{prob}[S]{1753}
$\int \sin^2 3x \sin^3 2x\,\mathrm{d}x$.
\end{prob}

\begin{prob}[S]{1754}
$\int \frac{\mathrm{d}x}{\sin^2 x \cos^2 x}$.
\probhint{$1 \equiv \sin^2 x + \cos^2 x$.}
\end{prob}

\begin{prob}[S]{1755}
$\int \frac{\mathrm{d}x}{\sin^2 x \cos x}$.
\end{prob}

\begin{prob}[S]{1756}
$\int \frac{\mathrm{d}x}{\sin x \cos^3 x}$.
\end{prob}

\begin{prob}[S]{1757}
$\int \frac{\cos^3 x}{\sin x}\,\mathrm{d}x$.
\end{prob}

\begin{prob}[S]{1758}
$\int \frac{\mathrm{d}x}{\cos^4 x}$.
\end{prob}

\begin{prob}[S]{1759}
$\int \frac{\mathrm{d}x}{1 + \mathrm{e}^x}$.
\end{prob}

\begin{prob}[S]{1760}
$\int \frac{(1 + \mathrm{e}^x)^2}{1 + \mathrm{e}^{2x}}\,\mathrm{d}x$.
\end{prob}

\begin{prob}[S]{1761}
$\int \sinh^2 x\,\mathrm{d}x$.
\end{prob}

\begin{prob}[S]{1762}
$\int \cosh^2 x\,\mathrm{d}x$.
\end{prob}

\begin{prob}[S]{1763}
$\int \sinh x \sinh 2x\,\mathrm{d}x$.
\end{prob}

\begin{prob}[S]{1764}
$\int \cosh x \cdot \cosh 3x\,\mathrm{d}x$.
\end{prob}

\begin{prob}[S]{1765}
$\int \frac{\mathrm{d}x}{\sinh^2 x \cosh^2 x}$.
\end{prob}

\probgroup{用适当的代换，求下列积分：}

\begin{prob}[S]{1766}
$\int x^2 \sqrt[3]{1 - x}\,\mathrm{d}x$.
\end{prob}

\begin{prob}[S]{1767}
$\int x^3 (1 - 5x^2)^{10}\,\mathrm{d}x$.
\end{prob}

\begin{prob}[S]{1768}
$\int \frac{x^2}{\sqrt{2 - x}}\,\mathrm{d}x$.
\end{prob}

\begin{prob}[S]{1769}
$\int \frac{x^5}{\sqrt{1 - x^2}}\,\mathrm{d}x$.
\end{prob}

\begin{prob}[S]{1770}
$\int x^5 (2 - 5x^3)^{\frac{2}{3}}\,\mathrm{d}x$.
\end{prob}

\begin{prob}[S]{1771}
$\int \cos^5 x \cdot \sqrt{\sin x}\,\mathrm{d}x$.
\end{prob}

\begin{prob}[S]{1772}
$\int \frac{\sin x \cos^3 x}{1 + \cos^2 x}\,\mathrm{d}x$.
\end{prob}

\begin{prob}[S]{1773}
$\int \frac{\sin^2 x}{\cos^6 x}\,\mathrm{d}x$.
\end{prob}

\begin{prob}[S]{1774}
$\int \frac{\ln x\,\mathrm{d}x}{x\sqrt{1 + \ln x}}$.
\end{prob}

\begin{prob}[S]{1775}
$\int \frac{\mathrm{d}x}{\mathrm{e}^{\frac{x}{2}} + \mathrm{e}^x}$.
\end{prob}

\begin{prob}[S]{1776}
$\int \frac{\mathrm{d}x}{\sqrt{1 + \mathrm{e}^x}}$.
\end{prob}

\begin{prob}[S]{1777}
$\int \frac{\arctan\sqrt{x}}{\sqrt{x}} \cdot \frac{\mathrm{d}x}{1 + x}$.
\end{prob}

\probgroup{用三角函数代换 $x = a\sin t, x = a\tan t, x = a\sin^2 t$ 等，求下列积分（参数为正的）：}

\begin{prob}[S]{1778}
$\int \frac{\mathrm{d}x}{(1 - x^2)^{\frac{3}{2}}}$.
\end{prob}

\begin{prob}[S]{1779}
$\int \frac{x^2\,\mathrm{d}x}{\sqrt{x^2 - 2}}$.
\end{prob}

\begin{prob}[S]{1780}
$\int \sqrt{1 - x^2}\,\mathrm{d}x$.
\end{prob}

\begin{prob}[S]{1781}
$\int \frac{\mathrm{d}x}{(x^2 + a^2)^{\frac{3}{2}}}$.
\end{prob}

\begin{prob}[S]{1782}
$\int \sqrt{\frac{a + x}{a - x}}\,\mathrm{d}x$.
\end{prob}

\begin{prob}[S]{1783}
$\int x\sqrt{\frac{x}{2a - x}}\,\mathrm{d}x$.
\end{prob}

\begin{prob}[S]{1784}
$\int \frac{\mathrm{d}x}{\sqrt{(x - a)(b - x)}}$.
\probhint{用代换 $x - a = (b - a)\sin^2 t$.}
\end{prob}

\begin{prob}[S]{1785}
$\int \sqrt{(x - a)(b - a)}\,\mathrm{d}x$.
\end{prob}

\probgroup{用双曲函数代换 $x = a\sinh t, x = a\cosh t$ 等，求下列积分（参数为正的）：}

\begin{prob}[S]{1786}
$\int \sqrt{a^2 + x^2}\,\mathrm{d}x$.
\end{prob}

\begin{prob}[S]{1787}
$\int \frac{x^2\,\mathrm{d}x}{\sqrt{a^2 + x^2}}$.
\end{prob}

\begin{prob}[S]{1788}
$\int \sqrt{\frac{x - a}{x + a}}\,\mathrm{d}x$.
\end{prob}

\begin{prob}[S]{1789}
$\int \frac{\mathrm{d}x}{\sqrt{(x + a)(x + b)}}$.
\end{prob}

\begin{prob}[S]{1790}
$\int \sqrt{(x + a)(x + b)}\,\mathrm{d}x$.
\probhint{令 $x + a = (b - a)\sinh^2 t$.}
\end{prob}

\probgroup{用分部积分法，求下列积分：}

\begin{prob}[S]{1791}
$\int \ln x\,\mathrm{d}x$.
\end{prob}

\begin{prob}[S]{1792}
$\int x^n \ln x\,\mathrm{d}x \quad (n \neq -1)$.
\end{prob}

\begin{prob}[S]{1793}
$\int \left(\frac{\ln x}{x}\right)^2\,\mathrm{d}x$.
\end{prob}

\begin{prob}[S]{1794}
$\int \sqrt{x}\ln^2 x\,\mathrm{d}x$.
\end{prob}

\begin{prob}[S]{1795}
$\int x\mathrm{e}^{-x}\,\mathrm{d}x$.
\end{prob}

\begin{prob}[S]{1796}
$\int x^2 \mathrm{e}^{-2x}\,\mathrm{d}x$.
\end{prob}

\begin{prob}[S]{1797}
$\int x^3 \mathrm{e}^{-x^2}\,\mathrm{d}x$.
\end{prob}

\begin{prob}[S]{1798}
$\int x\cos x\,\mathrm{d}x$.
\end{prob}

\begin{prob}[S]{1799}
$\int x^2 \sin 2x\,\mathrm{d}x$.
\end{prob}

\begin{prob}[S]{1800}
$\int x\sinh x\,\mathrm{d}x$.
\end{prob}

\begin{prob}[S]{1801}
$\int x^3 \cosh 3x\,\mathrm{d}x$.
\end{prob}

\begin{prob}[S]{1802}
$\int \arctan x\,\mathrm{d}x$.
\end{prob}

\begin{prob}[S]{1803}
$\int \arcsin x\,\mathrm{d}x$.
\end{prob}

\begin{prob}[S]{1804}
$\int x\arctan x\,\mathrm{d}x$.
\end{prob}

\begin{prob}[S]{1805}
$\int x^2 \arccos x\,\mathrm{d}x$.
\end{prob}

\begin{prob}[S]{1806}
$\int \frac{\arcsin x}{x^2}\,\mathrm{d}x$.
\end{prob}

\begin{prob}[S]{1807}
$\int \ln(x + \sqrt{1 + x^2})\,\mathrm{d}x$.
\end{prob}

\begin{prob}[S]{1808}
$\int x\ln\frac{1 + x}{1 - x}\,\mathrm{d}x$.
\end{prob}

\begin{prob}[S]{1809}
$\int \arctan\sqrt{x}\,\mathrm{d}x$.
\end{prob}

\begin{prob}[S]{1810}
$\int \sin x \cdot \ln(\tan x)\,\mathrm{d}x$.
\end{prob}

\probgroup{求下列积分：}

\begin{prob}[S]{1811}
$\int x^5 \mathrm{e}^{x^3}\,\mathrm{d}x$.
\end{prob}

\begin{prob}[S]{1812}
$\int (\arcsin x)^2\,\mathrm{d}x$.
\end{prob}

\begin{prob}[S]{1813}
$\int x(\arctan x)^2\,\mathrm{d}x$.
\end{prob}

\begin{prob}[S]{1814}
$\int x^2 \ln\frac{1 - x}{1 + x}\,\mathrm{d}x$.
\end{prob}

\begin{prob}[S]{1815}
$\int \frac{x\ln(x + \sqrt{1 + x^2})}{\sqrt{1 + x^2}}\,\mathrm{d}x$.
\end{prob}

\begin{prob}[S]{1816}
$\int \frac{x^2}{(1 + x^2)^2}\,\mathrm{d}x$.
\end{prob}

\begin{prob}[S]{1817}
$\int \frac{\mathrm{d}x}{(a^2 + x^2)^2}$.
\end{prob}

\begin{prob}[S]{1818}
$\int \sqrt{a^2 - x^2}\,\mathrm{d}x$.
\end{prob}

\begin{prob}[S]{1819}
$\int \sqrt{x^2 + a}\,\mathrm{d}x$.
\end{prob}

\begin{prob}[S]{1820}
$\int x^2 \sqrt{a^2 + x^2}\,\mathrm{d}x$.
\end{prob}

\begin{prob}[S]{1821}
$\int x\sin^2 x\,\mathrm{d}x$.
\end{prob}

\begin{prob}[S]{1822}
$\int \mathrm{e}^{\sqrt{x}}\,\mathrm{d}x$.
\end{prob}

\begin{prob}[S]{1823}
$\int x\sin\sqrt{x}\,\mathrm{d}x$.
\end{prob}

\begin{prob}[S]{1824}
$\int \frac{x\mathrm{e}^{\arctan x}}{(1 + x^2)^{\frac{3}{2}}}\,\mathrm{d}x$.
\end{prob}

\begin{prob}[S]{1825}
$\int \frac{\mathrm{e}^{\arctan x}}{(1 + x^2)^{\frac{3}{2}}}\,\mathrm{d}x$.
\end{prob}

\begin{prob}[S]{1826}
$\int \sin(\ln x)\,\mathrm{d}x$.
\end{prob}

\begin{prob}[S]{1827}
$\int \cos(\ln x)\,\mathrm{d}x$.
\end{prob}

\begin{prob}[S]{1828}
$\int \mathrm{e}^{ax} \cos bx\,\mathrm{d}x$.
\end{prob}

\begin{prob}[S]{1829}
$\int \mathrm{e}^{ax} \sin bx\,\mathrm{d}x$.
\end{prob}

\begin{prob}[S]{1830}
$\int \mathrm{e}^{2x} \sin^2 x\,\mathrm{d}x$.
\end{prob}

\begin{prob}[S]{1831}
$\int (\mathrm{e}^x - \cos x)^2\,\mathrm{d}x$.
\end{prob}

\begin{prob}[S]{1832}
$\int \frac{\operatorname{arccot} \mathrm{e}^x}{\mathrm{e}^x}\,\mathrm{d}x$.
\end{prob}

\begin{prob}[S]{1833}
$\int \frac{\ln(\sin x)}{\sin^2 x}\,\mathrm{d}x$.
\end{prob}

\begin{prob}[S]{1834}
$\int \frac{x\,\mathrm{d}x}{\cos^2 x}$.
\end{prob}

\begin{prob}[S]{1835}
$\int \frac{x\mathrm{e}^x}{(x + 1)^2}\,\mathrm{d}x$.
\end{prob}

\probgroup{求下列积分：}

\begin{prob}[S]{1836}
$\int \frac{\mathrm{d}x}{a + bx^2} \quad (ab \neq 0)$.
\end{prob}

\begin{prob}[S]{1837}
$\int \frac{\mathrm{d}x}{x^2 - x + 2}$.
\end{prob}

\begin{prob}[S]{1838}
$\int \frac{\mathrm{d}x}{3x^2 - 2x - 1}$.
\end{prob}

\begin{prob}[S]{1839}
$\int \frac{x\,\mathrm{d}x}{x^4 - 2x^2 - 1}$.
\end{prob}

\begin{prob}[S]{1840}
$\int \frac{x + 1}{x^2 + x + 1}\,\mathrm{d}x$.
\end{prob}

\begin{prob}[S]{1841}
$\int \frac{x\,\mathrm{d}x}{x^2 - 2x\cos\alpha + 1}$.
\end{prob}

\begin{prob}[S]{1842}
$\int \frac{x^3\,\mathrm{d}x}{x^4 - x^2 + 2}$.
\end{prob}

\begin{prob}[S]{1843}
$\int \frac{x^5\,\mathrm{d}x}{x^6 - x^3 - 2}$.
\end{prob}

\begin{prob}[S]{1844}
$\int \frac{\mathrm{d}x}{3\sin^2 x - 8\sin x\cos x + 5\cos^2 x}$.
\end{prob}

\begin{prob}[S]{1845}
$\int \frac{\mathrm{d}x}{\sin x + 2\cos x + 3}$.
\end{prob}

\begin{prob}[S]{1846}
$\int \frac{\mathrm{d}x}{\sqrt{a + bx^2}} \quad (b \neq 0)$.
\end{prob}

\begin{prob}[S]{1847}
$\int \frac{\mathrm{d}x}{\sqrt{1 - 2x - x^2}}$.
\end{prob}

\begin{prob}[S]{1848}
$\int \frac{\mathrm{d}x}{\sqrt{x + x^2}}$.
\end{prob}

\begin{prob}[S]{1849}
$\int \frac{\mathrm{d}x}{\sqrt{2x^2 - x + 2}}$.
\end{prob}

\begin{prob}[M]{1850}
证明：若
\[
y = ax^2 + bx + c \quad (a \neq 0),
\]
则
\begin{align*}
\text{当 } a > 0 \text{ 时，} & \int \frac{\mathrm{d}x}{\sqrt{y}} = \frac{1}{\sqrt{a}} \ln\left|\frac{y'}{2} + \sqrt{ay}\right| + C, \\
\text{当 } a < 0 \text{ 时，} & \int \frac{\mathrm{d}x}{\sqrt{y}} = \frac{1}{\sqrt{-a}} \arcsin \frac{-y'}{\sqrt{b^2 - 4ac}} + C.
\end{align*}
\end{prob}

\begin{prob}[S]{1851}
$\int \frac{x\,\mathrm{d}x}{\sqrt{5 + x - x^2}}$.
\end{prob}

\begin{prob}[S]{1852}
$\int \frac{x + 1}{\sqrt{x^2 + x + 1}}\,\mathrm{d}x$.
\end{prob}

\begin{prob}[M]{1853}
(a) $\int \frac{x\,\mathrm{d}x}{\sqrt{1 - 3x^2 - 2x^4}}$; \qquad (b) $\int \frac{\cos x\,\mathrm{d}x}{\sqrt{1 + \sin x + \cos^2 x}}$.
\end{prob}

\begin{prob}[S]{1854}
$\int \frac{x^3\,\mathrm{d}x}{\sqrt{x^4 - 2x^2 - 1}}$.
\end{prob}

\begin{prob}[S]{1855}
$\int \frac{x + x^3}{\sqrt{1 + x^2 - x^4}}\,\mathrm{d}x$.
\end{prob}

\begin{prob}[S]{1856}
$\int \frac{\mathrm{d}x}{x\sqrt{x^2 + x + 1}}$.
\end{prob}

\begin{prob}[S]{1857}
$\int \frac{\mathrm{d}x}{x^2\sqrt{x^2 + x - 1}}$.
\end{prob}

\begin{prob}[S]{1858}
$\int \frac{\mathrm{d}x}{(x + 1)\sqrt{x^2 + 1}}$.
\end{prob}

\begin{prob}[S]{1859}
$\int \frac{\mathrm{d}x}{(x - 1)\sqrt{x^2 - 2}}$.
\end{prob}

\begin{prob}[S]{1860}
$\int \frac{\mathrm{d}x}{(x + 2)^2 \sqrt{x^2 + 2x - 5}}$.
\end{prob}

\begin{prob}[S]{1861}
$\int \sqrt{2 + x - x^2}\,\mathrm{d}x$.
\end{prob}

\begin{prob}[S]{1862}
$\int \sqrt{2 + x + x^2}\,\mathrm{d}x$.
\end{prob}

\begin{prob}[S]{1863}
$\int \sqrt{x^4 + 2x^2 - 1}\,x\,\mathrm{d}x$.
\end{prob}

\begin{prob}[S]{1864}
$\int \frac{1 - x + x^2}{x\sqrt{1 + x - x^2}}\,\mathrm{d}x$.
\end{prob}

\begin{prob}[S]{1865}
$\int \frac{x^2 + 1}{x\sqrt{x^4 + 1}}\,\mathrm{d}x$.
\end{prob}

\probsec{§2. 有理函数的积分法}

\probgroup{利用待定系数法，求下列积分：}

\begin{prob}[S]{1866}
$\int \frac{2x + 3}{(x - 2)(x + 5)}\,\mathrm{d}x$.
\end{prob}

\begin{prob}[S]{1867}
$\int \frac{x\,\mathrm{d}x}{(x + 1)(x + 2)(x + 3)}$.
\end{prob}

\begin{prob}[S]{1868}
$\int \frac{x^{10}\,\mathrm{d}x}{x^2 + x - 2}$.
\end{prob}

\begin{prob}[S]{1869}
$\int \frac{x^3 + 1}{x^3 - 5x^2 + 6x}\,\mathrm{d}x$.
\end{prob}

\begin{prob}[S]{1870}
$\int \frac{x^4}{x^4 + 5x^2 + 4}\,\mathrm{d}x$.
\end{prob}

\begin{prob}[S]{1871}
$\int \frac{x\,\mathrm{d}x}{x^3 - 3x + 2}$.
\end{prob}

\begin{prob}[S]{1872}
$\int \frac{x^2 + 1}{(x + 1)^2 (x - 1)}\,\mathrm{d}x$.
\end{prob}

\begin{prob}[S]{1873}
$\int \left(\frac{x}{x^2 - 3x + 2}\right)^2\,\mathrm{d}x$.
\end{prob}

\begin{prob}[S]{1874}
$\int \frac{\mathrm{d}x}{(x + 1)(x + 2)^2 (x + 3)^3}$.
\end{prob}

\begin{prob}[S]{1875}
$\int \frac{\mathrm{d}x}{x^5 + x^4 - 2x^3 - 2x^2 + x + 1}$.
\end{prob}

\begin{prob}[S]{1876}
$\int \frac{x^2 + 5x + 4}{x^4 + 5x^2 + 4}\,\mathrm{d}x$.
\end{prob}

\begin{prob}[S]{1877}
$\int \frac{\mathrm{d}x}{(x + 1)(x^2 + 1)}$.
\end{prob}

\begin{prob}[S]{1878}
$\int \frac{\mathrm{d}x}{(x^2 - 4x + 4)(x^2 - 4x + 5)}$.
\end{prob}

\begin{prob}[S]{1879}
$\int \frac{x\,\mathrm{d}x}{(x - 1)^2 (x^2 + 2x + 2)}$.
\end{prob}

\begin{prob}[S]{1880}
$\int \frac{\mathrm{d}x}{x(1 + x)(1 + x + x^2)}$.
\end{prob}

\begin{prob}[S]{1881}
$\int \frac{\mathrm{d}x}{x^3 + 1}$.
\end{prob}

\begin{prob}[S]{1882}
$\int \frac{x\,\mathrm{d}x}{x^3 - 1}$.
\end{prob}

\begin{prob}[S]{1883}
$\int \frac{\mathrm{d}x}{x^4 - 1}$.
\end{prob}

\begin{prob}[S]{1884}
$\int \frac{\mathrm{d}x}{x^4 + 1}$.
\end{prob}

\begin{prob}[S]{1885}
$\int \frac{\mathrm{d}x}{x^4 + x^2 + 1}$.
\end{prob}

\begin{prob}[S]{1886}
$\int \frac{\mathrm{d}x}{x^6 + 1}$.
\end{prob}

\begin{prob}[S]{1887}
$\int \frac{\mathrm{d}x}{(1 + x)(1 + x^2)(1 + x^3)}$.
\end{prob}

\begin{prob}[S]{1888}
$\int \frac{\mathrm{d}x}{x^5 - x^4 + x^3 - x^2 + x - 1}$.
\end{prob}

\begin{prob}[S]{1889}
$\int \frac{x^2\,\mathrm{d}x}{x^4 + 3x^3 + \frac{9}{2}x^2 + 3x + 1}$.
\end{prob}

\begin{prob}[M]{1890}
在什么条件下，积分
\[
\int \frac{ax^2 + bx + c}{x^3(x - 1)^2}\,\mathrm{d}x
\]
为有理函数？
\end{prob}

\probgroup{利用奥斯特罗格拉茨基方法求积分：}

\begin{prob}[S]{1891}
$\displaystyle \int \frac{x\,\mathrm{d}x}{(x - 1)^2(x + 1)^3}$.
\end{prob}

\begin{prob}[S]{1892}
$\displaystyle \int \frac{\mathrm{d}x}{(x^3 + 1)^2}$.
\end{prob}

\begin{prob}[S]{1893}
$\displaystyle \int \frac{\mathrm{d}x}{(x^2 + 1)^3}$.
\end{prob}

\begin{prob}[S]{1894}
$\displaystyle \int \frac{x^2\,\mathrm{d}x}{(x^2 + 2x + 2)^2}$.
\end{prob}

\begin{prob}[S]{1895}
$\displaystyle \int \frac{\mathrm{d}x}{(x^4 + 1)^2}$.
\end{prob}

\begin{prob}[S]{1896}
$\displaystyle \int \frac{x^2 + 3x - 2}{(x - 1)(x^2 + x + 1)^2}\,\mathrm{d}x$.
\end{prob}

\begin{prob}[S]{1897}
$\displaystyle \int \frac{\mathrm{d}x}{(x^4 - 1)^3}$.
\end{prob}

\probgroup{分出下列积分的代数部分：}

\begin{prob}[S]{1898}
$\displaystyle \int \frac{x^2 + 1}{(x^4 + x^2 + 1)^2}\,\mathrm{d}x$.
\end{prob}

\begin{prob}[S]{1899}
$\displaystyle \int \frac{\mathrm{d}x}{(x^3 + x + 1)^3}$.
\end{prob}

\begin{prob}[S]{1900}
$\displaystyle \int \frac{4x^5 - 1}{(x^5 + x + 1)^2}\,\mathrm{d}x$.
\end{prob}

\begin{prob}[M]{1901}
计算积分
\[
\int \frac{\mathrm{d}x}{x^4 + 2x^3 + 3x^2 + 2x + 1}.
\]
\end{prob}

\begin{prob}[M]{1902}
在什么条件下，积分
\[
\int \frac{\alpha x^2 + 2\beta x + \gamma}{(ax^2 + 2bx + c)^2}\,\mathrm{d}x
\]
为有理函数？
\end{prob}

\probgroup{利用不同方法，计算下列积分：}

\begin{prob}[S]{1903}
$\displaystyle \int \frac{x^3}{(x - 1)^{100}}\,\mathrm{d}x$.
\end{prob}

\begin{prob}[S]{1904}
$\displaystyle \int \frac{x\,\mathrm{d}x}{x^8 - 1}$.
\end{prob}

\begin{prob}[S]{1905}
$\displaystyle \int \frac{x^3\,\mathrm{d}x}{x^8 + 3}$.
\end{prob}

\begin{prob}[S]{1906}
$\displaystyle \int \frac{x^2 + x}{x^6 + 1}\,\mathrm{d}x$.
\end{prob}

\begin{prob}[S]{1907}
$\displaystyle \int \frac{x^4 - 3}{x(x^8 + 3x^4 + 2)}\,\mathrm{d}x$.
\end{prob}

\begin{prob}[S]{1908}
$\displaystyle \int \frac{x^4\,\mathrm{d}x}{(x^{10} - 10)^2}$.
\end{prob}

\begin{prob}[S]{1909}
$\displaystyle \int \frac{x^{11}\,\mathrm{d}x}{x^8 + 3x^4 + 2}$.
\end{prob}

\begin{prob}[S]{1910}
$\displaystyle \int \frac{x^9\,\mathrm{d}x}{(x^{10} + 2x^5 + 2)^2}$.
\end{prob}

\begin{prob}[S]{1911}
$\displaystyle \int \frac{x^{2n - 1}}{x^n + 1}\,\mathrm{d}x$.
\end{prob}

\begin{prob}[S]{1912}
$\displaystyle \int \frac{x^{3n - 1}}{(x^{2n} + 1)^2}\,\mathrm{d}x$.
\end{prob}

\begin{prob}[S]{1913}
$\displaystyle \int \frac{\mathrm{d}x}{x(x^{10} + 2)}$.
\end{prob}

\begin{prob}[S]{1914}
$\displaystyle \int \frac{\mathrm{d}x}{x(x^{10} + 1)^2}$.
\end{prob}

\begin{prob}[S]{1915}
$\displaystyle \int \frac{1 - x^7}{x(1 + x^7)}\,\mathrm{d}x$.
\end{prob}

\begin{prob}[S]{1916}
$\displaystyle \int \frac{x^4 - 1}{x(x^4 - 5)(x^5 - 5x + 1)}\,\mathrm{d}x$.
\end{prob}

\begin{prob}[S]{1917}
$\displaystyle \int \frac{x^2 + 1}{x^4 + x^2 + 1}\,\mathrm{d}x$.
\end{prob}

\begin{prob}[S]{1918}
$\displaystyle \int \frac{x^2 - 1}{x^4 + x^3 + x^2 + x + 1}\,\mathrm{d}x$.
\end{prob}

\begin{prob}[S]{1919}
$\displaystyle \int \frac{x^5 - x}{x^8 + 1}\,\mathrm{d}x$.
\end{prob}

\begin{prob}[S]{1920}
$\displaystyle \int \frac{x^4 + 1}{x^6 + 1}\,\mathrm{d}x$.
\end{prob}

\begin{prob}[M]{1921}
试导出用于计算积分
\[
I_n = \int \frac{\mathrm{d}x}{(ax^2 + bx + c)^n} \quad (a \neq 0)
\]
的递推公式。利用这个公式计算：
\[
I_3 = \int \frac{\mathrm{d}x}{(x^2 + x + 1)^3}.
\]
\probhint{利用恒等式 $4a(ax^2 + bx + c) \equiv (2ax + b)^2 + (4ac - b^2)$。}
\end{prob}

\begin{prob}[M]{1922}
利用代换 $t = \dfrac{x + a}{x + b}$ 计算积分：
\[
I = \int \frac{\mathrm{d}x}{(x + a)^m (x + b)^n} \quad (m \text{ 和 } n \text{ 为正整数}).
\]
利用这个代换，求
\[
\int \frac{\mathrm{d}x}{(x - 2)^2(x + 3)^3}.
\]
\end{prob}

\begin{prob}[M]{1923}
若 $P_n(x)$ 为 $x$ 的 $n$ 次多项式，计算
\[
\int \frac{P_n(x)}{(x - a)^{n + 1}}\,\mathrm{d}x.
\]
\probhint{利用泰勒公式。}
\end{prob}

\begin{prob}[M]{1924}
设 $R(x) = R^*(x^2)$，其中 $R^*$ 为有理函数。函数 $R(x)$ 分解为有理公式时有什么特性？
\end{prob}

\begin{prob}[M]{1925}
计算
\[
\int \frac{\mathrm{d}x}{1 + x^{2n}},
\]
式中 $n$ 为正整数。
\end{prob}

\probsec{§3. 无理函数的积分法}

\probgroup{利用化被积函数为有理函数的方法，求下列积分：}

\begin{prob}[S]{1926}
$\displaystyle \int \frac{\mathrm{d}x}{1 + \sqrt{x}}$.
\end{prob}

\begin{prob}[S]{1927}
$\displaystyle \int \frac{\mathrm{d}x}{x(1 + 2\sqrt{x} + \sqrt[3]{x})}$.
\end{prob}

\begin{prob}[S]{1928}
$\displaystyle \int \frac{x\sqrt[3]{2 + x}}{x + \sqrt[3]{2 + x}}\,\mathrm{d}x$.
\end{prob}

\begin{prob}[S]{1929}
$\displaystyle \int \frac{1 - \sqrt{x + 1}}{1 + \sqrt[3]{x + 1}}\,\mathrm{d}x$.
\end{prob}

\begin{prob}[S]{1930}
$\displaystyle \int \frac{\mathrm{d}x}{\sqrt{x}(1 + \sqrt[4]{x})^3}$.
\end{prob}

\begin{prob}[S]{1931}
$\displaystyle \int \frac{\sqrt{x + 1} - \sqrt{x - 1}}{\sqrt{x + 1} + \sqrt{x - 1}}\,\mathrm{d}x$.
\end{prob}

\begin{prob}[S]{1932}
$\displaystyle \int \frac{\mathrm{d}x}{\sqrt[3]{(x + 1)^2(x - 1)^4}}$.
\end{prob}

\begin{prob}[S]{1933}
$\displaystyle \int \frac{x\,\mathrm{d}x}{\sqrt[4]{x^3(a - x)}} \quad (a > 0)$.
\end{prob}

\begin{prob}[S]{1934}
$\displaystyle \int \frac{\mathrm{d}x}{\sqrt[n]{(x - a)^{n + 1}(x - b)^{n - 1}}} \quad (n \text{ 为正整数})$.
\end{prob}

\begin{prob}[M]{1935}
$\displaystyle \int \frac{\mathrm{d}x}{1 + \sqrt{x} + \sqrt{1 + x}}$.
\probhint{令 $x = \left(\dfrac{u^2 - 1}{2u}\right)^2$。}
\end{prob}

\begin{prob}[M]{1936}
考虑积分
\[
\int R\left[x(x - a)^{\frac{p}{n}}, (x - b)^{\frac{q}{n}}\right] \mathrm{d}x,
\]
其中 $R$ 为有理函数，$p, q, n$ 为整数。证明：若 $p + q = kn$，其中 $k$ 为整数，则该积分为初等函数。
\end{prob}

\probgroup{求最简单二次无理式的积分：}

\begin{prob}[S]{1937}
$\displaystyle \int \frac{x^2}{\sqrt{1 + x + x^2}}\,\mathrm{d}x$.
\end{prob}

\begin{prob}[S]{1938}
$\displaystyle \int \frac{\mathrm{d}x}{(x + 1)\sqrt{x^2 + x + 1}}$.
\end{prob}

\begin{prob}[S]{1939}
$\displaystyle \int \frac{\mathrm{d}x}{(1 - x)^2\sqrt{1 - x^2}}$.
\end{prob}

\begin{prob}[S]{1940}
$\displaystyle \int \frac{\sqrt{x^2 + 2x + 2}}{x}\,\mathrm{d}x$.
\end{prob}

\begin{prob}[S]{1941}
$\displaystyle \int \frac{x\,\mathrm{d}x}{(1 + x)\sqrt{1 - x - x^2}}$.
\end{prob}

\begin{prob}[S]{1942}
$\displaystyle \int \frac{1 - x + x^2}{\sqrt{1 + x - x^2}}\,\mathrm{d}x$.
\end{prob}

\probgroup{利用公式
\[
\int \frac{P_n(x)}{y}\,\mathrm{d}x = Q_{n - 1}(x)y + \lambda \int \frac{\mathrm{d}x}{y},
\]
式中 $y = \sqrt{ax^2 + bx + c}$，$P_n(x)$ 为 $n$ 次多项式，$Q_{n - 1}(x)$ 为 $n - 1$ 次多项式，$\lambda$ 为常数，求下列积分：}

\begin{prob}[S]{1943}
$\displaystyle \int \frac{x^3}{\sqrt{1 + 2x - x^2}}\,\mathrm{d}x$.
\end{prob}

\begin{prob}[S]{1944}
$\displaystyle \int \frac{x^{10}\,\mathrm{d}x}{\sqrt{1 + x^2}}$.
\end{prob}

\begin{prob}[S]{1945}
$\displaystyle \int x^4\sqrt{a^2 - x^2}\,\mathrm{d}x$.
\end{prob}

\begin{prob}[S]{1946}
$\displaystyle \int \frac{x^3 - 6x^2 + 11x - 6}{\sqrt{x^2 + 4x + 3}}\,\mathrm{d}x$.
\end{prob}

\begin{prob}[S]{1947}
$\displaystyle \int \frac{\mathrm{d}x}{x^3\sqrt{x^2 + 1}}$.
\end{prob}

\begin{prob}[S]{1948}
$\displaystyle \int \frac{\mathrm{d}x}{x^4\sqrt{x^2 - 1}}$.
\end{prob}

\begin{prob}[S]{1949}
$\displaystyle \int \frac{\mathrm{d}x}{(x - 1)^3\sqrt{x^2 + 3x + 1}}$.
\end{prob}

\begin{prob}[S]{1950}
$\displaystyle \int \frac{\mathrm{d}x}{(x + 1)^5\sqrt{x^2 + 2x}}$.
\end{prob}

\begin{prob}[M]{1951}
在什么条件下，积分
\[
\int \frac{a_1 x^2 + b_1 x + c_1}{\sqrt{ax^2 + bx + c}}\,\mathrm{d}x
\]
是代数函数？
\end{prob}

\probgroup{分解有理函数 $\dfrac{P(x)}{Q(x)}$ 为最简分式，求积分 $\displaystyle \int \frac{P(x)}{Q(x)y}\,\mathrm{d}x$，式中 $y = \sqrt{ax^2 + bx + c}$：}

\begin{prob}[S]{1952}
$\displaystyle \int \frac{x\,\mathrm{d}x}{(x - 1)^2\sqrt{1 + 2x - x^2}}$.
\end{prob}

\begin{prob}[S]{1953}
$\displaystyle \int \frac{x\,\mathrm{d}x}{(x^2 - 1)\sqrt{x^2 - x - 1}}$.
\end{prob}

\begin{prob}[S]{1954}
$\displaystyle \int \frac{\sqrt{x^2 + x + 1}}{(x + 1)^2}\,\mathrm{d}x$.
\end{prob}

\begin{prob}[S]{1955}
$\displaystyle \int \frac{x^3}{(1 + x)\sqrt{1 + 2x - x^2}}\,\mathrm{d}x$.
\end{prob}

\begin{prob}[S]{1956}
$\displaystyle \int \frac{x\,\mathrm{d}x}{(x^2 - 3x + 2)\sqrt{x^2 - 4x + 3}}$.
\end{prob}

\begin{prob}[S]{1957}
$\displaystyle \int \frac{\mathrm{d}x}{(1 + x^2)\sqrt{1 - x^2}}$.
\end{prob}

\begin{prob}[S]{1958}
$\displaystyle \int \frac{\mathrm{d}x}{(x^2 + 1)\sqrt{x^2 - 1}}$.
\end{prob}

\begin{prob}[S]{1959}
$\displaystyle \int \frac{\mathrm{d}x}{(1 - x^4)\sqrt{1 + x^2}}$.
\end{prob}

\begin{prob}[S]{1960}
$\displaystyle \int \frac{\sqrt{x^2 + 2}}{x^2 + 1}\,\mathrm{d}x$.
\end{prob}

\probgroup{化二次三项式为标准形式，计算下列积分：}

\begin{prob}[S]{1961}
$\displaystyle \int \frac{\mathrm{d}x}{(x^2 + x + 1)\sqrt{x^2 + x - 1}}$.
\end{prob}

\begin{prob}[S]{1962}
$\displaystyle \int \frac{x^2\,\mathrm{d}x}{(4 - 2x + x^2)\sqrt{2 + 2x - x^2}}$.
\end{prob}

\begin{prob}[S]{1963}
$\displaystyle \int \frac{(x + 1)\,\mathrm{d}x}{(x^2 + x + 1)\sqrt{x^2 + x + 1}}$.
\end{prob}

\begin{prob}[M]{1964}
利用分式线性代换 $x = \dfrac{\alpha + \beta t}{1 + t}$ 计算积分：
\[
\int \frac{\mathrm{d}x}{(x^2 - x + 1)\sqrt{x^2 + x + 1}}.
\]
\end{prob}

\begin{prob}[M]{1965}
求
\[
\int \frac{\mathrm{d}x}{(x^2 + 2)\sqrt{2x^2 - 2x + 5}}.
\]
\end{prob}

\probgroup{利用欧拉代换：\\
(1) 若 $a > 0$，$\sqrt{ax^2 + bx + c} = \pm\sqrt{a}x + z$；\\
(2) 若 $c > 0$，$\sqrt{ax^2 + bx + c} = xz \pm \sqrt{c}$；\\
(3) $\sqrt{a(x - x_1)(x - x_2)} = z(x - x_1)$，\\
求下列积分：}

\begin{prob}[S]{1966}
$\displaystyle \int \frac{\mathrm{d}x}{x + \sqrt{x^2 + x + 1}}$.
\end{prob}

\begin{prob}[S]{1967}
$\displaystyle \int \frac{\mathrm{d}x}{1 + \sqrt{1 - 2x - x^2}}$.
\end{prob}

\begin{prob}[S]{1968}
$\displaystyle \int x\sqrt{x^2 - 2x + 2}\,\mathrm{d}x$.
\end{prob}

\begin{prob}[S]{1969}
$\displaystyle \int \frac{x - \sqrt{x^2 + 3x + 2}}{x + \sqrt{x^2 + 3x + 2}}\,\mathrm{d}x$.
\end{prob}

\begin{prob}[S]{1970}
$\displaystyle \int \frac{\mathrm{d}x}{\left[1 + \sqrt{x(1 + x)}\right]^2}$.
\end{prob}

\probgroup{利用不同方法，计算下列积分：}

\begin{prob}[S]{1971}
$\displaystyle \int \frac{\mathrm{d}x}{\sqrt{x^2 + 1} - \sqrt{x^2 - 1}}$.
\end{prob}

\begin{prob}[S]{1972}
$\displaystyle \int \frac{x\,\mathrm{d}x}{(1 - x^3)\sqrt{1 - x^2}}$.
\end{prob}

\begin{prob}[S]{1973}
$\displaystyle \int \frac{\mathrm{d}x}{\sqrt{2} + \sqrt{1 - x} + \sqrt{1 + x}}$.
\end{prob}

\begin{prob}[S]{1974}
$\displaystyle \int \frac{x + \sqrt{1 + x + x^2}}{1 + x + \sqrt{1 + x + x^2}}\,\mathrm{d}x$.
\end{prob}

\begin{prob}[S]{1975}
$\displaystyle \int \frac{\sqrt{x(x + 1)}}{\sqrt{x} + \sqrt{x + 1}}\,\mathrm{d}x$.
\end{prob}

\begin{prob}[S]{1976}
$\displaystyle \int \frac{(x^2 - 1)\,\mathrm{d}x}{(x^2 + 1)\sqrt{x^4 + 1}}$.
\end{prob}

\begin{prob}[S]{1977}
$\displaystyle \int \frac{(x^2 + 1)\,\mathrm{d}x}{(x^2 - 1)\sqrt{x^4 + 1}}$.
\end{prob}

\begin{prob}[S]{1978}
$\displaystyle \int \frac{\mathrm{d}x}{x\sqrt{x^4 + 2x^2 - 1}}$.
\end{prob}

\begin{prob}[S]{1979}
$\displaystyle \int \frac{(x^2 + 1)\,\mathrm{d}x}{x\sqrt{x^4 + x^2 + 1}}$.
\end{prob}

\begin{prob}[M]{1980}
证明：积分
\[
\int R(x, \sqrt{ax + b}, \sqrt{cx + d})\,\mathrm{d}x \quad (R \text{ 为有理函数})
\]
的求法归结为有理函数的积分法。
\end{prob}

\probgroup{计算下列积分：}

\begin{prob}[S]{1981}
$\displaystyle \int \sqrt{x^3 + x^4}\,\mathrm{d}x$.
\end{prob}

\begin{prob}[S]{1982}
$\displaystyle \int \frac{\sqrt{x}}{(1 + \sqrt[3]{x})^2}\,\mathrm{d}x$.
\end{prob}

\begin{prob}[S]{1983}
$\displaystyle \int \frac{x\,\mathrm{d}x}{\sqrt{1 + \sqrt[3]{x^2}}}$.
\end{prob}

\begin{prob}[S]{1984}
$\displaystyle \int \frac{x^5\,\mathrm{d}x}{\sqrt{1 - x^2}}$.
\end{prob}

\begin{prob}[S]{1985}
$\displaystyle \int \frac{\mathrm{d}x}{\sqrt[3]{1 + x^3}}$.
\end{prob}

\begin{prob}[S]{1986}
$\displaystyle \int \frac{\mathrm{d}x}{\sqrt[4]{1 + x^4}}$.
\end{prob}

\begin{prob}[S]{1987}
$\displaystyle \int \frac{\mathrm{d}x}{x\sqrt[6]{1 + x^6}}$.
\end{prob}

\begin{prob}[S]{1988}
$\displaystyle \int \frac{\mathrm{d}x}{x^3\sqrt[5]{1 + \frac{1}{x}}}$.
\end{prob}

\begin{prob}[S]{1989}
$\displaystyle \int \sqrt[3]{3x - x^3}\,\mathrm{d}x$.
\end{prob}

\begin{prob}[M]{1990}
在什么情形下，积分
\[
\int \sqrt{1 + x^m}\,\mathrm{d}x \quad (m \text{ 为有理数})
\]
为初等函数？
\end{prob}

\probsec{§4. 三角函数的积分法}

\probgroup{求下列积分：}

\begin{prob}[S]{1991}
$\displaystyle \int \cos^5 x\,\mathrm{d}x$.
\end{prob}

\begin{prob}[S]{1992}
$\displaystyle \int \sin^6 x\,\mathrm{d}x$.
\end{prob}

\begin{prob}[S]{1993}
$\int \cos^6 x \,\mathrm{d}x$.
\end{prob}

\begin{prob}[S]{1994}
$\int \sin^2 x \cos^4 x \,\mathrm{d}x$.
\end{prob}

\begin{prob}[S]{1995}
$\int \sin^4 x \cos^5 x \,\mathrm{d}x$.
\end{prob}

\begin{prob}[S]{1996}
$\int \sin^5 x \cos^5 x \,\mathrm{d}x$.
\end{prob}

\begin{prob}[S]{1997}
$\int \frac{\sin^3 x}{\cos^4 x} \,\mathrm{d}x$.
\end{prob}

\begin{prob}[S]{1998}
$\int \frac{\cos^4 x}{\sin^3 x} \,\mathrm{d}x$.
\end{prob}

\begin{prob}[S]{1999}
$\int \frac{\mathrm{d}x}{\sin^3 x}$.
\end{prob}

\begin{prob}[S]{2000}
$\int \frac{\mathrm{d}x}{\cos^3 x}$.
\end{prob}

\begin{prob}[S]{2001}
$\int \frac{\mathrm{d}x}{\sin^4 x \cos^4 x}$.
\end{prob}

\begin{prob}[S]{2002}
$\int \frac{\mathrm{d}x}{\sin^3 x \cos^5 x}$.
\end{prob}

\begin{prob}[S]{2003}
$\int \frac{\mathrm{d}x}{\sin x \cos^4 x}$.
\end{prob}

\begin{prob}[S]{2004}
$\int \tan^5 x \,\mathrm{d}x$.
\end{prob}

\begin{prob}[S]{2005}
$\int \cot^6 x \,\mathrm{d}x$.
\end{prob}

\begin{prob}[S]{2006}
$\int \frac{\sin^4 x}{\cos^6 x} \,\mathrm{d}x$.
\end{prob}

\begin{prob}[S]{2007}
$\int \frac{\mathrm{d}x}{\sqrt{\sin^3 x \cos^5 x}}$.
\end{prob}

\begin{prob}[S]{2008}
$\int \frac{\mathrm{d}x}{\cos x \sqrt[3]{\sin^2 x}}$.
\end{prob}

\begin{prob}[S]{2009}
$\int \frac{\mathrm{d}x}{\sqrt{\tan x}}$.
\end{prob}

\begin{prob}[S]{2010}
$\int \frac{\mathrm{d}x}{\sqrt[3]{\tan x}}$.
\end{prob}

\begin{prob}[M]{2011}
推出下列积分的递推公式：
\begin{enumerate}[label=(\alph*), itemsep=2pt, topsep=3pt]
\item $I_n = \int \sin^n x \,\mathrm{d}x$;
\item $K_n = \int \cos^n x \,\mathrm{d}x \quad (n > 2)$.
\end{enumerate}
利用这些公式计算
\[
\int \sin^6 x \,\mathrm{d}x \quad \text{和} \quad \int \cos^8 x \,\mathrm{d}x.
\]
\end{prob}

\begin{prob}[M]{2012}
推出下列积分的递推公式：
\begin{enumerate}[label=(\alph*), itemsep=2pt, topsep=3pt]
\item $I_n = \int \frac{\mathrm{d}x}{\sin^n x}$;
\item $K_n = \int \frac{\mathrm{d}x}{\cos^n x} \quad (n > 2)$.
\end{enumerate}
利用这些公式计算
\[
\int \frac{\mathrm{d}x}{\sin^5 x} \quad \text{和} \quad \int \frac{\mathrm{d}x}{\cos^7 x}.
\]
\end{prob}

\probgroup{为了计算下面的积分，可以运用公式：\\[3pt]
I. $\sin\alpha \sin\beta = \frac{1}{2}[\cos(\alpha - \beta) - \cos(\alpha + \beta)]$;\\[3pt]
II. $\cos\alpha \cos\beta = \frac{1}{2}[\cos(\alpha - \beta) + \cos(\alpha + \beta)]$;\\[3pt]
III. $\sin\alpha \cos\beta = \frac{1}{2}[\sin(\alpha - \beta) + \sin(\alpha + \beta)]$.\\[3pt]
求积分：}

\begin{prob}[S]{2013}
$\int \sin 5x \cos x \,\mathrm{d}x$.
\end{prob}

\begin{prob}[S]{2014}
$\int \cos x \cos 2x \cos 3x \,\mathrm{d}x$.
\end{prob}

\begin{prob}[S]{2015}
$\int \sin x \sin \frac{x}{2} \sin \frac{x}{3} \,\mathrm{d}x$.
\end{prob}

\begin{prob}[S]{2016}
$\int \sin x \sin(x + a) \sin(x + b) \,\mathrm{d}x$.
\end{prob}

\begin{prob}[S]{2017}
$\int \cos^2 ax \cos^2 bx \,\mathrm{d}x$.
\end{prob}

\begin{prob}[S]{2018}
$\int \sin^3 2x \cdot \cos^2 3x \,\mathrm{d}x$.
\end{prob}

\probgroup{为了计算下面的积分，可以运用恒等式：\\[3pt]
$\sin(\alpha - \beta) \equiv \sin[(x + \alpha) - (x + \beta)]$,\\[3pt]
$\cos(\alpha - \beta) \equiv \cos[(x + \alpha) - (x + \beta)]$.\\[3pt]
求积分：}

\begin{prob}[S]{2019}
$\int \frac{\mathrm{d}x}{\sin(x + a)\sin(x + b)}$.
\end{prob}

\begin{prob}[S]{2020}
$\int \frac{\mathrm{d}x}{\sin(x + a)\cos(x + b)}$.
\end{prob}

\begin{prob}[S]{2021}
$\int \frac{\mathrm{d}x}{\cos(x + a)\cos(x + b)}$.
\end{prob}

\begin{prob}[S]{2022}
$\int \frac{\mathrm{d}x}{\sin x - \sin a}$.
\end{prob}

\begin{prob}[S]{2023}
$\int \frac{\mathrm{d}x}{\cos x + \cos a}$.
\end{prob}

\begin{prob}[S]{2024}
$\int \tan x \tan(x + a) \,\mathrm{d}x$.
\end{prob}

\probgroup{形如 $\int R(\sin x, \cos x)\,\mathrm{d}x$ ($R$ 为有理函数) 的积分，在一般情形下可利用代换 $\tan\frac{x}{2} = t$ 化为有理函数的积分。\\[3pt]
(a) 若等式 $R(-\sin x, \cos x) \equiv -R(\sin x, \cos x)$ 或 $R(\sin x, -\cos x) \equiv -R(\sin x, \cos x)$ 成立，则最好利用相应的代换 $\cos x = t$ 或 $\sin x = t$.\\[3pt]
(b) 若等式 $R(-\sin x, -\cos x) \equiv R(\sin x, \cos x)$ 成立，则最好利用代换 $\tan x = t$.\\[3pt]
求积分：}

\begin{prob}[S]{2025}
$\int \frac{\mathrm{d}x}{2\sin x - \cos x + 5}$.
\end{prob}

\begin{prob}[S]{2026}
$\int \frac{\mathrm{d}x}{(2 + \cos x)\sin x}$.
\end{prob}

\begin{prob}[S]{2027}
$\int \frac{\sin^2 x}{\sin x + 2\cos x}\,\mathrm{d}x$.
\end{prob}

\begin{prob}[M]{2028}
$\int \frac{\mathrm{d}x}{1 + \varepsilon \cos x}$, \quad (a) $0 < \varepsilon < 1$; \qquad (b) $\varepsilon > 1$.
\end{prob}

\begin{prob}[S]{2029}
$\int \frac{\sin^2 x}{1 + \sin^2 x}\,\mathrm{d}x$.
\end{prob}

\begin{prob}[S]{2030}
$\int \frac{\mathrm{d}x}{a^2 \sin^2 x + b^2 \cos^2 x}$.
\end{prob}

\begin{prob}[S]{2031}
$\int \frac{\cos^2 x \,\mathrm{d}x}{(a^2 \sin^2 x + b^2 \cos^2 x)^2}$.
\end{prob}

\begin{prob}[S]{2032}
$\int \frac{\sin x \cos x}{\sin x + \cos x}\,\mathrm{d}x$.
\end{prob}

\begin{prob}[S]{2033}
$\int \frac{\mathrm{d}x}{(a\sin x + b\cos x)^2}$.
\end{prob}

\begin{prob}[S]{2034}
$\int \frac{\sin x \,\mathrm{d}x}{\sin^3 x + \cos^3 x}$.
\end{prob}

\begin{prob}[S]{2035}
$\int \frac{\mathrm{d}x}{\sin^4 x + \cos^4 x}$.
\end{prob}

\begin{prob}[S]{2036}
$\int \frac{\sin^2 x \cos^2 x}{\sin^8 x + \cos^8 x}\,\mathrm{d}x$.
\end{prob}

\begin{prob}[S]{2037}
$\int \frac{\sin^2 x - \cos^2 x}{\sin^4 x + \cos^4 x}\,\mathrm{d}x$.
\end{prob}

\begin{prob}[S]{2038}
$\int \frac{\sin x \cos x}{1 + \sin^4 x}\,\mathrm{d}x$.
\end{prob}

\begin{prob}[S]{2039}
$\int \frac{\mathrm{d}x}{\sin^6 x + \cos^6 x}$.
\end{prob}

\begin{prob}[S]{2040}
$\int \frac{\mathrm{d}x}{(\sin^2 x + 2\cos^2 x)^2}$.
\end{prob}

\begin{prob}[S]{2041}
把分母化为对数的形式，求积分
\[
\int \frac{\mathrm{d}x}{a\sin x + b\cos x}.
\]
\end{prob}

\begin{prob}[M]{2042}
证明：
\[
\int \frac{a_1 \sin x + b_1 \cos x}{a\sin x + b\cos x}\,\mathrm{d}x = Ax + B\ln |a\sin x + b\cos x| + C,
\]
式中 $A, B, C$ 为常数。
\probhint{令 $a_1 \sin x + b_1 \cos x = A(a\sin x + b\cos x) + B(a\cos x - b\sin x)$，式中 $A$ 和 $B$ 为常数。}
\end{prob}

\probgroup{求积分：}

\begin{prob}[M]{2043}
(a) $\int \frac{\sin x - \cos x}{\sin x + 2\cos x}\,\mathrm{d}x$; \qquad (b) $\int \frac{\sin x}{\sin x - 3\cos x}\,\mathrm{d}x$.
\end{prob}

\begin{prob}[S]{2044}
$\int \frac{\mathrm{d}x}{3 + 5\tan x}$.
\end{prob}

\begin{prob}[S]{2045}
$\int \frac{a_1 \sin x + b_1 \cos x}{(a\sin x + b\cos x)^2}\,\mathrm{d}x$.
\end{prob}

\begin{prob}[M]{2046}
证明：
\[
\int \frac{a_1 \sin x + b_1 \cos x + c_1}{a\sin x + b\cos x + c}\,\mathrm{d}x = Ax + B\ln |a\sin x + b\cos x + c| + C\int \frac{\mathrm{d}x}{a\sin x + b\cos x + c},
\]
式中 $A, B, C$ 都是常系数。
\end{prob}

\probgroup{求积分：}

\begin{prob}[S]{2047}
$\int \frac{\sin x + 2\cos x - 3}{\sin x - 2\cos x + 3}\,\mathrm{d}x$.
\end{prob}

\begin{prob}[S]{2048}
$\int \frac{\sin x}{\sqrt{2} + \sin x + \cos x}\,\mathrm{d}x$.
\end{prob}

\begin{prob}[S]{2049}
$\int \frac{2\sin x + \cos x}{3\sin x + 4\cos x - 2}\,\mathrm{d}x$.
\end{prob}

\begin{prob}[M]{2050}
证明：
\[
\int \frac{a_1 \sin^2 x + 2b_1 \sin x \cos x + c_1 \cos^2 x}{a\sin x + b\cos x}\,\mathrm{d}x = A\sin x + B\cos x + C\int \frac{\mathrm{d}x}{a\sin x + b\cos x},
\]
式中 $A, B, C$ 都是常系数。
\end{prob}

\probgroup{求积分：}

\begin{prob}[S]{2051}
$\int \frac{\sin^2 x - 4\sin x \cos x + 3\cos^2 x}{\sin x + \cos x}\,\mathrm{d}x$.
\end{prob}

\begin{prob}[S]{2052}
$\int \frac{\sin^2 x - \sin x \cos x + 2\cos^2 x}{\sin x + 2\cos x}\,\mathrm{d}x$.
\end{prob}

\begin{prob}[M]{2053}
证明：若 $(a - c)^2 + b^2 \neq 0$，则
\[
\int \frac{a_1 \sin x + b_1 \cos x}{a\sin^2 x + 2b\sin x \cos x + c\cos^2 x}\,\mathrm{d}x = A\int \frac{\mathrm{d}u_1}{k_1 u_1^2 + \lambda_1} + B\int \frac{\mathrm{d}u_2}{k_2 u_2^2 + \lambda_2},
\]
式中 $A, B$ 为待定系数，$\lambda_1, \lambda_2$ 为方程
\[
\begin{vmatrix}
a - \lambda & b \\
b & c - \lambda
\end{vmatrix} = 0 \quad (\lambda_1 \neq \lambda_2)
\]
的根，而
\[
u_i = (a - \lambda_i)\sin x + b\cos x, \quad k_i = \frac{1}{a - \lambda_i} \quad (i = 1, 2).
\]
\end{prob}

\probgroup{求积分：}

\begin{prob}[S]{2054}
$\int \frac{2\sin x - \cos x}{3\sin^2 x + 4\cos^2 x}\,\mathrm{d}x$.
\end{prob}

\begin{prob}[S]{2055}
$\int \frac{(\sin x + \cos x)\,\mathrm{d}x}{2\sin^2 x - 4\sin x \cos x + 5\cos^2 x}$.
\end{prob}

\begin{prob}[S]{2056}
$\int \frac{\sin x - 2\cos x}{1 + 4\sin x \cos x}\,\mathrm{d}x$.
\end{prob}

\begin{prob}[M]{2057}
证明：
\[
\int \frac{\mathrm{d}x}{(a\sin x + b\cos x)^n} = \frac{A\sin x + B\cos x}{(a\sin x + b\cos x)^{n-1}} + C\int \frac{\mathrm{d}x}{(a\sin x + b\cos x)^{n-2}},
\]
式中 $A, B, C$ 为待定系数。
\end{prob}

\begin{prob}[S]{2058}
求 $\int \frac{\mathrm{d}x}{(\sin x + 2\cos x)^3}$.
\end{prob}

\begin{prob}[M]{2059}
若 $n$ 为大于 $1$ 的正整数，证明：
\[
\int \frac{\mathrm{d}x}{(a + b\cos x)^n} = \frac{A\sin x}{(a + b\cos x)^{n-1}} + B\int \frac{\mathrm{d}x}{(a + b\cos x)^{n-1}} + C\int \frac{\mathrm{d}x}{(a + b\cos x)^{n-2}},
\]
其中 $|a| \neq |b|$，并求出系数 $A, B$ 和 $C$。
\end{prob}

\probgroup{求积分：}

\begin{prob}[S]{2060}
$\int \frac{\sin x \,\mathrm{d}x}{\cos x \sqrt{1 + \sin^2 x}}$.
\end{prob}

\begin{prob}[S]{2061}
$\int \frac{\sin^2 x}{\cos^2 x \sqrt{\tan x}}\,\mathrm{d}x$.
\end{prob}

\begin{prob}[S]{2062}
$\int \frac{\sin x \,\mathrm{d}x}{\sqrt{2 + \sin 2x}}$.
\end{prob}

\begin{prob}[S]{2063}
$\int \frac{\mathrm{d}x}{(1 + \varepsilon \cos x)^2} \quad (0 < \varepsilon < 1)$.
\end{prob}

\begin{prob}[M]{2064}
$\int \frac{\cos^{n-1} \frac{x + a}{2}}{\sin^{n+1} \frac{x - a}{2}}\,\mathrm{d}x$.
\probhint{令 $t = \frac{\cos \frac{x + a}{2}}{\sin \frac{x - a}{2}}$。}
\end{prob}

\begin{prob}[M]{2065}
推出积分
\[
I_n = \int \left(\frac{\sin \frac{x - a}{2}}{\sin \frac{x + a}{2}}\right)^n \mathrm{d}x \quad (n \text{ 为正整数})
\]
的递推公式。
\end{prob}

\probsec{§5. 各种超越函数的积分法}

\begin{prob}[M]{2066}
证明：若 $P(x)$ 为 $n$ 次多项式，则
\[
\int P(x)\mathrm{e}^{ax}\,\mathrm{d}x = \mathrm{e}^{ax}\left[\frac{P(x)}{a} - \frac{P'(x)}{a^2} + \cdots + (-1)^n \frac{P^{(n)}(x)}{a^{n+1}}\right] + C.
\]
\end{prob}

\begin{prob}[M]{2067}
证明：若 $P(x)$ 为 $n$ 次多项式，则
\begin{align*}
\int P(x)\cos ax\,\mathrm{d}x &= \frac{\sin ax}{a}\left[P(x) - \frac{P''(x)}{a^2} + \frac{P^{(4)}(x)}{a^4} - \cdots\right] \\
&\quad + \frac{\cos ax}{a^2}\left[P'(x) - \frac{P'''(x)}{a^2} + \frac{P^{(5)}(x)}{a^4} - \cdots\right] + C, \\
\int P(x)\sin ax\,\mathrm{d}x &= -\frac{\cos ax}{a}\left[P(x) - \frac{P''(x)}{a^2} + \frac{P^{(4)}(x)}{a^4} - \cdots\right] \\
&\quad + \frac{\sin ax}{a^2}\left[P'(x) - \frac{P'''(x)}{a^2} + \frac{P^{(5)}(x)}{a^4} - \cdots\right] + C.
\end{align*}
\end{prob}

\probgroup{求积分：}

\begin{prob}[S]{2068}
$\int x^3 \mathrm{e}^{3x}\,\mathrm{d}x$.
\end{prob}

\begin{prob}[S]{2069}
$\int (x^2 - 2x + 2)\mathrm{e}^{-x}\,\mathrm{d}x$.
\end{prob}

\begin{prob}[S]{2070}
$\int x^5 \sin 5x \,\mathrm{d}x$.
\end{prob}

\begin{prob}[S]{2071}
$\int (1 + x^2)^2 \cos x \,\mathrm{d}x$.
\end{prob}

\begin{prob}[S]{2072}
$\int x^7 \mathrm{e}^{-x^2}\,\mathrm{d}x$.
\end{prob}

\begin{prob}[S]{2073}
$\int x^2 \mathrm{e}^{\sqrt{x}}\,\mathrm{d}x$.
\end{prob}

\begin{prob}[S]{2074}
$\int \mathrm{e}^{ax} \cos^2 bx \,\mathrm{d}x$.
\end{prob}

\begin{prob}[S]{2075}
$\int \mathrm{e}^{ax} \sin^3 bx \,\mathrm{d}x$.
\end{prob}

\begin{prob}[S]{2076}
$\int x \mathrm{e}^x \sin x \,\mathrm{d}x$.
\end{prob}

\begin{prob}[S]{2077}
$\int x^2 \mathrm{e}^x \cos x \,\mathrm{d}x$.
\end{prob}

\begin{prob}[S]{2078}
$\int x \mathrm{e}^x \sin^2 x \,\mathrm{d}x$.
\end{prob}

\begin{prob}[S]{2079}
$\int (x - \sin x)^3 \,\mathrm{d}x$.
\end{prob}

\begin{prob}[S]{2080}
$\int \cos^2 \sqrt{x}\,\mathrm{d}x$.
\end{prob}

\begin{prob}[M]{2081}
证明：若 $R$ 为有理函数，数 $a_1, a_2, \cdots, a_n$ 为可公约的，则积分
\[
\int R(\mathrm{e}^{a_1 x}, \mathrm{e}^{a_2 x}, \cdots, \mathrm{e}^{a_n x})\,\mathrm{d}x
\]
是初等函数。
\end{prob}

\probgroup{求下列积分：}

\begin{prob}[S]{2082}
$\int \frac{\mathrm{d}x}{(1 + \mathrm{e}^x)^2}$.
\end{prob}

\begin{prob}[S]{2083}
$\int \frac{\mathrm{e}^{2x}}{1 + \mathrm{e}^x}\,\mathrm{d}x$.
\end{prob}

\begin{prob}[S]{2084}
$\int \frac{\mathrm{d}x}{\mathrm{e}^{2x} + \mathrm{e}^x - 2}$.
\end{prob}

\begin{prob}[S]{2085}
$\int \frac{\mathrm{d}x}{1 + \mathrm{e}^{\frac{x}{2}} + \mathrm{e}^{\frac{x}{3}} + \mathrm{e}^{\frac{x}{6}}}$.
\end{prob}

\begin{prob}[S]{2086}
$\int \frac{1 + \mathrm{e}^{\frac{x}{2}}}{\left(1 + \mathrm{e}^{\frac{x}{4}}\right)^2}\,\mathrm{d}x$.
\end{prob}

\begin{prob}[S]{2087}
$\int \frac{\mathrm{d}x}{\sqrt{\mathrm{e}^x - 1}}$.
\end{prob}

\begin{prob}[S]{2088}
$\int \sqrt{\frac{\mathrm{e}^x - 1}{\mathrm{e}^x + 1}} \,\mathrm{d}x$.
\end{prob}

\begin{prob}[S]{2089}
$\int \sqrt{\mathrm{e}^{2x} + 4\mathrm{e}^x - 1} \,\mathrm{d}x$.
\end{prob}

\begin{prob}[S]{2090}
$\int \frac{\mathrm{d}x}{\sqrt{1 + \mathrm{e}^x} + \sqrt{1 - \mathrm{e}^x}}$.
\end{prob}

\begin{prob}[M]{2091}
证明: 若 $R$ 为有理函数, 其分母仅有实根, 则积分
\[
\int R(x)\mathrm{e}^{ax} \,\mathrm{d}x
\]
可用初等函数和超越函数
\[
\int \frac{\mathrm{e}^{ax}}{x} \,\mathrm{d}x = \operatorname{li}(\mathrm{e}^{ax}) + C, \quad \text{式中 } \operatorname{li} x = \int \frac{\mathrm{d}x}{\ln x}
\]
来表示.
\end{prob}

\begin{prob}[M]{2092}
若 $P\left(\frac{1}{x}\right) = a_0 + \frac{a_1}{x} + \cdots + \frac{a_n}{x^n}$, $a_0, a_1, \cdots, a_n$ 为常数, 则在什么情形下, 积分
\[
\int P\left(\frac{1}{x}\right) \mathrm{e}^x \,\mathrm{d}x
\]
为初等函数?
\end{prob}

\probgroup{求积分：}

\begin{prob}[S]{2093}
$\int \left(1 - \frac{2}{x}\right)^2 \mathrm{e}^x \,\mathrm{d}x$.
\end{prob}

\begin{prob}[S]{2094}
$\int \left(1 - \frac{1}{x}\right) \mathrm{e}^{-x} \,\mathrm{d}x$.
\end{prob}

\begin{prob}[S]{2095}
$\int \frac{\mathrm{e}^{2x}}{x^2 - 3x + 2} \,\mathrm{d}x$.
\end{prob}

\begin{prob}[S]{2096}
$\int \frac{x\mathrm{e}^x}{(x + 1)^2} \,\mathrm{d}x$.
\end{prob}

\begin{prob}[S]{2097}
$\int \frac{x^4 \mathrm{e}^{2x}}{(x - 2)^2} \,\mathrm{d}x$.
\end{prob}

\probgroup{求含有 $\ln f(x)$, $\arctan f(x)$, $\arcsin f(x)$, $\arccos f(x)$ 等函数的积分, 其中 $f(x)$ 为代数函数：}

\begin{prob}[S]{2098}
$\int \ln^n x \,\mathrm{d}x \ (n \text{ 为正整数})$.
\end{prob}

\begin{prob}[S]{2099}
$\int x^3 \ln^3 x \,\mathrm{d}x$.
\end{prob}

\begin{prob}[S]{2100}
$\int \left(\frac{\ln x}{x}\right)^3 \,\mathrm{d}x$.
\end{prob}

\begin{prob}[S]{2101}
$\int \ln \left[(x + a)^{x + a}(x + b)^{x + b}\right] \frac{\mathrm{d}x}{(x + a)(x + b)}$.
\end{prob}

\begin{prob}[S]{2102}
$\int \ln^2(x + \sqrt{1 + x^2}) \,\mathrm{d}x$.
\end{prob}

\begin{prob}[S]{2103}
$\int \ln(\sqrt{1 - x} + \sqrt{1 + x}) \,\mathrm{d}x$.
\end{prob}

\begin{prob}[S]{2104}
$\int \frac{\ln x}{(1 + x^2)^{\frac{3}{2}}} \,\mathrm{d}x$.
\end{prob}

\begin{prob}[S]{2105}
$\int x \arctan(x + 1) \,\mathrm{d}x$.
\end{prob}

\begin{prob}[S]{2106}
$\int \sqrt{x} \arctan \sqrt{x} \,\mathrm{d}x$.
\end{prob}

\begin{prob}[S]{2107}
$\int x \arcsin(1 - x) \,\mathrm{d}x$.
\end{prob}

\begin{prob}[S]{2108}
$\int \arcsin \sqrt{x} \,\mathrm{d}x$.
\end{prob}

\begin{prob}[S]{2109}
$\int x \arccos \frac{1}{x} \,\mathrm{d}x$.
\end{prob}

\begin{prob}[S]{2110}
$\int \arcsin \frac{2\sqrt{x}}{1 + x} \,\mathrm{d}x$.
\end{prob}

\begin{prob}[S]{2111}
$\int \frac{\arccos x}{(1 - x^2)^{\frac{3}{2}}} \,\mathrm{d}x$.
\end{prob}

\begin{prob}[S]{2112}
$\int \frac{x \arccos x}{(1 - x^2)^{\frac{3}{2}}} \,\mathrm{d}x$.
\end{prob}

\begin{prob}[S]{2113}
$\int x \arctan x \ln(1 + x^2) \,\mathrm{d}x$.
\end{prob}

\begin{prob}[S]{2114}
$\int x \ln \frac{1 + x}{1 - x} \,\mathrm{d}x$.
\end{prob}

\begin{prob}[S]{2115}
$\int \frac{\ln(x + \sqrt{1 + x^2}) \,\mathrm{d}x}{(1 + x^2)^{\frac{3}{2}}}$.
\end{prob}

\probgroup{求含有双曲函数的积分：}

\begin{prob}[S]{2116}
$\int \sinh^2 x \cosh^2 x \,\mathrm{d}x$.
\end{prob}

\begin{prob}[S]{2117}
$\int \cosh^4 x \,\mathrm{d}x$.
\end{prob}

\begin{prob}[S]{2118}
$\int \sinh^3 x \,\mathrm{d}x$.
\end{prob}

\begin{prob}[S]{2119}
$\int \sinh x \sinh 2x \sinh 3x \,\mathrm{d}x$.
\end{prob}

\begin{prob}[S]{2120}
$\int \tanh x \,\mathrm{d}x$.
\end{prob}

\begin{prob}[S]{2121}
$\int \coth^2 x \,\mathrm{d}x$.
\end{prob}

\begin{prob}[S]{2122}
$\int \sqrt{\tanh x} \,\mathrm{d}x$.
\end{prob}

\begin{prob}[L]{2123}
\begin{enumerate}[label=(\alph*), itemsep=2pt, topsep=3pt]
  \item $\int \frac{\mathrm{d}x}{\sinh x + 2\cosh x}$;
  \item $\int \frac{\mathrm{d}x}{\sinh^2 x - 4\sinh x \cosh x + 9\cosh^2 x}$;
  \item $\int \frac{\mathrm{d}x}{0.1 + \cosh x}$;
  \item $\int \frac{\cosh x \,\mathrm{d}x}{3\sinh x - 4\cosh x}$.
\end{enumerate}
\end{prob}

\begin{prob}[S]{2124}
$\int \sinh ax \sin bx \,\mathrm{d}x$.
\end{prob}

\begin{prob}[S]{2125}
$\int \sinh ax \cos bx \,\mathrm{d}x$.
\end{prob}

\probsec{§6. 求函数积分的各种例子}

\probgroup{求积分：}

\begin{prob}[S]{2126}
$\int \frac{\mathrm{d}x}{x^6(1 + x^2)}$.
\end{prob}

\begin{prob}[S]{2127}
$\int \frac{x^2 \,\mathrm{d}x}{(1 - x^2)^3}$.
\end{prob}

\begin{prob}[S]{2128}
$\int \frac{\mathrm{d}x}{1 + x^4 + x^8}$.
\end{prob}

\begin{prob}[S]{2129}
$\int \frac{\mathrm{d}x}{\sqrt{x} + \sqrt[3]{x}}$.
\end{prob}

\begin{prob}[S]{2130}
$\int x^2 \sqrt{\frac{x}{1 - x}} \,\mathrm{d}x$.
\end{prob}

\begin{prob}[S]{2131}
$\int \frac{x + 2}{x^2 \sqrt{1 - x^2}} \,\mathrm{d}x$.
\end{prob}

\begin{prob}[S]{2132}
$\int \sqrt{\frac{x}{1 - x\sqrt{x}}} \,\mathrm{d}x$.
\end{prob}

\begin{prob}[S]{2133}
$\int \frac{x^5 \,\mathrm{d}x}{\sqrt{1 + x^2}}$.
\end{prob}

\begin{prob}[S]{2134}
$\int \frac{\mathrm{d}x}{\sqrt[3]{x^2(1 - x)}}$.
\end{prob}

\begin{prob}[S]{2135}
$\int \frac{\mathrm{d}x}{x\sqrt{1 + x^3 + x^6}}$.
\end{prob}

\begin{prob}[S]{2136}
$\int \frac{\mathrm{d}x}{x\sqrt{x^4 - 2x^2 - 1}}$.
\end{prob}

\begin{prob}[S]{2137}
$\int \frac{1 + \sqrt{1 - x^2}}{1 - \sqrt{1 - x^2}} \,\mathrm{d}x$.
\end{prob}

\begin{prob}[S]{2138}
$\int \frac{(1 + x) \,\mathrm{d}x}{x + \sqrt{x + x^2}}$.
\end{prob}

\begin{prob}[S]{2139}
$\int \frac{\ln(1 + x + x^2)}{(1 + x)^2} \,\mathrm{d}x$.
\end{prob}

\begin{prob}[S]{2140}
$\int (2x + 3)\arccos(2x - 3) \,\mathrm{d}x$.
\end{prob}

\begin{prob}[S]{2141}
$\int x \ln(4 + x^4) \,\mathrm{d}x$.
\end{prob}

\begin{prob}[S]{2142}
$\int \frac{\arcsin x}{x^2} \cdot \frac{1 + x^2}{\sqrt{1 - x^2}} \,\mathrm{d}x$.
\end{prob}

\begin{prob}[S]{2143}
$\int \frac{x \ln(1 + \sqrt{1 + x^2})}{\sqrt{1 + x^2}} \,\mathrm{d}x$.
\end{prob}

\begin{prob}[S]{2144}
$\int x\sqrt{x^2 + 1} \ln\sqrt{x^2 - 1} \,\mathrm{d}x$.
\end{prob}

\begin{prob}[S]{2145}
$\int \frac{x}{\sqrt{1 - x^2}} \ln \frac{x}{\sqrt{1 - x}} \,\mathrm{d}x$.
\end{prob}

\begin{prob}[S]{2146}
$\int \frac{\mathrm{d}x}{(2 + \sin x)^2}$.
\end{prob}

\begin{prob}[S]{2147}
$\int \frac{\sin 4x}{\sin^8 x + \cos^8 x} \,\mathrm{d}x$.
\end{prob}

\begin{prob}[S]{2148}
$\int \frac{\mathrm{d}x}{\sin x \sqrt{1 + \cos x}}$.
\end{prob}

\begin{prob}[S]{2149}
$\int \frac{ax^2 + b}{x^2 + 1} \arctan x \,\mathrm{d}x$.
\end{prob}

\begin{prob}[S]{2150}
$\int \frac{ax^2 + b}{x^2 - 1} \ln \left|\frac{x - 1}{x + 1}\right| \,\mathrm{d}x$.
\end{prob}

\begin{prob}[S]{2151}
$\int \frac{x \ln x}{(1 + x^2)^2} \,\mathrm{d}x$.
\end{prob}

\begin{prob}[S]{2152}
$\int \frac{x \arctan x}{\sqrt{1 + x^2}} \,\mathrm{d}x$.
\end{prob}

\begin{prob}[S]{2153}
$\int \frac{\sin 2x}{\sqrt{1 + \cos^4 x}} \,\mathrm{d}x$.
\end{prob}

\begin{prob}[S]{2154}
$\int \frac{x^3 \arccos x}{\sqrt{1 - x^2}} \,\mathrm{d}x$.
\end{prob}

\begin{prob}[S]{2155}
$\int \frac{x^4 \arctan x}{1 + x^2} \,\mathrm{d}x$.
\end{prob}

\begin{prob}[S]{2156}
$\int \frac{x \arctan x}{(1 + x^2)^2} \,\mathrm{d}x$.
\end{prob}

\begin{prob}[S]{2157}
$\int \frac{x \ln(x + \sqrt{1 + x^2})}{(1 - x^2)^2} \,\mathrm{d}x$.
\end{prob}

\begin{prob}[S]{2158}
$\int \sqrt{1 - x^2} \arcsin x \,\mathrm{d}x$.
\end{prob}

\begin{prob}[S]{2159}
$\int x(1 + x^2)\operatorname{arccot} x \,\mathrm{d}x$.
\end{prob}

\begin{prob}[S]{2160}
$\int x^x(1 + \ln x) \,\mathrm{d}x$.
\end{prob}

\begin{prob}[S]{2161}
$\int \frac{\arcsin \mathrm{e}^x}{\mathrm{e}^x} \,\mathrm{d}x$.
\end{prob}

\begin{prob}[S]{2162}
$\int \frac{\arctan \mathrm{e}^{\frac{x}{2}}}{\mathrm{e}^{\frac{x}{2}}(1 + \mathrm{e}^x)} \,\mathrm{d}x$.
\end{prob}

\begin{prob}[S]{2163}
$\int \frac{\mathrm{d}x}{(\mathrm{e}^{x + 1} + 1)^2 - (\mathrm{e}^{x - 1} + 1)^2}$.
\end{prob}

\begin{prob}[S]{2164}
$\int \sqrt{\tanh^2 x + 1} \,\mathrm{d}x$.
\end{prob}

\begin{prob}[S]{2165}
$\int \frac{1 + \sin x}{1 + \cos x} \mathrm{e}^x \,\mathrm{d}x$.
\end{prob}

\begin{prob}[S]{2166}
$\int |x| \,\mathrm{d}x$.
\end{prob}

\begin{prob}[S]{2167}
$\int x|x| \,\mathrm{d}x$.
\end{prob}

\begin{prob}[S]{2168}
$\int (x + |x|)^2 \,\mathrm{d}x$.
\end{prob}

\begin{prob}[S]{2169}
$\int \{|1 + x| - |1 - x|\} \,\mathrm{d}x$.
\end{prob}

\begin{prob}[S]{2170}
$\int \mathrm{e}^{-|x|} \,\mathrm{d}x$.
\end{prob}

\begin{prob}[S]{2171}
$\int \max(1, x^2) \,\mathrm{d}x$.
\end{prob}

\begin{prob}[M]{2172}
$\int \varphi(x) \,\mathrm{d}x$, 其中 $\varphi(x)$ 为数 $x$ 到最近整数的距离.
\end{prob}

\begin{prob}[M]{2173}
$\int [x]|\sin \pi x| \,\mathrm{d}x \ (x \geqslant 0)$.
\end{prob}

\begin{prob}[M]{2174}
$\int f(x) \,\mathrm{d}x$, 其中
$f(x) = \begin{cases} 1 - x^2, & |x| \leqslant 1, \\ 1 - |x|, & |x| > 1. \end{cases}$
\end{prob}

\begin{prob}[M]{2175}
$\int f(x) \,\mathrm{d}x$, 式中
$f(x) = \begin{cases} 1, & -\infty < x < 0, \\ x + 1, & 0 \leqslant x \leqslant 1, \\ 2x, & 1 < x < +\infty. \end{cases}$
\end{prob}

\begin{prob}[S]{2176}
求 $\int x f''(x) \,\mathrm{d}x$.
\end{prob}

\begin{prob}[S]{2177}
求 $\int f'(2x) \,\mathrm{d}x$.
\end{prob}

\begin{prob}[S]{2178}
设 $f'(x^2) = \frac{1}{x} \ (x > 0)$, 求 $f(x)$.
\end{prob}

\begin{prob}[M]{2179}
\begin{enumerate}[label=(\alph*), itemsep=2pt, topsep=3pt]
  \item 设 $f'(\sin^2 x) = \cos^2 x$, 求 $f(x)$.
  \item 设 $f'(\ln x) = \begin{cases} 1, & 0 < x \leqslant 1, \\ x, & 1 < x < +\infty, \end{cases}$ 且 $f(0) = 0$, 求 $f(x)$.
\end{enumerate}
\end{prob}

\begin{prob}[L]{2180}
设 $f(x)$ 为单调的连续函数, $f^{-1}(x)$ 为其反函数. 证明: 若
\[
\int f(x) \,\mathrm{d}x = F(x) + C,
\]
则
\[
\int f^{-1}(x) \,\mathrm{d}x = x f^{-1}(x) - F(f^{-1}(x)) + C.
\]
研究例子:
\begin{enumerate}[label=(\alph*), itemsep=2pt, topsep=3pt]
  \item $f(x) = x^n \ (n > 0)$;
  \item $f(x) = \mathrm{e}^x$;
  \item $f(x) = \arcsin x$;
  \item $f(x) = \operatorname{artanh} x$.
\end{enumerate}
\end{prob}

\probchapter{第四章\quad 定积分}

\probsec{§1. 定积分是积分和的极限}

\begin{prob}[M]{2181}
把区间 $[-1, 4]$ 分为 $n$ 个相等的子区间, 并取这些子区间中点的坐标作自变量 $\xi_i$ 的值 ($i = 0, 1, \cdots, n-1$). 求函数 $f(x) = 1 + x$ 在此区间上的积分和 $S_n$.
\end{prob}

\begin{prob}[L]{2182}
把所给区间分为 $n$ 个相等的部分, 求下列函数 $f(x)$ 在相应区间上的下积分和 $\underline{S}_n$ 与上积分和 $\overline{S}_n$:
\begin{enumerate}[label=(\alph*), itemsep=2pt, topsep=3pt]
  \item $f(x) = x^3 \quad (-2 \leqslant x \leqslant 3)$;
  \item $f(x) = \sqrt{x} \quad (0 \leqslant x \leqslant 1)$;
  \item $f(x) = 2^x \quad (0 \leqslant x \leqslant 10)$.
\end{enumerate}
\end{prob}

\begin{prob}[M]{2183}
把闭区间 $[1, 2]$ 分为 $n$ 份, 使这 $n$ 份的长构成一等比数列, 求函数 $f(x) = x^4$ 在此区间上的下积分和. 当 $n \to \infty$ 时此和的极限等于什么?
\end{prob}

\begin{prob}[M]{2184}
从积分的定义出发, 求
\[
\int_{0}^{T} (v_0 + gt)\,\mathrm{d}t,
\]
其中 $T, v_0, g$ 为常数.
\end{prob}

\probgroup{以适当的方法分割积分区间, 并视积分为相应积分和的极限, 计算下列定积分:}

\begin{prob}[S]{2185}
$\displaystyle\int_{-1}^{2} x^2\,\mathrm{d}x$.
\end{prob}

\begin{prob}[S]{2186}
$\displaystyle\int_{0}^{1} a^x\,\mathrm{d}x \ (a > 0)$.
\end{prob}

\begin{prob}[S]{2187}
$\displaystyle\int_{0}^{\frac{\pi}{2}} \sin x\,\mathrm{d}x$.
\end{prob}

\begin{prob}[S]{2188}
$\displaystyle\int_{0}^{x} \cos t\,\mathrm{d}t$.
\end{prob}

\begin{prob}[M]{2189}
$\displaystyle\int_{a}^{b} \frac{\mathrm{d}x}{x^2} \ (0 < a < b)$.
\probhint{令 $\xi_i = \sqrt{x_i x_{i+1}} \ (i = 0, 1, \cdots, n)$.}
\end{prob}

\begin{prob}[M]{2190}
$\displaystyle\int_{a}^{b} x^m\,\mathrm{d}x \ (0 < a < b; \, m \neq -1)$.
\probhint{选择诸分点, 使它们的横坐标 $x_i$ 构成一等比数列.}
\end{prob}

\begin{prob}[S]{2191}
$\displaystyle\int_{a}^{b} \frac{\mathrm{d}x}{x} \ (0 < a < b)$.
\end{prob}

\begin{prob}[M]{2192}
计算泊松积分
\[
\int_{0}^{\pi} \ln(1 - 2\alpha\cos x + \alpha^2)\,\mathrm{d}x,
\]
考虑两种情形: (a) $|\alpha| < 1$; \qquad (b) $|\alpha| > 1$.
\probhint{分解多项式 $\alpha^{2n} - 1$ 为二次因式.}
\end{prob}

\begin{prob}[M]{2193.1}
设函数 $f(x)$ 及 $\varphi(x)$ 在 $[a, b]$ 上连续, 证明:
\[
\lim_{\max|\Delta x_i| \to 0} \sum_{i=0}^{n-1} f(\xi_i)\varphi(\theta_i)\Delta x_i = \int_{a}^{b} f(x)\varphi(x)\,\mathrm{d}x,
\]
其中 $x_i \leqslant \xi_i \leqslant x_{i+1}, \, x_i \leqslant \theta_i \leqslant x_{i+1} \ (i = 0, 1, \cdots, n-1)$, 且 $\Delta x_i = x_{i+1} - x_i \ (x_0 = a, \, x_n = b)$.
\end{prob}

\begin{prob}[M]{2193.2}
设函数 $f(x)$ 在 $[0, 1]$ 上有界且单调, 证明:
\[
\int_{0}^{1} f(x)\,\mathrm{d}x - \frac{1}{n} \sum_{k=1}^{n} f\left(\frac{k}{n}\right) = O\left(\frac{1}{n}\right).
\]
\end{prob}

\begin{prob}[M]{2193.3}
设函数 $f(x)$ 在闭区间 $[a, b]$ 上有界并凹 (参考习题 1312), 证明:
\[
(b - a)\frac{f(a) + f(b)}{2} \leqslant \int_{a}^{b} f(x)\,\mathrm{d}x \leqslant (b - a)f\left(\frac{a + b}{2}\right).
\]
\end{prob}

\begin{prob}[M]{2193.4}
设 $f(x) \in C^{(2)}[1, +\infty)$, 且当 $x \in [1, +\infty)$ 时 $f(x) \geqslant 0, \, f'(x) \geqslant 0, \, f''(x) \leqslant 0$. 证明: 当 $n \to \infty$ 时
\[
\sum_{k=1}^{n} f(k) = \frac{1}{2}f(n) + \int_{1}^{n} f(x)\,\mathrm{d}x + O(1).
\]
\end{prob}

\begin{prob}[M]{2193.5}
设 $f(x) \in C^{(1)}[a, b]$,
\[
\Delta_n = \int_{a}^{b} f(x)\,\mathrm{d}x - \frac{b - a}{n} \sum_{k=1}^{n} f\left(a + k\frac{b - a}{n}\right),
\]
求 $\displaystyle\lim_{n \to \infty} n\Delta_n$.
\end{prob}

\begin{prob}[M]{2194}
证明: 不连续函数
\[
f(x) = \operatorname{sgn}\left(\sin\frac{\pi}{x}\right)
\]
在区间 $[0, 1]$ 上可积.
\end{prob}

\begin{prob}[M]{2195}
证明: 黎曼函数
\[
\varphi(x) = \begin{cases} 0, & x \text{ 为无理数}, \\ \dfrac{1}{n}, & x = \dfrac{m}{n} \ (m \text{ 及 } n \ (n \geqslant 1) \text{ 为互素的整数}) \end{cases}
\]
在任何有限的区间上可积.
\end{prob}

\begin{prob}[M]{2196}
证明: 函数
\[
f(x) = \begin{cases} \dfrac{1}{x} - \left[\dfrac{1}{x}\right], & x \neq 0, \\ 0, & x = 0 \end{cases}
\]
在闭区间 $[0, 1]$ 上可积.
\end{prob}

\begin{prob}[M]{2197}
证明: 狄利克雷函数
\[
\chi(x) = \begin{cases} 0, & x \text{ 为无理数}, \\ 1, & x \text{ 为有理数} \end{cases}
\]
在任何区间上不可积.
\end{prob}

\begin{prob}[M]{2198}
设函数 $f(x)$ 在 $[a, b]$ 上可积, 且
\[
f_n(x) = \sup f(x) \quad (x_i \leqslant x \leqslant x_{i+1}),
\]
其中 $x_i = a + \dfrac{i}{n}(b - a) \ (i = 0, 1, \cdots, n; \, n = 1, 2, \cdots)$. 证明:
\[
\lim_{n \to \infty} \int_{a}^{b} f_n(x)\,\mathrm{d}x = \int_{a}^{b} f(x)\,\mathrm{d}x.
\]
\end{prob}

\begin{prob}[M]{2199}
证明: 若函数 $f(x)$ 在 $[a, b]$ 上可积, 则存在连续函数 $\varphi_n(x) \ (n = 1, 2, \cdots)$ 的序列, 使得在 $a \leqslant c \leqslant b$ 时
\[
\int_{a}^{c} f(x)\,\mathrm{d}x = \lim_{n \to \infty} \int_{a}^{c} \varphi_n(x)\,\mathrm{d}x.
\]
\end{prob}

\begin{prob}[M]{2200}
证明: 若有界函数 $f(x)$ 在闭区间 $[a, b]$ 上可积, 则其绝对值 $|f(x)|$ 在 $[a, b]$ 上也可积, 并且
\[
\left|\int_{a}^{b} f(x)\,\mathrm{d}x\right| \leqslant \int_{a}^{b} |f(x)|\,\mathrm{d}x.
\]
\end{prob}

\begin{prob}[M]{2201}
设 $f(x)$ 在闭区间 $[a, b]$ 上绝对可积, 即积分 $\displaystyle\int_{a}^{b} |f(x)|\,\mathrm{d}x$ 存在. 这个函数在 $[a, b]$ 上是否为可积函数? 研究例子:
\[
f(x) = \begin{cases} 1, & x \text{ 为有理数}, \\ -1, & x \text{ 为无理数}. \end{cases}
\]
\end{prob}

\begin{prob}[M]{2202}
设函数 $\varphi(x)$ 在闭区间 $[A, B]$ 上有定义并连续, 函数 $f(x)$ 在 $[a, b]$ 上可积, 并且当 $a \leqslant x \leqslant b$ 时 $A \leqslant f(x) \leqslant B$. 证明: 函数 $\varphi(f(x))$ 在 $[a, b]$ 上可积.
\end{prob}

\begin{prob}[M]{2203}
若函数 $f(x)$ 及 $\varphi(x)$ 可积, 则函数 $f(\varphi(x))$ 是否也必定可积? 研究例子:
\[
f(x) = \begin{cases} 0, & x = 0, \\ 1, & x \neq 0, \end{cases}
\]
$\varphi(x)$ 为黎曼函数 (见习题 2195).
\end{prob}

\begin{prob}[M]{2204}
设函数 $f(x)$ 在闭区间 $[A, B]$ 上可积, 证明: 函数 $f(x)$ 有积分连续性, 即
\[
\lim_{h \to 0} \int_{a}^{b} |f(x + h) - f(x)|\,\mathrm{d}x = 0,
\]
其中 $[a, b] \subset [A, B]$.
\end{prob}

\begin{prob}[M]{2205}
设函数 $f(x)$ 在闭区间 $[a, b]$ 上可积, 证明: 等式
\[
\int_{a}^{b} f^2(x)\,\mathrm{d}x = 0
\]
当且仅当对属于闭区间 $[a, b]$ 内函数 $f(x)$ 连续的一切点有 $f(x) = 0$ 时方成立.
\end{prob}

\probsec{§2. 利用不定积分计算定积分的方法}

\probgroup{利用牛顿--莱布尼茨公式, 求下列定积分并绘出相应的曲边图形面积:}

\begin{prob}[S]{2206}
$\displaystyle\int_{-1}^{8} \sqrt[3]{x}\,\mathrm{d}x$.
\end{prob}

\begin{prob}[S]{2207}
$\displaystyle\int_{0}^{\pi} \sin x\,\mathrm{d}x$.
\end{prob}

\begin{prob}[S]{2208}
$\displaystyle\int_{\frac{1}{\sqrt{3}}}^{\sqrt{3}} \frac{\mathrm{d}x}{1 + x^2}$.
\end{prob}

\begin{prob}[S]{2209}
$\displaystyle\int_{-\frac{1}{2}}^{\frac{1}{2}} \frac{\mathrm{d}x}{\sqrt{1 - x^2}}$.
\end{prob}

\begin{prob}[S]{2210}
$\displaystyle\int_{\sinh 1}^{\sinh 2} \frac{\mathrm{d}x}{\sqrt{1 + x^2}}$.
\end{prob}

\begin{prob}[S]{2211}
$\displaystyle\int_{0}^{2} |1 - x|\,\mathrm{d}x$.
\end{prob}

\begin{prob}[S]{2212}
$\displaystyle\int_{-1}^{1} \frac{\mathrm{d}x}{x^2 - 2x\cos\alpha + 1} \ (0 < \alpha < \pi)$.
\end{prob}

\begin{prob}[S]{2213}
$\displaystyle\int_{0}^{2\pi} \frac{\mathrm{d}x}{1 + \varepsilon\cos x} \ (0 \leqslant \varepsilon < 1)$.
\end{prob}

\begin{prob}[S]{2214}
$\displaystyle\int_{-1}^{1} \frac{\mathrm{d}x}{\sqrt{(1 - 2ax + a^2)(1 - 2bx + b^2)}} \ (|a| < 1, \, |b| < 1, \, ab > 0)$.
\end{prob}

\begin{prob}[S]{2215}
$\displaystyle\int_{0}^{\frac{\pi}{2}} \frac{\mathrm{d}x}{a^2 \sin^2 x + b^2 \cos^2 x} \ (ab \neq 0)$.
\end{prob}

\begin{prob}[L]{2216}
对于下列定积分:
\begin{enumerate}[label=(\alph*), itemsep=2pt, topsep=3pt]
  \item $\displaystyle\int_{-1}^{1} \frac{\mathrm{d}x}{x}$;
  \item $\displaystyle\int_{0}^{2\pi} \frac{\sec^2 x\,\mathrm{d}x}{2 + \tan^2 x}$;
  \item $\displaystyle\int_{-1}^{1} \frac{\mathrm{d}}{\mathrm{d}x}\left(\arctan\frac{1}{x}\right)\mathrm{d}x$,
\end{enumerate}
说明为什么运用牛顿--莱布尼茨公式会得到不正确的结果.
\end{prob}

\begin{prob}[S]{2217}
求 $\displaystyle\int_{-1}^{1} \frac{\mathrm{d}}{\mathrm{d}x}\left(\frac{1}{1 + 2^{\frac{1}{x}}}\right)\mathrm{d}x$.
\end{prob}

\begin{prob}[S]{2218}
求 $\displaystyle\int_{0}^{100\pi} \sqrt{1 - \cos 2x}\,\mathrm{d}x$.
\end{prob}

\probgroup{利用定积分求下列和的极限值:}

\begin{prob}[S]{2219}
$\displaystyle\lim_{n \to \infty} \left(\frac{1}{n^2} + \frac{2}{n^2} + \cdots + \frac{n-1}{n^2}\right)$.
\end{prob}

\begin{prob}[S]{2220}
$\displaystyle\lim_{n \to \infty} \left(\frac{1}{n+1} + \frac{1}{n+2} + \cdots + \frac{1}{n+n}\right)$.
\end{prob}

\begin{prob}[S]{2221}
$\displaystyle\lim_{n \to \infty} \left(\frac{n}{n^2 + 1^2} + \frac{n}{n^2 + 2^2} + \cdots + \frac{n}{n^2 + n^2}\right)$.
\end{prob}

\begin{prob}[S]{2222}
$\displaystyle\lim_{n \to \infty} \frac{1}{n}\left(\sin\frac{\pi}{n} + \sin\frac{2\pi}{n} + \cdots + \sin\frac{(n-1)\pi}{n}\right)$.
\end{prob}

\begin{prob}[S]{2223}
$\displaystyle\lim_{n \to \infty} \frac{1^p + 2^p + \cdots + n^p}{n^{p+1}} \ (p > 0)$.
\end{prob}

\begin{prob}[S]{2224}
$\displaystyle\lim_{n \to \infty} \frac{1}{n}\left(\sqrt{1 + \frac{1}{n}} + \sqrt{1 + \frac{2}{n}} + \cdots + \sqrt{1 + \frac{n}{n}}\right)$.
\end{prob}

\begin{prob}[S]{2225}
$\displaystyle\lim_{n \to \infty} \frac{\sqrt[n]{n!}}{n}$.
\end{prob}

\begin{prob}[S]{2226}
$\displaystyle\lim_{n \to \infty} \left[\frac{1}{n} \sum_{k=1}^{n} f\left(a + k\frac{b - a}{n}\right)\right]$.
\end{prob}

\probgroup{弃掉高阶无穷小量, 求下列和的极限值:}

\begin{prob}[S]{2227}
$\displaystyle\lim_{n \to \infty} \left[\left(1 + \frac{1}{n}\right)\sin\frac{\pi}{n^2} + \left(1 + \frac{2}{n}\right)\sin\frac{2\pi}{n^2} + \cdots + \left(1 + \frac{n-1}{n}\right)\sin\frac{(n-1)\pi}{n^2}\right]$.
\end{prob}

\begin{prob}[S]{2228}
$\displaystyle\lim_{n \to \infty} \sin\frac{\pi}{n} \cdot \sum_{k=1}^{n} \frac{1}{2 + \cos\frac{k\pi}{n}}$.
\end{prob}

\begin{prob}[S]{2229}
$\lim_{n\to\infty} \dfrac{\sum\limits_{k=1}^{n} \sqrt{(nx+k)(nx+k+1)}}{n^2}\ (x > 0)$.
\end{prob}

\begin{prob}[S]{2230}
$\lim_{n\to\infty} \left( \dfrac{2^{\frac{1}{n}}}{n+1} + \dfrac{2^{\frac{2}{n}}}{n+\frac{1}{2}} + \cdots + \dfrac{2^{\frac{n}{n}}}{n+\frac{1}{n}} \right)$.
\end{prob}

\begin{prob}[M]{2231}
求：
\[ \frac{\mathrm{d}}{\mathrm{d}x} \int_a^b \sin x^2 \,\mathrm{d}x; \qquad \frac{\mathrm{d}}{\mathrm{d}a} \int_a^b \sin x^2 \,\mathrm{d}x; \qquad \frac{\mathrm{d}}{\mathrm{d}b} \int_a^b \sin x^2 \,\mathrm{d}x. \]
\end{prob}

\begin{prob}[L]{2232}
求：
\begin{enumerate}[label=(\alph*), itemsep=2pt, topsep=3pt]
  \item $\dfrac{\mathrm{d}}{\mathrm{d}x} \int_0^{x^2} \sqrt{1+t^2}\,\mathrm{d}t$;
  \item $\dfrac{\mathrm{d}}{\mathrm{d}x} \int_{x^2}^{x^3} \dfrac{\mathrm{d}t}{\sqrt{1+t^4}}$;
  \item $\dfrac{\mathrm{d}}{\mathrm{d}x} \int_{\sin x}^{\cos x} \cos(\pi t^2)\,\mathrm{d}t$.
\end{enumerate}
\end{prob}

\begin{prob}[L]{2233}
求：
\begin{enumerate}[label=(\alph*), itemsep=2pt, topsep=3pt]
  \item $\lim_{x\to 0} \dfrac{\int_0^x \cos t^2 \,\mathrm{d}t}{x}$;
  \item $\lim_{x\to +\infty} \dfrac{\int_0^x (\arctan t)^2 \,\mathrm{d}t}{\sqrt{x^2+1}}$;
  \item $\lim_{x\to +\infty} \dfrac{\left(\int_0^x \mathrm{e}^{t^2}\,\mathrm{d}t\right)^2}{\int_0^x \mathrm{e}^{2t^2}\,\mathrm{d}t}$;
  \item $\lim_{n\to\infty} \int_0^1 f(nx)\,\mathrm{d}x$，其中 $f(x) \in C[0,+\infty)$，且当 $x\to +\infty$ 时 $f(x) \to A$.
\end{enumerate}
\end{prob}

\begin{prob}[M]{2234}
证明：当 $x\to \infty$ 时
\[ \int_0^x \mathrm{e}^{t^2}\,\mathrm{d}t \sim \frac{1}{2x}\mathrm{e}^{x^2}. \]
\end{prob}

\begin{prob}[M]{2235}
求
\[ \lim_{x\to +0} \frac{\int_0^{\sin x} \sqrt{\tan t}\,\mathrm{d}t}{\int_0^{\tan x} \sqrt{\sin t}\,\mathrm{d}t}. \]
\end{prob}

\begin{prob}[M]{2236}
设 $f(x)$ 为连续正值函数. 证明：当 $x \geqslant 0$ 时，函数
\[ \varphi(x) = \frac{\int_0^x t f(t)\,\mathrm{d}t}{\int_0^x f(t)\,\mathrm{d}t} \]
递增.
\end{prob}

\begin{prob}[M]{2237}
求：
\begin{enumerate}[label=(\alph*), itemsep=2pt, topsep=3pt]
  \item $\int_0^2 f(x)\,\mathrm{d}x$，设 $f(x) = \begin{cases} x^2, & 0 \leqslant x \leqslant 1, \\ 2-x, & 1 < x \leqslant 2; \end{cases}$
  \item $\int_0^1 f(x)\,\mathrm{d}x$，设 $f(x) = \begin{cases} x, & 0 \leqslant x \leqslant t, \\ t \cdot \dfrac{1-x}{1-t}, & t \leqslant x \leqslant 1. \end{cases}$
\end{enumerate}
\end{prob}

\begin{prob}[L]{2238}
计算下列积分并作出这些积分对参数 $\alpha$ 的函数关系 $I = I(\alpha)$ 的图像：
\begin{enumerate}[label=(\alph*), itemsep=2pt, topsep=3pt]
  \item $I = \int_0^1 x|x-\alpha|\,\mathrm{d}x$;
  \item $I = \int_0^\pi \dfrac{\sin^2 x}{1 + 2\alpha \cos x + \alpha^2}\,\mathrm{d}x$;
  \item $I = \int_0^\pi \dfrac{\sin x\,\mathrm{d}x}{\sqrt{1 - 2\alpha \cos x + \alpha^2}}$.
\end{enumerate}
\end{prob}

\probgroup{利用分部积分公式，求下列定积分：}

\begin{prob}[S]{2239}
$\int_0^{\ln 2} x\mathrm{e}^{-x}\,\mathrm{d}x$.
\end{prob}

\begin{prob}[S]{2240}
$\int_0^\pi x\sin x\,\mathrm{d}x$.
\end{prob}

\begin{prob}[S]{2241}
$\int_0^{2\pi} x^2 \cos x\,\mathrm{d}x$.
\end{prob}

\begin{prob}[S]{2242}
$\int_{\frac{1}{\mathrm{e}}}^{\mathrm{e}} |\ln x|\,\mathrm{d}x$.
\end{prob}

\begin{prob}[S]{2243}
$\int_0^1 \arccos x\,\mathrm{d}x$.
\end{prob}

\begin{prob}[S]{2244}
$\int_0^{\sqrt{3}} x\arctan x\,\mathrm{d}x$.
\end{prob}

\probgroup{利用适当的变量代换，求下列定积分：}

\begin{prob}[S]{2245}
$\int_{-1}^1 \dfrac{x\,\mathrm{d}x}{\sqrt{5-4x}}$.
\end{prob}

\begin{prob}[S]{2246}
$\int_0^a x^2\sqrt{a^2-x^2}\,\mathrm{d}x$.
\end{prob}

\begin{prob}[S]{2247}
$\int_0^{0.75} \dfrac{\mathrm{d}x}{(x+1)\sqrt{x^2+1}}$.
\end{prob}

\begin{prob}[S]{2248}
$\int_0^{\ln 2} \sqrt{\mathrm{e}^x - 1}\,\mathrm{d}x$.
\end{prob}

\begin{prob}[S]{2249}
$\int_0^1 \dfrac{\arcsin \sqrt{x}}{\sqrt{x(1-x)}}\,\mathrm{d}x$.
\end{prob}

\begin{prob}[M]{2250}
令 $x - \dfrac{1}{x} = t$，计算积分
\[ \int_{-1}^1 \frac{1+x^2}{1+x^4}\,\mathrm{d}x. \]
\end{prob}

\begin{prob}[L]{2251}
对于下列定积分和代换 $x = \varphi(t)$：
\begin{enumerate}[label=(\alph*), itemsep=2pt, topsep=3pt]
  \item $\int_{-1}^1 \mathrm{d}x,\ t = x^{\frac{2}{3}}$;
  \item $\int_{-1}^1 \dfrac{\mathrm{d}x}{1+x^2},\ x = \dfrac{1}{t}$;
  \item $\int_0^\pi \dfrac{\mathrm{d}x}{1+\sin^2 x},\ \tan x = t$.
\end{enumerate}
说明为什么用 $\varphi(t)$ 代换 $x$ 会引致不正确的结果.
\end{prob}

\begin{prob}[M]{2252}
在积分
\[ \int_0^3 x \sqrt[3]{1-x^2}\,\mathrm{d}x \]
中，令 $x = \sin t$ 是否可以？
\end{prob}

\begin{prob}[M]{2253}
若在积分 $\int_0^1 \sqrt{1-x^2}\,\mathrm{d}x$ 中作变量代换 $x = \sin t$，可否取数 $\pi$ 和 $\dfrac{\pi}{2}$ 作为新的积分上下限？
\end{prob}

\begin{prob}[M]{2254}
证明：若函数 $f(x)$ 在闭区间 $[a, b]$ 内连续，则：
\[ \int_a^b f(x)\,\mathrm{d}x = (b-a) \int_0^1 f(a + (b-a)x)\,\mathrm{d}x. \]
\end{prob}

\begin{prob}[M]{2255}
证明等式：
\[ \int_0^a x^3 f(x^2)\,\mathrm{d}x = \frac{1}{2} \int_0^{a^2} x f(x)\,\mathrm{d}x \quad (a > 0). \]
\end{prob}

\begin{prob}[M]{2256}
设 $f(x)$ 为闭区间 $[A, B] \supset [a, b]$ 上的连续函数，当 $[a-x, b-x] \subset [A, B]$ 时，求 $\dfrac{\mathrm{d}}{\mathrm{d}x} \int_a^b f(x+y)\,\mathrm{d}y$.
\end{prob}

\begin{prob}[M]{2257}
证明：若函数 $f(x)$ 在闭区间 $[0, 1]$ 上连续，则：
\begin{enumerate}[label=(\alph*), itemsep=2pt, topsep=3pt]
  \item $\int_0^{\frac{\pi}{2}} f(\sin x)\,\mathrm{d}x = \int_0^{\frac{\pi}{2}} f(\cos x)\,\mathrm{d}x$;
  \item $\int_0^\pi x f(\sin x)\,\mathrm{d}x = \dfrac{\pi}{2} \int_0^\pi f(\sin x)\,\mathrm{d}x$.
\end{enumerate}
\end{prob}

\begin{prob}[M]{2258}
证明：若函数 $f(x)$ 在闭区间 $[-l, l]$ 上连续，则
\begin{enumerate}[label=(\arabic*), itemsep=2pt, topsep=3pt]
  \item 当函数 $f(x)$ 为偶函数时
  \[ \int_{-l}^l f(x)\,\mathrm{d}x = 2\int_0^l f(x)\,\mathrm{d}x, \]
  \item 当函数 $f(x)$ 奇函数时
  \[ \int_{-l}^l f(x)\,\mathrm{d}x = 0. \]
\end{enumerate}
给出这些事实的几何解释.
\end{prob}

\begin{prob}[M]{2259}
证明：偶函数的原函数中之一个为奇函数，而奇函数的一切原函数皆为偶函数.
\end{prob}

\begin{prob}[M]{2260}
引入新变量
\[ t = x + \frac{1}{x}, \]
计算积分
\[ \int_{\frac{1}{2}}^2 \left(1 + x - \frac{1}{x}\right)\mathrm{e}^{x+\frac{1}{x}}\,\mathrm{d}x. \]
\end{prob}

\begin{prob}[M]{2261}
在积分
\[ \int_0^{2\pi} f(x)\cos x\,\mathrm{d}x \]
中进行变量代换 $\sin x = t$.
\end{prob}

\begin{prob}[M]{2262}
计算积分
\[ \int_{\mathrm{e}^{-2\pi n}}^1 \left| \left[ \cos x \left( \ln \frac{1}{x} \right) \right]' \right| \,\mathrm{d}x \quad (n\ \text{为正整数}). \]
\end{prob}

\begin{prob}[M]{2263}
求积分
\[ \int_0^\pi \frac{x\sin x}{1+\cos^2 x}\,\mathrm{d}x. \]
\end{prob}

\begin{prob}[M]{2264}
设
\[ f(x) = \frac{(x+1)^2(x-1)}{x^3(x-2)}, \]
求积分
\[ \int_{-1}^3 \frac{f'(x)}{1+f^2(x)}\,\mathrm{d}x. \]
\end{prob}

\begin{prob}[M]{2265}
证明：若 $f(x)$ 为定义于 $-\infty < x < +\infty$ 而周期为 $T$ 的连续的周期函数，则
\[ \int_a^{a+T} f(x)\,\mathrm{d}x = \int_0^T f(x)\,\mathrm{d}x, \]
式中 $a$ 为任意的数.
\end{prob}

\begin{prob}[M]{2266}
证明：当 $n$ 为奇数时，函数
\[ F(x) = \int_0^x \sin^n t\,\mathrm{d}t \quad \text{和} \quad G(x) = \int_0^x \cos^n t\,\mathrm{d}t \]
为以 $2\pi$ 为周期的周期函数；而当 $n$ 为偶数时，其中的每一个皆为线性函数与周期函数之和.
\end{prob}

\begin{prob}[M]{2267}
证明：若 $f(x)$ 为以 $T$ 为周期的连续的周期函数，则函数
\[ F(x) = \int_{x_0}^x f(t)\,\mathrm{d}t \]
在一般情形下是线性函数与周期等于 $T$ 的周期函数的和，
\end{prob}

\probgroup{计算下列积分：}

\begin{prob}[S]{2268}
$\int_0^1 x(2-x^2)^{12}\,\mathrm{d}x$.
\end{prob}

\begin{prob}[S]{2269}
$\int_{-1}^1 \dfrac{x\,\mathrm{d}x}{x^2+x+1}$.
\end{prob}

\begin{prob}[S]{2270}
$\int_1^{\mathrm{e}} (x\ln x)^2\,\mathrm{d}x$.
\end{prob}

\begin{prob}[S]{2271}
$\int_1^9 x\sqrt[3]{1-x}\,\mathrm{d}x$.
\end{prob}

\begin{prob}[S]{2272}
$\int_{-2}^{-1} \dfrac{\mathrm{d}x}{x\sqrt{x^2-1}}$.
\end{prob}

\begin{prob}[S]{2273}
$\int_0^1 x^{15}\sqrt{1+3x^8}\,\mathrm{d}x$.
\end{prob}

\begin{prob}[S]{2274}
$\int_0^3 \arcsin \sqrt{\dfrac{x}{1+x}}\,\mathrm{d}x$.
\end{prob}

\begin{prob}[S]{2275}
$\int_0^{2\pi} \dfrac{\mathrm{d}x}{(2+\cos x)(3+\cos x)}$.
\end{prob}

\begin{prob}[S]{2276}
$\int_0^{2\pi} \dfrac{\mathrm{d}x}{\sin^4 x + \cos^4 x}$.
\end{prob}

\begin{prob}[S]{2277}
$\int_0^{\frac{\pi}{2}} \sin x \sin 2x \sin 3x\,\mathrm{d}x$.
\end{prob}

\begin{prob}[S]{2278}
$\int_0^\pi (x\sin x)^2\,\mathrm{d}x$.
\end{prob}

\begin{prob}[S]{2279}
$\int_0^\pi \mathrm{e}^x \cos^2 x\,\mathrm{d}x$.
\end{prob}

\begin{prob}[S]{2280}
$\int_0^{\ln 2} \sinh^4 x\,\mathrm{d}x$.
\end{prob}

\probgroup{利用递推公式来计算下列依赖于取正整数值的参数 $n$ 的积分：}

\begin{prob}[S]{2281}
$I_n = \int_0^{\frac{\pi}{2}} \sin^n x\,\mathrm{d}x$.
\end{prob}

\begin{prob}[S]{2282}
$I_n = \int_0^{\frac{\pi}{2}} \cos^n x\,\mathrm{d}x$.
\end{prob}

\begin{prob}[S]{2283}
$I_n = \int_0^{\frac{\pi}{4}} \tan^{2n} x\,\mathrm{d}x$.
\end{prob}

\begin{prob}[S]{2284}
$I_n = \int_0^1 (1-x^2)^n\,\mathrm{d}x$.
\end{prob}

\begin{prob}[S]{2285}
$I_n = \int_0^1 \dfrac{x^n\,\mathrm{d}x}{\sqrt{1-x^2}}$.
\end{prob}

\begin{prob}[S]{2286}
$I_n = \int_0^1 x^m (\ln x)^n\,\mathrm{d}x$.
\end{prob}

\begin{prob}[S]{2287}
$I_n = \int_0^{\frac{\pi}{4}} \left( \dfrac{\sin x - \cos x}{\sin x + \cos x} \right)^{2n+1}\,\mathrm{d}x$.
\end{prob}

\begin{prob}[M]{2288}
利用欧拉公式
\[ \mathrm{e}^{\mathrm{i}x} = \cos x + \mathrm{i}\sin x, \]
证明：
\[ \int_0^{2\pi} \mathrm{e}^{\mathrm{i}nx}\mathrm{e}^{-\mathrm{i}mx}\,\mathrm{d}x = \begin{cases} 0, & m \neq n, \\ 2\pi, & m = n \end{cases} \quad (n\ \text{及}\ m\ \text{为整数}). \]
\end{prob}

\begin{prob}[M]{2289}
证明：
\[ \int_a^b \mathrm{e}^{(\alpha+\mathrm{i}\beta)x}\,\mathrm{d}x = \frac{\mathrm{e}^{b(\alpha+\mathrm{i}\beta)} - \mathrm{e}^{a(\alpha+\mathrm{i}\beta)}}{\alpha+\mathrm{i}\beta} \quad (\alpha\ \text{及}\ \beta\ \text{为常数}). \]
\end{prob}

\probgroup{利用欧拉公式：$\cos x = \frac{1}{2}(\mathrm{e}^{\mathrm{i}x} + \mathrm{e}^{-\mathrm{i}x}),\ \sin x = \frac{1}{2\mathrm{i}}(\mathrm{e}^{\mathrm{i}x} - \mathrm{e}^{-\mathrm{i}x})$，计算下列积分（$m$ 及 $n$ 为正整数）：}

\begin{prob}[S]{2290}
$\int_0^{\frac{\pi}{2}} \sin^{2m} x \cos^{2n} x\,\mathrm{d}x$.
\end{prob}

\begin{prob}[S]{2291}
$\int_0^\pi \dfrac{\sin nx}{\sin x}\,\mathrm{d}x$.
\end{prob}

\begin{prob}[S]{2292}
$\int_0^\pi \dfrac{\cos(2n+1)x}{\cos x}\,\mathrm{d}x$.
\end{prob}

\begin{prob}[S]{2293}
$\int_0^\pi \cos^n x \cos nx\,\mathrm{d}x$.
\end{prob}

\begin{prob}[S]{2294}
$\int_0^\pi \sin^n x \sin nx\,\mathrm{d}x$.
\end{prob}

\probgroup{求下列积分（$n$ 为正整数）：}

\begin{prob}[S]{2295}
$\int_0^\pi \sin^{n-1} x \cos(n+1)x\,\mathrm{d}x$.
\end{prob}

\begin{prob}[S]{2296}
$\int_0^\pi \cos^{n-1} x \sin(n+1)x\,\mathrm{d}x$.
\end{prob}

\begin{prob}[S]{2297}
$\int_0^{2\pi} \mathrm{e}^{-ax}\cos^{2n} x\,\mathrm{d}x$.
\end{prob}

\begin{prob}[S]{2298}
$\int_0^{\frac{\pi}{2}} \ln \cos x \cdot \cos 2nx\,\mathrm{d}x$.
\end{prob}

\begin{prob}[M]{2299}
多次利用分部积分法，计算欧拉积分：
\[ B(m, n) = \int_0^1 x^{m-1}(1-x)^{n-1}\,\mathrm{d}x \quad (m\ \text{及}\ n\ \text{为正整数}). \]
\end{prob}

\begin{prob}[M]{2300}
勒让德多项式 $P_n(x)$ 可由以下公式定义：
\[ P_n(x) = \frac{1}{2^n n!} \frac{\mathrm{d}^n}{\mathrm{d}x^n}\left[(x^2-1)^n\right] \quad (n = 0, 1, 2, \cdots). \]
证明：
\[ \int_{-1}^1 P_m(x)P_n(x)\,\mathrm{d}x = \begin{cases} 0, & m \neq n, \\ \dfrac{2}{2n+1}, & m = n. \end{cases} \]
\end{prob}

\begin{prob}[M]{2301}
设函数 $f(x)$ 在 $[a, b]$ 上可积，函数 $F(x)$ 在 $[a, b]$ 内除了有限个内点 $c_i\ (i = 1, \cdots, p)$ 及点 $a$ 与 $b$ 外皆满足等式 $F'(x) = f(x)$，而在这些点上 $F(x)$ 有第一类间断点
(广义原函数). 证明:
\[
\int_{a}^{b} f(x)\,\mathrm{d}x = F(b-0) - F(a+0) - \sum_{i=1}^{p} [F(c_i+0) - F(c_i-0)].
\]
\end{prob}

\begin{prob}[L]{2302}
设函数 $f(x)$ 在闭区间 $[a, b]$ 上可积, 而
\[
F(x) = C + \int_{a}^{x} f(\xi)\,\mathrm{d}\xi
\]
为 $f(x)$ 的不定积分. 证明: 函数 $F(x)$ 连续, 且在函数 $f(x)$ 连续的一切点处成立等式
\[
F'(x) = f(x).
\]
问在函数 $f(x)$ 的间断点处函数 $F(x)$ 的导数是什么?

研究例子:
\begin{enumerate}[label=(\alph*), itemsep=2pt, topsep=3pt]
  \item $f\left(\dfrac{1}{n}\right) = 1$ ($n = \pm 1, \pm 2, \cdots$), 当 $x \neq \dfrac{1}{n}$ 时 $f(x) = 0$;
  \item $f(x) = \operatorname{sgn} x$.
\end{enumerate}
\end{prob}

\probgroup{求下列有界不连续函数的不定积分:}

\begin{prob}[S]{2303}
$\int \operatorname{sgn} x\,\mathrm{d}x$.
\end{prob}

\begin{prob}[S]{2304}
$\int \operatorname{sgn}(\sin x)\,\mathrm{d}x$.
\end{prob}

\begin{prob}[S]{2305}
$\int [x]\,\mathrm{d}x$ ($x \geqslant 0$).
\end{prob}

\begin{prob}[S]{2306}
$\int x[x]\,\mathrm{d}x$ ($x \geqslant 0$).
\end{prob}

\begin{prob}[S]{2307}
$\int (-1)^{[x]}\,\mathrm{d}x$.
\end{prob}

\begin{prob}[S]{2308}
$\int_{0}^{x} f(x)\,\mathrm{d}x$, 其中 $f(x) = \begin{cases} 1, & |x| < l, \\ 0, & |x| > l. \end{cases}$
\end{prob}

\probgroup{计算下列有界不连续函数的定积分:}

\begin{prob}[S]{2309}
$\int_{0}^{3} \operatorname{sgn}(x - x^3)\,\mathrm{d}x$.
\end{prob}

\begin{prob}[S]{2310}
$\int_{0}^{2} [\mathrm{e}^x]\,\mathrm{d}x$.
\end{prob}

\begin{prob}[S]{2311}
$\int_{0}^{6} [x]\sin\dfrac{\pi x}{6}\,\mathrm{d}x$.
\end{prob}

\begin{prob}[S]{2312}
$\int_{0}^{\pi} x\operatorname{sgn}(\cos x)\,\mathrm{d}x$.
\end{prob}

\begin{prob}[S]{2313}
$\int_{1}^{n+1} \ln[x]\,\mathrm{d}x$ ($n$ 为正整数).
\end{prob}

\begin{prob}[S]{2314}
$\int_{0}^{1} \operatorname{sgn}[\sin(\ln x)]\,\mathrm{d}x$.
\end{prob}

\begin{prob}[M]{2315}
求 $\int_{E} |\cos x|\sqrt{\sin x}\,\mathrm{d}x$, 其中 $E$ 为闭区间 $[0, 4\pi]$ 中使被积函数有意义的一切值的集合.
\end{prob}

\probsec{§3. 中值定理}

\begin{prob}[L]{2316}
确定下列定积分的符号:
\begin{enumerate}[label=(\alph*), itemsep=2pt, topsep=3pt]
  \item $\int_{0}^{2\pi} x\sin x\,\mathrm{d}x$;
  \item $\int_{0}^{2\pi} \dfrac{\sin x}{x}\,\mathrm{d}x$;
  \item $\int_{-2}^{2} x^3 2^x\,\mathrm{d}x$;
  \item $\int_{\frac{1}{2}}^{1} x^2 \ln x\,\mathrm{d}x$.
\end{enumerate}
\end{prob}

\begin{prob}[L]{2317}
确定哪个积分较大:
\begin{enumerate}[label=(\alph*), itemsep=2pt, topsep=3pt]
  \item $\int_{0}^{\frac{\pi}{2}} \sin^{10} x\,\mathrm{d}x$ 或 $\int_{0}^{\frac{\pi}{2}} \sin^2 x\,\mathrm{d}x$;
  \item $\int_{0}^{1} \mathrm{e}^{-x}\,\mathrm{d}x$ 或 $\int_{0}^{1} \mathrm{e}^{-x^2}\,\mathrm{d}x$;
  \item $\int_{0}^{\pi} \mathrm{e}^{-x^2}\cos^2 x\,\mathrm{d}x$ 或 $\int_{\pi}^{2\pi} \mathrm{e}^{-x^2}\cos^2 x\,\mathrm{d}x$.
\end{enumerate}
\end{prob}

\begin{prob}[L]{2318}
求下列函数在所给区间内的平均值:
\begin{enumerate}[label=(\alph*), itemsep=2pt, topsep=3pt]
  \item $f(x) = x^2$, 在 $[0, 1]$ 上;
  \item $f(x) = \sqrt{x}$, 在 $[0, 100]$ 上;
  \item $f(x) = 10 + 2\sin x + 3\cos x$, 在 $[0, 2\pi]$ 上;
  \item $f(x) = \sin x\sin(x + \varphi)$, 在 $[0, 2\pi]$ 上.
\end{enumerate}
\end{prob}

\begin{prob}[M]{2319.1}
求椭圆之焦径
\[
r = \frac{\rho}{1 - \varepsilon\cos\varphi} \quad (0 < \varepsilon < 1)
\]
之长的平均值.
\end{prob}

\begin{prob}[S]{2319.2}
求初速度为 $v_0$ 之自由落体的速度之平均值.
\end{prob}

\begin{prob}[M]{2320}
电流强度依规律
\[
i = i_0 \sin\left(\frac{2\pi t}{T} + \varphi\right)
\]
变化, 其中 $i_0$ 为振幅, $t$ 为时间, $T$ 为周期, $\varphi$ 为初相. 求电流强度之平方的平均值.
\end{prob}

\begin{prob}[M]{2321}
设 $f(x) \in C[0, +\infty)$ 且 $\lim\limits_{x\to+\infty} f(x) = A$. 求 $\lim\limits_{x\to+\infty} \dfrac{1}{x}\int_{0}^{x} f(t)\,\mathrm{d}t$. 研究例子 $f(x) = \arctan x$.
\end{prob}

\begin{prob}[L]{2322}
命
\[
\int_{0}^{x} f(t)\,\mathrm{d}t = x f(\theta x).
\]
求 $\theta$, 设: (a) $f(t) = t^n$ ($n > -1$); (b) $f(t) = \ln t$; (c) $f(t) = \mathrm{e}^t$. $\lim\limits_{x\to+0} \theta$ 及 $\lim\limits_{x\to+\infty} \theta$ 等于什么?
\end{prob}

\probgroup{利用第一中值定理, 估计积分:}

\begin{prob}[S]{2323}
$\int_{0}^{2\pi} \dfrac{\mathrm{d}x}{1 + 0.5\cos x}$.
\end{prob}

\begin{prob}[S]{2324}
$\int_{0}^{1} \dfrac{x^9}{\sqrt{1 + x}}\,\mathrm{d}x$.
\end{prob}

\begin{prob}[S]{2325}
$\int_{0}^{100} \dfrac{\mathrm{e}^{-x}}{x + 100}\,\mathrm{d}x$.
\end{prob}

\begin{prob}[M]{2326.1}
证明等式:

(a) $\lim\limits_{n\to\infty} \int_{0}^{1} \dfrac{x^n}{1 + x}\,\mathrm{d}x = 0$; \qquad (b) $\lim\limits_{n\to\infty} \int_{0}^{\frac{\pi}{2}} \sin^n x\,\mathrm{d}x = 0$.
\end{prob}

\begin{prob}[M]{2326.2}
求:

(a) $\lim\limits_{\varepsilon\to 0} \int_{0}^{1} \dfrac{\mathrm{d}x}{\varepsilon x^2 + 1}$; \qquad (b) $\lim\limits_{\varepsilon\to+0} \int_{a\varepsilon}^{b\varepsilon} f(x)\,\dfrac{\mathrm{d}x}{x}$ ($a > 0, b > 0, f(x) \in C[0, 1]$).
\end{prob}

\begin{prob}[M]{2327}
设函数 $f(x)$ 在 $[a, b]$ 上连续, 而 $\varphi(x)$ 在 $[a, b]$ 上连续且在 $(a, b)$ 上可微, 并且当 $a < x < b$ 时 $\varphi'(x) \geqslant 0$. 应用分部积分法及第一中值定理, 证明第二中值定理.
\end{prob}

\probgroup{利用第二中值定理, 估计积分:}

\begin{prob}[S]{2328}
$\int_{100\pi}^{200\pi} \dfrac{\sin x}{x}\,\mathrm{d}x$.
\end{prob}

\begin{prob}[S]{2329}
$\int_{a}^{b} \dfrac{\mathrm{e}^{-\alpha x}}{x}\sin x\,\mathrm{d}x$ ($\alpha \geqslant 0; 0 < a < b$).
\end{prob}

\begin{prob}[S]{2330}
$\int_{a}^{b} \sin x^2\,\mathrm{d}x$ ($0 < a < b$).
\end{prob}

\begin{prob}[M]{2331}
设函数 $\varphi(x)$ 及 $\psi(x)$ 和它们的平方在区间 $[a, b]$ 上可积. 证明柯西--布尼亚科夫斯基不等式:
\[
\left\{ \int_{a}^{b} \varphi(x)\psi(x)\,\mathrm{d}x \right\}^2 \leqslant \int_{a}^{b} \varphi^2(x)\,\mathrm{d}x \int_{a}^{b} \psi^2(x)\,\mathrm{d}x.
\]
\end{prob}

\begin{prob}[M]{2332}
设函数 $f(x)$ 在闭区间 $[a, b]$ 上连续可微且 $f(a) = 0$. 证明不等式:
\[
M^2 \leqslant (b - a) \int_{a}^{b} f'^2(x)\,\mathrm{d}x,
\]
其中 $M = \sup\limits_{a \leqslant x \leqslant b} |f(x)|$.
\end{prob}

\begin{prob}[M]{2333}
证明等式:
\[
\lim_{n\to\infty} \int_{n}^{n+p} \frac{\sin x}{x}\,\mathrm{d}x = 0 \quad (p > 0).
\]
\end{prob}

\probsec{§4. 广义积分}

\probgroup{计算下列积分:}

\begin{prob}[S]{2334}
$\int_{a}^{+\infty} \dfrac{\mathrm{d}x}{x^2}$ ($a > 0$).
\end{prob}

\begin{prob}[S]{2335}
$\int_{0}^{1} \ln x\,\mathrm{d}x$.
\end{prob}

\begin{prob}[S]{2336}
$\int_{-\infty}^{+\infty} \dfrac{\mathrm{d}x}{1 + x^2}$.
\end{prob}

\begin{prob}[S]{2337}
$\int_{-1}^{1} \dfrac{\mathrm{d}x}{\sqrt{1 - x^2}}$.
\end{prob}

\begin{prob}[S]{2338}
$\int_{2}^{+\infty} \dfrac{\mathrm{d}x}{x^2 + x - 2}$.
\end{prob}

\begin{prob}[S]{2339}
$\int_{-\infty}^{+\infty} \dfrac{\mathrm{d}x}{(x^2 + x + 1)^2}$.
\end{prob}

\begin{prob}[S]{2340}
$\int_{0}^{+\infty} \dfrac{\mathrm{d}x}{1 + x^3}$.
\end{prob}

\begin{prob}[S]{2341}
$\int_{0}^{+\infty} \dfrac{x^2 + 1}{x^4 + 1}\,\mathrm{d}x$.
\end{prob}

\begin{prob}[S]{2342}
$\int_{0}^{1} \dfrac{\mathrm{d}x}{(2 - x)\sqrt{1 - x}}$.
\end{prob}

\begin{prob}[S]{2343}
$\int_{1}^{+\infty} \dfrac{\mathrm{d}x}{x\sqrt{1 + x^5 + x^{10}}}$.
\end{prob}

\begin{prob}[S]{2344}
$\int_{0}^{+\infty} \dfrac{x\ln x}{(1 + x^2)^2}\,\mathrm{d}x$.
\end{prob}

\begin{prob}[S]{2345}
$\int_{0}^{+\infty} \dfrac{\arctan x}{(1 + x^2)^{\frac{3}{2}}}\,\mathrm{d}x$.
\end{prob}

\begin{prob}[S]{2346}
$\int_{0}^{+\infty} \mathrm{e}^{-ax}\cos bx\,\mathrm{d}x$ ($a > 0$).
\end{prob}

\begin{prob}[S]{2347}
$\int_{0}^{+\infty} \mathrm{e}^{-ax}\sin bx\,\mathrm{d}x$ ($a > 0$).
\end{prob}

\probgroup{利用递推公式计算下列广义积分 ($n$ 为正整数):}

\begin{prob}[S]{2348}
$I_n = \int_{0}^{+\infty} x^n \mathrm{e}^{-x}\,\mathrm{d}x$.
\end{prob}

\begin{prob}[S]{2349}
$I_n = \int_{-\infty}^{+\infty} \dfrac{\mathrm{d}x}{(ax^2 + 2bx + c)^n}$ ($ac - b^2 > 0$).
\end{prob}

\begin{prob}[S]{2350}
$I_n = \int_{1}^{+\infty} \dfrac{\mathrm{d}x}{x(x + 1)\cdots(x + n)}$.
\end{prob}

\begin{prob}[S]{2351}
$I_n = \int_{0}^{1} \dfrac{x^n\,\mathrm{d}x}{\sqrt{(1 - x)(1 + x)}}$.
\end{prob}

\begin{prob}[S]{2352}
$I_n = \int_{0}^{+\infty} \dfrac{\mathrm{d}x}{\cosh^{n+1} x}$.
\end{prob}

\begin{prob}[M]{2353}
(a) $\int_{0}^{\frac{\pi}{2}} \ln\sin x\,\mathrm{d}x$; \qquad (b) $\int_{0}^{\frac{\pi}{2}} \ln\cos x\,\mathrm{d}x$.
\end{prob}

\begin{prob}[M]{2354}
求
\[
\int_E \frac{|\sin x - \cos x|}{\sqrt{\sin x}} \mathrm{e}^{-\frac{x}{2}}\,\mathrm{d}x,
\]
其中 $E$ 为区间 $(0, +\infty)$ 中使被积函数有意义的一切 $x$ 值的集合.
\end{prob}

\begin{prob}[M]{2355}
证明等式:
\[
\int_0^{+\infty} f\left(ax + \frac{b}{x}\right)\,\mathrm{d}x = \frac{1}{a} \int_0^{+\infty} f\left(\sqrt{x^2 + 4ab}\right)\,\mathrm{d}x,
\]
其中 $a > 0, b > 0$ (假定等式左端的积分有意义).
\end{prob}

\begin{prob}[L]{2356}
数
\[
M[f] = \lim_{x \to +\infty} \frac{1}{x} \int_0^x f(\xi)\,\mathrm{d}\xi
\]
称为函数 $f(x)$ 在区间 $(0, +\infty)$ 上的平均值. 求下列函数在此区间上的平均值:
\begin{enumerate}[label=(\alph*), itemsep=2pt, topsep=3pt]
\item $f(x) = \sin^2 x + \cos^2(x\sqrt{2})$;
\item $f(x) = \arctan x$;
\item $f(x) = \sqrt{x} \sin x$.
\end{enumerate}
\end{prob}

\begin{prob}[L]{2357}
求:
\begin{enumerate}[label=(\alph*), itemsep=2pt, topsep=3pt]
\item $\displaystyle\lim_{x \to 0} x \int_x^1 \frac{\cos t}{t^2}\,\mathrm{d}t$;
\item $\displaystyle\lim_{x \to \infty} \frac{\int_0^x \sqrt{1 + t^4}\,\mathrm{d}t}{x^3}$;
\item $\displaystyle\lim_{x \to +0} \frac{\int_x^{+\infty} t^{-1} \mathrm{e}^{-t}\,\mathrm{d}t}{\ln \frac{1}{x}}$;
\item $\displaystyle\lim_{x \to 0} x^\alpha \int_x^1 \frac{f(t)}{t^{\alpha + 1}}\,\mathrm{d}t$,
\end{enumerate}
其中 $\alpha > 0$, $f(t)$ 为闭区间 $[0, 1]$ 上的连续函数.
\end{prob}

\probgroup{研究下列积分的收敛性:}

\begin{prob}[S]{2358}
$\displaystyle\int_0^{+\infty} \frac{x^2\,\mathrm{d}x}{x^4 - x^2 + 1}$.
\end{prob}

\begin{prob}[S]{2359}
$\displaystyle\int_1^{+\infty} \frac{\mathrm{d}x}{x \sqrt[3]{x^2 + 1}}$.
\end{prob}

\begin{prob}[S]{2360}
$\displaystyle\int_0^2 \frac{\mathrm{d}x}{\ln x}$.
\end{prob}

\begin{prob}[S]{2361}
$\displaystyle\int_0^{+\infty} x^{p - 1} \mathrm{e}^{-x}\,\mathrm{d}x$.
\end{prob}

\begin{prob}[S]{2362}
$\displaystyle\int_0^1 x^p \ln^q \frac{1}{x}\,\mathrm{d}x$.
\end{prob}

\begin{prob}[S]{2363}
$\displaystyle\int_0^{+\infty} \frac{x^m}{1 + x^n}\,\mathrm{d}x\ (n \geqslant 0)$.
\end{prob}

\begin{prob}[S]{2364}
$\displaystyle\int_0^{+\infty} \frac{\arctan ax}{x^n}\,\mathrm{d}x\ (a \neq 0)$.
\end{prob}

\begin{prob}[S]{2365}
$\displaystyle\int_0^{+\infty} \frac{\ln(1 + x)}{x^n}\,\mathrm{d}x$.
\end{prob}

\begin{prob}[S]{2366}
$\displaystyle\int_0^{+\infty} \frac{x^m \arctan x}{2 + x^n}\,\mathrm{d}x\ (n \geqslant 0)$.
\end{prob}

\begin{prob}[S]{2367}
$\displaystyle\int_0^{+\infty} \frac{\cos ax}{1 + x^n}\,\mathrm{d}x\ (n \geqslant 0)$.
\end{prob}

\begin{prob}[S]{2368}
$\displaystyle\int_0^{+\infty} \frac{\sin^2 x}{x}\,\mathrm{d}x$.
\end{prob}

\begin{prob}[S]{2369}
$\displaystyle\int_0^{\frac{\pi}{2}} \frac{\mathrm{d}x}{\sin^p x \cos^q x}$.
\end{prob}

\begin{prob}[M]{2370}
(a) $\displaystyle\int_0^1 \frac{x^n\,\mathrm{d}x}{\sqrt{1 - x^2}}$; \qquad (b) $\displaystyle\int_0^{+\infty} \frac{\mathrm{d}x}{\sqrt{x^3 + x}}$.
\end{prob}

\begin{prob}[S]{2371}
$\displaystyle\int_0^{+\infty} \frac{\mathrm{d}x}{x^p + x^q}$.
\end{prob}

\begin{prob}[S]{2372}
$\displaystyle\int_0^1 \frac{\ln x}{1 - x^2}\,\mathrm{d}x$.
\end{prob}

\begin{prob}[S]{2373}
$\displaystyle\int_0^{\frac{\pi}{2}} \frac{\ln(\sin x)}{\sqrt{x}}\,\mathrm{d}x$.
\end{prob}

\begin{prob}[S]{2374}
$\displaystyle\int_1^{+\infty} \frac{\mathrm{d}x}{x^p \ln^q x}$.
\end{prob}

\begin{prob}[S]{2375}
$\displaystyle\int_{\mathrm{e}}^{+\infty} \frac{\mathrm{d}x}{x^p (\ln x)^q (\ln \ln x)^r}$.
\end{prob}

\begin{prob}[M]{2376}
\begin{enumerate}[label=(\alph*), itemsep=2pt, topsep=3pt]
\item $\displaystyle\int_{-\infty}^{+\infty} \frac{\mathrm{d}x}{|x - a_1|^{p_1} |x - a_2|^{p_2} \cdots |x - a_n|^{p_n}}\ (a_1 < a_2 < \cdots < a_n)$;
\item $\displaystyle\int_0^{+\infty} x^\alpha |x - 1|^\beta\,\mathrm{d}x$.
\end{enumerate}
\end{prob}

\begin{prob}[M]{2377}
$\displaystyle\int_0^{+\infty} \frac{P_m(x)}{P_n(x)}\,\mathrm{d}x$, 式中 $P_m(x)$ 和 $P_n(x)$ 为次数分别为 $m$ 及 $n$ 的互素的多项式.
\end{prob}

\probgroup{研究下列积分的绝对收敛性和条件收敛性:}

\begin{prob}[S]{2378}
$\displaystyle\int_0^{+\infty} \frac{\sin x}{x}\,\mathrm{d}x$. \probhint{$|\sin x| \geqslant \sin^2 x$.}
\end{prob}

\begin{prob}[S]{2379}
$\displaystyle\int_0^{+\infty} \frac{\sqrt{x} \cos x}{x + 100}\,\mathrm{d}x$.
\end{prob}

\begin{prob}[L]{2380}
\begin{enumerate}[label=(\alph*), itemsep=2pt, topsep=3pt]
\item $\displaystyle\int_0^{+\infty} x^p \sin(x^q)\,\mathrm{d}x\ (q \neq 0)$;
\item $\displaystyle\int_0^{\frac{\pi}{2}} \sin(\sec x)\,\mathrm{d}x$;
\item $\displaystyle\int_0^{+\infty} x^2 \cos(\mathrm{e}^x)\,\mathrm{d}x$.
\end{enumerate}
\end{prob}

\begin{prob}[S]{2381}
$\displaystyle\int_0^{+\infty} \frac{x^p \sin x}{1 + x^q}\,\mathrm{d}x\ (q \geqslant 0)$.
\end{prob}

\begin{prob}[S]{2382}
$\displaystyle\int_0^{+\infty} \frac{\sin\left(x + \frac{1}{x}\right)}{x^n}\,\mathrm{d}x$.
\end{prob}

\begin{prob}[M]{2383}
$\displaystyle\int_a^{+\infty} \frac{P_m(x)}{P_n(x)} \sin x\,\mathrm{d}x$, 式中 $P_m(x)$ 及 $P_n(x)$ 为整多项式, 且若 $x \geqslant a \geqslant 0$, 则 $P_n(x) > 0$.
\end{prob}

\begin{prob}[M]{2384.1}
若 $\displaystyle\int_a^{+\infty} f(x)\,\mathrm{d}x$ 收敛, 则当 $x \to +\infty$ 时是否必有 $f(x) \to 0$? 研究例子:
\begin{enumerate}[label=(\alph*), itemsep=2pt, topsep=3pt]
\item $\displaystyle\int_0^{+\infty} \sin(x^2)\,\mathrm{d}x$;
\item $\displaystyle\int_0^{+\infty} (-1)^{[x^2]}\,\mathrm{d}x$.
\end{enumerate}
\end{prob}

\begin{prob}[M]{2384.2}
设 $f(x) \in C^{(1)}[x_0, +\infty)$, 当 $x_0 \leqslant x < +\infty$ 时 $|f'(x)| < C$, 且
\[
\int_{x_0}^{+\infty} |f(x)|\,\mathrm{d}x
\]
收敛. 证明: 当 $x \to +\infty$ 时 $f(x) \to 0$. \probhint{研究积分 $\displaystyle\int_{x_0}^{+\infty} f(x) f'(x)\,\mathrm{d}x$.}
\end{prob}

\begin{prob}[M]{2385}
设函数在 $[a, b]$ 上有定义且无界. 可否把函数 $f(x)$ 的收敛广义积分
\[
\int_a^b f(x)\,\mathrm{d}x
\]
看作相应积分和
\[
\sum_{i=0}^{n-1} f(\xi_i) \Delta x_i \quad (x_i \leqslant \xi_i \leqslant x_{i+1},\ \Delta x_i = x_{i+1} - x_i)
\]
的极限?
\end{prob}

\begin{prob}[M]{2386}
设:
\[
\int_a^{+\infty} f(x)\,\mathrm{d}x \tag{1}
\]
收敛, 函数 $\varphi(x)$ 有界, 则积分
\[
\int_a^{+\infty} f(x)\varphi(x)\,\mathrm{d}x \tag{2}
\]
是否必定收敛? 举出适当的例子.

若积分 $(1)$ 绝对收敛, 问积分 $(2)$ 的收敛性如何?
\end{prob}

\begin{prob}[M]{2387}
证明: 若 $\displaystyle\int_a^{+\infty} f(x)\,\mathrm{d}x$ 收敛, $f(x)$ 为单调函数, 则 $f(x) = O\left(\frac{1}{x}\right)$.
\end{prob}

\begin{prob}[M]{2388}
设函数 $f(x)$ 在区间 $0 < x \leqslant 1$ 内是单调函数, 且在点 $x = 0$ 的邻域内是无界的. 证明: 若
\[
\int_0^1 f(x)\,\mathrm{d}x
\]
存在, 则
\[
\lim_{n \to \infty} \frac{1}{n} \sum_{k=1}^n f\left(\frac{k}{n}\right) = \int_0^1 f(x)\,\mathrm{d}x.
\]
\end{prob}

\begin{prob}[M]{2389}
证明: 若函数 $f(x)$ 在区间 $0 < x < a$ 内单调且有界, 且
\[
\int_0^a x^p f(x)\,\mathrm{d}x
\]
存在, 则
\[
\lim_{x \to +0} x^{p+1} f(x) = 0.
\]
\end{prob}

\begin{prob}[L]{2390}
证明:
\begin{enumerate}[label=(\alph*), itemsep=2pt, topsep=3pt]
\item $\mathrm{v.p.} \displaystyle\int_{-1}^1 \frac{\mathrm{d}x}{x} = 0$;
\item $\mathrm{v.p.} \displaystyle\int_0^{+\infty} \frac{\mathrm{d}x}{1 - x^2} = 0$;
\item $\mathrm{v.p.} \displaystyle\int_{-\infty}^{+\infty} \sin x\,\mathrm{d}x = 0$.
\end{enumerate}
\end{prob}

\begin{prob}[M]{2391}
证明: 当 $x \geqslant 0$ 且 $x \neq 1$ 时存在
\[
\mathrm{li}\,x = \mathrm{v.p.} \int_0^x \frac{\mathrm{d}\xi}{\ln \xi}.
\]
\end{prob}

\probgroup{求下列积分:}

\begin{prob}[S]{2392}
$\mathrm{v.p.} \displaystyle\int_0^{+\infty} \frac{\mathrm{d}x}{x^2 - 3x + 2}$.
\end{prob}

\begin{prob}[S]{2393}
$\mathrm{v.p.} \displaystyle\int_{\frac{1}{2}}^2 \frac{\mathrm{d}x}{x \ln x}$.
\end{prob}

\begin{prob}[S]{2394}
$\mathrm{v.p.} \displaystyle\int_{-\infty}^{+\infty} \frac{1 + x}{1 + x^2}\,\mathrm{d}x$.
\end{prob}

\begin{prob}[S]{2395}
$\mathrm{v.p.} \displaystyle\int_{-\infty}^{+\infty} \arctan x\,\mathrm{d}x$.
\end{prob}

\probsec{§5. 面积的计算法}

\begin{prob}[M]{2396}
证明: 正抛物线弓形的面积等于
\[
S = \frac{2}{3} bh,
\]
式中 $b$ 为弓形的底, $h$ 为高 (图 13).
\probfig{fig_2396.png}{图 13}
\end{prob}

\probgroup{求下列直角坐标方程所给曲线为界的图形的面积\footnote{在第四章的这一节和以后各节都把一切的参数当作是正的.}:}

\begin{prob}[S]{2397}
$ax = y^2, ay = x^2$.
\end{prob}

\begin{prob}[S]{2398}
$y = x^2, x + y = 2$.
\end{prob}

\begin{prob}[S]{2399}
$y = 2x - x^2, x + y = 0$.
\end{prob}

\begin{prob}[L]{2400}
\begin{enumerate}[label=(\alph*), itemsep=2pt, topsep=3pt]
\item $y = |\lg x|, y = 0, x = 0, 1, x = 10$;
\item $y = 2^x, y = 2, x = 0$;
\item $y = (x + 1)^2, y = \sin \pi y, y = 0\ (0 \leqslant y \leqslant 1)$.
\end{enumerate}
\end{prob}

\begin{prob}[S]{2401}
$y = x, y = x + \sin^2 x\ (0 \leqslant x \leqslant \pi)$.
\end{prob}

\begin{prob}[S]{2402}
$y = \dfrac{a^3}{a^2 + x^2}, y = 0$.
\end{prob}

\begin{prob}[S]{2403}
$\dfrac{x^2}{a^2} + \dfrac{y^2}{b^2} = 1$.
\end{prob}

\begin{prob}[S]{2404}
$y^2 = x^2(a^2 - x^2)$.
\end{prob}

\begin{prob}[S]{2405}
$y^2 = 2px, 27py^2 = 8(x - p)^3$.
\end{prob}

\begin{prob}[S]{2406}
$Ax^2 + 2Bxy + Cy^2 = 1\ (A > 0, AC - B^2 > 0)$.
\end{prob}

\begin{prob}[S]{2407}
$y^2 = \dfrac{x^3}{2a - x}$ (蔓叶线), $x = 2a$.
\end{prob}

\begin{prob}[S]{2408}
$x = a \ln \dfrac{a + \sqrt{a^2 - y^2}}{y} - \sqrt{a^2 - y^2}$ (曳物线), $y = 0$.
\end{prob}

\begin{prob}[S]{2409}
$y^2 = \dfrac{x^n}{(1 + x^{n+2})^2}\ (x > 0; n > -2)$.
\end{prob}

\begin{prob}[S]{2410}
$y = \mathrm{e}^{-x} |\sin x|, y = 0\ (x \geqslant 0)$.
\end{prob}

\begin{prob}[M]{2411}
抛物线 $y^2 = 2x$ 把圆 $x^2 + y^2 = 8$ 的面积分为两部分, 这两部分的比如何?
\end{prob}

\begin{prob}[M]{2412}
把双曲线 $x^2 - y^2 = 1$ 上的点 $M(x, y)$ 的坐标表示为双曲线扇形 $OM'M$ 的面积 $S$ 的函数, 此扇形以双曲线的弧 $M'M$ 与二射线 $OM$ 及 $OM'$ 为界, 其中 $M'(x, -y)$ 是点 $M$ 相对于 $Ox$ 轴的对称点.
\end{prob}

\probgroup{求下列参数方程所给曲线所围图形的面积:}

\begin{prob}[S]{2413}
$x = a(t - \sin t), y = a(1 - \cos t)\ (0 \leqslant t \leqslant 2\pi)$ (摆线) 及 $y = 0$.
\end{prob}

\begin{prob}[S]{2414}
$x = 2t - t^2, y = 2t^2 - t^3$.
\end{prob}

\begin{prob}[S]{2415}
$x = a(\cos t + t \sin t), y = a(\sin t - t \cos t)\ (0 \leqslant t \leqslant 2\pi)$ (圆的渐伸线) 及 $x = a, y \leqslant 0$.
\end{prob}

\begin{prob}[S]{2416}
$x = a(2\cos t - \cos 2t), y = a(2\sin t - \sin 2t)$.
\end{prob}

\begin{prob}[M]{2417}
\begin{enumerate}[label=(\alph*), itemsep=2pt, topsep=3pt]
\item $x = \dfrac{c^2}{a} \cos^3 t, y = \dfrac{c^2}{b} \sin^3 t\ (c^2 = a^2 - b^2)$ (椭圆的渐屈线);
\item $x = a \cos t, y = \dfrac{a \sin^2 t}{2 + \sin t}$.
\end{enumerate}
\end{prob}

\probgroup{求下列极坐标方程所给曲线所围图形 $S$ 的面积:}

\begin{prob}[S]{2418}
$r^2 = a^2 \cos 2\varphi$ (双纽线).
\end{prob}

\begin{prob}[S]{2419}
$r = a(1 + \cos \varphi)$ (心脏线).
\end{prob}

\begin{prob}[S]{2420}
$r = a \sin 3\varphi$ (三叶线).
\end{prob}

\begin{prob}[S]{2421}
$r = \dfrac{p}{1 - \cos \varphi}$ (抛物线), $\varphi = \dfrac{\pi}{4}, \varphi = \dfrac{\pi}{2}$.
\end{prob}

\begin{prob}[L]{2422}
\begin{enumerate}[label=(\alph*), itemsep=2pt, topsep=3pt]
\item $r = \dfrac{p}{1 + \varepsilon \cos \varphi}\ (0 < \varepsilon < 1)$ (椭圆);
\item $r = 3 + 2\cos \varphi$;
\item $r = \dfrac{1}{\varphi}, r = \dfrac{1}{\sin \varphi}\ \left(0 < \varphi \leqslant \frac{\pi}{2}\right)$.
\end{enumerate}
\end{prob}

\begin{prob}[S]{2423}
$r = a\cos\varphi$, $r = a(\cos\varphi + \sin\varphi)$ $\left(M\left(\frac{a}{2}, 0\right) \in S\right)$.
\end{prob}

\begin{prob}[S]{2424}
求由曲线 $\varphi = r\arctan r$ 及二射线 $\varphi = 0$ 及 $\varphi = \frac{\pi}{\sqrt{3}}$ 所围成之扇形的面积.
\end{prob}

\begin{prob}[L]{2425}
求下列曲线所围图形的面积:
\begin{enumerate}[label=(\alph*), itemsep=2pt, topsep=3pt]
\item $r^2 + \varphi^2 = 1$;
\item 曲线 $\varphi = \sin(\pi r)$ ($0 \leqslant r \leqslant 1$) 的蔓叶线;
\item $\varphi = 4r - r^3$, $\varphi = 0$;
\item $\varphi = r - \sin r$, $\varphi = \pi$;
\item 封闭曲线 $r = \frac{2at}{1+t^2}$, $\varphi = \frac{\pi t}{1+t}$.
\end{enumerate}
\end{prob}

\probgroup{变换为极坐标, 求下列曲线所围图形的面积:}

\begin{prob}[S]{2426}
$x^3 + y^3 = 3axy$ (笛卡儿叶形线).
\end{prob}

\begin{prob}[S]{2427}
$x^4 + y^4 = a^2(x^2 + y^2)$.
\end{prob}

\begin{prob}[S]{2428}
$(x^2 + y^2)^2 = 2a^2 xy$ (双纽线).
\end{prob}

\probgroup{用参数方程的形式给出下列曲线, 再求曲线所围图形的面积:}

\begin{prob}[S]{2429}
$x^{\frac{2}{3}} + y^{\frac{2}{3}} = a^{\frac{2}{3}}$ (星形线).
\end{prob}

\begin{prob}[M]{2430}
$x^4 + y^4 = ax^2 y$.
\probhint{令 $y = tx$.}
\end{prob}

\probsec{§6. 弧长的计算法}

\probgroup{求下列曲线的弧长:}

\begin{prob}[S]{2431}
$y = x^{\frac{3}{2}}$ ($0 \leqslant x \leqslant 4$).
\end{prob}

\begin{prob}[S]{2432}
$y^2 = 2px$ ($0 \leqslant x \leqslant x_0$).
\end{prob}

\begin{prob}[S]{2433}
$y = a\cosh\frac{x}{a}$, 从点 $A(0, a)$ 至点 $B(b, h)$.
\end{prob}

\begin{prob}[S]{2434}
$y = \mathrm{e}^x$ ($0 \leqslant x \leqslant x_0$).
\end{prob}

\begin{prob}[S]{2435}
$x = \frac{1}{4}y^2 - \frac{1}{2}\ln y$ ($1 \leqslant y \leqslant \mathrm{e}$).
\end{prob}

\begin{prob}[S]{2436}
$y = a\ln\frac{a^2}{a^2 - x^2}$ ($0 \leqslant x \leqslant b < a$).
\end{prob}

\begin{prob}[S]{2437}
$y = \ln\cos x$ $\left(0 \leqslant x \leqslant a < \frac{\pi}{2}\right)$.
\end{prob}

\begin{prob}[S]{2438}
$x = a\ln\frac{a + \sqrt{a^2 - y^2}}{y} - \sqrt{a^2 - y^2}$ ($0 < b \leqslant y \leqslant a$).
\end{prob}

\begin{prob}[S]{2439}
$y^2 = \frac{x^3}{2a - x}$ $\left(0 \leqslant x \leqslant \frac{5}{3}a\right)$.
\end{prob}

\begin{prob}[S]{2440}
$x^{\frac{2}{3}} + y^{\frac{2}{3}} = a^{\frac{2}{3}}$ (星形线).
\end{prob}

\begin{prob}[S]{2441}
$x = \frac{c^2}{a}\cos^3 t$, $y = \frac{c^2}{b}\sin^3 t$, $c^2 = a^2 - b^2$ (椭圆的渐屈线).
\end{prob}

\begin{prob}[S]{2442}
$x = a\cos^4 t$, $y = a\sin^4 t$.
\end{prob}

\begin{prob}[S]{2443}
$x = a(t - \sin t)$, $y = a(1 - \cos t)$ ($0 \leqslant t \leqslant 2\pi$).
\end{prob}

\begin{prob}[S]{2444}
$x = a(\cos t + t\sin t)$, $y = a(\sin t - t\cos t)$ ($0 \leqslant t \leqslant 2\pi$) (圆的渐伸线).
\end{prob}

\begin{prob}[M]{2445}
(a) $x = a(\sinh t - t)$, $y = a(\cosh t - 1)$ ($0 \leqslant t \leqslant T$); \qquad (b) $x = \cosh^3 t$, $y = \sinh^3 t$ ($0 \leqslant t \leqslant T$).
\end{prob}

\begin{prob}[S]{2446}
$r = a\varphi$ (阿基米德螺线) ($0 \leqslant \varphi \leqslant 2\pi$).
\end{prob}

\begin{prob}[S]{2447}
$r = a\mathrm{e}^{m\varphi}$ ($m > 0, 0 < r < a$).
\end{prob}

\begin{prob}[S]{2448}
$r = a(1 + \cos\varphi)$.
\end{prob}

\begin{prob}[S]{2449}
$r = \frac{p}{1 + \cos\varphi}$ $\left(|\varphi| \leqslant \frac{\pi}{2}\right)$.
\end{prob}

\begin{prob}[S]{2450}
$r = a\sin^3\frac{\varphi}{3}$.
\end{prob}

\begin{prob}[S]{2451}
$r = a\tanh\frac{\varphi}{2}$ ($0 \leqslant \varphi \leqslant 2\pi$).
\end{prob}

\begin{prob}[L]{2452}
\begin{enumerate}[label=(\alph*), itemsep=2pt, topsep=3pt]
\item $\varphi = \frac{1}{2}\left(r + \frac{1}{r}\right)$ ($1 \leqslant r \leqslant 3$);
\item $\varphi = \sqrt{r}$ ($0 \leqslant r \leqslant 5$);
\item $\varphi = \int_0^r \frac{\sinh\rho}{\rho}\,\mathrm{d}\rho$ ($0 \leqslant r \leqslant R$);
\item $r = 1 + \cos t$, $\varphi = t - \tan\frac{t}{2}$ ($0 \leqslant t \leqslant T < \pi$).
\end{enumerate}
\end{prob}

\begin{prob}[M]{2453}
证明: 椭圆
\[
x = a\cos t, \quad y = b\sin t
\]
的弧长等于正弦曲线 $y = c\sin\frac{x}{b}$ 的一波之长. 其中 $c = \sqrt{a^2 - b^2}$.
\end{prob}

\begin{prob}[M]{2454}
抛物线 $4ay = x^2$ 沿 $Ox$ 轴滚动. 证明: 抛物线的焦点的轨迹是悬链线.
\end{prob}

\begin{prob}[M]{2455}
求曲线
\[
y = \pm\left(\frac{1}{3} - x\right)\sqrt{x}
\]
的封闭部分与等周长圆周所分别围成的面积之比.
\end{prob}

\probsec{§7. 体积的计算法}

\begin{prob}[M]{2456}
求顶楼的体积, 其底是边长等于 $a$ 及 $b$ 的矩形, 其顶的棱边等于 $c$, 而高等于 $h$.
\end{prob}

\begin{prob}[M]{2457}
求截楔形的体积, 其平行的上下底为边长分别等于 $A, B$ 和 $a, b$ 的矩形, 而高等于 $h$.
\end{prob}

\begin{prob}[M]{2458}
求截锥体的体积, 其上下底为半轴长分别等于 $A, B$ 和 $a, b$ 的椭圆, 而高等于 $h$.
\end{prob}

\begin{prob}[S]{2459}
求旋转抛物体的体积, 其底为 $S$, 而高等于 $H$.
\end{prob}

\begin{prob}[M]{2460}
设一物体之垂直于 $Ox$ 轴的横截面的面积 $S = S(x)$ 依下面的二次式规律变化:
\[
S(x) = Ax^2 + Bx + C \quad (a \leqslant x \leqslant b),
\]
其中 $A, B$ 及 $C$ 为常数. 证明: 此物体之体积等于
\[
V = \frac{H}{6}\left[S(a) + 4S\left(\frac{a+b}{2}\right) + S(b)\right],
\]
其中 $H = b - a$ (辛普森公式).
\end{prob}

\begin{prob}[M]{2461}
物体是点 $M(x, y, z)$ 的集合, 其中 $0 \leqslant z \leqslant 1$, 而且当 $z$ 为有理数时, $0 \leqslant x \leqslant 1$, $0 \leqslant y \leqslant 1$; 当 $z$ 为无理数时, $-1 \leqslant x \leqslant 0$, $-1 \leqslant y \leqslant 0$. 证明: 此物体的体积不存在, 尽管相应积分
\[
\int_0^1 S(z)\,\mathrm{d}x = 1.
\]
\end{prob}

\probgroup{求下列曲面所围成的体积:}

\begin{prob}[S]{2462}
$\frac{x^2}{a^2} + \frac{y^2}{b^2} = 1$, $z = \frac{c}{a}x$, $z = 0$.
\end{prob}

\begin{prob}[S]{2463}
$\frac{x^2}{a^2} + \frac{y^2}{b^2} + \frac{z^2}{c^2} = 1$ (椭球面).
\end{prob}

\begin{prob}[S]{2464}
$\frac{x^2}{a^2} + \frac{y^2}{b^2} - \frac{z^2}{c^2} = 1$, $z = \pm c$.
\end{prob}

\begin{prob}[S]{2465}
$x^2 + z^2 = a^2$, $y^2 + z^2 = a^2$.
\end{prob}

\begin{prob}[S]{2466}
$x^2 + y^2 + z^2 = a^2$, $x^2 + y^2 = ax$.
\end{prob}

\begin{prob}[S]{2467}
$z^2 = b(a - x)$, $x^2 + y^2 = ax$.
\end{prob}

\begin{prob}[S]{2468}
$\frac{x^2}{a^2} + \frac{y^2}{z^2} = 1$ ($0 < z < a$).
\end{prob}

\begin{prob}[S]{2469}
$x + y + z^2 = 1$, $x = 0$, $y = 0$, $z = 0$.
\end{prob}

\begin{prob}[S]{2470}
$x^2 + y^2 + z^2 + xy + yz + zx = a^2$.
\end{prob}

\begin{prob}[M]{2471}
证明: 将平面图形
\[
a \leqslant x \leqslant b, \quad 0 \leqslant y \leqslant y(x)
\]
绕 $Oy$ 轴旋转所成的旋转体的体积等于
\[
V_y = 2\pi\int_a^b xy(x)\,\mathrm{d}x,
\]
这里 $y(x)$ 为单值连续函数.
\end{prob}

\probgroup{求下列曲线段旋转所成旋转体的体积:}

\begin{prob}[S]{2472}
$y = b\left(\frac{x}{a}\right)^{\frac{2}{3}}$ ($0 \leqslant x \leqslant a$) 绕 $Ox$ 轴 (半立方抛物线).
\end{prob}

\begin{prob}[M]{2473}
$y = 2x - x^2$, $y = 0$: (a) 绕 $Ox$ 轴; \qquad (b) 绕 $Oy$ 轴.
\end{prob}

\begin{prob}[M]{2474}
$y = \sin x$, $y = 0$ ($0 \leqslant x \leqslant \pi$): (a) 绕 $Ox$ 轴; \qquad (b) 绕 $Oy$ 轴.
\end{prob}

\begin{prob}[M]{2475}
$y = b\left(\frac{x}{a}\right)^2$, $y = b\left|\frac{x}{a}\right|$: (a) 绕 $Ox$ 轴; \qquad (b) 绕 $Oy$ 轴.
\end{prob}

\begin{prob}[M]{2476}
$y = \mathrm{e}^{-x}$, $y = 0$ ($0 \leqslant x < +\infty$): (a) 绕 $Ox$ 轴; \qquad (b) 绕 $Oy$ 轴.
\end{prob}

\begin{prob}[S]{2477}
$x^2 + (y - b)^2 = a^2$ ($0 < a \leqslant b$), 绕 $Ox$ 轴.
\end{prob}

\begin{prob}[S]{2478}
$x^2 - xy + y^2 = a^2$, 绕 $Ox$ 轴.
\end{prob}

\begin{prob}[S]{2479}
$y = \mathrm{e}^{-x}\sqrt{\sin x}$ ($0 \leqslant x < +\infty$), 绕 $Ox$ 轴.
\end{prob}

\begin{prob}[L]{2480}
$x = a(t - \sin t)$, $y = a(1 - \cos t)$ ($0 \leqslant t \leqslant 2\pi$), $y = 0$:
\begin{enumerate}[label=(\alph*), itemsep=2pt, topsep=3pt]
\item 绕 $Ox$ 轴;
\item 绕 $Oy$ 轴;
\item 绕直线 $y = 2a$.
\end{enumerate}
\end{prob}

\begin{prob}[M]{2481}
$x = a\sin^3 t$, $y = b\cos^3 t$ ($0 \leqslant t \leqslant 2\pi$): (a) 绕 $Ox$ 轴; \qquad (b) 绕 $Oy$ 轴.
\end{prob}

\begin{prob}[M]{2482.1}
求曲线
\[
x = 2t - t^2, \quad y = 4t - t^3
\]
所围图形绕 (a) $Ox$ 轴; (b) $Oy$ 轴旋转所成旋转体的体积.
\end{prob}

\begin{prob}[M]{2482.2}
证明: 把平面图形
\[
0 \leqslant \alpha \leqslant \varphi \leqslant \beta \leqslant \pi, \quad 0 \leqslant r \leqslant r(\varphi) \quad (\varphi\text{ 与 }r\text{ 为极坐标})
\]
绕极轴旋转所成旋转体的体积等于
\[
V = \frac{2\pi}{3}\int_\alpha^\beta r^3(\varphi)\sin\varphi\,\mathrm{d}\varphi.
\]
\end{prob}

\probgroup{求下列由极坐标或直角坐标给的平面图形经旋转后所成旋转体的体积:}

\begin{prob}[M]{2483.1}
$r = a(1 + \cos\varphi)$ ($0 \leqslant \varphi \leqslant 2\pi$): (a) 绕极轴; \qquad (b) 绕直线 $r\cos\varphi = -\frac{a}{4}$.
\end{prob}

\begin{prob}[L]{2483.2}
$(x^2 + y^2)^2 = a^2(x^2 - y^2)$:
\begin{enumerate}[label=(\alph*), itemsep=2pt, topsep=3pt]
\item 绕 $Ox$ 轴;
\item 绕 $Oy$ 轴;
\item 绕直线 $y = x$.
\end{enumerate}
\probhint{化为极坐标.}
\end{prob}

\begin{prob}[M]{2484.1}
求半圈阿基米德螺线
\[
r = a\varphi \quad (a > 0; 0 \leqslant \varphi \leqslant \pi)
\]
绕极轴旋转所成旋转体的体积.
\end{prob}

\begin{prob}[M]{2484.2}
求曲线
\[
\varphi = \pi r^3, \quad \varphi = \pi
\]
所围图形绕极轴旋转所成旋转体的体积.
\end{prob}

\begin{prob}[M]{2485}
求图形
\[
a \leqslant r \leqslant a\sqrt{2\sin 2\varphi}
\]
绕极轴旋转而成的旋转体的体积.
\end{prob}

\probsec{§8. 旋转曲面表面积的计算法}

\probgroup{求旋转下列曲线所成曲面的面积:}

\begin{prob}[S]{2486}
$y = x\sqrt{\frac{x}{a}}$ ($0 \leqslant x \leqslant a$), 绕 $Ox$ 轴.
\end{prob}

\begin{prob}[S]{2487}
$y = a\cos\frac{\pi x}{2b}$ ($|x| \leqslant b$), 绕 $Ox$ 轴.
\end{prob}

\begin{prob}[S]{2488}
$y = \tan x$ $\left(0 \leqslant x \leqslant \frac{\pi}{4}\right)$, 绕 $Ox$ 轴.
\end{prob}

\begin{prob}[M]{2489}
$y^2 = 2px$ ($0 \leqslant x \leqslant x_0$), (a) 绕 $Ox$ 轴; \qquad (b) 绕 $Oy$ 轴.
\end{prob}

\begin{prob}[M]{2490}
$\frac{x^2}{a^2} + \frac{y^2}{b^2} = 1$ ($0 < b \leqslant a$), (a) 绕 $Ox$ 轴; \qquad (b) 绕 $Oy$ 轴.
\end{prob}

\begin{prob}[S]{2491}
$x^2 + (y - b)^2 = a^2$ ($b \geqslant a$), 绕 $Ox$ 轴.
\end{prob}

\begin{prob}[S]{2492}
$x^{\frac{2}{3}} + y^{\frac{2}{3}} = a^{\frac{2}{3}}$, 绕 $Ox$ 轴.
\end{prob}

\begin{prob}[M]{2493}
$y = a\cosh\frac{x}{a}$ ($|x| \leqslant b$), (a) 绕 $Ox$ 轴; \qquad (b) 绕 $Oy$ 轴.
\end{prob}

\begin{prob}[S]{2494}
$\pm x = a\ln\frac{a + \sqrt{a^2 - y^2}}{y} - \sqrt{a^2 - y^2}$, 绕 $Ox$ 轴.
\end{prob}

\begin{prob}[L]{2495}
$x = a(t - \sin t)$, $y = a(1 - \cos t)$ ($0 \leqslant t \leqslant 2\pi$),
\begin{enumerate}[label=(\alph*), itemsep=2pt, topsep=3pt]
\item 绕 $Ox$ 轴;
\item 绕 $Oy$ 轴;
\item 绕直线 $y = 2a$.
\end{enumerate}
\end{prob}

\begin{prob}[S]{2496}
$x = a\cos^3 t$, $y = a\sin^3 t$, 绕直线 $y = x$.
\end{prob}

\begin{prob}[S]{2497}
$r = a(1 + \cos\varphi)$, 绕极轴.
\end{prob}

\begin{prob}[L]{2498}
$r^2 = a^2\cos 2\varphi$,
\begin{enumerate}[label=(\alph*), itemsep=2pt, topsep=3pt]
\item 绕极轴;
\item 绕轴 $\varphi = \frac{\pi}{2}$;
\item 绕轴 $\varphi = \frac{\pi}{4}$.
\end{enumerate}
\end{prob}

\begin{prob}[M]{2499}
由抛物线 $ay = a^2 - x^2$ 及 $Ox$ 轴围成的图形绕 $Ox$ 轴旋转而构成一旋转体, 求其表面积与等体积球的表面积之比.
\end{prob}

\begin{prob}[M]{2500}
由直线 $x = \frac{p}{2}$ 与抛物线 $y^2 = 2px$ 围成的图形绕直线 $y = p$ 旋转而构成一旋转体, 求其体积和表面积.
\end{prob}

\probsec{§9. 矩的计算法. 质心的坐标}

\begin{prob}[M]{2501.1}
求半径为 $a$ 的半圆弧对过此弧两端点的直径的静矩和转动惯量.
\end{prob}

\begin{prob}[M]{2501.2}
求抛物线弧 $y^2 = 2px \left(0 \leqslant x \leqslant \dfrac{p}{2}\right)$ 对直线 $x = \dfrac{p}{2}$ 的静矩.
\end{prob}

\begin{prob}[M]{2502.1}
求底为 $b$, 高为 $h$ 的均质三角形平板对其底边的静矩和转动惯量 ($\rho = 1$).
\end{prob}

\begin{prob}[M]{2502.2}
求曲线
\[
ay = 2ax - x^2 \quad (a > 0), \quad y = 0
\]
所围抛物线弓形对 $Ox$ 轴和 $Oy$ 轴的转动惯量 $I_x = M_2^{(x)}$ 和 $I_y = M_2^{(y)}$. 由关系式 $I_x = S r_x^2$, $I_y = S r_y^2$ ($S$ 为弓形面积) 定义的惯性半径 $r_x$ 和 $r_y$ 等于什么？
\end{prob}

\begin{prob}[M]{2503}
求半轴长为 $a$ 和 $b$ 的均质椭圆形平板对其主轴的转动惯量 ($\rho = 1$).
\end{prob}

\begin{prob}[M]{2504.1}
求底半径为 $r$ 高为 $h$ 的均质圆锥对其底平面的静矩和转动惯量 ($\rho = 1$).
\end{prob}

\begin{prob}[M]{2504.2}
求半径为 $R$ 质量为 $M$ 的均质球体对其直径的转动惯量.
\end{prob}

\begin{prob}[M]{2505}
证明古尔丹第一定理：平面曲线弧 $C$ 绕此弧所在平面上不与它相交的轴旋转而成的旋转面，其面积等于此弧的长度与它的质心所画出的圆周之长的乘积.
\end{prob}

\begin{prob}[M]{2506}
证明古尔丹第二定理：平面图形 $S$ 绕此图形所在平面上不与它相交的轴旋转而成的旋转体，其体积等于图形 $S$ 的面积与此图形的质心所画出的圆周之长的乘积.
\end{prob}

\begin{prob}[M]{2507}
求圆弧 $x = a\cos\varphi, y = a\sin\varphi \ (|\varphi| \leqslant \alpha \leqslant \pi)$ 质心的坐标.
\end{prob}

\begin{prob}[M]{2508}
求抛物线 $ax = y^2, ay = x^2 \ (a > 0)$ 所围图形的质心的坐标.
\end{prob}

\begin{prob}[M]{2509}
求图形
\[
\frac{x^2}{a^2} + \frac{y^2}{b^2} \leqslant 1 \quad (0 \leqslant x \leqslant a, 0 \leqslant y \leqslant b)
\]
的质心的坐标.
\end{prob}

\begin{prob}[M]{2510}
求半径为 $a$ 的均质半球的质心坐标.
\end{prob}

\begin{prob}[M]{2511}
求对数螺线
\[
r = a\mathrm{e}^{m\varphi} \quad (m > 0)
\]
上由点 $O(-\infty, 0)$ 到点 $P(\varphi, r)$ 的弧 $OP$ 的质心 $C(\varphi_0, r_0)$ 之坐标. 当点 $P$ 移动时，点 $C$ 画出怎样的曲线？
\end{prob}

\begin{prob}[M]{2512}
求曲线 $r = a(1 + \cos\varphi)$ 所围图形的质心坐标.
\end{prob}

\begin{prob}[M]{2513}
求摆线 $x = a(t - \sin t), y = a(1 - \cos t) \ (0 \leqslant t \leqslant 2\pi)$ 的第一拱与 $Ox$ 轴所围图形的质心的坐标.
\end{prob}

\begin{prob}[M]{2514}
求图形 $0 \leqslant x \leqslant a; y^2 \leqslant 2px$ 绕 $Ox$ 轴旋转所成旋转体的质心的坐标.
\end{prob}

\begin{prob}[M]{2515}
求半球 $x^2 + y^2 + z^2 = a^2 \ (z \geqslant 0)$ 的质心的坐标.
\end{prob}

\probsec{§10. 力学和物理学中的问题}

\probgroup{组成适当的积分并求出其极限，以便求解下列问题：}

\begin{prob}[M]{2516}
杆的长度 $l = 10\,\mathrm{m}$, 若该杆的线密度按规律 $\delta = 6 + 0.3x\,\mathrm{kg/m}$ 而变，其中 $x$ 为到杆的一个端点的距离，求杆的质量.
\end{prob}

\begin{prob}[M]{2517}
把质量为 $m$ 的物体从地球 (其半径为 $R$) 表面抬升到高度为 $h$ 的地方，需要对它作多少功？若物体远离至无穷远处，则功等于多少？
\end{prob}

\begin{prob}[M]{2518}
若 $10\,\mathrm{N}$ 的力能使弹簧伸长 $1\,\mathrm{cm}$, 现在要使这弹簧伸长 $10\,\mathrm{cm}$, 问要作多少功？
\probhint{利用胡克定律。}
\end{prob}

\begin{prob}[M]{2519}
直径为 $20\,\mathrm{cm}$, 长为 $80\,\mathrm{cm}$ 的圆柱形汽缸充满压强为 $100\,\mathrm{N/cm}^2$ 的蒸汽. 假定蒸汽的温度保持不变，要使其体积减小一半，需要作多少功？
\end{prob}

\begin{prob}[M]{2520}
求水对竖直放置的半圆形档板的压力，该档板的半径为 $a$, 而水面位于档板顶部直径的位置.\footnote{此题是流体静力学问题，在实际求解时还应知道水面上的压强 $p_0$, 以及水的密度 $\rho$. 原书所给答案对应着 $p_0 = 0$ 的特例. 类似情形还出现在习题 2521, 4070, 4071, 4072. —— 译注}
\end{prob}

\begin{prob}[M]{2521}
求水对竖直壁面的压力，此壁面的形状为梯形，其下底 $a = 10\,\mathrm{m}$, 上底 $b = 6\,\mathrm{m}$, 高 $h = 5\,\mathrm{m}$, 下底位于水面以下 $c = 20\,\mathrm{m}$ 的深度.\footnote{此题应考虑水面处大气压 $p_0$ 的影响. 书后的答案对应 $p_0 = 0$ 的情形，与实际结果有较大差别. —— 译注}
\end{prob}

\probgroup{写出微分方程并求解下列问题：}

\begin{prob}[M]{2522}
一质点的速度按规律
\[
v = v_0 + at \quad (v_0 = \text{常数}, \ a = \text{常数})
\]
变化，问在时间间隔 $[0, T]$ 内此质点经过怎样的路程？
\end{prob}

\begin{prob}[M]{2523}
半径为 $R$ 而密度为 $\delta$ 的均质球体以角速度 $\omega$ 绕其直径旋转. 求此球的动能.
\end{prob}

\begin{prob}[M]{2524}
线密度 $\mu_0$ 为常数的无穷直线以怎样的力吸引距此直线距离为 $a$ 质量为 $m$ 的质点？
\end{prob}

\begin{prob}[M]{2525}
计算半径为 $a$ 且面密度 $\delta_0$ 为常数的圆形薄板以怎样的力吸引质量为 $m$ 的质点 $P$, 此质点位于通过薄板中心 $Q$ 且垂直于薄板平面的直线上，距离 $PQ$ 等于 $b$.
\end{prob}

\begin{prob}[M]{2526}
根据托里拆利定律，液体通过小孔从容器中流出的速度等于
\[
v = c\sqrt{2gh},
\]
式中 $g$ 为重力加速度，$h$ 为液体表面在小孔以上的高度，$c = 0.6$ 为实验系数. 直径为 $D = 1\,\mathrm{m}$ 高为 $H = 2\,\mathrm{m}$ 的直立圆柱形大桶，其底部有一个直径为 $d = 1\,\mathrm{cm}$ 的圆孔，问此桶充满液体后经过多长时间方可完全排空？
\end{prob}

\begin{prob}[M]{2527}
旋转体容器应当具有什么形状，才能使液体从容器底部流出时，液体表面的下降是均匀的？
\end{prob}

\begin{prob}[M]{2528}
镭在每一时刻的衰变速度与其现存的量成正比. 设镭的量在初始时刻 $t = 0$ 为 $Q_0$, 经过时间 $T = 1600$ 年它的量减少了一半. 求镭的衰变规律.
\end{prob}

\begin{prob}[M]{2529}
在一种把物质 $A$ 变为物质 $B$ 的二阶化学反应中，反应速度与此二物质浓度之积成正比. 若在 $t = 0$ 时在容器中有 $20\%$ 的物质 $B$, 而当 $t = 15\,\mathrm{min}$ 时其浓度变为 $80\%$, 问在 $t = 1\,\mathrm{h}$ 时其浓度如何？
\end{prob}

\begin{prob}[M]{2530}
根据胡克定律，杆的相对伸长 $\varepsilon$ 与相应横截面上的应力 $\sigma$ 成正比，即
\[
\varepsilon = \frac{\sigma}{E},
\]
式中 $E$ 为杨氏模量. 求圆锥形杆在自重作用下的伸长量，此锥形的顶向下而底固定，设底半径等于 $R$, 圆锥的高为 $H$, 密度为 $\rho$.
\end{prob}

\probsec{§11. 定积分的近似计算法}

\begin{prob}[M]{2531}
利用矩形公式 ($n = 12$), 近似地计算
\[
\int_0^{2\pi} x\sin x\,\mathrm{d}x,
\]
并把结果同精确答案进行比较.
\end{prob}

\probgroup{利用梯形公式计算下列积分并估计它们的误差：}

\begin{prob}[S]{2532}
$\displaystyle \int_0^1 \frac{\mathrm{d}x}{1+x} \ (n = 8)$.
\end{prob}

\begin{prob}[S]{2533}
$\displaystyle \int_0^1 \frac{\mathrm{d}x}{1+x^3} \ (n = 12)$.
\end{prob}

\begin{prob}[S]{2534}
$\displaystyle \int_0^{\frac{\pi}{2}} \sqrt{1 - \frac{1}{4}\sin^2 x}\,\mathrm{d}x \ (n = 6)$.
\end{prob}

\probgroup{利用辛普森公式计算下列积分：}

\begin{prob}[S]{2535}
$\displaystyle \int_1^9 \sqrt{x}\,\mathrm{d}x \ (n = 4)$.
\end{prob}

\begin{prob}[S]{2536}
$\displaystyle \int_0^{\pi} \sqrt{3 + \cos x}\,\mathrm{d}x \ (n = 6)$.
\end{prob}

\begin{prob}[S]{2537}
$\displaystyle \int_0^{\frac{\pi}{2}} \frac{\sin x}{x}\,\mathrm{d}x \ (n = 10)$.
\end{prob}

\begin{prob}[S]{2538}
$\displaystyle \int_0^1 \frac{x\,\mathrm{d}x}{\ln(1+x)} \ (n = 6)$.
\end{prob}

\begin{prob}[M]{2539}
取 $n = 10$, 计算卡塔兰常数
\[
G = \int_0^1 \frac{\arctan x}{x}\,\mathrm{d}x.
\]
\end{prob}

\begin{prob}[M]{2540}
利用公式
\[
\frac{\pi}{4} = \int_0^1 \frac{\mathrm{d}x}{1+x^2}
\]
计算数 $\pi$, 精确到 $10^{-5}$.
\end{prob}

\begin{prob}[M]{2541}
精确到 $0.001$, 计算
\[
\int_0^1 \mathrm{e}^{x^2}\,\mathrm{d}x.
\]
\end{prob}

\begin{prob}[M]{2542}
精确到 $10^{-4}$, 计算
\[
\int_0^1 (\mathrm{e}^x - 1)\ln\frac{1}{x}\,\mathrm{d}x.
\]
\end{prob}

\begin{prob}[M]{2543}
精确到 $0.001$, 计算概率积分
\[
\int_0^{+\infty} \mathrm{e}^{-x^2}\,\mathrm{d}x.
\]
\end{prob}

\begin{prob}[M]{2544}
近似地求出半轴为 $a = 10$ 及 $b = 6$ 的椭圆的周长.
\end{prob}

\begin{prob}[M]{2545}
取 $\Delta x = \dfrac{\pi}{3}$, 描点作出函数
\[
y = \int_0^x \frac{\sin t}{t}\,\mathrm{d}t \quad (0 \leqslant x \leqslant 2\pi)
\]
的图像.
\end{prob}

\probchapter{第五章\quad 级数}

\probsec{§1. 数项级数. 同号级数收敛性的判别法}

\probgroup{直接证明下列级数的收敛性并求它们的和：}

\begin{prob}[S]{2546}
$1 - \dfrac{1}{2} + \dfrac{1}{4} - \dfrac{1}{8} + \cdots + \dfrac{(-1)^{n-1}}{2^{n-1}} + \cdots$.
\end{prob}

\begin{prob}[S]{2547}
$\left(\dfrac{1}{2} + \dfrac{1}{3}\right) + \left(\dfrac{1}{2^2} + \dfrac{1}{3^2}\right) + \cdots + \left(\dfrac{1}{2^n} + \dfrac{1}{3^n}\right) + \cdots$.
\end{prob}

\begin{prob}[S]{2548}
$\dfrac{1}{2} + \dfrac{3}{2^2} + \dfrac{5}{2^3} + \cdots + \dfrac{2n - 1}{2^n} + \cdots$.
\end{prob}

\begin{prob}[S]{2549}
$\dfrac{1}{1 \cdot 2} + \dfrac{1}{2 \cdot 3} + \dfrac{1}{3 \cdot 4} + \cdots + \dfrac{1}{n(n+1)} + \cdots$.
\end{prob}

\begin{prob}[S]{2550}
$\dfrac{1}{1 \cdot 4} + \dfrac{1}{4 \cdot 7} + \cdots + \dfrac{1}{(3n - 2)(3n + 1)} + \cdots$.
\end{prob}

\begin{prob}[M]{2551}
\begin{enumerate}[label=(\alph*), itemsep=2pt, topsep=3pt]
  \item $q\sin\alpha + q^2\sin 2\alpha + \cdots + q^n\sin n\alpha + \cdots$ ($|q| < 1$);
  \item $q\cos\alpha + q^2\cos 2\alpha + \cdots + q^n\cos n\alpha + \cdots$ ($|q| < 1$).
\end{enumerate}
\end{prob}

\begin{prob}[S]{2552}
$\sum_{n=1}^{\infty} (\sqrt{n+2} - 2\sqrt{n+1} + \sqrt{n})$.
\end{prob}

\begin{prob}[M]{2553}
研究级数 $\sum_{n=1}^{\infty} \sin nx$ 的收敛性.
\probhint{证明：当 $x \neq k\pi$ ($k$ 是整数) 时，要 $\sin nx \to 0$ ($n \to \infty$) 是不可能的.}
\end{prob}

\begin{prob}[M]{2554}
证明：若级数 $\sum_{n=1}^{\infty} a_n$ 收敛，则把该级数的各项在不变更其先后次序的情况下分别组合起来，所得级数
\[
\sum_{n=1}^{\infty} A_n \quad \left(\text{其中 } A_n = \sum_{i=p_n}^{p_{n+1}-1} a_i \ (p_1 = 1, p_1 < p_2 < \cdots)\right),
\]
也收敛且有相同的和. 反之不真. 举出例子.
\end{prob}

\begin{prob}[M]{2555}
证明：若级数 $\sum_{n=1}^{\infty} a_n$ 的各项是正的，且把这级数的各项分别组合而得到的级数 $\sum_{n=1}^{\infty} A_n$ 收敛，则原来的级数 $\sum_{n=1}^{\infty} a_n$ 也收敛.
\end{prob}

\probgroup{研究下列级数的收敛性：}

\begin{prob}[S]{2556}
$1 - 1 + 1 - 1 + 1 - 1 + \cdots$.
\end{prob}

\begin{prob}[S]{2557}
$0.001 + \sqrt{0.001} + \sqrt[3]{0.001} + \cdots$.
\end{prob}

\begin{prob}[S]{2558}
$\dfrac{1}{1!} + \dfrac{1}{2!} + \dfrac{1}{3!} + \cdots + \dfrac{1}{n!} + \cdots$.
\end{prob}

\begin{prob}[S]{2559}
$1 + \dfrac{1}{3} + \dfrac{1}{5} + \dfrac{1}{7} + \cdots + \dfrac{1}{2n - 1} + \cdots$.
\end{prob}

\begin{prob}[S]{2560}
$\dfrac{1}{1001} + \dfrac{1}{2001} + \dfrac{1}{3001} + \cdots + \dfrac{1}{1000n + 1} + \cdots$.
\end{prob}

\begin{prob}[S]{2561}
$1 + \dfrac{2}{3} + \dfrac{3}{5} + \cdots + \dfrac{n}{2n - 1} + \cdots$.
\end{prob}

\begin{prob}[S]{2562}
$1 + \dfrac{1}{3^2} + \dfrac{1}{5^2} + \cdots + \dfrac{1}{(2n - 1)^2} + \cdots$.
\end{prob}

\begin{prob}[S]{2563}
$\dfrac{1}{\sqrt{2}} + \dfrac{1}{2\sqrt{3}} + \dfrac{1}{3\sqrt{4}} + \cdots + \dfrac{1}{n\sqrt{n+1}} + \cdots$.
\end{prob}

\begin{prob}[S]{2564}
$\dfrac{1}{\sqrt{1 \cdot 3}} + \dfrac{1}{\sqrt{3 \cdot 5}} + \cdots + \dfrac{1}{\sqrt{(2n - 1)(2n + 1)}} + \cdots$.
\end{prob}

\begin{prob}[M]{2565}
证明：由等差数列各项的倒数组成的级数是发散的.
\end{prob}

\begin{prob}[M]{2566}
证明：若级数 $\sum_{n=1}^{\infty} a_n$ (A) 及 $\sum_{n=1}^{\infty} b_n$ (B) 皆收敛且 $a_n \leqslant c_n \leqslant b_n$ ($n = 1, 2, \cdots$)，则级数 $\sum_{n=1}^{\infty} c_n$ (C) 也收敛. 若级数 (A) 和 (B) 皆发散，问级数 (C) 的收敛性若何？
\end{prob}

\begin{prob}[M]{2567}
设已知二发散级数
\[
\sum_{n=1}^{\infty} a_n \quad \text{及} \quad \sum_{n=1}^{\infty} b_n
\]
的各项不为负数，问下列级数的收敛性若何：
(a) $\sum_{n=1}^{\infty} \min(a_n, b_n)$ 及 (b) $\sum_{n=1}^{\infty} \max(a_n, b_n)$？
\end{prob}

\begin{prob}[M]{2568}
证明：若级数 $\sum_{n=1}^{\infty} a_n$ ($a_n \geqslant 0$) 收敛，则级数 $\sum_{n=1}^{\infty} a_n^2$ 也收敛. 逆命题不成立，举出例子.
\end{prob}

\begin{prob}[M]{2569}
证明：若级数 $\sum_{n=1}^{\infty} a_n^2$ 及 $\sum_{n=1}^{\infty} b_n^2$ 收敛，则级数
\[
\sum_{n=1}^{\infty} |a_n b_n|, \quad \sum_{n=1}^{\infty} (a_n + b_n)^2, \quad \sum_{n=1}^{\infty} \frac{|a_n|}{n}
\]
也收敛.
\end{prob}

\begin{prob}[M]{2570}
证明：若
\[
\lim_{n \to \infty} n a_n = a \neq 0,
\]
则级数 $\sum_{n=1}^{\infty} a_n$ 发散.
\end{prob}

\begin{prob}[M]{2571}
证明：若各项为正且其值单调递减的级数 $\sum_{n=1}^{\infty} a_n$ 收敛，则
\[
\lim_{n \to \infty} n a_n = 0.
\]
\end{prob}

\begin{prob}[M]{2572}
若当 $p = 1, 2, 3, \cdots$ 时 $\lim_{n \to \infty} (a_{n+1} + a_{n+2} + \cdots + a_{n+p}) = 0$，问级数 $\sum_{n=1}^{\infty} a_n$ 是否收敛？
\end{prob}

\probgroup{利用柯西准则，证明下列正项级数的收敛性：}

\begin{prob}[S]{2573}
$a_0 + \dfrac{a_1}{10} + \cdots + \dfrac{a_n}{10^n} + \cdots$ ($|a_n| < 10$).
\end{prob}

\begin{prob}[S]{2574}
$\dfrac{\sin x}{2} + \dfrac{\sin 2x}{2^2} + \cdots + \dfrac{\sin nx}{2^n} + \cdots$.
\end{prob}

\begin{prob}[S]{2575.1}
$\dfrac{\cos x - \cos 2x}{1} + \dfrac{\cos 2x - \cos 3x}{2} + \cdots + \dfrac{\cos nx - \cos(n+1)x}{n} + \cdots$.
\end{prob}

\begin{prob}[M]{2575.2}
$\dfrac{\cos x}{1^2} + \dfrac{\cos x^2}{2^2} + \cdots + \dfrac{\cos x^n}{n^2} + \cdots$.
\probhint{利用不等式
\[
\frac{1}{n^2} < \frac{1}{n(n-1)} = \frac{1}{n-1} - \frac{1}{n} \quad (n = 2, 3, \cdots).
\]}
\end{prob}

\probgroup{利用柯西准则，证明下列级数的发散性：}

\begin{prob}[S]{2576}
$1 + \dfrac{1}{2} + \dfrac{1}{3} + \cdots + \dfrac{1}{n} + \cdots$.
\end{prob}

\begin{prob}[M]{2577}
\begin{enumerate}[label=(\alph*), itemsep=2pt, topsep=3pt]
  \item $1 + \dfrac{1}{2} - \dfrac{1}{3} + \dfrac{1}{4} + \dfrac{1}{5} - \dfrac{1}{6} + \cdots$;
  \item $\dfrac{1}{\sqrt{1 \cdot 2}} + \dfrac{1}{\sqrt{2 \cdot 3}} + \cdots + \dfrac{1}{\sqrt{n(n+1)}} + \cdots$.
\end{enumerate}
\end{prob}

\probgroup{运用比较判别法、达朗贝尔判别法或柯西判别法，研究下列级数的收敛性：}

\begin{prob}[S]{2578}
$\dfrac{1000}{1!} + \dfrac{1000^2}{2!} + \dfrac{1000^3}{3!} + \cdots + \dfrac{1000^n}{n!} + \cdots$.
\end{prob}

\begin{prob}[S]{2579}
$\dfrac{(1!)^2}{2!} + \dfrac{(2!)^2}{4!} + \cdots + \dfrac{(n!)^2}{(2n)!} + \cdots$.
\end{prob}

\begin{prob}[S]{2580}
$\dfrac{1!}{1} + \dfrac{2!}{2^2} + \dfrac{3!}{3^3} + \cdots + \dfrac{n!}{n^n} + \cdots$.
\end{prob}

\begin{prob}[M]{2581}
\begin{enumerate}[label=(\alph*), itemsep=2pt, topsep=3pt]
  \item $\dfrac{2 \cdot 1!}{1} + \dfrac{2^2 \cdot 2!}{2^2} + \dfrac{2^3 \cdot 3!}{3^3} + \cdots + \dfrac{2^n n!}{n^n} + \cdots$;
  \item $\dfrac{3 \cdot 1!}{1} + \dfrac{3^2 \cdot 2!}{2^2} + \dfrac{3^3 \cdot 3!}{3^3} + \cdots + \dfrac{3^n n!}{n^n} + \cdots$.
\end{enumerate}
\end{prob}

\begin{prob}[S]{2582}
$\dfrac{(1!)^2}{2} + \dfrac{(2!)^2}{2^4} + \dfrac{(3!)^2}{2^9} + \cdots + \dfrac{(n!)^2}{2^{n^2}} + \cdots$.
\end{prob}

\begin{prob}[S]{2583}
$\dfrac{1000}{1} + \dfrac{1000 \cdot 1001}{1 \cdot 3} + \dfrac{1000 \cdot 1001 \cdot 1002}{1 \cdot 3 \cdot 5} + \cdots$.
\end{prob}

\begin{prob}[S]{2584}
$\dfrac{4}{2} + \dfrac{4 \cdot 7}{2 \cdot 6} + \dfrac{4 \cdot 7 \cdot 10}{2 \cdot 6 \cdot 10} + \cdots$.
\end{prob}

\begin{prob}[L]{2585}
\begin{enumerate}[label=(\alph*), itemsep=2pt, topsep=3pt]
  \item $\sum_{n=1}^{\infty} (\sqrt{2} - \sqrt[3]{2})(\sqrt{2} - \sqrt[5]{2}) \cdots (\sqrt{2} - \sqrt[2n+1]{2})$;
  \item $\sum_{n=1}^{\infty} a_n$，其中 $a_n = \begin{cases} \dfrac{1}{n}, & n = m^2, \\ \dfrac{1}{n^2}, & n \neq m^2 \end{cases}$ ($m$ 为正整数);
  \item $\sum_{n=1}^{\infty} nx \prod_{k=1}^{n} \dfrac{\sin^2 k\alpha}{1 + x^2 + \cos^2 k\alpha}$.
\end{enumerate}
\end{prob}

\begin{prob}[S]{2586}
$\sum_{n=1}^{\infty} \dfrac{n^2}{\left(2 + \dfrac{1}{n}\right)^n}$.
\end{prob}

\begin{prob}[S]{2587}
$\sum_{n=1}^{\infty} \dfrac{n^{n + \frac{1}{n}}}{\left(n + \dfrac{1}{n}\right)^n}$.
\end{prob}

\begin{prob}[S]{2588}
$\sum_{n=1}^{\infty} \dfrac{1}{\sqrt[n]{\ln n}}$.
\end{prob}

\begin{prob}[L]{2589}
\begin{enumerate}[label=(\alph*), itemsep=2pt, topsep=3pt]
  \item $\sum_{n=1}^{\infty} \dfrac{n^{n-1}}{(2n^2 + n + 1)^{\frac{n+1}{2}}}$;
  \item $\sum_{n=1}^{\infty} \dfrac{n^5}{2^n + 3^n}$;
  \item $\sum_{n=2}^{\infty} \left(\dfrac{n - 1}{n + 1}\right)^{n(n-1)}$;
  \item $\sqrt{2} + \sqrt{2 - \sqrt{2}} + \sqrt{2 - \sqrt{2 + \sqrt{2}}} + \sqrt{2 - \sqrt{2 + \sqrt{2 + \sqrt{2}}}} + \cdots$.
\end{enumerate}
\probhint{$\sqrt{2} = 2\cos\dfrac{\pi}{4}$.}
\end{prob}

\begin{prob}[M]{2590}
证明：若
\[
\lim_{n \to \infty} \frac{a_{n+1}}{a_n} = q \quad (a_n > 0),
\]
则 $a_n = o(q_1^n)$，其中 $q_1 > q$.
\end{prob}

\begin{prob}[M]{2591.1}
设正项级数 $\sum_{n=1}^{\infty} a_n$ ($a_n > 0$) 的各项在 $n \geqslant n_0$ 时满足不等式
\[
\frac{a_{n+1}}{a_n} \leqslant \rho < 1.
\]
证明：对于级数的余项
\[
R_n = a_{n+1} + a_{n+2} + \cdots
\]
在 $n \geqslant n_0$ 时有以下估计：
\[
R_n \leqslant a_{n_0} \cdot \frac{\rho^{n - n_0 + 1}}{1 - \rho}.
\]
\end{prob}

\begin{prob}[M]{2591.2}
已知级数
\[
\sum_{n=1}^{\infty} \frac{[(2n)!!]^2}{(4n)!!},
\]
其中 $(2n)!! = 2 \cdot 4 \cdots 2n$. 为使此级数的和与部分和 $S_n$ 之差小于 $\varepsilon = 10^{-6}$, 应取多少项方可?
\end{prob}

\begin{prob}[M]{2592}
证明: 若
\[
\overline{\lim}_{n\to\infty} \frac{a_{n+1}}{a_n} = q < 1 \quad (a_n > 0),
\]
则级数 $\sum_{n=1}^{\infty} a_n$ 收敛.

逆命题不成立. 研究例子
\[
\frac{1}{2} + \frac{1}{3} + \frac{1}{2^2} + \frac{1}{3^2} + \frac{1}{2^3} + \frac{1}{3^3} + \cdots.
\]
\end{prob}

\begin{prob}[M]{2593}
证明: 若对于级数 $\sum_{n=1}^{\infty} a_n$ ($a_n > 0$) 存在极限
\[
\lim_{n\to\infty} \frac{a_{n+1}}{a_n} = q, \tag{A}
\]
则
\[
\lim_{n\to\infty} \sqrt[n]{a_n} = q \tag{B}
\]
也存在.

逆命题不成立: 若极限 (B) 存在, 则极限 (A) 可以不存在. 研究例子
\[
\sum_{n=1}^{\infty} \frac{3 + (-1)^n}{2^{n+1}}.
\]
\end{prob}

\begin{prob}[M]{2594}
证明: 若
\[
\overline{\lim}_{n\to\infty} \sqrt[n]{a_n} = q \quad (a_n \geqslant 0),
\]
则 (a) 当 $q < 1$ 时级数 $\sum_{n=1}^{\infty} a_n$ 收敛; (b) 当 $q > 1$ 时该级数发散 (柯西判别法的推广).
\end{prob}

\probgroup{研究下列级数的收敛性：}

\begin{prob}[S]{2595}
$\sum_{n=1}^{\infty} \frac{2 + (-1)^n}{2^n}$.
\end{prob}

\begin{prob}[S]{2596}
$\sum_{n=1}^{\infty} \frac{n \cos^2 \frac{n\pi}{3}}{2^n}$.
\end{prob}

\begin{prob}[M]{2597}
(a) $\sum_{n=1}^{\infty} \frac{n^3 [\sqrt{2} + (-1)^n]^n}{3^n}$; \qquad (b) $\sum_{n=1}^{\infty} \left(\frac{1 + \cos n}{2 + \cos n}\right)^{2n - \ln n}$.
\end{prob}

\probgroup{利用拉比判别法和高斯判别法，研究下列级数的收敛性：}

\begin{prob}[S]{2598}
$\left(\frac{1}{2}\right)^p + \left(\frac{1 \cdot 3}{2 \cdot 4}\right)^p + \left(\frac{1 \cdot 3 \cdot 5}{2 \cdot 4 \cdot 6}\right)^p + \cdots$.
\end{prob}

\begin{prob}[S]{2599}
$\frac{a}{b} + \frac{a(a+d)}{b(b+d)} + \frac{a(a+d)(a+2d)}{b(b+d)(b+2d)} + \cdots$ ($a > 0, b > 0, d > 0$).
\end{prob}

\begin{prob}[S]{2600}
$\sum_{n=1}^{\infty} \frac{n!\mathrm{e}^n}{n^{n+p}}$.
\end{prob}

\begin{prob}[S]{2601}
$\sum_{n=1}^{\infty} \frac{\sqrt{n!}}{(2+\sqrt{1})(2+\sqrt{2})\cdots(2+\sqrt{n})}$.
\end{prob}

\begin{prob}[S]{2602}
$\sum_{n=1}^{\infty} \frac{n! n^{-p}}{q(q+1)\cdots(q+n)}$ ($p > 0, q > 0$).
\end{prob}

\begin{prob}[S]{2603}
$\sum_{n=1}^{\infty} \frac{p(p+1)\cdots(p+n-1)}{n!} \cdot \frac{1}{n^q}$ ($p > 0, q > 0$).
\end{prob}

\begin{prob}[S]{2604}
$\sum_{n=1}^{\infty} \left[\frac{1 \cdot 3 \cdot 5 \cdots (2n-1)}{2 \cdot 4 \cdot 6 \cdots (2n)}\right]^p \cdot \frac{1}{n^q}$ ($p > 0, q > 0$).
\end{prob}

\begin{prob}[S]{2605}
$\sum_{n=1}^{\infty} \left[\frac{p(p+1)\cdots(p+n-1)}{q(q+1)\cdots(q+n-1)}\right]^{\alpha}$ ($p > 0, q > 0$).
\end{prob}

\begin{prob}[M]{2606}
证明: 对于正项级数 $\sum_{n=1}^{\infty} a_n$ ($a_n > 0$), 若在 $n \to \infty$ 时成立条件
\[
\frac{a_n}{a_{n+1}} = 1 + \frac{p}{n} + o\left(\frac{1}{n}\right),
\]
则
\[
a_n = o\left(\frac{1}{n^{p-\varepsilon}}\right),
\]
其中 $\varepsilon > 0$ 是任意小量; 并且若 $p > 0$, 则当 $n \to \infty$ 时 $a_n \downarrow 0$, 即从 $n \geqslant n_0$ 开始, $a_n$ 单调递减, 且在 $n \to \infty$ 时 $a_n \to 0$.
\end{prob}

\probgroup{求出通项 $a_n$ 的减小的阶，从而研究级数 $\sum_{n=1}^{\infty} a_n$ 的收敛性：}

\begin{prob}[S]{2607}
$a_n = \frac{n^p + a_1 n^{p-1} + \cdots + a_p}{n^q + b_1 n^{q-1} + \cdots + b_q}$, 其中 $n^q + b_1 n^{q-1} + \cdots + b_q > 0$.
\end{prob}

\begin{prob}[S]{2608}
$a_n = \frac{1}{n^p} \sin \frac{\pi}{n}$.
\end{prob}

\begin{prob}[S]{2609}
$a_n = (\sqrt{n+1} - \sqrt{n})^p \ln \frac{n-1}{n+1}$ ($n > 1$).
\end{prob}

\begin{prob}[S]{2610}
$a_n = \ln^p \left(\sec \frac{\pi}{n}\right)$.
\end{prob}

\begin{prob}[S]{2611}
$a_n = \log_{b^n} \left(1 + \frac{\sqrt[n]{a}}{n}\right)$ ($a > 0, b > 0$).
\end{prob}

\begin{prob}[S]{2612}
$a_n = \left[\mathrm{e} - \left(1 + \frac{1}{n}\right)^n\right]^p$.
\end{prob}

\begin{prob}[M]{2613}
(a) $a_n = \frac{1}{n^{1+\frac{k}{\ln n}}}$; \qquad (b) $a_n = \frac{1}{n^{1+\frac{1}{n}}}$.
\end{prob}

\begin{prob}[M]{2614}
证明热梅判别法: 对于正项级数 $\sum_{n=1}^{\infty} a_n$ ($a_n > 0$), 若当 $n > n_0$ 时
\[
(1 - \sqrt[n]{a_n}) \frac{n}{\ln n} \geqslant p > 1,
\]
则此级数收敛; 若当 $n > n_0$ 时
\[
(1 - \sqrt[n]{a_n}) \frac{n}{\ln n} \leqslant 1,
\]
则此级数发散.
\end{prob}

\begin{prob}[M]{2615}
证明: 若存在 $\alpha > 0$ 使当 $n \geqslant n_0$ 时 $\frac{\ln \frac{1}{a_n}}{\ln n} \geqslant 1 + \alpha$ ($a_n > 0$), 则级数 $\sum_{n=1}^{\infty} a_n$ ($a_n > 0$) 收敛; 若 $n \geqslant n_0$ 时 $\frac{\ln \frac{1}{a_n}}{\ln n} \leqslant 1$, 则此级数发散 (对数差别法).
\end{prob}

\probgroup{研究具有如下通项的级数的收敛性：}

\begin{prob}[S]{2616}
$a_n = n^{\ln x}$ ($x > 0$).
\end{prob}

\begin{prob}[S]{2617}
$a_n = \frac{1}{(\ln \ln n)^{\ln n}}$ ($n > 1$).
\end{prob}

\begin{prob}[S]{2618}
$a_n = \frac{1}{(\ln n)^{\ln \ln n}}$ ($n > 1$).
\end{prob}

\begin{prob}[M]{2619}
利用柯西积分判别法, 研究具有如下通项的级数的收敛性:

(a) $a_n = \frac{1}{n \ln^p n}$ ($n > 1$); \qquad (b) $a_n = \frac{1}{n (\ln n)^p (\ln \ln n)^q}$ ($n > 2$).
\end{prob}

\begin{prob}[L]{2620}
研究下列级数的收敛性:
\begin{enumerate}[label=(\alph*), itemsep=2pt, topsep=3pt]
\item $\sum_{n=1}^{\infty} \frac{\ln 2 \cdot \ln 3 \cdots \ln(n+1)}{\ln(2+p) \cdot \ln(3+p) \cdots \ln(n+1+p)}$ ($p > 0$);
\item $\sum_{n=1}^{\infty} \frac{\nu(n)}{n^2}$, 其中 $\nu(n)$ 为数 $n$ 的位数;
\item $\sum_{n=2}^{\infty} \frac{1}{\ln(n!)}$.
\end{enumerate}
\end{prob}

\begin{prob}[M]{2621}
设 $\lambda_n$ ($n = 1, 2, \cdots$) 依次是方程 $\tan x = x$ 的各个正根. 研究级数 $\sum_{n=1}^{\infty} \lambda_n^{-2}$ 的收敛性.
\end{prob}

\begin{prob}[M]{2622}
证明: 设正项级数 $\sum_{n=1}^{\infty} a_n$ 的项单调递减, 则级数 $\sum_{n=1}^{\infty} a_n$ 与级数 $\sum_{n=0}^{\infty} 2^n a_{2^n}$ 同时收敛或同时发散.
\end{prob}

\begin{prob}[M]{2623}
设 $f(x)$ 为单调不增的正值函数. 证明: 若级数 $\sum_{n=1}^{\infty} f(n)$ 收敛, 则对于它的余项 $R_n = \sum_{k=n+1}^{\infty} f(k)$ 有以下估计:
\[
\int_{n+1}^{+\infty} f(x)\,\mathrm{d}x \leqslant R_n \leqslant f(n+1) + \int_{n+1}^{+\infty} f(x)\,\mathrm{d}x.
\]
利用此式, 求级数 $\sum_{n=1}^{\infty} \frac{1}{n^3}$ 的和, 精确到 $0.01$.
\end{prob}

\begin{prob}[M]{2624}
证明叶尔马科夫判别法: 若 $f(x)$ 为单调递减的正值函数, 且
\[
\lim_{x\to+\infty} \frac{\mathrm{e}^x f(\mathrm{e}^x)}{f(x)} = \lambda.
\]
则级数 $\sum_{n=1}^{\infty} f(n)$ 在 $\lambda < 1$ 时收敛, 在 $\lambda > 1$ 时发散.
\end{prob}

\begin{prob}[M]{2625}
证明罗巴切夫斯基判别法: 若正项级数 $\sum_{n=1}^{\infty} a_n$ 的项单调趋于零, 则级数 $\sum_{n=1}^{\infty} a_n$ 与级数 $\sum_{m=0}^{\infty} p_m \cdot 2^{-m}$ 同时收敛或者同时发散, 其中 $p_m$ 是满足不等式
\[
a_n \geqslant 2^{-m} \quad (n = 1, 2, \cdots, p_m)
\]
的项 $a_n$ 的最大的序号.
\end{prob}

\probgroup{研究下列级数的收敛性：}

\begin{prob}[S]{2626}
$\sum_{n=2}^{\infty} \frac{\sqrt{n+2} - \sqrt{n-2}}{n^{\alpha}}$.
\end{prob}

\begin{prob}[S]{2627}
$\sum_{n=1}^{\infty} (\sqrt{n+a} - \sqrt[4]{n^2+n+b})$.
\end{prob}

\begin{prob}[S]{2628}
$\sum_{n=1}^{\infty} \left(\cot \frac{n\pi}{4n-2} - \sin \frac{n\pi}{2n+1}\right)$.
\end{prob}

\begin{prob}[S]{2629}
$\sum_{n=1}^{\infty} \left(\frac{1}{\sqrt{n}} - \sqrt{\ln \frac{n+1}{n}}\right)$.
\end{prob}

\begin{prob}[S]{2630}
$\sum_{n=1}^{\infty} \frac{\ln(n!)}{n^{\alpha}}$.
\end{prob}

\begin{prob}[S]{2631}
$\sum_{n=1}^{\infty} \mathrm{e}^{-\sqrt[3]{n}}$.
\end{prob}

\begin{prob}[S]{2632}
$\sum_{n=1}^{\infty} n^2 \mathrm{e}^{-\sqrt{n}}$.
\end{prob}

\begin{prob}[S]{2633}
$\sum_{n=1}^{\infty} (n^{\frac{1}{n^2+1}} - 1)$.
\end{prob}

\begin{prob}[S]{2634}
$\sum_{n=1}^{\infty} \mathrm{e}^{\frac{a\ln n + b}{c\ln n + d}}$.
\end{prob}

\begin{prob}[S]{2635}
$\sum_{n=1}^{\infty} \frac{1}{\ln^2 \left(\sin \frac{1}{n}\right)}$.
\end{prob}

\begin{prob}[S]{2636}
$\sum_{n=1}^{\infty} \left(\cos \frac{a}{n}\right)^{n^3}$.
\end{prob}

\begin{prob}[S]{2637}
$\sum_{n=3}^{\infty} \ln \left(\frac{\cosh \frac{\pi}{n}}{\cos \frac{\pi}{n}}\right)$.
\end{prob}

\begin{prob}[S]{2638}
$\sum_{n=1}^{\infty} \frac{n!}{n^{\sqrt{n}}}$.
\end{prob}

\begin{prob}[S]{2639}
$\sum_{n=2}^{\infty} \frac{n^{\ln n}}{(\ln n)^n}$.
\end{prob}

\begin{prob}[S]{2640}
$\sum_{n=1}^{\infty} \left(a^{\frac{1}{n}} - \frac{b^{\frac{1}{n}} + c^{\frac{1}{n}}}{2}\right)$ ($a > 0, b > 0, c > 0$).
\end{prob}

\begin{prob}[S]{2641}
$\sum_{n=1}^{\infty} (n^{n^{\alpha}} - 1)$.
\end{prob}

\begin{prob}[S]{2642}
$\sum_{n=1}^{\infty} \left[\ln \frac{1}{n^{\alpha}} - \ln \left(\sin \frac{1}{n^{\alpha}}\right)\right]$.
\end{prob}

\begin{prob}[S]{2643}
$\sum_{n=1}^{\infty} a^{-(b\ln n + c\ln^2 n)}$ ($a > 0$).
\end{prob}

\begin{prob}[S]{2644}
$\sum_{n=1}^{\infty} \frac{n^{2n}}{(n+a)^{n+b}(n+b)^{n+a}}$ ($a > 0, b > 0$).
\end{prob}

\begin{prob}[S]{2645}
$\sum_{n=1}^{\infty} \frac{[(n+1)!]^n}{2! \cdot 4 \cdots (2n)!}$.
\end{prob}

\probgroup{研究级数 $\sum_{n=1}^{\infty} u_n$ 的收敛性，其通项如下：}

\begin{prob}[S]{2646}
$u_n = \int_0^{\frac{1}{n}} \frac{\sqrt{x}\,\mathrm{d}x}{1 + x^2}$.
\end{prob}

\begin{prob}[S]{2647}
$u_n = \frac{1}{\int_0^n \sqrt[4]{1 + x^4}\,\mathrm{d}x}$.
\end{prob}

\begin{prob}[S]{2648}
$u_n = \int_{n\pi}^{(n+1)\pi} \frac{\sin^2 x}{x}\,\mathrm{d}x$.
\end{prob}

\begin{prob}[S]{2649}
$u_n = \int_n^{n+1} \mathrm{e}^{-\sqrt{x}}\,\mathrm{d}x$.
\end{prob}

\begin{prob}[S]{2650}
$u_n = \int_0^{\frac{\pi}{n}} \frac{\sin^3 x}{1 + x}\,\mathrm{d}x$.
\end{prob}

\begin{prob}[S]{2651}
$u_n = \frac{1! + 2! + \cdots + n!}{(2n)!}$.
\end{prob}

\begin{prob}[S]{2652}
$u_n = \frac{\sum_{k=1}^n \ln^2 k}{n^{\alpha}}$.
\end{prob}

\probgroup{用相应的级数来代替数列 $x_n$ ($n = 1, 2, \cdots$)，然后研究它们的收敛性，设：}

\begin{prob}[S]{2653}
$x_n = 1 + \frac{1}{\sqrt{2}} + \cdots + \frac{1}{\sqrt{n}} - 2\sqrt{n}$.
\end{prob}

\begin{prob}[S]{2654}
$x_n = \sum_{k=1}^n \frac{\ln k}{k} - \frac{(\ln n)^2}{2}$.
\end{prob}

\begin{prob}[L]{2655}
对于下列级数:
\begin{enumerate}[label=(\alph*), itemsep=2pt, topsep=3pt]
\item $\sum_{n=1}^{\infty} \frac{1}{n^2}$;
\item $\sum_{n=1}^{\infty} \frac{2^n}{(n+1)!}$;
\item $\sum_{n=1}^{\infty} \frac{1}{(2n-1)!}$,
\end{enumerate}
大约应取多少项来求级数的和方可精确到 $10^{-5}$?
\end{prob}

\probsec{§2. 变号级数收敛性的判别法}

\begin{prob}[M]{2656}
证明：可把非绝对收敛级数的各项在不变更其顺序的情况下分别组合起来，使所得的新级数绝对收敛。
\end{prob}

\begin{prob}[M]{2657}
设有级数 $\sum_{n=1}^{\infty} a_n$，若 (a) 当 $n \to \infty$ 时，此级数的通项 $a_n$ 趋于零；(b) 在不变更其顺序的情况下分别组合该级数的各项，所得级数 $\sum_{n=1}^{\infty} A_n$ 收敛；(c) 在项 $A_n = \sum_{i=p_n}^{p_{n+1}-1} a_i$ ($1 = p_1 < p_2 < \cdots$) 中相加项 $a_i$ 的数目是有界的，证明：级数 $\sum_{n=1}^{\infty} a_n$ 是收敛的。
\end{prob}

\begin{prob}[M]{2658}
证明：若将收敛级数的各项重新排列，使每一项离开原有的位置不超过 $m$ 个位置（$m$ 为预先给定的数），则级数的和不变。
\end{prob}

\probgroup{证明下列级数的收敛性并求它们的和：}

\begin{prob}[S]{2659}
$1 - \frac{3}{2} + \frac{5}{4} - \frac{7}{8} + \cdots$.
\end{prob}

\begin{prob}[S]{2660}
$1 + \frac{1}{2} - \frac{1}{4} + \frac{1}{8} + \frac{1}{16} - \frac{1}{32} + \cdots$.
\end{prob}

\begin{prob}[M]{2661}
$1 - \frac{1}{2} + \frac{1}{3} - \frac{1}{4} + \frac{1}{5} - \frac{1}{6} + \cdots$.
\probhint{运用公式 $1 + \frac{1}{2} + \cdots + \frac{1}{n} = C + \ln n + \varepsilon_n$，其中 $C$ 为欧拉常数，并且 $\lim_{n \to \infty} \varepsilon_n = 0$。}
\end{prob}

\begin{prob}[M]{2662}
已知 $\sum_{n=1}^{\infty} \frac{(-1)^{n+1}}{n} = \ln 2$。将该级数的各项重排，得到下列级数：

(a) $1 + \frac{1}{3} - \frac{1}{2} + \frac{1}{5} + \frac{1}{7} - \frac{1}{4} + \cdots$； \qquad (b) $1 - \frac{1}{2} - \frac{1}{4} + \frac{1}{3} - \frac{1}{6} - \frac{1}{8} + \cdots$。

求这些级数的和。
\end{prob}

\begin{prob}[M]{2663}
把收敛级数 $\sum_{n=1}^{\infty} \frac{(-1)^{n+1}}{\sqrt{n}}$ 的各项重排，使它成为发散的。
\end{prob}

\probgroup{研究下列变号级数的收敛性：}

\begin{prob}[S]{2664}
$\sum_{n=1}^{\infty} \frac{(-1)^{\frac{n(n-1)}{2}}}{2^n}$.
\end{prob}

\begin{prob}[M]{2665}
(a) $\sum_{n=1}^{\infty} (-1)^n \left(\frac{2n+100}{3n+1}\right)^n$； \qquad (b) $1 + \frac{1}{2} + \frac{1}{3} - \frac{1}{4} - \frac{1}{5} - \frac{1}{6} + \frac{1}{7} + \frac{1}{8} + \frac{1}{9} - \cdots$.
\end{prob}

\begin{prob}[M]{2666}
设
\[\sum_{n=1}^{\infty} (-1)^n b_n, \tag{1}\]
其中 $b_n > 0$，且当 $n \to \infty$ 时 $b_n \to 0$。由此是否可知级数 $(1)$ 收敛？考察例子
\[\sum_{n=1}^{\infty} (-1)^n \cdot \frac{2 + (-1)^n}{n}.\]
\end{prob}

\probgroup{研究下列变号级数的收敛性：}

\begin{prob}[S]{2667}
$\sum_{n=1}^{\infty} \frac{\ln^{100} n}{n} \sin \frac{n\pi}{4}$.
\end{prob}

\begin{prob}[S]{2668}
$\sum_{n=1}^{\infty} (-1)^n \frac{\sin^2 n}{n}$.
\end{prob}

\begin{prob}[S]{2669}
$\sum_{n=1}^{\infty} (-1)^n \frac{\sqrt{n}}{n+100}$.
\end{prob}

\begin{prob}[S]{2670}
$\sum_{n=2}^{\infty} \frac{(-1)^n}{\sqrt{n} + (-1)^n}$.
\end{prob}

\begin{prob}[S]{2671}
$\sum_{n=1}^{\infty} \sin \left(\pi \sqrt{n^2 + k^2}\right)$.
\end{prob}

\begin{prob}[S]{2672}
$\sum_{n=1}^{\infty} \frac{(-1)^{[\sqrt{n}]}}{n}$.
\end{prob}

\begin{prob}[M]{2673}
(a) $\sum_{n=1}^{\infty} \frac{(-1)^n}{\sqrt[n]{n}}$； \qquad (b) $\sum_{n=2}^{\infty} \frac{1}{\ln^2 n} \cos \frac{\pi n^2}{n+1}$.
\end{prob}

\begin{prob}[M]{2674}
证明：若
\[\frac{b_n}{b_{n+1}} = 1 + \frac{p}{n} + o\left(\frac{1}{n}\right),\]
其中 $p > 0$（参阅习题 2606），则交错级数 $b_1 - b_2 + b_3 - b_4 + \cdots + (-1)^{n-1} b_n + \cdots$（$b_n > 0$）收敛。
\end{prob}

\probgroup{研究下列级数的绝对收敛性（除了习题 2690）和条件收敛性：}

\begin{prob}[S]{2675}
$\sum_{n=1}^{\infty} \frac{(-1)^{n-1}}{n^p}$.
\end{prob}

\begin{prob}[S]{2676}
$\sum_{n=1}^{\infty} \frac{(-1)^{n-1}}{n^{p + \frac{1}{n}}}$.
\end{prob}

\begin{prob}[S]{2677}
$\sum_{n=2}^{\infty} \ln \left[1 + \frac{(-1)^n}{n^p}\right]$.
\end{prob}

\begin{prob}[S]{2678}
$\sum_{n=1}^{\infty} (-1)^{n-1} \frac{2^n \sin^{2n} x}{n}$.
\end{prob}

\begin{prob}[S]{2679}
$\sum_{n=1}^{\infty} \frac{(-1)^n}{x + n}$.
\end{prob}

\begin{prob}[S]{2680}
$\sum_{n=2}^{\infty} \frac{(-1)^n}{[n + (-1)^n]^p}$.
\end{prob}

\begin{prob}[S]{2681}
$\sum_{n=1}^{\infty} \frac{(-1)^{n-1}}{[\sqrt{n} + (-1)^{n-1}]^p}$.
\end{prob}

\begin{prob}[S]{2682}
$\sum_{n=1}^{\infty} \frac{\sin \frac{n\pi}{4}}{n^p + \sin \frac{n\pi}{4}}$.
\end{prob}

\begin{prob}[S]{2683}
$\sum_{n=1}^{\infty} (-1)^n \frac{n-1}{n+1} \frac{1}{\sqrt[100]{n}}$.
\end{prob}

\begin{prob}[S]{2684}
$\sum_{n=1}^{\infty} (-1)^{\frac{n(n-1)}{2}} \frac{n^{100}}{2^n}$.
\end{prob}

\begin{prob}[S]{2685}
$\sum_{n=1}^{\infty} \frac{(-1)^n}{\sqrt[n^2]{n}}$.
\end{prob}

\begin{prob}[S]{2686}
$\sum_{n=2}^{\infty} \frac{\sin \frac{n\pi}{12}}{\ln n}$.
\end{prob}

\begin{prob}[S]{2687}
$\sum_{n=1}^{\infty} \frac{(-1)^{[\sqrt{n}]}}{n^p}$.
\end{prob}

\begin{prob}[S]{2688}
$\sum_{n=1}^{\infty} \frac{(-1)^{[\ln n]}}{n}$.
\end{prob}

\begin{prob}[S]{2689}
$\sum_{n=1}^{\infty} (-1)^{n-1} \left[\frac{1 \cdot 3 \cdot 5 \cdots (2n-1)}{2 \cdot 4 \cdot 6 \cdots (2n)}\right]^p$.
\end{prob}

\begin{prob}[S]{2690}
$\sum_{n=1}^{\infty} \frac{\sin n \cdot \sin n^2}{n}$.
\end{prob}

\begin{prob}[M]{2691}
$\sum_{n=1}^{\infty} \sin n^2$.
\probhint{证明 $\lim_{n \to \infty} \sin n^2 \neq 0$。}
\end{prob}

\begin{prob}[M]{2692}
设
\[R(x) = \frac{a_0 x^p + a_1 x^{p-1} + \cdots + a_p}{b_0 x^q + b_1 x^{q-1} + \cdots + b_q}\]
为有理函数，式中 $a_0 \neq 0, b_0 \neq 0$，且当 $x \geqslant n_0$ 时 $|b_0 x^q + b_1 x^{q-1} + \cdots + b_q| > 0$。研究级数
\[\sum_{n=n_0}^{\infty} (-1)^n R(n)\]
的绝对收敛性和条件收敛性。
\end{prob}

\probgroup{研究下列级数的收敛性：}

\begin{prob}[S]{2693}
$\frac{1}{1^p} - \frac{1}{2^q} + \frac{1}{3^p} - \frac{1}{4^q} + \frac{1}{5^p} - \frac{1}{6^q} + \cdots$.
\end{prob}

\begin{prob}[S]{2694}
$1 + \frac{1}{3^p} - \frac{1}{2^p} + \frac{1}{5^p} + \frac{1}{7^p} - \frac{1}{4^p} + \cdots$.
\end{prob}

\begin{prob}[M]{2695}
(a) $1 + \frac{1}{3^p} - \frac{1}{1^p} + \frac{1}{5^p} + \frac{1}{7^p} - \frac{1}{3^p} + \frac{1}{9^p} + \frac{1}{11^p} - \frac{1}{5^p} + \cdots$；

(b) $1 - \frac{2}{2^q} + \frac{1}{3^p} + \frac{1}{4^p} - \frac{2}{5^q} + \frac{1}{6^p} + \frac{1}{7^p} - \frac{2}{8^q} + \frac{1}{9^p} + \cdots$.
\end{prob}

\begin{prob}[M]{2696}
证明：级数

(a) $\sin x + \frac{\sin 2x}{2} + \frac{\sin 3x}{3} + \cdots$； \qquad (b) $\cos x + \frac{\cos 2x}{2} + \frac{\cos 3x}{3} + \cdots$

在区间 $(0, \pi)$ 内不绝对收敛。
\end{prob}

\begin{prob}[M]{2697}
对于级数
\[\sum_{n=1}^{\infty} \frac{\cos nx}{n^p}, \qquad \sum_{n=1}^{\infty} \frac{\sin nx}{n^p} \quad (0 < x < \pi)\]
对全体参数 $(p, x)$ 定出：(a) 绝对收敛域；(b) 非绝对收敛域。
\end{prob}

\begin{prob}[L]{2698}
研究下列级数的收敛性：
\begin{enumerate}[label=(\alph*), itemsep=2pt, topsep=3pt]
\item $\sum_{n=2}^{\infty} \frac{(-1)^n \sqrt[n]{n}}{\ln n}$；
\item $\sum_{n=2}^{\infty} \frac{\sin\left(n + \frac{1}{n}\right)}{\ln(\ln n)}$；
\item $\sum_{n=10}^{\infty} \frac{\sin n}{n + 10 \sin n}$.
\end{enumerate}
\end{prob}

\begin{prob}[M]{2699}
对于级数
\[\sum_{n=1}^{\infty} (-1)^{n-1} \frac{(1+p)(2+p)\cdots(n+p)}{n! n^q}\]
定出：(a) 绝对收敛域；(b) 条件收敛域。
\end{prob}

\begin{prob}[M]{2700}
研究级数
\[\sum_{n=1}^{\infty} \binom{m}{n}\]
的收敛性，其中 $\binom{m}{n} = \frac{m(m-1)\cdots(m-n+1)}{n!}$。
\end{prob}

\begin{prob}[M]{2701}
若级数 $\sum_{n=1}^{\infty} a_n$ 收敛且
\[\lim_{n \to \infty} \frac{b_n}{a_n} = 1,\]
则可否断定级数 $\sum_{n=1}^{\infty} b_n$ 也收敛？研究例子
\[\sum_{n=1}^{\infty} \frac{(-1)^n}{\sqrt{n}} \quad \text{和} \quad \sum_{n=1}^{\infty} \left[\frac{(-1)^n}{\sqrt{n}} + \frac{1}{n}\right].\]
\end{prob}

\begin{prob}[M]{2702.1}
设 $\sum_{n=1}^{\infty} a_n$ 为非绝对收敛的级数，
\[P_n = \sum_{i=1}^{n} \frac{|a_i| + a_i}{2}, \qquad N_n = \sum_{i=1}^{n} \frac{|a_i| - a_i}{2}.\]
证明：
\[\lim_{n \to \infty} \frac{N_n}{P_n} = 1.\]
\end{prob}

\begin{prob}[M]{2702.2}
证明：对于每一个 $p > 0$，级数
\[\sum_{n=1}^{\infty} \frac{(-1)^{n+1}}{n^p}\]
的和在 $\frac{1}{2}$ 与 $1$ 之间。
\end{prob}

\begin{prob}[M]{2703}
在下列级数中：

(a) $\sum_{n=1}^{\infty} \frac{(-1)^{n+1}}{\sqrt{n^2 + 1}}$； \qquad (b) $\sum_{n=1}^{\infty} \frac{\sin n}{\sqrt{n}}$，

应取多少项来计算级数的和，方可使其精度达到 $\varepsilon = 10^{-6}$？
\end{prob}

\begin{prob}[M]{2704}
证明：若把级数
\[1 - \frac{1}{2} + \frac{1}{3} - \frac{1}{4} + \frac{1}{5} - \cdots\]
的各项重排，使依次 $p$ 个正项的一组与依次 $q$ 个负项的一组相交替，则新级数的和为
\[\ln 2 + \frac{1}{2} \ln \frac{p}{q}.\]
\end{prob}

\begin{prob}[M]{2705}
证明：若改变调和级数
\[1 + \frac{1}{2} + \frac{1}{3} + \frac{1}{4} + \cdots,\]
部分项的符号，使得 $p$ 个正项之后跟随着 $q$ 个负项（$p \neq q$），但不变更原来的顺序，则此级数仍是发散的。仅当 $p = q$ 时得到收敛级数。
\end{prob}

\probsec{§3. 级数的运算}

\begin{prob}[M]{2706}
若两个级数，(a) 一个收敛，而另一个发散；(b) 两个都发散，则关于这两个级数的和可下何种断言？
\end{prob}

\begin{prob}[M]{2707}
求二级数的和：
\[\sum_{n=1}^{\infty} \left[\frac{1}{3^n} + \frac{(-1)^n}{n^3}\right] + \sum_{n=1}^{\infty} \left[\frac{1}{3^{n+1}} + \frac{(-1)^{n+1}}{n^3}\right].\]
\end{prob}

\probgroup{求下列级数的和：}

\begin{prob}[S]{2708}
$\sum_{n=1}^{\infty} \left[\frac{1}{2^n} + \frac{(-1)^n}{3^n}\right]$.
\end{prob}

\begin{prob}[S]{2709}
$\sum_{n=1}^{\infty} \frac{\cos \frac{2n\pi}{3}}{2^n}$.
\end{prob}

\begin{prob}[S]{2710}
$\sum_{n=0}^{\infty} x^{\left[\frac{n}{2}\right]} y^{\left[\frac{n+1}{2}\right]}$ ($|xy| < 1$).
\end{prob}

\begin{prob}[M]{2711}
证明：
\[\sum_{n=0}^{\infty} \frac{1}{n!} \cdot \sum_{n=0}^{\infty} \frac{(-1)^n}{n!} = 1.\]
\end{prob}

\begin{prob}[M]{2712}
证明：
\[\left(\sum_{n=0}^{\infty} q^n\right)^2 = \sum_{n=0}^{\infty} (n+1) q^n \quad (|q| < 1).\]
\end{prob}

\begin{prob}[M]{2713}
证明：收敛级数
\[\sum_{n=1}^{\infty} \frac{(-1)^{n+1}}{\sqrt{n}}\]
的平方是发散级数。
\end{prob}

\begin{prob}[M]{2714}
证明：下面二收敛级数
\[\sum_{n=1}^{\infty} \frac{(-1)^{n-1}}{n^\alpha} \quad (\alpha > 0) \quad \text{及} \quad \sum_{n=1}^{\infty} \frac{(-1)^{n-1}}{n^\beta} \quad (\beta > 0)\]
之积当 $\alpha + \beta > 1$ 时是收敛级数，而当 $\alpha + \beta < 1$ 时是发散级数。
\end{prob}

\begin{prob}[M]{2715}
验证下面二发散级数
\[1 - \sum_{n=1}^{\infty} \left(\frac{3}{2}\right)^n \quad \text{和} \quad 1 + \sum_{n=1}^{\infty} \left(\frac{3}{2}\right)^{n-1} \left(2^n + \frac{1}{2^{n+1}}\right)\]
之积是绝对收敛级数。
\end{prob}

\probsec{§4. 函数项级数}

\probgroup{定出下列函数项级数的绝对收敛域和条件收敛域.}

\begin{prob}[S]{2716}
$\sum_{n=1}^{\infty} \frac{n}{x^n}$.
\end{prob}

\begin{prob}[S]{2717}
$\sum_{n=1}^{\infty} \frac{(-1)^n}{2n-1}\left(\frac{1-x}{1+x}\right)^n$.
\end{prob}

\begin{prob}[S]{2718}
$\sum_{n=1}^{\infty} \frac{n}{n+1}\left(\frac{x}{2x+1}\right)^n$.
\end{prob}

\begin{prob}[S]{2719}
$\sum_{n=1}^{\infty} \frac{1\cdot 3\cdot\cdots\cdot (2n-1)}{2\cdot 4\cdots (2n)}\left(\frac{2x}{1+x^2}\right)^n$.
\end{prob}

\begin{prob}[S]{2720}
$\sum_{n=1}^{\infty} \frac{n\cdot 3^{2n}}{2^n}x^n(1-x)^n$.
\end{prob}

\begin{prob}[S]{2721}
$\sum_{n=1}^{\infty} \frac{2^n \sin^n x}{n^2}$.
\end{prob}

\begin{prob}[S]{2722}
$\sum_{n=1}^{\infty} \frac{(-1)^n}{(x+n)^p}$.
\end{prob}

\begin{prob}[S]{2723}
$\sum_{n=1}^{\infty} \frac{n^p \sin nx}{1+n^q} \ (q > 0;\ 0 < x < \pi)$.
\end{prob}

\begin{prob}[S]{2724}
$\sum_{n=1}^{\infty} \frac{x^n}{1-x^n}$ (兰伯特级数).
\end{prob}

\begin{prob}[S]{2725}
$\sum_{n=1}^{\infty} \left[\frac{x(x+n)}{n}\right]^n$.
\end{prob}

\begin{prob}[S]{2726}
$\sum_{n=1}^{\infty} \frac{x^n}{1+x^{2n}}$.
\end{prob}

\begin{prob}[S]{2727}
$\sum_{n=1}^{\infty} \frac{x^n}{(1+x)(1+x^2)\cdots (1+x^n)}$.
\end{prob}

\begin{prob}[S]{2728}
$\sum_{n=1}^{\infty} n\mathrm{e}^{-nx}$.
\end{prob}

\begin{prob}[S]{2729}
$\sum_{n=1}^{\infty} \frac{1}{\sqrt[n]{n!}}\cdot \frac{1}{1+a^{2n}x^2}$.
\end{prob}

\begin{prob}[S]{2730}
$\sum_{n=1}^{\infty} (2-x)(2-x^{\frac{1}{2}})(2-x^{\frac{1}{3}})\cdots (2-x^{\frac{1}{n}}) \ (x > 0)$.
\end{prob}

\begin{prob}[S]{2731}
$\sum_{n=1}^{\infty} \frac{(n+x)^n}{n^{n+x}}$.
\end{prob}

\begin{prob}[S]{2732}
$\sum_{n=1}^{\infty} \frac{x^n y^n}{x^n+y^n} \ (x > 0;\ y > 0)$.
\end{prob}

\begin{prob}[S]{2733}
$\sum_{n=1}^{\infty} \frac{x^n}{n+y^n} \ (y \geqslant 0)$.
\end{prob}

\begin{prob}[S]{2734}
$\sum_{n=1}^{\infty} \sqrt[n]{|x|^{n^2}+|y|^{n^2}}$.
\end{prob}

\begin{prob}[S]{2735}
$\sum_{n=1}^{\infty} \frac{\ln(1+x^n)}{n^y} \ (x \geqslant 0)$.
\end{prob}

\begin{prob}[S]{2736}
$\sum_{n=1}^{\infty} \tan^n\left(x+\frac{y}{n}\right)$.
\end{prob}

\begin{prob}[M]{2737}
证明: 若洛朗级数 $\sum_{n=-\infty}^{+\infty} a_n x^n$ 当 $x=x_1$ 和 $x=x_2$ ($|x_1| < |x_2|$) 时收敛, 则此级数当 $|x_1| < |x| < |x_2|$ 时也收敛.
\end{prob}

\begin{prob}[M]{2738}
求洛朗级数
\[
\sum_{n=-\infty}^{+\infty} \frac{n}{2^{|n|}} x^n
\]
的收敛域并求它的和.
\end{prob}

\begin{prob}[L]{2739}
求牛顿级数的绝对收敛域和条件收敛域:
\begin{enumerate}[label=(\alph*), itemsep=2pt, topsep=3pt]
\item $\sum_{n=1}^{\infty} \frac{x^{[n]}}{n!}$;
\item $\sum_{n=1}^{\infty} \frac{1}{n^p}\frac{x^{[n]}}{n!}$;
\item $\sum_{n=1}^{\infty} \frac{(\mathrm{e}x)^n y^{[n]}}{n^n}$.
\end{enumerate}
其中 $x^{[n]} = x(x-1)\cdots [x-(n-1)]$.
\end{prob}

\begin{prob}[M]{2740}
证明: 若狄利克雷级数 $\sum_{n=1}^{\infty} \frac{a_n}{n^x}$ 当 $x=x_0$ 时收敛, 则此级数当 $x > x_0$ 时也收敛.
\end{prob}

\begin{prob}[M]{2741}
证明: 序列 $f_n(x)$ ($n=1,2,\cdots$) 在集合 $X$ 上一致收敛于极限函数 $f(x)$ 的充分必要条件是
\[
\lim_{n\to\infty} \left\{\sup_{x\in X} r_n(x)\right\} = 0,
\]
式中 $r_n(x) = |f(x)-f_n(x)|$.
\end{prob}

\begin{prob}[L]{2742}
序列 $f_n(x)$ ($n=1,2,\cdots$):
\begin{enumerate}[label=(\alph*), itemsep=2pt, topsep=3pt]
\item 在区间 $(x_0, +\infty)$ 上收敛;
\item 在每一个有限的区间 $(a,b) \subset (x_0, +\infty)$ 上一致收敛;
\item 在区间 $(x_0, +\infty)$ 上一致收敛,
\end{enumerate}
这分别是什么意思?
\end{prob}

\begin{prob}[M]{2743}
对于序列
\[
f_n(x) = x^n \ (n=1,2,\cdots) \ (0 < x < 1),
\]
求出其项的最小序号 $N=N(\varepsilon, x)$, 使从这项起序列的项在已知点 $x$ 与极限函数的差不超过 $0.001$, 设 $x = \frac{1}{10}, \frac{1}{\sqrt{10}}, \cdots, \frac{1}{\sqrt[m]{10}}, \cdots$. 此序列在区间 $(0,1)$ 内是否一致收敛?
\end{prob}

\begin{prob}[L]{2744}
应当取级数
\[
\sum_{n=1}^{\infty} \frac{\sin nx}{n(n+1)}
\]
的多少项方可使部分和 $S_n(x)$ 当 $-\infty < x < +\infty$ 时与级数的和之差小于 $\varepsilon$? 对于下列 $\varepsilon$ 值, 给出具体的计算结果:
\begin{enumerate}[label=(\alph*), itemsep=2pt, topsep=3pt]
\item $\varepsilon = 0.1$;
\item $\varepsilon = 0.01$;
\item $\varepsilon = 0.001$.
\end{enumerate}
\end{prob}

\begin{prob}[M]{2745}
对怎样的 $n$, 不等式
\[
\left|\mathrm{e}^x - \sum_{i=0}^{n} \frac{x^i}{i!}\right| < 0.001 \quad (0 \leqslant x \leqslant 10)
\]
能保证成立?
\end{prob}

\probgroup{研究序列在所给区间上的一致收敛性:}

\begin{prob}[M]{2746}
$f_n(x) = x^n$; (a) $0 \leqslant x \leqslant \frac{1}{2}$; \qquad (b) $0 \leqslant x \leqslant 1$.
\end{prob}

\begin{prob}[S]{2747}
$f_n(x) = x^n - x^{n+1}$; $0 \leqslant x \leqslant 1$.
\end{prob}

\begin{prob}[S]{2748}
$f_n(x) = x^n - x^{2n}$; $0 \leqslant x \leqslant 1$.
\end{prob}

\begin{prob}[S]{2749}
$f_n(x) = \frac{1}{x+n}$; $0 < x < +\infty$.
\end{prob}

\begin{prob}[S]{2750}
$f_n(x) = \frac{nx}{1+n+x}$; $0 \leqslant x \leqslant 1$.
\end{prob}

\begin{prob}[L]{2751}
$f_n(x) = \frac{x^n}{1+x^n}$;
\begin{enumerate}[label=(\alph*), itemsep=2pt, topsep=3pt]
\item $0 \leqslant x \leqslant 1-\varepsilon$;
\item $1-\varepsilon \leqslant x \leqslant 1+\varepsilon$;
\item $1+\varepsilon \leqslant x < +\infty$, 其中 $\varepsilon > 0$.
\end{enumerate}
\end{prob}

\begin{prob}[M]{2752}
$f_n(x) = \frac{2nx}{1+n^2 x^2}$; (a) $0 \leqslant x \leqslant 1$; \qquad (b) $1 < x < +\infty$.
\end{prob}

\begin{prob}[S]{2753}
$f_n(x) = \sqrt{x^2+\frac{1}{n^2}}$; $-\infty < x < +\infty$.
\end{prob}

\begin{prob}[S]{2754}
$f_n(x) = n\left(\sqrt{x+\frac{1}{n}}-\sqrt{x}\right)$; $1 < x < +\infty$.
\end{prob}

\begin{prob}[M]{2755}
(a) $f_n(x) = \frac{\sin nx}{n}$; $-\infty < x < +\infty$; \qquad (b) $f_n(x) = \sin\frac{x}{n}$; $-\infty < x < +\infty$.
\end{prob}

\begin{prob}[M]{2756}
(a) $f_n(x) = \arctan nx$; $0 < x < +\infty$; \qquad (b) $f_n(x) = x\arctan nx$; $0 < x < +\infty$.
\end{prob}

\begin{prob}[S]{2757}
$f_n(x) = \mathrm{e}^{n(x-1)}$; $0 < x < 1$.
\end{prob}

\begin{prob}[M]{2758}
$f_n(x) = \mathrm{e}^{-(x-n)^2}$; (a) $-l < x < l$, 其中 $l$ 为任意正数; \qquad (b) $-\infty < x < +\infty$.
\end{prob}

\begin{prob}[S]{2759}
$f_n(x) = \frac{x}{n}\ln\frac{x}{n}$; $0 < x < 1$.
\end{prob}

\begin{prob}[M]{2760}
$f_n(x) = \left(1+\frac{x}{n}\right)^n$; (a) 在有限的区间 $(a,b)$ 上; \qquad (b) 在区间 $(-\infty, +\infty)$ 上.
\end{prob}

\begin{prob}[S]{2761}
$f_n(x) = n\left(x^{\frac{1}{n}}-1\right)$; $1 \leqslant x \leqslant a$.
\end{prob}

\begin{prob}[S]{2762}
$f_n(x) = \sqrt[n]{1+x^n}$; $0 \leqslant x \leqslant 2$.
\end{prob}

\begin{prob}[M]{2763}
$f_n(x) = \begin{cases} n^2 x, & 0 \leqslant x \leqslant \dfrac{1}{n}, \\ n^2\left(\dfrac{2}{n}-x\right), & \dfrac{1}{n} < x < \dfrac{2}{n}, \\ 0, & x \geqslant \dfrac{2}{n}; \end{cases}$ 在闭区间 $0 \leqslant x \leqslant 1$ 上.
\end{prob}

\begin{prob}[M]{2764}
设 $f(x)$ 为定义于闭区间 $[a,b]$ 上的任意函数, 且
\[
f_n(x) = \frac{[nf(x)]}{n} \ (n=1,2,\cdots),
\]
证明: 当 $n\to\infty$ 时
\[
f_n(x) \rightrightarrows f(x) \ (a \leqslant x \leqslant b).
\]
\end{prob}

\begin{prob}[M]{2765}
设函数 $f(x)$ 在区间 $(a,b)$ 内有连续的导数 $f'(x)$, 且
\[
f_n(x) = n\left[f\left(x+\frac{1}{n}\right) - f(x)\right].
\]
证明: 在闭区间 $\alpha \leqslant x \leqslant \beta$ 上 (其中 $a < \alpha < \beta < b$), $f_n(x) \rightrightarrows f'(x)$.
\end{prob}

\begin{prob}[M]{2766}
设 $f_n(x) = \sum_{i=0}^{n-1} \frac{1}{n} f\left(x+\frac{i}{n}\right)$, 其中 $f(x)$ 为 $(-\infty, +\infty)$ 上的连续函数. 证明: 序列 $f_n(x)$ 在任何有限闭区间 $[a,b]$ 上一致收敛.
\end{prob}

\probgroup{研究下列级数的收敛性:}

\begin{prob}[M]{2767}
$\sum_{n=0}^{\infty} x^n$; (a) 在区间 $|x| < q$ 内, 此外 $q < 1$; \qquad (b) 在区间 $|x| < 1$ 内.
\end{prob}

\begin{prob}[M]{2768}
(a) $\sum_{n=1}^{\infty} \frac{x^n}{n^2}$; 在闭区间 $-1 \leqslant x \leqslant 1$ 上; \qquad (b) $\sum_{n=0}^{\infty} \frac{x^n}{n!}$; 在区间 $(0, +\infty)$ 内.
\end{prob}

\begin{prob}[S]{2769}
$\sum_{n=0}^{\infty} (1-x)x^n$; 在闭区间 $0 \leqslant x \leqslant 1$ 上.
\end{prob}

\begin{prob}[S]{2770}
$\sum_{n=1}^{\infty} \left(\frac{x^n}{n} - \frac{x^{n+1}}{n+1}\right)$; $-1 \leqslant x \leqslant 1$.
\end{prob}

\begin{prob}[S]{2771}
$\sum_{n=1}^{\infty} \frac{x}{[(n-1)x+1](nx+1)}$; $0 < x < +\infty$.
\end{prob}

\begin{prob}[S]{2772}
$\sum_{n=1}^{\infty} \frac{1}{(x+n)(x+n+1)}$; $0 < x < +\infty$.
\end{prob}

\begin{prob}[M]{2773}
$\sum_{n=1}^{\infty} \frac{nx}{(1+x)(1+2x)\cdots (1+nx)}$; (a) $0 \leqslant x \leqslant \varepsilon$, 其中 $\varepsilon > 0$; \qquad (b) $\varepsilon \leqslant x < +\infty$.
\end{prob}

\begin{prob}[L]{2774}
利用魏尔斯特拉斯判别法, 证明下列函数项级数在所指区间内的一致收敛性:
\begin{enumerate}[label=(\alph*), itemsep=2pt, topsep=3pt]
\item $\sum_{n=1}^{\infty} \frac{1}{x^2+n^2}$, $-\infty < x < +\infty$;
\item $\sum_{n=1}^{\infty} \frac{(-1)^n}{x+2^n}$, $-2 < x < +\infty$;
\item $\sum_{n=1}^{\infty} \frac{x}{1+n^4 x^2}$, $0 \leqslant x < +\infty$;
\item $\sum_{n=1}^{\infty} \frac{nx}{1+n^5 x^2}$, $|x| < +\infty$;
\item $\sum_{n=1}^{\infty} \frac{n^2}{\sqrt{n!}}(x^n+x^{-n})$, $\frac{1}{2} \leqslant |x| \leqslant 2$;
\item $\sum_{n=1}^{\infty} \frac{x^n}{\left[\frac{n}{2}\right]!}$, $|x| < a$, $a$ 为任意正数;
\item $\sum_{n=1}^{\infty} \frac{\sin nx}{\sqrt[3]{n^4+x^4}}$, $|x| < +\infty$;
\item $\sum_{n=1}^{\infty} \frac{\cos nx}{n^2}$, $|x| < +\infty$;
\end{enumerate}
\begin{enumerate}[label=(\alph*), start=9, itemsep=2pt, topsep=3pt]
  \item $\displaystyle\sum_{n=1}^{\infty} \frac{\sin nx}{n\sqrt{n}}$，$ |x| < +\infty $；
  \item $\displaystyle\sum_{n=2}^{\infty} \ln\left(1 + \frac{x^2}{n\ln^2 n}\right)$，$ |x| < a $；
  \item $\displaystyle\sum_{n=1}^{\infty} x^2 \mathrm{e}^{-nx}$，$ 0 \leqslant x < +\infty $；
  \item $\displaystyle\sum_{n=1}^{\infty} \arctan\frac{2x}{x^2+n^3}$，$ |x| < +\infty $。
\end{enumerate}
\end{prob}

\probgroup{研究下列函数项级数在指定区间上的一致收敛性：}

\begin{prob}[M]{2775}
$\displaystyle\sum_{n=1}^{\infty} \frac{\sin nx}{n}$；(a) 在闭区间 $\varepsilon \leqslant x \leqslant 2\pi - \varepsilon$ 上，其中 $\varepsilon > 0$；(b) 在闭区间 $0 \leqslant x \leqslant 2\pi$ 上。
\end{prob}

\begin{prob}[S]{2776}
$\displaystyle\sum_{n=1}^{\infty} 2^n \sin\frac{1}{3^n x}$；$0 < x < +\infty$。
\end{prob}

\begin{prob}[M]{2777}
$\displaystyle\sum_{n=1}^{\infty} \frac{(-1)^n}{x+n}$；$0 < x < +\infty$。
\probhint{估计级数的余项。}
\end{prob}

\begin{prob}[S]{2778}
$\displaystyle\sum_{n=2}^{\infty} \frac{(-1)^n}{n+\sin x}$；$0 \leqslant x \leqslant 2\pi$。
\end{prob}

\begin{prob}[S]{2779}
$\displaystyle\sum_{n=1}^{\infty} \frac{(-1)^{\frac{n(n-1)}{2}}}{\sqrt[3]{n^2+\mathrm{e}^x}}$；$|x| \leqslant 10$。
\end{prob}

\begin{prob}[S]{2780}
$\displaystyle\sum_{n=1}^{\infty} \frac{\cos\frac{2n\pi}{3}}{\sqrt{n^2+x^2}}$；$-\infty < x < +\infty$。
\end{prob}

\begin{prob}[S]{2781}
$\displaystyle\sum_{n=1}^{\infty} \frac{\sin x \sin nx}{\sqrt{n+x}}$；$0 \leqslant x < +\infty$。
\end{prob}

\begin{prob}[S]{2782}
$\displaystyle\sum_{n=1}^{\infty} \frac{(-1)^{[\sqrt{n}]}}{\sqrt{n(n+x)}}$；$0 \leqslant x < +\infty$。
\end{prob}

\begin{prob}[M]{2783}
不连续函数序列可否一致收敛于连续函数？研究例子
\[
f_n(x) = \frac{1}{n}\psi(x) \quad (n=1,2,\cdots),
\]
其中
\[
\psi(x) =
\begin{cases}
0, & x\text{ 为无理数}, \\
1, & x\text{ 为有理数}.
\end{cases}
\]
\end{prob}

\begin{prob}[M]{2784}
证明：若级数
\[
\sum_{n=1}^{\infty} |f_n(x)|
\]
在 $[a,b]$ 上一致收敛，则级数
\[
\sum_{n=1}^{\infty} f_n(x)
\]
在 $[a,b]$ 上也一致收敛。
\end{prob}

\begin{prob}[M]{2785}
若级数
\[
\sum_{n=1}^{\infty} f_n(x)
\]
在 $[a,b]$ 上绝对并一致收敛，则级数
\[
\sum_{n=1}^{\infty} |f_n(x)|
\]
在 $[a,b]$ 上是否必定一致收敛？研究例子
\[
\sum_{n=0}^{\infty} (-1)^n (1-x)x^n,
\]
其中 $0 \leqslant x \leqslant 1$。
\end{prob}

\begin{prob}[M]{2786}
证明：绝对且一致收敛的级数
\[
\sum_{n=1}^{\infty} f_n(x) \quad (0 \leqslant x \leqslant 1),
\]
其中
\[
f_n(x) =
\begin{cases}
0, & 0 \leqslant x \leqslant 2^{-(n+1)}, \\
\dfrac{1}{n}\sin^2(2^{n+1}\pi x), & 2^{-(n+1)} < x < 2^{-n}, \\
0, & 2^{-n} \leqslant x \leqslant 1,
\end{cases}
\]
不能用非负项的收敛数项级数作为其强级数。
\end{prob}

\begin{prob}[M]{2787}
证明：若函数项级数
\[
\sum_{n=1}^{\infty} \varphi_n(x)
\]
的各项是闭区间 $[a,b]$ 上的单调函数，此级数在闭区间 $[a,b]$ 的两个端点绝对收敛，则此级数在闭区间 $[a,b]$ 上绝对并一致收敛。
\end{prob}

\begin{prob}[M]{2788}
证明：幂级数
\[
\sum_{n=0}^{\infty} a_n x^n
\]
在全部位于其收敛区间内的任何闭区间上一致收敛。
\end{prob}

\begin{prob}[M]{2789}
设 $a_n \to \infty$ 且级数 $\displaystyle\sum_{n=1}^{\infty} \left|\frac{1}{a_n}\right|$ 收敛。证明：级数
\[
\sum_{n=1}^{\infty} \frac{1}{x - a_n}
\]
在不包含点 $a_n$ $(n=1,2,\cdots)$ 的任何有界闭集合上绝对并一致收敛。
\end{prob}

\begin{prob}[M]{2790}
证明：若级数 $\displaystyle\sum_{n=1}^{\infty} a_n$ 收敛，则狄利克雷级数
\[
\sum_{n=1}^{\infty} \frac{a_n}{n^x}
\]
在 $x \geqslant 0$ 时一致收敛。
\end{prob}

\begin{prob}[M]{2791}
设级数 $\displaystyle\sum_{n=1}^{\infty} a_n$ 收敛。证明：级数
\[
\sum_{n=1}^{\infty} a_n \mathrm{e}^{-nx}
\]
在区域 $x \geqslant 0$ 内一致收敛。
\end{prob}

\begin{prob}[M]{2792}
证明：函数
\[
f(x) = \sum_{n=1}^{\infty} \frac{\sin nx}{n^3}
\]
在区域 $-\infty < x < +\infty$ 内连续并有连续的导数。
\end{prob}

\begin{prob}[M]{2793}
证明：函数
\[
f(x) = \sum_{n=-\infty}^{+\infty} \frac{1}{(n-x)^2}
\]
\begin{enumerate}[label=(\alph*), itemsep=2pt, topsep=3pt]
  \item 在除整数点 $x = 0, \pm 1, \pm 2, \cdots$ 外的一切点有定义并且连续；
  \item 为周期函数，其周期等于 $1$。
\end{enumerate}
\end{prob}

\begin{prob}[M]{2794}
证明：级数
\[
\sum_{n=1}^{\infty} \left[ nx\mathrm{e}^{-nx} - (n-1)x\mathrm{e}^{-(n-1)x} \right]
\]
在闭区间 $0 \leqslant x \leqslant 1$ 上收敛但不一致收敛，而它的和是此区间上的连续函数。
\end{prob}

\begin{prob}[L]{2795}
确定函数 $f(x)$ 的存在域并研究其连续性，设：
\begin{enumerate}[label=(\alph*), itemsep=2pt, topsep=3pt]
  \item $\displaystyle f(x) = \sum_{n=1}^{\infty} \left(x + \frac{1}{n}\right)^n$；
  \item $\displaystyle f(x) = \sum_{n=1}^{\infty} \frac{x + n(-1)^n}{x^2 + n^2}$；
  \item $\displaystyle f(x) = \sum_{n=1}^{\infty} \frac{x}{(1 + x^2)^n}$。
\end{enumerate}
\end{prob}

\begin{prob}[M]{2796}
设 $r_k$ $(k=1,2,\cdots)$ 是闭区间 $[0,1]$ 上的有理数。证明函数
\[
f(x) = \sum_{k=1}^{\infty} \frac{|x - r_k|}{3^k} \quad (0 \leqslant x \leqslant 1)
\]
具有下列性质：1) 连续；2) 在无理点可微而在有理点不可微。
\end{prob}

\begin{prob}[M]{2797}
证明：黎曼 $\zeta$ 函数
\[
\zeta(x) = \sum_{n=1}^{\infty} \frac{1}{n^x}
\]
在区域 $x > 1$ 内是连续的并且在此区域内有连续的各阶导数。
\end{prob}

\begin{prob}[M]{2798}
证明：$\theta$ 函数
\[
\theta(x) = \sum_{n=-\infty}^{+\infty} \mathrm{e}^{-\pi n^2 x}
\]
当 $x > 0$ 时有定义并无穷多次可微。
\end{prob}

\begin{prob}[M]{2799}
确定函数 $f(x)$ 的存在域并研究它们的可微性，设：
(a) $\displaystyle f(x) = \sum_{n=1}^{\infty} \frac{(-1)^n x}{n+x}$； \qquad (b) $\displaystyle f(x) = \sum_{n=1}^{\infty} \frac{|x|}{n^2 + x^2}$。
\end{prob}

\begin{prob}[M]{2800}
证明：序列
\[
f_n(x) = \frac{1}{n}\arctan x^n \quad (n=1,2,\cdots)
\]
在区间 $(-\infty, +\infty)$ 内一致收敛，但
\[
\left[ \lim_{n\to\infty} f_n(x) \right]_{x=1}' \neq \lim_{n\to\infty} f_n'(1).
\]
\end{prob}

\begin{prob}[M]{2801}
证明：序列
\[
f_n(x) = x^2 + \frac{1}{n}\sin n\left(x + \frac{\pi}{2}\right)
\]
在区间 $(-\infty, +\infty)$ 内一致收敛，但
\[
\left[ \lim_{n\to\infty} f_n(x) \right]' \neq \lim_{n\to\infty} f_n'(x).
\]
\end{prob}

\begin{prob}[L]{2802}
试确定参数 $\alpha$ 取何值时下列命题为真：
\begin{enumerate}[label=(\alph*), itemsep=2pt, topsep=3pt]
  \item 序列
  \[
  f_n(x) = n^\alpha x \mathrm{e}^{-nx} \quad (n=1,2,\cdots) \tag{1}
  \]
  $(n=1,2,\cdots)$ 在闭区间 $[0,1]$ 上收敛；
  \item 序列 $(1)$ 在 $[0,1]$ 上一致收敛；
  \item $\displaystyle\lim_{n\to\infty} \int_0^1 f_n(x)\,\mathrm{d}x$ 可在积分号下取极限。
\end{enumerate}
\end{prob}

\begin{prob}[M]{2803}
证明：序列
\[
f_n(x) = nx\mathrm{e}^{-nx^2} \quad (n=1,2,\cdots)
\]
在闭区间 $[0,1]$ 上收敛，但
\[
\int_0^1 \left[ \lim_{n\to\infty} f_n(x) \right]\,\mathrm{d}x \neq \lim_{n\to\infty} \int_0^1 f_n(x)\,\mathrm{d}x.
\]
\end{prob}

\begin{prob}[M]{2804}
证明：序列
\[
f_n(x) = nx(1-x)^n \quad (n=1,2,\cdots)
\]
在闭区间 $[0,1]$ 上收敛而不一致收敛，但
\[
\lim_{n\to\infty} \int_0^1 f_n(x)\,\mathrm{d}x = \int_0^1 \lim_{n\to\infty} f_n(x)\,\mathrm{d}x.
\]
\end{prob}

\begin{prob}[M]{2805}
在表达式
\[
\lim_{n\to\infty} \int_0^1 \frac{nx}{1 + n^2 x^4}\,\mathrm{d}x
\]
中，在积分符号下取极限是否合理？
\end{prob}

\probgroup{求极限：}

\begin{prob}[S]{2806}
$\displaystyle\lim_{x\to 1-0} \sum_{n=1}^{\infty} \frac{(-1)^{n+1}}{n} \cdot \frac{x^n}{x^n + 1}$。
\end{prob}

\begin{prob}[S]{2807}
$\displaystyle\lim_{x\to 1-0} \sum_{n=1}^{\infty} (x^n - x^{n+1})$。
\end{prob}

\begin{prob}[M]{2808}
(a) $\displaystyle\lim_{x\to +0} \sum_{n=1}^{\infty} \frac{1}{2^n n^x}$； \qquad (b) $\displaystyle\lim_{x\to \infty} \sum_{n=1}^{\infty} \frac{x^2}{1 + n^2 x^2}$。
\end{prob}

\begin{prob}[S]{2809}
逐项微分级数 $\displaystyle\sum_{n=1}^{\infty} \arctan\frac{x}{n^2}$ 合理否？
\end{prob}

\begin{prob}[M]{2810}
在闭区间 $[0,1]$ 上逐项积分级数
\[
\sum_{n=1}^{\infty} \left( x^{\frac{1}{2n+1}} - x^{\frac{1}{2n-1}} \right)
\]
合理否？
\end{prob}

\begin{prob}[M]{2811.1}
设 $f(x)$ $(-\infty < x < +\infty)$ 是无穷多次可微的函数，且其导数 $f^{(n)}(x)$ $(n=1,2,\cdots)$ 的序列在每一个有限区间 $(a,b)$ 内一致收敛于函数 $\varphi(x)$。证明：$\varphi(x) = C\mathrm{e}^x$，其中 $C$ 为常数。考察例子 $f_n(x) = \mathrm{e}^{-(x-n)^2}$，$n=1,2,\cdots$。
\end{prob}

\begin{prob}[M]{2811.2}
设函数 $f_n(x)$ $(n=1,2,\cdots)$ 在 $(-\infty, +\infty)$ 上有定义且有界，且在任何闭区间 $[a,b]$ 上 $f_n(x) \rightrightarrows \varphi(x)$。由此是否可知
\[
\lim_{n\to\infty} \sup_x f_n(x) = \sup_x \varphi(x) ?
\]
\end{prob}

\probsec{§5. 幂级数}

\probgroup{求下列幂级数的收敛半径和收敛区间并研究其在收敛区间端点的性质：}

\begin{prob}[S]{2812}
$\displaystyle\sum_{n=1}^{\infty} \frac{x^n}{n^p}$.
\end{prob}

\begin{prob}[S]{2813}
$\displaystyle\sum_{n=1}^{\infty} \frac{3^n + (-2)^n}{n}(x+1)^n$.
\end{prob}

\begin{prob}[S]{2814}
$\displaystyle\sum_{n=1}^{\infty} \frac{(n!)^2}{(2n)!}x^n$.
\end{prob}

\begin{prob}[S]{2815}
$\displaystyle\sum_{n=1}^{\infty} \alpha^{n^2}x^n \ (0 < \alpha < 1)$.
\end{prob}

\begin{prob}[S]{2816}
$\displaystyle\sum_{n=1}^{\infty} \left(1 + \frac{1}{n}\right)^{n^2} x^n$.
\end{prob}

\begin{prob}[S]{2817}
$\displaystyle\sum_{n=1}^{\infty} \frac{n!}{a^{n^2}}x^n \ (a > 1)$.
\end{prob}

\begin{prob}[S]{2818}
$\displaystyle\sum_{n=1}^{\infty} \left[\frac{1 \cdot 3 \cdot 5 \cdots (2n-1)}{2 \cdot 4 \cdot 6 \cdots (2n)}\right]^p \left(\frac{x-1}{2}\right)^n$.
\end{prob}

\begin{prob}[S]{2819}
$\displaystyle\sum_{n=1}^{\infty} (-1)^n \left[\frac{2^n(n!)^2}{(2n+1)!}\right]^p x^n$.
\end{prob}

\begin{prob}[S]{2820}
$\displaystyle\sum_{n=1}^{\infty} \frac{m(m-1)\cdots(m-n+1)}{n!}x^n$.
\end{prob}

\begin{prob}[S]{2821}
$\displaystyle\sum_{n=1}^{\infty} \left(\frac{a^n}{n} + \frac{b^n}{n^2}\right)x^n \ (a > 0, b > 0)$.
\end{prob}

\begin{prob}[S]{2822}
$\displaystyle\sum_{n=1}^{\infty} \frac{x^n}{a^n + b^n} \ (a > 0, b > 0)$.
\end{prob}

\begin{prob}[S]{2823}
$\displaystyle\sum_{n=1}^{\infty} \frac{x^n}{a^{\sqrt{n}}} \ (a > 0)$.
\end{prob}

\begin{prob}[S]{2824}
$\displaystyle\sum_{n=1}^{\infty} \frac{3^{-\sqrt{n}}x^n}{\sqrt{n^2+1}}$.
\end{prob}

\begin{prob}[S]{2825}
$\displaystyle\sum_{n=1}^{\infty} \frac{(2n)!!}{(2n+1)!!}x^n$.
\end{prob}

\begin{prob}[S]{2826}
$\displaystyle\sum_{n=1}^{\infty} \frac{(-1)^n}{n!}\left(\frac{n}{\mathrm{e}}\right)^n x^n$.
\end{prob}

\begin{prob}[S]{2827}
$\displaystyle\sum_{n=1}^{\infty} \left(1 + \frac{1}{2} + \cdots + \frac{1}{n}\right)x^n$.
\end{prob}

\begin{prob}[S]{2828}
$\displaystyle\sum_{n=1}^{\infty} \frac{[3 + (-1)^n]^n}{n}x^n$.
\end{prob}

\begin{prob}[S]{2829}
$\displaystyle\sum_{n=2}^{\infty} \frac{\left(1 + 2\cos\frac{\pi n}{4}\right)^n}{\ln n}x^n$.
\end{prob}

\begin{prob}[S]{2830}
$\displaystyle\sum_{n=1}^{\infty} \frac{x^{n^2}}{2^n}$.
\end{prob}

\begin{prob}[L]{2831}
\begin{enumerate}[label=(\alph*), itemsep=2pt, topsep=3pt]
\item $\displaystyle\sum_{n=1}^{\infty} \frac{(-1)^{[\sqrt{n}]}}{n}x^n$ (普林斯海姆级数).
\item $\displaystyle\sum_{n=1}^{\infty} \frac{10^{\nu(n)}}{n}(1-x)^n$, 其中 $\nu(n)$ 为数 $n$ 的位数;
\item $\displaystyle\sum_{n=1}^{\infty} \left(\frac{x}{\sin n}\right)^n$.
\end{enumerate}
\end{prob}

\begin{prob}[M]{2832}
求超几何级数
\[
1 + \frac{\alpha \cdot \beta}{1 \cdot \gamma}x + \frac{\alpha(\alpha+1)\beta(\beta+1)}{1 \cdot 2 \cdot \gamma(\gamma+1)}x^2 + \cdots + \frac{\alpha(\alpha+1)\cdots(\alpha+n-1)\beta(\beta+1)\cdots(\beta+n-1)}{1 \cdot 2 \cdots n \cdot \gamma(\gamma+1)\cdots(\gamma+n-1)}x^n + \cdots
\]
的收敛域.
\end{prob}

\probgroup{求下列广义幂级数的收敛域：}

\begin{prob}[S]{2833}
$\displaystyle\sum_{n=0}^{\infty} \frac{1}{2n+1}\left(\frac{1-x}{1+x}\right)^n$.
\end{prob}

\begin{prob}[S]{2834}
$\displaystyle\sum_{n=1}^{\infty} \frac{1}{x^n}\sin\frac{\pi}{2^n}$.
\end{prob}

\begin{prob}[S]{2835}
$\displaystyle\sum_{n=-\infty}^{\infty} \frac{x^n}{2^{n^2}}$.
\end{prob}

\begin{prob}[S]{2836}
$\displaystyle\sum_{n=1}^{\infty} \left(1 + \frac{1}{n}\right)^{-n^2} \mathrm{e}^{-nx}$.
\end{prob}

\begin{prob}[S]{2837}
$\displaystyle\sum_{n=1}^{\infty} \frac{3^{3n}(n!)^3}{(3n)!}\tan^n x$.
\end{prob}

\begin{prob}[S]{2838}
按二项式 $x+1$ 的非负整数次幂展开函数
\[
f(x) = x^3.
\]
\end{prob}

\begin{prob}[L]{2839}
把函数
\[
f(x) = \frac{1}{a-x} \quad (a \neq 0)
\]
按以下方式展开为幂级数：(a) 依 $x$ 的幂展开；(b) 依二项式 $x-b$ 的幂展开，此处 $b \neq a$；(c) 依 $\frac{1}{x}$ 的幂展开. 求出相应的收敛域.
\end{prob}

\begin{prob}[M]{2840}
按差 $x-1$ 的非负整数次幂展开函数 $f(x) = \ln x$，并说明展开式的收敛区间. 求级数 $\displaystyle\sum_{n=1}^{\infty} \frac{(-1)^{n+1}}{n}$ 的和.
\end{prob}

\probgroup{写出下列函数按变量 $x$ 的非负整数次幂的展开式，并求出相应的收敛区间：}

\begin{prob}[S]{2841}
$f(x) = \sinh x$.
\end{prob}

\begin{prob}[S]{2842}
$f(x) = \cosh x$.
\end{prob}

\begin{prob}[S]{2843}
$f(x) = \sin^2 x$.
\end{prob}

\begin{prob}[S]{2844}
$f(x) = a^x \ (a > 0, a \neq 1)$.
\end{prob}

\begin{prob}[S]{2845}
$f(x) = \sin(\mu \arcsin x)$.
\end{prob}

\begin{prob}[S]{2846}
$f(x) = \cos(\mu \arcsin x)$.
\end{prob}

\begin{prob}[S]{2847}
写出函数 $f(x) = x^x$ 按差 $x-1$ 的非负整数次幂展开式的前三项.
\end{prob}

\begin{prob}[M]{2848}
写出函数 $f(x) = (1+x)^{\frac{1}{x}} \ (x \neq 0)$ 和 $f(0) = \mathrm{e}$ 按变量 $x$ 的非负整数次幂展开式的前三项.
\end{prob}

\begin{prob}[S]{2849}
将函数 $\sin(x+h)$ 和 $\cos(x+h)$ 按变量 $h$ 的非负整数次幂展开.
\end{prob}

\begin{prob}[M]{2850.1}
不进行实际的展开工作而求函数 $f(x) = \frac{x}{x^2 - 5x + 6}$ 的幂级数展开式的收敛区间：(a) 依 $x$ 的幂展开；(b) 依二项式 $x-5$ 的幂展开.
\end{prob}

\begin{prob}[M]{2850.2}
能否下结论说：在 $(-\infty, +\infty)$ 上当 $N \to \infty$ 时
\[
\sum_{n=1}^{N} (-1)^{n-1}\frac{x^{2n-1}}{(2n-1)!} \rightrightarrows \sin x ?
\]
\end{prob}

\probgroup{利用基本展开式 $\mathrm{I}-\mathrm{V}$，写出下列函数关于 $x$ 的幂级数展开式：}

\begin{prob}[S]{2851}
$\mathrm{e}^{-x^2}$.
\end{prob}

\begin{prob}[S]{2852}
$\cos^2 x$.
\end{prob}

\begin{prob}[S]{2853}
$\sin^3 x$.
\end{prob}

\begin{prob}[S]{2854}
$\frac{x^{10}}{1-x}$.
\end{prob}

\begin{prob}[S]{2855}
$\frac{1}{(1-x)^2}$.
\end{prob}

\begin{prob}[S]{2856}
$\frac{x}{\sqrt{1-2x}}$.
\end{prob}

\begin{prob}[S]{2857}
$\ln\sqrt{\frac{1+x}{1-x}}$.
\end{prob}

\begin{prob}[S]{2858}
$\frac{x}{1+x-2x^2}$.
\probhint{把所给分式分解为简单分式.}
\end{prob}

\begin{prob}[S]{2859}
$\frac{12-5x}{6-5x-x^2}$.
\end{prob}

\begin{prob}[S]{2860}
$\frac{x}{(1-x)(1-x^2)}$.
\end{prob}

\begin{prob}[S]{2861}
$\frac{1}{1-x-x^2}$.
\end{prob}

\begin{prob}[M]{2862}
(a) $\frac{1}{1+x+x^2}$; \qquad (b) $f(x) = \frac{1}{1+x+x^2+x^3}$. 计算 $f^{(1000)}(0)$.
\end{prob}

\begin{prob}[S]{2863}
$\frac{x\cos\alpha - x^2}{1 - 2x\cos\alpha + x^2}$.
\end{prob}

\begin{prob}[S]{2864}
$\frac{x\sin\alpha}{1 - 2x\cos\alpha + x^2}$.
\end{prob}

\begin{prob}[S]{2865}
$\frac{x\sinh\alpha}{1 - 2x\cosh\alpha + x^2}$.
\end{prob}

\begin{prob}[S]{2866}
$\frac{1}{(1-x^2)\sqrt{1-x^2}}$.
\end{prob}

\begin{prob}[S]{2867}
$\ln(1+x+x^2+x^3)$.
\end{prob}

\begin{prob}[S]{2868}
$\mathrm{e}^{x\cos\alpha}\cos(x\sin\alpha)$. \probhint{运用欧拉公式.}
\end{prob}

\probgroup{首先展开导数，然后用逐项积分的方法求下列函数的幂级数展开式.}

\begin{prob}[M]{2869}
$f(x) = \arctan x$. 求级数 $\displaystyle\sum_{n=1}^{\infty} \frac{(-1)^{n+1}}{2n-1}$ 的和.
\end{prob}

\begin{prob}[S]{2870}
$f(x) = \arcsin x$.
\end{prob}

\begin{prob}[S]{2871}
$f(x) = \ln\left(x + \sqrt{1+x^2}\right)$.
\end{prob}

\begin{prob}[S]{2872}
$f(x) = \ln\left(1 - 2x\cos\alpha + x^2\right)$.
\end{prob}

\begin{prob}[L]{2873}
利用不同方法，求下列函数的幂级数展开式：
\begin{enumerate}[label=(\alph*), itemsep=2pt, topsep=3pt]
\item $f(x) = (1+x)\ln(1+x)$;
\item $f(x) = \frac{1}{4}\ln\frac{1+x}{1-x} + \frac{1}{2}\arctan x$;
\item $f(x) = \arctan\frac{2-2x}{1+4x}$;
\item $f(x) = \arctan\frac{2x}{2-x^2}$;
\item $f(x) = x\arctan x - \ln\sqrt{1+x^2}$;
\item $f(x) = \arccos(1-2x^2)$;
\item $f(x) = x\arcsin x + \sqrt{1-x^2}$;
\item $f(x) = x\ln\left(x+\sqrt{1+x^2}\right) - \sqrt{1+x^2}$.
\end{enumerate}
\end{prob}

\begin{prob}[L]{2874}
利用展开式
\[
f(x+h) - f(x) = hf'(x) + \frac{h^2}{2!}f''(x) + \cdots
\]
的唯一性，求下列函数的 $n$ 阶导数：
\begin{enumerate}[label=(\alph*), itemsep=2pt, topsep=3pt]
\item $f(x) = \mathrm{e}^{x^2}$;
\item $f(x) = \mathrm{e}^{\frac{a}{x}}$;
\item $f(x) = \arctan x$.
\end{enumerate}
\end{prob}

\begin{prob}[S]{2875}
依二项式 $x+1$ 的正整数次幂展开函数
\[
f(x) = \ln\frac{1}{2+2x+x^2}.
\]
\end{prob}

\begin{prob}[S]{2876}
把函数
\[
f(x) = \frac{1}{1-x}
\]
按变量 $x$ 的负整数次幂展开成幂级数.
\end{prob}

\begin{prob}[S]{2877}
把函数
\[
f(x) = \ln x
\]
按分式 $\frac{x-1}{x+1}$ 的正整数次幂展开成幂级数.
\end{prob}

\begin{prob}[S]{2878}
把函数
\[
f(x) = \frac{x}{\sqrt{1+x}}
\]
按分式 $\frac{x}{1+x}$ 的正整数次幂展开成幂级数.
\end{prob}

\begin{prob}[M]{2879}
设
\[
f(x) = \sum_{n=0}^{\infty} \frac{x^n}{n!},
\]
直接证明：
\[
f(x)f(y) = f(x+y).
\]
\end{prob}

\begin{prob}[M]{2880}
若我们定义
\[
\sin x = \sum_{n=0}^{\infty} (-1)^n \frac{x^{2n+1}}{(2n+1)!}, \quad \cos x = \sum_{n=0}^{\infty} (-1)^n \frac{x^{2n}}{(2n)!},
\]
证明：(a) $\sin x\cos x = \frac{1}{2}\sin 2x$; (b) $\sin^2 x + \cos^2 x = 1$.
\end{prob}

\begin{prob}[M]{2881}
写出函数 $f(x) = \left[\displaystyle\sum_{n=0}^{\infty} \left(\frac{x^n}{n+1}\right)\right]^{-1}$ 的幂级数展开式中的若干项.
\end{prob}

\probgroup{对幂级数进行相应的运算，从而求出下列函数的幂级数展开式：}

\begin{prob}[S]{2882}
$f(x) = (1+x)\mathrm{e}^{-x}$.
\end{prob}

\begin{prob}[S]{2883}
$f(x) = (1-x)^2\cosh\sqrt{x}$.
\end{prob}

\begin{prob}[S]{2884}
$f(x) = \ln^2(1-x)$.
\end{prob}

\begin{prob}[S]{2885}
$f(x) = (1+x^2)\arctan x$.
\end{prob}

\begin{prob}[S]{2886}
$f(x) = \mathrm{e}^x\cos x$.
\end{prob}

\begin{prob}[S]{2887}
$f(x) = \mathrm{e}^x\sin x$.
\end{prob}

\begin{prob}[S]{2888}
$f(x) = \frac{\ln(1+x)}{1+x}$.
\end{prob}

\begin{prob}[S]{2889}
$f(x) = (\arctan x)^2$.
\end{prob}

\begin{prob}[S]{2890}
$f(x) = \left(\frac{\arcsin x}{x}\right)^2$.
\end{prob}

\probgroup{将下列函数按变量 $x$ 的正整数次幂展开成幂级数，写出展开式 (异于零) 的前三项：}

\begin{prob}[S]{2891}
$f(x) = \tan x$.
\end{prob}

\begin{prob}[S]{2892}
$f(x) = \tanh x$.
\end{prob}

\begin{prob}[S]{2893}
$f(x) = \cot x - \frac{1}{x}$.
\end{prob}

\begin{prob}[M]{2894}
设 $\sec x$ 的展开式写为以下形式：
\[
\sec x = \sum_{n=0}^{\infty} \frac{E_n}{(2n)!}x^{2n},
\]
求出系数 $E_n$ (欧拉数)的递推公式.
\end{prob}

\begin{prob}[S]{2895}
把函数
\[
f(x) = \frac{1}{\sqrt{1 - 2tx + x^2}} \quad (|x| < 1)
\]
展开成幂级数.
\end{prob}

\begin{prob}[M]{2896}
设 $f(x) = \displaystyle\sum_{n=0}^{\infty} a_n x^n$. 写出函数
\[
F(x) = \frac{f(x)}{1-x}
\]
的展开式.
\end{prob}

\begin{prob}[M]{2897}
若级数 $\displaystyle\sum_{n=0}^{\infty} a_n x^n$ 有收敛半径为 $R_1$，而级数 $\displaystyle\sum_{n=0}^{\infty} b_n x^n$ 有收敛半径 $R_2$，则级数：
\[
\text{(a) } \sum_{n=0}^{\infty} (a_n \pm b_n)x^n; \qquad \text{(b) } \sum_{n=0}^{\infty} a_n b_n x^n
\]
的收敛半径 $R$ 是怎样的？
\end{prob}

\begin{prob}[M]{2898}
设
\[
l = \varliminf_{n \to \infty} \left|\frac{a_n}{a_{n+1}}\right| \quad \text{和} \quad L = \varlimsup_{n \to \infty} \left|\frac{a_n}{a_{n+1}}\right|,
\]
证明：幂级数 $\displaystyle\sum_{n=0}^{\infty} a_n x^n$ 的收敛半径 $R$ 满足不等式
\[
l \leqslant R \leqslant L.
\]
\end{prob}

\begin{prob}[M]{2899.1}
证明：若函数 $f(x) = \displaystyle\sum_{n=0}^{\infty} a_n x^n$ 且 $|n!a_n| < M \ (n=1,2,\cdots)$，其中 $M$ 是常数，则：(1) $f(x)$ 在任何点 $a$ 无穷次可微；(2) 下述展开式成立：
\[
f(x) = \sum_{n=0}^{\infty} \frac{f^{(n)}(a)}{n!}(x-a)^n \quad (|x| < +\infty).
\]
\end{prob}

\begin{prob}[M]{2899.2}
设 $f(x) \in C^{(\infty)}(a, b)$，且当 $x \in (a, b)$ 时 $|f^{(n)}(x)| \leqslant c^n$ ($n = 0, 1, 2, \cdots$)，证明：函数 $f(x)$ 可展开为幂级数
\[
f(x) = \sum_{n=0}^{\infty} a_n (x - x_0)^n \quad (x_0 \in (a, b)),
\]
此幂级数在区间 $(a, b)$ 上收敛．
\end{prob}

\begin{prob}[M]{2899.3}
设 $f(x) \in C^{(\infty)}[-1, 1]$，且当 $x \in [-1, 1]$ 时 $f^{(n)}(x) \geqslant 0$ ($n = 0, 1, 2, \cdots$)．证明：函数 $f(x)$ 在区间 $(-1, 1)$ 内可展开为幂级数
\[
f(x) = \sum_{n=0}^{\infty} a_n x^n.
\]
\probhint{对于函数 $f(x)$ 的泰勒级数余项 $R_n(x)$，利用导数 $f^{(n)}(x)$ 的单调性得到以下估计：
\[
|R_n(x)| \leqslant |x|^{n+1} f(1).
\]}
\end{prob}

\begin{prob}[M]{2900}
证明：若 1) $a_n \geqslant 0$；2) 存在
\[
\lim_{x \to R - 0} \sum_{n=0}^{\infty} a_n x^n = S, \quad \text{则 } \sum_{n=0}^{\infty} a_n R^n = S.
\]
\end{prob}

\probgroup{将下列函数展开成幂级数：}

\begin{prob}[S]{2901}
$\displaystyle \int_0^x \mathrm{e}^{-t^2} \,\mathrm{d}t$．
\end{prob}

\begin{prob}[S]{2902}
$\displaystyle \int_0^x \frac{\,\mathrm{d}t}{\sqrt{1 - t^4}}$．
\end{prob}

\begin{prob}[S]{2903}
$\displaystyle \int_0^x \frac{\sin t}{t} \,\mathrm{d}t$．
\end{prob}

\begin{prob}[S]{2904}
$\displaystyle \int_0^x \frac{\arctan t}{t} \,\mathrm{d}t$．
\end{prob}

\begin{prob}[S]{2905}
$\displaystyle \int_0^x \frac{t \,\mathrm{d}t}{\ln(1 + t)}$（写出四项）．
\end{prob}

\probgroup{运用逐项微分法计算下列级数的和：}

\begin{prob}[S]{2906}
$\displaystyle x + \frac{x^3}{3} + \frac{x^5}{5} + \cdots$．
\end{prob}

\begin{prob}[S]{2907}
$\displaystyle x - \frac{x^3}{3} + \frac{x^5}{5} - \cdots$．
\end{prob}

\begin{prob}[S]{2908}
$\displaystyle 1 + \frac{x^2}{2!} + \frac{x^4}{4!} + \cdots$．
\end{prob}

\begin{prob}[S]{2909}
$\displaystyle \frac{x}{1 \cdot 2} + \frac{x^2}{2 \cdot 3} + \frac{x^3}{3 \cdot 4} + \cdots$．
\end{prob}

\begin{prob}[S]{2910}
$\displaystyle 1 + \frac{1}{2} x + \frac{1 \cdot 3}{2 \cdot 4} x^2 + \frac{1 \cdot 3 \cdot 5}{2 \cdot 4 \cdot 6} x^3 + \cdots$．
\probhint{以 $1 - x$ 乘级数的导数．}
\end{prob}

\probgroup{运用逐项积分法计算下列级数的和：}

\begin{prob}[S]{2911}
$x + 2x^2 + 3x^3 + \cdots$．
\end{prob}

\begin{prob}[S]{2912}
$x - 4x^2 + 9x^3 - 16x^4 + \cdots$．
\end{prob}

\begin{prob}[S]{2913}
$1 \cdot 2x + 2 \cdot 3x^2 + 3 \cdot 4x^3 + \cdots$．
\end{prob}

\begin{prob}[M]{2914}
证明：级数
\[
y = \sum_{n=0}^{\infty} \frac{x^{4n}}{(4n)!}
\]
满足方程
\[
y^{(4)} = y.
\]
\end{prob}

\begin{prob}[M]{2915}
证明：级数
\[
y = \sum_{n=0}^{\infty} \frac{x^n}{(n!)^2}
\]
满足方程
\[
x y'' + y' - y = 0.
\]
\end{prob}

\probgroup{求在复域内（$z = x + \mathrm{i}y$）下列幂级数的收敛半径及收敛圆：}

\begin{prob}[S]{2916}
$\displaystyle \sum_{n=1}^{\infty} \frac{(z - 1 - \mathrm{i})^n}{n \cdot 2^n}$．
\end{prob}

\begin{prob}[S]{2917}
$\displaystyle \sum_{n=1}^{\infty} \frac{(1 + \mathrm{i})^n z^n}{(n + 1)(n + 2)}$．
\end{prob}

\begin{prob}[S]{2918}
$\displaystyle \sum_{n=1}^{\infty} \frac{n! z^n}{(1 + \mathrm{i})(1 + 2\mathrm{i}) \cdots (1 + n\mathrm{i})}$．
\end{prob}

\begin{prob}[S]{2919}
$\displaystyle \sum_{n=1}^{\infty} \frac{z^n}{n^{\alpha + \mathrm{i}\beta}}$．
\end{prob}

\begin{prob}[S]{2920}
$\displaystyle \sum_{n=1}^{\infty} \frac{(z - \mathrm{e}^{\mathrm{i}\alpha})^n}{n(1 - \mathrm{e}^{\mathrm{i}\alpha})^n}$．
\end{prob}

\begin{prob}[M]{2921}
利用牛顿二项式公式近似地计算 $\sqrt[3]{9}$，并且估计当只取展开式三项时的误差．
\end{prob}

\begin{prob}[L]{2922}
近似地计算并估计相应误差：
\begin{enumerate}[label=(\alph*), itemsep=2pt, topsep=3pt]
  \item $\arctan 1.2$；
  \item $\sqrt[10]{1000}$；
  \item $\displaystyle \frac{1}{\sqrt{\mathrm{e}}}$；
  \item $\ln 1.25$．
\end{enumerate}
\end{prob}

\probgroup{利用适当的展开式，计算下列函数的值：}

\begin{prob}[S]{2923}
$\sin 18^\circ$，精确到 $10^{-5}$．
\end{prob}

\begin{prob}[S]{2924}
$\cos 1^\circ$，精确到 $10^{-6}$．
\end{prob}

\begin{prob}[S]{2925}
$\tan 9^\circ$，精确到 $10^{-3}$．
\end{prob}

\begin{prob}[S]{2926}
$\mathrm{e}$，精确到 $10^{-6}$．
\end{prob}

\begin{prob}[S]{2927}
$\ln 1.2$，精确到 $10^{-4}$．
\end{prob}

\begin{prob}[M]{2928}
由等式
\[
\frac{\pi}{6} = \arcsin \frac{1}{2}
\]
求数 $\pi$，精确到 $10^{-4}$．
\end{prob}

\begin{prob}[M]{2929}
利用等式
\[
\frac{\pi}{4} = \arctan \frac{1}{2} + \arctan \frac{1}{3},
\]
计算数 $\pi$，精确到 $0.001$．
\end{prob}

\begin{prob}[M]{2930}
利用等式
\[
\frac{\pi}{4} = 4\arctan \frac{1}{5} - \arctan \frac{1}{239},
\]
求数 $\pi$，精确到 $10^{-9}$．
\end{prob}

\begin{prob}[M]{2931}
利用公式
\[
\ln(n + 1) = \ln n + 2\left[\frac{1}{2n + 1} + \frac{1}{3(2n + 1)^3} + \cdots\right],
\]
求 $\ln 2$ 和 $\ln 3$，精确到 $10^{-5}$．
\end{prob}

\begin{prob}[L]{2932}
将被积函数展开成级数，从而计算下列积分之值，精确到 $0.001$：
\begin{enumerate}[label=(\alph*), itemsep=2pt, topsep=3pt]
  \item $\displaystyle \int_0^1 \mathrm{e}^{-x^2} \,\mathrm{d}x$；
  \item $\displaystyle \int_2^4 \mathrm{e}^{\frac{1}{x}} \,\mathrm{d}x$；
  \item $\displaystyle \int_0^2 \frac{\sin x}{x} \,\mathrm{d}x$；
  \item $\displaystyle \int_0^1 \cos x^2 \,\mathrm{d}x$；
  \item $\displaystyle \int_0^1 \frac{\sinh x}{x} \,\mathrm{d}x$；
  \item $\displaystyle \int_2^{+\infty} \frac{\,\mathrm{d}x}{1 + x^3}$；
  \item $\displaystyle \int_0^{\frac{1}{3}} \frac{\,\mathrm{d}x}{\sqrt[3]{1 - x^2}}$；
  \item $\displaystyle \int_0^1 \frac{\,\mathrm{d}x}{\sqrt{1 + x^4}}$；
  \item $\displaystyle \int_{10}^{100} \frac{\ln(1 + x)}{x} \,\mathrm{d}x$；
  \item $\displaystyle \int_0^{\frac{1}{2}} \frac{\arctan x}{x} \,\mathrm{d}x$；
  \item $\displaystyle \int_0^{\frac{1}{2}} \frac{\arcsin x}{x} \,\mathrm{d}x$；
  \item $\displaystyle \int_0^1 x^x \,\mathrm{d}x$．
\end{enumerate}
\end{prob}

\begin{prob}[M]{2933}
求一段正弦曲线
\[
y = \sin x \quad (0 \leqslant x \leqslant \pi)
\]
之弧长，精确到 $0.01$．
\end{prob}

\begin{prob}[M]{2934}
椭圆之半轴为 $a = 1$ 及 $b = \frac{1}{2}$，求椭圆的周长，并精确到 $0.01$．
\end{prob}

\begin{prob}[M]{2935}
两根电线杆相距 $2l = 20\,\mathrm{m}$，电线成抛物线的形状\footnote{严格地说，若不计电线的伸长，电线形状应为悬链线．请读者研究抛物线近似的误差．——译注}．若电线下垂高度 $h = 40\,\mathrm{cm}$，计算电线的长度，并精确到 $1\,\mathrm{cm}$．
\end{prob}

\probsec{§6. 傅里叶级数}

\begin{prob}[M]{2936}
将函数
\[
f(x) = \sin^4 x
\]
展开成傅里叶级数．
\end{prob}

\begin{prob}[M]{2937}
三角多项式
\[
P_n(x) = \sum_{i=0}^{n} (\alpha_i \cos ix + \beta_i \sin ix)
\]
的傅里叶级数是怎样的？
\end{prob}

\begin{prob}[L]{2938}
将函数
\[
f(x) = \operatorname{sgn} x \quad (-\pi < x < \pi)
\]
展开成傅里叶级数．绘出函数的图像及此函数之傅里叶级数之若干部分和的图像．

利用该展开式，求莱布尼茨级数
\[
\sum_{n=1}^{\infty} \frac{(-1)^{n-1}}{2n - 1}.
\]
的和．
\end{prob}

\probgroup{把下列函数展开为傅里叶级数：}

\begin{prob}[M]{2939}
在区间 $(0, 2l)$ 内展开 $f(x) = \begin{cases} A, & 0 < x < l, \\ 0, & l < x < 2l, \end{cases}$ 其中 $A$ 为常数．
\end{prob}

\begin{prob}[S]{2940}
在区间 $(-\pi, \pi)$ 内展开 $f(x) = x$．
\end{prob}

\begin{prob}[S]{2941}
在区间 $(0, 2\pi)$ 内展开 $f(x) = \frac{\pi - x}{2}$．
\end{prob}

\begin{prob}[S]{2942}
在区间 $(-\pi, \pi)$ 内展开 $f(x) = |x|$．
\end{prob}

\begin{prob}[M]{2943}
在区间 $(-\pi, \pi)$ 内展开 $f(x) = \begin{cases} ax, & -\pi < x < 0, \\ bx, & 0 < x < \pi, \end{cases}$ 其中 $a$ 和 $b$ 为常数．
\end{prob}

\begin{prob}[S]{2944}
在区间 $(-\pi, \pi)$ 内展开 $f(x) = \pi^2 - x^2$．
\end{prob}

\begin{prob}[S]{2945}
在区间 $(-\pi, \pi)$ 内展开 $f(x) = \cos ax$（$a$ 不是整数）．
\end{prob}

\begin{prob}[S]{2946}
在区间 $(-\pi, \pi)$ 内展开 $f(x) = \sin ax$（$a$ 不是整数）．
\end{prob}

\begin{prob}[S]{2947}
在区间 $(-\pi, \pi)$ 内展开 $f(x) = \sinh ax$．
\end{prob}

\begin{prob}[S]{2948}
在区间 $(-h, h)$ 内展开 $f(x) = \mathrm{e}^{ax}$．
\end{prob}

\begin{prob}[S]{2949}
在区间 $(a, a + 2l)$ 内展开 $f(x) = x$．
\end{prob}

\begin{prob}[S]{2950}
在区间 $(-\pi, \pi)$ 内展开 $f(x) = x \sin x$．
\end{prob}

\begin{prob}[S]{2951}
在区间 $\left(-\frac{\pi}{2}, \frac{\pi}{2}\right)$ 内展开 $f(x) = x \cos x$．
\end{prob}

\probgroup{将下列周期函数展开成傅里叶级数：}

\begin{prob}[S]{2952}
$f(x) = \operatorname{sgn}(\cos x)$．
\end{prob}

\begin{prob}[S]{2953}
$f(x) = \arcsin(\sin x)$．
\end{prob}

\begin{prob}[S]{2954}
$f(x) = \arcsin(\cos x)$．
\end{prob}

\begin{prob}[S]{2955}
$f(x) = x - [x]$．
\end{prob}

\begin{prob}[S]{2956}
$f(x) = (x)$，其中 $(x)$ 是 $x$ 到与它最近的整数的距离．
\end{prob}

\begin{prob}[S]{2957}
$f(x) = |\sin x|$．
\end{prob}

\begin{prob}[S]{2958}
$f(x) = |\cos x|$．
\end{prob}

\begin{prob}[S]{2959}
$\displaystyle f(x) = \sum_{n=1}^{\infty} \alpha^n \frac{\sin nx}{\sin x}$ ($|\alpha| < 1$)．
\end{prob}

\begin{prob}[M]{2960}
把函数
\[
f(x) = \sec x \quad \left(\frac{-\pi}{4} < x < \frac{\pi}{4}\right)
\]
展开为傅里叶级数．
\probhint{推出系数 $a_n$ 与 $a_{n-2}$ 之间的关系．}
\end{prob}

\begin{prob}[L]{2961}
将函数 $f(x) = x^2$ 展开成傅里叶级数：(a) 在区间 $(-\pi, \pi)$ 内按倍角的余弦展开；(b) 在区间 $(0, \pi)$ 内按倍角的正弦展开；(c) 在区间 $(0, 2\pi)$ 内展开．绘出函数的图像及情形 (a)，(b) 与 (c) 的傅里叶级数之和的图像．

利用这些展开式，求下列级数的和：
\[
\sum_{n=1}^{\infty} \frac{1}{n^2}, \quad \sum_{n=1}^{\infty} \frac{(-1)^{n+1}}{n^2} \quad \text{及} \quad \sum_{n=1}^{\infty} \frac{1}{(2n - 1)^2}.
\]
\end{prob}

\begin{prob}[M]{2962}
由展开式
\[
x = 2\sum_{n=1}^{\infty} (-1)^{n+1} \frac{\sin nx}{n} \quad (-\pi < x < \pi)
\]
用逐项积分的方法，求函数 $x^2, x^3$ 和 $x^4$ 在区间 $(-\pi, \pi)$ 内的傅里叶级数．
\end{prob}

\begin{prob}[L]{2963}
写出函数
\[
f(x) = \begin{cases} 1, & |x| < \alpha, \\ 0, & \alpha < |x| < \pi \end{cases}
\]
的李雅普诺夫等式．利用李雅普诺夫等式，求下列级数之和：
\[
\sum_{n=1}^{\infty} \frac{\sin^2 n\alpha}{n^2} \quad \text{及} \quad \sum_{n=1}^{\infty} \frac{\cos^2 n\alpha}{n^2}.
\]
\end{prob}

\begin{prob}[M]{2964}
将函数
\[
f(x) = \begin{cases}
x, & 0 \leqslant x \leqslant 1, \\
1, & 1 < x < 2, \\
3 - x, & 2 \leqslant x \leqslant 3
\end{cases}
\]
展开成傅里叶级数.
\end{prob}

\probgroup{利用公式
\[
\cos x = \frac{1}{2}(t + \bar{t}), \quad \sin x = \frac{1}{2\mathrm{i}}(t - \bar{t}),
\]
式中 $t = \mathrm{e}^{\mathrm{i}x}$ 及 $\bar{t} = \mathrm{e}^{-\mathrm{i}x}$.}

\begin{prob}[S]{2965}
$\cos^{2m} x$ ($m$ 为正整数).
\end{prob}

\begin{prob}[S]{2966}
$\dfrac{q\sin x}{1 - 2q\cos x + q^2}$ ($|q| < 1$).
\end{prob}

\begin{prob}[S]{2967}
$\dfrac{1 - q^2}{1 - 2q\cos x + q^2}$ ($|q| < 1$).
\end{prob}

\begin{prob}[S]{2968}
$\dfrac{1 - q\cos x}{1 - 2q\cos x + q^2}$ ($|q| < 1$).
\end{prob}

\begin{prob}[S]{2969}
$\ln(1 - 2q\cos x + q^2)$ ($|q| < 1$).
\end{prob}

\probgroup{将下列无界周期函数展开成傅里叶级数：}

\begin{prob}[S]{2970}
$f(x) = \ln\left|\sin\dfrac{x}{2}\right|$.
\end{prob}

\begin{prob}[S]{2971}
$f(x) = \ln\left|\cos\dfrac{x}{2}\right|$.
\end{prob}

\begin{prob}[S]{2972}
$f(x) = \ln\left|\tan\dfrac{x}{2}\right|$.
\end{prob}

\begin{prob}[M]{2973}
将函数
\[
f(x) = \int_0^x \ln\sqrt{\left|\cot\frac{t}{2}\right|}\,\mathrm{d}t \quad (-\pi \leqslant x \leqslant \pi)
\]
展开成傅里叶级数.
\end{prob}

\begin{prob}[M]{2974}
函数
\[
x = x(s), \quad y = y(s) \quad (0 \leqslant s \leqslant 4a)
\]
是正方形 $0 \leqslant x \leqslant a, \, 0 \leqslant y \leqslant a$ 的围线的参数方程，其中 $s$ 为依逆时针方向从点 $O(0,0)$ 起计算的弧长. 试将这两个函数展开成傅里叶级数.
\end{prob}

\begin{prob}[M]{2975}
应当如何把给定在区间 $\left(0, \dfrac{\pi}{2}\right)$ 内的可积函数 $f(x)$ 延拓到区间 $(-\pi, \pi)$ 内，使得它展开成傅里叶级数后具有以下形式：
\[
f(x) = \sum_{n=1}^{\infty} a_n \cos(2n-1)x \quad (-\pi < x < \pi)?
\]
\end{prob}

\begin{prob}[M]{2976}
应当如何把给定在区间 $\left(0, \dfrac{\pi}{2}\right)$ 内的可积函数 $f(x)$ 延拓到区间 $(-\pi, \pi)$ 内，使得它展开成傅里叶级数后具有以下形式：
\[
f(x) = \sum_{n=1}^{\infty} b_n \sin(2n-1)x \quad (-\pi < x < \pi)?
\]
\end{prob}

\begin{prob}[L]{2977}
在区间 $\left(0, \dfrac{\pi}{2}\right)$ 内把函数
\[
f(x) = x\left(\frac{\pi}{2} - x\right)
\]
展开：(a) 依角的奇倍数的余弦展开；(b) 依角的奇倍数的正弦展开. 给出情形 (a) 与 (b) 的傅里叶级数之和的图像.
\end{prob}

\begin{prob}[M]{2978}
设 $f(x)$ 是以 $\pi$ 为周期的反周期函数，即
\[
f(x + \pi) \equiv -f(x).
\]
问此函数在区间 $(-\pi, \pi)$ 内的傅里叶级数具有怎样的特性？
\end{prob}

\begin{prob}[M]{2979}
若 $f(x + \pi) \equiv f(x)$，则函数 $f(x)$ 在区间 $(-\pi, \pi)$ 内的傅里叶级数具有怎样的特性？
\end{prob}

\begin{prob}[M]{2980}
对于一个周期为 $2\pi$ 的函数 $y = f(x)$，若函数的图像：(a) 以点 $(0,0)$，$\left(\pm\dfrac{\pi}{2}, 0\right)$ 为对称中心；(b) 以坐标原点为对称中心并以 $x = \pm\dfrac{\pi}{2}$ 为对称轴，问其傅里叶系数 $a_n, b_n$ ($n = 1, 2, \cdots$) 具有怎样的特性？
\end{prob}

\begin{prob}[M]{2981}
若
\[
\varphi(-x) \equiv \psi(x),
\]
问函数 $\varphi(x)$ 与 $\psi(x)$ 的傅里叶系数 $a_n, b_n$ 与 $\alpha_n, \beta_n$ ($n = 0, 1, 2, \cdots$) 之间有何关系？
\end{prob}

\begin{prob}[M]{2982}
若
\[
\varphi(-x) \equiv -\psi(x),
\]
问函数 $\varphi(x)$ 与 $\psi(x)$ 的傅里叶系数 $a_n, b_n$ 与 $\alpha_n, \beta_n$ ($n = 0, 1, 2, \cdots$) 之间有何关系？
\end{prob}

\begin{prob}[M]{2983}
已知周期为 $2\pi$ 的可积函数 $f(x)$ 的傅里叶系数为 $a_n, b_n$ ($n = 0, 1, 2, \cdots$)，试计算“平移”后的函数 $f(x+h)$ ($h$ 为常数) 的傅里叶系数 $\bar{a}_n, \bar{b}_n$ ($n = 0, 1, 2, \cdots$).
\end{prob}

\begin{prob}[M]{2984}
已知周期为 $2\pi$ 的可积函数 $f(x)$ 的傅里叶系数为 $a_n, b_n$ ($n = 0, 1, 2, \cdots$)，试计算斯捷克洛夫函数
\[
f_h(x) = \frac{1}{2h} \int_{x-h}^{x+h} f(\xi)\,\mathrm{d}\xi
\]
的傅里叶系数 $A_n, B_n$ ($n = 0, 1, 2, \cdots$).
\end{prob}

\begin{prob}[L]{2985}
设 $f(x)$ 是以 $2\pi$ 为周期的连续函数，而 $a_n, b_n$ ($n = 0, 1, 2, \cdots$) 为其傅里叶系数. 求卷积函数
\[
F(x) = \frac{1}{\pi} \int_{-\pi}^{\pi} f(t)f(x+t)\,\mathrm{d}t
\]
的傅里叶系数 $A_n, B_n$ ($n = 0, 1, 2, \cdots$).

利用所得的结果，推出李雅普诺夫等式.
\end{prob}

\probsec{§7. 级数求和法}

\probgroup{求下列级数的和：}

\begin{prob}[S]{2986}
$\dfrac{1}{1\cdot 3} + \dfrac{1}{3\cdot 5} + \dfrac{1}{5\cdot 7} + \cdots$.
\end{prob}

\begin{prob}[S]{2987}
$\dfrac{1}{1\cdot 2\cdot 3} + \dfrac{1}{2\cdot 3\cdot 4} + \dfrac{1}{3\cdot 4\cdot 5} + \cdots$.
\end{prob}

\begin{prob}[S]{2988}
$\dfrac{1}{1\cdot 2} - \dfrac{1}{2\cdot 3} + \dfrac{1}{3\cdot 4} - \dfrac{1}{4\cdot 5} + \cdots$.
\end{prob}

\begin{prob}[S]{2989}
$\sum\limits_{n=1}^{\infty} \dfrac{n}{(n+1)(n+2)(n+3)}$.
\end{prob}

\begin{prob}[S]{2990}
$\sum\limits_{n=1}^{\infty} \dfrac{1}{n(n+m)}$ ($m$ 为正整数).
\end{prob}

\begin{prob}[S]{2991}
$\dfrac{1}{1\cdot 2\cdot 3} + \dfrac{1}{3\cdot 4\cdot 5} + \dfrac{1}{5\cdot 6\cdot 7} + \cdots$.
\end{prob}

\begin{prob}[S]{2992}
$\sum\limits_{n=2}^{\infty} \dfrac{1}{n^2 - 1}$.
\end{prob}

\begin{prob}[S]{2993}
$\sum\limits_{n=1}^{\infty} \dfrac{2n+1}{n^2(n+1)^2}$.
\end{prob}

\begin{prob}[S]{2994}
$\sum\limits_{n=1}^{\infty} \dfrac{1}{n(2n+1)}$.
\end{prob}

\begin{prob}[S]{2995}
$\sum\limits_{n=1}^{\infty} \dfrac{n^2}{n!}$.
\end{prob}

\begin{prob}[S]{2996}
$\sum\limits_{n=0}^{\infty} \dfrac{2^n(n+1)}{n!}$.
\end{prob}

\begin{prob}[S]{2997}
$\sum\limits_{n=1}^{\infty} \dfrac{1}{n^2(n+1)^2}$.
\end{prob}

\begin{prob}[S]{2998}
$\sum\limits_{n=1}^{\infty} \dfrac{1}{n^2(n+1)^2(n+2)^2}$.
\end{prob}

\begin{prob}[S]{2999}
$\sum\limits_{n=0}^{\infty} \dfrac{(-1)^n n}{(2n+1)!}$.
\end{prob}

\begin{prob}[S]{3000}
$\sum\limits_{n=2}^{\infty} \dfrac{(-1)^n}{n^2 + n - 2}$.
\end{prob}

\begin{prob}[M]{3001}
设 $P(x) = a_0 + a_1 x + \cdots + a_m x^m$. 求级数 $\sum\limits_{n=1}^{\infty} \dfrac{P(n)}{n!} x^n$ 的和.
\end{prob}

\probgroup{求下列级数的和：}

\begin{prob}[S]{3002}
$\sum\limits_{n=0}^{\infty} \dfrac{n^2 + 1}{2^n n!} x^n$.
\end{prob}

\begin{prob}[S]{3003}
$\sum\limits_{n=0}^{\infty} \dfrac{(-1)^n n^3}{(n+1)!} x^n$.
\end{prob}

\begin{prob}[S]{3004}
$\sum\limits_{n=0}^{\infty} \dfrac{(-1)^n(2n^2 + 1)}{(2n)!} x^{2n}$.
\end{prob}

\begin{prob}[S]{3005}
$\sum\limits_{n=0}^{\infty} \dfrac{n^2 x^n}{(2n+1)!}$.
\end{prob}

\probgroup{利用逐项微分法求级数的和：}

\begin{prob}[S]{3006}
$\sum\limits_{n=1}^{\infty} \dfrac{x^n}{n}$.
\end{prob}

\begin{prob}[S]{3007}
$\sum\limits_{n=1}^{\infty} \dfrac{(-1)^{n-1} x^{2n}}{n(2n-1)}$.
\end{prob}

\begin{prob}[S]{3008}
$\sum\limits_{n=0}^{\infty} \dfrac{x^{4n+1}}{4n+1}$.
\end{prob}

\begin{prob}[M]{3009}
$\sum\limits_{n=1}^{\infty} \dfrac{a(a+d)\cdots [a+(n-1)d]}{d \cdot 2d \cdots nd} x^n$ ($d > 0$). \probhint{用 $1-x$ 去乘级数的导数.}
\end{prob}

\begin{prob}[S]{3010}
$\dfrac{1}{3} \cdot \dfrac{x}{2} + \dfrac{1\cdot 4}{3\cdot 6} \cdot \left(\dfrac{x}{2}\right)^2 + \dfrac{1\cdot 4\cdot 7}{3\cdot 6\cdot 9} \cdot \left(\dfrac{x}{2}\right)^3 + \cdots$.
\end{prob}

\probgroup{利用逐项积分法求级数的和：}

\begin{prob}[S]{3011}
$\sum\limits_{n=1}^{\infty} n^2 x^{n-1}$.
\end{prob}

\begin{prob}[S]{3012}
$\sum\limits_{n=1}^{\infty} n(n+2) x^n$.
\end{prob}

\begin{prob}[S]{3013}
$\sum\limits_{n=0}^{\infty} \dfrac{(2n+1)x^{2n}}{n!}$.
\end{prob}

\probgroup{利用阿贝尔方法，求下列级数的和：}

\begin{prob}[S]{3014}
$1 - \dfrac{1}{4} + \dfrac{1}{7} - \dfrac{1}{10} + \cdots$.
\end{prob}

\begin{prob}[S]{3015}
$1 - \dfrac{1}{3} + \dfrac{1}{5} - \dfrac{1}{7} + \cdots$.
\end{prob}

\begin{prob}[S]{3016}
$1 - \dfrac{1}{2} + \dfrac{1\cdot 3}{2\cdot 4} - \dfrac{1\cdot 3\cdot 5}{2\cdot 4\cdot 6} + \cdots$.
\end{prob}

\begin{prob}[S]{3017}
$1 + \dfrac{1}{2} \cdot \dfrac{1}{3} + \dfrac{1\cdot 3}{2\cdot 4} \cdot \dfrac{1}{5} + \cdots$.
\end{prob}

\probgroup{求下列三角级数的和：}

\begin{prob}[S]{3018}
$\sum\limits_{n=1}^{\infty} \dfrac{\sin nx}{n}$.
\end{prob}

\begin{prob}[S]{3019}
$\sum\limits_{n=1}^{\infty} \dfrac{\cos nx}{n}$.
\end{prob}

\begin{prob}[S]{3020}
$\sum\limits_{n=1}^{\infty} \dfrac{\sin n\alpha \sin nx}{n}$.
\end{prob}

\begin{prob}[S]{3021}
$\sum\limits_{n=1}^{\infty} \dfrac{\sin^2 n\alpha \sin nx}{n}$ $\left(0 < \alpha < \dfrac{\pi}{2}\right)$.
\end{prob}

\begin{prob}[S]{3022}
$\sum\limits_{n=1}^{\infty} \dfrac{\sin(2n-1)x}{2n-1}$.
\end{prob}

\begin{prob}[S]{3023}
$\sum\limits_{n=2}^{\infty} (-1)^n \dfrac{\cos nx}{n^2 - 1}$.
\end{prob}

\begin{prob}[S]{3024}
$\sum\limits_{n=1}^{\infty} \dfrac{\cos(2n-1)x}{(2n-1)^2}$.
\end{prob}

\begin{prob}[S]{3025}
$\sum\limits_{n=1}^{\infty} (-1)^{n-1} \dfrac{\sin nx}{n(n+1)}$.
\end{prob}

\begin{prob}[S]{3026}
$\sum\limits_{n=0}^{\infty} \dfrac{\cos nx}{n!}$.
\end{prob}

\begin{prob}[M]{3027}
画出曲线 $\sum\limits_{n=1}^{\infty} \dfrac{\sin nx \cdot \sin ny}{n^2} = 0$.
\end{prob}

\probgroup{求下列级数的和：}

\begin{prob}[S]{3028}
$\sum\limits_{n=1}^{\infty} \dfrac{[(n-1)!]^2}{(2n)!} (2x)^{2n}$.
\end{prob}

\begin{prob}[S]{3029}
$\sum\limits_{n=0}^{\infty} \dfrac{(n!)^2}{(2n)!} x^n$.
\end{prob}

\begin{prob}[S]{3030}
$\dfrac{1!}{x+1} + \dfrac{2!}{(x+1)(x+2)} + \dfrac{3!}{(x+1)(x+2)(x+3)} + \cdots$.
\end{prob}

\begin{prob}[M]{3031}
$\dfrac{a_1}{a_2 + x} + \dfrac{a_1}{a_2 + x} \cdot \dfrac{a_2}{a_3 + x} + \cdots$，在 $x > 0, \, a_n > 0$ ($n = 1, 2, \cdots$)，且级数 $\sum\limits_{n=1}^{\infty} \dfrac{1}{a_n}$ 发散的条件下.
\end{prob}

\begin{prob}[M]{3032}
$\dfrac{x}{1-x^2} + \dfrac{x^2}{1-x^4} + \dfrac{x^4}{1-x^8} + \cdots$，若 (a) $|x| < 1$；(b) $|x| > 1$.
\end{prob}

\begin{prob}[M]{3033}
$\sum\limits_{n=1}^{\infty} \dfrac{x^{n+1}}{(1-x^n)(1-x^{n+1})}$，若 (a) $|x| < 1$；(b) $|x| > 1$.
\end{prob}

\probsec{§8. 利用级数求定积分}

\probgroup{利用被积函数的级数展开式计算下列积分：}

\begin{prob}[S]{3034}
$\int_0^1 \ln\dfrac{1}{1-x}\,\mathrm{d}x$.
\end{prob}

\begin{prob}[S]{3035}
$\int_0^1 \dfrac{\ln(x + \sqrt{1+x^2})}{x}\,\mathrm{d}x$.
\end{prob}

\begin{prob}[S]{3036}
$\int_0^1 \dfrac{\ln(1+x)}{x}\,\mathrm{d}x$.
\end{prob}

\begin{prob}[S]{3037}
$\int_0^1 x^{p-1} \ln(1 - x^q)\,\mathrm{d}x$ ($p > 0, \, q > 0$).
\end{prob}

\begin{prob}[S]{3038}
$\int_0^1 \ln x \cdot \ln(1-x)\,\mathrm{d}x$.
\end{prob}

\begin{prob}[S]{3039}
$\int_0^{+\infty} \dfrac{x\,\mathrm{d}x}{\mathrm{e}^{2\pi x} - 1}$.
\end{prob}

\begin{prob}[S]{3040}
$\int_0^{+\infty} \dfrac{x\,\mathrm{d}x}{\mathrm{e}^x + 1}$.
\end{prob}

\begin{prob}[M]{3041}
按模数 $k$ ($0 \leqslant k < 1$) 的正整数次幂展开第一类完全椭圆积分
\[
F(k) = \int_0^{\frac{\pi}{2}} \frac{\mathrm{d}\varphi}{\sqrt{1 - k^2 \sin^2\varphi}}.
\]
\end{prob}

\begin{prob}[M]{3042}
按模数 $k$ ($0 \leqslant k < 1$) 的正整数次幂展开第二类完全椭圆积分
\[
E(k) = \int_0^{\frac{\pi}{2}} \sqrt{1 - k^2 \sin^2\varphi}\,\mathrm{d}\varphi.
\]
\end{prob}

\begin{prob}[M]{3043}
利用按椭圆离心率的正整数次幂展开的级数表示椭圆
\[
x = a\cos t, \quad x = b\sin t \quad (0 \leqslant t \leqslant 2\pi)
\]
的弧长.
\end{prob}

\probgroup{证明下列等式：}

\begin{prob}[S]{3044}
$\displaystyle \int_{0}^{1} \frac{\mathrm{d}x}{x^{x}} = \sum_{n=1}^{\infty} \frac{1}{n^{n}}$.
\end{prob}

\begin{prob}[S]{3045}
$\displaystyle \int_{0}^{+\infty} \mathrm{e}^{-x^{2}} \sin ax\,\mathrm{d}x = \frac{1}{2} \sum_{n=0}^{\infty} \frac{(-1)^{n} n!}{(2n+1)!} a^{2n+1}$.
\end{prob}

\begin{prob}[S]{3046}
$\displaystyle \int_{0}^{2\pi} \mathrm{e}^{\cos x} \cos(\sin x) \cos nx\,\mathrm{d}x = \frac{\pi}{n!} \ (n = 0, 1, 2, \cdots)$.
\end{prob}

\probgroup{求：}

\begin{prob}[S]{3047}
$\displaystyle \int_{0}^{2\pi} \mathrm{e}^{a\cos x} \cos(a\sin x - nx)\,\mathrm{d}x$ ($n$ 是正整数).
\end{prob}

\begin{prob}[M]{3048}
$\displaystyle \int_{0}^{\pi} \frac{x\sin x}{1 - 2\alpha\cos x + \alpha^{2}}\,\mathrm{d}x$. \probhint{参考习题 2864.}
\end{prob}

\begin{prob}[S]{3049}
$\displaystyle \int_{0}^{\pi} \ln(1 - 2\alpha\cos x + \alpha^{2})\,\mathrm{d}x$.
\end{prob}

\begin{prob}[M]{3050}
证明公式：
\[
\int_{0}^{+\infty} \frac{\mathrm{e}^{-x}}{a+x}\,\mathrm{d}x = \frac{1}{a} - \frac{1!}{a^{2}} + \frac{2!}{a^{3}} - \cdots + (-1)^{n-1} \frac{(n-1)!}{a^{n}} + (-1)^{n} \frac{\theta_{n} n!}{a^{n+1}}, \qquad (1)
\]
其中 $a > 0$ 且 $0 < \theta_{n} < 1$. 若于公式 (1) 中取两项来表示积分
\[
\int_{0}^{+\infty} \frac{\mathrm{e}^{-x}}{100+x}\,\mathrm{d}x,
\]
其精确程度如何？
\end{prob}

\probsec{§9. 无穷乘积}

\probgroup{证明下列等式：}

\begin{prob}[S]{3051}
$\displaystyle \prod_{n=2}^{\infty} \left(1 - \frac{1}{n^{2}}\right) = \frac{1}{2}$.
\end{prob}

\begin{prob}[S]{3052}
$\displaystyle \prod_{n=2}^{\infty} \frac{n^{3}-1}{n^{3}+1} = \frac{2}{3}$.
\end{prob}

\begin{prob}[S]{3053}
$\displaystyle \prod_{n=2}^{\infty} \left[1 - \frac{2}{n(n+1)}\right] = \frac{1}{3}$.
\end{prob}

\begin{prob}[S]{3054}
$\displaystyle \prod_{n=0}^{\infty} \left[1 + \left(\frac{1}{2}\right)^{2^{n}}\right] = 2$.
\end{prob}

\begin{prob}[S]{3055}
$\displaystyle \prod_{n=1}^{\infty} \cos \frac{\pi}{2^{n+1}} = \frac{2}{\pi}$.
\end{prob}

\begin{prob}[S]{3056}
$\displaystyle \prod_{n=1}^{\infty} \cos \frac{x}{2^{n}} = \frac{\sin x}{x}$.
\end{prob}

\begin{prob}[S]{3057}
$\displaystyle \prod_{n=1}^{\infty} \cosh \frac{x}{2^{n}} = \frac{\sinh x}{x}$.
\end{prob}

\begin{prob}[S]{3058}
$\displaystyle \prod_{n=0}^{\infty} \left(1 + x^{2^{n}}\right) = \frac{1}{1-x} \ (|x| < 1)$.
\end{prob}

\begin{prob}[S]{3059}
$\displaystyle \frac{\pi}{2} = \frac{2}{\sqrt{2}} \cdot \frac{2}{\sqrt{2+\sqrt{2}}} \cdot \frac{2}{\sqrt{2+\sqrt{2+\sqrt{2}}}} \cdots$.
\end{prob}

\begin{prob}[S]{3060}
$\displaystyle \prod_{n=1}^{\infty} \frac{3n}{3n-1} \cdot \frac{3n}{3n+1} = \frac{2\pi}{3\sqrt{3}}$.
\end{prob}

\probgroup{试证下列无穷乘积的收敛性并求出其值：}

\begin{prob}[S]{3061}
$\displaystyle \prod_{n=3}^{\infty} \frac{n^{2}-4}{n^{2}-1}$.
\end{prob}

\begin{prob}[S]{3062}
$\displaystyle \prod_{n=1}^{\infty} \left[1 + \frac{1}{n(n+2)}\right]$.
\end{prob}

\begin{prob}[S]{3063}
$\displaystyle \prod_{n=1}^{\infty} \frac{(2n+1)(2n+7)}{(2n+3)(2n+5)}$.
\end{prob}

\begin{prob}[S]{3064}
$\displaystyle \prod_{n=1}^{\infty} a^{\frac{(-1)^{n}}{n}} \ (a > 0)$.
\end{prob}

\begin{prob}[L]{3065}
可否由乘积 $\displaystyle \prod_{n=1}^{\infty} p_{n}$ 与 $\displaystyle \prod_{n=1}^{\infty} q_{n}$ 的收敛性得出下列乘积：
\begin{enumerate}[label=(\alph*), itemsep=2pt, topsep=3pt]
\item $\displaystyle \prod_{n=1}^{\infty} (p_{n} + q_{n})$;
\item $\displaystyle \prod_{n=1}^{\infty} p_{n}^{2}$;
\item $\displaystyle \prod_{n=1}^{\infty} p_{n} q_{n}$;
\item $\displaystyle \prod_{n=1}^{\infty} \frac{p_{n}}{q_{n}}$
\end{enumerate}
的收敛性？
\end{prob}

\probgroup{研究下列无穷乘积的收敛性：}

\begin{prob}[S]{3066}
$\displaystyle \prod_{n=1}^{\infty} \frac{1}{n}$.
\end{prob}

\begin{prob}[S]{3067}
$\displaystyle \prod_{n=1}^{\infty} \frac{(n+1)^{2}}{n(n+2)}$.
\end{prob}

\begin{prob}[S]{3068}
$\displaystyle \prod_{n=1}^{\infty} \left(1 + \frac{1}{n^{p}}\right)$.
\end{prob}

\begin{prob}[S]{3069}
$\displaystyle \prod_{n=2}^{\infty} \left(1 - \frac{1}{n}\right)$.
\end{prob}

\begin{prob}[S]{3070}
$\displaystyle \prod_{n=2}^{\infty} \left(\frac{n^{2}-1}{n^{2}+1}\right)^{p}$.
\end{prob}

\begin{prob}[M]{3071}
$\displaystyle \prod_{n=n_{0}}^{\infty} \frac{n^{2} + a_{1}n + b_{1}}{n^{2} + an + b}$，其中当 $n \geqslant n_{0}$ 时 $n^{2} + an + b > 0$.
\end{prob}

\begin{prob}[M]{3072}
$\displaystyle \prod_{n=n_{0}}^{\infty} \frac{(n-a_{1})(n-a_{2})\cdots(n-a_{p})}{(n-b_{1})(n-b_{2})\cdots(n-b_{p})}$，其中 $n_{0} > b_{i} \ (i = 1, 2, \cdots, p)$.
\end{prob}

\begin{prob}[S]{3073}
$\displaystyle \prod_{n=0}^{\infty} \sqrt{\frac{n+1}{n+2}}$.
\end{prob}

\begin{prob}[S]{3074}
$\displaystyle \prod_{n=1}^{\infty} \frac{n}{\sqrt{n^{2}+1}}$.
\end{prob}

\begin{prob}[S]{3075}
$\displaystyle \prod_{n=1}^{\infty} \sqrt[n]{1 + \frac{1}{n}}$.
\end{prob}

\begin{prob}[S]{3076}
$\displaystyle \prod_{n=1}^{\infty} \sqrt[n^{2}]{n}$.
\end{prob}

\begin{prob}[S]{3077}
$\displaystyle \prod_{n=1}^{\infty} \left(1 + \frac{x}{n}\right) \mathrm{e}^{-\frac{x}{n}}$.
\end{prob}

\begin{prob}[M]{3078}
$\displaystyle \prod_{n=1}^{\infty} \left(1 - \frac{x}{c+n}\right) \mathrm{e}^{\frac{x}{n}}$，其中 $c > 0$.
\end{prob}

\begin{prob}[S]{3079}
$\displaystyle \prod_{n=1}^{\infty} (1 - x^{n})$.
\end{prob}

\begin{prob}[S]{3080}
$\displaystyle \prod_{n=1}^{\infty} \left(1 + \frac{x^{n}}{2^{n}}\right)$.
\end{prob}

\begin{prob}[S]{3081}
$\displaystyle \prod_{n=1}^{\infty} \left[1 + \frac{\left(1 + \frac{1}{n}\right)^{n^{2}}}{x^{n}}\right]$.
\end{prob}

\begin{prob}[S]{3082}
$\displaystyle \prod_{n=1}^{\infty} \left(1 - \frac{x}{\sqrt{n}}\right) \mathrm{e}^{\frac{x}{\sqrt{n}} + \frac{x^{2}}{2n}}$.
\end{prob}

\begin{prob}[S]{3083}
$\displaystyle \prod_{n=1}^{\infty} \left(1 + \frac{x^{n}}{n^{p}}\right) \cos \frac{x^{n}}{n^{q}}$.
\end{prob}

\begin{prob}[S]{3084}
$\displaystyle \prod_{n=1}^{\infty} \left(\frac{\sin \frac{x}{n}}{\frac{x}{n}}\right)^{p}$.
\end{prob}

\begin{prob}[S]{3085}
$\displaystyle \prod_{n=1}^{\infty} \sqrt[n]{\ln(n+x) - \ln n}$.
\end{prob}

\begin{prob}[M]{3086}
证明：若级数 $\displaystyle \sum_{n=1}^{\infty} x_{n}^{2}$ 收敛，则乘积 $\displaystyle \prod_{n=1}^{\infty} \cos x_{n}$ 收敛.
\end{prob}

\begin{prob}[M]{3087}
证明：若级数 $\displaystyle \sum_{n=1}^{\infty} \alpha_{n}$ 绝对收敛，则乘积 $\displaystyle \prod_{n=1}^{\infty} \tan\left(\frac{\pi}{4} + \alpha_{n}\right) \ \left(|\alpha_{n}| < \frac{\pi}{4}\right)$ 收敛.
\end{prob}

\probgroup{研究下列无穷乘积的绝对收敛性和条件收敛性：}

\begin{prob}[S]{3088}
$\displaystyle \prod_{n=1}^{\infty} \left[1 + \frac{(-1)^{n+1}}{n}\right]$.
\end{prob}

\begin{prob}[S]{3089}
$\displaystyle \prod_{n=1}^{\infty} \left[1 + \frac{(-1)^{n+1}}{\sqrt{n}}\right]$.
\end{prob}

\begin{prob}[S]{3090}
$\displaystyle \prod_{n=1}^{\infty} \left[1 + \frac{(-1)^{n+1}}{n^{p}}\right]$.
\end{prob}

\begin{prob}[S]{3091}
$\displaystyle \prod_{n=2}^{\infty} \left[1 + \frac{(-1)^{n}}{\ln n}\right]$.
\end{prob}

\begin{prob}[S]{3092}
$\displaystyle \prod_{n=2}^{\infty} \frac{\sqrt{n}}{\sqrt{n} + (-1)^{n}}$.
\end{prob}

\begin{prob}[S]{3093}
$\displaystyle \prod_{n=1}^{\infty} n^{(-1)^{n}}$.
\end{prob}

\begin{prob}[S]{3094}
$\displaystyle \prod_{n=1}^{\infty} \sqrt[n]{n^{(-1)^{n}}}$.
\end{prob}

\begin{prob}[S]{3095}
$\displaystyle \prod_{n=1}^{\infty} \left[1 + \frac{(-1)^{\frac{n(n-1)}{2}}}{n}\right]$.
\end{prob}

\begin{prob}[M]{3096}
$\displaystyle \left(1 + \frac{1}{\sqrt{1}}\right)\left(1 - \frac{1}{\sqrt{3}}\right)\left(1 - \frac{1}{\sqrt{5}}\right)\left(1 + \frac{1}{\sqrt{2}}\right)\left(1 - \frac{1}{\sqrt{7}}\right)\left(1 - \frac{1}{\sqrt{9}}\right)\left(1 + \frac{1}{\sqrt{3}}\right)\cdots$.
\end{prob}

\begin{prob}[M]{3097}
$\displaystyle \left(1 + \frac{1}{1^{\alpha}}\right)\left(1 - \frac{1}{2^{\alpha}}\right)^{2}\left(1 + \frac{1}{3^{\alpha}}\right)\left(1 + \frac{1}{4^{\alpha}}\right)\left(1 - \frac{1}{5^{\alpha}}\right)^{2}\left(1 + \frac{1}{6^{\alpha}}\right)\cdots$.
\end{prob}

\begin{prob}[M]{3098}
证明：乘积
\[
\left(1 + \frac{1}{\sqrt{2}} + \frac{1}{2}\right)\left(1 - \frac{1}{\sqrt{2}}\right)\left(1 + \frac{1}{\sqrt{3}} + \frac{1}{3}\right)\left(1 - \frac{1}{\sqrt{3}}\right)\cdots
\]
收敛，尽管级数
\[
\left(\frac{1}{\sqrt{2}} + \frac{1}{2}\right) + \left(-\frac{1}{\sqrt{2}}\right) + \left(\frac{1}{\sqrt{3}} + \frac{1}{3}\right) + \left(-\frac{1}{\sqrt{3}}\right) + \cdots
\]
发散.
\end{prob}

\begin{prob}[M]{3099}
证明：乘积 $\displaystyle \prod_{n=1}^{\infty} (1 + \alpha_{n})$ 收敛，其中
\[
\alpha_{n} = \begin{cases}
\dfrac{-1}{\sqrt{k}}, & n = 2k-1, \\
\dfrac{1}{\sqrt{k}} + \dfrac{1}{k} + \dfrac{1}{k\sqrt{k}}, & n = 2k,
\end{cases}
\]
尽管级数 $\displaystyle \sum_{n=1}^{\infty} \alpha_{n}$ 和 $\displaystyle \sum_{n=1}^{\infty} \alpha_{n}^{2}$ 二者发散.
\end{prob}

\begin{prob}[M]{3100}
设
\[
\zeta(x) = \sum_{n=1}^{\infty} \frac{1}{n^{x}}
\]
（黎曼 $\zeta$ 函数），而 $p_{n} \ (n = 1, 2, \cdots)$ 是素数数列. 证明：
\[
\prod_{n=1}^{\infty} \left(1 - \frac{1}{p_{n}^{x}}\right)^{-1} = \zeta(x).
\]
\end{prob}

\begin{prob}[M]{3101}
设 $p_{n} \ (n = 1, 2, \cdots)$ 是素数数列，证明：乘积
\[
\prod_{n=1}^{\infty} \left(1 - \frac{1}{p_{n}}\right)^{-1}
\]
和级数
\[
\sum_{n=1}^{\infty} \frac{1}{p_{n}}
\]
发散（欧拉）.
\end{prob}

\begin{prob}[M]{3102}
设 $a_{n} > 0 \ (n = 1, 2, \cdots)$，且
\[
\frac{a_{n}}{a_{n+1}} = 1 + \frac{p}{n} + O\left(\frac{1}{n^{1+\varepsilon}}\right) \quad (\varepsilon > 0),
\]
证明：
\[
a_{n} = O^{*}\left(\frac{1}{n^{p}}\right).
\]
\probhint{研究
\[
\lim_{n\to\infty} a_{n} n^{p} = a_{1} \prod_{n=1}^{\infty} \frac{a_{n+1}}{a_{n}} \left(1 + \frac{1}{n}\right)^{p}.
\]}
\end{prob}

\begin{prob}[S]{3103}
利用沃利斯公式证明：
\[
\frac{1 \cdot 3 \cdot 5 \cdots (2n-1)}{2 \cdot 4 \cdot 6 \cdots (2n)} \sim \frac{1}{\sqrt{\pi n}}.
\]
\end{prob}

\begin{prob}[L]{3104}
证明：表达式
\[
a_{n} = \frac{n!\mathrm{e}^{n}}{n^{n+\frac{1}{2}}}
\]
当 $n \to \infty$ 时有异于零的极限 $A$. 由此推出斯特林公式
\[
n! = A n^{n+\frac{1}{2}} \mathrm{e}^{-n} (1 + \varepsilon_{n}),
\]
其中 $\lim_{n\to\infty} \varepsilon_{n} = 0$，$A = \sqrt{2\pi}$.
\probhint{把所求极限表示为无穷乘积的形式：
\[
A = \lim_{n\to\infty} a_{n} = a_{1} \prod_{n=1}^{\infty} \frac{a_{n+1}}{a_{n}}.
\]
为了确定常数 $A$，可利用沃利斯公式.}
\end{prob}

\begin{prob}[L]{3105}
根据欧拉的定义，$\Gamma$ 函数 $\Gamma(x)$ 由下面的公式来确定：
\[
\Gamma(x) = \lim_{n\to\infty} \frac{n! n^{x}}{x(x+1)\cdots(x+n)}.
\]
由这个公式出发：
\begin{enumerate}[label=(\alph*), itemsep=2pt, topsep=3pt]
\item 把函数 $\Gamma(x)$ 表示为无穷乘积的形式；
\item 证明：$\Gamma(x)$ 对于不为负整数的一切实数 $x$ 皆有意义；
\item 推出下面这个性质：$\Gamma(x+1) = x\Gamma(x)$；
\item 对于正整数 $n$ 求 $\Gamma(n)$ 之值.
\end{enumerate}
\end{prob}

\begin{prob}[M]{3106}
设函数 $f(x)$ 在闭区间 $[a, b]$ 上可积，且
\[
\delta_{n} = \frac{b-a}{n}, \quad f_{in} = f(a + i\delta_{n}) \quad (i = 1, 2, \cdots, n),
\]
证明：
\[
\lim_{n\to\infty} \prod_{i=1}^{n} (1 + \delta_{n} f_{in}) = \mathrm{e}^{\int_{a}^{b} f(x)\,\mathrm{d}x}.
\]
\end{prob}

\begin{prob}[M]{3107}
证明:
\[
\lim_{n \to \infty} \frac{\sqrt[n]{\displaystyle\sum_{i=0}^{n-1} (a + ib)}}{\displaystyle\sum_{i=0}^{n-1} (a + ib)} n = \frac{2}{\mathrm{e}},
\]
其中 $a > 0$ 和 $b > 0$.
\end{prob}

\begin{prob}[M]{3108}
设 $f_n(x)$ ($n = 1, 2, \cdots$) 在区间 $(a, b)$ 内为连续函数且 $|f_n(x)| \leqslant c_n$ ($n = 1, 2, \cdots$), 其中级数 $\sum_{n=1}^{\infty} c_n$ 收敛. 证明: 函数
\[
F(x) = \prod_{n=1}^{\infty} [1 + f_n(x)] \quad (|f_n(x)| < 1)
\]
在区间 $(a, b)$ 上是连续的.
\end{prob}

\begin{prob}[M]{3109}
求函数
\[
F(x) = \prod_{n=1}^{\infty} [1 + f_n(x)]
\]
的导数之表达式. $F'(x)$ 存在的充分条件为何?
\end{prob}

\begin{prob}[M]{3110}
证明: 若 $0 < x < y$, 则
\[
\lim_{n \to \infty} \frac{x(x+1)\cdots(x+n)}{y(y+1)\cdots(y+n)} = 0.
\]
\end{prob}

\probsec{§10. 斯特林公式}

\probgroup{利用斯特林公式，近似地计算：}

\begin{prob}[S]{3111}
$\lg 100!$.
\end{prob}

\begin{prob}[S]{3112}
$1 \cdot 3 \cdot 5 \cdots 1999$.
\end{prob}

\begin{prob}[S]{3113}
$\dfrac{1 \cdot 3 \cdot 5 \cdots 99}{2 \cdot 4 \cdot 6 \cdots 100}$.
\end{prob}

\begin{prob}[S]{3114}
$\mathrm{C}_{100}^{40}$.
\end{prob}

\begin{prob}[S]{3115}
$\dfrac{100!}{20!30!50!}$.
\end{prob}

\begin{prob}[S]{3116}
$\displaystyle\int_{0}^{1} (1 - x^2)^{50} \,\mathrm{d}x$.
\end{prob}

\begin{prob}[S]{3117}
$\displaystyle\int_{0}^{2\pi} \sin^{200} x \,\mathrm{d}x$.
\end{prob}

\begin{prob}[M]{3118}
推出乘积
\[
(2n - 1)!! = 1 \cdot 3 \cdot 5 \cdots (2n - 1)
\]
的渐近公式.
\end{prob}

\begin{prob}[S]{3119}
若 $n$ 甚大，近似地计算 $\mathrm{C}_{2n}^{n}$.
\end{prob}

\begin{prob}[L]{3120}
利用斯特林公式求下列极限:
\begin{enumerate}[label=(\alph*), itemsep=2pt, topsep=3pt]
\item $\displaystyle\lim_{n \to \infty} \sqrt[n^2]{n!}$;
\item $\displaystyle\lim_{n \to \infty} \frac{n}{\sqrt[n]{n!}}$;
\item $\displaystyle\lim_{n \to \infty} \frac{n}{\sqrt[n]{(2n-1)!!}}$;
\item $\displaystyle\lim_{n \to \infty} \frac{\ln n!}{\ln n^n}$.
\end{enumerate}
\end{prob}

\probsec{§11. 用多项式逼近连续函数}

\begin{prob}[M]{3121}
求在给定点按下表取值的最低次的 $n$ 次多项式 $P_n(x)$:
\begin{center}
\begin{tabular}{|c|c|c|c|c|}
\hline
$x$ & $-2$ & $0$ & $4$ & $5$ \\
\hline
$y$ & $5$ & $1$ & $-3$ & $1$ \\
\hline
\end{tabular}
\end{center}
$P_n(-1), P_n(1), P_n(6)$ 近似地等于什么?
\end{prob}

\begin{prob}[M]{3122}
写出经过 $A(x_0 - h, y_{-1}), B(x_0, y_0), C(x_0 + h, y_1)$ 这三点的抛物线方程
\[
y = ax^2 + bx + c.
\]
\end{prob}

\begin{prob}[M]{3123}
利用数值 $x_0 = 1, y_0 = 1$; $x_1 = 25, y_1 = 5$; $x_2 = 100, y_2 = 10$ 推出开平方根 $y = \sqrt{x}$ ($1 \leqslant x \leqslant 100$) 的近似公式.
\end{prob}

\begin{prob}[M]{3124}
利用数值
\[
\sin 0^\circ = 0, \quad \sin 30^\circ = \frac{1}{2}, \quad \sin 90^\circ = 1
\]
推出如下形式的近似公式:
\[
\sin x^\circ \approx ax + bx^3 \quad (0 \leqslant x \leqslant 90; x = \operatorname{arc} x^\circ)
\]
利用这个公式，近似地求:
\[
\sin 20^\circ, \quad \sin 40^\circ, \quad \sin 80^\circ
\]
\end{prob}

\begin{prob}[M]{3125}
取点 $x_i = 0, \pm \dfrac{1}{2}, \pm 1$ 为拉格朗日多项式的插值节点，对函数 $f(x) = |x|$ 作出在闭区间 $[-1, 1]$ 上的拉格朗日插值多项式.
\end{prob}

\begin{prob}[M]{3126}
以拉格朗日多项式代换函数 $y(x)$，近似地计算
\[
\int_{0}^{2} y(x) \,\mathrm{d}x
\]
其中
\begin{center}
\begin{tabular}{|c|c|c|c|c|c|}
\hline
$x$ & $0$ & $0.5$ & $1$ & $1.5$ & $2$ \\
\hline
$y(x)$ & $5$ & $4.5$ & $3$ & $2.5$ & $5$ \\
\hline
\end{tabular}
\end{center}
\end{prob}

\begin{prob}[M]{3127}
对于函数 $x, x^2, x^3$，试在闭区间 $[0, 1]$ 上作出伯恩斯坦多项式 $B_n(x)$.
\end{prob}

\begin{prob}[M]{3128}
对于在闭区间 $[a, b]$ 上的已知函数 $f(x)$ 写出伯恩斯坦多项式 $B_n(x)$ 的公式.
\end{prob}

\begin{prob}[M]{3129}
在闭区间 $[-1, 1]$ 上用伯恩斯坦多项式 $B_4(x)$ 逼近函数 $f(x) = \dfrac{|x| + x}{2}$. 作出函数 $y = \dfrac{|x| + x}{2}$ 和 $y = B_4(x)$ 的图像.
\end{prob}

\begin{prob}[M]{3130}
在 $-1 \leqslant x \leqslant 1$ 时用偶数次伯恩斯坦多项式逼近函数，
\[
f(x) = |x|.
\]
\end{prob}

\begin{prob}[M]{3131}
对于函数
\[
f(x) = \mathrm{e}^{kx} \quad (a \leqslant x \leqslant b)
\]
写出伯恩斯坦多项式 $B_n(x)$.
\end{prob}

\begin{prob}[M]{3132}
在闭区间 $-\dfrac{\pi}{2} \leqslant x \leqslant \dfrac{\pi}{2}$ 上，对于函数 $f(x) = \cos x$ 计算多项式 $B_n(x)$.
\end{prob}

\begin{prob}[M]{3133.1}
证明: 在闭区间 $[-1, 1]$ 上 $|x| = \lim_{n \to \infty} P_n(x)$，其中
\[
P_n(x) = 1 - \frac{1 - x^2}{2} - \sum_{i=2}^{n} \frac{1 \cdot 3 \cdots (2i - 3)}{2 \cdot 4 \cdots (2i)} (1 - x^2)^i.
\]
\end{prob}

\begin{prob}[M]{3133.2}
设 $f(x) \in C[a, b]$，且
\[
M_k = \int_{a}^{b} x^k f(x) \,\mathrm{d}x = 0 \quad (k = 0, 1, 2, \cdots),
\]
证明: 当 $x \in [a, b]$ 时有 $f(x) \equiv 0$.
\probhint{利用关于用多项式逼近连续函数的魏尔斯特拉斯定理.}
\end{prob}

\begin{prob}[L]{3134}
设 $f(x)$ 是以 $2\pi$ 为周期的函数，并且在闭区间 $-\pi \leqslant x \leqslant \pi$ 上连续，而 $a_n, b_n$ ($n = 0, 1, 2, \cdots$) 是它的傅里叶系数. 证明: 费耶尔三角多项式
\[
\sigma_n(x) = \frac{a_0}{2} + \sum_{i=1}^{n-1} \left(1 - \frac{i}{n}\right) (a_i \cos ix + b_i \sin ix)
\]
在区间 $[-\pi, \pi]$ 上一致收敛于函数 $f(x)$.
\end{prob}

\begin{prob}[S]{3135}
对于函数
\[
-\pi \leqslant x \leqslant \pi
\]
作出费耶尔多项式 $\sigma_{2n-1}(x)$.
\end{prob}


\probpart{第二部分\quad 多元函数}
\probchapter{第六章\quad 多元函数微分学}

\probsec{§1. 函数的极限. 连续性}

\probgroup{确定并画出下列函数的存在域：}

\begin{prob}[S]{3136}
$u = x + \sqrt{y}$.
\end{prob}

\begin{prob}[S]{3137}
$u = \sqrt{1 - x^2} + \sqrt{y^2 - 1}$.
\end{prob}

\begin{prob}[S]{3138}
$u = \sqrt{1 - x^2 - y^2}$.
\end{prob}

\begin{prob}[S]{3139}
$u = \dfrac{1}{\sqrt{x^2 + y^2 - 1}}$.
\end{prob}

\begin{prob}[S]{3140}
$u = \sqrt{(x^2 + y^2 - 1)(4 - x^2 - y^2)}$.
\end{prob}

\begin{prob}[S]{3141}
$u = \sqrt{\dfrac{x^2 + y^2 - x}{2x - x^2 - y^2}}$.
\end{prob}

\begin{prob}[S]{3142}
$u = \sqrt{1 - (x^2 + y^2)}$.
\end{prob}

\begin{prob}[S]{3143}
$u = \ln(-x - y)$.
\end{prob}

\begin{prob}[S]{3144}
$u = \arcsin \dfrac{y}{x}$.
\end{prob}

\begin{prob}[S]{3145}
$u = \arccos \dfrac{x}{x + y}$.
\end{prob}

\begin{prob}[S]{3146}
$u = \arcsin \dfrac{x}{y^2} + \arcsin(1 - y)$.
\end{prob}

\begin{prob}[S]{3147}
$u = \sqrt{\sin(x^2 + y^2)}$.
\end{prob}

\begin{prob}[S]{3148}
$u = \arccos \dfrac{z}{\sqrt{x^2 + y^2}}$.
\end{prob}

\begin{prob}[S]{3149}
$u = \ln(xyz)$.
\end{prob}

\begin{prob}[S]{3150}
$u = \ln(-1 - x^2 - y^2 + z^2)$.
\end{prob}

\probgroup{作出下列函数的等值线：}

\begin{prob}[S]{3151}
$z = x + y$.
\end{prob}

\begin{prob}[S]{3152}
$z = x^2 + y^2$.
\end{prob}

\begin{prob}[S]{3153}
$z = x^2 - y^2$.
\end{prob}

\begin{prob}[S]{3154}
$z = (x + y)^2$.
\end{prob}

\begin{prob}[S]{3155}
$z = \dfrac{y}{x}$.
\end{prob}

\begin{prob}[S]{3156}
$z = \dfrac{1}{x^2 + 2y^2}$.
\end{prob}

\begin{prob}[S]{3157}
$z = \sqrt{xy}$.
\end{prob}

\begin{prob}[S]{3158}
$z = |x| + y$.
\end{prob}

\begin{prob}[L]{3159}
\begin{enumerate}[label=(\alph*), itemsep=2pt, topsep=3pt]
\item $z = |x| + |y| - |x + y|$;
\item $z = \min(x, y)$;
\item $z = \max(|x|, |y|)$;
\item $z = \min(x^2, y)$.
\end{enumerate}
\end{prob}

\begin{prob}[S]{3160}
$z = \mathrm{e}^{\frac{2x}{x^2 + y^2}}$.
\end{prob}

\begin{prob}[S]{3161}
$z = x^y \ (x > 0)$.
\end{prob}

\begin{prob}[S]{3162}
$z = x^y \mathrm{e}^{-x} \ (x > 0)$.
\end{prob}

\begin{prob}[S]{3163}
$z = \ln\sqrt{\dfrac{(x - a)^2 + y^2}{(x + a)^2 + y^2}} \ (a > 0)$.
\end{prob}

\begin{prob}[S]{3164}
$z = \arctan \dfrac{2ay}{x^2 + y^2 - a^2} \ (a > 0)$.
\end{prob}

\begin{prob}[S]{3165}
$z = \operatorname{sgn}(\sin x \sin y)$.
\end{prob}

\probgroup{求下列函数的等值面：}

\begin{prob}[S]{3166}
$u = x + y + z$.
\end{prob}

\begin{prob}[S]{3167}
$u = x^2 + y^2 + z^2$.
\end{prob}

\begin{prob}[S]{3168}
$u = x^2 + y^2 - z^2$.
\end{prob}

\begin{prob}[S]{3169}
$u = (x + y)^2 + z^2$.
\end{prob}

\begin{prob}[S]{3170}
$u = \operatorname{sgn}\sin(x^2 + y^2 + z^2)$.
\end{prob}

\probgroup{根据曲面的已知方程研究其性质：}

\begin{prob}[S]{3171}
$z = f(y - ax)$.
\end{prob}

\begin{prob}[S]{3172}
$z = f(\sqrt{x^2 + y^2})$.
\end{prob}

\begin{prob}[S]{3173}
$z = xf\left(\dfrac{y}{x}\right)$.
\end{prob}

\begin{prob}[S]{3174}
$z = f\left(\dfrac{y}{x}\right)$.
\end{prob}

\begin{prob}[M]{3175}
作出函数
\[
F(t) = f(\cos t, \sin t)
\]
的图像，式中
\[
f(x, y) = \begin{cases}
1, & y \geqslant x, \\
0, & y < x.
\end{cases}
\]
\end{prob}

\begin{prob}[S]{3176}
若
\[
f(x, y) = \frac{2xy}{x^2 + y^2},
\]
求 $f\left(1, \dfrac{y}{x}\right)$.
\end{prob}

\begin{prob}[S]{3177}
若
\[
f\left(\frac{y}{x}\right) = \frac{\sqrt{x^2 + y^2}}{x} \quad (x > 0),
\]
求 $f(x)$.
\end{prob}

\begin{prob}[S]{3178}
设
\[
z = \sqrt{y} + f(\sqrt{x} - 1).
\]
若当 $y = 1$ 时 $z = x$，求函数 $f$ 和 $z$.
\end{prob}

\begin{prob}[S]{3179}
设
\[
z = x + y + f(x - y).
\]
若当 $y = 0$ 时 $z = x^2$，求函数 $f$ 和 $z$.
\end{prob}

\begin{prob}[S]{3180}
若 $f\left(x + y, \dfrac{y}{x}\right) = x^2 - y^2$，求 $f(x, y)$.
\end{prob}

\begin{prob}[M]{3181}
证明：对于函数
\[
f(x, y) = \frac{x - y}{x + y},
\]
有
\[
\lim_{x \to 0} \left\{ \lim_{y \to 0} f(x, y) \right\} = 1, \qquad \lim_{y \to 0} \left\{ \lim_{x \to 0} f(x, y) \right\} = -1,
\]
从而 $\lim\limits_{\substack{x \to 0 \\ y \to 0}} f(x, y)$ 不存在.
\end{prob}

\begin{prob}[M]{3182}
证明：对于函数
\[
f(x, y) = \frac{x^2 y^2}{x^2 y^2 + (x - y)^2},
\]
有
\[
\lim_{x \to 0} \left\{ \lim_{y \to 0} f(x, y) \right\} = \lim_{y \to 0} \left\{ \lim_{x \to 0} f(x, y) \right\} = 0,
\]
然而 $\lim\limits_{\substack{x \to 0 \\ y \to 0}} f(x, y)$ 不存在.
\end{prob}

\begin{prob}[M]{3183.1}
证明：对于函数
\[
f(x, y) = (x + y)\sin\frac{1}{x}\sin\frac{1}{y},
\]
累次极限 $\lim\limits_{x \to 0} \left\{ \lim\limits_{y \to 0} f(x, y) \right\}$ 和 $\lim\limits_{y \to 0} \left\{ \lim\limits_{x \to 0} f(x, y) \right\}$ 不存在，但存在 $\lim\limits_{\substack{x \to 0 \\ y \to 0}} f(x, y) = 0$.
\end{prob}

\begin{prob}[S]{3183.2}
是否存在极限
\[
\lim_{\substack{x \to 0 \\ y \to 0}} \frac{2xy}{x^2 + y^2}?
\]
\end{prob}

\begin{prob}[M]{3183.3}
当 $t \to +\infty$ 时，函数
\[
f(x, y) = x^2 \mathrm{e}^{-(x^2 - y)}
\]
沿任意射线 $x = t\cos\alpha, y = t\sin\alpha$ ($0 \leqslant t < +\infty$) 的极限等于什么？当 $x \to +\infty, y \to +\infty$ 时，可否称此函数为无穷小量？
\end{prob}

\begin{prob}[L]{3184}
求 $\lim\limits_{x \to a} \left\{ \lim\limits_{y \to b} f(x, y) \right\}$ 及 $\lim\limits_{y \to b} \left\{ \lim\limits_{x \to a} f(x, y) \right\}$，设：
\begin{enumerate}[label=(\alph*), itemsep=2pt, topsep=3pt]
\item $f(x, y) = \dfrac{x^2 + y^2}{x^2 + y^4}, \ a = \infty, b = \infty$;
\item $f(x, y) = \dfrac{x^y}{1 + x^y}, \ a = \infty, b = +0$;
\item $f(x, y) = \sin\dfrac{\pi x}{2x + y}, \ a = \infty, b = \infty$;
\item $f(x, y) = \dfrac{1}{xy}\tan\dfrac{xy}{1 + xy}, \ a = 0, b = \infty$;
\item $f(x, y) = \log_x(x + y), \ a = 1, b = 0$.
\end{enumerate}
\end{prob}

\probgroup{求下列二重极限：}

\begin{prob}[S]{3185}
$\lim\limits_{\substack{x \to \infty \\ y \to \infty}} \dfrac{x + y}{x^2 - xy + y^2}$.
\end{prob}

\begin{prob}[S]{3186}
$\lim\limits_{\substack{x \to \infty \\ y \to \infty}} \dfrac{x^2 + y^2}{x^4 + y^4}$.
\end{prob}

\begin{prob}[S]{3187}
$\lim\limits_{\substack{x \to 0 \\ y \to a}} \dfrac{\sin xy}{x}$.
\end{prob}

\begin{prob}[S]{3188}
$\lim\limits_{\substack{x \to +\infty \\ y \to +\infty}} (x^2 + y^2)\mathrm{e}^{-(x + y)}$.
\end{prob}

\begin{prob}[S]{3189}
$\lim\limits_{\substack{x \to +\infty \\ y \to +\infty}} \left(\dfrac{xy}{x^2 + y^2}\right)^{x^2}$.
\end{prob}

\begin{prob}[S]{3190}
$\lim\limits_{\substack{x \to 0 \\ y \to 0}} (x^2 + y^2)^{x^2 y^2}$.
\end{prob}

\begin{prob}[S]{3191}
$\lim\limits_{\substack{x \to \infty \\ y \to a}} \left(1 + \dfrac{1}{x}\right)^{\frac{x^2}{x + y}}$.
\end{prob}

\begin{prob}[S]{3192}
$\lim\limits_{\substack{x \to 1 \\ y \to 0}} \dfrac{\ln(x + \mathrm{e}^y)}{\sqrt{x^2 + y^2}}$.
\end{prob}

\begin{prob}[M]{3193}
若 $x = \rho\cos\varphi, y = \rho\sin\varphi$，问沿怎样的方向 $\varphi$ 存在有限的极限值：

(a) $\lim\limits_{\rho \to +0} \mathrm{e}^{\frac{x}{x^2 + y^2}}$; \qquad (b) $\lim\limits_{\rho \to +\infty} \mathrm{e}^{x^2 - y^2} \cdot \sin 2xy$.
\end{prob}

\probgroup{求下列函数的间断点：}

\begin{prob}[S]{3194}
$u = \dfrac{1}{\sqrt{x^2 + y^2}}$.
\end{prob}

\begin{prob}[S]{3195}
$u = \dfrac{xy}{x + y}$.
\end{prob}

\begin{prob}[S]{3196}
$u = \dfrac{x + y}{x^3 + y^3}$.
\end{prob}

\begin{prob}[S]{3197}
$u = \sin\dfrac{1}{xy}$.
\end{prob}

\begin{prob}[S]{3198}
$u = \dfrac{1}{\sin x \sin y}$.
\end{prob}

\begin{prob}[S]{3199}
$u = \ln(1 - x^2 - y^2)$.
\end{prob}

\begin{prob}[S]{3200}
$u = \dfrac{1}{xyz}$.
\end{prob}

\begin{prob}[S]{3201}
$u = \ln\dfrac{1}{\sqrt{(x - a)^2 + (y - b)^2 + (z - c)^2}}$.
\end{prob}

\begin{prob}[L]{3202}
证明：
\begin{enumerate}[label=(\alph*), itemsep=2pt, topsep=3pt]
\item 函数
\[
f(x, y) = \begin{cases}
\dfrac{2xy}{x^2 + y^2}, & x^2 + y^2 \neq 0, \\
0, & x^2 + y^2 = 0
\end{cases}
\]
分别对每一个变量 $x$ 或 $y$（当另一变量的值固定时）是连续的，但对这些变量的总体不是连续的；
\item 函数
\[
f(x, y) = \begin{cases}
\dfrac{x^2 y}{x^4 + y^4}, & x^2 + y^2 \neq 0, \\
0, & x^2 + y^2 = 0
\end{cases}
\]
在点 $O(0, 0)$ 处沿过此点的每一射线
\[
x = t\cos\alpha, \quad y = t\sin\alpha \quad (0 \leqslant t < +\infty)
\]
连续，即存在
\[
\lim_{t \to 0} f(t\cos\alpha, t\sin\alpha) = f(0, 0),
\]
但此函数在点 $(0, 0)$ 并非连续的.
\end{enumerate}
\end{prob}

\begin{prob}[M]{3203.1}
研究线性函数
\[
u = 2x - 3y + 5
\]
在平面 $E^2 = \{|x| < +\infty, |y| < +\infty\}$ 上的一致连续性.
\end{prob}

\begin{prob}[M]{3203.2}
研究函数
\[
u = \sqrt{x^2 + y^2}
\]
在平面 $E^2 = \{|x| < +\infty, |y| < +\infty\}$ 上的一致连续性.
\end{prob}

\begin{prob}[M]{3203.3}
函数
\[
f(x, y) = \sin\frac{\pi}{1 - x^2 - y^2}
\]
在区域 $x^2 + y^2 < 1$ 内是否一致连续？
\end{prob}

\begin{prob}[M]{3203.4}
给定函数
\[
u = \arcsin\frac{x}{y}.
\]
此函数在其定义域 $E$ 内是否连续？是否一致连续？
\end{prob}

\begin{prob}[M]{3204}
证明：函数
\[
f(x, y) = \begin{cases}
x\sin\dfrac{1}{y}, & y \neq 0, \\
0, & y = 0
\end{cases}
\]
的间断点的集合不是闭集.
\end{prob}

\begin{prob}[M]{3205}
证明：若函数 $f(x, y)$ 在某区域 $G$ 内对变量 $x$ 是连续的，而关于 $x$ 对变量 $y$ 是一致连续的，则此函数在该区域内是连续的.
\end{prob}

\begin{prob}[M]{3206}
证明：若在某区域 $G$ 内函数 $f(x, y)$ 对变量 $x$ 是连续的，并满足对变量 $y$ 的利普希茨条件，即
\[
|f(x, y') - f(x, y'')| \leqslant L|y' - y''|,
\]
式中 $(x, y') \in G, (x, y'') \in G$，而 $L$ 为常数，则此函数在该区域内是连续的.
\end{prob}

\begin{prob}[M]{3207}
证明：若函数 $f(x, y)$ 分别对每一个变量 $x$ 和 $y$ 是连续的，并对其中的一个是单调的，则此函数对两个变量的总体是连续的（杨定理）.
\end{prob}

\begin{prob}[M]{3208}
设函数 $f(x, y)$ 在区域 $a \leqslant x \leqslant A, b \leqslant y \leqslant B$ 上是连续的，而函数序列 $\varphi_n(x)$ ($n = 1, 2, \cdots$) 在 $[a, A]$ 上一致收敛并满足条件 $b \leqslant \varphi_n(x) \leqslant B$. 证明：函数序列
\[
F_n(x) = f(x, \varphi_n(x)) \quad (n = 1, 2, \cdots)
\]
也在 $[a, A]$ 上一致收敛.
\end{prob}

\begin{prob}[M]{3209}
设：1) 函数 $f(x, y)$ 在区域 $R$ ($a < x < A; b < y < B$) 内连续；2) 函数 $\varphi(x)$ 在区间 $(a, A)$ 内连续，且函数值属于区间 $(b, B)$. 证明：函数
\[
F(x) = f(x, \varphi(x))
\]
在区间 $(a, A)$ 内连续.
\end{prob}

\begin{prob}[M]{3210}
设：1) 函数 $f(x,y)$ 在区域 $R$ ($a < x < A; b < y < B$) 内连续；2) 函数 $x = \varphi(u,v)$ 及 $y = \psi(u,v)$ 在区域 $R'$ ($a' < u < A'; b' < v < B'$) 内连续，且函数值分别属于区间 $(a, A)$ 和 $(b, B)$．证明：函数
\[
F(u,v) = f(\varphi(u,v), \psi(u,v))
\]
在区域 $R'$ 内连续．
\end{prob}

\probsec{§2. 偏导数. 函数的微分}

\begin{prob}[S]{3211.1}
证明：
\[
f'_x(x,b) = \frac{\mathrm{d}}{\mathrm{d}x}[f(x,b)].
\]
\end{prob}

\begin{prob}[M]{3211.2}
设
\[
f(x,y) = x + (y-1)\arcsin\sqrt{\frac{x}{y}},
\]
求 $f'_x(x,1)$．
\end{prob}

\begin{prob}[M]{3212.1}
设 $f(x,y) = \sqrt[3]{xy}$，求 $f'_x(0,0)$ 和 $f'_y(0,0)$．此函数在点 $O(0,0)$ 是否可微？
\end{prob}

\begin{prob}[S]{3212.2}
函数 $f(x,y) = \sqrt[3]{x^3 + y^3}$ 在点 $O(0,0)$ 是否可微？
\end{prob}

\begin{prob}[M]{3212.3}
研究函数
\[
f(x,y) = \begin{cases} \mathrm{e}^{-\frac{1}{x^2+y^2}}, & x^2+y^2 > 0, \\ 0, & x^2+y^2 = 0 \end{cases}
\]
在点 $O(0,0)$ 的可微性．
\end{prob}

\probgroup{求下列函数的一阶和二阶偏导数：}

\begin{prob}[S]{3213}
$u = x^4 + y^4 - 4x^2 y^2$．
\end{prob}

\begin{prob}[S]{3214}
$u = xy + \dfrac{x}{y}$．
\end{prob}

\begin{prob}[S]{3215}
$u = \dfrac{x}{y^2}$．
\end{prob}

\begin{prob}[S]{3216}
$u = \dfrac{x}{\sqrt{x^2+y^2}}$．
\end{prob}

\begin{prob}[S]{3217}
$u = x\sin(x+y)$．
\end{prob}

\begin{prob}[S]{3218}
$u = \dfrac{\cos x^2}{y}$．
\end{prob}

\begin{prob}[S]{3219}
$u = \tan\dfrac{x^2}{y}$．
\end{prob}

\begin{prob}[S]{3220}
$u = x^y$．
\end{prob}

\begin{prob}[S]{3221}
$u = \ln(x+y^2)$．
\end{prob}

\begin{prob}[S]{3222}
$u = \arctan\dfrac{y}{x}$．
\end{prob}

\begin{prob}[S]{3223}
$u = \arctan\dfrac{x+y}{1-xy}$．
\end{prob}

\begin{prob}[S]{3224}
$u = \arcsin\dfrac{x}{\sqrt{x^2+y^2}}$．
\end{prob}

\begin{prob}[S]{3225}
$u = \dfrac{1}{\sqrt{x^2+y^2+z^2}}$．
\end{prob}

\begin{prob}[S]{3226}
$u = \left(\dfrac{x}{y}\right)^z$．
\end{prob}

\begin{prob}[S]{3227}
$u = x^{\frac{y}{z}}$．
\end{prob}

\begin{prob}[S]{3228}
$u = x^{y^z}$．
\end{prob}

\begin{prob}[L]{3229}
设
\begin{enumerate}[label=(\alph*), itemsep=2pt, topsep=3pt]
\item $u = x^2 - 2xy - 3y^2$；
\item $u = x^{y^2}$；
\item $u = \arccos\sqrt{\dfrac{x}{y}}$，
\end{enumerate}
验证等式
\[
\frac{\partial^2 u}{\partial x\partial y} = \frac{\partial^2 u}{\partial y\partial x}.
\]
\end{prob}

\begin{prob}[M]{3230.1}
设
\[
f(x,y) = \begin{cases} xy\dfrac{x^2-y^2}{x^2+y^2}, & x^2+y^2 \neq 0, \\ 0, & x^2+y^2 = 0. \end{cases}
\]
证明：
\[
f''_{xy}(0,0) \neq f''_{yx}(0,0).
\]
\end{prob}

\begin{prob}[M]{3230.2}
设
\[
f(x,y) = \begin{cases} \dfrac{2xy}{x^2+y^2}, & x^2+y^2 > 0, \\ 0, & x^2+y^2 = 0, \end{cases}
\]
是否存在 $f''_{xy}(0,0)$？
\end{prob}

\begin{prob}[L]{3231}
设 $u = f(x,y,z)$ 为 $n$ 次齐次函数，就下列各题验证关于齐次函数的欧拉定理：
\begin{enumerate}[label=(\alph*), itemsep=2pt, topsep=3pt]
\item $u = (x-2y+3z)^2$；
\item $u = \dfrac{x}{\sqrt{x^2+y^2+z^2}}$；
\item $u = \left(\dfrac{x}{y}\right)^{\frac{y}{x}}$．
\end{enumerate}
\end{prob}

\begin{prob}[M]{3232}
证明：若可微函数 $u = f(x,y,z)$ 满足方程
\[
x\frac{\partial u}{\partial x} + y\frac{\partial u}{\partial y} + z\frac{\partial u}{\partial z} = nu,
\]
则它为 $n$ 次齐次函数．
\probhint{研究辅助函数
\[
F(t) = \frac{f(tx,ty,tz)}{t^n}.
\]}
\end{prob}

\begin{prob}[M]{3233}
证明：若 $f(x,y,z)$ 是可微的 $n$ 次齐次函数，则其偏导数 $f'_x(x,y,z)$，$f'_y(x,y,z)$，$f'_z(x,y,z)$ 是 $n-1$ 次齐次函数．
\end{prob}

\begin{prob}[M]{3234}
设 $u = f(x,y,z)$ 是二阶可微的 $n$ 次齐次函数，证明：
\[
\left(x\frac{\partial}{\partial x} + y\frac{\partial}{\partial y} + z\frac{\partial}{\partial z}\right)^2 u = n(n-1)u.
\]
\end{prob}

\probgroup{求下列函数的一阶和二阶微分（$x,y,z$ 为自变量）：}

\begin{prob}[S]{3235}
$u = x^m y^n$．
\end{prob}

\begin{prob}[S]{3236}
$u = \dfrac{x}{y}$．
\end{prob}

\begin{prob}[S]{3237}
$u = \sqrt{x^2+y^2}$．
\end{prob}

\begin{prob}[S]{3238}
$u = \ln\sqrt{x^2+y^2}$．
\end{prob}

\begin{prob}[S]{3239}
$u = \mathrm{e}^{xy}$．
\end{prob}

\begin{prob}[S]{3240}
$u = xy + yz + zx$．
\end{prob}

\begin{prob}[S]{3241}
$u = \dfrac{z}{x^2+y^2}$．
\end{prob}

\begin{prob}[M]{3242}
设 $f(x,y,z) = \sqrt[z]{\dfrac{x}{y}}$，求 $\mathrm{d}f(1,1,1)$ 及 $\mathrm{d}^2 f(1,1,1)$．
\end{prob}

\begin{prob}[M]{3243}
证明：若 $u = \sqrt{x^2+y^2+z^2}$，则 $\mathrm{d}^2 u \geqslant 0$．
\end{prob}

\begin{prob}[L]{3244}
假设 $x, y$ 的绝对值很小，对下列各式推出近似公式：
\begin{enumerate}[label=(\alph*), itemsep=2pt, topsep=3pt]
\item $(1+x)^m (1+y)^n$；
\item $\ln(1+x)\cdot \ln(1+y)$；
\item $\arctan\dfrac{x+y}{1+xy}$．
\end{enumerate}
\end{prob}

\begin{prob}[L]{3245}
用微分来代替函数的增量，近似地计算：
\begin{enumerate}[label=(\alph*), itemsep=2pt, topsep=3pt]
\item $1.002 \cdot 2.003^2 \cdot 3.004^3$；
\item $\dfrac{1.03^2}{\sqrt[3]{0.98} \cdot \sqrt[4]{1.05^3}}$；
\item $\sqrt{1.02^3 + 1.97^3}$；
\item $\sin 29^\circ \cdot \tan 46^\circ$；
\item $0.97^{1.05}$．
\end{enumerate}
\end{prob}

\begin{prob}[M]{3246}
设矩形的边 $x = 6\,\mathrm{m}$ 和 $y = 8\,\mathrm{m}$，若第一个边增加 $2\,\mathrm{mm}$，而第二个边减少 $5\,\mathrm{mm}$，则矩形的对角线和面积变化多少？
\end{prob}

\begin{prob}[M]{3247}
扇形的中心角 $\alpha = 60^\circ$ 增加 $\Delta\alpha = 1^\circ$，为了使扇形的面积仍然不变，则应当把扇形的半径 $R = 20\,\mathrm{cm}$ 减少多少？
\end{prob}

\begin{prob}[M]{3248}
证明：乘积的相对误差近似地等于乘数的相对误差之和．
\end{prob}

\begin{prob}[M]{3249}
当测量圆柱的底半径 $R$ 和高 $H$ 时得到以下结果：
\[
R = 2.5\,\mathrm{m} \pm 0.1\,\mathrm{m}; \quad H = 4.0\,\mathrm{m} \pm 0.2\,\mathrm{m},
\]
则所计算出的圆柱体积会有怎样的绝对误差 $\Delta$ 和相对误差 $\delta$？
\end{prob}

\begin{prob}[M]{3250}
三角形的边 $a = 200\,\mathrm{m} \pm 2\,\mathrm{m}$，$b = 300\,\mathrm{m} \pm 5\,\mathrm{m}$，它们之间的角 $C = 60^\circ \pm 1^\circ$．则所计算出的三角形第三边 $c$ 会有怎样的绝对误差？
\end{prob}

\begin{prob}[M]{3251}
证明：在点 $(0,0)$ 连续的函数
\[
f(x,y) = \sqrt{|xy|},
\]
在点 $(0,0)$ 有两个偏导数 $f'_x(0,0)$ 和 $f'_y(0,0)$，但在点 $(0,0)$ 并非可微的．

说明导数 $f'_x(x,y)$ 和 $f'_y(x,y)$ 在点 $(0,0)$ 的邻域中的性质．
\end{prob}

\begin{prob}[M]{3252}
证明：函数
\[
f(x,y) = \begin{cases} \dfrac{xy}{\sqrt{x^2+y^2}}, & x^2+y^2 \neq 0, \\ 0, & x^2+y^2 = 0 \end{cases}
\]
在点 $(0,0)$ 的邻域中连续且有界的偏导数 $f'_x(x,y)$ 和 $f'_y(x,y)$，但此函数在点 $(0,0)$ 不可微．
\end{prob}

\begin{prob}[M]{3253}
证明：函数
\[
f(x,y) = \begin{cases} (x^2+y^2)\sin\dfrac{1}{x^2+y^2}, & x^2+y^2 \neq 0, \\ 0, & x^2+y^2 \neq 0 \end{cases}
\]
在点 $(0,0)$ 的邻域中有偏导数 $f'_x(x,y)$ 和 $f'_y(x,y)$，这些偏导数在点 $(0,0)$ 是不连续的，且在此点的任何邻域中是无界的；然而，此函数在点 $(0,0)$ 可微．
\end{prob}

\begin{prob}[M]{3254}
证明：在某凸区域 $E$ 内具有有界偏导数 $f'_x(x,y)$ 和 $f'_y(x,y)$ 的函数 $f(x,y)$ 在此区域 $E$ 内一致连续．
\end{prob}

\begin{prob}[M]{3255}
证明：若函数 $f(x,y)$ 对变量 $x$ 连续（对每一个固定的值 $y$）且有对变量 $y$ 的有界的导数 $f'_y(x,y)$，则此函数对变量 $x$ 和 $y$ 的总体是连续的．
\end{prob}

\probgroup{在下列问题中求所列偏导数：}

\begin{prob}[M]{3256}
$\dfrac{\partial^4 u}{\partial x^4}$，$\dfrac{\partial^4 u}{\partial x^3 \partial y}$，$\dfrac{\partial^4 u}{\partial x^2 \partial y^2}$，若
\[
u = x - y + x^2 + 2xy + y^2 + x^3 - 3x^2 y - y^3 + x^4 - 4x^2 y^2 + y^4.
\]
\end{prob}

\begin{prob}[S]{3257}
$\dfrac{\partial^3 u}{\partial x^2 \partial y}$，若 $u = x\ln(xy)$．
\end{prob}

\begin{prob}[S]{3258}
$\dfrac{\partial^6 u}{\partial x^3 \partial y^3}$，若 $u = x^3 \sin y + y^3 \sin x$．
\end{prob}

\begin{prob}[M]{3259}
$\dfrac{\partial^3 u}{\partial x\partial y\partial z}$，若 $u = \arctan\dfrac{x+y+z-xyz}{1-xy-xz-yz}$．
\end{prob}

\begin{prob}[S]{3260}
$\dfrac{\partial^3 u}{\partial x\partial y\partial z}$，若 $u = \mathrm{e}^{xyz}$．
\end{prob}

\begin{prob}[M]{3261}
$\dfrac{\partial^4 u}{\partial x\partial y\partial\xi\partial\eta}$，若 $u = \ln\dfrac{1}{\sqrt{(x-\xi)^2 + (y-\eta)^2}}$．
\end{prob}

\begin{prob}[S]{3262}
$\dfrac{\partial^{p+q} u}{\partial x^p \partial y^q}$，若 $u = (x-x_0)^p (y-y_0)^q$．
\end{prob}

\begin{prob}[S]{3263}
$\dfrac{\partial^{m+n} u}{\partial x^m \partial y^n}$，若 $u = \dfrac{x+y}{x-y}$．
\end{prob}

\begin{prob}[S]{3264}
$\dfrac{\partial^{m+n} u}{\partial x^m \partial y^n}$，若 $u = (x^2+y^2)\mathrm{e}^{x+y}$．
\end{prob}

\begin{prob}[S]{3265}
$\dfrac{\partial^{p+q+r} u}{\partial x^p \partial y^q \partial z^r}$，若 $u = xyz\mathrm{e}^{x+y+z}$．
\end{prob}

\begin{prob}[S]{3266}
若 $f(x,y) = \mathrm{e}^x \sin y$，求 $f^{(m+n)}_{x^m y^n}(0,0)$．
\end{prob}

\begin{prob}[M]{3267}
证明：若 $u = f(xyz)$，则 $\dfrac{\partial^3 u}{\partial x\partial y\partial z} = F(t)$，式中 $t = xyz$，并求函数 $F$．
\end{prob}

\begin{prob}[M]{3268}
设
\[
u = x^4 - 2x^3 y - 2xy^3 + y^4 + x^3 - 3x^2 y - 3xy^2 + y^3 + 2x^2 - xy + 2y^2 + x + y + 1,
\]
求 $\mathrm{d}^4 u$．导数 $\dfrac{\partial^4 u}{\partial x^4}$，$\dfrac{\partial^4 u}{\partial x^3 \partial y}$，$\dfrac{\partial^4 u}{\partial x^2 \partial y^2}$，$\dfrac{\partial^4 u}{\partial x\partial y^3}$ 和 $\dfrac{\partial^4 u}{\partial y^4}$ 等于什么？
\end{prob}

\probgroup{在下列各题中求所指出的阶的全微分：}

\begin{prob}[S]{3269}
$\mathrm{d}^3 u$，若 $u = x^3 + y^3 - 3xy(x-y)$．
\end{prob}

\begin{prob}[S]{3270}
$\mathrm{d}^3 u$，若 $u = \sin(x^2+y^2)$．
\end{prob}

\begin{prob}[S]{3271}
$\mathrm{d}^{10} u$，若 $u = \ln(x+y)$．
\end{prob}

\begin{prob}[S]{3272}
$\mathrm{d}^6 u$，若 $u = \cos x \cosh y$．
\end{prob}

\begin{prob}[S]{3273}
$\mathrm{d}^3 u$，若 $u = xyz$．
\end{prob}

\begin{prob}[S]{3274}
$\mathrm{d}^4 u$，若 $u = \ln(x^x y^y z^z)$．
\end{prob}

\begin{prob}[S]{3275}
$\mathrm{d}^n u$，若 $u = \mathrm{e}^{ax+by}$．
\end{prob}

\begin{prob}[S]{3276}
$\mathrm{d}^n u$，若 $u = X(x)Y(y)$．
\end{prob}

\begin{prob}[S]{3277}
$\mathrm{d}^n u$，若 $u = f(x+y+z)$．
\end{prob}

\begin{prob}[S]{3278}
$\mathrm{d}^n u$，若 $u = \mathrm{e}^{ax+by+cz}$．
\end{prob}

\begin{prob}[M]{3279}
设 $P_n(x,y,z)$ 为 $n$ 次齐次多项式．证明：
\[
\mathrm{d}^n P_n(x,y,z) = n! P_n(\mathrm{d}x, \mathrm{d}y, \mathrm{d}z).
\]
\end{prob}

\begin{prob}[M]{3280}
设
\[
Au = x\frac{\partial u}{\partial x} + y\frac{\partial u}{\partial y}.
\]
求 $Au$ 和 $A^2 u = A(Au)$，若 (a) $u = \dfrac{x}{x^2+y^2}$；\qquad (b) $u = \ln\sqrt{x^2+y^2}$．
\end{prob}

\begin{prob}[M]{3281}
设
\[
\Delta u = \frac{\partial^2 u}{\mathrm{d}x^2} + \frac{\partial^2 u}{\partial y^2},
\]
求 $\Delta u$, 若 (a) $u = \sin x \cosh y$; (b) $u = \ln \sqrt{x^2 + y^2}$.
\end{prob}

\begin{prob}[M]{3282}
设
\[
\Delta_1 u = \left(\frac{\partial u}{\partial x}\right)^2 + \left(\frac{\partial u}{\partial y}\right)^2 + \left(\frac{\partial u}{\partial z}\right)^2, \quad \Delta_2 u = \frac{\partial^2 u}{\partial x^2} + \frac{\partial^2 u}{\partial y^2} + \frac{\partial^2 u}{\partial z^2}.
\]
求 $\Delta_1 u$ 和 $\Delta_2 u$, 若 (a) $u = x^3 + y^3 + z^3 - 3xyz$; (b) $u = \dfrac{1}{\sqrt{x^2 + y^2 + z^2}}$.
\end{prob}

\probgroup{求下列复合函数的一阶和二阶导数：}

\begin{prob}[S]{3283}
$u = f(x^2 + y^2 + z^2)$.
\end{prob}

\begin{prob}[S]{3284}
$u = f\left(x, \frac{x}{y}\right)$.
\end{prob}

\begin{prob}[S]{3285}
$u = f(x, xy, xyz)$.
\end{prob}

\begin{prob}[S]{3286}
设 $u = f(x + y, xy)$, 求 $\frac{\partial^2 u}{\partial x \partial y}$.
\end{prob}

\begin{prob}[S]{3287}
设 $u = f(x + y + z, x^2 + y^2 + z^2)$, 求 $\Delta u = \frac{\partial^2 u}{\partial x^2} + \frac{\partial^2 u}{\partial y^2} + \frac{\partial^2 u}{\partial z^2}$.
\end{prob}

\probgroup{求下列复合函数的一阶和二阶全微分 ($x, y$ 及 $z$ 为自变量)：}

\begin{prob}[S]{3288}
$u = f(t)$, 其中 $t = x + y$.
\end{prob}

\begin{prob}[S]{3289}
$u = f(t)$, 其中 $t = \frac{y}{x}$.
\end{prob}

\begin{prob}[S]{3290}
$u = f(\sqrt{x^2 + y^2})$.
\end{prob}

\begin{prob}[S]{3291}
$u = f(t)$, 其中 $t = xyz$.
\end{prob}

\begin{prob}[S]{3292}
$u = f(x^2 + y^2 + z^2)$.
\end{prob}

\begin{prob}[S]{3293}
$u = f(\xi, \eta)$, 其中 $\xi = ax, \eta = by$.
\end{prob}

\begin{prob}[S]{3294}
$u = f(\xi, \eta)$, 其中 $\xi = x + y, \eta = x - y$.
\end{prob}

\begin{prob}[S]{3295}
$u = f(\xi, \eta)$, 其中 $\xi = xy, \eta = \frac{x}{y}$.
\end{prob}

\begin{prob}[S]{3296}
$u = f(x + y, z)$.
\end{prob}

\begin{prob}[S]{3297}
$u = f(x + y + z, x^2 + y^2 + z^2)$.
\end{prob}

\begin{prob}[S]{3298}
$u = f\left(\frac{x}{y}, \frac{y}{z}\right)$.
\end{prob}

\begin{prob}[S]{3299}
$u = f(x, y, z)$, 其中 $x = t, y = t^2, z = t^3$.
\end{prob}

\begin{prob}[S]{3300}
$u = f(\xi, \eta, \zeta)$, 其中 $\xi = ax, \eta = by, \zeta = cz$.
\end{prob}

\begin{prob}[S]{3301}
$u = f(\xi, \eta, \zeta)$, 其中 $\xi = x^2 + y^2, \eta = x^2 - y^2, \zeta = 2xy$.
\end{prob}

\probgroup{求 $\mathrm{d}^n u$, 设：}

\begin{prob}[S]{3302}
$u = f(ax + by + cz)$.
\end{prob}

\begin{prob}[S]{3303}
$u = f(ax, by, cz)$.
\end{prob}

\begin{prob}[M]{3304}
$u = f(\xi, \eta, \zeta)$, 其中
\[
\xi = a_1 x + b_1 y + c_1 z, \quad \eta = a_2 x + b_2 y + c_2 z, \quad \zeta = a_3 x + b_3 y + c_3 z.
\]
\end{prob}

\begin{prob}[M]{3305}
设 $u = f(r)$, 其中 $r = \sqrt{x^2 + y^2 + z^2}$ 和 $f$ 为二阶可微的函数. 证明:
\[
\Delta u = F(r),
\]
其中 $\Delta u = \frac{\partial^2 u}{\partial x^2} + \frac{\partial^2 u}{\partial y^2} + \frac{\partial^2 u}{\partial z^2}$, $\Delta$ 为拉普拉斯算子, 并求函数 $F$.
\end{prob}

\begin{prob}[M]{3306}
设 $u$ 和 $v$ 为二阶可微的函数, $\Delta$ 为拉普拉斯算子 (参阅习题 3305). 证明:
\[
\Delta(uv) = u\Delta v + v\Delta u + 2\Delta(u, v),
\]
其中
\[
\Delta(u, v) = \frac{\partial u}{\partial x}\frac{\partial v}{\partial x} + \frac{\partial u}{\partial y}\frac{\partial v}{\partial y} + \frac{\partial u}{\partial z}\frac{\partial v}{\partial z}.
\]
\end{prob}

\begin{prob}[M]{3307}
证明: 函数
\[
u = \ln\sqrt{(x - a)^2 + (y - b)^2}
\]
($a$ 和 $b$ 为常数) 满足拉普拉斯方程
\[
\frac{\partial^2 u}{\partial x^2} + \frac{\partial^2 u}{\partial y^2} = 0.
\]
\end{prob}

\begin{prob}[M]{3308}
证明: 若函数 $u = u(x, y)$ 满足拉普拉斯方程 (参阅习题 3307), 则函数
\[
v = u\left(\frac{x}{x^2 + y^2}, \frac{y}{x^2 + y^2}\right)
\]
也满足此方程.
\end{prob}

\begin{prob}[M]{3309}
证明: 函数
\[
u = \frac{1}{2a\sqrt{\pi t}}\mathrm{e}^{-\frac{(x - b)^2}{4a^2 t}}
\]
($a$ 和 $b$ 为常数) 满足热传导方程
\[
\frac{\partial u}{\partial t} = a^2\frac{\partial^2 u}{\partial x^2}.
\]
\end{prob}

\begin{prob}[M]{3310}
证明: 若函数 $u = u(x, t)$ 满足热传导方程 (参阅 3309 题), 则函数
\[
v = \frac{1}{a\sqrt{t}}\mathrm{e}^{-\frac{x^2}{4a^2 t}} u\left(\frac{x}{a^2 t}, -\frac{1}{a^4 t}\right) \quad (t > 0)
\]
也满足该方程.
\end{prob}

\begin{prob}[M]{3311}
证明: 函数
\[
u = \frac{1}{r} \quad \left(\text{式中 } r = \sqrt{(x - a)^2 + (y - b)^2 + (z - c)^2}\right)
\]
当 $r \neq 0$ 时满足拉普拉斯方程
\[
\Delta u \equiv \frac{\partial^2 u}{\partial x^2} + \frac{\partial^2 u}{\partial y^2} + \frac{\partial^2 u}{\partial z^2} = 0.
\]
\end{prob}

\begin{prob}[M]{3312}
证明: 若函数 $u = u(x, y, z)$ 满足拉普拉斯方程 (参阅 3311 题), 则函数
\[
v = \frac{1}{r} u\left(\frac{k^2 x}{r^2}, \frac{k^2 y}{r^2}, \frac{k^2 z}{r^2}\right)
\]
(式中 $k$ 为常数, $r = \sqrt{x^2 + y^2 + z^2}$) 也满足该方程.
\end{prob}

\begin{prob}[M]{3313}
证明: 函数
\[
u = \frac{C_1\mathrm{e}^{-ar} + C_2\mathrm{e}^{ar}}{r}
\]
(式中 $r = \sqrt{x^2 + y^2 + z^2}$, $C_1, C_2$ 为常数) 满足亥姆霍兹方程
\[
\frac{\partial^2 u}{\partial x^2} + \frac{\partial^2 u}{\partial y^2} + \frac{\partial^2 u}{\partial z^2} = a^2 u.
\]
\end{prob}

\begin{prob}[M]{3314}
设函数 $u_1 = u_1(x, y, z)$ 及 $u_2 = u_2(x, y, z)$ 满足拉普拉斯方程 $\Delta u = 0$. 证明: 函数
\[
v = u_1(x, y, z) + (x^2 + y^2 + z^2) u_2(x, y, z)
\]
满足双调和方程
\[
\Delta(\Delta v) = 0.
\]
\end{prob}

\begin{prob}[M]{3315}
设 $f(x, y, z)$ 是 $m$ 阶可微的 $n$ 次齐次函数. 证明:
\[
\left(x\frac{\partial}{\partial x} + y\frac{\partial}{\partial y} + z\frac{\partial}{\partial z}\right)^m f(x, y, z) = n(n - 1)\cdots(n - m + 1) f(x, y, z).
\]
\end{prob}

\begin{prob}[M]{3316}
若
\[
z = \sin y + f(\sin x - \sin y),
\]
其中 $f$ 为可微函数, 试简化表达式
\[
\sec x\frac{\partial z}{\partial x} + \sec y\frac{\partial z}{\partial y}.
\]
\end{prob}

\begin{prob}[M]{3317}
证明: 函数
\[
z = x^n f\left(\frac{y}{x^2}\right)
\]
(其中 $f$ 为任意的可微函数) 满足方程
\[
x\frac{\partial z}{\partial x} + 2y\frac{\partial z}{\partial y} = nz.
\]
\end{prob}

\begin{prob}[M]{3318}
证明:
\[
z = y f(x^2 - y^2)
\]
(其中 $f$ 为任意的可微函数) 满足方程
\[
y^2\frac{\partial z}{\partial x} + xy\frac{\partial z}{\partial y} = xz.
\]
\end{prob}

\begin{prob}[M]{3319}
若
\[
u = \frac{1}{12}x^4 - \frac{1}{6}x^3(y + z) + \frac{1}{2}x^2 yz + f(y - x, z - x),
\]
其中 $f$ 为可微函数, 试简化表达式
\[
\frac{\partial u}{\partial x} + \frac{\partial u}{\partial y} + \frac{\partial u}{\partial z}.
\]
\end{prob}

\begin{prob}[M]{3320}
设
\[
x^2 = vw, \quad y^2 = uw, \quad z^2 = uv,
\]
\[
f(x, y, z) = F(u, v, w),
\]
证明:
\[
x f'_x + y f'_y + z f'_z = u F'_u + v F'_v + w F'_w.
\]
\end{prob}

\probgroup{假定任意函数 $\varphi, \psi$ 等为足够多次可微的函数, 验证下列等式：}

\begin{prob}[S]{3321}
$y\frac{\partial z}{\partial x} - x\frac{\partial z}{\partial y} = 0$, 若 $z = \varphi(x^2 + y^2)$.
\end{prob}

\begin{prob}[S]{3322}
$x^2\frac{\partial z}{\partial x} - xy\frac{\partial z}{\partial y} + y^2 = 0$, 若 $z = \frac{y^2}{3x} + \varphi(xy)$.
\end{prob}

\begin{prob}[S]{3323}
$(x^2 - y^2)\frac{\partial z}{\partial x} + xy\frac{\partial z}{\partial y} = xyz$, 若 $z = \mathrm{e}^y \varphi\left(y\mathrm{e}^{\frac{x^2}{2y^2}}\right)$.
\end{prob}

\begin{prob}[S]{3324}
$x\frac{\partial u}{\partial x} + \alpha y\frac{\partial u}{\partial y} + \beta z\frac{\partial u}{\partial z} = nu$, 若 $u = x^n \varphi\left(\frac{y}{x^\alpha}, \frac{z}{x^\beta}\right)$.
\end{prob}

\begin{prob}[S]{3325}
$x\frac{\partial u}{\partial x} + y\frac{\partial u}{\partial y} + z\frac{\partial u}{\partial z} = u + \frac{xy}{z}$, 若 $u = \frac{xy}{z}\ln x + x\varphi\left(\frac{y}{x}, \frac{z}{x}\right)$.
\end{prob}

\begin{prob}[S]{3326}
$\frac{\partial^2 u}{\partial t^2} = a^2\frac{\partial^2 u}{\partial x^2}$, 若 $u = \varphi(x - at) + \psi(x + at)$.
\end{prob}

\begin{prob}[S]{3327}
$\frac{\partial^2 u}{\partial x^2} - 2\frac{\partial^2 u}{\partial x\partial y} + \frac{\partial^2 u}{\partial y^2} = 0$, 若 $u = x\varphi(x + y) + y\psi(x + y)$.
\end{prob}

\begin{prob}[S]{3328}
$x^2\frac{\partial^2 u}{\partial x^2} + 2xy\frac{\partial^2 u}{\partial x\partial y} + y^2\frac{\partial^2 u}{\partial y^2} = 0$, 若 $u = \varphi\left(\frac{y}{x}\right) + x\psi\left(\frac{y}{x}\right)$.
\end{prob}

\begin{prob}[S]{3329}
$x^2\frac{\partial^2 u}{\partial x^2} + 2xy\frac{\partial^2 u}{\partial x\partial y} + y^2\frac{\partial^2 u}{\partial y^2} = n(n - 1)u$, 若 $u = x^n\varphi\left(\frac{y}{x}\right) + x^{1 - n}\psi\left(\frac{y}{x}\right)$.
\end{prob}

\begin{prob}[S]{3330}
$\frac{\partial u}{\partial x}\frac{\partial^2 u}{\partial x\partial y} = \frac{\partial u}{\partial y}\frac{\partial^2 u}{\partial x^2}$, 若 $u = \varphi[x + \psi(y)]$.
\end{prob}

\probgroup{用逐次微分的方法消去任意函数 $\varphi$ 和 $\psi$：}

\begin{prob}[S]{3331}
$z = x + \varphi(xy)$.
\end{prob}

\begin{prob}[S]{3332}
$z = x\varphi\left(\frac{x}{y^2}\right)$.
\end{prob}

\begin{prob}[S]{3333}
$z = \varphi(\sqrt{x^2 + y^2})$.
\end{prob}

\begin{prob}[S]{3334}
$u = \varphi(x - y, y - z)$.
\end{prob}

\begin{prob}[S]{3335}
$u = \varphi\left(\frac{x}{y}, \frac{y}{x}\right)$.
\end{prob}

\begin{prob}[S]{3336}
$z = \varphi(x) + \psi(y)$.
\end{prob}

\begin{prob}[S]{3337}
$z = \varphi(x)\psi(y)$.
\end{prob}

\begin{prob}[S]{3338}
$z = \varphi(x + y) + \psi(x - y)$.
\end{prob}

\begin{prob}[S]{3339}
$z = x\varphi\left(\frac{x}{y}\right) + y\psi\left(\frac{x}{y}\right)$.
\end{prob}

\begin{prob}[S]{3340}
$z = \varphi(xy) + \psi\left(\frac{x}{y}\right)$.
\end{prob}

\begin{prob}[M]{3341}
求函数
\[
z = x^2 - y^2
\]
在点 $M(1, 1)$ 沿与 $Ox$ 轴的正向组成角 $\alpha = 60^\circ$ 的方向 $l$ 的导数.
\end{prob}

\begin{prob}[L]{3342}
求函数
\[
z = x^2 - xy + y^2
\]
在点 $M(1, 1)$ 沿与 $Ox$ 轴的正向组成 $\alpha$ 角的方向 $l$ 的导数. 在怎样的方向上此导数: (a) 有最大值; (b) 有最小值; (c) 等于 $0$.
\end{prob}

\begin{prob}[M]{3343}
求函数
\[
z = \ln(x^2 + y^2)
\]
在点 $M(x_0, y_0)$ 沿与过此点的等值线垂直方向的导数.
\end{prob}

\begin{prob}[M]{3344}
求函数
\[
z = 1 - \left(\frac{x^2}{a^2} + \frac{y^2}{b^2}\right)
\]
在点 $M\left(\frac{a}{\sqrt{2}}, \frac{b}{\sqrt{2}}\right)$ 沿曲线 $\frac{x^2}{a^2} + \frac{y^2}{b^2} = 1$ 在此点的内法线方向的导数.
\end{prob}

\begin{prob}[M]{3345}
求函数
\[
u = xyz
\]
在点 $M(1, 1, 1)$ 沿方向 $l\{\cos\alpha, \cos\beta, \cos\gamma\}$ 的导数. 函数在该点的梯度的大小是什么?
\end{prob}

\begin{prob}[M]{3346}
求函数
\[
u = \frac{1}{r}
\]
(式中 $r = \sqrt{x^2 + y^2 + z^2}$) 在点 $M_0(x_0, y_0, z_0)$ 的梯度的大小和方向.
\end{prob}

\begin{prob}[M]{3347}
求函数
\[
u = x^2 + y^2 - z^2
\]
在点 $A(\varepsilon, 0, 0)$ 及 $B(0, \varepsilon, 0)$ 二点的梯度之间的角度.
\end{prob}

\begin{prob}[M]{3348}
函数
\[
u = x + y + z
\]
和
\[
v = x + y + z + 0.001\sin\left(10^6\pi\sqrt{x^2 + y^2 + z^2}\right)
\]
的梯度的大小相差多少?
\end{prob}

\begin{prob}[M]{3349}
证明: 函数
\[
u = ax^2 + by^2 + cz^2
\]
和
\[
v = ax^2 + by^2 + cz^2 + 2mx + 2ny + 2pz
\]
($a, b, c, m, n, p$ 为常数且 $a^2 + b^2 + c^2 \neq 0$) 在点 $M_0(x_0, y_0, z_0)$ 的梯度之间的角度当点 $M_0$ 无限远移时趋于零.
\end{prob}

\begin{prob}[M]{3350}
设 $u = f(x, y, z)$ 为二阶可微的函数. 若 $\cos\alpha, \cos\beta, \cos\gamma$ 为方向 $l$ 的方向余弦, 求 $\frac{\partial^2 u}{\partial l^2} = \frac{\partial}{\partial l}\left(\frac{\partial u}{\partial l}\right)$.
\end{prob}

\begin{prob}[M]{3351}
设 $u = f(x, y, z)$ 为二阶可微的函数,
\[
l_1\{\cos\alpha_1, \cos\beta_1, \cos\gamma_1\}, \quad l_2\{\cos\alpha_2, \cos\beta_2, \cos\gamma_2\}, \quad l_3\{\cos\alpha_3, \cos\beta_3, \cos\gamma_3\}
\]
为三个互相垂直的方向. 证明:
\begin{enumerate}[label=(\alph*), itemsep=2pt, topsep=3pt]
\item $\left(\frac{\partial u}{\partial l_1}\right)^2 + \left(\frac{\partial u}{\partial l_2}\right)^2 + \left(\frac{\partial u}{\partial l_3}\right)^2 = \left(\frac{\partial u}{\partial x}\right)^2 + \left(\frac{\partial u}{\partial y}\right)^2 + \left(\frac{\partial u}{\partial z}\right)^2$;
\item $\frac{\partial^2 u}{\partial l_1^2} + \frac{\partial^2 u}{\partial l_2^2} + \frac{\partial^2 u}{\partial l_3^2} = \frac{\partial^2 u}{\partial x^2} + \frac{\partial^2 u}{\partial y^2} + \frac{\partial^2 u}{\partial z^2}$.
\end{enumerate}
\end{prob}

\begin{prob}[M]{3352}
设 $u = u(x, y)$ 为可微的函数, 且当 $y = x^2$ 时有
\[
u(x, y) = 1, \quad \frac{\partial u}{\partial x} = x,
\]
求当 $y = x^2$ 时的 $\frac{\partial u}{\partial y}$.
\end{prob}

\begin{prob}[M]{3353}
设函数 $u = u(x, y)$ 满足方程
\[
\frac{\partial^2 u}{\partial x^2} - \frac{\partial^2 u}{\partial y^2} = 0,
\]
以及下列条件:
\[
u(x, 2x) = x, \quad u'_x(x, 2x) = x^2.
\]
求 $u''_{xx}(x, 2x), u''_{xy}(x, 2x), u''_{yy}(x, 2x)$.
\end{prob}

\probgroup{假设 $z = z(x, y)$, 解下列方程：}

\begin{prob}[S]{3354}
$\frac{\partial^2 z}{\partial x^2} = 0$.
\end{prob}

\begin{prob}[S]{3355}
$\frac{\partial^2 z}{\partial x\partial y} = 0$.
\end{prob}

\begin{prob}[S]{3356}
$\frac{\partial^n z}{\partial y^n} = 0$.
\end{prob}

\begin{prob}[S]{3357}
假定 $u = u(x, y, z)$ 解方程
\[
\frac{\partial^3 u}{\partial x \partial y \partial z} = 0.
\]
\end{prob}

\begin{prob}[M]{3358}
求方程
\[
\frac{\partial z}{\partial y} = x^2 + 2y
\]
的解 $z = z(x, y)$，使它满足条件 $z(x, x^2) = 1$。
\end{prob}

\begin{prob}[M]{3359}
求方程
\[
\frac{\partial^2 z}{\partial y^2} = 2
\]
的解 $z = z(x, y)$，使它满足条件 $z(x, 0) = 1$，$z'_y(x, 0) = x$。
\end{prob}

\begin{prob}[M]{3360}
求方程
\[
\frac{\partial^2 z}{\partial x \partial y} = x + y
\]
的解 $z = z(x, y)$，使它满足条件 $z(x, 0) = x$，$z(0, y) = y^2$。
\end{prob}

\probsec{§3. 隐函数的微分法}

\begin{prob}[M]{3361}
证明：在每一点都不连续的狄利克雷函数
\[
y(x) = \begin{cases}
1, & x\text{ 为有理数}, \\
0, & x\text{ 为无理数}
\end{cases}
\]
满足方程 $y^2 - y = 0$。
\end{prob}

\begin{prob}[M]{3362}
设函数 $f(x)$ 定义于区间 $(a, b)$ 内。在怎样的情况下，方程
\[
f(x)y = 0
\]
在 $a < x < b$ 时有唯一连续的解 $y = 0$？
\end{prob}

\begin{prob}[M]{3363}
设函数 $f(x)$ 和 $g(x)$ 在区间 $(a, b)$ 内有定义且连续。在怎样的情况下，方程
\[
f(x)y = g(x)
\]
在区间 $(a, b)$ 内有唯一连续的解？
\end{prob}

\begin{prob}[L]{3364}
已知方程
\[
x^2 + y^2 = 1, \tag{1}
\]
设
\[
y = y(x) \quad (-1 \leqslant x \leqslant 1) \tag{2}
\]
为满足方程 (1) 的单值函数。
\begin{enumerate}[label=\arabic*), itemsep=2pt, topsep=3pt]
  \item 有多少单值函数 (2) 满足方程 (1)？
  \item 有多少单值连续函数 (2) 满足方程 (1)？
  \item 设：(a) $y(0) = 1$；(b) $y(1) = 0$，则有多少单值连续函数 (2) 满足方程 (1)？
\end{enumerate}
\end{prob}

\begin{prob}[L]{3365}
已知方程
\[
x^2 = y^2, \tag{1}
\]
设
\[
y = y(x) \quad (-\infty < x < +\infty) \tag{2}
\]
是满足方程 (1) 的单值函数。
\begin{enumerate}[label=\arabic*), itemsep=2pt, topsep=3pt]
  \item 有多少单值函数 (2) 满足方程 (1)？
  \item 有多少单值连续函数 (2) 满足方程 (1)？
  \item 有多少单值可微函数 (2) 满足方程 (1)？
  \item 设：(a) $y(1) = 1$；(b) $y(0) = 0$，则有多少单值连续函数 (2) 满足方程 (1)？
  \item 设 $y(1) = 1$，且 $\delta$ 为充分小的数，则有多少单值连续函数 $y = y(x)\ (1 - \delta < x < 1 + \delta)$ 满足方程 (1)？
\end{enumerate}
\end{prob}

\begin{prob}[L]{3366}
方程
\[
x^2 + y^2 = x^4 + y^4
\]
定义出多值函数 $y(x)$。函数在怎样的区域内：1) 单值，2) 有二个值，3) 有三个值，4) 有四个值？求此函数的分支点及单值连续的各分支。
\end{prob}

\begin{prob}[M]{3367}
求由方程
\[
(x^2 + y^2)^2 = x^2 - y^2
\]
定义的多值函数 $y$ 的分支点和单值连续的各分支 $y = y(x)\ (-1 \leqslant x \leqslant 1)$。
\end{prob}

\begin{prob}[L]{3368}
设函数 $f(x)$ 当 $a < x < b$ 时连续，函数 $\varphi(y)$ 当 $c < y < d$ 时单调增加而且连续。在怎样的条件下，方程
\[
\varphi(y) = f(x)
\]
定义出单值函数
\[
y = \varphi^{-1}(f(x))?
\]
研究例子：(a) $\sin y + \sinh y = x$；\qquad (b) $\mathrm{e}^{-y} = -\sin^2 x$。
\end{prob}

\begin{prob}[M]{3369}
设
\[
x = y + \varphi(y), \tag{1}
\]
其中 $\varphi(0) = 0$，且当 $-a < y < a$ 时 $|\varphi'(y)| \leqslant k < 1$。证明：当 $-\varepsilon < x < \varepsilon$ 时存在唯一的可微函数 $y = y(x)$ 满足方程 (1) 且 $y(0) = 0$。
\end{prob}

\begin{prob}[M]{3370}
设 $y = y(x)$ 为由方程
\[
x = ky + \varphi(y)
\]
定义的隐函数，其中常数 $k \neq 0$，$\varphi(y)$ 为以 $\omega$ 为周期的周期函数，且 $|\varphi'(y)| < |k|$。证明：
\[
y = \frac{x}{k} + \psi(x),
\]
其中 $\psi(x)$ 为以 $|k|\omega$ 为周期的周期函数。
\end{prob}

\probgroup{对于由下列各方程定义的函数 $y$，求 $y'$ 和 $y''$：}

\begin{prob}[S]{3371}
$x^2 + 2xy - y^2 = a^2$。
\end{prob}

\begin{prob}[S]{3372}
$\ln\sqrt{x^2 + y^2} = \arctan\dfrac{y}{x}$。
\end{prob}

\begin{prob}[S]{3373}
$y - \varepsilon\sin y = x\ (0 < \varepsilon < 1)$。
\end{prob}

\begin{prob}[S]{3374}
$x^y = y^x\ (x \neq y)$。
\end{prob}

\begin{prob}[S]{3375}
$y = 2x\arctan\dfrac{y}{x}$。
\end{prob}

\begin{prob}[M]{3376}
证明：若
\[
1 + xy = k(x - y),
\]
式中 $k$ 为常数，则成立等式
\[
\frac{\mathrm{d}x}{1 + x^2} = \frac{\mathrm{d}y}{1 + y^2}.
\]
\end{prob}

\begin{prob}[M]{3377}
证明：若
\[
x^2 y^2 + x^2 + y^2 - 1 = 0,
\]
则当 $xy > 0$ 时成立等式
\[
\frac{\mathrm{d}x}{\sqrt{1 - x^4}} + \frac{\mathrm{d}y}{\sqrt{1 - y^4}} = 0.
\]
\end{prob}

\begin{prob}[M]{3378}
证明：方程
\[
(x^2 + y^2)^2 = a^2(x^2 - y^2) \quad (a \neq 0)
\]
在点 $x = 0, y = 0$ 的邻域中定义出两个可微函数：$y = y_1(x)$ 和 $y = y_2(x)$。求 $y_1'(0)$ 及 $y_2'(0)$。
\end{prob}

\begin{prob}[M]{3379}
设
\[
(x^2 + y^2)^2 = 3x^2 y - y^3,
\]
求 $y'$ 当 $x = 0$ 和 $y = 0$ 时的值。
\end{prob}

\begin{prob}[S]{3380}
设 $x^2 + xy + y^2 = 3$，求 $y', y''$ 及 $y'''$。
\end{prob}

\begin{prob}[M]{3381}
设
\[
x^2 - xy + 2y^2 + x - y - 1 = 0,
\]
求 $y', y''$ 及 $y'''$ 当 $x = 0, y = 1$ 时的值。
\end{prob}

\begin{prob}[M]{3382}
证明：对于二次曲线
\[
ax^2 + 2bxy + cy^2 + 2dx + 2ey + f = 0
\]
成立等式
\[
\frac{\mathrm{d}^3}{\mathrm{d}x^3}\left[(y'')^{-\frac{2}{3}}\right] = 0.
\]
\end{prob}

\probgroup{求函数 $z = z(x, y)$ 的一阶和二阶偏导数，设：}

\begin{prob}[S]{3383}
$x^2 + y^2 + z^2 = a^2$。
\end{prob}

\begin{prob}[S]{3384}
$z^3 - 3xyz = a^3$。
\end{prob}

\begin{prob}[S]{3385}
$x + y + z = \mathrm{e}^z$。
\end{prob}

\begin{prob}[S]{3386}
$z = \sqrt{x^2 - y^2}\cdot\tan\dfrac{x}{\sqrt{x^2 - y^2}}$。
\end{prob}

\begin{prob}[S]{3387}
$x + y + z = \mathrm{e}^{-(x + y + z)}$。
\end{prob}

\begin{prob}[M]{3388}
设
\[
x^2 + y^2 + z^2 - 3xyz = 0 \tag{1}
\]
且
\[
f(x, y, z) = xy^2 z^3.
\]
求 $f'_x(1, 1, 1)$，若：(a) $z = z(x, y)$ 是由方程 (1) 定义的隐函数，(b) 若 $y = y(x, z)$ 是由方程 (1) 定义的隐函数。说明为什么这些导数相异。
\end{prob}

\begin{prob}[M]{3389}
设 $x^2 + 2y^2 + 3z^2 + xy - z - 9 = 0$，求 $\dfrac{\partial^2 z}{\partial x^2}, \dfrac{\partial^2 z}{\partial x \partial y}, \dfrac{\partial^2 z}{\partial y^2}$ 当 $x = 1, y = -2, z = 1$ 时的值。
\end{prob}

\probgroup{求 $\mathrm{d}z$ 和 $\mathrm{d}^2 z$，设}

\begin{prob}[S]{3390}
$\dfrac{x^2}{a^2} + \dfrac{y^2}{b^2} + \dfrac{z^2}{c^2} = 1$。
\end{prob}

\begin{prob}[S]{3391}
$xyz = x + y + z$。
\end{prob}

\begin{prob}[S]{3392}
$\dfrac{x}{z} = \ln\dfrac{z}{y} + 1$。
\end{prob}

\begin{prob}[S]{3393}
$z = x + \arctan\dfrac{y}{z - x}$。
\end{prob}

\begin{prob}[S]{3394}
设 $u^3 - 3(x + y)u^2 + z^3 = 0$，求 $\mathrm{d}u$。
\end{prob}

\begin{prob}[S]{3395}
设 $F(x + y + z, x^2 + y^2 + z^2) = 0$，求 $\dfrac{\partial^2 z}{\partial x \partial y}$。
\end{prob}

\begin{prob}[S]{3396}
设 $F(x - y, y - z, z - x) = 0$，求 $\dfrac{\partial z}{\partial x}$ 和 $\dfrac{\partial z}{\partial y}$。
\end{prob}

\begin{prob}[S]{3397}
设 $F(x, x + y, x + y + z) = 0$，求 $\dfrac{\partial z}{\partial x}, \dfrac{\partial z}{\partial y}$ 和 $\dfrac{\partial^2 z}{\partial x^2}$。
\end{prob}

\begin{prob}[S]{3398.1}
设 $F(xz, yz) = 0$，求 $\dfrac{\partial^2 z}{\partial x^2}$。
\end{prob}

\begin{prob}[M]{3398.2}
设 (a) $F(x + z, y + z) = 0$；\qquad (b) $F\left(\dfrac{x}{z}, \dfrac{y}{z}\right) = 0$，求 $\mathrm{d}^2 z$。
\end{prob}

\begin{prob}[M]{3399}
设 $z = z(x, y)$ 为由方程 $z^3 - xz + y = 0$ 定义的可微函数，且当 $x = 3, y = -2$ 时 $z = 2$。求 $\mathrm{d}z(3, -2)$ 和 $\mathrm{d}^2 z(3, -2)$。
\end{prob}

\begin{prob}[M]{3400}
设 $x = x(y, z), y = y(x, z), z = z(x, y)$ 为由方程 $F(x, y, z) = 0$ 定义的函数。证明：
\[
\frac{\partial x}{\partial y}\frac{\partial y}{\partial z}\frac{\partial z}{\partial x} = -1.
\]
\end{prob}

\begin{prob}[S]{3401}
设 $x + y + z = 0, x^2 + y^2 + z^2 = 1$，求 $\dfrac{\mathrm{d}x}{\mathrm{d}z}$ 和 $\dfrac{\mathrm{d}y}{\mathrm{d}z}$。
\end{prob}

\begin{prob}[M]{3402}
\begin{enumerate}[label=(\alph*), itemsep=2pt, topsep=3pt]
  \item 设 $x^2 + y^2 = \dfrac{1}{2}z^2, x + y + z = 2$，求 $\dfrac{\mathrm{d}x}{\mathrm{d}z}, \dfrac{\mathrm{d}y}{\mathrm{d}z}, \dfrac{\mathrm{d}^2 x}{\mathrm{d}z^2}$ 和 $\dfrac{\mathrm{d}^2 y}{\mathrm{d}z^2}$ 当 $x = 1, y = -1, z = 2$ 时的值。
  \item 设 $xu - yv = 0, yu + xv = 1$，求 $\dfrac{\partial u}{\partial x}, \dfrac{\partial u}{\partial y}, \dfrac{\partial v}{\partial x}$ 和 $\dfrac{\partial v}{\partial y}$。
\end{enumerate}
\end{prob}

\begin{prob}[M]{3403}
设由方程组
\[
\begin{cases}
x\mathrm{e}^{u + v} + 2uv = 1, \\
y\mathrm{e}^{u - v} - \dfrac{u}{1 + v} = 2x
\end{cases}
\]
定义的可微函数 $u = u(x, y)$ 和 $v = v(x, y)$ 满足 $u(1, 2) = 0$ 和 $v(1, 2) = 0$，求 $\mathrm{d}u(1, 2)$ 和 $\mathrm{d}v(1, 2)$。
\end{prob}

\begin{prob}[S]{3404}
设 $u + v = x + y,\quad \dfrac{\sin u}{\sin v} = \dfrac{x}{y}$，求 $\mathrm{d}u, \mathrm{d}v, \mathrm{d}^2 u$ 和 $\mathrm{d}^2 v$。
\end{prob}

\begin{prob}[M]{3405}
设
\[
\mathrm{e}^{\frac{u}{x}}\cos\frac{v}{y} = \frac{x}{\sqrt{2}}, \quad \mathrm{e}^{\frac{u}{x}}\sin\frac{v}{y} = \frac{y}{\sqrt{2}},
\]
求 $\mathrm{d}u, \mathrm{d}v, \mathrm{d}^2 u$ 和 $\mathrm{d}^2 v$ 当 $x = 1, y = 1, u = 0, v = \dfrac{\pi}{4}$ 时的表达式。
\end{prob}

\begin{prob}[M]{3406}
设
\[
x = t + t^{-1}, \quad y = t^2 + t^{-2}, \quad z = t^3 + t^{-3},
\]
求 $\frac{\mathrm{d}y}{\mathrm{d}x}, \frac{\mathrm{d}z}{\mathrm{d}x}, \frac{\mathrm{d}^2 y}{\mathrm{d}x^2}$ 和 $\frac{\mathrm{d}^2 z}{\mathrm{d}x^2}$.
\end{prob}

\begin{prob}[M]{3407}
在 $Oxy$ 平面上怎样的区域内，方程组
\[
x = u + v, \quad y = u^2 + v^2, \quad z = u^3 + v^3
\]
（式中参数 $u$ 和 $v$ 取一切可能的实数值）定义 $z$ 为变量 $x$ 和 $y$ 的函数？求导数 $\frac{\partial z}{\partial x}$ 和 $\frac{\partial z}{\partial y}$.
\end{prob}

\begin{prob}[L]{3408}
\begin{enumerate}[label=(\alph*), itemsep=2pt, topsep=3pt]
\item 设
\[
\begin{cases}
x = u + \ln v, \\
y = v - \ln u, \\
z = 2u + v,
\end{cases}
\]
求 $\frac{\partial z}{\partial x}$ 和 $\frac{\partial x}{\partial y}$ 在点 $u = 1, v = 1$ 的值.
\item 设
\[
\begin{cases}
x = u + v^2, \\
y = u^2 - v^3, \\
z = 2uv,
\end{cases}
\]
求 $\frac{\partial^2 x}{\partial x \partial y}$ 在点 $u = 2, v = 1$ 的值.
\item 设 $x = \cos\varphi\cos\psi, y = \cos\varphi\sin\psi, z = \sin\varphi$，求 $\frac{\partial^2 z}{\partial x^2}$.
\end{enumerate}
\end{prob}

\begin{prob}[S]{3409}
设 $x = u\cos v, y = u\sin v, z = v$，求 $\frac{\partial^2 z}{\partial x^2}, \frac{\partial^2 z}{\partial x\partial y}$ 及 $\frac{\partial^2 z}{\partial y^2}$.
\end{prob}

\begin{prob}[M]{3410}
设函数 $z = z(x, y)$ 由方程组
\[
\begin{cases}
x = \mathrm{e}^{u+v}, \\
y = \mathrm{e}^{u-v}, \\
z = uv
\end{cases}
\]
（$u$ 及 $v$ 为参数）定义，求当 $u = 0$ 及 $v = 0$ 时的 $\mathrm{d}z$ 及 $\mathrm{d}^2 z$.
\end{prob}

\begin{prob}[S]{3411}
设 $z = x^2 + y^2$，其中 $y = y(x)$ 为由方程 $x^2 - xy + y^2 = 1$ 定义的函数，求 $\frac{\mathrm{d}z}{\mathrm{d}x}$ 及 $\frac{\mathrm{d}^2 z}{\mathrm{d}x^2}$.
\end{prob}

\begin{prob}[S]{3412}
设 $u = \frac{x+z}{y+z}$，其中 $z$ 为由方程 $z\mathrm{e}^z = x\mathrm{e}^x + y\mathrm{e}^y$ 定义的函数，求 $\frac{\partial u}{\partial x}$ 及 $\frac{\partial u}{\partial y}$.
\end{prob}

\begin{prob}[M]{3413}
设方程
\[
x = \varphi(u, v), \quad y = \psi(u, v), \quad z = \chi(u, v)
\]
定义 $z$ 为 $x$ 和 $y$ 的函数. 求 $\frac{\partial z}{\partial x}$ 和 $\frac{\partial z}{\partial y}$.
\end{prob}

\begin{prob}[M]{3414}
设
\[
x = \varphi(u, v), \quad y = \psi(u, v),
\]
求反函数 $u = u(x, y)$ 和 $v = v(x, y)$ 的一阶和二阶偏导数.
\end{prob}

\begin{prob}[M]{3415}
设 (a) $x = u\cos\frac{v}{u}, y = u\sin\frac{v}{u}$；\qquad (b) $x = \mathrm{e}^u + u\sin v, y = \mathrm{e}^u - u\cos v$，求 $\frac{\partial u}{\partial x}, \frac{\partial u}{\partial y}, \frac{\partial v}{\partial x}, \frac{\partial v}{\partial y}$.
\end{prob}

\begin{prob}[M]{3416}
函数 $u = u(x)$ 由方程组
\[
\begin{cases}
u = f(x, y, z), \\
g(x, y, z) = 0, \\
h(x, y, z) = 0
\end{cases}
\]
定义，求 $\frac{\mathrm{d}u}{\mathrm{d}x}$ 和 $\frac{\mathrm{d}^2 u}{\mathrm{d}x^2}$.
\end{prob}

\begin{prob}[M]{3417}
函数 $u = u(x, y)$ 由方程组
\[
\begin{cases}
u = f(x, y, z, t), \\
g(y, z, t) = 0, \\
h(z, t) = 0
\end{cases}
\]
定义，求 $\frac{\partial u}{\partial x}$ 和 $\frac{\partial u}{\partial y}$.
\end{prob}

\begin{prob}[M]{3418}
设
\[
x = f(u, v, w), \quad y = g(u, v, w), \quad z = h(u, v, w),
\]
求 $\frac{\partial u}{\partial x}, \frac{\partial u}{\partial y}$ 和 $\frac{\partial u}{\partial z}$.
\end{prob}

\begin{prob}[M]{3419}
设函数 $z = z(x, y)$ 满足方程组
\[
\begin{cases}
f(x, y, z, t) = 0, \\
g(x, y, z, t) = 0,
\end{cases}
\]
式中 $t$ 为参变量，求 $\mathrm{d}z$.
\end{prob}

\begin{prob}[L]{3420}
设 $u = f(z)$，其中 $z = z(x, y)$ 为由方程 $z = x + y\varphi(z)$ 定义的隐函数. 证明拉格朗日公式：
\[
\frac{\partial^n u}{\partial y^n} = \frac{\partial^{n-1}}{\partial x^{n-1}}\left\{[\varphi(z)]^n \frac{\partial u}{\partial x}\right\}.
\]
\probhint{对 $n = 1$ 证明公式并运用数学归纳法.}
\end{prob}

\begin{prob}[M]{3421}
证明：由方程
\[
\Phi(x - az, y - bz) = 0 \qquad (1)
\]
（其中 $\Phi(u, v)$ 是变量 $u, v$ 的任意可微函数，$a$ 和 $b$ 为常数）定义的函数 $z = z(x, y)$ 为方程
\[
a\frac{\partial z}{\partial x} + b\frac{\partial z}{\partial y} = 1
\]
的解. 说明曲面 $(1)$ 的几何性质.
\end{prob}

\begin{prob}[M]{3422}
证明：由方程
\[
\Phi\left(\frac{x - x_0}{z - z_0}, \frac{y - y_0}{z - z_0}\right) = 0 \qquad (2)
\]
（其中 $\Phi(u, v)$ 是变量 $u$ 和 $v$ 的任意可微函数）定义的函数 $z = z(x, y)$ 满足方程
\[
(x - x_0)\frac{\partial z}{\partial x} + (y - y_0)\frac{\partial z}{\partial y} = z - z_0.
\]
说明曲面 $(2)$ 的几何性质.
\end{prob}

\begin{prob}[M]{3423}
证明：由方程
\[
ax + by + cz = \Phi(x^2 + y^2 + z^2) \qquad (3)
\]
（其中 $\Phi(u)$ 是变量 $u$ 的任意可微函数，$a, b$ 和 $c$ 为常数）定义的函数 $z = z(x, y)$ 满足方程
\[
(cy - bz)\frac{\partial z}{\partial x} + (az - cx)\frac{\partial z}{\partial y} = bx - ay.
\]
说明曲面 $(3)$ 的几何性质.
\end{prob}

\begin{prob}[M]{3424}
函数 $z = z(x, y)$ 由方程
\[
x^2 + y^2 + z^2 = yf\left(\frac{z}{y}\right)
\]
给出，证明：
\[
(x^2 - y^2 - z^2)\frac{\partial z}{\partial x} + 2xy\frac{\partial z}{\partial y} = 2xz.
\]
\end{prob}

\begin{prob}[M]{3425}
函数 $z = z(x, y)$ 由方程
\[
F(x + zy^{-1}, y + zx^{-1}) = 0
\]
给出. 证明：
\[
x\frac{\partial z}{\partial x} + y\frac{\partial z}{\partial y} = z - xy.
\]
\end{prob}

\begin{prob}[M]{3426}
证明：由方程组
\[
\begin{cases}
x\cos\alpha + y\sin\alpha + \ln z = f(\alpha), \\
-x\sin\alpha + y\cos\alpha = f'(\alpha)
\end{cases}
\]
（其中 $\alpha = \alpha(x, y)$ 为参变量，$f(\alpha)$ 为任意可微函数）定义的函数 $z = z(x, y)$ 满足方程
\[
\left(\frac{\partial z}{\partial x}\right)^2 + \left(\frac{\partial z}{\partial y}\right)^2 = z^2.
\]
\end{prob}

\begin{prob}[M]{3427}
证明：由方程组
\[
\begin{cases}
z = \alpha x + \dfrac{y}{\alpha} + f(\alpha), \\
0 = x - \dfrac{y}{\alpha^2} + f'(\alpha)
\end{cases}
\]
给出的函数 $z = z(x, y)$ 满足方程
\[
\frac{\partial z}{\partial x}\frac{\partial z}{\partial y} = 1.
\]
\end{prob}

\begin{prob}[M]{3428}
证明：由方程组
\[
\begin{cases}
[z - f(\alpha)]^2 = x^2(y^2 - \alpha^2), \\
[z - f(\alpha)] f'(\alpha) = \alpha x^2
\end{cases}
\]
给出的函数 $z = z(x, y)$ 满足方程
\[
\frac{\partial z}{\partial x}\frac{\partial z}{\partial y} = xy.
\]
\end{prob}

\begin{prob}[M]{3429}
证明：由方程组
\[
\begin{cases}
z = \alpha x + y\varphi(\alpha) + \psi(\alpha), \\
0 = x + y\varphi'(\alpha) + \psi'(\alpha)
\end{cases}
\]
给出的函数 $z = z(x, y)$ 满足方程
\[
\frac{\partial^2 z}{\partial x^2}\frac{\partial^2 z}{\partial y^2} - \left(\frac{\partial^2 z}{\partial x\partial y}\right)^2 = 0.
\]
\end{prob}

\begin{prob}[M]{3430}
证明：由方程
\[
y = x\varphi(z) + \psi(z)
\]
定义的隐函数 $z = z(x, y)$ 满足方程
\[
\left(\frac{\partial z}{\partial y}\right)^2 \frac{\partial^2 z}{\partial x^2} - 2\frac{\partial z}{\partial x}\frac{\partial z}{\partial y}\frac{\partial^2 z}{\partial x\partial y} + \left(\frac{\partial z}{\partial x}\right)^2 \frac{\partial^2 z}{\partial y^2} = 0.
\]
\end{prob}

\probsec{§4. 变量代换}

\begin{prob}[S]{3431}
把 $y$ 看作新的自变量，变换方程
\[
y'y''' - 3y''^2 = x.
\]
\end{prob}

\begin{prob}[S]{3432}
用同样的方法变换方程
\[
y'^2 y^{(4)} - 10y'y''y''' + 15y''^3 = 0.
\]
\end{prob}

\begin{prob}[S]{3433}
把 $x$ 看作函数，把 $t = xy$ 看作自变量，变换方程
\[
y'' + \frac{2}{x}y' + y = 0.
\]
\end{prob}

\probgroup{引入新变量，变换下列常微分方程：}

\begin{prob}[S]{3434}
$x^2 y'' + xy' + y = 0$，若 $x = \mathrm{e}^t$.
\end{prob}

\begin{prob}[S]{3435}
$y''' = \frac{6y}{x^3}$，若 $t = \ln |x|$.
\end{prob}

\begin{prob}[S]{3436}
$(1 - x^2)y'' - xy' + n^2 y = 0$，若 $x = \cos t$.
\end{prob}

\begin{prob}[S]{3437}
$y'' + y'\tanh x + \frac{m^2}{\cosh^2 x}y = 0$，若 $x = \ln\tan\frac{t}{2}$.
\end{prob}

\begin{prob}[S]{3438}
$y'' + p(x)y' + q(x)y = 0$，若 $y = u\mathrm{e}^{-\frac{1}{2}\int_{x_0}^x p(\xi)\,\mathrm{d}\xi}$，其中 $p(x) \in C^{(1)}$.
\end{prob}

\begin{prob}[S]{3439}
$x^4 y'' + xyy' - 2y^2 = 0$，若 $x = \mathrm{e}^t, y = u\mathrm{e}^{2t}$，其中 $u = u(t)$.
\end{prob}

\begin{prob}[S]{3440}
$(1 + x^2)^2 y'' = y$，若 $x = \tan t, y = \frac{u}{\cos t}$，其中 $u = u(t)$.
\end{prob}

\begin{prob}[S]{3441}
$(1 - x^2)^2 y'' = -y$，若 $x = \tanh t, y = \frac{u}{\cosh t}$，其中 $u = u(t)$.
\end{prob}

\begin{prob}[S]{3442}
$y'' + (x + y)(1 + y')^3 = 0$，若 $x = u + t, y = u - t$，其中 $u = u(t)$.
\end{prob}

\begin{prob}[S]{3443}
$y''' - x^3 y'' + xy' - y = 0$，若 $x = \frac{1}{t}, y = \frac{u}{t}$，其中 $u = u(t)$.
\end{prob}

\begin{prob}[M]{3444}
令
\[
u = \frac{y}{x - b}, \quad t = \ln\left|\frac{x - a}{x - b}\right|,
\]
并把 $u$ 看作变量 $t$ 的函数，以变换斯托克斯方程
\[
y'' = \frac{Ay}{(x - a)^2 (x - b)^2}.
\]
\end{prob}

\begin{prob}[M]{3445}
证明：若方程
\[
\frac{\mathrm{d}^2 y}{\mathrm{d}x^2} + p(x)\frac{\mathrm{d}y}{\mathrm{d}x} + q(x)y = 0
\]
由代换 $x = \varphi(\xi)$ 变换为方程
\[
\frac{\mathrm{d}^2 y}{\mathrm{d}\xi^2} + P(\xi)\frac{\mathrm{d}y}{\mathrm{d}\xi} + Q(\xi)y = 0,
\]
即
\[
[2P(\xi)Q(\xi) + Q'(\xi)][Q(\xi)]^{-\frac{3}{2}} = [2p(x)q(x) + q'(x)][q(x)]^{-\frac{3}{2}}.
\]
\end{prob}

\begin{prob}[S]{3446}
在方程
\[
\Phi(y, y', y'') = 0
\]
（其中 $\Phi$ 为变量 $y, y', y''$ 的齐次函数）中令 $y = \mathrm{e}^{\int_{x_0}^x u\,\mathrm{d}x}$.
\end{prob}

\begin{prob}[M]{3447}
在方程
\[
F(x^2 y'', xy', y) = 0
\]
(其中 $F$ 为其自变量的齐次函数) 中令 $u = x\frac{y'}{y}$.
\end{prob}

\begin{prob}[L]{3448}
证明: 方程 $y'''(1+y'^2)-3y'y''^2=0$ 的形式在单应变换
\[
x = \frac{a_1\xi+b_1\eta+c_1}{a\xi+b\eta+c}, \quad y = \frac{a_2\xi+b_2\eta+c_2}{a\xi+b\eta+c}
\]
下保持不变.
\probhint{该变换是由以下最简单的变换构成的:
\[
x = \alpha X + \beta Y + \gamma, \quad y = Y;
\]
\[
X = \frac{1}{X_1}, \quad Y = \frac{Y_1}{X_1};
\]
\[
X_1 = a\xi + b\eta + c, \quad Y_1 = a_2\xi + b_2\eta + c_2.
\]}
\end{prob}

\begin{prob}[M]{3449}
证明: 施瓦茨函数
\[
S[x(t)] = \frac{x'''(t)}{x'(t)} - \frac{3}{2}\left[\frac{x''(t)}{x'(t)}\right]^2
\]
的值在分式线性变换
\[
y = \frac{ax(t)+b}{cx(t)+d} \quad (ad - bc \neq 0)
\]
下保持不变.
\end{prob}

\probgroup{令 $x = r\cos\varphi, y = r\sin\varphi$, 写出下列方程在极坐标 $r, \varphi$ 下的形式:}

\begin{prob}[S]{3450}
$\frac{\mathrm{d}y}{\mathrm{d}x} = \frac{x+y}{x-y}$.
\end{prob}

\begin{prob}[S]{3451}
$(xy'-y)^2 = 2xy(1+y'^2)$.
\end{prob}

\begin{prob}[S]{3452}
$(x^2+y^2)^2 y'' = (x+yy')^3$.
\end{prob}

\begin{prob}[S]{3453}
写出表达式 $\frac{x+yy'}{xy'-y}$ 在极坐标下的形式.
\end{prob}

\begin{prob}[M]{3454}
把平面曲线的曲率
\[
K = \frac{|y''_{xx}|}{(1+y'^2_x)^{\frac{3}{2}}}
\]
用极坐标 $r$ 和 $\varphi$ 表示出来.
\end{prob}

\begin{prob}[M]{3455}
写出方程组
\[
\frac{\mathrm{d}x}{\mathrm{d}t} = y + kx(x^2+y^2), \quad \frac{\mathrm{d}y}{\mathrm{d}t} = -x + ky(x^2+y^2)
\]
在极坐标下的形式.
\end{prob}

\begin{prob}[M]{3456}
引入新函数 $r = \sqrt{x^2+y^2}, \varphi = \arctan\frac{y}{x}$, 以便变换表达式
\[
W = x\frac{\mathrm{d}^2 y}{\mathrm{d}t^2} - y\frac{\mathrm{d}^2 x}{\mathrm{d}t^2}.
\]
\end{prob}

\begin{prob}[M]{3457}
在勒让德变换中, 曲线 $y=y(x)$ 的每一点 $(x, y)$ 对应于点 $(X, Y)$, 其中
\[
X = y', \quad Y = xy' - y.
\]
求 $Y', Y''$ 及 $Y'''$.
\end{prob}

\probgroup{引入新变量 $\xi$ 及 $\eta$, 解下列方程:}

\begin{prob}[S]{3458}
$\frac{\partial z}{\partial x} = \frac{\partial z}{\partial y}$, 令 $\xi = x+y$ 和 $\eta = x-y$.
\end{prob}

\begin{prob}[S]{3459}
$y\frac{\partial z}{\partial x} - x\frac{\partial z}{\partial y} = 0$, 令 $\xi = x, \eta = x^2+y^2$.
\end{prob}

\begin{prob}[S]{3460}
$a\frac{\partial z}{\partial x} + b\frac{\partial z}{\partial y} = 1$ ($a \neq 0$), 令 $\xi = x, \eta = y - bz$.
\end{prob}

\begin{prob}[S]{3461}
$x\frac{\partial z}{\partial x} + y\frac{\partial z}{\partial y} = z$, 令 $\xi = x$ 及 $\eta = \frac{y}{x}$.
\end{prob}

\probgroup{把 $u$ 与 $v$ 看作新的自变量, 变换下列方程:}

\begin{prob}[S]{3462}
$x\frac{\partial z}{\partial x} + \sqrt{1+y^2}\frac{\partial z}{\partial y} = xy$, 若 $u = \ln x, v = \ln(y+\sqrt{1+y^2})$.
\end{prob}

\begin{prob}[S]{3463}
$(x+y)\frac{\partial z}{\partial x} - (x-y)\frac{\partial z}{\partial y} = 0$, 若 $u = \ln\sqrt{x^2+y^2}, v = \arctan\frac{y}{x}$.
\end{prob}

\begin{prob}[S]{3464}
$x\frac{\partial z}{\partial x} + y\frac{\partial z}{\partial y} = z + \sqrt{x^2+y^2+z^2}$, 若 $u = \frac{y}{x}, v = z + \sqrt{x^2+y^2+z^2}$.
\end{prob}

\begin{prob}[S]{3465}
$x\frac{\partial z}{\partial x} + y\frac{\partial z}{\partial y} = \frac{x}{z}$, 若 $u = 2x - z^2, v = \frac{y}{z}$.
\end{prob}

\begin{prob}[S]{3466}
$(x+z)\frac{\partial z}{\partial x} + (y+z)\frac{\partial z}{\partial y} = x+y+z$, 若 $u = x+z, v = y+z$.
\end{prob}

\begin{prob}[M]{3467}
取
\[
\xi = y + z\mathrm{e}^{-x}, \quad \eta = x + z\mathrm{e}^{-y}
\]
作为新的自变量, 变换表达式
\[
(z+\mathrm{e}^x)\frac{\partial z}{\partial x} + (z+\mathrm{e}^y)\frac{\partial z}{\partial y} - (z^2 - \mathrm{e}^{x+y}).
\]
\end{prob}

\begin{prob}[M]{3468}
令
\[
x = uv, \quad y = \frac{1}{2}(u^2 - v^2),
\]
变换表达式
\[
\left(\frac{\partial z}{\partial x}\right)^2 + \left(\frac{\partial z}{\partial y}\right)^2.
\]
\end{prob}

\begin{prob}[M]{3469}
令 $\xi = x, \eta = y - x, \zeta = z - x$, 变换方程
\[
\frac{\partial u}{\partial x} + \frac{\partial u}{\partial y} + \frac{\partial u}{\partial z} = 0.
\]
\end{prob}

\begin{prob}[M]{3470}
取 $x$ 作为函数, 而 $y$ 和 $z$ 作为自变量, 变换方程
\[
(x-z)\frac{\partial z}{\partial x} + y\frac{\partial z}{\partial y} = 0.
\]
\end{prob}

\begin{prob}[M]{3471}
取 $x$ 作为函数, 而 $u = y-z, v = y+z$ 作为自变量, 变换方程
\[
(y-z)\frac{\partial z}{\partial x} + (y+z)\frac{\partial z}{\partial y} = 0.
\]
\end{prob}

\begin{prob}[M]{3472}
取 $x$ 作为函数, 而 $u = xz, v = yz$ 作为自变量, 变换表达式
\[
A = \left(\frac{\partial z}{\partial x}\right)^2 + \left(\frac{\partial z}{\partial y}\right)^2.
\]
\end{prob}

\begin{prob}[M]{3473}
令 $\mathrm{e}^{\xi} = x - u, \mathrm{e}^{\eta} = y - u, \mathrm{e}^{\zeta} = z - u$, 变换方程
\[
(y+z+u)\frac{\partial u}{\partial x} + (x+z+u)\frac{\partial u}{\partial y} + (x+y+u)\frac{\partial u}{\partial z} = x+y+z.
\]
\end{prob}

\probgroup{在下列方程中, 代入新的变量 $u, v, w$, 其中 $w = w(u, v)$:}

\begin{prob}[S]{3474}
$y\frac{\partial z}{\partial x} - x\frac{\partial z}{\partial y} = (y-x)z$, 令 $u = x^2+y^2, v = \frac{1}{x} + \frac{1}{y}, w = \ln z - (x+y)$.
\end{prob}

\begin{prob}[S]{3475}
$x^2\frac{\partial z}{\partial x} + y^2\frac{\partial z}{\partial y} = z^2$, 令 $u = x, v = \frac{1}{y} - \frac{1}{x}, w = \frac{1}{z} - \frac{1}{x}$.
\end{prob}

\begin{prob}[S]{3476}
$(xy+z)\frac{\partial z}{\partial x} + (1-y^2)\frac{\partial z}{\partial y} = x+yz$, 令 $u = yz-x, v = xz-y, w = xy-z$.
\end{prob}

\begin{prob}[S]{3477}
$\left(x\frac{\partial z}{\partial x}\right)^2 + \left(y\frac{\partial z}{\partial y}\right)^2 = z^2 \frac{\partial z}{\partial x}\cdot\frac{\partial z}{\partial y}$, 令 $x = u\mathrm{e}^w, y = v\mathrm{e}^w, z = w\mathrm{e}^w$.
\end{prob}

\begin{prob}[M]{3478}
令 $u = \ln\sqrt{x^2+y^2}, v = \arctan z, w = x+y+z$, 其中 $w = w(u, v)$, 变换表达式
\[
\frac{x-y}{\dfrac{\partial z}{\partial x} - \dfrac{\partial z}{\partial y}}.
\]
\end{prob}

\begin{prob}[M]{3479}
令 $u = x\mathrm{e}^z, v = y\mathrm{e}^z, w = z\mathrm{e}^z$, 其中 $w = w(u, v)$, 变换表达式
\[
A = \frac{\dfrac{\partial z}{\partial x}}{\dfrac{\partial z}{\partial y}}.
\]
\end{prob}

\begin{prob}[M]{3480}
令 $\xi = \frac{x}{z}, \eta = \frac{y}{z}, \zeta = z, w = \frac{u}{z}$, 其中 $w = w(\xi, \eta, \zeta)$, 变换方程
\[
x\frac{\partial u}{\partial x} + y\frac{\partial u}{\partial y} + z\frac{\partial u}{\partial z} = u + \frac{xy}{z}.
\]
\end{prob}

\probgroup{令 $x = r\cos\varphi, y = r\sin\varphi$, 写出下列各式在极坐标 $r$ 和 $\varphi$ 下的形式:}

\begin{prob}[S]{3481}
$w = x\frac{\partial u}{\partial y} - y\frac{\partial u}{\partial x}$.
\end{prob}

\begin{prob}[S]{3482}
$w = x\frac{\partial u}{\partial x} + y\frac{\partial u}{\partial y}$.
\end{prob}

\begin{prob}[S]{3483}
$w = \left(\frac{\partial u}{\partial x}\right)^2 + \left(\frac{\partial u}{\partial y}\right)^2$.
\end{prob}

\begin{prob}[S]{3484}
$w = \frac{\partial^2 u}{\partial x^2} + \frac{\partial^2 u}{\partial y^2}$.
\end{prob}

\begin{prob}[S]{3485}
$w = x^2\frac{\partial^2 u}{\partial x^2} + 2xy\frac{\partial^2 u}{\partial x\partial y} + y^2\frac{\partial^2 u}{\partial y^2}$.
\end{prob}

\begin{prob}[S]{3486}
$w = y^2\frac{\partial^2 z}{\partial x^2} - 2xy\frac{\partial^2 z}{\partial x\partial y} + x^2\frac{\partial^2 z}{\partial y^2} - \left(x\frac{\partial z}{\partial x} + y\frac{\partial z}{\partial y}\right)$.
\end{prob}

\begin{prob}[M]{3487}
令 $x = r\cos\varphi, y = r\sin\varphi$, 变换表达式
\[
I = \frac{\partial u}{\partial x}\frac{\partial v}{\partial y} - \frac{\partial u}{\partial y}\frac{\partial v}{\partial x}.
\]
\end{prob}

\begin{prob}[M]{3488}
引入新的自变量
\[
\xi = x - at, \quad \eta = x + at,
\]
解方程
\[
\frac{\partial^2 u}{\partial t^2} = a^2\frac{\partial^2 u}{\partial x^2}.
\]
\end{prob}

\probgroup{取 $u$ 及 $v$ 作新的自变量, 变换下列方程:}

\begin{prob}[S]{3489}
$2\frac{\partial^2 z}{\partial x^2} + \frac{\partial^2 z}{\partial x\partial y} - \frac{\partial^2 z}{\partial y^2} + \frac{\partial z}{\partial x} + \frac{\partial z}{\partial y} = 0$, 设 $u = x + 2y + 2$ 及 $v = x - y - 1$.
\end{prob}

\begin{prob}[S]{3490}
$(1+x^2)\frac{\partial^2 z}{\partial x^2} + (1+y^2)\frac{\partial^2 z}{\partial y^2} + x\frac{\partial z}{\partial x} + y\frac{\partial z}{\partial y} = 0$, 设 $u = \ln(x+\sqrt{1+x^2})$ 及 $v = \ln(y+\sqrt{1+y^2})$.
\end{prob}

\begin{prob}[S]{3491}
$ax^2\frac{\partial^2 z}{\partial x^2} + 2bxy\frac{\partial^2 z}{\partial x\partial y} + cy^2\frac{\partial^2 z}{\partial y^2} = 0$ ($a, b, c$ 为常数), 设 $u = \ln x, v = \ln y$.
\end{prob}

\begin{prob}[S]{3492}
$\frac{\partial^2 z}{\partial x^2} + \frac{\partial^2 z}{\partial y^2} = 0$, 设 $u = \frac{x}{x^2+y^2}, v = -\frac{y}{x^2+y^2}$.
\end{prob}

\begin{prob}[S]{3493}
$\frac{\partial^2 z}{\partial x^2} + \frac{\partial^2 z}{\partial y^2} + m^2 z = 0$, 设 $x = \mathrm{e}^u\cos v, y = \mathrm{e}^u\sin v$.
\end{prob}

\begin{prob}[S]{3494}
$\frac{\partial^2 z}{\partial x^2} - y\frac{\partial^2 z}{\partial y^2} = \frac{1}{2}\frac{\partial z}{\partial y}$ ($y > 0$), 设 $u = x - 2\sqrt{y}$ 及 $v = x + 2\sqrt{y}$.
\end{prob}

\begin{prob}[S]{3495}
$x^2\frac{\partial^2 z}{\partial x^2} - y^2\frac{\partial^2 z}{\partial y^2} = 0$, 设 $u = xy, v = \frac{x}{y}$.
\end{prob}

\begin{prob}[S]{3496}
$x^2\frac{\partial^2 z}{\partial x^2} - (x^2+y^2)\frac{\partial^2 z}{\partial x\partial y} + y^2\frac{\partial^2 z}{\partial y^2} = 0$, 设 $u = x+y, v = \frac{1}{x} + \frac{1}{y}$.
\end{prob}

\begin{prob}[S]{3497}
$xy\frac{\partial^2 z}{\partial x^2} - (x^2+y^2)\frac{\partial^2 z}{\partial x\partial y} + xy\frac{\partial^2 z}{\partial y^2} + y\frac{\partial z}{\partial x} + x\frac{\partial z}{\partial y} = 0$, 设 $u = \frac{1}{2}(x^2+y^2), v = xy$.
\end{prob}

\begin{prob}[S]{3498}
$x^2\frac{\partial^2 z}{\partial x^2} - 2x\sin y\frac{\partial^2 z}{\partial x\partial y} + \sin^2 y\frac{\partial^2 z}{\partial y^2} = 0$, 设 $u = x\tan\frac{y}{2}, v = x$.
\end{prob}

\begin{prob}[S]{3499}
$x\frac{\partial^2 z}{\partial x^2} - y\frac{\partial^2 z}{\partial y^2} = 0$ ($x > 0, y > 0$), 设 $x = (u+v)^2$ 及 $y = (u-v)^2$.
\end{prob}

\begin{prob}[S]{3500}
$\frac{\partial^2 z}{\partial x\partial y} = \left(1+\frac{\partial z}{\partial y}\right)^3$, 设 $u = x$ 及 $v = y+z$.
\end{prob}

\begin{prob}[M]{3501}
利用线性变换
\[
\xi = x + \lambda_1 y, \quad \eta = x + \lambda_2 y,
\]
把方程
\[
A\frac{\partial^2 u}{\partial x^2} + 2B\frac{\partial^2 u}{\partial x\partial y} + C\frac{\partial^2 u}{\partial y^2} = 0 \tag{1}
\]
(其中 $A, B$ 和 $C$ 为常数, 且 $AC - B^2 < 0$) 变换为以下形式:
\[
\frac{\partial^2 u}{\partial\xi\partial\eta} = 0.
\]
求满足方程 (1) 的函数的一般形式.
\end{prob}

\begin{prob}[M]{3502}
证明: 拉普拉斯方程
\[
\Delta z \equiv \frac{\partial^2 z}{\partial x^2} + \frac{\partial^2 z}{\partial y^2} = 0
\]
的形式在满足条件
\[
\frac{\partial \varphi}{\partial u} = \frac{\partial \psi}{\partial v}, \quad \frac{\partial \varphi}{\partial v} = -\frac{\partial \psi}{\partial u}
\]
的任何非退化变换
\[
x = \varphi(u, v), \quad y = \psi(u, v)
\]
下保持不变.
\end{prob}

\begin{prob}[M]{3503}
令 $u = f(r)$, 其中 $r = \sqrt{x^2+y^2}$, 变换方程:

(a) $\Delta u \equiv \frac{\partial^2 u}{\partial x^2} + \frac{\partial^2 u}{\partial y^2} = 0$; \qquad (b) $\Delta(\Delta u) = 0$.
\end{prob}

\begin{prob}[M]{3504}
若令
\[
w = f(u)
\]
其中 $u = (x-x_0)(y-y_0)$, 则方程
\[
\frac{\partial^2 w}{\partial x\partial y} + cw = 0
\]
化为何种形式?
\end{prob}

\begin{prob}[M]{3505}
令
\[
x+y = X, \quad y = XY,
\]
变换表达式
\[
A = x\frac{\partial^2 u}{\partial x^2} + y\frac{\partial^2 u}{\partial x\partial y} + \frac{\partial u}{\partial x}.
\]
\end{prob}

\begin{prob}[M]{3506}
证明: 方程
\[
\frac{\partial^2 z}{\partial x^2} + 2xy^2\frac{\partial z}{\partial x} + 2(y-y^3)\frac{\partial z}{\partial y} + x^2 y^2 z = 0
\]
的形式在变换
\[
x = uv, \quad y = \frac{1}{v}
\]
下保持不变.
\end{prob}

\begin{prob}[M]{3507}
证明: 方程
\[
\frac{\partial^2 z}{\partial x^2} + 2\frac{\partial^2 z}{\partial x\partial y} + \frac{\partial^2 z}{\partial y^2} = 0
\]
的形式在变换
\[
u = x+z, \quad v = y+z
\]
下保持不变.
\end{prob}

\begin{prob}[M]{3508}
令
\[
x = \eta\zeta, \quad y = \xi\zeta, \quad z = \xi\eta,
\]
变换方程
\[
xy\frac{\partial^2 u}{\partial x\partial y} + yz\frac{\partial^2 u}{\partial y\partial z} + xz\frac{\partial^2 u}{\partial x\partial z} = 0.
\]
\end{prob}

\begin{prob}[M]{3509}
令
\[
y_1 = x_2+x_3-x_1, \quad y_2 = x_1+x_3-x_2, \quad y_3 = x_1+x_2-x_3,
\]
变换方程
\[
\frac{\partial^2 z}{\partial x_1^2} + \frac{\partial^2 z}{\partial x_2^2} + \frac{\partial^2 z}{\partial x_3^2} + \frac{\partial^2 z}{\partial x_1\partial x_2} + \frac{\partial^2 z}{\partial x_1\partial x_3} + \frac{\partial^2 z}{\partial x_2\partial x_3} = 0.
\]
\end{prob}

\begin{prob}[M]{3510}
令
\[
\xi = \frac{y}{x}, \quad \eta = \frac{z}{x}, \quad \zeta = y - z,
\]
变换方程
\[
x^2\frac{\partial^2 u}{\partial x^2} + y^2\frac{\partial^2 u}{\partial y^2} + z^2\frac{\partial^2 u}{\partial z^2} + 2xy\frac{\partial^2 u}{\partial x\partial y} + 2xz\frac{\partial^2 u}{\partial x\partial z} + 2yz\frac{\partial^2 u}{\partial y\partial z} = 0.
\]
\probhint{把方程写为 $A^2 u - Au = 0$ 的形式，其中
\[
A = x\frac{\partial}{\partial x} + y\frac{\partial}{\partial y} + z\frac{\partial}{\partial z}.
\]}
\end{prob}

\begin{prob}[M]{3511}
令
\[
x = r\sin\theta\cos\varphi, \quad y = r\sin\theta\sin\varphi, \quad z = r\cos\theta,
\]
把表达式
\[
\Delta_1 u = \left(\frac{\partial u}{\partial x}\right)^2 + \left(\frac{\partial u}{\partial y}\right)^2 + \left(\frac{\partial u}{\partial z}\right)^2, \quad \Delta_2 u = \frac{\partial^2 u}{\partial x^2} + \frac{\partial^2 u}{\partial y^2} + \frac{\partial^2 u}{\partial z^2}
\]
变换为球坐标下的形式.
\probhint{变换由两个特殊的变换构成：
\[
\begin{aligned}
x &= R\cos\varphi, & y &= R\sin\varphi, & z &= z; \\
R &= r\sin\theta, & \varphi &= \varphi, & z &= r\cos\theta.
\end{aligned}
\]}
\end{prob}

\begin{prob}[M]{3512}
引入新函数 $w$，令 $w = z^2$，变换方程
\[
z\left(\frac{\partial^2 z}{\partial x^2} + \frac{\partial^2 z}{\partial y^2}\right) = \left(\frac{\partial z}{\partial x}\right)^2 + \left(\frac{\partial z}{\partial y}\right)^2.
\]
\end{prob}

\probgroup{取 $u$ 和 $v$ 为新的自变量，而 $w = w(u, v)$ 为新函数，变换下列方程：}

\begin{prob}[S]{3513}
$y\dfrac{\partial^2 z}{\partial y^2} + 2\dfrac{\partial z}{\partial y} = \dfrac{2}{x}$，设 $u = \dfrac{x}{y}, v = x, w = xz - y$.
\end{prob}

\begin{prob}[S]{3514}
$\dfrac{\partial^2 z}{\partial x^2} - 2\dfrac{\partial^2 z}{\partial x\partial y} + \dfrac{\partial^2 z}{\partial y^2} = 0$，设 $u = x + y, v = \dfrac{y}{x}, w = \dfrac{z}{x}$.
\end{prob}

\begin{prob}[S]{3515}
$\dfrac{\partial^2 z}{\partial x^2} + 2\dfrac{\partial^2 z}{\partial x\partial y} + \dfrac{\partial^2 z}{\partial y^2} = 0$，设 $u = x + y, v = x - y, w = xy - z$.
\end{prob}

\begin{prob}[S]{3516}
$\dfrac{\partial^2 z}{\partial x^2} + \dfrac{\partial^2 z}{\partial x\partial y} + \dfrac{\partial z}{\partial x} = z$，设 $u = \dfrac{x+y}{2}, v = \dfrac{x-y}{2}, w = z\mathrm{e}^y$.
\end{prob}

\begin{prob}[S]{3517}
$\dfrac{\partial^2 z}{\partial x^2} - 2\dfrac{\partial^2 z}{\partial x\partial y} + \left(1 + \dfrac{y}{x}\right)\dfrac{\partial^2 z}{\partial y^2} = 0$，设 $u = x, v = x + y, w = x + y + z$.
\end{prob}

\begin{prob}[S]{3518}
$(1 - x^2)\dfrac{\partial^2 z}{\partial x^2} + (1 - y^2)\dfrac{\partial^2 z}{\partial y^2} = x\dfrac{\partial z}{\partial x} + y\dfrac{\partial z}{\partial y}$，设 $x = \sin u, y = \sin v, z = \mathrm{e}^w$.
\end{prob}

\begin{prob}[M]{3519}
$(1 - x^2)\dfrac{\partial^2 z}{\partial x^2} - \dfrac{\partial^2 z}{\partial y^2} - 2x\dfrac{\partial z}{\partial x} - \dfrac{1}{4}z = 0$ ($|x| < 1$)，设
\[
u = \frac{1}{2}(y + \arccos x), \quad v = \frac{1}{2}(y - \arccos x), \quad w = z\sqrt[4]{1 - x^2}.
\]
\end{prob}

\begin{prob}[M]{3520}
$\dfrac{\partial^2 z}{\partial x^2} + \dfrac{\partial^2 z}{\partial y^2} = 2\dfrac{x\dfrac{\partial z}{\partial x} - y\dfrac{\partial z}{\partial y}}{x^2 - y^2} - \dfrac{3(x^2 + y^2)z}{(x^2 - y^2)^2}$ ($|x| > |y|$)，设
\[
u = x + y, \quad v = x - y, \quad w = \frac{z}{\sqrt{x^2 - y^2}}.
\]
\end{prob}

\begin{prob}[M]{3521}
证明：任何方程
\[
\frac{\partial^2 z}{\partial x\partial y} + a\frac{\partial z}{\partial x} + b\frac{\partial z}{\partial y} + cz = 0
\]
($a, b, c$ 为常数) 用代换
\[
z = u\mathrm{e}^{\alpha x + \beta y}
\]
(其中 $\alpha$ 与 $\beta$ 为常量，$u = u(x, y)$) 可以化为以下形式：
\[
\frac{\partial^2 u}{\partial x\partial y} + c_1 u = 0 \quad (c_1 \text{ 为常数}).
\]
\end{prob}

\begin{prob}[M]{3522}
证明：方程
\[
\frac{\partial^2 u}{\partial x^2} = \frac{\partial u}{\partial y}
\]
的形式在变换
\[
x' = \frac{x}{y}, \quad y' = -\frac{1}{y}, \quad u' = \frac{u}{\sqrt{y}}\mathrm{e}^{-\frac{x^2}{4y}}
\]
($u'$ 为变量 $x'$ 与 $y'$ 的函数) 下保持不变.
\end{prob}

\begin{prob}[M]{3523}
令 $u = x + z, v = y + z, w = x + y + z$，且认为 $w = w(u, v)$，变换方程
\[
q(1 + q)\frac{\partial^2 z}{\partial x^2} - (1 + p + q + 2pq)\frac{\partial^2 z}{\partial x\partial y} + p(1 + p)\frac{\partial^2 z}{\partial y^2} = 0,
\]
其中 $p = \dfrac{\partial z}{\partial x}, q = \dfrac{\partial z}{\partial y}$.
\end{prob}

\begin{prob}[M]{3524}
令 $x = \mathrm{e}^\xi, y = \mathrm{e}^\eta, z = \mathrm{e}^\zeta, u = \mathrm{e}^w$，其中 $w = w(\xi, \eta, \zeta)$，变换方程
\[
x^2\frac{\partial^2 u}{\partial x^2} + y^2\frac{\partial^2 u}{\partial y^2} + z^2\frac{\partial^2 u}{\partial z^2} = \left(x\frac{\partial u}{\partial x}\right)^2 + \left(y\frac{\partial u}{\partial y}\right)^2 + \left(z\frac{\partial u}{\partial z}\right)^2.
\]
\end{prob}

\begin{prob}[M]{3525}
证明：方程
\[
\frac{\partial^2 z}{\partial x^2}\frac{\partial^2 z}{\partial y^2} - \left(\frac{\partial^2 z}{\partial x\partial y}\right)^2 = 0
\]
的形式与变量 $x, y$ 和 $z$ 所分别担任的角色无关.
\end{prob}

\begin{prob}[M]{3526}
取 $x$ 作为变量 $y$ 和 $z$ 的函数，解方程
\[
\left(\frac{\partial z}{\partial y}\right)^2 \frac{\partial^2 z}{\partial x^2} - 2\frac{\partial z}{\partial x}\frac{\partial z}{\partial y}\frac{\partial^2 z}{\partial x\partial y} + \left(\frac{\partial z}{\partial x}\right)^2 \frac{\partial^2 z}{\partial y^2} = 0.
\]
\end{prob}

\begin{prob}[M]{3527}
运用勒让德变换
\[
X = \frac{\partial z}{\partial x}, \quad Y = \frac{\partial z}{\partial y}, \quad Z = x\frac{\partial z}{\partial x} + y\frac{\partial z}{\partial y} - z,
\]
其中 $Z = Z(X, Y)$，变换方程
\[
A\left(\frac{\partial z}{\partial x}, \frac{\partial z}{\partial y}\right)\frac{\partial^2 z}{\partial x^2} + 2B\left(\frac{\partial z}{\partial x}, \frac{\partial z}{\partial y}\right)\frac{\partial^2 z}{\partial x\partial y} + C\left(\frac{\partial z}{\partial x}, \frac{\partial z}{\partial y}\right)\frac{\partial^2 z}{\partial y^2} = 0.
\]
\end{prob}

\probsec{§5. 几何上的应用}

\probgroup{写出下列曲线在已知点的切线和法平面方程：}

\begin{prob}[S]{3528}
$x = a\cos\alpha\cos t, y = a\sin\alpha\cos t, z = a\sin t$；在点 $t = t_0$.
\end{prob}

\begin{prob}[S]{3529}
$x = a\sin^2 t, y = b\sin t\cos t, z = c\cos^2 t$；在点 $t = \dfrac{\pi}{4}$.
\end{prob}

\begin{prob}[S]{3530}
$y = x, z = x^2$；在点 $M(1, 1, 1)$.
\end{prob}

\begin{prob}[S]{3531}
$x^2 + z^2 = 10, y^2 + z^2 = 10$；在点 $M(1, 1, 3)$.
\end{prob}

\begin{prob}[S]{3532}
$x^2 + y^2 + z^2 = 6, x + y + z = 0$；在点 $M(1, -2, 1)$.
\end{prob}

\begin{prob}[M]{3533}
在曲线 $x = t, y = t^2, z = t^3$ 上求一点，此处的切线平行于平面 $x + 2y + z = 4$.
\end{prob}

\begin{prob}[M]{3534}
证明：螺旋线 $x = a\cos t, y = a\sin t, z = bt$ 的切线与 $Oz$ 轴形成定角.
\end{prob}

\begin{prob}[M]{3535}
证明：曲线 $x = a\mathrm{e}^t\cos t, y = a\mathrm{e}^t\sin t, z = a\mathrm{e}^t$ 与锥面 $x^2 + y^2 = z^2$ 的各母线相交的角度相同.
\end{prob}

\begin{prob}[M]{3536}
证明：斜驶线
\[
\tan\left(\frac{\pi}{4} + \frac{\psi}{2}\right) = \mathrm{e}^{k\varphi} \quad (k = \text{常数})
\]
(其中 $\varphi$ 是地球上点的经度，$\psi$ 是地球上点的纬度) 与地球的一切子午线相交成定角.
\end{prob}

\begin{prob}[M]{3537}
求曲线
\[
z = f(x, y), \quad \frac{x - x_0}{\cos\alpha} = \frac{y - y_0}{\sin\alpha}
\]
(其中 $f$ 为可微函数) 在点 $M_0(x_0, y_0)$ 的切线与 $Oxy$ 平面所成角的正切.
\end{prob}

\begin{prob}[M]{3538}
求函数
\[
u = \frac{x}{\sqrt{x^2 + y^2 + z^2}}
\]
在点 $M(1, 2, -2)$ 沿曲线
\[
x = t, \quad y = 2t^2, \quad z = -2t^4
\]
在此点的切线方向的导数.
\end{prob}

\probgroup{写出下列曲面在点 $M_0$ 的切面和法线方程：}

\begin{prob}[S]{3539}
$z = x^2 + y^2$；$M_0(1, 2, 5)$.
\end{prob}

\begin{prob}[S]{3540}
$x^2 + y^2 + z^2 = 169$；$M_0(3, 4, 12)$.
\end{prob}

\begin{prob}[S]{3541}
$z = \arctan\dfrac{y}{x}$；$M_0\left(1, 1, \dfrac{\pi}{4}\right)$.
\end{prob}

\begin{prob}[S]{3542}
$ax^2 + by^2 + cz^2 = 1$；$M_0(x_0, y_0, z_0)$.
\end{prob}

\begin{prob}[S]{3543}
$z = y + \ln\dfrac{x}{z}$；$M_0(1, 1, 1)$.
\end{prob}

\begin{prob}[S]{3544}
$2^{\frac{x}{z}} + 2^{\frac{y}{z}} = 8$；$M_0(2, 2, 1)$.
\end{prob}

\begin{prob}[S]{3545}
$x = a\cos\psi\cos\varphi, y = b\cos\psi\sin\varphi, z = c\sin\psi$；$M_0(\varphi_0, \psi_0)$.
\end{prob}

\begin{prob}[S]{3546}
$x = r\cos\varphi, y = r\sin\varphi, z = r\cot\alpha$；$M_0(\varphi_0, r_0)$.
\end{prob}

\begin{prob}[S]{3547}
$x = u\cos v, y = u\sin v, z = av$；$M_0(u_0, v_0)$.
\end{prob}

\begin{prob}[M]{3548}
求曲面
\[
x = u + v, \quad y = u^2 + v^2, \quad z = u^3 + v^3
\]
的切平面当切点 $M(u, v) (u \neq v)$ 无限接近于曲面的边界线 $u = v$ 上的点 $M_0(u_0, v_0)$ 时的极限位置.
\end{prob}

\begin{prob}[M]{3549}
在曲面
\[
x^2 + 2y^2 + 3z^2 + 2xy + 2xz + 4yz = 8
\]
上求出切平面平行于坐标平面的诸切点.
\end{prob}

\begin{prob}[M]{3550}
在椭球面
\[
\frac{x^2}{a^2} + \frac{y^2}{b^2} + \frac{z^2}{c^2} = 1
\]
上怎样的点，椭球面的法线与各坐标轴成等角？
\end{prob}

\begin{prob}[M]{3551}
求曲面
\[
x^2 + 2y^2 + 3z^2 = 21
\]
的平行于平面 $x + 4y + 6z = 0$ 的各切平面.
\end{prob}

\begin{prob}[M]{3552}
证明：曲面 $xyz = a^3$ ($a > 0$) 的切平面与坐标面形成体积一定的四面体.
\end{prob}

\begin{prob}[M]{3553}
证明：曲面
\[
\sqrt{x} + \sqrt{y} + \sqrt{z} = \sqrt{a} \quad (a > 0)
\]
的切平面在坐标轴上割下的诸线段，其和为常量.
\end{prob}

\begin{prob}[M]{3554}
证明：锥面
\[
z = x f\left(\frac{y}{x}\right)
\]
的切平面经过其顶点.
\end{prob}

\begin{prob}[M]{3555}
证明：旋转面
\[
z = f(\sqrt{x^2 + y^2}) \quad (f' \neq 0)
\]
的法线与旋转轴相交.
\end{prob}

\begin{prob}[M]{3556}
求椭球面
\[
x^2 + y^2 + z^2 - xy = 1
\]
在坐标面上的投影.
\end{prob}

\begin{prob}[M]{3557}
分正方形 $\{0 \leqslant x \leqslant 1, 0 \leqslant y \leqslant 1\}$ 为直径不大于 $\delta$ 的有限个部分 $\sigma$. 若曲面
\[
z = 1 - x^2 - y^2
\]
在属于同一部分 $\sigma$ 的任何两点 $P(x, y)$ 及 $P_1(x_1, y_1)$ 的法线方向相差小于 $1^\circ$，求数 $\delta$ 的上界.
\end{prob}

\begin{prob}[L]{3558}
设
\[
z = f(x, y), \quad \text{其中 } (x, y) \in D \tag{1}
\]
为曲面的方程，$\varphi(P_1, P)$ 为曲面 (1) 在点 $P(x, y) \in D$ 及 $P_1(x_1, y_1) \in D$ 二点的法线之间的夹角. 证明：若 $D$ 为有界闭区域，函数 $f(x, y)$ 在区域 $D$ 内具有有界的二阶导数，则李雅普诺夫不等式
\[
\varphi(P_1, P) < C\rho(P_1, P) \tag{2}
\]
成立，其中 $C$ 为常数，$\rho(P_1, P)$ 为点 $P$ 与 $P_1$ 之间的距离.
\end{prob}

\begin{prob}[M]{3559}
圆柱面 $x^2 + y^2 = a^2$ 与曲面 $bz = xy$ 在公共点 $M_0(x_0, y_0, z_0)$ 相交成怎样的角？
\end{prob}

\begin{prob}[M]{3560}
证明：球坐标的坐标面 $x^2 + y^2 + z^2 = r^2, y = x\tan\varphi, x^2 + y^2 = z^2\tan^2\theta$ 两两正交.
\end{prob}

\begin{prob}[M]{3561}
证明：球面
\[
x^2 + y^2 + z^2 = 2ax, \quad x^2 + y^2 + z^2 = 2by, \quad x^2 + y^2 + z^2 = 2cz
\]
形成三重正交坐标系.
\end{prob}

\begin{prob}[M]{3562}
当 $\lambda = \lambda_1, \lambda = \lambda_2, \lambda = \lambda_3$ 时，经过每一点 $M(x, y, z)$ 有三个二次曲面：
\[
\frac{x^2}{a^2 - \lambda^2} + \frac{y^2}{b^2 - \lambda^2} + \frac{z^2}{c^2 - \lambda^2} = -1 \quad (a > b > c > 0).
\]
证明：这些曲面是正交的.
\end{prob}

\begin{prob}[L]{3563}
求函数 $u = x + y + z$ 沿球面 $x^2 + y^2 + z^2 = 1$ 在点 $M_0(x_0, y_0, z_0)$ 的外法线方向的导数. 在球面上怎样的点，函数 $u$ 的上述法向导数有：
\begin{enumerate}[label=(\alph*), itemsep=2pt, topsep=3pt]
\item 最大值，
\item 最小值，
\item 等于零？
\end{enumerate}
\end{prob}

\begin{prob}[M]{3564}
求函数 $u = x^2 + y^2 + z^2$ 沿椭球面 $\dfrac{x^2}{a^2} + \dfrac{y^2}{b^2} + \dfrac{z^2}{c^2} = 1$ 在点 $M_0(x_0, y_0, z_0)$ 的外法线方向的导数.
\end{prob}

\begin{prob}[M]{3565}
设 $\dfrac{\partial u}{\partial n}$ 和 $\dfrac{\partial v}{\partial n}$ 为函数 $u$ 和 $v$ 在曲面 $F(x, y, z) = 0$ 上的点的法向导数. 证
明：
\[
\frac{\partial}{\partial n}(uv) = u\frac{\partial v}{\partial n} + v\frac{\partial u}{\partial n}.
\]
\end{prob}

\probgroup{求含一个参变量的平面曲线族的包络线：}

\begin{prob}[S]{3566}
$x\cos\alpha + y\sin\alpha = p$ ($p = \text{常数}$).
\end{prob}

\begin{prob}[S]{3567}
$(x - a)^2 + y^2 = \dfrac{a^2}{2}$.
\end{prob}

\begin{prob}[S]{3568}
$y = kx + \dfrac{a}{k}$ ($a = \text{常数}$).
\end{prob}

\begin{prob}[S]{3569}
$y^2 = 2px + p^2$.
\end{prob}

\begin{prob}[M]{3570}
设有长为 $l$ 的线段，其两端点沿坐标轴滑动，求如此产生的线段族的包络线.
\end{prob}

\begin{prob}[M]{3571}
求面积 $S$ 为常数的椭圆族 $\dfrac{x^2}{a^2} + \dfrac{y^2}{b^2} = 1$ 的包络线.
\end{prob}

\begin{prob}[M]{3572}
炮弹在真空中以初速度 $v_0$ 射出，当投射角 $\alpha$ 在竖直平面中变化时，求炮弹轨迹的包络线.
\end{prob}

\begin{prob}[M]{3573}
证明：平面曲线的法线的包络线是此曲线的渐屈线.
\end{prob}

\begin{prob}[L]{3574}
研究下列曲线族的判别曲线的性质（$c$ 是参变量）：
\begin{enumerate}[label=(\alph*), itemsep=2pt, topsep=3pt]
\item 立方抛物线 $y = (x - c)^3$；
\item 半立方抛物线 $y^2 = (x - c)^3$；
\item 尼尔抛物线 $y^3 = (x - c)^2$；
\item 环索线 $(y - c)^2 = x^2\dfrac{a - x}{a + x}$.
\end{enumerate}
\end{prob}

\begin{prob}[M]{3575}
求半径为 $r$，中心在圆周 $x = R\cos t, y = R\sin t, z = 0$（$t$ 是参数，$R > r$）上的球族的包络面.
\end{prob}

\begin{prob}[M]{3576}
求球族
\[
(x - t\cos\alpha)^2 + (y - t\cos\beta)^2 + (z - t\cos\gamma)^2 = 1
\]
（其中 $\cos^2\alpha + \cos^2\beta + \cos^2\gamma = 1$，$t$ 是参变量）的包络面.
\end{prob}

\begin{prob}[M]{3577}
求相应体积 $V$ 是常数的椭球面族 $\dfrac{x^2}{a^2} + \dfrac{y^2}{b^2} + \dfrac{z^2}{c^2} = 1$ 的包络面.
\end{prob}

\begin{prob}[M]{3578}
求半径为 $\rho$，中心在圆锥面 $x^2 + y^2 = z^2$ 上的球族的包络面.
\end{prob}

\begin{prob}[M]{3579}
有一发光点位于坐标原点．若 $x_0^2 + y_0^2 + z_0^2 > R^2$，求由球
\[
(x - x_0)^2 + (y - y_0)^2 + (z - z_0)^2 \leqslant R^2
\]
投影所生成的阴影圆锥.
\end{prob}

\begin{prob}[M]{3580}
若参量 $p$ 和 $q$ 满足方程
\[
p^2 + q^2 = 1,
\]
求平面族
\[
z - z_0 = p(x - x_0) + q(y - y_0)
\]
的包络面.
\end{prob}

\probsec{§6. 泰勒公式}

\begin{prob}[M]{3581}
在点 $A(1, -2)$ 的邻域内根据泰勒公式展开函数
\[
f(x, y) = 2x^2 - xy - y^2 - 6x - 3y + 5.
\]
\end{prob}

\begin{prob}[M]{3582}
在点 $A(1, 1, 1)$ 的邻域内根据泰勒公式展开函数
\[
f(x, y, z) = x^3 + y^3 + z^3 - 3xyz.
\]
\end{prob}

\begin{prob}[M]{3583}
当自变量值从 $x = 1, y = -1$ 变到 $x_1 = 1 + h, y_1 = -1 + k$ 时，求函数 $f(x, y) = x^2 y + xy^2 - 2xy$ 的增量.
\end{prob}

\begin{prob}[M]{3584}
设
\[
f(x, y, z) = Ax^2 + By^2 + Cz^2 + 2Dxy + 2Exz + 2Fyz,
\]
按数 $h, k$ 和 $l$ 的正整数次幂展开 $f(x + h, y + k, z + l)$.
\end{prob}

\begin{prob}[M]{3585}
写出函数
\[
f(x, y) = x^y
\]
在点 $A(1, 1)$ 的邻域内的展开式，到二次项为止.
\end{prob}

\begin{prob}[M]{3586}
根据麦克劳林公式展开函数
\[
f(x, y) = \sqrt{1 - x^2 - y^2},
\]
到四次项为止.
\end{prob}

\begin{prob}[M]{3587}
若 $|x|$ 和 $|y|$ 同 $1$ 比较为很小的量，对于下列表达式
(a) $\dfrac{\cos x}{\cos y}$；\qquad (b) $\arctan\dfrac{1 + x + y}{1 - x + y}$
推出精确到二次项的近似公式.
\end{prob}

\begin{prob}[M]{3588}
假定 $x, y, z$ 的绝对值是很小的量，简化表达式
\[
\cos(x + y + z) - \cos x \cos y \cos z.
\]
\end{prob}

\begin{prob}[M]{3589}
按 $h$ 的幂次展开函数
\[
F(x, y) = \frac{1}{4}[f(x + h, y) + f(x, y + h) + f(x - h, y) + f(x, y - h)] - f(x, y),
\]
精确到 $h^4$.
\end{prob}

\begin{prob}[M]{3590}
已知中心在点 $P(x, y)$ 半径为 $\rho$ 的圆周，设 $f(P) = f(x, y)$ 及 $P_i(x_i, y_i)$（$i = 1, 2, 3$）为该圆周的内接正三角形的顶点，并且 $x_1 = x + \rho, y_1 = y$．按 $\rho$ 的正整数次幂展开函数
\[
F(\rho) = \frac{1}{3}[f(P_1) + f(P_2) + f(P_3)],
\]
精确到 $\rho^2$.
\end{prob}

\begin{prob}[M]{3591}
按 $h$ 与 $k$ 的幂次展开函数
\[
\Delta_{xy}f(x, y) = f(x + h, y + k) - f(x + h, y) - f(x, y + k) + f(x, y).
\]
\end{prob}

\begin{prob}[M]{3592}
依 $\rho$ 的幂次展开函数
\[
F(\rho) = \frac{1}{2\pi}\int_{0}^{2\pi} f(x + \rho\cos\varphi, y + \rho\sin\varphi)\,\mathrm{d}\varphi.
\]
\end{prob}

\probgroup{将下列函数展开成麦克劳林级数：}

\begin{prob}[S]{3593}
$f(x, y) = (1 + x)^m (1 + y)^n$.
\end{prob}

\begin{prob}[S]{3594}
$f(x, y) = \ln(1 + x + y)$.
\end{prob}

\begin{prob}[S]{3595}
$f(x, y) = \mathrm{e}^x \sin y$.
\end{prob}

\begin{prob}[S]{3596}
$f(x, y) = \mathrm{e}^x \cos y$.
\end{prob}

\begin{prob}[S]{3597}
$f(x, y) = \sin x \sinh y$.
\end{prob}

\begin{prob}[S]{3598}
$f(x, y) = \cos x \cosh y$.
\end{prob}

\begin{prob}[S]{3599}
$f(x, y) = \sin(x^2 + y^2)$.
\end{prob}

\begin{prob}[S]{3600}
$f(x, y) = \ln(1 + x)\ln(1 + y)$.
\end{prob}

\begin{prob}[M]{3601}
写出函数
\[
f(x, y) = \int_{0}^{1} (1 + x)^{t^2 y}\,\mathrm{d}t
\]
的麦克劳林级数的前三项.
\end{prob}

\begin{prob}[M]{3602}
按二项式 $x - 1$ 和 $y + 1$ 的正整数次幂将函数 $\mathrm{e}^{x + y}$ 展开成幂级数.
\end{prob}

\begin{prob}[M]{3603}
写出函数 $f(x, y) = \dfrac{x}{y}$ 在点 $M(1, 1)$ 的邻域内的泰勒级数展开式.
\end{prob}

\begin{prob}[M]{3604}
设 $z$ 为由方程 $z^3 - 2xz + y = 0$ 定义的 $x$ 和 $y$ 的隐函数，且当 $x = 1$ 和 $y = 1$ 时 $z = 1$．写出函数 $z$ 按二项式 $x - 1$ 和 $y - 1$ 的升幂排列的展开式中的若干项.
\end{prob}

\probgroup{研究下列曲线的奇点的种类并大略地画出这些曲线：}

\begin{prob}[S]{3605}
$y^2 = ax^2 + x^3$.
\end{prob}

\begin{prob}[S]{3606}
$x^3 + y^3 - 3xy = 0$.
\end{prob}

\begin{prob}[S]{3607}
$x^2 + y^2 = x^4 + y^4$.
\end{prob}

\begin{prob}[S]{3608}
$x^2 + y^4 = x^6$.
\end{prob}

\begin{prob}[S]{3609}
$(x^2 + y^2)^2 = a^2(x^2 - y^2)$.
\end{prob}

\begin{prob}[S]{3610}
$(y - x^2)^2 = x^5$.
\end{prob}

\begin{prob}[S]{3611}
$(a + x)y^2 = (a - x)x^2$.
\end{prob}

\begin{prob}[M]{3612}
研究参变量 $a, b, c$ ($a \leqslant b \leqslant c$) 的值与曲线 $y^2 = (x - a)(x - b)(x - c)$ 的形状之关系.
\end{prob}

\probgroup{研究超越曲线的奇点：}

\begin{prob}[S]{3613}
$y^2 = 1 - \mathrm{e}^{-x^2}$.
\end{prob}

\begin{prob}[S]{3614}
$y^2 = 1 - \mathrm{e}^{-x^3}$.
\end{prob}

\begin{prob}[S]{3615}
$y = x\ln x$.
\end{prob}

\begin{prob}[S]{3616}
$y = \dfrac{x}{1 + \mathrm{e}^{\frac{1}{x}}}$.
\end{prob}

\begin{prob}[S]{3617}
$y = \arctan\left(\dfrac{1}{\sin x}\right)$.
\end{prob}

\begin{prob}[S]{3618}
$y^2 = \sin\dfrac{\pi}{x}$.
\end{prob}

\begin{prob}[S]{3619}
$y^2 = \sin x^2$.
\end{prob}

\begin{prob}[S]{3620}
$y^2 = \sin^3 x$.
\end{prob}

\probsec{§7. 多元函数的极值}

\probgroup{研究下列多元函数的极值：}

\begin{prob}[S]{3621}
$z = x^2 + (y - 1)^2$.
\end{prob}

\begin{prob}[S]{3622}
$z = x^2 - (y - 1)^2$.
\end{prob}

\begin{prob}[S]{3623}
$z = (x - y + 1)^2$.
\end{prob}

\begin{prob}[S]{3624}
$z = x^2 - xy + y^2 - 2x + y$.
\end{prob}

\begin{prob}[S]{3625}
$z = x^2 y^3(6 - x - y)$.
\end{prob}

\begin{prob}[S]{3626}
$z = x^3 + y^3 - 3xy$.
\end{prob}

\begin{prob}[M]{3627}
(a) $z = x^4 + y^4 - x^2 - 2xy - y^2$；\qquad (b) $z = 2x^4 + y^4 - x^2 - 2y^2$.
\end{prob}

\begin{prob}[S]{3628}
$z = xy + \dfrac{50}{x} + \dfrac{20}{y}$ ($x > 0, y > 0$).
\end{prob}

\begin{prob}[S]{3629}
$z = xy\sqrt{1 - \dfrac{x^2}{a^2} - \dfrac{y^2}{b^2}}$ ($a > 0, b > 0$).
\end{prob}

\begin{prob}[S]{3630}
$z = \dfrac{ax + by + c}{\sqrt{x^2 + y^2 + 1}}$ ($a^2 + b^2 + c^2 \neq 0$).
\end{prob}

\begin{prob}[S]{3631}
$z = 1 - \sqrt{x^2 + y^2}$.
\end{prob}

\begin{prob}[S]{3632}
$z = \mathrm{e}^{2x + 3y}(8x^2 - 6xy + 3y^2)$.
\end{prob}

\begin{prob}[S]{3633}
$z = \mathrm{e}^{x^2 - y}(5 - 2x + y)$.
\end{prob}

\begin{prob}[S]{3634}
$z = (5x + 7y - 25)\mathrm{e}^{-(x^2 + xy + y^2)}$.
\end{prob}

\begin{prob}[S]{3635}
$z = x^2 + xy + y^2 - 4\ln x - 10\ln y$.
\end{prob}

\begin{prob}[S]{3636}
$z = \sin x + \cos y + \cos(x - y)$ $\left(0 \leqslant x \leqslant \dfrac{\pi}{2}; 0 \leqslant y \leqslant \dfrac{\pi}{2}\right)$.
\end{prob}

\begin{prob}[S]{3637}
$z = \sin x \sin y \sin(x + y)$ ($0 \leqslant x \leqslant \pi; 0 \leqslant y \leqslant \pi$).
\end{prob}

\begin{prob}[S]{3638}
$z = x - 2y + \ln\sqrt{x^2 + y^2} + 3\arctan\dfrac{y}{x}$.
\end{prob}

\begin{prob}[S]{3639}
$z = xy\ln(x^2 + y^2)$.
\end{prob}

\begin{prob}[S]{3640}
$z = x + y + 4\sin x \sin y$.
\end{prob}

\begin{prob}[S]{3641}
$z = (x^2 + y^2)\mathrm{e}^{-(x^2 + y^2)}$.
\end{prob}

\begin{prob}[S]{3642}
$u = x^2 + y^2 + z^2 + 2x + 4y - 6z$.
\end{prob}

\begin{prob}[S]{3643}
$u = x^3 + y^2 + z^2 + 12xy + 2z$.
\end{prob}

\begin{prob}[S]{3644}
$u = x + \dfrac{y^2}{4x} + \dfrac{x^2}{y} + \dfrac{2}{x}$ ($x > 0, y > 0, z > 0$).
\end{prob}

\begin{prob}[S]{3645}
$u = xy^2 z^3(a - x - 2y - 3z)$ ($a > 0$).
\end{prob}

\begin{prob}[S]{3646}
$u = \dfrac{a^2}{x} + \dfrac{x^2}{y} + \dfrac{y^2}{z} + \dfrac{z^2}{b}$ ($x > 0, y > 0, z > 0, a > 0, b > 0$).
\end{prob}

\begin{prob}[S]{3647}
$u = \sin x + \sin y + \sin z - \sin(x + y + z)$ ($0 \leqslant x \leqslant \pi; 0 \leqslant y \leqslant \pi; 0 \leqslant z \leqslant \pi$).
\end{prob}

\begin{prob}[S]{3648}
$u = x_1 x_2^2 \cdots x_n^n (1 - x_1 - 2x_2 - \cdots - nx_n)$ ($x_1 > 0, x_2 > 0, \dots, x_n > 0$).
\end{prob}

\begin{prob}[S]{3649}
$u = x_1 + \dfrac{x_2}{x_1} + \dfrac{x_3}{x_2} + \cdots + \dfrac{x_n}{x_{n-1}} + \dfrac{2}{x_n}$ ($x_i > 0, i = 1, 2, \dots, n$).
\end{prob}

\begin{prob}[M]{3650}
惠更斯问题．在 $a$ 和 $b$ 二正数间插入 $n$ 个数 $x_1, x_2, \dots, x_n$，使分数
\[
u = \frac{x_1 x_2 \cdots x_n}{(a + x_1)(x_1 + x_2)\cdots(x_n + b)}
\]
的值最大.
\end{prob}

\probgroup{求变量 $x$ 和 $y$ 的隐函数 $z$ 的极值：}

\begin{prob}[S]{3651}
$x^2 + y^2 + z^2 - 2x + 2y - 4z - 10 = 0$.
\end{prob}

\begin{prob}[S]{3652}
$x^2 + y^2 + z^2 - xz - yz + 2x + 2y + 2z - 2 = 0$.
\end{prob}

\begin{prob}[S]{3653}
$(x^2 + y^2 + z^2)^2 = a^2(x^2 + y^2 - z^2)$.
\end{prob}

\probgroup{求下列函数的条件极值点：}

\begin{prob}[S]{3654}
$z = xy$, 若 $x + y = 1$.
\end{prob}

\begin{prob}[S]{3655}
$z = \frac{x}{a} + \frac{y}{b}$, 若 $x^2 + y^2 = 1$.
\end{prob}

\begin{prob}[S]{3656}
$z = x^2 + y^2$, 若 $\frac{x}{a} + \frac{y}{b} = 1$.
\end{prob}

\begin{prob}[M]{3657}
(a) $z = Ax^2 + 2Bxy + Cy^2$, 若 $x^2 + y^2 = 1$; \qquad (b) $z = x^2 + 12xy + 2y^2$, 若 $4x^2 + y^2 = 25$.
\end{prob}

\begin{prob}[S]{3658}
$z = \cos^2 x + \cos^2 y$, 若 $x - y = \frac{\pi}{4}$.
\end{prob}

\begin{prob}[S]{3659}
$u = x - 2y + 2z$, 若 $x^2 + y^2 + z^2 = 1$.
\end{prob}

\begin{prob}[S]{3660}
$u = x^m y^n z^p$, 若 $x + y + z = a$ ($m > 0, n > 0, p > 0, a > 0$).
\end{prob}

\begin{prob}[S]{3661}
$u = x^2 + y^2 + z^2$, 若 $\frac{x^2}{a^2} + \frac{y^2}{b^2} + \frac{z^2}{c^2} = 1$ ($a > b > c > 0$).
\end{prob}

\begin{prob}[S]{3662}
$u = xy^2 z^3$, 若 $x + 2y + 3z = a$ ($x > 0, y > 0, z > 0, a > 0$).
\end{prob}

\begin{prob}[M]{3663}
\begin{enumerate}[label=(\alph*), itemsep=2pt, topsep=3pt]
\item $u = xyz$, 若 $x^2 + y^2 + z^2 = 1, x + y + z = 0$;
\item $u = xy + yz$, 若 $x^2 + y^2 = 2, y + z = 2$ ($x > 0, y > 0, z > 0$).
\end{enumerate}
\end{prob}

\begin{prob}[S]{3664}
$u = \sin x \sin y \sin z$, 若 $x + y + z = \frac{\pi}{2}$ ($x > 0, y > 0, z > 0$).
\end{prob}

\begin{prob}[M]{3665}
$u = \frac{x^2}{a^2} + \frac{y^2}{b^2} + \frac{z^2}{c^2}$, 若 $x^2 + y^2 + z^2 = 1, x\cos\alpha + y\cos\beta + z\cos\gamma = 0$ ($a > b > c > 0, \cos^2\alpha + \cos^2\beta + \cos^2\gamma = 1$).
\end{prob}

\begin{prob}[M]{3666}
$u = (x - \xi)^2 + (y - \eta)^2 + (z - \zeta)^2$, 若 $Ax + By + Cz = 0, x^2 + y^2 + z^2 = R^2, \frac{\xi}{\cos\alpha} = \frac{\eta}{\cos\beta} = \frac{\zeta}{\cos\gamma}$, 其中 $\cos^2\alpha + \cos^2\beta + \cos^2\gamma = 1$.
\end{prob}

\begin{prob}[S]{3667}
$u = x_1^2 + x_2^2 + \cdots + x_n^2$, 若 $\frac{x_1}{a_1} + \frac{x_2}{a_2} + \cdots + \frac{x_n}{a_n} = 1$ ($a_i > 0; i = 1, 2, \cdots, n$).
\end{prob}

\begin{prob}[S]{3668}
$u = x_1^p + x_2^p + \cdots + x_n^p$ ($p > 1$), 若 $x_1 + x_2 + \cdots + x_n = a$ ($a > 0$).
\end{prob}

\begin{prob}[M]{3669}
$u = \frac{\alpha_1}{x_1} + \frac{\alpha_2}{x_2} + \cdots + \frac{\alpha_n}{x_n}$, 若 $\beta_1 x_1 + \beta_2 x_2 + \cdots + \beta_n x_n = 1$ ($\alpha_i > 0, \beta_i > 0, x_i > 0; i = 1, 2, \cdots, n$).
\end{prob}

\begin{prob}[S]{3670}
$u = x_1^{\alpha_1} x_2^{\alpha_2} \cdots x_n^{\alpha_n}$, 若 $x_1 + x_2 + \cdots + x_n = a$ ($a > 0, \alpha_i > 1, i = 1, 2, \cdots, n$).
\end{prob}

\begin{prob}[M]{3671}
在 $\sum_{i=1}^{n} x_i^2 = 1$ 的条件下, 求二次型 $u = \sum_{i,j=1}^{n} a_{ij} x_i x_j$ ($a_{ij} = a_{ji}$) 的极值.
\end{prob}

\begin{prob}[M]{3672}
若 $n \geqslant 1$ 及 $x \geqslant 0, y \geqslant 0$, 证明不等式: $\frac{x^n + y^n}{2} \geqslant \left(\frac{x+y}{2}\right)^n$.
\probhint{在 $x + y = S$ 的条件下, 求函数 $z = \frac{1}{2}(x^n + y^n)$ 的极小值.}
\end{prob}

\begin{prob}[M]{3673}
证明赫尔德不等式:
\[
\sum_{i=1}^{n} a_i x_i \leqslant \left(\sum_{i=1}^{n} a_i^k\right)^{\frac{1}{k}} \left(\sum_{i=1}^{n} x_i^{k'}\right)^{\frac{1}{k'}}
\]
\[
\left(a_i \geqslant 0, x_i \geqslant 0, i = 1, 2, \cdots, n; k > 1, \frac{1}{k} + \frac{1}{k'} = 1\right).
\]
\probhint{在 $\sum_{i=1}^{n} a_i x_i = A$ 的条件下, 求函数
\[
u = \left(\sum_{i=1}^{n} a_i^k\right)^{\frac{1}{k}} \left(\sum_{i=1}^{n} x_i^{k'}\right)^{\frac{1}{k'}}
\]
的极小值.}
\end{prob}

\begin{prob}[M]{3674}
对于 $n$ 阶行列式 $A = |a_{ij}|$, 证明阿达马不等式:
\[
A^2 \leqslant \prod_{i=1}^{n} \left(\sum_{j=1}^{n} a_{ij}^2\right).
\]
\probhint{在关系式
\[
\sum_{j=1}^{n} a_{ij}^2 = S_i \quad (i = 1, 2, \cdots, n)
\]
存在的条件下, 研究行列式 $A = |a_{ij}|$ 的极值.}
\end{prob}

\probgroup{求下列函数在指定区域内的上确界 ($\sup$) 和下确界 ($\inf$)：}

\begin{prob}[S]{3675}
$z = x - 2y - 3$, 若 $0 \leqslant x \leqslant 1, 0 \leqslant y \leqslant 1, 0 \leqslant x + y \leqslant 1$.
\end{prob}

\begin{prob}[S]{3676}
$z = x^2 + y^2 - 12x + 16y$, 若 $x^2 + y^2 \leqslant 25$.
\end{prob}

\begin{prob}[S]{3677}
$z = x^2 - xy + y^2$, 若 $|x| + |y| \leqslant 1$.
\end{prob}

\begin{prob}[S]{3678}
$u = x^2 + 2y^2 + 3z^2$, 若 $x^2 + y^2 + z^2 \leqslant 100$.
\end{prob}

\begin{prob}[S]{3679}
$u = x + y + z$, 若 $x^2 + y^2 \leqslant z \leqslant 1$.
\end{prob}

\begin{prob}[M]{3680}
求函数
\[
u = (x + y + z)\mathrm{e}^{-(x + 2y + 3z)}
\]
在区域 $x > 0, y > 0, z > 0$ 内的下确界 ($\inf$) 与上确界 ($\sup$).
\end{prob}

\begin{prob}[M]{3681}
证明: 函数 $z = (1 + \mathrm{e}^y)\cos x - y\mathrm{e}^y$ 有无穷多个极大值而无一极小值.
\end{prob}

\begin{prob}[M]{3682}
函数 $f(x, y)$ 在点 $M_0(x_0, y_0)$ 有极小值的充分条件是否为此函数在通过点 $M_0$ 的每一条直线上有极小值呢? 研究例子 $f(x, y) = (x - y^2)(2x - y^2)$.
\end{prob}

\begin{prob}[S]{3683}
分解已知正数 $a$ 为 $n$ 个正的因数, 使得它们的倒数的和为最小.
\end{prob}

\begin{prob}[S]{3684}
分解已知正数 $a$ 为 $n$ 个相加数, 使得它们的平方和为最小.
\end{prob}

\begin{prob}[S]{3685}
分解已知正数 $a$ 为 $n$ 个正的因数, 使得它们的已知正数次幂的和为最小.
\end{prob}

\begin{prob}[M]{3686}
已知在平面上的 $n$ 个质点 $P_1(x_1, y_1), P_2(x_2, y_2), \cdots, P_n(x_n, y_n)$, 其质量分别为 $m_1, m_2, \cdots, m_n$. 点 $P(x, y)$ 位于何处时, 该质点系对此点的转动惯量为最小?
\end{prob}

\begin{prob}[S]{3687}
容积为 $V$ 的无盖长方浴盆在何种尺寸下有最小的表面积?
\end{prob}

\begin{prob}[M]{3688}
横截面为半圆形的无盖柱形浴盆, 其表面积等于 $S$, 在何种尺寸下此盆有最大的容积?
\end{prob}

\begin{prob}[M]{3689}
在球面 $x^2 + y^2 + z^2 = 1$ 上求一点, 这点到 $n$ 个已知点 $M_i(x_i, y_i, z_i)$ ($i = 1, 2, \cdots, n$) 距离的平方和为最小.
\end{prob}

\begin{prob}[M]{3690}
底面相同的直圆柱体与直圆锥体拼接在一起构成一个物体, 其总表面积 $Q$ 取给定值. 为了使此物体的体积为最大, 求其尺寸大小.
\end{prob}

\begin{prob}[M]{3691}
一长方体的上下两底均为正方形, 分别与同样的两个正四角锥体拼接在一起构成一个物体, 其体积 $V$ 取给定值. 当四角锥的侧面对它们的底成怎样的倾角时, 该物体的总表面积为最小?
\end{prob}

\begin{prob}[M]{3692}
将周长为 $2p$ 的矩形绕其一边旋转, 矩形所扫过的区域构成一旋转体, 求使该旋转体体积为最大的那个矩形.
\end{prob}

\begin{prob}[M]{3693}
将周长为 $2p$ 的三角形绕其一边旋转, 三角形所扫过的区域构成一旋转体, 求使该旋转体体积为最大的那个三角形.
\end{prob}

\begin{prob}[S]{3694}
在半径为 $R$ 的半球内作出具有最大体积的内接长方体.
\end{prob}

\begin{prob}[S]{3695}
在已知的直圆锥内作出具有最大体积的内接长方体.
\end{prob}

\begin{prob}[M]{3696}
在椭球
\[
\frac{x^2}{a^2} + \frac{y^2}{b^2} + \frac{z^2}{c^2} = 1
\]
内作出具有最大体积的内接长方体.
\end{prob}

\begin{prob}[M]{3697}
直圆锥的母线 $l$ 与底面成倾角 $\alpha$, 试在此直圆锥中作出具有最大全表面积的内接长方体.
\end{prob}

\begin{prob}[M]{3698}
在椭圆抛物面 $\frac{z}{c} = \frac{x^2}{a^2} + \frac{y^2}{b^2}$ 和 $z = c$ 所围区域内作出具有最大体积的内接长方体.
\end{prob}

\begin{prob}[S]{3699}
求点 $M_0(x_0, y_0, z_0)$ 至平面 $Ax + By + Cz + D = 0$ 上的点的最短距离.
\end{prob}

\begin{prob}[M]{3700}
求空间二直线
\[
\frac{x - x_1}{m_1} = \frac{y - y_1}{n_1} = \frac{z - z_1}{p_1}, \quad \frac{x - x_2}{m_2} = \frac{y - y_2}{n_2} = \frac{z - z_2}{p_2}
\]
之间的最短距离.
\end{prob}

\begin{prob}[S]{3701}
求抛物线 $y = x^2$ 和直线 $x - y - 2 = 0$ 之间的最短距离.
\end{prob}

\begin{prob}[M]{3702}
求有心二次曲线
\[
Ax^2 + 2Bxy + Cy^2 = 1
\]
的半轴.
\end{prob}

\begin{prob}[M]{3703}
求有心二次曲面
\[
Ax^2 + By^2 + Cz^2 + 2Dxy + 2Eyz + 2Fxz = 1
\]
的半轴.
\end{prob}

\begin{prob}[M]{3704}
求用平面
\[
Ax + By + Cz = 0
\]
与柱体
\[
\frac{x^2}{a^2} + \frac{y^2}{b^2} \leqslant 1
\]
相交所成椭圆的面积.
\end{prob}

\begin{prob}[M]{3705}
求用平面
\[
x\cos\alpha + y\cos\beta + z\cos\gamma = 0
\]
(其中 $\cos^2\alpha + \cos^2\beta + \cos^2\gamma = 1$) 与椭球
\[
\frac{x^2}{a^2} + \frac{y^2}{b^2} + \frac{z^2}{c^2} \leqslant 1
\]
相交所成截面的面积.
\end{prob}

\begin{prob}[M]{3706}
根据费马原理, 光在最短时间内从一点传播至另一点.

假定这两点位于交界面为平面的不同的光介质中, 并且光的传播速度在第一种介质中等于 $v_1$, 而在第二种介质中等于 $v_2$, 试推出光的折射定律.
\end{prob}

\begin{prob}[M]{3707}
一折射棱镜的折射角\footnote{折射棱镜的折射角指两折射面之间的二面角. —— 译注}为 $\alpha$, 折射率为 $n$. 光线以怎样的入射角射向此棱镜侧面, 其偏向角 (即入射光线与出射光线之间的角) 为最小? 求此最小偏向角.
\end{prob}

\begin{prob}[M]{3708}
变量 $x$ 和 $y$ 满足系数待定的线性方程
\[
y = ax + b.
\]
经过一系列精度相同的测量, 对于量 $x$ 和 $y$ 得到值 $x_i, y_i$ ($i = 1, 2, \cdots, n$). 利用最小二乘法, 求系数 $a$ 和 $b$ 的最可靠数值.
\probhint{根据最小二乘法, 系数 $a$ 和 $b$ 的最可靠数值所对应的误差的平方和
\[
\sum_{i=1}^{n} \Delta_i^2 = \sum_{i=1}^{n} (ax_i + b - y_i)^2
\]
为最小.}
\end{prob}

\begin{prob}[M]{3709}
在平面上已知 $n$ 个点 $M_i(x_i, y_i)$ ($i = 1, 2, \cdots, n$). 直线
\[
x\cos\alpha + y\sin\alpha - p = 0
\]
在怎样的位置时, 这些点与此直线的偏差的平方和为最小?
\end{prob}

\begin{prob}[M]{3710}
在区间 $(1, 3)$ 内用线性函数 $ax + b$ 来近似地代替函数 $x^2$, 使得绝对偏差
\[
\Delta = \sup |x^2 - (ax + b)| \quad (1 \leqslant x \leqslant 3)
\]
为最小.
\end{prob}

\probchapter{第七章\quad 带参数的积分}

\probsec{§1. 带参数的常义积分}

\begin{prob}[M]{3711}
证明: 不连续函数 $f(x, y) = \operatorname{sgn}(x - y)$ 的积分
\[
F(y) = \int_0^1 f(x, y) \,\mathrm{d}x
\]
为连续函数. 作出函数 $u = F(y)$ 的图像.
\end{prob}

\begin{prob}[M]{3712}
研究函数
\[
F(y) = \int_0^1 \frac{y f(x)}{x^2 + y^2} \,\mathrm{d}x
\]
的连续性, 其中 $f(x)$ 在闭区间 $[0, 1]$ 上是正的连续函数.
\end{prob}

\begin{prob}[L]{3713}
求:
\begin{enumerate}[label=(\alph*), itemsep=2pt, topsep=3pt]
\item $\displaystyle \lim_{\alpha \to 0} \int_\alpha^{1+\alpha} \frac{\mathrm{d}x}{1 + x^2 + \alpha^2}$;
\item $\displaystyle \lim_{\alpha \to 0} \int_{-1}^1 \sqrt{x^2 + \alpha^2} \,\mathrm{d}x$;
\item $\displaystyle \lim_{\alpha \to 0} \int_0^2 x^2 \cos \alpha x \,\mathrm{d}x$;
\item $\displaystyle \lim_{n \to \infty} \int_0^1 \frac{\mathrm{d}x}{1 + \left(1 + \frac{x}{n}\right)^n}$;
\item $\displaystyle \lim_{R \to \infty} \int_0^{\frac{\pi}{2}} \mathrm{e}^{-R \sin \theta} \,\mathrm{d}\theta$.
\end{enumerate}
\end{prob}

\begin{prob}[M]{3714.1}
设函数 $f(x)$ 在闭区间 $[A, B]$ 上连续. 证明:
\[
\lim_{h \to 0} \frac{1}{h} \int_a^x [f(t + h) - f(t)] \,\mathrm{d}t = f(x) - f(a) \quad (A < a < x < B).
\]
\end{prob}

\begin{prob}[M]{3714.2}
设: (1) 在 $[-1, 1]$ 上 $\varphi_n(x) \geqslant 0$ ($n = 1, 2, \cdots$); (2) 对于 $0 < \varepsilon \leqslant |x| \leqslant 1$, 当 $n \to \infty$ 时 $\varphi_n(x) \rightrightarrows 0$; (3) 当 $n \to \infty$ 时 $\displaystyle \int_{-1}^1 \varphi_n(x) \,\mathrm{d}x \to 1$. 证明: 若 $f(x) \in C[-1, 1]$, 则
\[
\lim_{n \to \infty} \int_{-1}^1 f(x) \varphi_n(x) \,\mathrm{d}x = f(0).
\]
\end{prob}

\begin{prob}[M]{3715}
在表达式
\[
\lim_{y \to 0} \int_0^1 \frac{x}{y^2} \mathrm{e}^{-\frac{x^2}{y^2}} \,\mathrm{d}x
\]
中, 可否在积分符号下完成极限运算?
\end{prob}

\begin{prob}[M]{3716}
当 $y = 0$ 时可否根据莱布尼茨法则计算函数
\[
F(y) = \int_0^1 \ln \sqrt{x^2 + y^2} \,\mathrm{d}x
\]
的导数?
\end{prob}

\begin{prob}[M]{3717}
若
\[
F(x) = \int_x^{x^2} \mathrm{e}^{-x y^2} \,\mathrm{d}y,
\]
计算 $F'(x)$.
\end{prob}

\begin{prob}[L]{3718}
设:
\begin{enumerate}[label=(\alph*), itemsep=2pt, topsep=3pt]
\item $\displaystyle F(\alpha) = \int_{\sin \alpha}^{\cos \alpha} \mathrm{e}^{\alpha \sqrt{1 - x^2}} \,\mathrm{d}x$;
\item $\displaystyle F(\alpha) = \int_{a+\alpha}^{b+\alpha} \frac{\sin \alpha x}{x} \,\mathrm{d}x$;
\item $\displaystyle F(\alpha) = \int_0^\alpha \frac{\ln(1 + \alpha x)}{x} \,\mathrm{d}x$;
\item $\displaystyle F(\alpha) = \int_0^\alpha f(x + \alpha, x - \alpha) \,\mathrm{d}x$;
\item $\displaystyle F(\alpha) = \int_0^{\alpha^2} \mathrm{d}x \int_{x-\alpha}^{x+\alpha} \sin(x^2 + y^2 - \alpha^2) \,\mathrm{d}y$.
\end{enumerate}
求 $F'(\alpha)$.
\end{prob}

\begin{prob}[M]{3719}
若
\[
F(x) = \int_0^x (x + y) f(y) \,\mathrm{d}y,
\]
其中 $f(x)$ 为可微函数, 求 $F''(x)$.
\end{prob}

\begin{prob}[M]{3720}
设
\[
F(x) = \int_a^b f(y) |x - y| \,\mathrm{d}y,
\]
其中 $a < b$, $f(y)$ 为可微函数, 求 $F''(x)$.
\end{prob}

\begin{prob}[M]{3721.1}
设
\[
F(x) = \frac{1}{h^2} \int_0^h \mathrm{d}\xi \int_0^h f(x + \xi + \eta) \,\mathrm{d}\eta \quad (h > 0),
\]
其中 $f(x)$ 为连续函数, 求 $F''(x)$.
\end{prob}

\begin{prob}[M]{3721.2}
设
\[
F(x) = \int_0^x f(t)(x - t)^{n-1} \,\mathrm{d}t,
\]
求 $F^{(n)}(x)$.
\end{prob}

\begin{prob}[L]{3722}
证明公式:
\[
\frac{\mathrm{d}^n}{\mathrm{d}x^n} \left( \frac{\sin x}{x} \right) = \frac{1}{x^{n+1}} \int_0^x y^n \cos \left( y + \frac{n\pi}{2} \right) \,\mathrm{d}y \quad (n = 1, 2, \cdots). \tag{1}
\]
利用公式 (1) 获得以下估计:
\[
\left| \frac{\mathrm{d}^n}{\mathrm{d}x^n} \left( \frac{\sin x}{x} \right) \right| \leqslant \frac{1}{n + 1}, \quad x \in (-\infty, +\infty).
\]
\end{prob}

\begin{prob}[M]{3723}
在区间 $1 \leqslant x \leqslant 3$ 上用线性函数 $a + bx$ 近似地代替函数 $f(x) = x^2$, 使得
\[
\int_1^3 (a + bx - x^2)^2 \,\mathrm{d}x = \min.
\]
\end{prob}

\begin{prob}[M]{3724}
依条件: 函数 $a + bx$ 及 $\sqrt{1 + x^2}$ 在已知区间 $[0, 1]$ 上的均方偏差为最小, 求近似公式
\[
\sqrt{1 + x^2} \approx a + bx \quad (0 \leqslant x \leqslant 1).
\]
\end{prob}

\begin{prob}[L]{3725}
求完全椭圆积分
\[
E(k) = \int_0^{\frac{\pi}{2}} \sqrt{1 - k^2 \sin^2 \varphi} \,\mathrm{d}\varphi, \quad F(k) = \int_0^{\frac{\pi}{2}} \frac{\mathrm{d}\varphi}{\sqrt{1 - k^2 \sin^2 \varphi}} \quad (0 < k < 1)
\]
的导数, 并把它们用函数 $E(k)$ 和 $F(k)$ 表示出来.

证明: $E(k)$ 满足微分方程
\[
E''(k) + \frac{1}{k} E'(k) + \frac{E(k)}{1 - k^2} = 0.
\]
\end{prob}

\begin{prob}[M]{3726}
证明: 阶数 $n$ 为整数的贝塞尔函数
\[
J_n(x) = \frac{1}{\pi} \int_0^\pi \cos(n\varphi - x \sin \varphi) \,\mathrm{d}\varphi
\]
满足贝塞尔方程
\[
x^2 J_n''(x) + x J_n'(x) + (x^2 - n^2) J_n(x) = 0.
\]
\end{prob}

\begin{prob}[M]{3727}
设
\[
I(\alpha) = \int_0^\alpha \frac{\varphi(x) \,\mathrm{d}x}{\sqrt{\alpha - x}},
\]
其中函数 $\varphi(x)$ 及其导数 $\varphi'(x)$ 在闭区间 $0 \leqslant x \leqslant a$ 上连续. 证明: 当 $0 < \alpha < a$ 时有
\[
I'(\alpha) = \frac{\varphi(0)}{\sqrt{\alpha}} + \int_0^\alpha \frac{\varphi'(x)}{\sqrt{\alpha - x}} \,\mathrm{d}x.
\]
\probhint{令 $x = \alpha t$.}
\end{prob}

\begin{prob}[M]{3728}
设
\[
u(x) = \int_0^1 K(x, y) v(y) \,\mathrm{d}y,
\]
其中
\[
K(x, y) = \begin{cases} x(1 - y), & x \leqslant y, \\ y(1 - x), & x > y, \end{cases}
\]
及 $v(y)$ 都是连续的, 证明: 函数 $u(x)$ 满足方程
\[
u''(x) = -v(x) \quad (0 \leqslant x \leqslant 1).
\]
\end{prob}

\begin{prob}[M]{3729}
设
\[
F(x, y) = \int_{\frac{x}{y}}^{xy} (x - yz) f(z) \,\mathrm{d}z,
\]
其中 $f(z)$ 为可微函数, 求 $F''_{xy}(x, y)$.
\end{prob}

\begin{prob}[M]{3730}
设 $f(x)$ 为二阶可微函数, $F(x)$ 为可微函数. 证明: 函数
\[
u(x, t) = \frac{1}{2} [f(x - at) + f(x + at)] + \frac{1}{2a} \int_{x-at}^{x+at} F(z) \,\mathrm{d}z
\]
满足弦的振动方程
\[
\frac{\partial^2 u}{\partial t^2} = a^2 \frac{\partial^2 u}{\partial x^2}
\]
及初始条件: $u(x, 0) = f(x)$, $u'_t(x, 0) = F(x)$.
\end{prob}

\begin{prob}[M]{3731}
证明: 若函数 $f(x)$ 在闭区间 $[0, l]$ 上连续, 且当 $0 \leqslant \xi \leqslant l$ 时 $(x - \xi)^2 + y^2 + z^2 \neq 0$, 则函数
\[
u(x, y, z) = \int_0^l \frac{f(\xi) \,\mathrm{d}\xi}{\sqrt{(x - \xi)^2 + y^2 + z^2}}
\]
满足拉普拉斯方程
\[
\frac{\partial^2 u}{\partial x^2} + \frac{\partial^2 u}{\partial y^2} + \frac{\partial^2 u}{\partial z^2} = 0.
\]
\end{prob}

\probgroup{应用对参数的微分法, 计算下列积分:}

\begin{prob}[M]{3732}
$\displaystyle \int_0^{\frac{\pi}{2}} \ln(a^2 \sin^2 x + b^2 \cos^2 x) \,\mathrm{d}x$.
\end{prob}

\begin{prob}[M]{3733}
$\displaystyle \int_0^\pi \ln(1 - 2a \cos x + a^2) \,\mathrm{d}x$.
\end{prob}

\begin{prob}[M]{3734}
$\displaystyle \int_0^{\frac{\pi}{2}} \frac{\arctan(a \tan x)}{\tan x} \,\mathrm{d}x$.
\end{prob}

\begin{prob}[M]{3735}
$\displaystyle \int_0^{\frac{\pi}{2}} \ln \frac{1 + a \cos x}{1 - a \cos x} \cdot \frac{\mathrm{d}x}{\cos x} \quad (|a| < 1)$.
\end{prob}

\begin{prob}[M]{3736}
利用公式
\[
\frac{\arctan x}{x} = \int_0^1 \frac{\mathrm{d}y}{1 + x^2 y^2},
\]
计算积分
\[
\int_0^1 \frac{\arctan x}{x} \frac{\mathrm{d}x}{\sqrt{1 - x^2}}.
\]
\end{prob}

\begin{prob}[M]{3737}
应用积分号下的积分法, 计算积分
\[
\int_0^1 \frac{x^b - x^a}{\ln x} \,\mathrm{d}x \quad (a > 0, b > 0).
\]
\end{prob}

\begin{prob}[M]{3738}
计算积分:
\begin{enumerate}[label=(\alph*), itemsep=2pt, topsep=3pt]
\item $\displaystyle \int_0^1 \sin\left(\ln \frac{1}{x}\right) \frac{x^b - x^a}{\ln x} \,\mathrm{d}x$;
\item $\displaystyle \int_0^1 \cos\left(\ln \frac{1}{x}\right) \frac{x^b - x^a}{\ln x} \,\mathrm{d}x \quad (a > 0, b > 0)$.
\end{enumerate}
\end{prob}

\begin{prob}[M]{3739}
设 $F(k)$ 和 $E(k)$ 为完全椭圆积分 (参阅习题 3725). 证明公式:
\begin{enumerate}[label=(\alph*), itemsep=2pt, topsep=3pt]
\item $\displaystyle \int_0^k F(k) k \,\mathrm{d}k = E(k) - k_1^2 F(k)$;
\item $\displaystyle \int_0^k E(k) k \,\mathrm{d}k = \frac{1}{3} [(1 + k^2) E(k) - k_1^2 F(k)]$,
\end{enumerate}
其中 $k_1^2 = 1 - k^2$.
\end{prob}

\begin{prob}[M]{3740}
证明公式:
\[
\int_0^x x J_0(x) \,\mathrm{d}x = x J_1(x),
\]
其中 $J_0(x)$ 及 $J_1(x)$ 为阶数是 $0$ 与 $1$ 的贝塞尔函数 (参阅习题 3726).
\end{prob}

\probsec{§2. 带参数的广义积分. 积分的一致收敛性}

\probgroup{求积分的收敛域:}

\begin{prob}[S]{3741}
$\int_{0}^{+\infty} \frac{\mathrm{e}^{-ax}}{1+x^2}\,\mathrm{d}x$.
\end{prob}

\begin{prob}[S]{3742}
$\int_{\pi}^{+\infty} \frac{x\cos x}{x^p + x^q}\,\mathrm{d}x$.
\end{prob}

\begin{prob}[S]{3743}
$\int_{0}^{+\infty} \frac{\sin x^q}{x^p}\,\mathrm{d}x$.
\end{prob}

\begin{prob}[S]{3744}
$\int_{0}^{2} \frac{\mathrm{d}x}{|\ln x|^p}$.
\end{prob}

\begin{prob}[S]{3745}
$\int_{0}^{1} \frac{\cos \frac{1}{1-x}}{\sqrt[n]{1-x^2}}\,\mathrm{d}x$.
\end{prob}

\begin{prob}[S]{3746}
$\int_{0}^{+\infty} \frac{\sin x}{x^p + \sin x}\,\mathrm{d}x \ (p > 0)$.
\end{prob}

\probgroup{利用与级数比较的方法研究下列积分的收敛性：}

\begin{prob}[S]{3747}
$\int_{0}^{+\infty} \frac{\cos x}{x+a}\,\mathrm{d}x$.
\end{prob}

\begin{prob}[S]{3748}
$\int_{0}^{+\infty} \frac{x\,\mathrm{d}x}{1 + x^n \sin^2 x} \ (n > 0)$.
\end{prob}

\begin{prob}[S]{3749}
$\int_{\pi}^{+\infty} \frac{\mathrm{d}x}{x^p \sqrt[3]{\sin^2 x}}$.
\end{prob}

\begin{prob}[S]{3750}
$\int_{0}^{+\infty} \frac{\sin(x+x^2)}{x^n}\,\mathrm{d}x$.
\end{prob}

\begin{prob}[M]{3751}
用肯定的方式陈述，什么是积分
\[
\int_{a}^{+\infty} f(x, y)\,\mathrm{d}x
\]
在已知区间 $(y_1, y_2)$ 内不一致收敛？
\end{prob}

\begin{prob}[M]{3752}
证明：若 1) 积分
\[
\int_{a}^{+\infty} f(x, y)\,\mathrm{d}x
\]
收敛，2) 函数 $\varphi(x, y)$ 有界，且关于 $x$ 是单调的，则积分
\[
\int_{a}^{+\infty} f(x)\varphi(x, y)\,\mathrm{d}x
\]
一致收敛（在相应区域内）。
\end{prob}

\begin{prob}[M]{3753}
证明：一致收敛的积分
\[
I = \int_{1}^{+\infty} \mathrm{e}^{-\frac{1}{y^2}\left(x-\frac{1}{y}\right)^2}\,\mathrm{d}x \quad (0 < y < 1)
\]
不能以与参数无关的收敛积分为强函数。
\end{prob}

\begin{prob}[M]{3754}
证明：积分
\[
I = \int_{0}^{+\infty} \alpha \mathrm{e}^{-\alpha x}\,\mathrm{d}x
\]
(a) 在任何区间 $0 < a \leqslant \alpha \leqslant b$ 内一致收敛；(b) 在区间 $0 \leqslant \alpha \leqslant b$ 内非一致收敛。
\end{prob}

\begin{prob}[M]{3755.1}
证明：狄利克雷积分
\[
I = \int_{0}^{+\infty} \frac{\sin \alpha x}{x}\,\mathrm{d}x
\]
(a) 在每一个不含数值 $\alpha = 0$ 的闭区间 $[a, b]$ 上一致收敛，(b) 在含数值 $\alpha = 0$ 的每一个闭区间 $[a, b]$ 上非一致收敛。
\end{prob}

\begin{prob}[M]{3755.2}
研究积分 $\int_{1}^{+\infty} \frac{\mathrm{d}x}{x^\alpha}$ 在以下区间上的一致收敛性：
\begin{enumerate}[label=(\alph*), itemsep=2pt, topsep=3pt]
  \item $1 < \alpha_0 \leqslant \alpha < +\infty$；
  \item $1 < \alpha < +\infty$。
\end{enumerate}
\end{prob}

\begin{prob}[M]{3755.3}
证明：积分
\[
\int_{0}^{+\infty} \frac{\mathrm{d}x}{x^\alpha + 1}
\]
在区间 $1 < \alpha < +\infty$ 上一致收敛。
\end{prob}

\probgroup{研究下列积分在所指定区间内的一致收敛性：}

\begin{prob}[M]{3756}
(a) $\int_{0}^{1} \frac{\mathrm{d}x}{x^\alpha} \ (0 < \alpha < 1)$；\qquad (b) $\int_{0}^{+\infty} \mathrm{e}^{-\alpha x}\sin x\,\mathrm{d}x \ (0 < \alpha_0 \leqslant \alpha < +\infty)$。
\end{prob}

\begin{prob}[S]{3757}
$\int_{1}^{+\infty} x^\alpha \mathrm{e}^{-x}\,\mathrm{d}x \ (a \leqslant \alpha \leqslant b)$。
\end{prob}

\begin{prob}[S]{3758}
$\int_{-\infty}^{+\infty} \frac{\cos \alpha x}{1+x^2}\,\mathrm{d}x \ (-\infty < \alpha < +\infty)$。
\end{prob}

\begin{prob}[S]{3759}
$\int_{0}^{+\infty} \frac{\mathrm{d}x}{(x+\alpha)^2 + 1} \ (0 \leqslant \alpha < +\infty)$。
\end{prob}

\begin{prob}[M]{3760}
(a) $\int_{0}^{+\infty} \frac{\sin x}{x} \mathrm{e}^{-\alpha x}\,\mathrm{d}x \ (0 \leqslant \alpha < +\infty)$；\qquad (b) $\int_{1}^{+\infty} \frac{\ln^p x}{x\sqrt{x}}\,\mathrm{d}x \ (0 \leqslant p \leqslant 10)$。
\end{prob}

\begin{prob}[S]{3761}
$\int_{1}^{+\infty} \mathrm{e}^{-\alpha x}\frac{\cos x}{x^p}\,\mathrm{d}x \ (0 \leqslant \alpha < +\infty)$，其中 $p > 0$ 是常数。
\end{prob}

\begin{prob}[S]{3762}
$\int_{0}^{+\infty} \sqrt{\alpha}\mathrm{e}^{-\alpha x^2}\,\mathrm{d}x \ (0 \leqslant \alpha < +\infty)$。
\end{prob}

\begin{prob}[M]{3763}
$\int_{-\infty}^{+\infty} \mathrm{e}^{-(x-\alpha)^2}\,\mathrm{d}x$；(a) $a < \alpha < b$；(b) $-\infty < \alpha < +\infty$。
\end{prob}

\begin{prob}[M]{3764}
(a) $\int_{0}^{+\infty} \mathrm{e}^{-x^2(1+y^2)} \sin x\,\mathrm{d}y \ (-\infty < x < +\infty)$；\qquad (b) $\int_{0}^{+\infty} \frac{\sin x^2}{1+x^p}\,\mathrm{d}x \ (p \geqslant 0)$。
\end{prob}

\begin{prob}[M]{3765}
选取数 $b > 0$，使得在 $1.1 \leqslant p \leqslant 10$ 时
\[
\int_{b}^{+\infty} \frac{\mathrm{d}x}{1+x^p} < \varepsilon,
\]
式中 $\varepsilon = 10^{-6}$。
\end{prob}

\begin{prob}[M]{3766}
$\int_{0}^{1} x^{p-1} \ln^q \frac{1}{x}\,\mathrm{d}x$；(a) $p \geqslant p_0 > 0$；(b) $p > 0 \ (q > -1)$。
\end{prob}

\begin{prob}[S]{3767}
$\int_{0}^{1} \frac{x^n}{\sqrt{1-x^2}}\,\mathrm{d}x \ (0 \leqslant n < +\infty)$。
\end{prob}

\begin{prob}[S]{3768}
$\int_{0}^{1} \sin \frac{1}{x}\frac{\mathrm{d}x}{x^n} \ (0 < n < 2)$。
\end{prob}

\begin{prob}[S]{3769}
$\int_{0}^{2} \frac{x^\alpha\,\mathrm{d}x}{\sqrt[3]{(x-1)(x-2)^2}} \ \left(|\alpha| < \frac{1}{2}\right)$。
\end{prob}

\begin{prob}[S]{3770}
$\int_{0}^{1} \frac{\sin \alpha x}{\sqrt{|x-\alpha|}}\,\mathrm{d}x \ (0 \leqslant \alpha \leqslant 1)$。
\end{prob}

\begin{prob}[M]{3771}
若积分在参数的已知值的某邻域内一致收敛，则称此积分对参数的已知值一致收敛。

证明：积分
\[
I = \int_{0}^{+\infty} \frac{\alpha\,\mathrm{d}x}{1+\alpha^2 x^2}
\]
对每一个 $\alpha \neq 0$ 的值一致收敛，而对 $\alpha = 0$ 非一致收敛。
\end{prob}

\begin{prob}[M]{3772}
在下式中
\[
\lim_{\alpha \to +0} \int_{0}^{+\infty} \alpha \mathrm{e}^{-\alpha x}\,\mathrm{d}x
\]
把极限移到积分符号内合理吗？
\end{prob}

\begin{prob}[M]{3773}
函数 $f(x)$ 在区间 $(0, +\infty)$ 内可积，证明公式：
\[
\lim_{\alpha \to +0} \int_{0}^{+\infty} \mathrm{e}^{-\alpha x} f(x)\,\mathrm{d}x = \int_{0}^{+\infty} f(x)\,\mathrm{d}x.
\]
\end{prob}

\begin{prob}[M]{3774.1}
证明：若 $f'(x)$ 在 $[a, +\infty)$ 上绝对可积，则存在 $\lim_{x \to \infty} f(x)$。
\end{prob}

\begin{prob}[M]{3774.2}
证明：若 $f(x)$ 在区间 $(0, +\infty)$ 内绝对可积，则
\[
\lim_{n \to \infty} \int_{0}^{+\infty} f(x) \sin nx\,\mathrm{d}x = 0.
\]
\end{prob}

\begin{prob}[M]{3775}
证明：若 (1) 在每一个有限区间 $(a, b)$ 内 $f(x, y) \rightrightarrows f(x, y_0)$；(2) $|f(x, y)| \leqslant F(x)$，其中 $\int_{a}^{+\infty} F(x)\,\mathrm{d}x < +\infty$，则
\[
\lim_{y \to y_0} \int_{a}^{+\infty} f(x, y)\,\mathrm{d}x = \int_{a}^{+\infty} \lim_{y \to y_0} f(x, y)\,\mathrm{d}x.
\]
\end{prob}

\begin{prob}[M]{3776.1}
利用积分号与极限号互换，计算积分
\[
\int_{0}^{+\infty} \mathrm{e}^{-x^2}\,\mathrm{d}x = \int_{0}^{+\infty} \lim_{n \to \infty} \left[\left(1 + \frac{x^2}{n}\right)^{-n}\right]\,\mathrm{d}x.
\]
\end{prob}

\begin{prob}[M]{3776.2}
设 $f(x)$ 在 $[0, +\infty)$ 上连续且有界。证明：
\[
\lim_{y \to 0} \frac{2}{\pi} \int_{0}^{+\infty} \frac{y f(x)}{x^2 + y^2}\,\mathrm{d}x = f(0).
\]
\end{prob}

\begin{prob}[M]{3777.1}
求
\[
\lim_{n \to \infty} \int_{0}^{+\infty} \frac{\mathrm{d}x}{x^n + 1}.
\]
\end{prob}

\begin{prob}[M]{3777.2}
证明：积分
\[
F(a) = \int_{0}^{+\infty} \mathrm{e}^{-(x-a)^2}\,\mathrm{d}x
\]
是参数 $a$ 的连续函数。
\end{prob}

\begin{prob}[M]{3778.1}
证明：
\[
F(\alpha) = \int_{0}^{1} \frac{\sin \frac{\alpha}{x}}{x^\alpha}\,\mathrm{d}x
\]
在区间 $0 < \alpha < 1$ 上是连续函数。
\end{prob}

\begin{prob}[M]{3778.2}
求函数
\[
F(a) = \int_{0}^{+\infty} \frac{\sin(1-a^2)x}{x}\,\mathrm{d}x
\]
的间断点。作出函数 $y = F(a)$ 的图像。
\end{prob}

\probgroup{研究下列函数在所指定区间内的连续性：}

\begin{prob}[S]{3779}
$F(\alpha) = \int_{0}^{+\infty} \frac{x\,\mathrm{d}x}{2 + x^\alpha}$，当 $\alpha > 2$。
\end{prob}

\begin{prob}[S]{3780}
$F(\alpha) = \int_{1}^{+\infty} \frac{\cos x}{x^\alpha}\,\mathrm{d}x$，当 $\alpha > 0$。
\end{prob}

\begin{prob}[S]{3781}
$F(\alpha) = \int_{0}^{\pi} \frac{\sin x}{x^\alpha (\pi - x)^\alpha}\,\mathrm{d}x$，当 $0 < \alpha < 2$。
\end{prob}

\begin{prob}[S]{3782}
$F(\alpha) = \int_{0}^{+\infty} \frac{\mathrm{e}^{-x}}{|\sin x|^\alpha}\,\mathrm{d}x$，当 $0 < \alpha < 1$。
\end{prob}

\begin{prob}[S]{3783}
$F(\alpha) = \int_{0}^{+\infty} \alpha \mathrm{e}^{-x\alpha^2}\,\mathrm{d}x$，其中 $-\infty < \alpha < +\infty$。
\end{prob}

\probsec{§3. 广义积分号下的微分法和积分法}

\begin{prob}[M]{3784}
利用公式
\[
\int_{0}^{1} x^{n-1}\,\mathrm{d}x = \frac{1}{n} \quad (n > 0)
\]
计算积分
\[
I = \int_{0}^{1} x^{n-1} \ln^m x\,\mathrm{d}x, \quad \text{其中 } m \text{ 为正整数}.
\]
\end{prob}

\begin{prob}[M]{3785}
利用公式
\[
\int_{0}^{+\infty} \frac{\mathrm{d}x}{x^2 + a} = \frac{\pi}{2\sqrt{a}} \quad (a > 0)
\]
计算积分
\[
I = \int_{0}^{+\infty} \frac{\mathrm{d}x}{(x^2 + a)^{n+1}}, \quad \text{其中 } n \text{ 为正整数}.
\]
\end{prob}

\begin{prob}[M]{3786}
证明：狄利克雷积分
\[
I(\alpha) = \int_{0}^{+\infty} \frac{\sin \alpha x}{x}\,\mathrm{d}x
\]
当 $\alpha \neq 0$ 时有导数，但是不能利用莱布尼茨法则来求它。
\probhint{令 $\alpha x = y$。}
\end{prob}

\begin{prob}[M]{3787}
证明：函数
\[
F(\alpha) = \int_{0}^{+\infty} \frac{\cos x}{1 + (x + \alpha)^2}\,\mathrm{d}x
\]
在区域 $-\infty < \alpha < +\infty$ 内连续并且可微。
\end{prob}

\begin{prob}[M]{3788}
从等式
\[
\frac{\mathrm{e}^{-ax} - \mathrm{e}^{-bx}}{x} = \int_{a}^{b} \mathrm{e}^{-xy}\,\mathrm{d}y
\]
出发，计算积分
\[
\int_{0}^{+\infty} \frac{\mathrm{e}^{-ax} - \mathrm{e}^{-bx}}{x}\,\mathrm{d}x \quad (a > 0, b > 0).
\]
\end{prob}

\begin{prob}[M]{3789}
证明傅茹兰公式：
\[
\int_{0}^{+\infty} \frac{f(ax) - f(bx)}{x}\,\mathrm{d}x = f(0) \ln \frac{b}{a} \quad (a > 0, b > 0),
\]
式中 $f(x)$ 为连续函数，积分 $\int_{A}^{+\infty} \frac{f(x)}{x}\,\mathrm{d}x$ 对任何 $A > 0$ 都有意义。
\end{prob}

\probgroup{利用傅茹兰公式，计算积分：}

\begin{prob}[S]{3790}
$\int_{0}^{+\infty} \frac{\cos ax - \cos bx}{x}\,\mathrm{d}x \ (a > 0, b > 0)$。
\end{prob}

\begin{prob}[S]{3791}
$\int_{0}^{+\infty} \frac{\sin ax - \sin bx}{x}\,\mathrm{d}x \ (a > 0, b > 0)$。
\end{prob}

\begin{prob}[S]{3792}
$\int_{0}^{+\infty} \frac{\arctan ax - \arctan bx}{x}\,\mathrm{d}x \ (a > 0, b > 0)$。
\end{prob}

\probgroup{利用对参数的微分法计算下列积分：}

\begin{prob}[S]{3793}
$\int_{0}^{+\infty} \frac{\mathrm{e}^{-\alpha x^2} - \mathrm{e}^{-\beta x^2}}{x}\,\mathrm{d}x \ (\alpha > 0, \beta > 0)$。
\end{prob}

\begin{prob}[S]{3794}
$\int_{0}^{+\infty} \left(\frac{\mathrm{e}^{-\alpha x} - \mathrm{e}^{-\beta x}}{x}\right)^2\,\mathrm{d}x \ (\alpha > 0, \beta > 0)$。
\end{prob}

\begin{prob}[S]{3795}
$\int_{0}^{+\infty} \frac{\mathrm{e}^{-\alpha x} - \mathrm{e}^{-\beta x}}{x} \sin mx\,\mathrm{d}x \ (\alpha > 0, \beta > 0)$。
\end{prob}

\begin{prob}[S]{3796}
$\int_{0}^{+\infty} \frac{\mathrm{e}^{-\alpha x} - \mathrm{e}^{-\beta x}}{x} \cos mx\,\mathrm{d}x \ (\alpha > 0, \beta > 0)$。
\end{prob}

\probgroup{计算下列积分：}

\begin{prob}[S]{3797}
$\int_{0}^{1} \frac{\ln(1 - \alpha^2 x^2)}{x^2 \sqrt{1 - x^2}}\,\mathrm{d}x \ (|\alpha| \leqslant 1)$。
\end{prob}

\begin{prob}[S]{3798}
$\int_{0}^{1} \frac{\ln(1 - \alpha^2 x^2)}{\sqrt{1 - x^2}}\,\mathrm{d}x \ (|\alpha| \leqslant 1)$。
\end{prob}

\begin{prob}[S]{3799}
$\int_{1}^{+\infty} \frac{\arctan \alpha x}{x^2 \sqrt{x^2 - 1}}\,\mathrm{d}x$。
\end{prob}

\begin{prob}[S]{3800}
$\int_{0}^{+\infty} \frac{\ln(\alpha^2 + x^2)}{\beta^2 + x^2}\,\mathrm{d}x$。
\end{prob}

\begin{prob}[S]{3801}
$\int_{0}^{+\infty} \frac{\arctan \alpha x \arctan \beta x}{x^2}\,\mathrm{d}x$。
\end{prob}

\begin{prob}[S]{3802}
$\int_{0}^{+\infty} \frac{\ln(1 + \alpha^2 x^2)\ln(1 + \beta^2 x^2)}{x^4}\,\mathrm{d}x$。
\end{prob}

\begin{prob}[M]{3803}
从公式
\[
I^2 = \int_{0}^{+\infty} \mathrm{e}^{-x^2}\,\mathrm{d}x \int_{0}^{+\infty} x\mathrm{e}^{-x^2 y^2}\,\mathrm{d}y
\]
出发，计算欧拉--泊松积分
\[
I = \int_{0}^{+\infty} \mathrm{e}^{-x^2}\,\mathrm{d}x.
\]
\end{prob}

\probgroup{利用欧拉--泊松积分，求下列积分：}

\begin{prob}[S]{3804}
$\displaystyle \int_{-\infty}^{+\infty} \mathrm{e}^{-(ax^2+2bx+c)} \,\mathrm{d}x \quad (a>0,\, ac-b^2>0)$.
\end{prob}

\begin{prob}[S]{3805}
$\displaystyle \int_{-\infty}^{+\infty} (a_1 x^2 + 2b_1 x + c_1)\mathrm{e}^{-(ax^2+2bx+c)} \,\mathrm{d}x \quad (a>0,\, ac-b^2>0)$.
\end{prob}

\begin{prob}[S]{3806}
$\displaystyle \int_{-\infty}^{+\infty} \mathrm{e}^{-ax^2} \cosh bx \,\mathrm{d}x \quad (a>0)$.
\end{prob}

\begin{prob}[S]{3807}
$\displaystyle \int_0^{+\infty} \mathrm{e}^{-\left(x^2+\frac{a^2}{x^2}\right)} \,\mathrm{d}x \quad (a>0)$.
\end{prob}

\begin{prob}[S]{3808}
$\displaystyle \int_0^{+\infty} \frac{\mathrm{e}^{-\alpha x^2}-\mathrm{e}^{-\beta x^2}}{x^2} \,\mathrm{d}x \quad (\alpha>0,\, \beta>0)$.
\end{prob}

\begin{prob}[S]{3809}
$\displaystyle \int_0^{+\infty} \mathrm{e}^{-ax^2} \cos bx \,\mathrm{d}x \quad (a>0)$.
\end{prob}

\begin{prob}[M]{3810}
\begin{enumerate}[label=(\alph*), itemsep=2pt, topsep=3pt]
\item $\displaystyle \int_0^{+\infty} x\mathrm{e}^{-ax^2} \sin bx \,\mathrm{d}x \quad (a>0)$；
\item $\displaystyle \int_0^{+\infty} x^{2n}\mathrm{e}^{-x^2} \cos 2bx \,\mathrm{d}x$ ($n$ 为正整数).
\end{enumerate}
\end{prob}

\begin{prob}[M]{3811}
证明：
\[
\lim_{x\to+\infty} \sqrt{x} \int_{-\delta}^{\delta} \mathrm{e}^{-axt^2} \,\mathrm{d}t = \sqrt{\frac{\pi}{a}} \quad (a>0,\, \delta>0).
\]
\end{prob}

\begin{prob}[M]{3812.1}
从积分
\[
I(\alpha) = \int_0^{+\infty} \mathrm{e}^{-\alpha x} \frac{\sin\beta x}{x} \,\mathrm{d}x \quad (\alpha\geqslant 0)
\]
出发，计算狄利克雷积分
\[
D(\beta) = \int_0^{+\infty} \frac{\sin\beta x}{x} \,\mathrm{d}x.
\]
\end{prob}

\begin{prob}[M]{3812.2}
积分正弦被定义为
\[
\mathrm{Si}\,x = \int_0^x \frac{\sin t}{t} \,\mathrm{d}t.
\]
函数 $y = \mathrm{Si}\,x$ 的图像大致具有怎样的形状？
\end{prob}

\probgroup{利用狄利克雷积分和傅茹兰积分，求下列积分：}

\begin{prob}[S]{3813}
$\displaystyle \int_0^{+\infty} \frac{\mathrm{e}^{-\alpha x^2}-\cos\beta x}{x^2} \,\mathrm{d}x \quad (\alpha>0)$.
\end{prob}

\begin{prob}[S]{3814}
$\displaystyle \int_0^{+\infty} \frac{\sin\alpha x \sin\beta x}{x} \,\mathrm{d}x \quad (|\alpha|\neq|\beta|)$.
\end{prob}

\begin{prob}[S]{3815}
$\displaystyle \int_0^{+\infty} \frac{\sin\alpha x \cos\beta x}{x} \,\mathrm{d}x$.
\end{prob}

\begin{prob}[S]{3816}
$\displaystyle \int_0^{+\infty} \frac{\sin^3 \alpha x}{x} \,\mathrm{d}x$.
\end{prob}

\begin{prob}[S]{3817}
$\displaystyle \int_0^{+\infty} \left(\frac{\sin\alpha x}{x}\right)^2 \,\mathrm{d}x$.
\end{prob}

\begin{prob}[S]{3818}
$\displaystyle \int_0^{+\infty} \left(\frac{\sin\alpha x}{x}\right)^3 \,\mathrm{d}x$.
\end{prob}

\begin{prob}[S]{3819}
$\displaystyle \int_0^{+\infty} \frac{\sin^4 x}{x^2} \,\mathrm{d}x$.
\end{prob}

\begin{prob}[S]{3820}
$\displaystyle \int_0^{+\infty} \frac{\sin^4 \alpha x - \sin^4 \beta x}{x} \,\mathrm{d}x \quad (\alpha\beta\neq 0)$.
\end{prob}

\begin{prob}[S]{3821}
$\displaystyle \int_0^{+\infty} \frac{\sin(x^2)}{x} \,\mathrm{d}x$.
\end{prob}

\begin{prob}[S]{3822}
$\displaystyle \int_0^{+\infty} \mathrm{e}^{-kx} \frac{\sin\alpha x \sin\beta x}{x^2} \,\mathrm{d}x \quad (k\geqslant 0,\, \alpha>0,\, \beta>0)$.
\end{prob}

\begin{prob}[M]{3823}
对于不同的 $x$ 值，求狄利克雷间断乘子
\[
D(x) = \frac{2}{\pi} \int_0^{+\infty} \sin\lambda \cos\lambda x \frac{\,\mathrm{d}\lambda}{\lambda}.
\]
作出函数 $y = D(x)$ 的图像.
\end{prob}

\begin{prob}[M]{3824}
计算积分：
\begin{enumerate}[label=(\alph*), itemsep=2pt, topsep=3pt]
\item $\displaystyle \mathrm{v.p.} \int_{-\infty}^{+\infty} \frac{\sin ax}{x+b} \,\mathrm{d}x$；
\item $\displaystyle \mathrm{v.p.} \int_{-\infty}^{+\infty} \frac{\cos ax}{x+b} \,\mathrm{d}x$.
\end{enumerate}
\end{prob}

\begin{prob}[M]{3825}
利用公式
\[
\frac{1}{1+x^2} = \int_0^{+\infty} \mathrm{e}^{-y(1+x^2)} \,\mathrm{d}y
\]
计算拉普拉斯积分
\[
L = \int_0^{+\infty} \frac{\cos ax}{1+x^2} \,\mathrm{d}x.
\]
\end{prob}

\begin{prob}[S]{3826}
计算积分
\[
L_1 = \int_0^{+\infty} \frac{x\sin ax}{1+x^2} \,\mathrm{d}x.
\]
\end{prob}

\probgroup{计算积分：}

\begin{prob}[S]{3827}
$\displaystyle \int_0^{+\infty} \frac{\sin^2 x}{1+x^2} \,\mathrm{d}x$.
\end{prob}

\begin{prob}[S]{3828}
$\displaystyle \int_0^{+\infty} \frac{\cos ax}{(1+x^2)^2} \,\mathrm{d}x$.
\end{prob}

\begin{prob}[S]{3829}
$\displaystyle \int_0^{+\infty} \frac{\cos ax}{ax^2+2bx+c} \,\mathrm{d}x \quad (a>0,\, ac-b^2>0)$.
\end{prob}

\begin{prob}[M]{3830}
利用公式
\[
\frac{1}{\sqrt{x}} = \frac{2}{\sqrt{\pi}} \int_0^{+\infty} \mathrm{e}^{-xy^2} \,\mathrm{d}y \quad (x>0)
\]
计算菲涅尔积分
\begin{align*}
\int_0^{+\infty} \sin(x^2) \,\mathrm{d}x &= \frac{1}{2} \int_0^{+\infty} \frac{\sin x}{\sqrt{x}} \,\mathrm{d}x, \\
\int_0^{+\infty} \cos(x^2) \,\mathrm{d}x &= \frac{1}{2} \int_0^{+\infty} \frac{\cos x}{\sqrt{x}} \,\mathrm{d}x.
\end{align*}
\end{prob}

\probgroup{求下列积分：}

\begin{prob}[S]{3831}
$\displaystyle \int_{-\infty}^{+\infty} \sin(ax^2+2bx+c) \,\mathrm{d}x \quad (a\neq 0)$.
\end{prob}

\begin{prob}[S]{3832}
$\displaystyle \int_{-\infty}^{+\infty} \sin x^2 \cdot \cos 2ax \,\mathrm{d}x$.
\end{prob}

\begin{prob}[S]{3833}
$\displaystyle \int_{-\infty}^{+\infty} \cos x^2 \cdot \cos 2ax \,\mathrm{d}x$.
\end{prob}

\begin{prob}[M]{3834}
证明公式：
1) $\displaystyle \int_0^{+\infty} \frac{\cos\alpha x}{a^2-x^2} \,\mathrm{d}x = \frac{\pi}{2a}\sin a\alpha$； \qquad
2) $\displaystyle \int_0^{+\infty} \frac{x\sin\alpha x}{a^2-x^2} \,\mathrm{d}x = -\frac{\pi}{2}\cos a\alpha$，
这里 $a\neq 0$，积分应理解为在柯西主值的意义上.
\end{prob}

\begin{prob}[L]{3835}
对于函数 $f(t)$，求拉普拉斯变换
\[
F(p) = \int_0^{+\infty} \mathrm{e}^{-pt} f(t) \,\mathrm{d}t \quad (p>0),
\]
设：
\begin{enumerate}[label=(\alph*), itemsep=2pt, topsep=3pt]
\item $f(t) = t^n$ ($n$ 为正整数)；
\item $f(t) = \sqrt{t}$；
\item $f(t) = \mathrm{e}^{\alpha t}$；
\item $f(t) = t\mathrm{e}^{-\alpha t}$；
\item $f(t) = \cos t$；
\item $\displaystyle f(t) = \frac{1-\mathrm{e}^{-t}}{t}$；
\item $f(t) = \sin\alpha\sqrt{t}$.
\end{enumerate}
\end{prob}

\begin{prob}[M]{3836}
证明公式：(利普希茨积分)
\[
\int_0^{+\infty} \mathrm{e}^{-at} J_0(bt) \,\mathrm{d}t = \frac{1}{\sqrt{a^2+b^2}} \quad (a>0),
\]
其中 $J_0(x) = \frac{1}{\pi} \int_0^{\pi} \cos(x\sin\varphi) \,\mathrm{d}\varphi$ 为 $0$ 阶贝塞尔函数（参阅习题 3726）.
\end{prob}

\begin{prob}[L]{3837}
求魏尔斯特拉斯变换
\[
F(x) = \frac{1}{\sqrt{\pi}} \int_{-\infty}^{+\infty} \mathrm{e}^{-(x-y)^2} f(y) \,\mathrm{d}y,
\]
设：
\begin{enumerate}[label=(\alph*), itemsep=2pt, topsep=3pt]
\item $f(y) = 1$；
\item $f(y) = y^2$；
\item $f(y) = \mathrm{e}^{2ay}$；
\item $f(y) = \cos ay$.
\end{enumerate}
\end{prob}

\begin{prob}[M]{3838}
切比雪夫--埃尔米特多项式由公式
\[
H_n(x) = (-1)^n \mathrm{e}^{x^2} \frac{\mathrm{d}^n}{\mathrm{d}x^n}(\mathrm{e}^{-x^2}) \quad (n=0,1,2,\cdots)
\]
定义，证明：
\[
\int_{-\infty}^{+\infty} H_m(x) H_n(x) \mathrm{e}^{-x^2} \,\mathrm{d}x =
\begin{cases}
0, & m\neq n, \\
2^n n! \sqrt{\pi}, & m=n.
\end{cases}
\]
\end{prob}

\begin{prob}[M]{3839}
计算在概率论中有重要意义的积分
\[
\varphi(x) = \frac{1}{2\pi\sigma_1\sigma_2} \int_{-\infty}^{+\infty} \mathrm{e}^{-\frac{1}{2}\left[\frac{\xi^2}{\sigma_1^2}+\frac{(x-\xi)^2}{\sigma_2^2}\right]} \,\mathrm{d}\xi \quad (\sigma_1>0,\, \sigma_2>0).
\]
\end{prob}

\begin{prob}[L]{3840}
设函数 $f(x)$ 在区间 $(-\infty,+\infty)$ 内连续且绝对可积．证明：积分
\[
u(x,t) = \frac{1}{2a\sqrt{\pi t}} \int_{-\infty}^{+\infty} f(\xi)\mathrm{e}^{-\frac{(\xi-x)^2}{4a^2 t}} \,\mathrm{d}\xi
\]
满足热传导方程
\[
\frac{\partial u}{\partial t} = \frac{1}{a^2}\frac{\partial^2 u}{\partial x^2}
\]
及初始条件
\[
\lim_{t\to+0} u(x,t) = f(x).
\]
\end{prob}

\probsec{§4. 欧拉积分}

\begin{prob}[M]{3841}
证明：$\Gamma$ 函数 $\Gamma(x)$ 在区域 $x>0$ 内连续，并且有各阶连续导数.
\end{prob}

\begin{prob}[M]{3842}
证明：$\beta$ 函数 $\beta(x,y)$ 在区域 $x>0, y>0$ 内连续，并且有各阶连续导数.
\end{prob}

\probgroup{利用欧拉积分计算下列积分：}

\begin{prob}[S]{3843}
$\displaystyle \int_0^1 \sqrt{x-x^2} \,\mathrm{d}x$.
\end{prob}

\begin{prob}[S]{3844}
$\displaystyle \int_0^a x^2\sqrt{a^2-x^2} \,\mathrm{d}x \quad (a>0)$.
\end{prob}

\begin{prob}[S]{3845}
$\displaystyle \int_0^{+\infty} \frac{\sqrt[4]{x}}{(1+x)^2} \,\mathrm{d}x$.
\end{prob}

\begin{prob}[S]{3846}
$\displaystyle \int_0^{+\infty} \frac{\mathrm{d}x}{1+x^3}$.
\end{prob}

\begin{prob}[S]{3847}
$\displaystyle \int_0^{+\infty} \frac{x^2\,\mathrm{d}x}{1+x^4}$.
\end{prob}

\begin{prob}[S]{3848}
$\displaystyle \int_0^{\frac{\pi}{2}} \sin^6 x \cdot \cos^4 x \,\mathrm{d}x$.
\end{prob}

\begin{prob}[S]{3849}
$\displaystyle \int_0^1 \frac{\mathrm{d}x}{\sqrt[n]{1-x^n}} \quad (n>0)$.
\end{prob}

\begin{prob}[S]{3850}
$\displaystyle \int_0^{+\infty} x^{2n}\mathrm{e}^{-x^2} \,\mathrm{d}x$ ($n$ 为正整数).
\end{prob}

\probgroup{求下列积分的存在域，并用欧拉积分表示这些积分：}

\begin{prob}[S]{3851}
$\displaystyle \int_0^{+\infty} \frac{x^{m-1}}{1+x^n} \,\mathrm{d}x \quad (n>0)$.
\end{prob}

\begin{prob}[S]{3852}
$\displaystyle \int_0^{+\infty} \frac{x^{m-1}}{(1+x)^n} \,\mathrm{d}x$.
\end{prob}

\begin{prob}[S]{3853}
$\displaystyle \int_0^{+\infty} \frac{x^m\,\mathrm{d}x}{(a+bx^n)^p} \quad (a>0,\, b>0,\, n>0)$.
\end{prob}

\begin{prob}[S]{3854}
$\displaystyle \int_a^b \frac{(x-a)^m(b-x)^n}{(x+c)^{m+n+2}} \,\mathrm{d}x \quad (0<a<b,\, c>0)$.
\end{prob}

\begin{prob}[S]{3855}
$\displaystyle \int_0^1 \frac{\mathrm{d}x}{\sqrt[n]{1-x^m}} \quad (m>0)$.
\end{prob}

\begin{prob}[S]{3856}
$\displaystyle \int_0^{\frac{\pi}{2}} \sin^m x \cos^n x \,\mathrm{d}x$.
\end{prob}

\begin{prob}[S]{3857}
$\displaystyle \int_0^{\frac{\pi}{2}} \tan^n x \,\mathrm{d}x$.
\end{prob}

\begin{prob}[S]{3858}
$\displaystyle \int_0^{\pi} \frac{\sin^{n-1} x}{(1+k\cos x)^n} \,\mathrm{d}x \quad (0<|k|<1)$.
\end{prob}

\begin{prob}[S]{3859}
$\displaystyle \int_0^{+\infty} \mathrm{e}^{-x^n} \,\mathrm{d}x \quad (n>0)$.
\end{prob}

\begin{prob}[S]{3860}
$\displaystyle \int_0^{+\infty} x^m \mathrm{e}^{-x^n} \,\mathrm{d}x$.
\end{prob}

\begin{prob}[S]{3861}
$\displaystyle \int_0^1 \left(\ln\frac{1}{x}\right)^p \,\mathrm{d}x$.
\end{prob}

\begin{prob}[S]{3862}
$\displaystyle \int_0^{+\infty} x^p \mathrm{e}^{-ax} \ln x \,\mathrm{d}x \quad (a>0)$.
\end{prob}

\begin{prob}[S]{3863}
$\displaystyle \int_0^{+\infty} \frac{x^{p-1}\ln x}{1+x} \,\mathrm{d}x \quad (p>0)$.
\end{prob}

\begin{prob}[L]{3864}
\begin{enumerate}[label=(\alph*), itemsep=2pt, topsep=3pt]
\item $\displaystyle \int_0^{+\infty} \frac{x^{p-1}\ln^2 x}{1+x} \,\mathrm{d}x \quad (p>0)$；
\item $\displaystyle \int_0^{\infty} \frac{x\ln x}{1+x^3} \,\mathrm{d}x$；
\item $\displaystyle \int_0^{\infty} \frac{\ln^2 x}{1+x^4} \,\mathrm{d}x$.
\end{enumerate}
\end{prob}

\begin{prob}[S]{3865}
$\displaystyle \int_0^{+\infty} \frac{x^{p-1}-x^{q-1}}{(1+x)\ln x} \,\mathrm{d}x$.
\end{prob}

\begin{prob}[M]{3866}
$\displaystyle \int_0^1 \frac{x^{p-1}-x^{-p}}{1-x} \,\mathrm{d}x \quad (0<p<1)$.
\probhint{此积分可看作 $\lim_{\varepsilon\to+0} [\beta(p,\varepsilon) - \beta(1-p,\varepsilon)]$.}
\end{prob}

\begin{prob}[S]{3867}
$\displaystyle \int_0^{+\infty} \frac{\sinh\alpha x}{\sinh\beta x} \,\mathrm{d}x \quad (0<\alpha<\beta)$.
\end{prob}

\begin{prob}[S]{3868}
$\displaystyle \int_0^1 \ln\Gamma(x) \,\mathrm{d}x$.
\end{prob}

\begin{prob}[S]{3869}
$\displaystyle \int_a^{a+1} \ln\Gamma(x) \,\mathrm{d}x \quad (a>0)$.
\end{prob}

\begin{prob}[S]{3870}
$\displaystyle \int_0^1 \ln\Gamma(x)\sin\pi x \,\mathrm{d}x$.
\end{prob}

\begin{prob}[S]{3871}
$\displaystyle \int_0^1 \ln\Gamma(x)\cos 2n\pi x \,\mathrm{d}x$ ($n$ 为正整数).
\end{prob}

\probgroup{证明等式：}

\begin{prob}[M]{3872}
$\displaystyle \int_0^1 \frac{\mathrm{d}x}{\sqrt{1-x^4}} \cdot \int_0^1 \frac{x^2\,\mathrm{d}x}{\sqrt{1-x^4}} = \frac{\pi}{4}$.
\end{prob}

\begin{prob}[M]{3873}
$\displaystyle \int_0^{+\infty} \mathrm{e}^{-x^4}\,\mathrm{d}x \cdot \int_0^{+\infty} x^2\mathrm{e}^{-x^4}\,\mathrm{d}x = \frac{\pi}{8\sqrt{2}}$.
\end{prob}

\begin{prob}[M]{3874}
$\displaystyle \prod_{m=1}^{n} \int_0^{+\infty} x^{m-1}\mathrm{e}^{-x^n}\,\mathrm{d}x = \left(\frac{1}{n}\right)^{n+\frac{1}{2}} (2\pi)^{\frac{n-1}{2}}$.
\end{prob}

\begin{prob}[M]{3875}
$\displaystyle \lim_{n\to\infty} \int_0^{+\infty} \mathrm{e}^{-x^n}\,\mathrm{d}x = 1$.
\end{prob}

\probgroup{利用等式 $\frac{1}{x^m} = \frac{1}{\Gamma(m)} \int_0^{+\infty} t^{m-1}\mathrm{e}^{-xt}\,\mathrm{d}t \ (x>0)$，求积分：}

\begin{prob}[S]{3876}
$\displaystyle \int_0^{+\infty} \frac{\cos ax}{x^m}\,\mathrm{d}x \quad (0<m<1)$.
\end{prob}

\begin{prob}[S]{3877}
$\displaystyle \int_0^{+\infty} \frac{\sin ax}{x^m}\,\mathrm{d}x \quad (0<m<2)$.
\end{prob}

\begin{prob}[M]{3878}
证明欧拉公式：
\begin{enumerate}[label=(\alph*), itemsep=2pt, topsep=3pt]
\item $\displaystyle \int_0^{+\infty} t^{x-1}\mathrm{e}^{-\lambda t\cos\alpha}\cos(\lambda t\sin\alpha)\,\mathrm{d}t = \frac{\Gamma(x)}{\lambda^x}\cos\alpha x$；
\item $\displaystyle \int_0^{+\infty} t^{x-1}\mathrm{e}^{-\lambda t\cos\alpha}\sin(\lambda t\sin\alpha)\,\mathrm{d}t = \frac{\Gamma(x)}{\lambda^x}\sin\alpha x$
\end{enumerate}
$\left(\lambda>0,\, x>0,\, -\frac{\pi}{2}<\alpha<\frac{\pi}{2}\right)$.
\end{prob}

\begin{prob}[M]{3879}
求曲线
\[
r^n = a^n \cos n\varphi \quad (a>0,\, n \text{ 为正整数})
\]
的弧长.
\end{prob}

\begin{prob}[M]{3880}
求曲线
\[
|x|^n + |y|^n = a^n \quad (n>0,\, a>0)
\]
所围的面积.
\end{prob}

\probsec{§5. 傅里叶积分公式}

\probgroup{用傅里叶积分表示下列函数：}

\begin{prob}[S]{3881}
$f(x) = \begin{cases} 1, & |x| < 1, \\ 0, & |x| > 1. \end{cases}$
\end{prob}

\begin{prob}[S]{3882}
$f(x) = \begin{cases} \operatorname{sgn} x, & |x| < 1, \\ 0, & |x| > 1. \end{cases}$
\end{prob}

\begin{prob}[S]{3883}
$f(x) = \operatorname{sgn}(x - a) - \operatorname{sgn}(x - b)\ (b > a).$
\end{prob}

\begin{prob}[S]{3884}
$f(x) = \begin{cases} h\left(1 - \dfrac{|x|}{a}\right), & |x| \leqslant a, \\ 0, & |x| > a. \end{cases}$
\end{prob}

\begin{prob}[S]{3885}
$f(x) = \dfrac{1}{a^2 + x^2}\ (a > 0).$
\end{prob}

\begin{prob}[S]{3886}
$f(x) = \dfrac{x}{a^2 + x^2}\ (a > 0).$
\end{prob}

\begin{prob}[S]{3887}
$f(x) = \begin{cases} \sin x, & |x| \leqslant \pi, \\ 0, & |x| > \pi. \end{cases}$
\end{prob}

\begin{prob}[S]{3888}
$f(x) = \begin{cases} \cos x, & |x| \leqslant \dfrac{\pi}{2}, \\ 0, & |x| > \dfrac{\pi}{2}. \end{cases}$
\end{prob}

\begin{prob}[S]{3889}
$f(x) = \begin{cases} A\sin \omega x, & |t| \leqslant \dfrac{2\pi n}{\omega}, \\ 0, & |t| > \dfrac{2\pi n}{\omega} \end{cases}\ (n\text{ 为正整数}).$
\end{prob}

\begin{prob}[S]{3890}
$f(x) = \mathrm{e}^{-\alpha |x|}\ (\alpha > 0).$
\end{prob}

\begin{prob}[S]{3891}
$f(x) = \mathrm{e}^{-\alpha |x|} \cos \beta x\ (\alpha > 0).$
\end{prob}

\begin{prob}[S]{3892}
$f(x) = \mathrm{e}^{-\alpha |x|} \sin \beta x\ (\alpha > 0).$
\end{prob}

\begin{prob}[S]{3893}
$f(x) = \mathrm{e}^{-x^2}.$
\end{prob}

\begin{prob}[S]{3894}
$f(x) = x\mathrm{e}^{-x^2}.$
\end{prob}

\begin{prob}[M]{3895}
用傅里叶积分来表示函数
\[
f(x) = \mathrm{e}^{-x}\ (0 < x < +\infty),
\]
(a) 用偶式延拓；\qquad (b) 用奇式延拓.
\end{prob}

\probgroup{对于下列函数 $f(t)$，求傅里叶变换
\[
F(x) = \frac{1}{\sqrt{2\pi}} \int_{-\infty}^{+\infty} f(t)\mathrm{e}^{-itx}\,\mathrm{d}t = \frac{1}{\sqrt{2\pi}} \lim_{l \to +\infty} \int_{-l}^{l} f(t)\mathrm{e}^{-itx}\,\mathrm{d}t :
\]}

\begin{prob}[S]{3896}
$f(x) = \mathrm{e}^{-\alpha |x|}\ (\alpha > 0).$
\end{prob}

\begin{prob}[S]{3897}
$f(x) = x\mathrm{e}^{-\alpha |x|}\ (\alpha > 0).$
\end{prob}

\begin{prob}[S]{3898}
$f(x) = \mathrm{e}^{-\frac{x^2}{2}}.$
\end{prob}

\begin{prob}[S]{3899}
$f(x) = \mathrm{e}^{-\frac{x^2}{2}} \cos \alpha x.$
\end{prob}

\begin{prob}[M]{3900}
求函数 $\varphi(x)$ 及 $\psi(x)$，设：
\begin{enumerate}[label=(\alph*), itemsep=2pt, topsep=3pt]
\item $\displaystyle \int_{0}^{+\infty} \varphi(y)\cos xy\,\mathrm{d}y = \frac{1}{1 + x^2};$
\item $\displaystyle \int_{0}^{+\infty} \psi(y)\sin xy\,\mathrm{d}y = \mathrm{e}^{-x}\ (x > 0).$
\end{enumerate}
\end{prob}

\probchapter{第八章\quad 多重积分和曲线积分}

\probsec{§1. 二重积分}

\begin{prob}[M]{3901}
把积分 $\iint\limits_{\substack{0 \leqslant x \leqslant 1 \\ 0 \leqslant y \leqslant 1}} xy\,\mathrm{d}x\,\mathrm{d}y$ 当作积分和的极限，用直线
\[
x = \frac{i}{n}, \quad y = \frac{j}{n} \quad (i, j = 1, 2, \cdots, n - 1)
\]
把积分域分为许多正方形，并选取被积函数在这些正方形之右顶点的值，计算此积分。
\end{prob}

\begin{prob}[M]{3902}
用直线
\[
x = 1 + \frac{i}{n}, \quad y = 1 + \frac{2j}{n} \quad (i, j = 0, 1, \cdots, n)
\]
把区域 $1 \leqslant x \leqslant 2$, $1 \leqslant y \leqslant 3$ 分为许多矩形，列出函数 $f(x, y) = x^2 + y^2$ 在此区域内的下积分和 $\underline{S}$ 与上积分和 $\overline{S}$。

当 $n \to \infty$ 时，上积分和与下积分和的极限等于什么？
\end{prob}

\begin{prob}[M]{3903}
用一组顶点 $A_{ij}$ 位于整数点的正方形作为积分域的近似域，并取被积函数在每个正方形距原点最远的顶点之值，近似地计算积分
\[
\iint\limits_{x^2 + y^2 \leqslant 25} \frac{\mathrm{d}x\,\mathrm{d}y}{\sqrt{24 + x^2 + y^2}},
\]
并与精确值加以比较。
\end{prob}

\begin{prob}[M]{3904}
用直线 $x = \text{常数}$, $y = \text{常数}$, $x + y = \text{常数}$ 把积分域 $S$ 分为四个相同的三角形，并取被积函数在每个三角形的质心之值，近似地计算积分
\[
\iint\limits_{S} \sqrt{x + y}\,\mathrm{d}S.
\]
其中 $S$ 是以直线 $x = 0, y = 0$ 及 $x + y = 1$ 为边的三角形。
\end{prob}

\begin{prob}[M]{3905}
把积分域 $S: \{x^2 + y^2 \leqslant 1\}$ 分为有限个直径小于 $\delta$ 的可求积的子区域 $\Delta S_i$ $(i = 1, 2, \cdots, n)$。怎样的值 $\delta$ 能保证不等式
\[
\left| \iint\limits_{S} \sin(x + y)\,\mathrm{d}S - \sum_{i=1}^{n} \sin(x_i + y_i)\Delta S_i \right| < 0.001
\]
成立？其中 $(x_i, y_i) \in \Delta S_i$。
\end{prob}

\probgroup{计算积分：}

\begin{prob}[S]{3906}
$\int_{0}^{1}\,\mathrm{d}x \int_{0}^{1} (x + y)\,\mathrm{d}y$。
\end{prob}

\begin{prob}[S]{3907}
$\int_{0}^{1}\,\mathrm{d}x \int_{x^2}^{x} xy^2\,\mathrm{d}y$。
\end{prob}

\begin{prob}[S]{3908}
$\int_{0}^{2\pi}\,\mathrm{d}\varphi \int_{0}^{a} r^2 \sin^2\varphi\,\mathrm{d}r$。
\end{prob}

\begin{prob}[M]{3909}
设 $R$ 为矩形
\[
a \leqslant x \leqslant A, \quad b \leqslant y \leqslant B,
\]
函数 $X(x)$ 和 $Y(y)$ 在相应区间上连续，证明等式：
\[
\iint\limits_{R} X(x)Y(y)\,\mathrm{d}x\,\mathrm{d}y = \int_{a}^{A} X(x)\,\mathrm{d}x \cdot \int_{b}^{B} Y(y)\,\mathrm{d}y.
\]
\end{prob}

\begin{prob}[M]{3910}
设
\[
f(x, y) = F''_{xy}(x, y),
\]
计算
\[
I = \int_{a}^{A}\,\mathrm{d}x \int_{b}^{B} f(x, y)\,\mathrm{d}y.
\]
\end{prob}

\begin{prob}[M]{3911}
设 $f(x)$ 为区间 $a \leqslant x \leqslant b$ 内的连续函数。证明不等式：
\[
\left[ \int_{a}^{b} f(x)\,\mathrm{d}x \right]^2 \leqslant (b - a) \int_{a}^{b} f^2(x)\,\mathrm{d}x,
\]
且仅当 $f(x) = \text{常数}$ 时等号成立。
\probhint{研究积分
\[
\int_{a}^{b}\,\mathrm{d}x \int_{a}^{b} [f(x) - f(y)]^2\,\mathrm{d}y.
\]}
\end{prob}

\begin{prob}[L]{3912}
下列积分有怎样的符号？
\begin{enumerate}[label=(\alph*), itemsep=2pt, topsep=3pt]
\item $\iint\limits_{|x| + |y| \leqslant 1} \ln(x^2 + y^2)\,\mathrm{d}x\,\mathrm{d}y$;
\item $\iint\limits_{x^2 + y^2 \leqslant 4} \sqrt[3]{1 - x^2 - y^2}\,\mathrm{d}x\,\mathrm{d}y$;
\item $\iint\limits_{\substack{0 \leqslant x \leqslant 1 \\ -1 \leqslant y \leqslant 1 - x}} \arcsin(x + y)\,\mathrm{d}x\,\mathrm{d}y$.
\end{enumerate}
\end{prob}

\begin{prob}[M]{3913}
求函数
\[
f(x, y) = \sin^2 x \sin^2 y
\]
在正方形 $0 \leqslant x \leqslant \pi, 0 \leqslant y \leqslant \pi$ 内的平均值。
\end{prob}

\begin{prob}[M]{3914}
利用中值定理估计积分
\[
I = \iint\limits_{|x| + |y| \leqslant 10} \frac{\mathrm{d}x\,\mathrm{d}y}{100 + \cos^2 x + \cos^2 y}.
\]
\end{prob}

\begin{prob}[M]{3915}
求圆 $(x - a)^2 + (y - b)^2 \leqslant R^2$ 上的点到原点的距离之平方的平均值。
\end{prob}

\probgroup{在问题 3916—3922 中，对二重积分 $\iint\limits_{\Omega} f(x, y)\,\mathrm{d}x\,\mathrm{d}y$ 按所给区域 $\Omega$ 依两个不同的顺序安置积分的上下限。}

\begin{prob}[S]{3916}
$\Omega$ — 以 $O(0, 0)$, $A(1, 0)$, $B(1, 1)$ 为顶点的三角形。
\end{prob}

\begin{prob}[S]{3917}
$\Omega$ — 以 $O(0, 0)$, $A(2, 1)$, $B(-2, 1)$ 为顶点的三角形。
\end{prob}

\begin{prob}[S]{3918}
$\Omega$ — 以 $O(0, 0)$, $A(1, 0)$, $B(1, 2)$, $C(0, 1)$ 为顶点的梯形。
\end{prob}

\begin{prob}[S]{3919}
$\Omega$ — 圆 $x^2 + y^2 \leqslant 1$。
\end{prob}

\begin{prob}[S]{3920}
$\Omega$ — 圆 $x^2 + y^2 \leqslant y$。
\end{prob}

\begin{prob}[S]{3921}
$\Omega$ — 被曲线 $y = x^2$ 和直线 $y = 1$ 所包围的区域。
\end{prob}

\begin{prob}[S]{3922}
$\Omega$ — 圆环 $1 \leqslant x^2 + y^2 \leqslant 4$。
\end{prob}

\begin{prob}[M]{3923}
证明狄利克雷公式：
\[
\int_{0}^{a}\,\mathrm{d}x \int_{0}^{x} f(x, y)\,\mathrm{d}y = \int_{0}^{a}\,\mathrm{d}y \int_{y}^{a} f(x, y)\,\mathrm{d}x \quad (a > 0).
\]
\end{prob}

\probgroup{在下列积分中改变积分的顺序：}

\begin{prob}[S]{3924}
$\int_{0}^{2}\,\mathrm{d}x \int_{x}^{2x} f(x, y)\,\mathrm{d}y$。
\end{prob}

\begin{prob}[S]{3925}
$\int_{-6}^{2}\,\mathrm{d}x \int_{\frac{x^2}{4} - 1}^{2 - x} f(x, y)\,\mathrm{d}y$。
\end{prob}

\begin{prob}[S]{3926}
$\int_{0}^{1}\,\mathrm{d}x \int_{x^3}^{x^2} f(x, y)\,\mathrm{d}y$。
\end{prob}

\begin{prob}[S]{3927}
$\int_{-1}^{1}\,\mathrm{d}x \int_{-\sqrt{1 - x^2}}^{1 - x^2} f(x, y)\,\mathrm{d}y$。
\end{prob}

\begin{prob}[S]{3928}
$\int_{1}^{2}\,\mathrm{d}x \int_{2 - x}^{\sqrt{2x - x^2}} f(x, y)\,\mathrm{d}y$。
\end{prob}

\begin{prob}[S]{3929}
$\int_{0}^{2a}\,\mathrm{d}x \int_{\sqrt{2ax - x^2}}^{\sqrt{2ax}} f(x, y)\,\mathrm{d}y \quad (a > 0)$。
\end{prob}

\begin{prob}[S]{3930}
$\int_{1}^{\mathrm{e}}\,\mathrm{d}x \int_{0}^{\ln x} f(x, y)\,\mathrm{d}y$。
\end{prob}

\begin{prob}[S]{3931}
$\int_{0}^{2\pi}\,\mathrm{d}x \int_{0}^{\sin x} f(x, y)\,\mathrm{d}y$。
\end{prob}

\probgroup{计算下列积分：}

\begin{prob}[M]{3932}
$\iint\limits_{\Omega} xy^2\,\mathrm{d}x\,\mathrm{d}y$, 其中 $\Omega$ 是被抛物线 $y^2 = 2px$ 和直线 $x = \frac{p}{2}$ $(p > 0)$ 所包围的区域。
\end{prob}

\begin{prob}[M]{3933}
$\iint\limits_{\Omega} \frac{\mathrm{d}x\,\mathrm{d}y}{\sqrt{2a - x}}$ $(a > 0)$, 其中 $\Omega$ 是以圆心在点 $(a, a)$ 半径为 $a$ 的圆周（它与坐标轴相切）的较短弧和坐标轴为界的区域。
\end{prob}

\begin{prob}[M]{3934}
$\iint\limits_{\Omega} |xy|\,\mathrm{d}x\,\mathrm{d}y$, 其中 $\Omega$ 是以 $a$ 为半径，以坐标原点为圆心的圆。
\end{prob}

\begin{prob}[M]{3935}
$\iint\limits_{\Omega} (x^2 + y^2)\,\mathrm{d}x\,\mathrm{d}y$, 其中 $\Omega$ 是以 $y = x, y = x + a, y = a$ 和 $y = 3a$ $(a > 0)$ 为边的平行四边形。
\end{prob}

\begin{prob}[M]{3936}
$\iint\limits_{\Omega} y^2\,\mathrm{d}x\,\mathrm{d}y$, 其中 $\Omega$ 是被横轴和摆线
\[
x = a(t - \sin t), \quad y = a(1 - \cos t) \quad (0 \leqslant t \leqslant 2\pi)
\]
的第一拱所包围的区域。
\end{prob}

\probgroup{在二重积分
\[
\iint\limits_{\Omega} f(x, y)\,\mathrm{d}x\,\mathrm{d}y
\]
中，令 $x = r\cos\varphi$ 和 $y = r\sin\varphi$，变换为极坐标 $r$ 和 $\varphi$，并配置积分的上下限，设：}

\begin{prob}[S]{3937}
$\Omega$ — 圆 $x^2 + y^2 \leqslant a^2$。
\end{prob}

\begin{prob}[S]{3938}
$\Omega$ — 圆 $x^2 + y^2 \leqslant ax$ $(a > 0)$。
\end{prob}

\begin{prob}[S]{3939}
$\Omega$ — 环 $a^2 \leqslant x^2 + y^2 \leqslant b^2$。
\end{prob}

\begin{prob}[S]{3940}
$\Omega$ — 三角形 $0 \leqslant x \leqslant 1; 0 \leqslant y \leqslant 1 - x$。
\end{prob}

\begin{prob}[S]{3941}
$\Omega$ — 区域 $-a \leqslant x \leqslant a; \frac{x^2}{a} \leqslant y \leqslant a$。
\end{prob}

\begin{prob}[M]{3942}
在怎样的情况下，当变换为极坐标之后，积分的上下限是常数？
\end{prob}

\probgroup{在下列积分中，令 $x = r\cos\varphi$ 和 $y = r\sin\varphi$，变换为极坐标 $r$ 和 $\varphi$，并依两种不同的顺序配置积分的上下限：}

\begin{prob}[S]{3943}
$\int_{0}^{1}\,\mathrm{d}x \int_{0}^{1} f(x, y)\,\mathrm{d}y$。
\end{prob}

\begin{prob}[S]{3944}
$\int_{0}^{1}\,\mathrm{d}x \int_{1 - x}^{\sqrt{1 - x^2}} f(x, y)\,\mathrm{d}y$。
\end{prob}

\begin{prob}[S]{3945}
$\int_{0}^{2}\,\mathrm{d}x \int_{x}^{\sqrt{3}x} f(\sqrt{x^2 + y^2})\,\mathrm{d}y$。
\end{prob}

\begin{prob}[S]{3946}
$\int_{0}^{1}\,\mathrm{d}x \int_{0}^{x^2} f(x, y)\,\mathrm{d}y$。
\end{prob}

\begin{prob}[M]{3947}
$\iint\limits_{\Omega} f(x, y)\,\mathrm{d}x\,\mathrm{d}y$, 其中区域 $\Omega$ 由曲线 $(x^2 + y^2)^2 = a^2(x^2 - y^2)$ $(x \geqslant 0)$ 围成。
\end{prob}

\probgroup{令 $r$ 和 $\varphi$ 为极坐标，在下列积分中变更积分的顺序：}

\begin{prob}[S]{3948}
$\int_{-\frac{\pi}{2}}^{\frac{\pi}{2}}\,\mathrm{d}\varphi \int_{0}^{a\cos\varphi} f(\varphi, r)\,\mathrm{d}r \quad (a > 0)$。
\end{prob}

\begin{prob}[S]{3949}
$\int_{0}^{\frac{\pi}{2}}\,\mathrm{d}\varphi \int_{0}^{a\sqrt{\sin 2\varphi}} f(\varphi, r)\,\mathrm{d}r \quad (a > 0)$。
\end{prob}

\begin{prob}[S]{3950}
$\int_{0}^{a}\,\mathrm{d}\varphi \int_{0}^{\varphi} f(\varphi, r)\,\mathrm{d}r \quad (0 < a < 2\pi)$。
\end{prob}

\probgroup{变换成极坐标，把二重积分化为一重积分：}

\begin{prob}[S]{3951}
$\iint\limits_{x^2 + y^2 \leqslant 1} f(\sqrt{x^2 + y^2})\,\mathrm{d}x\,\mathrm{d}y$。
\end{prob}

\begin{prob}[S]{3952}
$\iint\limits_{\Omega} f(\sqrt{x^2 + y^2})\,\mathrm{d}x\,\mathrm{d}y$, 其中 $\Omega = \{|y| \leqslant |x|; |x| \leqslant 1\}$。
\end{prob}

\begin{prob}[S]{3953}
$\iint\limits_{x^2 + y^2 \leqslant x} f\left(\frac{y}{x}\right)\,\mathrm{d}x\,\mathrm{d}y$。
\end{prob}

\probgroup{变换成极坐标，计算下列二重积分：}

\begin{prob}[S]{3954}
$\iint\limits_{x^2 + y^2 \leqslant a^2} \sqrt{x^2 + y^2}\,\mathrm{d}x\,\mathrm{d}y$。
\end{prob}

\begin{prob}[S]{3955}
$\iint\limits_{\pi^2 \leqslant x^2 + y^2 \leqslant 4\pi^2} \sin\sqrt{x^2 + y^2}\,\mathrm{d}x\,\mathrm{d}y$。
\end{prob}

\begin{prob}[M]{3956}
利用函数组
\[
u = \frac{y^2}{x}, \quad v = \sqrt{xy}
\]
把正方形 $S\{a < x < a + h, b < y < b + h\} (a > 0, b > 0)$ 变换为区域 $S'$。求区域 $S'$ 的面积与区域 $S$ 的面积之比。当 $h \to 0$ 时，此比值的极限等于什么？
\end{prob}

\probgroup{引入新的变量 $u$, $v$ 来代替 $x$, $y$，并确定下列二重积分中的积分限：}

\begin{prob}[S]{3957}
$\int_{a}^{b}\,\mathrm{d}x \int_{\alpha x}^{\beta x} f(x, y)\,\mathrm{d}y \quad (0 < a < b; 0 < \alpha < \beta)$, 令 $u = x, v = \frac{y}{x}$。
\end{prob}

\begin{prob}[S]{3958}
$\int_{0}^{2}\,\mathrm{d}x \int_{1 - x}^{2 - x} f(x, y)\,\mathrm{d}y$, 令 $u = x + y, v = x - y$。
\end{prob}

\begin{prob}[M]{3959}
$\iint\limits_{\Omega} f(x, y)\,\mathrm{d}x\,\mathrm{d}y$, 其中 $\Omega$ 是被曲线 $\sqrt{x} + \sqrt{y} = \sqrt{a}, x = 0, y = 0$ $(a > 0)$ 所
包围的区域，令
\[
x = u \cos^4 v, \quad y = u \sin^4 v.
\]
\end{prob}

\begin{prob}[M]{3960}
证明：变量代换
\[
x + y = \xi, \quad y = \xi\eta
\]
把三角形 $0 \leqslant x \leqslant 1, \, 0 \leqslant y \leqslant 1 - x$ 变为单位正方形 $0 \leqslant \xi \leqslant 1, \, 0 \leqslant \eta \leqslant 1$。
\end{prob}

\begin{prob}[M]{3961}
在怎样的变量代换下，由曲线 $xy = 1, \, xy = 2, \, x - y + 1 = 0, \, x - y - 1 = 0$ $(x > 0, \, y > 0)$ 围成的曲线四边形被变换成矩形，且其边平行于坐标轴？
\end{prob}

\probgroup{进行适当的变量代换，把二重积分化为一重积分：}

\begin{prob}[S]{3962}
$\displaystyle \iint\limits_{|x| + |y| \leqslant 1} f(x + y) \,\mathrm{d}x\,\mathrm{d}y$。
\end{prob}

\begin{prob}[S]{3963}
$\displaystyle \iint\limits_{x^2 + y^2 \leqslant 1} f(ax + by + c) \,\mathrm{d}x\,\mathrm{d}y \quad (a^2 + b^2 \neq 0)$。
\end{prob}

\begin{prob}[S]{3964}
$\displaystyle \iint\limits_{\Omega} f(xy) \,\mathrm{d}x\,\mathrm{d}y$，其中区域 $\Omega$ 由曲线 $xy = 1, \, xy = 2, \, y = x, \, y = 4x$ $(x > 0, \, y > 0)$ 围成。
\end{prob}

\probgroup{计算下列二重积分：}

\begin{prob}[S]{3965}
$\displaystyle \iint\limits_{\Omega} (x + y) \,\mathrm{d}x\,\mathrm{d}y$，其中区域 $\Omega$ 由曲线 $x^2 + y^2 = x + y$ 围成。
\end{prob}

\begin{prob}[S]{3966}
$\displaystyle \iint\limits_{|x| + |y| \leqslant 1} (|x| + |y|) \,\mathrm{d}x\,\mathrm{d}y$。
\end{prob}

\begin{prob}[S]{3967}
$\displaystyle \iint\limits_{\Omega} \sqrt{1 - \frac{x^2}{a^2} - \frac{y^2}{b^2}} \,\mathrm{d}x\,\mathrm{d}y$，其积分域 $\Omega$ 是椭圆区域 $\dfrac{x^2}{a^2} + \dfrac{y^2}{b^2} \leqslant 1$。
\end{prob}

\begin{prob}[S]{3968}
$\displaystyle \iint\limits_{x^4 + y^4 \leqslant 1} (x^2 + y^2) \,\mathrm{d}x\,\mathrm{d}y$。
\end{prob}

\begin{prob}[S]{3969}
$\displaystyle \iint\limits_{\Omega} (x + y) \,\mathrm{d}x\,\mathrm{d}y$，其积分域 $\Omega$ 由曲线 $y^2 = 2x, \, x + y = 4, \, x + y = 12$ 围成。
\end{prob}

\begin{prob}[S]{3970}
$\displaystyle \iint\limits_{\Omega} xy \,\mathrm{d}x\,\mathrm{d}y$，其中 $\Omega$ 是由曲线 $xy = 1, \, x + y = \dfrac{5}{2}$ 围成的区域。
\end{prob}

\begin{prob}[S]{3971}
$\displaystyle \iint\limits_{\substack{0 \leqslant x \leqslant \pi \\ 0 \leqslant y \leqslant \pi}} |\cos(x + y)| \,\mathrm{d}x\,\mathrm{d}y$。
\end{prob}

\begin{prob}[S]{3972}
$\displaystyle \iint\limits_{x^2 + y^2 \leqslant 1} \left| \frac{x + y}{\sqrt{2}} - x^2 - y^2 \right| \,\mathrm{d}x\,\mathrm{d}y$。
\end{prob}

\begin{prob}[S]{3973}
$\displaystyle \iint\limits_{\substack{|x| \leqslant 1 \\ 0 \leqslant y \leqslant 2}} \sqrt{|y - x^2|} \,\mathrm{d}x\,\mathrm{d}y$。
\end{prob}

\probgroup{计算不连续函数的积分：}

\begin{prob}[S]{3974}
$\displaystyle \iint\limits_{x^2 + y^2 \leqslant 4} \operatorname{sgn}(x^2 - y^2 + 2) \,\mathrm{d}x\,\mathrm{d}y$。
\end{prob}

\begin{prob}[S]{3975}
$\displaystyle \iint\limits_{\substack{0 \leqslant x \leqslant 2 \\ 0 \leqslant y \leqslant 2}} [x + y] \,\mathrm{d}x\,\mathrm{d}y$。
\end{prob}

\begin{prob}[S]{3976}
$\displaystyle \iint\limits_{x^2 \leqslant y \leqslant 4} \sqrt{[y - x^2]} \,\mathrm{d}x\,\mathrm{d}y$。
\end{prob}

\begin{prob}[M]{3977}
设 $m$ 及 $n$ 为正整数且其中至少有一个是奇数，证明：
\[
\iint\limits_{x^2 + y^2 \leqslant a^2} x^m y^n \,\mathrm{d}x\,\mathrm{d}y = 0.
\]
\end{prob}

\begin{prob}[M]{3978}
求
\[
\lim_{\rho \to 0} \frac{1}{\pi\rho^2} \iint\limits_{x^2 + y^2 \leqslant \rho^2} f(x, y) \,\mathrm{d}x\,\mathrm{d}y,
\]
其中 $f(x, y)$ 为连续函数。
\end{prob}

\begin{prob}[S]{3979}
设
\[
F(t) = \iint\limits_{\substack{0 \leqslant x \leqslant t \\ 0 \leqslant y \leqslant t}} \mathrm{e}^{\frac{tx}{y^2}} \,\mathrm{d}x\,\mathrm{d}y,
\]
求 $F'(t)$。
\end{prob}

\begin{prob}[S]{3980}
设
\[
F(t) = \iint\limits_{(x - t)^2 + (y - t)^2 \leqslant 1} \sqrt{x^2 + y^2} \,\mathrm{d}x\,\mathrm{d}y,
\]
求 $F'(t)$。
\end{prob}

\begin{prob}[S]{3981}
设
\[
F(t) = \iint\limits_{x^2 + y^2 \leqslant t^2} f(x, y) \,\mathrm{d}x\,\mathrm{d}y \quad (t > 0),
\]\footnote{要求 $f$ 在有关区域上连续。——译注}
求 $F'(t)$。
\end{prob}

\begin{prob}[M]{3982}
证明：若 $f(x, y)$ 连续，则函数
\[
u(x, y) = \frac{1}{2} \int_0^x \,\mathrm{d}\xi \int_{\xi - x + y}^{x + y - \xi} f(\xi, \eta) \,\mathrm{d}\eta
\]
满足方程
\[
\frac{\partial^2 u}{\partial x^2} - \frac{\partial^2 u}{\partial y^2} = f(x, y).
\]
\end{prob}

\begin{prob}[M]{3983}
设函数 $f(x, y)$ 的等值线是简单封闭曲线，区域 $S(v_1, v_2)$ 由曲线 $f(x, y) = v_1$ 及 $f(x, y) = v_2$ 围成。证明：
\[
\iint\limits_{S(v_1, v_2)} f(x, y) \,\mathrm{d}x\,\mathrm{d}y = \int_{v_1}^{v_2} v F'(v) \,\mathrm{d}v,
\]
其中 $F(v)$ 为曲线 $f(x, y) = v_1$ 与 $f(x, y) = v_2$ 所包围的面积。
\probhint{用函数 $f(x, y)$ 的无限接近的等值线把积分域分为许多部分。}
\end{prob}

\probsec{§2. 面积的计算法}

\probgroup{求下列曲线所围的面积：}

\begin{prob}[S]{3984}
$xy = a^2, \quad x + y = \dfrac{5}{2}a \quad (a > 0)$。
\end{prob}

\begin{prob}[S]{3985}
$y^2 = 2px + p^2, \quad y^2 = -2qx + q^2 \quad (p > 0, \, q > 0)$。
\end{prob}

\begin{prob}[S]{3986}
$(x - y)^2 + x^2 = a^2 \quad (a > 0)$。
\end{prob}

\probgroup{变换为极坐标，计算下列曲线所围的面积：}

\begin{prob}[S]{3987}
$(x^2 + y^2)^2 = 2a^2(x^2 - y^2), \quad x^2 + y^2 \geqslant a^2$。
\end{prob}

\begin{prob}[S]{3988}
$(x^3 + y^3)^2 = x^2 + y^2, \quad x \geqslant 0, \, y \geqslant 0$。
\end{prob}

\begin{prob}[S]{3989}
$(x^2 + y^2)^2 = a(x^3 - 3xy^2) \quad (a > 0)$。
\end{prob}

\begin{prob}[S]{3990}
$(x^2 + y^2)^2 = 8a^2 xy, \quad (x - a)^2 + (y - a)^2 \leqslant a^2 \quad (a > 0)$。
\end{prob}

\probgroup{根据以下公式引入广义极坐标 $r$ 和 $\varphi$：
\[
x = a r \cos^\alpha \varphi, \quad y = b r \sin^\alpha \varphi \quad (r \geqslant 0),
\]
其中 $a, b$ 和 $\alpha$ 为以适当的方法选出的常数，且考虑到 $\dfrac{D(x, y)}{D(r, \varphi)} = \alpha a b r \cos^{\alpha - 1} \varphi \sin^{\alpha - 1} \varphi$，求由下列曲线所围的面积（假定参数是正的）：}

\begin{prob}[S]{3991}
$\dfrac{x^2}{a^2} + \dfrac{y^2}{b^2} = \dfrac{x}{h} + \dfrac{y}{k}$。
\end{prob}

\begin{prob}[S]{3992}
$\dfrac{x^3}{a^3} + \dfrac{y^3}{b^3} = \dfrac{x^2}{h^2} + \dfrac{y^2}{k^2}; \quad x = 0, \, y = 0$。
\end{prob}

\begin{prob}[S]{3993}
$\left(\dfrac{x}{a} + \dfrac{y}{b}\right)^4 = \dfrac{x^2}{h^2} + \dfrac{y^2}{k^2} \quad (x > 0, \, y > 0)$。
\end{prob}

\begin{prob}[M]{3994}
(a) $\left(\dfrac{x}{a} + \dfrac{y}{b}\right)^4 = \dfrac{x^2}{h^2} - \dfrac{y^2}{k^2} \quad (x > 0, \, y > 0)$; \qquad (b) $\left(\dfrac{x}{a} + \dfrac{y}{b}\right)^5 = \dfrac{x^2 y^2}{c^4}$。
\end{prob}

\begin{prob}[S]{3995}
$\sqrt[4]{\dfrac{x}{a}} + \sqrt[4]{\dfrac{y}{b}} = 1; \quad x = 0, \, y = 0$。
\end{prob}

\probgroup{进行适当的变量代换，求下列曲线所围的面积：}

\begin{prob}[S]{3996}
$x + y = a, \quad x + y = b, \quad y = \alpha x, \quad y = \beta x \quad (0 < a < b, \, 0 < \alpha < \beta)$。
\end{prob}

\begin{prob}[S]{3997}
$xy = a^2, \quad xy = 2a^2, \quad y = x, \quad y = 2x \quad (x > 0, \, y > 0)$。
\end{prob}

\begin{prob}[L]{3998}
\begin{enumerate}[label=(\alph*), itemsep=2pt, topsep=3pt]
\item $y^2 = 2px, \quad y^2 = 2qx, \quad x^2 = 2ry, \quad x^2 = 2sy \quad (0 < p < q, \, 0 < r < s)$;
\item $x^2 = ay, \quad x^2 = by, \quad x^3 = cy^2, \quad x^3 = dy^2 \quad (0 < a < b, \, 0 < c < d)$;
\item $y = ax^p, \quad y = bx^p, \quad y = cx^q, \quad y = dx^q \quad (0 < p < q, \, 0 < a < b, \, 0 < c < d)$。
\end{enumerate}
\end{prob}

\begin{prob}[M]{3999}
\begin{enumerate}[label=(\alph*), itemsep=2pt, topsep=3pt]
\item $\sqrt{\dfrac{x}{a}} + \sqrt{\dfrac{y}{b}} = 1, \quad \sqrt{\dfrac{x}{a}} + \sqrt{\dfrac{y}{b}} = 2, \quad \dfrac{x}{a} = \dfrac{y}{b}, \quad \dfrac{4x}{a} = \dfrac{y}{b} \quad (a > 0, \, b > 0)$;
\item $\left(\dfrac{x}{a}\right)^{\frac{2}{3}} + \left(\dfrac{y}{b}\right)^{\frac{2}{3}} = 1, \quad \left(\dfrac{x}{a}\right)^{\frac{2}{3}} + \left(\dfrac{y}{b}\right)^{\frac{2}{3}} = 4, \quad \dfrac{x}{a} = \dfrac{y}{b}, \quad \dfrac{8x}{a} = \dfrac{y}{b} \quad (x > 0, \, y > 0)$。
\end{enumerate}
\end{prob}

\begin{prob}[M]{4000}
$\dfrac{x^2}{\lambda} + \dfrac{y^2}{\lambda - c^2} = 1$，其中 $\lambda$ 取下列各值：$\dfrac{1}{3}c^2, \, \dfrac{2}{3}c^2, \, \dfrac{4}{3}c^2, \, \dfrac{5}{3}c^2 \quad (x > 0, \, y > 0)$。
\end{prob}

\begin{prob}[M]{4001}
求椭圆 $(a_1 x + b_1 y + c_1)^2 + (a_2 x + b_2 y + c_2)^2 = 1$（其中 $\delta = a_1 b_2 - a_2 b_1 \neq 0$）的面积。
\end{prob}

\begin{prob}[M]{4002}
求椭圆
\[
\frac{x^2}{\cosh^2 u} + \frac{y^2}{\sinh^2 u} = c^2 \quad (u = u_1, u_2)
\]
和双曲线
\[
\frac{x^2}{\cos^2 v} - \frac{y^2}{\sin^2 v} = c^2 \quad (v = v_1, v_2)
\]
\[
(0 < u_1 < u_2; \quad 0 < v_1 < v_2; \quad x > 0, \, y > 0)
\]
所围的面积。
\probhint{令 $x = c \cosh u \cos v, \, y = c \sinh u \sin v$。}
\end{prob}

\begin{prob}[M]{4003}
求平面 $x + y + z = b$ 与曲面 $x^2 + y^2 + z^2 - xy - xz - yz = a^2$ 相截所得截断面之面积。
\end{prob}

\begin{prob}[M]{4004}
求平面 $z = 1 - 2(x + y)$ 与曲面 $\dfrac{1}{x} + \dfrac{1}{y} + \dfrac{1}{z} = 0$ 相截所得截断面之面积。
\end{prob}

\probsec{§3. 体积的计算法}

\begin{prob}[M]{4005}
试绘出一物体，其体积等于积分
\[
V = \int_0^1 \,\mathrm{d}x \int_0^{1 - x} (x^2 + y^2) \,\mathrm{d}y.
\]
\end{prob}

\begin{prob}[L]{4006}
绘出下列二重积分所表示的体积：
\begin{enumerate}[label=(\alph*), itemsep=2pt, topsep=3pt]
\item $\displaystyle \iint\limits_{\substack{0 \leqslant x + y \leqslant 1 \\ x \geqslant 0, y \geqslant 0}} (x + y) \,\mathrm{d}x\,\mathrm{d}y$;
\item $\displaystyle \iint\limits_{\frac{x^2}{4} + \frac{y^2}{9} \leqslant 1} \sqrt{1 - \frac{x^2}{4} - \frac{y^2}{9}} \,\mathrm{d}x\,\mathrm{d}y$;
\item $\displaystyle \iint\limits_{|x| + |y| \leqslant 1} (x^2 + y^2) \,\mathrm{d}x\,\mathrm{d}y$;
\item $\displaystyle \iint\limits_{x^2 + y^2 \leqslant x} \sqrt{x^2 + y^2} \,\mathrm{d}x\,\mathrm{d}y$;
\item $\displaystyle \iint\limits_{\substack{1 \leqslant x \leqslant 2 \\ x \leqslant y \leqslant 2x}} \sqrt{xy} \,\mathrm{d}x\,\mathrm{d}y$;
\item $\displaystyle \iint\limits_{x^2 + y^2 \leqslant 1} \sin \pi \sqrt{x^2 + y^2} \,\mathrm{d}x\,\mathrm{d}y$。
\end{enumerate}
\end{prob}

\probgroup{求下列曲面所围区域的体积：}

\begin{prob}[S]{4007}
$z = 1 + x + y, \quad z = 0, \quad x + y = 1, \quad x = 0, \quad y = 0$。
\end{prob}

\begin{prob}[S]{4008}
$x + y + z = a, \quad x^2 + y^2 = R^2, \quad x = 0, \quad y = 0, \quad z = 0 \quad (a \geqslant \sqrt{2}R)$。
\end{prob}

\begin{prob}[S]{4009}
$z = x^2 + y^2, \quad y = x^2, \quad y = 1, \quad z = 0$。
\end{prob}

\begin{prob}[S]{4010}
$z = \cos x \cos y, \quad z = 0, \quad |x + y| \leqslant \dfrac{\pi}{2}, \quad |x - y| \leqslant \dfrac{\pi}{2}$。
\end{prob}

\begin{prob}[S]{4011}
$z = \sin \dfrac{\pi y}{2x}, \quad z = 0, \quad y = x, \quad y = 0, \quad x = \pi$。
\end{prob}

\begin{prob}[S]{4012}
$z = xy, \quad x + y + z = 1, \quad z = 0$。
\end{prob}

\probgroup{变换成极坐标，求下列曲面所围区域的体积：}

\begin{prob}[S]{4013}
$z^2 = xy, \quad x^2 + y^2 = a^2$。
\end{prob}

\begin{prob}[S]{4014}
$z = x + y, \quad (x^2 + y^2)^2 = 2xy, \quad z = 0 \quad (x > 0, \, y > 0)$。
\end{prob}

\begin{prob}[S]{4015}
$z = x^2 + y^2, \quad x^2 + y^2 = x, \quad x^2 + y^2 = 2x, \quad z = 0$。
\end{prob}

\begin{prob}[S]{4016}
$x^2 + y^2 + z^2 = a^2, \quad x^2 + y^2 \geqslant a|x| \quad (a > 0)$。
\end{prob}

\begin{prob}[S]{4017}
$x^2 + y^2 - az = 0, \quad (x^2 + y^2)^2 = a^2(x^2 - y^2), \quad z = 0 \quad (a > 0)$。
\end{prob}

\begin{prob}[S]{4018}
$z = \mathrm{e}^{-(x^2 + y^2)}, \quad z = 0, \quad x^2 + y^2 = R^2$。
\end{prob}

\begin{prob}[S]{4019}
$z = c\cos\dfrac{\pi\sqrt{x^2 + y^2}}{2a}, \quad x^2 + y^2 = a^2, \quad y = x\tan\alpha, \quad y = x\tan\beta \quad (a > 0, \, c > 0, \, 0 \leqslant \alpha < \beta \leqslant 2\pi)$。
\end{prob}

\begin{prob}[S]{4020}
$z = x^2 + y^2, \quad z = x + y$。
\end{prob}

\probgroup{求下列曲面所围区域的体积（假定参数是正的）：}

\begin{prob}[S]{4021}
$\dfrac{x^2}{a^2} + \dfrac{y^2}{b^2} + \dfrac{z^2}{c^2} = 1, \quad \dfrac{x^2}{a^2} + \dfrac{y^2}{b^2} = \dfrac{z^2}{c^2} \quad (z > 0)$。
\end{prob}

\begin{prob}[S]{4022}
$\dfrac{x^2}{a^2} + \dfrac{y^2}{b^2} - \dfrac{z^2}{c^2} = -1, \quad \dfrac{x^2}{a^2} + \dfrac{y^2}{b^2} = 1$。
\end{prob}

\begin{prob}[S]{4023}
$\dfrac{x^2}{a^2} + \dfrac{y^2}{b^2} = \dfrac{z}{c}, \quad \dfrac{x^2}{a^2} + \dfrac{y^2}{b^2} = \dfrac{x}{a} + \dfrac{y}{b}, \quad z = 0$。
\end{prob}

\begin{prob}[S]{4024}
$\left(\dfrac{x^2}{a^2} + \dfrac{y^2}{b^2}\right)^2 + \dfrac{z}{c} = 1, \quad z = 0$。
\end{prob}

\begin{prob}[S]{4025}
$\left(\dfrac{x}{a} + \dfrac{y}{b}\right)^2 + \dfrac{z^2}{c^2} = 1, \quad x = 0, \quad y = 0, \quad z = 0$。
\end{prob}

\begin{prob}[S]{4026}
$\dfrac{x^2}{a^2} + \dfrac{y^2}{b^2} + \dfrac{z^2}{c^2} = 1, \quad \left(\dfrac{x^2}{a^2} + \dfrac{y^2}{b^2}\right)^2 = \dfrac{x^2}{a^2} - \dfrac{y^2}{b^2}$。
\end{prob}

\begin{prob}[S]{4027}
$z^2 = xy, \quad x + y = a, \quad x + y = b \quad (0 < a < b)$。
\end{prob}

\begin{prob}[S]{4028}
$z = x^2 + y^2, \quad xy = a^2, \quad xy = 2a^2, \quad y = \dfrac{x}{2}, \quad y = 2x, \quad z = 0$。
\end{prob}

\begin{prob}[S]{4029}
$z = xy, \quad x^2 = y, \quad x^2 = 2y, \quad y^2 = x, \quad y^2 = 2x, \quad z = 0$。
\end{prob}

\begin{prob}[S]{4030}
$z = c\sin\dfrac{\pi xy}{a^2}, \quad z = 0, \quad xy = a^2, \quad y = \alpha x, \quad y = \beta x \quad (0 < \alpha < \beta, \, x > 0)$。
\end{prob}

\begin{prob}[S]{4031}
$z = x^{\frac{3}{2}} + y^{\frac{3}{2}}, \quad z = 0, \quad x + y = 1, \quad x = 0, \quad y = 0$。
\end{prob}

\begin{prob}[S]{4032}
$\dfrac{x^2}{a^2} + \dfrac{y^2}{b^2} + \dfrac{z}{c} = 1$, $\left(\dfrac{x}{a}\right)^{\frac{2}{3}} + \left(\dfrac{y}{b}\right)^{\frac{2}{3}} = 1$, $z = 0$.
\end{prob}

\begin{prob}[M]{4033}
(a) $z = c\arctan\dfrac{y}{x}$, $z = 0$, $\sqrt{x^2 + y^2} = a\arctan\dfrac{y}{x}$ ($y \geqslant 0$); \qquad
(b) $z = y\mathrm{e}^{-\frac{xy}{a^2}}$, $xy = a^2$, $xy = 2a^2$, $y = m$, $y = n$, $z = 0$ ($0 < m < n$).
\end{prob}

\begin{prob}[S]{4034}
$\dfrac{x^n}{a^n} + \dfrac{y^n}{b^n} + \dfrac{z^n}{c^n} = 1$, $x = 0$, $y = 0$, $z = 0$ ($n > 0$).
\end{prob}

\begin{prob}[S]{4035}
$\left(\dfrac{x}{a} + \dfrac{y}{b}\right)^n + \left(\dfrac{z}{c}\right)^m = 1$, $x = 0$, $y = 0$, $z = 0$ ($n > 0, m > 0$).
\end{prob}

\probsec{§4. 曲面面积的计算}

\begin{prob}[M]{4036}
求曲面 $az = xy$ 包含在圆柱 $x^2 + y^2 = a^2$ 内那部分的面积.
\end{prob}

\begin{prob}[M]{4037}
求以曲面 $x^2 + z^2 = a^2$ 和 $y^2 + z^2 = a^2$ 为界的物体的表面积.
\end{prob}

\begin{prob}[M]{4038}
求球面 $x^2 + y^2 + z^2 = a^2$ 包含在柱面 $\dfrac{x^2}{a^2} + \dfrac{y^2}{b^2} = 1$ ($b \leqslant a$) 内那部分的面积.
\end{prob}

\begin{prob}[M]{4039}
求曲面 $z^2 = 2xy$ 被平面 $x + y = 1$, $x = 0$, $y = 0$ 所截部分的面积.
\end{prob}

\begin{prob}[M]{4040}
求曲面 $x^2 + y^2 + z^2 = a^2$ 在圆柱 $x^2 + y^2 = \pm ax$ 内那部分的面积（维维亚尼问题）.
\end{prob}

\begin{prob}[M]{4041}
求曲面 $z = \sqrt{x^2 + y^2}$ 包含在圆柱 $x^2 + y^2 = 2x$ 内那部分的面积.
\end{prob}

\begin{prob}[M]{4042}
求曲面 $z = \sqrt{x^2 - y^2}$ 包含在柱面 $(x^2 + y^2)^2 = a^2(x^2 - y^2)$ 内那部分的面积.
\end{prob}

\begin{prob}[M]{4043}
求曲面 $z = \dfrac{1}{2}(x^2 - y^2)$ 被平面 $x - y = \pm 1$, $x + y = \pm 1$ 所截部分的面积.
\end{prob}

\begin{prob}[M]{4044}
求曲面 $x^2 + y^2 = 2az$ 包含在柱面 $(x^2 + y^2)^2 = 2a^2 xy$ 内那部分的面积.
\end{prob}

\begin{prob}[L]{4045}
\begin{enumerate}[label=(\alph*), itemsep=2pt, topsep=3pt]
  \item 求曲面 $x^2 + y^2 = a^2$ 被平面 $x + z = 0$, $x - z = 0$ ($x > 0, y > 0$) 所截部分的面积;
  \item 求曲面 $(x^2 + y^2)^{\frac{3}{2}} + z = 1$ 被平面 $z = 0$ 所截部分的面积;
  \item 求曲面 $\left(\dfrac{x}{a} + \dfrac{y}{b}\right)^2 + \dfrac{2z}{c} = 1$ 被平面 $x = 0$, $y = 0$ 和 $z = 0$ 所截部分的面积;
  \item 求曲面 $\dfrac{x^2}{a} - \dfrac{y^2}{b} = 2z$ 被曲面 $\dfrac{x^2}{a^2} + \dfrac{x^2}{b^2} = 1$ ($z \geqslant 0$) 所截部分的面积;
  \item 求曲面 $\sin z = \sinh x \cdot \sinh y$ 被平面 $x = 1$ 和 $x = 2$ ($y \geqslant 0$) 所截部分的面积.
\end{enumerate}
\end{prob}

\begin{prob}[M]{4046}
求以曲面 $x^2 + y^2 = \dfrac{1}{3}z^2$ 和 $x + y + z = 2a$ ($a > 0$) 为界的物体的表面积和体积.
\end{prob}

\begin{prob}[M]{4047}
求球面在两条纬线和两条经线之间那部分的面积.
\end{prob}

\begin{prob}[M]{4048}
求螺旋面
\[
x = r\cos\varphi, \quad y = r\sin\varphi, \quad z = h\varphi
\]
在 $0 < r < a$, $0 < \varphi < 2\pi$ 那部分的面积.
\end{prob}

\begin{prob}[L]{4049}
求环面
\[
x = (b + a\cos\psi)\cos\varphi, \quad y = (b + a\cos\psi)\sin\varphi, \quad z = a\sin\psi \quad (0 < a \leqslant b)
\]
在两条经线 $\varphi = \varphi_1$, $\varphi = \varphi_2$ 和两条纬线 $\psi = \psi_1$, $\psi = \psi_2$ 之间那部分的面积. 整个环的表面积等于什么?
\end{prob}

\begin{prob}[M]{4050}
求矩形 $x = a > 0$, $0 \leqslant y \leqslant b$, $0 \leqslant z \leqslant c$ 对坐标原点的立体角 $\omega$. 若 $a$ 很大, 推出 $\omega$ 的近似公式.
\end{prob}

\probsec{§5. 二重积分在力学上的应用}

\begin{prob}[M]{4051}
求边长为 $a$ 的正方形薄板的质量, 设薄板上每一点的面密度与该点到正方形顶点之一的距离成正比, 且在正方形中心等于 $\rho_0$.
\end{prob}

\probgroup{求以下列曲线为界的均匀薄板的质心坐标：}

\begin{prob}[S]{4052}
$ay = x^2$, $x + y = 2a$ ($a > 0$).
\end{prob}

\begin{prob}[S]{4053}
$\sqrt{x} + \sqrt{y} = \sqrt{a}$, $x = 0$, $y = 0$.
\end{prob}

\begin{prob}[S]{4054}
$x^{\frac{2}{3}} + y^{\frac{2}{3}} = a^{\frac{2}{3}}$ ($x > 0, y > 0$).
\end{prob}

\begin{prob}[S]{4055}
$\left(\dfrac{x}{a} + \dfrac{y}{b}\right)^3 = \dfrac{xy}{c^2}$（线圈）.
\end{prob}

\begin{prob}[S]{4056}
$(x^2 + y^2)^2 = 2a^2 xy$ ($x > 0, y > 0$).
\end{prob}

\begin{prob}[S]{4057}
$r = a(1 + \cos\varphi)$, $\varphi = 0$.
\end{prob}

\begin{prob}[S]{4058}
$x = a(t - \sin t)$, $y = a(1 - \cos t)$ ($0 \leqslant t \leqslant 2\pi$), $y = 0$.
\end{prob}

\begin{prob}[M]{4059}
求圆形薄板 $x^2 + y^2 \leqslant a^2$ 的质心坐标. 设它在点 $M(x,y)$ 的面密度与点 $M$ 到点 $A(a,0)$ 的距离成正比.
\end{prob}

\begin{prob}[M]{4060}
求曲线
\[
y = \sqrt{2px}, \quad y = 0, \quad x = X
\]
所围图形的质心在参数 $X$ 变化时所描绘的曲线.
\end{prob}

\probgroup{求由下列曲线所围的面积（$\rho = 1$）对坐标轴 $Ox$ 和 $Oy$ 的转动惯量 $I_x$ 和 $I_y$：}

\begin{prob}[S]{4061}
$\dfrac{x}{b_1} + \dfrac{y}{h} = 1$, $\dfrac{x}{b_2} + \dfrac{y}{h} = 1$, $y = 0$ ($b_1 > 0, b_2 > 0, h > 0$).
\end{prob}

\begin{prob}[S]{4062}
$(x - a)^2 + (y - a)^2 = a^2$, $x = 0$, $y = 0$ ($0 \leqslant x \leqslant a$).
\end{prob}

\begin{prob}[S]{4063}
$r = a(1 + \cos\varphi)$.
\end{prob}

\begin{prob}[S]{4064}
$x^4 + y^4 = a^2(x^2 + y^2)$.
\end{prob}

\begin{prob}[S]{4065}
$xy = a^2$, $xy = 2a^2$, $x = 2y$, $2x = y$ ($x > 0, y > 0$).
\end{prob}

\begin{prob}[M]{4066.1}
求曲线
\[
(x^2 + y^2)^2 = a^2(x^2 - y^2)
\]
所围区域 $S$ 的极转动惯量
\[
I_0 = \iint\limits_S (x^2 + y^2)\,\mathrm{d}x\,\mathrm{d}y.
\]
\end{prob}

\begin{prob}[M]{4066.2}
求曲线
\[
ay = x^2, \quad ax = y^2 \quad (a > 0)
\]
所围平面图形的惯性积 $I_{xy}$.
\end{prob}

\begin{prob}[M]{4067}
证明公式:
\[
I_l = I_{l_0} + Sd^2,
\]
式中 $I_l$, $I_{l_0}$ 是图形 $S$ 对二平行轴 $l$ 和 $l_0$ 的转动惯量, 其中 $l_0$ 是通过图形的质心, 而 $d$ 为两轴间的距离.
\end{prob}

\begin{prob}[M]{4068}
证明: 平面图形 $S$ 对通过其质心 $O(0,0)$ 并与 $Ox$ 轴成 $\alpha$ 角的直线的转动惯量等于
\[
I = I_x\cos^2\alpha - 2I_{xy}\sin\alpha\cos\alpha + I_y\sin^2\alpha,
\]
其中 $I_x$ 和 $I_y$ 为图形 $S$ 对 $Ox$ 轴和 $Oy$ 轴的转动惯量, $I_{xy}$ 为惯性积:
\[
I_{xy} = \iint\limits_S \rho xy\,\mathrm{d}x\,\mathrm{d}y.
\]
\end{prob}

\begin{prob}[M]{4069}
求以 $a$ 为边的正三角形对通过三角形质心并与它的高成 $\alpha$ 角的直线的转动惯量.
\end{prob}

\begin{prob}[M]{4070}
设圆柱形容器 $x^2 + y^2 = a^2$, $z = 0$ 内盛有水, 水平面为 $z = h$, 求水对容器侧壁 $x \geqslant 0$ 的压力.\footnote{关于习题 4070\textasciitilde 4072 可参阅 166 页的脚注. —— 译注}
\end{prob}

\begin{prob}[M]{4071}
半径为 $a$ 的球体沉入密度为 $\delta$ 的液体中深度为 $h$（由球心算起）的地方, 这里 $h \geqslant a$. 求液体对球的上表面和下表面的压力.
\end{prob}

\begin{prob}[M]{4072}
底半径为 $a$ 高为 $b$ 的直圆柱体完全沉入密度为 $\delta$ 的液体中, 其中心在水面下的深度为 $h$, 而圆柱的轴与竖直方向成 $\alpha$ 角. 求液体对圆柱上底和下底的压力.
\end{prob}

\begin{prob}[M]{4073}
求均匀的圆柱体 $x^2 + y^2 \leqslant a^2$, $0 \leqslant z \leqslant h$ 对质点 $P(0,0,b)$ 的引力, 设圆柱的质量等于 $M$, 而质点的质量等于 $m$.
\end{prob}

\begin{prob}[M]{4074}
物体对挤压面
\[
\frac{x^2}{a^2} + \frac{y^2}{b^2} \leqslant 1
\]
的压强分布由公式
\[
p = p_0\left(1 - \frac{x^2}{a^2} - \frac{y^2}{b^2}\right)
\]
给出. 求物体对此面的平均压强.
\end{prob}

\begin{prob}[M]{4075}
以 $a$ 和 $b$ 为边的矩形草地上均匀地覆盖有已收割的干草, 其面密度为 $\rho$. 若运送质量为 $M$ 的货物到距离为 $r$ 的地方所需的功等于 $kMr$ ($0 < k < 1$)\footnote{原书如此，其实系数 $k$ 是有量纲的，它的值也不应限制在 $0$ 与 $1$ 之间. —— 译注}, 则为了把所有干草集中在草地的中心, 至少应消耗多少功?
\end{prob}

\probsec{§6. 三重积分}

\probgroup{计算下列三重积分：}

\begin{prob}[M]{4076}
$\displaystyle\iiint\limits_V xy^2 z^3\,\mathrm{d}x\,\mathrm{d}y\,\mathrm{d}z$, 其中 $V$ 是曲面 $z = xy$, $y = x$, $x = 1$, $z = 0$ 所围的区域.
\end{prob}

\begin{prob}[M]{4077}
$\displaystyle\iiint\limits_V \frac{\mathrm{d}x\,\mathrm{d}y\,\mathrm{d}z}{(1 + x + y + z)^3}$, 其中 $V$ 是曲面 $x + y + z = 1$, $x = 0$, $y = 0$, $z = 0$ 所围的区域.
\end{prob}

\begin{prob}[M]{4078}
$\displaystyle\iiint\limits_V xyz\,\mathrm{d}x\,\mathrm{d}y\,\mathrm{d}z$, 其中 $V$ 是曲面 $x^2 + y^2 + z^2 = 1$, $x = 0$, $y = 0$, $z = 0$ 所围的区域.
\end{prob}

\begin{prob}[M]{4079}
$\displaystyle\iiint\limits_V \left(\frac{x^2}{a^2} + \frac{y^2}{b^2} + \frac{z^2}{c^2}\right)\mathrm{d}x\,\mathrm{d}y\,\mathrm{d}z$, 其中 $V$ 是曲面 $\dfrac{x^2}{a^2} + \dfrac{y^2}{b^2} + \dfrac{z^2}{c^2} = 1$ 所围的区域.
\end{prob}

\begin{prob}[M]{4080}
$\displaystyle\iiint\limits_V \sqrt{x^2 + y^2}\,\mathrm{d}x\,\mathrm{d}y\,\mathrm{d}z$, 其中 $V$ 是曲面 $x^2 + y^2 = z^2$, $z = 1$ 所围的区域.
\end{prob}

\probgroup{在下列三重积分内，用不同方法配置积分的上下限：}

\begin{prob}[S]{4081}
$\displaystyle\int_0^1 \mathrm{d}x \int_0^{1-x} \mathrm{d}y \int_0^{x+y} f(x,y,z)\,\mathrm{d}z$.
\end{prob}

\begin{prob}[S]{4082}
$\displaystyle\int_{-1}^1 \mathrm{d}x \int_{-\sqrt{1-x^2}}^{\sqrt{1-x^2}} \mathrm{d}y \int_{\sqrt{x^2+y^2}}^1 f(x,y,z)\,\mathrm{d}z$.
\end{prob}

\begin{prob}[S]{4083}
$\displaystyle\int_0^1 \mathrm{d}x \int_0^1 \mathrm{d}y \int_0^{x^2+y^2} f(x,y,z)\,\mathrm{d}z$.
\end{prob}

\probgroup{用一重积分代替三重积分：}

\begin{prob}[S]{4084}
$\displaystyle \int_{0}^{x} \,\mathrm{d}\xi \int_{0}^{\xi} \,\mathrm{d}\eta \int_{0}^{\eta} f(\zeta)\,\mathrm{d}\zeta$.
\end{prob}

\begin{prob}[S]{4085}
$\displaystyle \int_{0}^{1} \,\mathrm{d}x \int_{0}^{1} \,\mathrm{d}y \int_{0}^{x+y} f(z)\,\mathrm{d}z$.
\end{prob}

\begin{prob}[M]{4086}
设 $f(x, y, z) = F'''_{xyz}(x, y, z)$，且 $a, b, c, A, B, C$ 为常数，求
\[
\int_{a}^{A} \,\mathrm{d}x \int_{b}^{B} \,\mathrm{d}y \int_{c}^{C} f(x, y, z)\,\mathrm{d}z.
\]
\end{prob}

\probgroup{变换为球坐标，计算积分：}

\begin{prob}[M]{4087}
$\displaystyle \iiint_{V} \sqrt{x^2+y^2+z^2} \,\mathrm{d}x\,\mathrm{d}y\,\mathrm{d}z$，其中 $V$ 是曲面 $x^2+y^2+z^2=z$ 所围的区域.
\end{prob}

\begin{prob}[M]{4088}
$\displaystyle \int_{0}^{1} \,\mathrm{d}x \int_{0}^{\sqrt{1-x^2}} \,\mathrm{d}y \int_{\sqrt{x^2+y^2}}^{\sqrt{2-x^2-y^2}} z^2 \,\mathrm{d}z$.
\end{prob}

\begin{prob}[M]{4089}
在积分中变换为球坐标：
\[
\iiint_{V} f\left(\sqrt{x^2+y^2+z^2}\right) \,\mathrm{d}x\,\mathrm{d}y\,\mathrm{d}z,
\]
其中 $V$ 是曲面 $z=x^2+y^2,\ x=y,\ x=1,\ y=0,\ z=0$ 所围的区域.
\end{prob}

\begin{prob}[M]{4090}
进行适当的变量代换，计算三重积分
\[
\iiint_{V} \sqrt{1 - \frac{x^2}{a^2} - \frac{y^2}{b^2} - \frac{z^2}{c^2}} \,\mathrm{d}x\,\mathrm{d}y\,\mathrm{d}z,
\]
其中 $V$ 为椭球 $\dfrac{x^2}{a^2}+\dfrac{y^2}{b^2}+\dfrac{z^2}{c^2}=1$.
\end{prob}

\begin{prob}[M]{4091}
变换为圆柱坐标，计算积分
\[
\iiint_{V} (x^2+y^2) \,\mathrm{d}x\,\mathrm{d}y\,\mathrm{d}z,
\]
其中 $V$ 是曲面 $x^2+y^2=2z,\ z=2$ 所围的区域.
\end{prob}

\begin{prob}[M]{4092}
计算积分
\[
\iiint_{V} x^2 \,\mathrm{d}x\,\mathrm{d}y\,\mathrm{d}z,
\]
其中 $V$ 是曲面 $z=ay^2,\ z=by^2,\ y>0\ (0<a<b),\ z=\alpha x,\ z=\beta x\ (0<\alpha<\beta),\ z=h\ (h>0)$ 所围的区域.
\end{prob}

\begin{prob}[M]{4093}
求积分
\[
\iiint_{V} xyz \,\mathrm{d}x\,\mathrm{d}y\,\mathrm{d}z,
\]
其中区域 $V$ 位于象限 $x>0, y>0, z>0$ 且由下列曲面围成：
\[
z=\frac{x^2+y^2}{m},\quad z=\frac{x^2+y^2}{n},\quad xy=a^2,\quad xy=b^2,\quad y=\alpha x,\quad y=\beta x
\]
\[
(0<a<b;\ 0<\alpha<\beta;\ 0<m<n).
\]
\end{prob}

\begin{prob}[M]{4094}
求函数 $f(x, y, z) = x^2+y^2+z^2$ 在区域 $x^2+y^2+z^2 \leqslant x+y+z$ 内的平均值.
\end{prob}

\begin{prob}[M]{4095}
求函数
\[
f(x, y, z) = \mathrm{e}^{\sqrt{\frac{x^2}{a^2}+\frac{y^2}{b^2}+\frac{z^2}{c^2}}}
\]
在区域 $\dfrac{x^2}{a^2}+\dfrac{y^2}{b^2}+\dfrac{z^2}{c^2} \leqslant 1$ 内的平均值.
\end{prob}

\begin{prob}[M]{4096}
利用中值定理，估计积分
\[
u = \iiint_{x^2+y^2+z^2 \leqslant R^2} \frac{\mathrm{d}x\,\mathrm{d}y\,\mathrm{d}z}{\sqrt{(x-a)^2+(y-b)^2+(z-c)^2}},
\]
其中 $a^2+b^2+c^2 > R^2$.
\end{prob}

\begin{prob}[M]{4097}
证明：若函数 $f(x, y, z)$ 在区域 $V$ 内是连续的，且对于任何区域 $\omega \subset V$ 有
\[
\iiint_{\omega} f(x, y, z) \,\mathrm{d}x\,\mathrm{d}y\,\mathrm{d}z = 0,
\]
则当 $(x, y, z) \in V$ 时 $f(x, y, z) \equiv 0$.
\end{prob}

\begin{prob}[M]{4098}
求 $F'(t)$，设：
\begin{enumerate}[label=(\alph*), itemsep=2pt, topsep=3pt]
\item $\displaystyle F(t) = \iiint_{x^2+y^2+z^2 \leqslant t^2} f(x^2+y^2+z^2) \,\mathrm{d}x\,\mathrm{d}y\,\mathrm{d}z$，其中 $f$ 为可微函数；
\item $\displaystyle F(t) = \iiint_{\substack{0 \leqslant x \leqslant t \\ 0 \leqslant y \leqslant t \\ 0 \leqslant z \leqslant t}} f(xyz) \,\mathrm{d}x\,\mathrm{d}y\,\mathrm{d}z$，其中 $f$ 为可微函数.
\end{enumerate}
\end{prob}

\begin{prob}[M]{4099}
求
\[
\iiint_{x^2+y^2+z^2 \leqslant 1} x^m y^n z^p \,\mathrm{d}x\,\mathrm{d}y\,\mathrm{d}z,
\]
其中 $m, n, p$ 为非负整数.
\end{prob}

\begin{prob}[M]{4100}
令 $x+y+z=\xi,\ y+z=\xi\eta,\ z=\xi\eta\zeta$，计算狄利克雷积分
\[
\iiint_{V} x^p y^q z^r (1-x-y-z)^s \,\mathrm{d}x\,\mathrm{d}y\,\mathrm{d}z \quad (p>0, q>0, r>0, s>0),
\]
其中 $V$ 是平面 $x+y+z=1,\ x=0,\ y=0,\ z=0$ 所围的区域.
\end{prob}

\probsec{§7. 利用三重积分计算体积}

\probgroup{求以下列曲面为界的物体的体积：}

\begin{prob}[S]{4101}
$z=x^2+y^2,\ z=2x^2+2y^2,\ y=x,\ y=x^2$.
\end{prob}

\begin{prob}[S]{4102}
$z=x+y,\ z=xy,\ x+y=1,\ x=0,\ y=0$.
\end{prob}

\begin{prob}[S]{4103}
$x^2+z^2=a^2,\ x+y=\pm a,\ x-y=\pm a$.
\end{prob}

\begin{prob}[S]{4104}
$az=x^2+y^2,\ z=\sqrt{x^2+y^2}\ (a>0)$.
\end{prob}

\begin{prob}[S]{4105}
$az=a^2-x^2-y^2,\ z=a-x-y,\ x=0,\ y=0,\ z=0\ (a>0)$.
\end{prob}

\begin{prob}[S]{4106}
$z=6-x^2-y^2,\ z=\sqrt{x^2+y^2}$.
\end{prob}

\probgroup{变换为球坐标或圆柱坐标，计算以下曲面所围的体积：}

\begin{prob}[S]{4107}
$x^2+y^2+z^2=2az,\ x^2+y^2 \leqslant z^2$.
\end{prob}

\begin{prob}[S]{4108}
$(x^2+y^2+z^2)^2 = a^2(x^2+y^2-z^2)$.
\end{prob}

\begin{prob}[S]{4109}
$(x^2+y^2+z^2)^3 = 3xyz$.
\end{prob}

\begin{prob}[S]{4110}
$x^2+y^2+z^2=a^2,\ x^2+y^2+z^2=b^2,\ x^2+y^2=z^2\ (z \geqslant 0)\ (0<a<b)$.
\end{prob}

\probgroup{在下列各题中最好利用广义球坐标
\[
r, \varphi, \psi \quad \left(r \geqslant 0;\ 0 \leqslant \varphi \leqslant 2\pi;\ -\frac{\pi}{2} \leqslant \psi \leqslant \frac{\pi}{2}\right),
\]
它们由以下公式引入：
\[
\begin{cases}
x = a r \cos^{\alpha}\varphi \cos^{\beta}\psi, \\
y = b r \sin^{\alpha}\varphi \cos^{\beta}\psi, \\
z = c r \sin^{\beta}\psi
\end{cases}
\]
($a, b, c, \alpha, \beta$ 为常数)，并且
\[
\frac{D(x, y, z)}{D(r, \varphi, \psi)} = \alpha\beta a b c r^2 \cos^{\alpha-1}\varphi \sin^{\alpha-1}\varphi \cos^{2\beta-1}\psi \sin^{\beta-1}\psi.
\]
求下列曲面为界的物体的体积：}

\begin{prob}[S]{4111}
$\displaystyle \left(\frac{x^2}{a^2}+\frac{y^2}{b^2}+\frac{z^2}{c^2}\right)^2 = \frac{x}{h}$.
\end{prob}

\begin{prob}[M]{4112}
(a) $\displaystyle \left(\frac{x^2}{a^2}+\frac{y^2}{b^2}+\frac{z^2}{c^2}\right)^2 = \frac{x^2}{a^2}+\frac{y^2}{b^2}$; \qquad
(b) $\displaystyle \left(\frac{x^2}{a^2}+\frac{y^2}{b^2}+\frac{z^2}{c^2}\right)^2 = \frac{x^2}{a^2}+\frac{y^2}{b^2}-\frac{z^2}{c^2}$.
\end{prob}

\begin{prob}[S]{4113}
$\displaystyle \frac{x^2}{a^2}+\frac{y^2}{b^2}+\frac{z^2}{c^2}=1,\ \frac{x^2}{a^2}+\frac{y^2}{b^2}=\frac{z}{c}$.
\end{prob}

\begin{prob}[S]{4114}
$\displaystyle \frac{x^2}{a^2}+\frac{y^2}{b^2}+\frac{z^4}{c^4}=1$.
\end{prob}

\begin{prob}[S]{4115}
$\displaystyle \left(\frac{x^2}{a^2}+\frac{y^2}{b^2}\right)^2+\frac{z^4}{c^4}=1$.
\end{prob}

\probgroup{利用适当的变量代换，计算以下列曲面为界的物体的体积（假定参数是正的）：}

\begin{prob}[M]{4116}
\begin{enumerate}[label=(\alph*), itemsep=2pt, topsep=3pt]
\item $\displaystyle \left(\frac{x}{a}+\frac{y}{b}+\frac{z}{c}\right)^2 = \frac{x}{h}+\frac{y}{k}\ (x \geqslant 0, y \geqslant 0, z \geqslant 0)$；
\item $\displaystyle \left(\frac{x}{a}+\frac{y}{b}+\frac{z}{c}\right)^2 = \frac{x}{h}-\frac{y}{k}\ (x \geqslant 0, y \geqslant 0, z \geqslant 0)$.
\end{enumerate}
\end{prob}

\begin{prob}[M]{4117}
\begin{enumerate}[label=(\alph*), itemsep=2pt, topsep=3pt]
\item $\displaystyle \left(\frac{x}{a}+\frac{y}{b}+\frac{z}{c}\right)^4 = \frac{xyz}{abc}\ (x \geqslant 0, y \geqslant 0, z \geqslant 0)$；
\item $\displaystyle \left(\frac{x}{a}+\frac{y}{b}\right)^2 + \left(\frac{z}{c}\right)^2 = 1\ (x \geqslant 0, y \geqslant 0, z \geqslant 0)$.
\end{enumerate}
\end{prob}

\begin{prob}[L]{4118}
\begin{enumerate}[label=(\alph*), itemsep=2pt, topsep=3pt]
\item $\displaystyle \sqrt{\frac{x}{a}}+\sqrt{\frac{y}{b}}+\sqrt{\frac{z}{c}} = 1 \quad (x \geqslant 0, y \geqslant 0, z \geqslant 0)$；
\item $\displaystyle \sqrt[3]{\frac{x}{a}}+\sqrt[3]{\frac{y}{b}}+\sqrt[3]{\frac{z}{c}} = 1 \quad (x \geqslant 0, y \geqslant 0, z \geqslant 0)$；
\item $\displaystyle \left(\frac{x}{a}\right)^{\frac{2}{3}}+\left(\frac{y}{b}\right)^{\frac{2}{3}}+\left(\frac{z}{c}\right)^{\frac{2}{3}} = 1$.
\end{enumerate}
\end{prob}

\begin{prob}[S]{4119}
$z=x^2+y^2,\ z=2(x^2+y^2),\ xy=a^2,\ xy=2a^2,\ x=2y,\ 2x=y\ (x>0, y>0)$.
\end{prob}

\begin{prob}[S]{4120}
$x^2+z^2=a^2,\ x^2+z^2=b^2,\ x^2-y^2-z^2=0\ (x>0)$.
\end{prob}

\begin{prob}[S]{4121}
$\displaystyle (x^2+y^2+z^2)^3 = \frac{a^6 x^2}{x^2+y^2}$.
\end{prob}

\begin{prob}[S]{4122}
$\displaystyle \left(\frac{x^2}{a^2}+\frac{y^2}{b^2}+\frac{z^2}{c^2}\right)^2 = \frac{z}{h} \cdot \mathrm{e}^{-\frac{\frac{z^2}{c^2}}{\frac{x^2}{a^2}+\frac{y^2}{b^2}+\frac{z^2}{c^2}}}$.
\end{prob}

\begin{prob}[S]{4123}
$\displaystyle \frac{\dfrac{x}{a}+\dfrac{y}{b}}{\dfrac{x}{a}+\dfrac{y}{b}+\dfrac{z}{c}} = \frac{2}{\pi} \arcsin\left(\frac{x}{a}+\frac{y}{b}+\frac{z}{c}\right),\ \frac{x}{a}+\frac{y}{b}=1,\ x=0,\ x=a$.
\end{prob}

\begin{prob}[S]{4124}
$\displaystyle \frac{x}{a}+\frac{y}{b}+\frac{z}{c} = \ln \frac{\dfrac{x}{a}+\dfrac{y}{b}+\dfrac{z}{c}}{\dfrac{x}{a}+\dfrac{y}{b}},\ x=0,\ z=0,\ \frac{y}{b}+\frac{z}{c}=0,\ \frac{x}{a}+\frac{y}{b}+\frac{z}{c}=1$.
\end{prob}

\begin{prob}[M]{4125}
曲面 $x^2+y^2+az=4a^2$ 将球 $x^2+y^2+z^2=4az$ 分成两部分，求这两部分的体积之比.
\end{prob}

\begin{prob}[M]{4126}
求以曲面
\[
x^2+y^2=az,\quad z=2a-\sqrt{x^2+y^2} \quad (a>0)
\]
为界的物体的体积和表面积.
\end{prob}

\begin{prob}[M]{4127}
求以平面
\[
\begin{aligned}
a_1 x + b_1 y + c_1 z &= \pm h_1, \\
a_2 x + b_2 y + c_2 z &= \pm h_2, \\
a_3 x + b_3 y + c_3 z &= \pm h_3
\end{aligned}
\]
为界的平行六面体的体积，设
\[
\Delta = \begin{vmatrix}
a_1 & b_1 & c_1 \\
a_2 & b_2 & c_2 \\
a_3 & b_3 & c_3
\end{vmatrix} \neq 0.
\]
\end{prob}

\begin{prob}[M]{4128}
求以曲面
\[
(a_1 x + b_1 y + c_1 z)^2 + (a_2 x + b_2 y + c_2 z)^2 + (a_3 x + b_3 y + c_3 z)^2 = h^2
\]
为界的物体的体积，设
\[
\Delta = \begin{vmatrix}
a_1 & b_1 & c_1 \\
a_2 & b_2 & c_2 \\
a_3 & b_3 & c_3
\end{vmatrix} \neq 0.
\]
\end{prob}

\begin{prob}[M]{4129}
求以曲面
\[
\left(\frac{x^2}{a^2}+\frac{y^2}{b^2}\right)^n + \frac{z^{2n}}{c^{2n}} = \frac{z}{h} \left(\frac{x^2}{a^2}+\frac{y^2}{b^2}\right)^{n-2} \quad (n>1)
\]
为界的物体的体积.
\end{prob}

\begin{prob}[M]{4130}
一物体位于正象限 $Oxyz\ (x>0, y>0, z>0)$，并以曲面
\[
\frac{x^m}{a^m} + \frac{y^n}{b^n} + \frac{z^p}{c^p} = 1 \quad (m>0, n>0, p>0),\quad x=0,\quad y=0,\quad z=0
\]
为界，求其体积.
\end{prob}

\probsec{§8. 三重积分在力学上的应用}

\begin{prob}[M]{4131}
设物体在点 $M(x,y,z)$ 的密度由公式 $\rho = x + y + z$ 给出，求占有单位体积 $0 \leqslant x \leqslant 1, 0 \leqslant y \leqslant 1, 0 \leqslant z \leqslant 1$ 之物体的质量．
\end{prob}

\begin{prob}[M]{4132}
若物体的密度按规律 $\rho = \rho_0 \mathrm{e}^{-k\sqrt{x^2+y^2+z^2}}$（其中 $\rho_0 > 0$ 及 $k > 0$ 为常数）而变化，求占有无限区域 $x^2 + y^2 + z^2 \geqslant 1$ 的物体的质量．
\end{prob}

\probgroup{求下列曲面为界的均匀物体的质心坐标：}

\begin{prob}[S]{4133}
$\dfrac{x^2}{a^2} + \dfrac{y^2}{b^2} = \dfrac{z^2}{c^2},\; z = c$．
\end{prob}

\begin{prob}[S]{4134}
$z = x^2 + y^2,\; x + y = a,\; x = 0,\; y = 0,\; z = 0$．
\end{prob}

\begin{prob}[S]{4135}
$x^2 = 2pz,\; y^2 = 2px,\; x = \dfrac{p}{2},\; z = 0$．
\end{prob}

\begin{prob}[S]{4136}
$\dfrac{x^2}{a^2} + \dfrac{y^2}{b^2} + \dfrac{z^2}{c^2} = 1,\; x \geqslant 0,\; y \geqslant 0,\; z \geqslant 0$．
\end{prob}

\begin{prob}[S]{4137}
$x^2 + z^2 = a^2,\; y^2 + z^2 = a^2$ ($z \geqslant 0$)．
\end{prob}

\begin{prob}[S]{4138}
$x^2 + y^2 = 2z,\; x + y = z$．
\end{prob}

\begin{prob}[S]{4139}
$\left(\dfrac{x^2}{a^2} + \dfrac{y^2}{b^2} + \dfrac{z^2}{c^2}\right)^2 = \dfrac{xyz}{abc}$ ($x \geqslant 0,\; y \geqslant 0,\; z \geqslant 0,\; a > 0,\; b > 0,\; c > 0$)．
\end{prob}

\begin{prob}[S]{4140}
$z = x^2 + y^2,\; z = \dfrac{1}{2}(x^2 + y^2),\; x + y = \pm 1,\; x - y = \pm 1$．
\end{prob}

\begin{prob}[S]{4141}
$\dfrac{x^n}{a^n} + \dfrac{y^n}{b^n} + \dfrac{z^n}{c^n} = 1,\; x = 0,\; y = 0,\; z = 0$ ($n > 0;\; x \geqslant 0,\; y \geqslant 0,\; z \geqslant 0$)．
\end{prob}

\begin{prob}[M]{4142}
求形状为立方体
\[
0 \leqslant x \leqslant 1, \quad 0 \leqslant y \leqslant 1, \quad 0 \leqslant z \leqslant 1
\]
的物体的质心坐标，设此物体在点 $(x,y,z)$ 的密度等于
\[
\rho = x^{\frac{2\alpha-1}{1-\alpha}} y^{\frac{2\beta-1}{1-\beta}} z^{\frac{2\gamma-1}{1-\gamma}},
\]
其中 $0 < \alpha < 1,\; 0 < \beta < 1,\; 0 < \gamma < 1$．
\end{prob}

\probgroup{求下列曲面（参量是正的）为界的均匀物体对坐标平面的转动惯量：}

\begin{prob}[S]{4143}
$\dfrac{x}{a} + \dfrac{y}{b} + \dfrac{z}{c} = 1,\; x = 0,\; y = 0,\; z = 0$．
\end{prob}

\begin{prob}[S]{4144}
$\dfrac{x^2}{a^2} + \dfrac{y^2}{b^2} + \dfrac{z^2}{c^2} = 1$．
\end{prob}

\begin{prob}[S]{4145}
$\dfrac{x^2}{a^2} + \dfrac{y^2}{b^2} = \dfrac{z^2}{c^2},\; z = c$．
\end{prob}

\begin{prob}[M]{4146}
(a) $\dfrac{x^2}{a^2} + \dfrac{y^2}{b^2} + \dfrac{z^2}{c^2} = 1,\; \dfrac{x^2}{a^2} + \dfrac{y^2}{b^2} = \dfrac{x}{a}$；\qquad
(b) $\dfrac{x^2}{a^2} + \dfrac{y^2}{b^2} = 2\dfrac{z}{c},\; \dfrac{x}{a} + \dfrac{y}{b} = \dfrac{z}{c}$．
\end{prob}

\begin{prob}[M]{4147}
\begin{enumerate}[label=(\alph*), itemsep=2pt, topsep=3pt]
  \item $\left(\dfrac{x^2}{a^2} + \dfrac{y^2}{b^2} + \dfrac{z^2}{c^2}\right)^2 = \dfrac{x^2}{a^2} + \dfrac{y^2}{b^2} - \dfrac{z^2}{c^2}$；
  \item $\left(\dfrac{x}{a}\right)^n + \left(\dfrac{y}{b}\right)^n + \left(\dfrac{z}{c}\right)^n = 1,\; x = 0,\; y = 0,\; z = 0$ ($n > 0;\; x \geqslant 0,\; y \geqslant 0,\; z \geqslant 0$)．
\end{enumerate}
\end{prob}

\probgroup{求下列曲面为界的均匀物体对 $Oz$ 轴的转动惯量：}

\begin{prob}[S]{4148}
$z = x^2 + y^2,\; x + y = \pm 1,\; x - y = \pm 1,\; z = 0$．
\end{prob}

\begin{prob}[M]{4149}
(a) $x^2 + y^2 + z^2 = 2,\; x^2 + y^2 = z^2$ ($z > 0$)；\qquad
(b) $(x^2 + y^2 + z^2)^3 = a^5 z$．
\end{prob}

\begin{prob}[M]{4150}
求质量为 $M$ 的非均匀球体 $x^2 + y^2 + z^2 \leqslant R^2$ 对其直径的转动惯量，设球内各点 $P(x,y,z)$ 的密度与该点至球心的距离成正比．
\end{prob}

\begin{prob}[M]{4151}
证明等式：
\[
I_l = I_{l_0} + Md^2,
\]
其中 $I_l$ 为物体对某轴 $l$ 的转动惯量，$I_{l_0}$ 为对平行于 $l$ 并通过物体质心的轴 $l_0$ 的转动惯量，$d$ 为此二轴之间的距离，$M$ 为物体的质量．
\end{prob}

\begin{prob}[L]{4152}
证明：占有区域 $V$ 的物体对过其质心 $O(0,0,0)$ 并与坐标轴成角 $\alpha, \beta, \gamma$ 的轴 $l$ 的转动惯量等于：
\[
\begin{aligned}
I_l &= I_x \cos^2\alpha + I_y \cos^2\beta + I_z \cos^2\gamma - 2K_{xy} \cos\alpha \cos\beta \\
&\quad - 2K_{xz} \cos\alpha \cos\gamma - 2K_{yz} \cos\beta \cos\gamma,
\end{aligned}
\]
其中 $I_x, I_y, I_z$ 为物体对坐标轴的转动惯量，而
\[
K_{xy} = \iiint\limits_V \rho xy \,\mathrm{d}x\,\mathrm{d}y\,\mathrm{d}z, \quad
K_{xz} = \iiint\limits_V \rho xz \,\mathrm{d}x\,\mathrm{d}y\,\mathrm{d}z, \quad
K_{yz} = \iiint\limits_V \rho yz \,\mathrm{d}x\,\mathrm{d}y\,\mathrm{d}z
\]
为惯性积．
\end{prob}

\begin{prob}[M]{4153}
求密度为 $\rho_0$ 的均匀圆柱体 $x^2 + y^2 \leqslant a^2,\; -h \leqslant z \leqslant h$ 对直线 $x = y = z$ 的转动惯量．
\end{prob}

\begin{prob}[M]{4154}
求密度为 $\rho_0$，以曲面
\[
(x^2 + y^2 + z^2)^2 = a^2(x^2 + y^2)
\]
为界的均匀物体对坐标原点的转动惯量．
\end{prob}

\begin{prob}[M]{4155}
求密度为 $\rho_0$ 的均匀球体 $\xi^2 + \eta^2 + \zeta^2 \leqslant R^2$ 在点 $P(x,y,z)$ 的牛顿引力势．
\probhint{令 $O\zeta$ 轴通过点 $P(x,y,z)$．}
\end{prob}

\begin{prob}[M]{4156}
求球壳层 $R_1^2 \leqslant \xi^2 + \eta^2 + \zeta^2 \leqslant R_2^2$ 在点 $P(x,y,z)$ 的牛顿引力势，设密度 $\rho = f(R)$，此处 $f$ 为已知函数，而
\[
R = \sqrt{\xi^2 + \eta^2 + \zeta^2}.
\]
\end{prob}

\begin{prob}[M]{4157}
求密度 $\rho_0$ 恒定的圆柱体 $\xi^2 + \eta^2 \leqslant a^2,\; 0 \leqslant \zeta \leqslant h$ 在点 $P(0,0,z)$ 的牛顿引力势．
\end{prob}

\begin{prob}[M]{4158}
半径为 $R$ 质量为 $M$ 的均匀球体 $\xi^2 + \eta^2 + \zeta^2 \leqslant R^2$ 以怎样的力吸引质量为 $m$ 的质点 $P(0,0,a)$？
\end{prob}

\begin{prob}[M]{4159}
求密度为 $\rho_0$ 的均匀圆柱体 $\xi^2 + \eta^2 \leqslant a^2,\; 0 \leqslant \zeta \leqslant h$ 对单位质量质点 $P(0,0,z)$ 的引力．
\end{prob}

\begin{prob}[M]{4160}
求密度为 $\rho_0$ 的均匀球锥体对位于其顶点的单位质量质点的引力，设球面半径为 $R$，而轴截面的扇形的角等于 $2\alpha$．
\end{prob}

\probsec{§9. 二重和三重广义积分}

\probgroup{研究下列具有无界积分域的广义积分的收敛性（$0 < m \leqslant |\varphi(x,y)| \leqslant M$）：}

\begin{prob}[S]{4161}
$\displaystyle\iint\limits_{x^2+y^2 \geqslant 1} \dfrac{\varphi(x,y)}{(x^2+y^2)^p} \,\mathrm{d}x\,\mathrm{d}y$．
\end{prob}

\begin{prob}[S]{4162}
$\displaystyle\int_{-\infty}^{+\infty} \int_{-\infty}^{+\infty} \dfrac{\mathrm{d}x\,\mathrm{d}y}{(1+|x|^p)(1+|y|^q)}$．
\end{prob}

\begin{prob}[S]{4163}
$\displaystyle\iint\limits_{0 \leqslant y \leqslant 1} \dfrac{\varphi(x,y)}{(1+x^2+y^2)^p} \,\mathrm{d}x\,\mathrm{d}y$．
\end{prob}

\begin{prob}[S]{4164}
$\displaystyle\iint\limits_{|x|+|y| \geqslant 1} \dfrac{\mathrm{d}x\,\mathrm{d}y}{|x|^p + |y|^q}$ ($p > 0,\; q > 0$)．
\end{prob}

\begin{prob}[S]{4165}
$\displaystyle\iint\limits_{x+y \geqslant 1} \dfrac{\sin x \sin y}{(x+y)^p} \,\mathrm{d}x\,\mathrm{d}y$．
\end{prob}

\begin{prob}[M]{4166}
证明：若连续函数 $f(x,y)$ 不为负，$S_n$ ($n = 1, 2, \cdots$) 为有界闭区域，并且组成区域 $S$ 的任意一个竭尽递增序列，则
\[
\iint\limits_S f(x,y) \,\mathrm{d}x\,\mathrm{d}y = \lim_{n \to \infty} \iint\limits_{S_n} f(x,y) \,\mathrm{d}x\,\mathrm{d}y,
\]
这里左端与右端同时有意义或同时无意义．
\end{prob}

\begin{prob}[M]{4167}
证明：
\[
\lim_{n \to \infty} \iint\limits_{\substack{|x| \leqslant n \\ |y| \leqslant n}} \sin(x^2 + y^2) \,\mathrm{d}x\,\mathrm{d}y = \pi,
\]
然而
\[
\lim_{n \to \infty} \iint\limits_{x^2+y^2 \leqslant 2\pi n} \sin(x^2 + y^2) \,\mathrm{d}x\,\mathrm{d}y = 0
\]
（$n$ 为正整数）．
\end{prob}

\begin{prob}[M]{4168}
证明：尽管二重积分
\[
\int_1^{+\infty} \mathrm{d}x \int_1^{+\infty} \dfrac{x^2 - y^2}{(x^2 + y^2)^2} \,\mathrm{d}y \quad \text{及} \quad \int_1^{+\infty} \mathrm{d}y \int_1^{+\infty} \dfrac{x^2 - y^2}{(x^2 + y^2)^2} \,\mathrm{d}x
\]
收敛，但积分
\[
\iint\limits_{x \geqslant 1,\, y \geqslant 1} \dfrac{x^2 - y^2}{(x^2 + y^2)^2} \,\mathrm{d}x\,\mathrm{d}y
\]
发散．
\end{prob}

\probgroup{计算下列积分（参数是正的）：}

\begin{prob}[S]{4169}
$\displaystyle\iint\limits_{\substack{xy \geqslant 1 \\ x \geqslant 1}} \dfrac{\mathrm{d}x\,\mathrm{d}y}{x^p y^q}$．
\end{prob}

\begin{prob}[S]{4170}
$\displaystyle\iint\limits_{\substack{x+y \geqslant 1 \\ 0 \leqslant x \leqslant 1}} \dfrac{\mathrm{d}x\,\mathrm{d}y}{(x+y)^p}$．
\end{prob}

\begin{prob}[S]{4171}
$\displaystyle\iint\limits_{x^2+y^2 \leqslant 1} \dfrac{\mathrm{d}x\,\mathrm{d}y}{\sqrt{1 - x^2 - y^2}}$．
\end{prob}

\begin{prob}[S]{4172}
$\displaystyle\iint\limits_{x^2+y^2 \geqslant 1} \dfrac{\mathrm{d}x\,\mathrm{d}y}{(x^2 + y^2)^p}$．
\end{prob}

\begin{prob}[S]{4173}
$\displaystyle\iint\limits_{y \geqslant x^2 + 1} \dfrac{\mathrm{d}x\,\mathrm{d}y}{x^4 + y^2}$．
\end{prob}

\begin{prob}[S]{4174}
$\displaystyle\iint\limits_{0 \leqslant x \leqslant y} \mathrm{e}^{-(x+y)} \,\mathrm{d}x\,\mathrm{d}y$．
\end{prob}

\probgroup{变换为极坐标，计算下列积分：}

\begin{prob}[S]{4175}
$\displaystyle\int_{-\infty}^{+\infty} \int_{-\infty}^{+\infty} \mathrm{e}^{-(x^2+y^2)} \,\mathrm{d}x\,\mathrm{d}y$．
\end{prob}

\begin{prob}[S]{4176}
$\displaystyle\int_{-\infty}^{+\infty} \int_{-\infty}^{+\infty} \mathrm{e}^{-(x^2+y^2)} \cos(x^2 + y^2) \,\mathrm{d}x\,\mathrm{d}y$．
\end{prob}

\begin{prob}[S]{4177}
$\displaystyle\int_{-\infty}^{+\infty} \int_{-\infty}^{+\infty} \mathrm{e}^{-(x^2+y^2)} \sin(x^2 + y^2) \,\mathrm{d}x\,\mathrm{d}y$．
\end{prob}

\probgroup{计算下列积分：}

\begin{prob}[S]{4178}
$\displaystyle\int_{-\infty}^{+\infty} \int_{-\infty}^{+\infty} \mathrm{e}^{ax^2 + 2bxy + cy^2 + 2dx + 2ey + f} \,\mathrm{d}x\,\mathrm{d}y$，其中 $a < 0,\; ac - b^2 > 0$．
\end{prob}

\begin{prob}[S]{4179}
$\displaystyle\iint\limits_{\frac{x^2}{a^2}+\frac{y^2}{b^2} \geqslant 1} \mathrm{e}^{-\left(\frac{x^2}{a^2}+\frac{y^2}{b^2}\right)} \,\mathrm{d}x\,\mathrm{d}y$．
\end{prob}

\begin{prob}[S]{4180}
$\displaystyle\int_{-\infty}^{+\infty} \int_{-\infty}^{+\infty} xy \mathrm{e}^{-\left(\frac{x^2}{a^2} + 2\varepsilon \frac{x}{a}\frac{y}{b} + \frac{y^2}{b^2}\right)} \,\mathrm{d}x\,\mathrm{d}y$ ($0 < |\varepsilon| < 1$)．
\end{prob}

\probgroup{研究下列不连续函数的二重广义积分的收敛性（$0 < m \leqslant |\varphi(x,y)| \leqslant M$）：}

\begin{prob}[S]{4181}
$\displaystyle\iint\limits_{\Omega} \dfrac{\mathrm{d}x\,\mathrm{d}y}{x^2 + y^2}$，式中区域 $\Omega$ 由条件 $|y| \leqslant x^2,\; x^2 + y^2 \leqslant 1$ 确定．
\end{prob}

\begin{prob}[S]{4182}
$\displaystyle\iint\limits_{x^2+y^2 \leqslant 1} \dfrac{\varphi(x,y)}{(x^2 + xy + y^2)^p} \,\mathrm{d}x\,\mathrm{d}y$．
\end{prob}

\begin{prob}[S]{4183}
$\displaystyle\iint\limits_{|x|+|y| \leqslant 1} \dfrac{\mathrm{d}x\,\mathrm{d}y}{|x|^p + |y|^q}$ ($p > 0,\; q > 0$)．
\end{prob}

\begin{prob}[S]{4184}
$\displaystyle\int_0^a \int_0^a \dfrac{\varphi(x,y)}{|x - y|^p} \,\mathrm{d}x\,\mathrm{d}y$．
\end{prob}

\begin{prob}[S]{4185}
$\displaystyle\iint\limits_{x^2+y^2 \leqslant 1} \dfrac{\varphi(x,y)}{(1 - x^2 - y^2)^p} \,\mathrm{d}x\,\mathrm{d}y$．
\end{prob}

\begin{prob}[M]{4186}
证明：若 1) 函数 $\varphi(x,y)$ 在有界区域 $a \leqslant x \leqslant A,\; b \leqslant y \leqslant B$ 内连续；2) 函数 $f(x)$ 在闭区间 $a \leqslant x \leqslant A$ 上连续；3) $p < 1$，则积分
\[
\int_a^A \mathrm{d}x \int_b^B \dfrac{\varphi(x,y)}{|f(x) - y|^p} \,\mathrm{d}y
\]
收敛．
\end{prob}

\probgroup{计算下列积分：}

\begin{prob}[S]{4187}
$\displaystyle\iint\limits_{x^2+y^2 \leqslant 1} \ln\dfrac{1}{\sqrt{x^2 + y^2}} \,\mathrm{d}x\,\mathrm{d}y$．
\end{prob}

\begin{prob}[S]{4188}
$\displaystyle\int_0^a \mathrm{d}x \int_0^x \dfrac{\mathrm{d}y}{\sqrt{(a - x)(x - y)}}$ ($a > 0$)．
\end{prob}

\begin{prob}[S]{4189}
$\displaystyle\iint\limits_{\Omega} \ln\sin(x - y) \,\mathrm{d}x\,\mathrm{d}y$，其中区域 $\Omega$ 是由直线 $y = 0,\; y = x,\; x = \pi$ 围成的．
\end{prob}

\begin{prob}[S]{4190}
$\displaystyle\iint\limits_{x^2+y^2 \leqslant x} \dfrac{\mathrm{d}x\,\mathrm{d}y}{\sqrt{x^2 + y^2}}$．
\end{prob}

\probgroup{研究下列三重积分的收敛性：}

\begin{prob}[S]{4191}
$\displaystyle\iiint\limits_{x^2+y^2+z^2 > 1} \dfrac{\varphi(x,y,z)}{(x^2 + y^2 + z^2)^p} \,\mathrm{d}x\,\mathrm{d}y\,\mathrm{d}z$，其中 $0 < m \leqslant |\varphi(x,y,z)| \leqslant M$．
\end{prob}

\begin{prob}[S]{4192}
$\displaystyle\iiint\limits_{x^2+y^2+z^2 \leqslant 1} \dfrac{\varphi(x,y,z)}{(x^2 + y^2 + z^2)^p} \,\mathrm{d}x\,\mathrm{d}y\,\mathrm{d}z$，其中 $0 < m \leqslant |\varphi(x,y,z)| \leqslant M$．
\end{prob}

\begin{prob}[S]{4193}
$\displaystyle\iiint\limits_{|x|+|y|+|z| > 1} \dfrac{\mathrm{d}x\,\mathrm{d}y\,\mathrm{d}z}{|x|^p + |y|^q + |z|^r}$ ($p > 0,\; q > 0,\; r > 0$)．
\end{prob}

\begin{prob}[M]{4194}
$\displaystyle\int_0^a \int_0^a \int_0^a \dfrac{f(x,y,z) \,\mathrm{d}x\,\mathrm{d}y\,\mathrm{d}z}{\{[y - \varphi(x)]^2 + [z - \psi(x)]^2\}^p}$，其中 $0 < m \leqslant |f(x,y,z)| \leqslant M$，而 $\varphi(x)$ 和 $\psi(x)$ 是闭区间 $[0, a]$ 上的连续函数．
\end{prob}

\begin{prob}[S]{4195}
$\displaystyle \iiint\limits_{\substack{|x|\leqslant 1 \\ |y|\leqslant 1 \\ |z|\leqslant 1}} \frac{\mathrm{d}x\,\mathrm{d}y\,\mathrm{d}z}{|x+y-z|^p}$.
\end{prob}

\probgroup{计算下列积分：}

\begin{prob}[S]{4196}
$\displaystyle \int_0^1 \int_0^1 \int_0^1 \frac{\mathrm{d}x\,\mathrm{d}y\,\mathrm{d}z}{x^p y^q z^r}$.
\end{prob}

\begin{prob}[S]{4197}
$\displaystyle \iiint\limits_{x^2+y^2+z^2>1} \frac{\mathrm{d}x\,\mathrm{d}y\,\mathrm{d}z}{(x^2+y^2+z^2)^3}$.
\end{prob}

\begin{prob}[S]{4198}
$\displaystyle \iiint\limits_{x^2+y^2+z^2\leqslant 1} \frac{\mathrm{d}x\,\mathrm{d}y\,\mathrm{d}z}{(1-x^2-y^2-z^2)^p}$.
\end{prob}

\begin{prob}[S]{4199}
$\displaystyle \int_{-\infty}^{+\infty} \int_{-\infty}^{+\infty} \int_{-\infty}^{+\infty} \mathrm{e}^{-(x^2+y^2+z^2)}\,\mathrm{d}x\,\mathrm{d}y\,\mathrm{d}z$.
\end{prob}

\begin{prob}[M]{4200}
计算积分
\[
\int_{-\infty}^{+\infty} \int_{-\infty}^{+\infty} \int_{-\infty}^{+\infty} \mathrm{e}^{-P(x_1, x_2, x_3)}\,\mathrm{d}x_1\,\mathrm{d}x_2\,\mathrm{d}x_3,
\]
其中 $P(x_1, x_2, x_3) = \sum_{i=1}^3 \sum_{j=1}^3 a_{ij} x_i x_j (a_{ij}=a_{ji})$ 为正定二次型.
\end{prob}

\probsec{§10. 多重积分}

\begin{prob}[M]{4201}
设 $K(x, y)$ 为区域 $R(a \leqslant x \leqslant b; a \leqslant y \leqslant b)$ 内的连续函数，且
\[
K_n(x, y) = \int_a^b \int_a^b \cdots \int_a^b K(x, t_1) K(t_1, t_2) \cdots K(t_n, y)\,\mathrm{d}t_1\,\mathrm{d}t_2 \cdots \mathrm{d}t_n,
\]
证明：
\[
K_{n+m+1}(x, y) = \int_a^b K_n(x, t) K_m(t, y)\,\mathrm{d}t.
\]
\end{prob}

\begin{prob}[M]{4202}
设 $f = f(x_1, x_2, \cdots, x_n)$ 为区域 $0 \leqslant x_i \leqslant x$ ($i = 1, 2, \cdots, n$) 内的连续函数. 证明等式：
\[
\int_0^x \mathrm{d}x_1 \int_0^{x_1} \mathrm{d}x_2 \cdots \int_0^{x_{n-1}} f\,\mathrm{d}x_n = \int_0^x \mathrm{d}x_n \int_{x_n}^x \mathrm{d}x_{n-1} \cdots \int_{x_2}^x f\,\mathrm{d}x_1 \quad (n \geqslant 2).
\]
\end{prob}

\begin{prob}[M]{4203}
证明：
\[
\int_0^t \mathrm{d}t_1 \int_0^{t_1} \mathrm{d}t_2 \cdots \int_0^{t_{n-1}} f(t_1) f(t_2) \cdots f(t_n)\,\mathrm{d}t_n = \frac{1}{n!} \left\{ \int_0^t f(\tau)\,\mathrm{d}\tau \right\}^n,
\]
其中 $f$ 为连续函数.
\end{prob}

\probgroup{计算下列多重积分：}

\begin{prob}[M]{4204}
\begin{enumerate}[label=(\alph*), itemsep=2pt, topsep=3pt]
  \item $\displaystyle \int_0^1 \int_0^1 \cdots \int_0^1 (x_1^2 + x_2^2 + \cdots + x_n^2)\,\mathrm{d}x_1\,\mathrm{d}x_2 \cdots \mathrm{d}x_n$;
  \item $\displaystyle \int_0^1 \int_0^1 \cdots \int_0^1 (x_1 + x_2 + \cdots + x_n)^2\,\mathrm{d}x_1\,\mathrm{d}x_2 \cdots \mathrm{d}x_n$.
\end{enumerate}
\end{prob}

\begin{prob}[S]{4205}
$\displaystyle I_n = \iint\limits_{\substack{x_1\geqslant 0, x_2\geqslant 0, \cdots, x_n\geqslant 0 \\ x_1+x_2+\cdots+x_n\leqslant a}} \cdots \int \mathrm{d}x_1\,\mathrm{d}x_2 \cdots \mathrm{d}x_n$.
\end{prob}

\begin{prob}[S]{4206}
$\displaystyle \int_0^1 \mathrm{d}x_1 \int_0^{x_1} \mathrm{d}x_2 \cdots \int_0^{x_{n-1}} x_1 x_2 \cdots x_n\,\mathrm{d}x_n$.
\end{prob}

\begin{prob}[S]{4207}
$\displaystyle \iint\limits_{\substack{x_1\geqslant 0, x_2\geqslant 0, \cdots, x_n\geqslant 0 \\ x_1+x_2+\cdots+x_n\leqslant 1}} \cdots \int \sqrt{x_1+x_2+\cdots+x_n}\,\mathrm{d}x_1\,\mathrm{d}x_2 \cdots \mathrm{d}x_n$.
\end{prob}

\begin{prob}[M]{4208}
求以平面
\[
a_{i1} x_1 + a_{i2} x_2 + \cdots + a_{in} x_n = \pm h_i \quad (i = 1, 2, \cdots, n)
\]
为界的 $n$ 维平行 $2n$ 面体的体积，这里设 $\Delta = |a_{ij}| \neq 0$.
\end{prob}

\begin{prob}[M]{4209}
求 $n$ 维角锥
\[
\frac{x_1}{a_1} + \frac{x_2}{a_2} + \cdots + \frac{x_n}{a_n} \leqslant 1, \quad x_i \geqslant 0 \ (i = 1, 2, \cdots, n) \quad (a_i > 0, i = 1, 2, \cdots, n)
\]
的体积.
\end{prob}

\begin{prob}[M]{4210}
求以曲面
\[
\frac{x_1^2}{a_1^2} + \frac{x_2^2}{a_2^2} + \cdots + \frac{x_{n-1}^2}{a_{n-1}^2} = \frac{x_n^2}{a_n^2}, \quad x_n = a_n
\]
为界的 $n$ 维圆锥的体积.
\end{prob}

\begin{prob}[M]{4211}
求 $n$ 维球体
\[
x_1^2 + x_2^2 + \cdots + x_n^2 \leqslant a^2
\]
的体积.
\end{prob}

\begin{prob}[M]{4212}
求 $\displaystyle \iint\limits_\Omega \cdots \int x_n^2\,\mathrm{d}x_1\,\mathrm{d}x_2 \cdots \mathrm{d}x_n$，其中区域 $\Omega$ 是由以下不等式确定的：
\[
x_1^2 + x_2^2 + \cdots + x_{n-1}^2 \leqslant a^2, \quad -\frac{h}{2} \leqslant x_n \leqslant \frac{h}{2}.
\]
\end{prob}

\begin{prob}[M]{4213}
计算
\[
\iint\limits_{x_1^2+x_2^2+\cdots+x_n^2\leqslant 1} \cdots \int \frac{\mathrm{d}x_1\,\mathrm{d}x_2 \cdots \mathrm{d}x_n}{\sqrt{1-x_1^2-x_2^2-\cdots-x_n^2}}.
\]
\end{prob}

\begin{prob}[M]{4214}
证明等式：
\[
\int_0^x \mathrm{d}x_1 \int_0^{x_1} \mathrm{d}x_2 \cdots \int_0^{x_{n-1}} f(x_n)\,\mathrm{d}x_n = \int_0^x f(u) \frac{(x-u)^{n-1}}{(n-1)!}\,\mathrm{d}u.
\]
\end{prob}

\begin{prob}[M]{4215}
证明等式：
\[
\int_0^x x_1\,\mathrm{d}x_1 \int_0^{x_1} x_2\,\mathrm{d}x_2 \cdots \int_0^{x_n} f(x_{n+1})\,\mathrm{d}x_{n+1} = \frac{1}{2^n n!} \int_0^x (x^2 - u^2)^n f(u)\,\mathrm{d}u.
\]
\end{prob}

\begin{prob}[M]{4216}
证明狄利克雷公式：
\[
\iint\limits_{\substack{x_1, x_2, \cdots, x_n \geqslant 0 \\ x_1+x_2+\cdots+x_n \leqslant 1}} \cdots \int x_1^{p_1-1} x_2^{p_2-1} \cdots x_n^{p_n-1}\,\mathrm{d}x_1\,\mathrm{d}x_2 \cdots \mathrm{d}x_n = \frac{\Gamma(p_1)\Gamma(p_2)\cdots\Gamma(p_n)}{\Gamma(p_1+p_2+\cdots+p_n+1)} \quad (p_1, p_2, \cdots, p_n > 0).
\]
\end{prob}

\begin{prob}[M]{4217}
证明刘维尔公式：
\[
\iint\limits_{\substack{x_1, x_2, \cdots, x_n \geqslant 0 \\ x_1+x_2+\cdots+x_n \leqslant 1}} \cdots \int f(x_1+x_2+\cdots+x_n) x_1^{p_1-1} x_2^{p_2-1} \cdots x_n^{p_n-1}\,\mathrm{d}x_1\,\mathrm{d}x_2 \cdots \mathrm{d}x_n = \frac{\Gamma(p_1)\Gamma(p_2)\cdots\Gamma(p_n)}{\Gamma(p_1+p_2+\cdots+p_n)} \int_0^1 f(u) u^{p_1+p_2+\cdots+p_n-1}\,\mathrm{d}u \quad (p_1, p_2, \cdots, p_n > 0),
\]
式中 $f(u)$ 为连续函数.
\probhint{运用数学归纳法.}
\end{prob}

\begin{prob}[M]{4218}
将区域 $x_1^2 + x_2^2 + \cdots + x_n^2 \leqslant R^2$ 上的 $n$ 重积分 $(n \geqslant 2)$
\[
\iint\limits_\Omega \cdots \int f\left(\sqrt{x_1^2+x_2^2+\cdots+x_n^2}\right)\,\mathrm{d}x_1\,\mathrm{d}x_2 \cdots \mathrm{d}x_n
\]
(其中 $f(u)$ 为连续函数) 化为一重积分.
\end{prob}

\begin{prob}[M]{4219}
计算半径为 $R$，密度为 $\rho_0$ 的均质球体对自身的引力势，即求积分
\[
u = \frac{\rho_0^2}{2} \iiint\iiint\limits_{\substack{x_1^2+y_1^2+z_1^2 \leqslant R^2 \\ x_2^2+y_2^2+z_2^2 \leqslant R^2}} \frac{\mathrm{d}x_1\,\mathrm{d}y_1\,\mathrm{d}z_1\,\mathrm{d}x_2\,\mathrm{d}y_2\,\mathrm{d}z_2}{r_{1,2}},
\]
式中 $r_{1,2} = \sqrt{(x_1-x_2)^2 + (y_1-y_2)^2 + (z_1-z_2)^2}$.
\end{prob}

\begin{prob}[M]{4220}
设 $\sum_{i,j=1}^n a_{ij} x_i x_j$ ($a_{ij}=a_{ji}$) 为正定二次型，计算 $n$ 重积分
\[
\int_{-\infty}^{+\infty} \int_{-\infty}^{+\infty} \cdots \int_{-\infty}^{+\infty} \mathrm{e}^{-\left\{\sum_{i,j=1}^n a_{ij} x_i x_j + 2\sum_{i=1}^n b_i x_i + c\right\}}\,\mathrm{d}x_1\,\mathrm{d}x_2 \cdots \mathrm{d}x_n.
\]
\end{prob}

\probsec{§11. 曲线积分}

\probgroup{计算下列第一型曲线积分：}

\begin{prob}[S]{4221}
$\displaystyle \int\limits_C (x+y)\,\mathrm{d}s$，其中 $C$ 为以 $O(0,0)$，$A(1,0)$ 和 $B(0,1)$ 为顶点的三角形围线.
\end{prob}

\begin{prob}[S]{4222}
$\displaystyle \int\limits_C y^2\,\mathrm{d}s$，其中 $C$ 为摆线 $x = a(t - \sin t)$，$y = a(1 - \cos t)$ ($0 \leqslant t \leqslant 2\pi$) 的一拱.
\end{prob}

\begin{prob}[S]{4223}
$\displaystyle \int\limits_C (x^2+y^2)\,\mathrm{d}s$，其中 $C$ 为曲线 $x = a(\cos t + t\sin t)$，$y = a(\sin t - t\cos t)$ ($0 \leqslant t \leqslant 2\pi$).
\end{prob}

\begin{prob}[S]{4224}
$\displaystyle \int\limits_C xy\,\mathrm{d}s$，其中 $C$ 为双曲线 $x = a\cosh t$，$y = a\sinh t$ ($0 \leqslant t \leqslant t_0$) 的弧.
\end{prob}

\begin{prob}[S]{4225}
$\displaystyle \int\limits_C (x^{\frac{4}{3}} + y^{\frac{4}{3}})\,\mathrm{d}s$，其中 $C$ 为内摆线 $x^{\frac{2}{3}} + y^{\frac{2}{3}} = a^{\frac{2}{3}}$ 的弧.
\end{prob}

\begin{prob}[S]{4226}
$\displaystyle \int\limits_C \mathrm{e}^{\sqrt{x^2+y^2}}\,\mathrm{d}s$，其中 $C$ 为由曲线 $r = a$，$\varphi = 0$，$\varphi = \frac{\pi}{4}$ ($r$ 和 $\varphi$ 为极坐标) 给出的凸围线.
\end{prob}

\begin{prob}[S]{4227}
$\displaystyle \int\limits_C |y|\,\mathrm{d}s$，其中 $C$ 为双纽线 $(x^2+y^2)^2 = a^2(x^2-y^2)$ 的弧.
\end{prob}

\begin{prob}[S]{4228}
$\displaystyle\int_C x\,\mathrm{d}s$，其中 $C$ 为对数螺线 $r = a\mathrm{e}^{k\varphi}$ ($k > 0$) 在圆 $r = a$ 内的部分．
\end{prob}

\begin{prob}[S]{4229}
$\displaystyle\int_C \sqrt{x^2+y^2}\,\mathrm{d}s$，其中 $C$ 为圆周 $x^2+y^2 = ax$．
\end{prob}

\begin{prob}[S]{4230}
$\displaystyle\int_C \frac{\mathrm{d}s}{y^2}$，其中 $C$ 为悬链线 $y = a\cosh\frac{x}{a}$．
\end{prob}

\probgroup{求下列空间曲线的弧长（参数是正的）：}

\begin{prob}[S]{4231}
$x = 3t, y = 3t^2, z = 2t^3$，从 $O(0,0,0)$ 到 $A(3,3,2)$．
\end{prob}

\begin{prob}[S]{4232}
$x = \mathrm{e}^{-t}\cos t, y = \mathrm{e}^{-t}\sin t, z = \mathrm{e}^{-t}$，当 $0 < t < +\infty$．
\end{prob}

\begin{prob}[S]{4233}
$y = a\arcsin\frac{x}{a}, z = \frac{a}{4}\ln\frac{a-x}{a+x}$，从 $O(0,0,0)$ 到 $A(x_0, y_0, z_0)$．
\end{prob}

\begin{prob}[S]{4234}
$(x-y)^2 = a(x+y), x^2-y^2 = \frac{9}{8}x^2$，从 $O(0,0,0)$ 到 $A(x_0, y_0, z_0)$．
\end{prob}

\begin{prob}[S]{4235}
$x^2+y^2 = cz, \frac{y}{x} = \tan\frac{z}{c}$，从 $O(0,0,0)$ 到 $A(x_0, y_0, z_0)$．
\end{prob}

\begin{prob}[S]{4236}
$x^2+y^2+z^2 = a^2, \sqrt{x^2+y^2}\cosh\left(\arctan\frac{y}{x}\right) = a$，从点 $A(a,0,0)$ 到点 $B(x,y,z)$．
\end{prob}

\probgroup{计算沿空间曲线的第一型曲线积分：}

\begin{prob}[S]{4237}
$\displaystyle\int_C (x^2+y^2+z^2)\,\mathrm{d}s$，其中 $C$ 为螺旋线 $x = a\cos t, y = a\sin t, z = bt$ ($0 \leqslant t \leqslant 2\pi$) 的一段．
\end{prob}

\begin{prob}[S]{4238}
$\displaystyle\int_C x^2\,\mathrm{d}s$，其中 $C$ 为圆周 $x^2+y^2+z^2 = a^2, x+y+z = 0$．
\end{prob}

\begin{prob}[S]{4239}
$\displaystyle\int_C z\,\mathrm{d}s$，其中 $C$ 为圆锥螺旋 $x = t\cos t, y = t\sin t, z = t$ ($0 \leqslant t \leqslant t_0$)．
\end{prob}

\begin{prob}[S]{4240}
$\displaystyle\int_C z\,\mathrm{d}s$，其中 $C$ 为曲线 $x^2+y^2 = z^2, y^2 = ax$ 上从点 $O(0,0,0)$ 到点 $A(a, a, \sqrt{2}a)$ 的弧．
\end{prob}

\begin{prob}[M]{4241.1}
设曲线 $x = a\cos t, y = b\sin t$ ($a \geqslant b > 0, 0 \leqslant t \leqslant 2\pi$) 在点 $(x,y)$ 的线密度等于 $|y|$，求其质量．
\end{prob}

\begin{prob}[M]{4241.2}
设抛物线 $y^2 = 2px$ $\left(0 \leqslant x \leqslant \frac{p}{2}\right)$ 在点 $M(x,y)$ 的线密度等于 $|y|$，求其质量．
\end{prob}

\begin{prob}[M]{4241.3}
求曲线 $x = at, y = \frac{a}{2}t^2, z = \frac{a}{3}t^3$ ($0 \leqslant t \leqslant 1$) 的弧之质量，其密度按规律 $\rho = \sqrt{\frac{2y}{a}}$ 而变化．
\end{prob}

\begin{prob}[M]{4242}
计算均匀曲线 $y = a\cosh\frac{x}{a}$ 从点 $A(0,a)$ 到点 $B(b,h)$ 的弧的质心的坐标．
\end{prob}

\begin{prob}[M]{4243}
求摆线 $x = a(t-\sin t), y = a(1-\cos t)$ ($0 \leqslant t \leqslant \pi$) 的弧的质心．
\end{prob}

\begin{prob}[M]{4244.1}
求星形线 $x^{\frac{2}{3}} + y^{\frac{2}{3}} = a^{\frac{2}{3}}$ ($x \geqslant 0, y \geqslant 0$) 的弧 $C$ 对坐标轴的静矩
\[
S_y = \int_C x\,\mathrm{d}s, \qquad S_x = \int_C y\,\mathrm{d}s.
\]
\end{prob}

\begin{prob}[M]{4244.2}
求圆周 $x^2+y^2 = a^2$ 对其直径的转动惯量．
\end{prob}

\begin{prob}[M]{4244.3}
求下列曲线对点 $O(0,0)$ 的极转动惯量：
\begin{enumerate}[label=(\alph*), itemsep=2pt, topsep=3pt]
\item 正方形 $\max\{|x|, |y|\} = a$ 的边界 $C$；
\item 以极坐标点 $P(a,0)$，$Q\left(a, \frac{2\pi}{3}\right)$，$R\left(a, \frac{4\pi}{3}\right)$ 为顶点的正三角形的边界 $C$．
\end{enumerate}
\end{prob}

\begin{prob}[M]{4244.4}
求星形线 $x^{\frac{2}{3}} + y^{\frac{2}{3}} = a^{\frac{2}{3}}$ 的平均极半径，即由以下公式给出的数 $r_0$ ($r_0 > 0$)：
\[
I_0 = sr_0^2,
\]
其中 $I_0$ 是星形线对坐标原点的极转动惯量（参见习题 4244.3），$s$ 是星形线的弧长．
\end{prob}

\begin{prob}[M]{4245}
计算球面上的三角形 $x^2+y^2+z^2 = a^2, x \geqslant 0, y \geqslant 0, z \geqslant 0$ 的围线的质心的坐标．
\end{prob}

\begin{prob}[M]{4246}
求均匀的弧 $x = \mathrm{e}^t\cos t, y = \mathrm{e}^t\sin t, z = \mathrm{e}^t$ ($-\infty < t \leqslant 0$) 的质心的坐标．
\end{prob}

\begin{prob}[M]{4247}
求螺线 $x = a\cos t, y = a\sin t, z = \frac{h}{2\pi}t$ ($0 \leqslant t \leqslant 2\pi$) 的一支对坐标轴的转动惯量．
\end{prob}

\begin{prob}[L]{4248}
计算第二型曲线积分
\[
\int_{OA} x\,\mathrm{d}y - y\,\mathrm{d}x,
\]
式中 $O$ 为坐标原点，点 $A$ 的坐标为 $(1,2)$，并设：
\begin{enumerate}[label=(\alph*), itemsep=2pt, topsep=3pt]
\item $OA$ 为直线段；
\item $OA$ 为抛物线，其轴为 $Oy$；
\item $OA$ 为由 $Ox$ 轴上的线段 $OB$ 和平行于 $Oy$ 轴的线段 $BA$ 组成的折线．
\end{enumerate}
\end{prob}

\begin{prob}[M]{4249}
对于上题中所指示的途径 (a), (b), (c)，计算
\[
\int_{OA} x\,\mathrm{d}y + y\,\mathrm{d}x.
\]
\end{prob}

\probgroup{在参数增加的方向，沿所指示的曲线计算下列第二型曲线积分：}

\begin{prob}[S]{4250}
$\displaystyle\int_C (x^2-2xy)\,\mathrm{d}x + (y^2-2xy)\,\mathrm{d}y$，其中 $C$ 为抛物线 $y = x^2$ ($-1 \leqslant x \leqslant 1$)．
\end{prob}

\begin{prob}[S]{4251}
$\displaystyle\int_C (x^2+y^2)\,\mathrm{d}x + (x^2-y^2)\,\mathrm{d}y$，其中 $C$ 为曲线 $y = 1 - |1-x|$ ($0 \leqslant x \leqslant 2$)．
\end{prob}

\begin{prob}[S]{4252}
$\displaystyle\oint_C (x+y)\,\mathrm{d}x + (x-y)\,\mathrm{d}y$，其中 $C$ 为依逆时针方向通过的椭圆周 $\frac{x^2}{a^2} + \frac{y^2}{b^2} = 1$．
\end{prob}

\begin{prob}[S]{4253}
$\displaystyle\int_C (2a-y)\,\mathrm{d}x + x\,\mathrm{d}y$，其中 $C$ 为摆线 $x = a(t-\sin t), y = a(1-\cos t)$ ($0 \leqslant t \leqslant 2\pi$) 的一拱．
\end{prob}

\begin{prob}[S]{4254}
$\displaystyle\oint_C \frac{(x+y)\,\mathrm{d}x - (x-y)\,\mathrm{d}y}{x^2+y^2}$，其中 $C$ 为依逆时针方向通过的圆周 $x^2+y^2 = a^2$．
\end{prob}

\begin{prob}[M]{4255}
$\displaystyle\oint_{ABCDA} \frac{\mathrm{d}x + \mathrm{d}y}{|x| + |y|}$，其中 $ABCDA$ 为以 $A(1,0), B(0,1), C(-1,0), D(0,-1)$ 为顶点的正方形围线．
\end{prob}

\begin{prob}[S]{4256}
$\displaystyle\int_{AB} \sin y\,\mathrm{d}x + \sin x\,\mathrm{d}y$，其中 $AB$ 为介于点 $A(0,\pi)$ 和点 $B(\pi,0)$ 之间的直线段．
\end{prob}

\begin{prob}[M]{4257}
$\displaystyle\oint_{OmAnO} \arctan\frac{y}{x}\,\mathrm{d}y - \mathrm{d}x$，其中 $OmA$ 为抛物线段 $y = x^2$，$OnA$ 为直线段 $y = x$．
\end{prob}

\probgroup{验证被积函数为全微分，并计算下列曲线积分：}

\begin{prob}[S]{4258}
$\displaystyle\int_{(-1,2)}^{(2,3)} x\,\mathrm{d}y + y\,\mathrm{d}x$．
\end{prob}

\begin{prob}[S]{4259}
$\displaystyle\int_{(0,1)}^{(3,-4)} x\,\mathrm{d}x + y\,\mathrm{d}y$．
\end{prob}

\begin{prob}[S]{4260}
$\displaystyle\int_{(0,1)}^{(2,3)} (x+y)\,\mathrm{d}x + (x-y)\,\mathrm{d}y$．
\end{prob}

\begin{prob}[S]{4261}
$\displaystyle\int_{(1,-1)}^{(1,1)} (x-y)(\mathrm{d}x - \mathrm{d}y)$．
\end{prob}

\begin{prob}[S]{4262}
$\displaystyle\int_{(0,0)}^{(a,b)} f(x+y)(\mathrm{d}x + \mathrm{d}y)$，其中 $f(u)$ 为连续函数．
\end{prob}

\begin{prob}[S]{4263}
$\displaystyle\int_{(2,1)}^{(1,2)} \frac{y\,\mathrm{d}x - x\,\mathrm{d}y}{x^2}$ 沿着不与 $Oy$ 轴相交的路径．
\end{prob}

\begin{prob}[S]{4264}
$\displaystyle\int_{(1,0)}^{(6,8)} \frac{x\,\mathrm{d}x + y\,\mathrm{d}y}{\sqrt{x^2+y^2}}$ 沿着不通过坐标原点的路径．
\end{prob}

\begin{prob}[S]{4265}
$\displaystyle\int_{(x_1,y_1)}^{(x_2,y_2)} \varphi(x)\,\mathrm{d}x + \psi(y)\,\mathrm{d}y$，其中 $\varphi$ 和 $\psi$ 为连续函数．
\end{prob}

\begin{prob}[S]{4266}
$\displaystyle\int_{(-2,-1)}^{(3,0)} (x^4 + 4xy^3)\,\mathrm{d}x + (6x^2y^2 - 5y^4)\,\mathrm{d}y$．
\end{prob}

\begin{prob}[S]{4267}
$\displaystyle\int_{(0,-1)}^{(1,0)} \frac{x\,\mathrm{d}y - y\,\mathrm{d}x}{(x-y)^2}$ 沿着不与直线 $y = x$ 相交的路径．
\end{prob}

\begin{prob}[S]{4268}
$\displaystyle\int_{(1,\pi)}^{(2,\pi)} \left(1 - \frac{y^2}{x^2}\cos\frac{y}{x}\right)\,\mathrm{d}x + \left(\sin\frac{y}{x} + \frac{y}{x}\cos\frac{y}{x}\right)\,\mathrm{d}y$ 沿着不与 $Oy$ 轴相交的路径．
\end{prob}

\begin{prob}[S]{4269}
$\displaystyle\int_{(0,0)}^{(a,b)} \mathrm{e}^x(\cos y\,\mathrm{d}x - \sin y\,\mathrm{d}y)$．
\end{prob}

\begin{prob}[M]{4270}
证明：若 $f(u)$ 为连续函数且 $C$ 为逐段光滑的封闭围线，则
\[
\oint_C f(x^2+y^2)(x\,\mathrm{d}x + y\,\mathrm{d}y) = 0.
\]
\end{prob}

\probgroup{求原函数 $z$，设：}

\begin{prob}[S]{4271}
$\mathrm{d}z = (x^2 + 2xy - y^2)\,\mathrm{d}x + (x^2 - 2xy - y^2)\,\mathrm{d}y$．
\end{prob}

\begin{prob}[S]{4272}
$\mathrm{d}z = \frac{y\,\mathrm{d}x - x\,\mathrm{d}y}{3x^2 - 2xy + 3y^2}$．
\end{prob}

\begin{prob}[S]{4273}
$\mathrm{d}z = \frac{(x^2 + 2xy + 5y^2)\,\mathrm{d}x + (x^2 - 2xy + y^2)\,\mathrm{d}y}{(x+y)^3}$．
\end{prob}

\begin{prob}[S]{4274}
$\mathrm{d}z = \mathrm{e}^x[\mathrm{e}^y(x - y + 2) + y]\,\mathrm{d}x + \mathrm{e}^x[\mathrm{e}^y(x - y) + 1]\,\mathrm{d}y$．
\end{prob}

\begin{prob}[S]{4275}
$\mathrm{d}z = \frac{\partial^{n+m+1}u}{\partial x^{n+1}\partial y^m}\,\mathrm{d}x + \frac{\partial^{n+m+1}u}{\partial x^n\partial y^{m+1}}\,\mathrm{d}y$．
\end{prob}

\begin{prob}[M]{4276}
$\mathrm{d}z = \frac{\partial^{n+m+1}}{\partial x^{n+2}\partial y^{m-1}}\left(\ln\frac{1}{r}\right)\,\mathrm{d}x - \frac{\partial^{n+m+1}}{\partial x^{n-1}\partial y^{m+2}}\left(\ln\frac{1}{r}\right)\,\mathrm{d}y$，其中 $r = \sqrt{x^2+y^2}$．
\end{prob}

\begin{prob}[M]{4277}
证明：下面的估计对于曲线积分是正确的：
\[
\left|\int_C P\,\mathrm{d}x + Q\,\mathrm{d}y\right| \leqslant LM,
\]
其中 $L$ 为积分路径的长，$M = \max \sqrt{P^2+Q^2}$（在弧 $C$ 上）．
\end{prob}

\begin{prob}[M]{4278}
估计积分
\[
I_R = \oint_{x^2+y^2=R^2} \frac{y\,\mathrm{d}x - x\,\mathrm{d}y}{(x^2+xy+y^2)^2}.
\]
证明：$\lim_{R\to\infty} I_R = 0$．
\end{prob}

\probgroup{计算沿空间曲线的积分（假定坐标系是右手的）：}

\begin{prob}[S]{4279}
$\displaystyle\int_C (y^2 - z^2)\,\mathrm{d}x + 2yz\,\mathrm{d}y - x^2\,\mathrm{d}z$，式中 $C$ 为曲线 $x = t, y = t^2, z = t^3$ ($0 \leqslant t \leqslant 1$)，以参数增加的方向为正．
\end{prob}

\begin{prob}[S]{4280}
$\displaystyle\int_C y\,\mathrm{d}x + z\,\mathrm{d}y + x\,\mathrm{d}z$，式中 $C$ 为纽形螺线 $x = a\cos t, y = a\sin t, z = bt$ ($0 \leqslant t \leqslant 2\pi$)，以参数增加的方向为正．
\end{prob}

\begin{prob}[M]{4281}
$\displaystyle\int_C (y-z)\,\mathrm{d}x + (z-x)\,\mathrm{d}y + (x-y)\,\mathrm{d}z$，式中 $C$ 为圆周 $x^2+y^2+z^2 = a^2, y = x\tan\alpha$ ($0 < \alpha < \pi$)，若从 $x$ 轴的正向看去，沿逆时针方向为正．
\end{prob}

\begin{prob}[M]{4282}
$\displaystyle\int_C y^2\,\mathrm{d}x + z^2\,\mathrm{d}y + x^2\,\mathrm{d}z$，式中 $C$ 为维维亚尼曲线 $x^2+y^2+z^2 = a^2, x^2+y^2 = ax$ ($z \geqslant 0, a > 0$)，若从 $Ox$ 轴的正向 ($x > a$) 看去，沿逆时针方向为正．
\end{prob}

\begin{prob}[M]{4283}
$\displaystyle\int_C (y^2-z^2)\,\mathrm{d}x + (z^2-x^2)\,\mathrm{d}y + (x^2-y^2)\,\mathrm{d}z$，式中 $C$ 为球面的一部分 $x^2+y^2+z^2 = 1, x \geqslant 0, y \geqslant 0, z \geqslant 0$ 的边界，当沿着它的正向移动时，这部分球面的外侧面保持在左方．
\end{prob}

\probgroup{计算下列全微分的曲线积分：}

\begin{prob}[S]{4284}
$\displaystyle\int_{(1,1,1)}^{(2,3,-4)} x\,\mathrm{d}x + y^2\,\mathrm{d}y - z^3\,\mathrm{d}z$．
\end{prob}

\begin{prob}[S]{4285}
$\displaystyle\int_{(1,2,3)}^{(6,1,1)} yz\,\mathrm{d}x + xz\,\mathrm{d}y + xy\,\mathrm{d}z$．
\end{prob}

\begin{prob}[M]{4286}
$\displaystyle\int_{(x_1,y_1,z_1)}^{(x_2,y_2,z_2)} \frac{x\,\mathrm{d}x + y\,\mathrm{d}y + z\,\mathrm{d}z}{\sqrt{x^2+y^2+z^2}}$，其中点 $(x_1,y_1,z_1)$ 位于球面 $x^2+y^2+z^2 = a^2$ 之上，而点 $(x_2,y_2,z_2)$ 位于球面 $x^2+y^2+z^2 = b^2$ 之上 ($a > 0, b > 0$)．
\end{prob}

\begin{prob}[S]{4287}
$\displaystyle\int_{(x_1,y_1,z_1)}^{(x_2,y_2,z_2)} \varphi(x)\,\mathrm{d}x + \psi(y)\,\mathrm{d}y + \chi(z)\,\mathrm{d}z$，式中 $\varphi, \psi, \chi$ 为连续函数．
\end{prob}

\begin{prob}[S]{4288}
$\displaystyle\int_{(x_1,y_1,z_1)}^{(x_2,y_2,z_2)} f(x+y+z)(\mathrm{d}x + \mathrm{d}y + \mathrm{d}z)$，其中 $f$ 为连续函数．
\end{prob}

\begin{prob}[S]{4289}
$\displaystyle\int_{(x_1,y_1,z_1)}^{(x_2,y_2,z_2)} f(\sqrt{x^2+y^2+z^2})(x\,\mathrm{d}x + y\,\mathrm{d}y + z\,\mathrm{d}z)$，式中 $f$ 为连续函数．
\end{prob}

\probgroup{求原函数 $u$，若}

\begin{prob}[S]{4290}
$\mathrm{d}u = (x^2 - 2yz)\,\mathrm{d}x + (y^2 - 2xz)\,\mathrm{d}y + (z^2 - 2xy)\,\mathrm{d}z$．
\end{prob}

\begin{prob}[S]{4291}
$\mathrm{d}u = \left(1 - \frac{1}{y} + \frac{y}{z}\right)\,\mathrm{d}x + \left(\frac{x}{z} + \frac{x}{y^2}\right)\,\mathrm{d}y - \frac{xy}{z^2}\,\mathrm{d}z$．
\end{prob}

\begin{prob}[S]{4292}
$\mathrm{d}u = \frac{(x+y-z)\,\mathrm{d}x + (x+y-z)\,\mathrm{d}y + (x+y+z)\,\mathrm{d}z}{x^2+y^2+z^2+2xy}$．
\end{prob}

\begin{prob}[M]{4293}
求当质量为 $m$ 的点从位置 $(x_1,y_1,z_1)$ 移动到位置 $(x_2,y_2,z_2)$ 时，重力所作的功（$Oz$ 轴的方向是竖直向上）．
\end{prob}

\begin{prob}[M]{4294}
弹性力指向坐标原点，力的大小与质点到坐标原点的距离成正比．设此点依逆时针方向描绘出椭圆 $\frac{x^2}{a^2} + \frac{y^2}{b^2} = 1$ 的正四分之一，求弹性力所作的功．
\end{prob}

\begin{prob}[M]{4295}
当单位质量从点 $M_1(x_1,y_1,z_1)$ 移动到点 $M_2(x_2,y_2,z_2)$ 时，求引力 $F = \frac{G}{r^2}$（其中 $r = \sqrt{x^2+y^2+z^2}$）对它所作的功，其中 $G$ 是引力常量．
\end{prob}

\probsec{§12. 格林公式}

\begin{prob}[M]{4296}
利用格林公式变换曲线积分
\[
I = \oint_C \sqrt{x^2 + y^2}\,\mathrm{d}x + y \left[ xy + \ln\left( x + \sqrt{x^2 + y^2} \right) \right] \mathrm{d}y,
\]
式中围线 $C$ 是有界区域 $S$ 的边界.
\end{prob}

\begin{prob}[M]{4297}
应用格林公式，计算曲线积分
\[
I = \oint_K (x + y)^2\,\mathrm{d}x - (x^2 + y^2)\,\mathrm{d}y,
\]
其中 $K$ 是以 $A(1,1)$, $B(3,2)$, $C(2,5)$ 为顶点的三角形围线 $ABC$，并且积分时的环绕方向为正. 直接计算积分，以验证所求得的结果.
\end{prob}

\probgroup{应用格林公式，计算下列曲线积分：}

\begin{prob}[S]{4298}
$\oint_C xy^2\,\mathrm{d}y - x^2 y\,\mathrm{d}x$，式中 $C$ 为圆周 $x^2 + y^2 = a^2$.
\end{prob}

\begin{prob}[S]{4299}
$\oint_C (x + y)\,\mathrm{d}x - (x - y)\,\mathrm{d}y$，式中 $C$ 为椭圆 $\dfrac{x^2}{a^2} + \dfrac{y^2}{b^2} = 1$.
\end{prob}

\begin{prob}[M]{4300}
$\oint_C \mathrm{e}^x [(1 - \cos y)\,\mathrm{d}x - (y - \sin y)\,\mathrm{d}y]$，其中 $C$ 为区域 $0 < x < \pi$, $0 < y < \sin x$ 的边界，并且积分时的环绕方向为正.
\end{prob}

\begin{prob}[S]{4301}
$\oint_{x^2+y^2=R^2} \mathrm{e}^{-(x^2-y^2)} (\cos 2xy\,\mathrm{d}x + \sin 2xy\,\mathrm{d}y)$.
\end{prob}

\begin{prob}[M]{4302}
设 $AmB$ 为连接点 $A(1,1)$ 和点 $B(2,6)$ 的直线段，$AnB$ 是连接点 $A, B$ 及坐标原点的抛物线段，且该抛物线的轴垂直于 $x$ 轴. 积分
\[
I_1 = \int_{AmB} (x + y)^2\,\mathrm{d}x - (x - y)^2\,\mathrm{d}y, \quad I_2 = \int_{AnB} (x + y)^2\,\mathrm{d}x - (x - y)^2\,\mathrm{d}y
\]
相差多少？
\end{prob}

\begin{prob}[M]{4303}
计算曲线积分
\[
\int_{AmO} (\mathrm{e}^x \sin y - my)\,\mathrm{d}x + (\mathrm{e}^x \cos y - m)\,\mathrm{d}y,
\]
其中 $AmO$ 为由点 $A(a,0)$ 至点 $O(0,0)$ 的上半圆周 $x^2 + y^2 = ax$.
\probhint{用 $Ox$ 轴的直线段 $OA$ 与路径 $AmO$ 相连组成封闭围线.}
\end{prob}

\begin{prob}[M]{4304}
计算曲线积分
\[
\int_{AmB} [\varphi(y)\mathrm{e}^x - my]\,\mathrm{d}x + [\varphi'(y)\mathrm{e}^x - m]\,\mathrm{d}y,
\]
式中 $\varphi(y)$ 和 $\varphi'(y)$ 为连续函数，$AmB$ 为连接点 $A(x_1, y_1)$ 和点 $B(x_2, y_2)$ 的任意路径，但要求此路径与线段 $AB$ 一起围成具有已知面积 $S$ 的区域 $AmBA$.
\end{prob}

\begin{prob}[M]{4305}
求二阶可微的二个连续函数 $P(x,y)$ 和 $Q(x,y)$，使得曲线积分
\[
I = \oint_C P(x + \alpha, y + \beta)\,\mathrm{d}x + Q(x + \alpha, y + \beta)\,\mathrm{d}y
\]
对于任何封闭围线 $C$ 与常数 $\alpha$ 和 $\beta$ 无关.
\end{prob}

\begin{prob}[M]{4306}
为了使曲线积分
\[
\int_{AmB} F(x,y)(y\,\mathrm{d}x + x\,\mathrm{d}y)
\]
与积分路径的形状无关，可微函数 $F(x,y)$ 应满足怎样的条件？
\end{prob}

\begin{prob}[M]{4307}
计算
\[
I = \oint_C \frac{x\,\mathrm{d}y - y\,\mathrm{d}x}{x^2 + y^2},
\]
其中 $C$ 不经过坐标原点的简单封闭围线，且积分时的环绕方向为正.
\probhint{研究两种情况：1) 坐标原点在围线之外；2) 围线 $C$ 包围坐标原点.}
\end{prob}

\probgroup{利用曲线积分计算由下列曲线所围的面积：}

\begin{prob}[S]{4308}
$x = a\cos t$, $y = b\sin t$ ($0 \leqslant t \leqslant 2\pi$) (椭圆).
\end{prob}

\begin{prob}[S]{4309}
$x = a\cos^3 t$, $y = b\sin^3 t$ ($0 \leqslant t \leqslant 2\pi$) (星形线).
\end{prob}

\begin{prob}[S]{4310}
轴 $Ox$ 与 $(x + y)^2 = ax$ ($a > 0$) (抛物线).
\end{prob}

\begin{prob}[S]{4311}
$x^3 + y^3 = 3axy$ ($a > 0$) (笛卡儿叶形线).
\probhint{令 $y = tx$.}
\end{prob}

\begin{prob}[S]{4312}
$(x^2 + y^2)^2 = a^2(x^2 - y^2)$ (双纽线).
\probhint{令 $y = x\tan\varphi$.}
\end{prob}

\begin{prob}[S]{4313}
曲线 $x^3 + y^3 = x^2 + y^2$ 及坐标轴.
\end{prob}

\begin{prob}[M]{4314}
计算由曲线
\[
(x + y)^{n+m+1} = ax^n y^m \quad (a > 0, n > 0, m > 0)
\]
所围的面积.
\end{prob}

\begin{prob}[M]{4315}
计算由曲线
\[
\left( \frac{x}{a} \right)^n + \left( \frac{y}{b} \right)^n = 1 \quad (a > 0, b > 0, n > 0)
\]
和坐标轴所围的面积.
\probhint{令 $\dfrac{x}{a} = \cos^{\frac{2}{n}}\varphi$, $\dfrac{y}{b} = \sin^{\frac{2}{n}}\varphi$.}
\end{prob}

\begin{prob}[M]{4316}
计算由曲线
\[
\left( \frac{x}{a} \right)^n + \left( \frac{y}{b} \right)^n = \left( \frac{x}{a} \right)^{n-1} + \left( \frac{y}{b} \right)^{n-1} \quad (a > 0, b > 0, n > 1)
\]
和坐标轴所围的面积.
\end{prob}

\begin{prob}[M]{4317}
计算由曲线
\[
\left( \frac{x}{a} \right)^{2n+1} + \left( \frac{y}{b} \right)^{2n+1} = c \left( \frac{x}{a} \right)^n \left( \frac{y}{b} \right)^n \quad (a > 0, b > 0, c > 0, n > 0)
\]
所围的面积.
\end{prob}

\begin{prob}[M]{4318}
一半径为 $r$ 的圆周沿半径为 $R$ 的固定圆周外部滚动而无滑动，动圆周上的一点所描绘的曲线称为外摆线.

假定比值 $\dfrac{R}{r} = n$ 是整数 ($n \geqslant 1$)，求外摆线所围的面积.

研究特殊情况 $r = R$ (心脏线).
\end{prob}

\begin{prob}[M]{4319}
一个半径为 $r$ 的圆周沿半径为 $R$ 的固定圆周内部滚动而无滑动时，动圆周上的一点所描绘的曲线称为内摆线.

假定比值 $\dfrac{R}{r} = n$ 是整数 ($n \geqslant 2$)，求内摆线所围的面积.

研究特殊情况 $r = \dfrac{R}{4}$ (星形线).
\end{prob}

\begin{prob}[M]{4320.1}
计算圆柱面 $x^2 + y^2 = ax$ 被曲面 $x^2 + y^2 + z^2 = a^2$ 所截那部分的面积.
\end{prob}

\begin{prob}[M]{4320.2}
一简单封闭围线 $C$ 位于上半平面 $y \geqslant 0$，它在绕 $Ox$ 轴旋转时确定了一个旋转体，证明：此旋转体的体积为
\[
V = -\pi \oint_C y^2\,\mathrm{d}x.
\]
\end{prob}

\begin{prob}[M]{4321}
计算
\[
I = \frac{1}{2\pi} \oint_C \frac{X\,\mathrm{d}Y - Y\,\mathrm{d}X}{X^2 + Y^2},
\]
若 $X = ax + by, Y = cx + dy$，且 $C$ 为包围坐标原点的简单封闭围线 ($ad - bc \neq 0$).
\end{prob}

\begin{prob}[M]{4322}
若简单的围线 $C$ 包围坐标原点，$X = \varphi(x,y), Y = \psi(x,y)$，而曲线 $\varphi(x,y) = 0$ 和 $\psi(x,y) = 0$ 在围线 $C$ 以内有几个单交点，计算积分 $I$ (参阅前题).
\end{prob}

\begin{prob}[M]{4323}
证明：若 $C$ 为封闭围线，$\boldsymbol{l}$ 为任意的方向，则有
\[
\oint_C \cos(\boldsymbol{l}, \boldsymbol{n})\,\mathrm{d}s = 0,
\]
式中 $\boldsymbol{n}$ 为围线 $C$ 的外法向量.
\end{prob}

\begin{prob}[M]{4324}
求积分
\[
I = \oint_C [x\cos(\boldsymbol{n}, x) + y\cos(\boldsymbol{n}, y)]\,\mathrm{d}s
\]
之值，式中 $C$ 是有界区域 $S$ 的边界，它是简单封闭曲线，而 $\boldsymbol{n}$ 是它的外法向量.
\end{prob}

\begin{prob}[M]{4325}
求
\[
\lim_{d(S) \to 0} \frac{1}{S} \oint_C (\boldsymbol{F} \cdot \boldsymbol{n})\,\mathrm{d}s,
\]
其中 $S$ 为包围点 $(x_0, y_0)$ 的围线 $C$ 所围的面积，$d(S)$ 为区域 $S$ 的直径，$\boldsymbol{n}$ 为围线 $C$ 的单位外法向量，$\boldsymbol{F}\{X, Y\}$ 为 $S + C$ 上的连续可微向量.
\end{prob}

\probsec{§13. 曲线积分在物理学上的应用}

\begin{prob}[M]{4326}
均匀分布在圆 $x^2 + y^2 = a^2, y \geqslant 0$ 的上半部的质量 $M$ 以怎样的力吸引位于 $(0,0)$ 质量为 $m$ 的质点？
\end{prob}

\begin{prob}[M]{4327}
计算单层的对数势
\[
u(x,y) = \oint_C \kappa \ln \frac{1}{r}\,\mathrm{d}s,
\]
式中 $\kappa = \text{常数}$，$r = \sqrt{(\xi - x)^2 + (\eta - y)^2}$，设围线 $C$ 是圆周 $\xi^2 + \eta^2 = R^2$.
\end{prob}

\begin{prob}[M]{4328}
采用极坐标 $\rho$ 和 $\varphi$，计算单层的对数势
\[
I_1 = \int_0^{2\pi} \cos m\psi \ln \frac{1}{r}\,\mathrm{d}\psi \quad \text{和} \quad I_2 = \int_0^{2\pi} \sin m\psi \ln \frac{1}{r}\,\mathrm{d}\psi,
\]
式中 $r$ 为点 $(\rho, \varphi)$ 与动点 $(1, \psi)$ 间的距离，$m$ 为正整数.
\end{prob}

\begin{prob}[M]{4329}
计算高斯积分
\[
u(x,y) = \oint_C \frac{\cos(\boldsymbol{r}, \boldsymbol{n})}{r}\,\mathrm{d}s,
\]
式中 $r = \sqrt{(\xi - x)^2 + (\eta - y)^2}$ 为向量 $\boldsymbol{r}$ 的长度，此向量连接点 $A(x,y)$ 和简单封闭光滑围线 $C$ 上的动点 $M(\xi, \eta)$，$(\boldsymbol{r}, \boldsymbol{n})$ 为向量 $\boldsymbol{r}$ 与曲线 $C$ 在点 $M$ 的外法向量 $\boldsymbol{n}$ 之间的夹角.
\end{prob}

\begin{prob}[M]{4330}
采用极坐标 $\rho$ 和 $\varphi$，计算双层的对数势
\[
K_1 = \int_0^{2\pi} \cos m\psi \frac{\cos(\boldsymbol{r}, \boldsymbol{n})}{r}\,\mathrm{d}\psi, \quad K_2 = \int_0^{2\pi} \sin m\psi \frac{\cos(\boldsymbol{r}, \boldsymbol{n})}{r}\,\mathrm{d}\psi,
\]
式中 $r$ 为点 $A(\rho, \varphi)$ 和动点 $M(1, \psi)$ 之间的距离，$(\boldsymbol{r}, \boldsymbol{n})$ 为方向 $\overrightarrow{AM} = \boldsymbol{r}$ 与引自点 $O(0,0)$ 的半径 $\overrightarrow{OM} = \boldsymbol{n}$ 之间的夹角，$m$ 为正整数.
\end{prob}

\begin{prob}[M]{4331}
若 $\Delta u \equiv \dfrac{\partial^2 u}{\partial x^2} + \dfrac{\partial^2 u}{\partial y^2} = 0$，则称二阶可微函数 $u = u(x,y)$ 为调和函数. 证明：当且仅当以下条件成立时 $u$ 才是调和函数：
\[
\oint_C \frac{\partial u}{\partial n}\,\mathrm{d}s = 0,
\]
式中 $C$ 为任意封闭围线，$\dfrac{\partial u}{\partial n}$ 为沿此围线之外法线方向的导数.
\end{prob}

\begin{prob}[M]{4332}
证明：
\[
\iint_S \left[ \left( \frac{\partial u}{\partial x} \right)^2 + \left( \frac{\partial u}{\partial y} \right)^2 \right] \mathrm{d}x\,\mathrm{d}y = -\iint_S u \Delta u\,\mathrm{d}x\,\mathrm{d}y + \oint_C u \frac{\partial u}{\partial n}\,\mathrm{d}s,
\]
式中光滑围线 $C$ 是有界区域 $S$ 的边界.
\end{prob}

\begin{prob}[M]{4333}
证明：若一函数在有界区域 $S$ 内及其边界 $C$ 上为调和函数，则此函数单值地由它在边界 $C$ 上的值确定 (参考习题 4332).
\end{prob}

\begin{prob}[M]{4334}
证明平面上的格林第二公式：
\[
\iint_S \begin{vmatrix} \Delta u & \Delta v \\ u & v \end{vmatrix} \mathrm{d}x\,\mathrm{d}y = \oint_C \begin{vmatrix} \dfrac{\partial u}{\partial n} & \dfrac{\partial v}{\partial n} \\ u & v \end{vmatrix} \mathrm{d}s,
\]
式中光滑围线 $C$ 是有界区域 $S$ 的边界，$\dfrac{\partial}{\partial n}$ 为沿 $C$ 的外法线方向的导数.
\end{prob}

\begin{prob}[M]{4335}
利用格林第二公式证明：若 $u = u(x,y)$ 是有界闭区域 $S$ 内的调和函数，则
\[
u(x,y) = \frac{1}{2\pi} \oint_C \left( u \frac{\partial \ln r}{\partial n} - \ln r \frac{\partial u}{\partial n} \right) \mathrm{d}s,
\]
式中 $C$ 为区域 $S$ 的边界，$\boldsymbol{n}$ 为围线 $C$ 的外法向量，$(x,y)$ 为区域 $S$ 的内点，
\[
r = \sqrt{(\xi - x)^2 + (\eta - y)^2}
\]
为点 $(x,y)$ 与围线 $C$ 上的动点 $(\xi, \eta)$ 之间的距离.
\probhint{从区域 $S$ 中除去点 $(x,y)$ 与该点的无穷小圆形邻域，并对区域 $S$ 的剩余部分应用格林第二公式.}
\end{prob}

\begin{prob}[M]{4336}
证明对于调和函数 $u(M) = u(x,y)$ 的中值定理：
\[
u(M) = \frac{1}{2\pi} \oint_C u(\xi, \eta)\,\mathrm{d}s,
\]
式中 $C$ 是以点 $M$ 为中心的圆周.
\end{prob}

\begin{prob}[M]{4337}
证明：有界闭区域内的非常数调和函数 $u(x,y)$ 在此区域内的点不能达到其最大或最小值 (极大值原理).
\end{prob}

\begin{prob}[M]{4338}
证明黎曼公式：
\[
\iint_S \begin{vmatrix} L[u] & M[v] \\ u & v \end{vmatrix} \mathrm{d}x\,\mathrm{d}y = \oint_C P\,\mathrm{d}x + Q\,\mathrm{d}y,
\]
式中
\[
L[u] = \frac{\partial^2 u}{\partial x \partial y} + a \frac{\partial u}{\partial x} + b \frac{\partial u}{\partial y} + cu, \quad M[v] = \frac{\partial^2 v}{\partial x \partial y} - a \frac{\partial v}{\partial x} - b \frac{\partial v}{\partial y} + cv
\]
($a, b, c$ 为常数)，$P$ 和 $Q$ 为某些确定的函数，围线 $C$ 是有界区域 $S$ 的边界.
\end{prob}

\begin{prob}[M]{4339}
设 $u = u(x,y)$ 和 $v = v(x,y)$ 为定常流的速度分量，$C$ 为区域 $S$ 的边界，求区域 $S$ 内的流体质量的变化率. 若流体是不可压缩的，且在区域 $S$ 内没有源和汇，则函数 $u$ 和 $v$ 满足怎样的方程？
\end{prob}

\begin{prob}[M]{4340}
根据毕奥–萨瓦尔定律，通过导线元 $\mathrm{d}\boldsymbol{s}$ 的电流 $i$ 在空间的点 $M(x,y,z)$ 处所对应的磁场强度为
\[
\mathrm{d}\boldsymbol{H} = ki \frac{(\boldsymbol{r} \times \mathrm{d}\boldsymbol{s})}{r^3},
\]
其中 $\boldsymbol{r}$ 为连接导线元 $\mathrm{d}\boldsymbol{s}$ 与点 $M$ 的向量，$k$ 为比例系数.

对于封闭导线 $C$ 的情形，求点 $M$ 的磁场强度 $\boldsymbol{H}$ 的投影 $H_x, H_y$ 和 $H_z$.
\end{prob}

\probsec{§14. 曲面积分}

\begin{prob}[M]{4341}
两个积分
\[
I_1 = \iint\limits_S (x^2 + y^2 + z^2)\,\mathrm{d}S, \quad I_2 = \iint\limits_P (x^2 + y^2 + z^2)\,\mathrm{d}P
\]
(式中 $S$ 为球面 $x^2 + y^2 + z^2 = a^2$，$P$ 为其内接八面体的表面 $|x| + |y| + |z| = a$) 相差若何？
\end{prob}

\begin{prob}[M]{4342}
计算 $\iint\limits_S z\,\mathrm{d}S$，式中 $S$ 为曲面 $x^2 + z^2 = 2az$ ($a > 0$) 被曲面 $z = \sqrt{x^2 + y^2}$ 所割下的部分.
\end{prob}

\probgroup{计算下列第一型曲面积分：}

\begin{prob}[S]{4343}
$\iint\limits_S (x + y + z)\,\mathrm{d}S$，式中 $S$ 为曲面 $x^2 + y^2 + z^2 = a^2, z \geqslant 0$.
\end{prob}

\begin{prob}[S]{4344}
$\iint\limits_S (x^2 + y^2)\,\mathrm{d}S$，式中 $S$ 为区域 $\sqrt{x^2 + y^2} \leqslant z \leqslant 1$ 的边界.
\end{prob}

\begin{prob}[S]{4345}
$\iint\limits_S \frac{\,\mathrm{d}S}{(1 + x + y)^2}$，式中 $S$ 为四面体 $x + y + z \leqslant 1, x \geqslant 0, y \geqslant 0, z \geqslant 0$ 的边界.
\end{prob}

\begin{prob}[S]{4346}
$\iint\limits_S |xyz|\,\mathrm{d}S$，式中 $S$ 为曲面 $z = x^2 + y^2$ 被平面 $z = 1$ 所截下的部分.
\end{prob}

\begin{prob}[M]{4347}
$\iint\limits_S \frac{\,\mathrm{d}S}{h}$，式中 $S$ 为椭球面，$h$ 为椭球中心到与椭球面微元 $\,\mathrm{d}S$ 相切的平面的距离.
\end{prob}

\begin{prob}[M]{4348}
$\iint\limits_S z\,\mathrm{d}S$，式中 $S$ 为螺旋面的一部分：$x = u\cos v, y = u\sin v, z = v$ ($0 < u < a; 0 < v < 2\pi$).
\end{prob}

\begin{prob}[M]{4349}
$\iint\limits_S z^2\,\mathrm{d}S$，式中 $S$ 为圆锥面的一部分：$x = r\cos\varphi\sin\alpha, y = r\sin\varphi\sin\alpha, z = r\cos\alpha$ ($0 \leqslant r \leqslant a; 0 \leqslant \varphi \leqslant 2\pi$)，$\alpha$ 为常数 $\left(0 < \alpha < \frac{\pi}{2}\right)$.
\end{prob}

\begin{prob}[M]{4350}
$\iint\limits_S (xy + yz + zx)\,\mathrm{d}S$，式中 $S$ 为圆锥面 $z = \sqrt{x^2 + y^2}$ 被曲面 $x^2 + y^2 = 2ax$ 所割下的部分.
\end{prob}

\begin{prob}[M]{4351}
证明泊松公式：
\[
\iint\limits_S f(ax + by + cz)\,\mathrm{d}S = 2\pi \int_{-1}^{1} f(u\sqrt{a^2 + b^2 + c^2})\,\mathrm{d}u,
\]
式中 $S$ 是球面 $x^2 + y^2 + z^2 = 1$.
\end{prob}

\begin{prob}[M]{4352.1}
求抛物面壳
\[
z = \frac{1}{2}(x^2 + y^2) \quad (0 \leqslant z \leqslant 1)
\]
的质量，此壳的面密度按规律 $\rho = z$ 而变.
\end{prob}

\begin{prob}[M]{4352.2}
求半球面
\[
x^2 + y^2 + z^2 = a^2 \quad (z \geqslant 0)
\]
的质量，其面密度在半球面上每一点 $M(x,y,z)$ 等于 $\frac{x}{a}$.
\end{prob}

\begin{prob}[M]{4352.3}
求均质三角形板
\[
x + y + z = a \quad (x \geqslant 0, y \geqslant 0, z \geqslant 0)
\]
对坐标平面的静矩.
\end{prob}

\begin{prob}[M]{4353}
求密度为 $\rho_0$ 的均质球面壳
\[
x^2 + y^2 + z^2 = a^2 \quad (z \geqslant 0)
\]
对 $Oz$ 轴的转动惯量.
\end{prob}

\begin{prob}[M]{4354}
求密度为 $\rho_0$ 的均质锥面壳
\[
\frac{x^2}{a^2} + \frac{y^2}{a^2} - \frac{z^2}{b^2} = 0 \quad (0 \leqslant z \leqslant b)
\]
对直线
\[
\frac{x}{1} = \frac{y}{0} = \frac{z - b}{0}
\]
的转动惯量.
\end{prob}

\begin{prob}[M]{4355}
求下列均质曲面的质心坐标：
\begin{enumerate}[label=(\alph*), itemsep=2pt, topsep=3pt]
    \item 曲面 $z = \sqrt{x^2 + y^2}$ 被曲面 $x^2 + y^2 = ax$ 所割下的部分；
    \item 曲面 $z = \sqrt{a^2 - x^2 - y^2}$ ($x \geqslant 0; y \geqslant 0; x + y \leqslant a$).
\end{enumerate}
\end{prob}

\begin{prob}[M]{4356.1}
求下列曲面 $S$ 的极转动惯量 $I_0 = \iint\limits_S (x^2 + y^2 + z^2)\,\mathrm{d}S$：
\begin{enumerate}[label=(\alph*), itemsep=2pt, topsep=3pt]
    \item 立方体表面 $\max\{|x|, |y|, |z|\} = a$；
    \item 圆柱体 $x^2 + y^2 \leqslant R^2, 0 \leqslant z \leqslant H$ 的全部表面.
\end{enumerate}
\end{prob}

\begin{prob}[M]{4356.2}
求均质三角形板
\[
x + y + z = 1 \quad (x \geqslant 0; y \geqslant 0; z \geqslant 0)
\]
对坐标平面的转动惯量.
\end{prob}

\begin{prob}[M]{4357}
密度为 $\rho_0$ 的均质截圆锥面
\[
x = r\cos\varphi, \quad y = r\sin\varphi, \quad z = r \quad (0 \leqslant \varphi \leqslant 2\pi, 0 < b \leqslant r \leqslant a)
\]
以怎样的力吸引质量为 $m$ 位于该圆锥面顶点的质点？
\end{prob}

\begin{prob}[M]{4358}
求密度为 $\rho_0$ 的均质球面 $S: x^2 + y^2 + z^2 = a^2$ 在点 $M_0(x_0, y_0, z_0)$ 的势，即计算积分
\[
u = \iint\limits_S \frac{\rho_0\,\mathrm{d}S}{r},
\]
式中 $r = \sqrt{(x - x_0)^2 + (y - y_0)^2 + (z - z_0)^2}$.
\end{prob}

\begin{prob}[L]{4359}
计算
\[
F(t) = \iint\limits_{x + y + z = t} f(x, y, z)\,\mathrm{d}S,
\]
式中
\[
f(x, y, z) = \begin{cases} 1 - x^2 - y^2 - z^2, & x^2 + y^2 + z^2 \leqslant 1, \\ 0, & x^2 + y^2 + z^2 > 1. \end{cases}
\]
作出函数 $u = F(t)$ 的图像.
\end{prob}

\begin{prob}[M]{4360}
计算积分
\[
F(t) = \iint\limits_{x^2 + y^2 + z^2 = t^2} f(x, y, z)\,\mathrm{d}S,
\]
式中
\[
f(x, y, z) = \begin{cases} x^2 + y^2, & z \geqslant \sqrt{x^2 + y^2}, \\ 0, & z < \sqrt{x^2 + y^2}. \end{cases}
\]
\end{prob}

\begin{prob}[M]{4361}
计算积分
\[
F(x, y, z, t) = \iint\limits_S f(\xi, \eta, \zeta)\,\mathrm{d}S,
\]
式中 $S$ 为变球面
\[
(\xi - x)^2 + (\eta - y)^2 + (\zeta - z)^2 = t^2,
\]
\[
f(\xi, \eta, \zeta) = \begin{cases} 1, & \xi^2 + \eta^2 + \zeta^2 < a^2, \\ 0, & \xi^2 + \eta^2 + \zeta^2 \geqslant a^2, \end{cases}
\]
且假设 $r = \sqrt{x^2 + y^2 + z^2} > a > 0$.
\end{prob}

\probgroup{计算下列第二型曲面积分：}

\begin{prob}[S]{4362}
$\iint\limits_S (x\,\mathrm{d}y\,\mathrm{d}z + y\,\mathrm{d}z\,\mathrm{d}x + z\,\mathrm{d}x\,\mathrm{d}y)$，式中 $S$ 为球面 $x^2 + y^2 + z^2 = a^2$ 的外侧.
\end{prob}

\begin{prob}[S]{4363}
$\iint\limits_S f(x)\,\mathrm{d}y\,\mathrm{d}z + g(y)\,\mathrm{d}z\,\mathrm{d}x + h(z)\,\mathrm{d}x\,\mathrm{d}y$，式中 $f(x), g(y), h(z)$ 为连续函数，$S$ 为平行六面体 $0 \leqslant x \leqslant a; 0 \leqslant y \leqslant b; 0 \leqslant z \leqslant c$ 的外表面.
\end{prob}

\begin{prob}[S]{4364}
$\iint\limits_S (y - z)\,\mathrm{d}y\,\mathrm{d}z + (z - x)\,\mathrm{d}z\,\mathrm{d}x + (x - y)\,\mathrm{d}x\,\mathrm{d}y$，式中 $S$ 为圆锥面 $x^2 + y^2 = z^2$ ($0 \leqslant z \leqslant h$) 的外侧.
\end{prob}

\begin{prob}[S]{4365}
$\iint\limits_S \left(\frac{\,\mathrm{d}y\,\mathrm{d}z}{x} + \frac{\,\mathrm{d}z\,\mathrm{d}x}{y} + \frac{\,\mathrm{d}x\,\mathrm{d}y}{z}\right)$，式中 $S$ 为椭球面 $\frac{x^2}{a^2} + \frac{y^2}{b^2} + \frac{z^2}{c^2} = 1$ 的外侧.
\end{prob}

\begin{prob}[S]{4366}
$\iint\limits_S x^2\,\mathrm{d}y\,\mathrm{d}z + y^2\,\mathrm{d}z\,\mathrm{d}x + z^2\,\mathrm{d}x\,\mathrm{d}y$，式中 $S$ 为球面 $(x - a)^2 + (y - b)^2 + (z - c)^2 = R^2$ 的外侧.
\end{prob}

\probsec{§15. 斯托克斯公式}

\begin{prob}[M]{4367}
应用斯托克斯公式，计算曲线积分
\[
\oint\limits_C y\,\mathrm{d}x + z\,\mathrm{d}y + x\,\mathrm{d}z,
\]
其中 $C$ 为圆周 $x^2 + y^2 + z^2 = a^2, x + y + z = 0$，并且从 $Ox$ 轴的正向来看，积分时此圆周的环绕方向是逆时针的. 用直接计算法检验结果.
\end{prob}

\begin{prob}[M]{4368}
计算积分
\[
\oint\limits_{AmB} (x^2 - yz)\,\mathrm{d}x + (y^2 - xz)\,\mathrm{d}y + (z^2 - xy)\,\mathrm{d}z,
\]
此积分是从点 $A(a, 0, 0)$ 至点 $B(a, 0, h)$ 沿着螺旋线 $x = a\cos\varphi, y = a\sin\varphi, z = \frac{h}{2\pi}\varphi$ 进行的.
\probhint{以直接段与曲线 $AmB$ 相连并应用斯托克斯公式.}
\end{prob}

\begin{prob}[M]{4369}
设 $C$ 为平面 $x\cos\alpha + y\cos\beta + z\cos\gamma - p = 0$ ($\cos\alpha, \cos\beta, \cos\gamma$ 为平面之法线的方向余弦) 上的封闭围线，所围面积为 $S$，求
\[
\oint\limits_C \begin{vmatrix} \,\mathrm{d}x & \,\mathrm{d}y & \,\mathrm{d}z \\ \cos\alpha & \cos\beta & \cos\gamma \\ x & y & z \end{vmatrix},
\]
积分沿围线 $C$ 的正方向进行.
\end{prob}

\probgroup{应用斯托克斯公式，计算积分：}

\begin{prob}[M]{4370}
$\oint\limits_C (y + z)\,\mathrm{d}x + (z + x)\,\mathrm{d}y + (x + y)\,\mathrm{d}z$，式中 $C$ 为椭圆周 $x = a\sin^2 t, y = 2a\sin t\cos t, z = a\cos^2 t$ ($0 \leqslant t \leqslant \pi$)，并且以参数 $t$ 增大的方向为积分时的正方向.
\end{prob}

\begin{prob}[M]{4371}
$\oint\limits_C (y - z)\,\mathrm{d}x + (z - x)\,\mathrm{d}y + (x - y)\,\mathrm{d}z$，式中 $C$ 为椭圆周 $x^2 + y^2 = a^2, \frac{x}{a} + \frac{y}{h} = 1$ ($a > 0, h > 0$). 并且若从 $Ox$ 轴正向看去，积分是沿此椭圆依逆时针方向进行的.
\end{prob}

\begin{prob}[M]{4372}
$\oint\limits_C (y^2 + z^2)\,\mathrm{d}x + (x^2 + z^2)\,\mathrm{d}y + (x^2 + y^2)\,\mathrm{d}z$，式中 $C$ 是曲线 $x^2 + y^2 + z^2 = 2Rx, x^2 + y^2 = 2rx$ ($0 < r < R, z > 0$)，并且在沿此曲线进行积分时，球面 $x^2 + y^2 + z^2 = 2Rx$ 外侧被该曲线所围的最小区域始终位于左边.
\end{prob}

\begin{prob}[M]{4373}
$\oint\limits_C (y^2 - z^2)\,\mathrm{d}x + (z^2 - x^2)\,\mathrm{d}y + (x^2 - y^2)\,\mathrm{d}z$，式中 $C$ 为用平面 $x + y + z = \frac{3}{2}a$ 截立方体 $0 \leqslant x \leqslant a, 0 \leqslant y \leqslant a, 0 \leqslant z \leqslant a$ 所得截面的边界，并且若从 $Ox$ 轴的正向看去，积分是沿 $C$ 依逆时针方向进行的.
\end{prob}

\begin{prob}[M]{4374}
$\oint\limits_C y^2z^2\,\mathrm{d}x + z^2x^2\,\mathrm{d}y + x^2y^2\,\mathrm{d}z$，式中 $C$ 为封闭曲线 $x = a\cos t, y = a\cos 2t, z = a\cos 3t$，并且以参数 $t$ 增大的方向为积分时的正方向.
\end{prob}

\begin{prob}[M]{4375}
有函数
\[
W(x, y, z) = ki \iint\limits_S \frac{\cos(\boldsymbol{r}, \boldsymbol{n})}{r^2}\,\mathrm{d}S \quad (k = \text{常数}).
\]
其中曲面 $S$ 的边界为围线 $C$，$\boldsymbol{n}$ 为曲面 $S$ 的法向量，$\boldsymbol{r}$ 为连接空间的点 $M(x, y, z)$ 与围线 $C$ 上的动点 $A(\xi, \eta, \zeta)$ 所成之径向量. 证明：此函数为通过围线 $C$ 的电流 $i$ 所产生磁场 $\boldsymbol{H}$ 的势 (参阅习题 4340).
\end{prob}

\probsec{§16. 奥斯特罗格拉茨基公式}

\probgroup{应用奥斯特罗格拉茨基公式变换下列曲面积分，设光滑曲面 $S$ 是有界区域 $V$ 的边界，$\cos\alpha, \cos\beta, \cos\gamma$ 为曲面 $S$ 的外法线的方向余弦：}

\begin{prob}[S]{4376}
\[\iint_{S} x^{3}\,\mathrm{d}y\,\mathrm{d}z + y^{3}\,\mathrm{d}z\,\mathrm{d}x + z^{3}\,\mathrm{d}x\,\mathrm{d}y.\]
\end{prob}

\begin{prob}[S]{4377}
\[\iint_{S} yz\,\mathrm{d}y\,\mathrm{d}z + zx\,\mathrm{d}z\,\mathrm{d}x + xy\,\mathrm{d}x\,\mathrm{d}y.\]
\end{prob}

\begin{prob}[S]{4378}
\[\iint_{S} \frac{x\cos\alpha + y\cos\beta + z\cos\gamma}{\sqrt{x^{2}+y^{2}+z^{2}}}\,\mathrm{d}S.\]
\end{prob}

\begin{prob}[S]{4379}
\[\iint_{S} \left(\frac{\partial u}{\partial x}\cos\alpha + \frac{\partial u}{\partial y}\cos\beta + \frac{\partial u}{\partial z}\cos\gamma\right)\mathrm{d}S.\]
\end{prob}

\begin{prob}[S]{4380}
\[\iint_{S} \left[\left(\frac{\partial R}{\partial y} - \frac{\partial Q}{\partial z}\right)\cos\alpha + \left(\frac{\partial P}{\partial z} - \frac{\partial R}{\partial x}\right)\cos\beta + \left(\frac{\partial Q}{\partial x} - \frac{\partial P}{\partial y}\right)\cos\gamma\right]\mathrm{d}S.\]
\end{prob}

\begin{prob}[M]{4381}
证明：若 $S$ 为封闭的简单曲面而 $\boldsymbol{l}$ 为任意的固定方向，则
\[\iint_{S} \cos(\boldsymbol{n}, \boldsymbol{l})\,\mathrm{d}S = 0,\]
式中 $\boldsymbol{n}$ 为曲面 $S$ 的外法向量。
\end{prob}

\begin{prob}[M]{4382}
证明：以曲面 $S$ 为界的物体的体积等于：
\[V = \frac{1}{3} \iint_{S} (x\cos\alpha + y\cos\beta + z\cos\gamma)\,\mathrm{d}S,\]
式中 $\cos\alpha, \cos\beta, \cos\gamma$ 为曲面 $S$ 的外法线的方向余弦。
\end{prob}

\begin{prob}[M]{4383}
证明：以光滑锥面 $F(x,y,z)=0$ 和平面 $Ax+By+Cz+D=0$ 为界的锥体的体积等于
\[V = \frac{1}{3}SH,\]
式中 $S$ 为位于该平面上的锥底之面积，$H$ 为锥体的高。
\end{prob}

\begin{prob}[M]{4384}
求以曲面 $z = \pm c$ 及
\[\begin{cases} x = a\cos u\cos v + b\sin u\sin v, \\ y = a\cos u\sin v - b\sin u\cos v, \\ z = c\sin u \end{cases}\]
为界的物体的体积。
\end{prob}

\begin{prob}[M]{4385.1}
求以曲面 $x = u\cos v, y = u\sin v, z = -u + a\cos v\ (u \geqslant 0)$ 及平面 $x = 0, x = a\ (a > 0)$ 为界的物体的体积。
\end{prob}

\begin{prob}[M]{4385.2}
求环面
\[\begin{cases} x = (b + a\cos\psi)\cos\varphi, \\ y = (b + a\cos\psi)\sin\varphi, \quad (0 < a \leqslant b) \\ z = a\sin\psi \end{cases}\]
所围区域的体积。
\end{prob}

\begin{prob}[M]{4386}
证明公式：
\[\begin{aligned}
&\frac{\mathrm{d}}{\mathrm{d}t}\left\{\iiint_{x^{2}+y^{2}+z^{2}\leqslant t^{2}} f(x,y,z,t)\,\mathrm{d}x\,\mathrm{d}y\,\mathrm{d}z\right\} \\
&= \iint_{x^{2}+y^{2}+z^{2}=t^{2}} f(x,y,z,t)\,\mathrm{d}S + \iiint_{x^{2}+y^{2}+z^{2}\leqslant t^{2}} \frac{\partial f}{\partial t}\,\mathrm{d}x\,\mathrm{d}y\,\mathrm{d}z \quad (t > 0).
\end{aligned}\]
\end{prob}

\probgroup{利用奥斯特罗格拉茨基公式计算下列曲面积分：}

\begin{prob}[S]{4387}
$\displaystyle \iint_{S} x^{2}\,\mathrm{d}y\,\mathrm{d}z + y^{2}\,\mathrm{d}z\,\mathrm{d}x + z^{2}\,\mathrm{d}x\,\mathrm{d}y$，式中 $S$ 为立方体 $0 \leqslant x \leqslant a, 0 \leqslant y \leqslant a, 0 \leqslant z \leqslant a$ 的外表面。
\end{prob}

\begin{prob}[S]{4388}
$\displaystyle \iint_{S} x^{3}\,\mathrm{d}y\,\mathrm{d}z + y^{3}\,\mathrm{d}z\,\mathrm{d}x + z^{3}\,\mathrm{d}x\,\mathrm{d}y$，式中 $S$ 为球 $x^{2}+y^{2}+z^{2}=a^{2}$ 的外表面。
\end{prob}

\begin{prob}[M]{4389}
$\displaystyle \iint_{S} (x-y+z)\,\mathrm{d}y\,\mathrm{d}z + (y-z+x)\,\mathrm{d}z\,\mathrm{d}x + (z-x+y)\,\mathrm{d}x\,\mathrm{d}y$，式中 $S$ 为曲面 $|x-y+z| + |y-z+x| + |z-x+y| = 1$ 的外侧。
\end{prob}

\begin{prob}[M]{4390}
计算
\[\iint_{S} (x^{2}\cos\alpha + y^{2}\cos\beta + z^{2}\cos\gamma)\,\mathrm{d}S,\]
式中 $S$ 为部分圆锥面 $x^{2}+y^{2}=z^{2}\ (0 \leqslant z \leqslant h)$，$\cos\alpha, \cos\beta, \cos\gamma$ 为此曲面外法线的方向余弦。
\probhint{与部分平面 $z = h, x^{2}+y^{2}\leqslant h^{2}$ 组成封闭曲面。}
\end{prob}

\begin{prob}[M]{4391}
证明公式：
\[\iiint_{V} \frac{\mathrm{d}\xi\,\mathrm{d}\eta\,\mathrm{d}\zeta}{r} = \frac{1}{2} \iint_{S} \cos(\boldsymbol{r}, \boldsymbol{n})\,\mathrm{d}S,\]
其中封闭曲面 $S$ 为区域 $V$ 的表面，$\boldsymbol{n}$ 为曲面 $S$ 上的点 $(\xi, \eta, \zeta)$ 处的外法向量，$r = \sqrt{(\xi-x)^{2}+(\eta-y)^{2}+(\zeta-z)^{2}}$，$\boldsymbol{r}$ 为从点 $(x,y,z)$ 到点 $(\xi, \eta, \zeta)$ 的径向量。
\end{prob}

\begin{prob}[M]{4392}
计算高斯积分
\[I(x,y,z) = \iint_{S} \frac{\cos(\boldsymbol{r}, \boldsymbol{n})}{r^{2}}\,\mathrm{d}S,\]
式中 $S$ 为简单封闭光滑曲面，它是区域 $V$ 的边界，$\boldsymbol{n}$ 为曲面 $S$ 上在点 $(\xi, \eta, \zeta)$ 处的外法向量，$\boldsymbol{r}$ 为连接点 $(x,y,z)$ 和点 $(\xi, \eta, \zeta)$ 的径向量，
\[r = \sqrt{(\xi-x)^{2}+(\eta-y)^{2}+(\zeta-z)^{2}}.\]
研究两种情形：

(a) 曲面 $S$ 不包围点 $(x,y,z)$；\qquad (b) 曲面 $S$ 包围点 $(x,y,z)$。
\end{prob}

\begin{prob}[L]{4393}
证明：若
\[\Delta u \equiv \frac{\partial^{2} u}{\partial x^{2}} + \frac{\partial^{2} u}{\partial y^{2}} + \frac{\partial^{2} u}{\partial z^{2}},\]
有界区域 $V$ 的边界 $S$ 为光滑曲面，则成立下列公式：
\begin{enumerate}[label=(\alph*), itemsep=2pt, topsep=3pt]
\item $\displaystyle \iint_{S} \frac{\partial u}{\partial n}\,\mathrm{d}S = \iiint_{V} \Delta u\,\mathrm{d}x\,\mathrm{d}y\,\mathrm{d}z$；
\item $\displaystyle \iint_{S} u \frac{\partial u}{\partial n}\,\mathrm{d}S = \iiint_{V} \left[\left(\frac{\partial u}{\partial x}\right)^{2} + \left(\frac{\partial u}{\partial y}\right)^{2} + \left(\frac{\partial u}{\partial z}\right)^{2}\right]\mathrm{d}x\,\mathrm{d}y\,\mathrm{d}z + \iiint_{V} u \Delta u\,\mathrm{d}x\,\mathrm{d}y\,\mathrm{d}z$，
\end{enumerate}
式中 $u$ 及其一阶和二阶偏导数是在区域 $V+S$ 内连续的函数，$\dfrac{\partial u}{\partial n}$ 为沿曲面 $S$ 的外法线的导数。
\end{prob}

\begin{prob}[M]{4394}
证明三维情形的格林第二公式：
\[\iiint_{V} \begin{vmatrix} \Delta u & \Delta v \\ u & v \end{vmatrix} \mathrm{d}x\,\mathrm{d}y\,\mathrm{d}z = \iint_{S} \begin{vmatrix} \dfrac{\partial u}{\partial n} & \dfrac{\partial v}{\partial n} \\ u & v \end{vmatrix} \mathrm{d}S,\]
式中区域 $V$ 以曲面 $S$ 为界，$\boldsymbol{n}$ 是曲面 $S$ 的外法向量，而函数 $u = u(x,y,z), v = v(x,y,z)$ 在区域 $V+S$ 内二阶可微。
\end{prob}

\begin{prob}[L]{4395}
设函数 $u = u(x,y,z)$ 在某区域内具有连续的一阶和二阶导数，若
\[\Delta u \equiv \frac{\partial^{2} u}{\partial x^{2}} + \frac{\partial^{2} u}{\partial y^{2}} + \frac{\partial^{2} u}{\partial z^{2}} = 0,\]
则 $u(x,y,z)$ 称为此区域内的调和函数。

证明：若有界闭区域 $V$ 以光滑曲面 $S$ 为界，$u$ 是此区域内的调和函数，则成立下列公式：
\begin{enumerate}[label=(\alph*), itemsep=2pt, topsep=3pt]
\item $\displaystyle \iint_{S} \frac{\partial u}{\partial n}\,\mathrm{d}S = 0$；
\item $\displaystyle \iiint_{V} \left[\left(\frac{\partial u}{\partial x}\right)^{2} + \left(\frac{\partial u}{\partial y}\right)^{2} + \left(\frac{\partial u}{\partial z}\right)^{2}\right]\mathrm{d}x\,\mathrm{d}y\,\mathrm{d}z = \iint_{S} u \frac{\partial u}{\partial n}\,\mathrm{d}S$，
\end{enumerate}
式中 $\boldsymbol{n}$ 为曲面 $S$ 的外法向量。

用公式 (b) 证明：区域 $V$ 内的调和函数由它在边界 $S$ 上的值唯一地确定。
\end{prob}

\begin{prob}[M]{4396}
证明：若函数 $u = u(x,y,z)$ 是以光滑曲面 $S$ 为界的有界闭区域 $V$ 内的调和函数，则
\[u(x,y,z) = \frac{1}{4\pi} \iint_{S} \left[u \frac{\cos(\boldsymbol{r}, \boldsymbol{n})}{r^{2}} + \frac{1}{r}\frac{\partial u}{\partial n}\right]\mathrm{d}S,\]
式中 $\boldsymbol{r}$ 是从区域 $V$ 的内点 $(x,y,z)$ 引至曲面 $S$ 上的点 $(\xi, \eta, \zeta)$ 的径向量，$r = \sqrt{(\xi-x)^{2}+(\eta-y)^{2}+(\zeta-z)^{2}}$，$\boldsymbol{n}$ 为曲面 $S$ 在点 $(\xi, \eta, \zeta)$ 的外法向量。
\end{prob}

\begin{prob}[M]{4397}
证明：若半径为 $R$ 的球的球心位于点 $(x_{0}, y_{0}, z_{0})$，$u = u(x,y,z)$ 为此球内的调和函数，则
\[u(x_{0}, y_{0}, z_{0}) = \frac{1}{4\pi R^{2}} \iint_{S} u(x,y,z)\,\mathrm{d}S\]
（中值定理）。
\end{prob}

\begin{prob}[M]{4398}
证明：若有界闭区域 $V$ 内的连续函数 $u = u(x,y,z)$ 在此区域内部是调和函数，并且它不是常数，则此函数在区域内的点不能达到最大值和最小值（极大值原理）。
\end{prob}

\begin{prob}[M]{4399}
设物体 $V$ 全部浸于液体中，利用帕斯卡定律证明：液体的浮力等于物体所排开的液体的重量，其方向是竖直向上（阿基米德定律）。
\end{prob}

\begin{prob}[L]{4400}
设 $S_{t}$ 是动球面 $(\xi-x)^{2}+(\eta-y)^{2}+(\zeta-z)^{2}=t^{2}$，而函数 $f(\xi, \eta, \zeta)$ 是连续的，证明：函数
\[u(x,y,z,t) = \frac{1}{4\pi} \iint_{S_{t}} \frac{f(\xi, \eta, \zeta)}{t}\,\mathrm{d}S_{t}\]
满足波动方程
\[\frac{\partial^{2} u}{\partial x^{2}} + \frac{\partial^{2} u}{\partial y^{2}} + \frac{\partial^{2} u}{\partial z^{2}} = \frac{\partial^{2} u}{\partial t^{2}}\]
和初始条件 $\left. u \right|_{t=0} = 0, \left. \dfrac{\partial u}{\partial t} \right|_{t=0} = f(x,y,z)$。
\probhint{用三重积分表示导数 $\dfrac{\partial u}{\partial t}$。}
\end{prob}

\probsec{§17. 场论初步}

\begin{prob}[L]{4401.1}
求场 $u = x^2 + 2y^2 + 3z^2 + xy + 3x - 2y - 6z$ 在下列各点的梯度的大小和方向：
\begin{enumerate}[label=(\alph*), itemsep=2pt, topsep=3pt]
  \item $O(0, 0, 0)$；
  \item $A(1, 1, 1)$；
  \item $B(2, 1, 1)$.
\end{enumerate}
在场中怎样的点，梯度等于零？
\end{prob}

\begin{prob}[M]{4401.2}
设
\[
u = xy - z^2,
\]
求 $\operatorname{grad} u$ 在点 $M(-9, 12, 10)$ 的大小和方向. 沿坐标角 $xOy$ 的等分线方向的导数 $\dfrac{\partial u}{\partial l}$ 等于什么？
\end{prob}

\begin{prob}[L]{4402}
在空间 $Oxyz$ 的哪些点，场
\[
u = x^3 + y^3 + z^3 - 3xyz
\]
的梯度：
\begin{enumerate}[label=(\alph*), itemsep=2pt, topsep=3pt]
  \item 垂直于 $Oz$ 轴；
  \item 平行于 $Oz$ 轴；
  \item 等于零？
\end{enumerate}
\end{prob}

\begin{prob}[M]{4403}
给定标量场
\[
u = \ln \frac{1}{r},
\]
其中 $r = \sqrt{(x - a)^2 + (y - b)^2 + (z - c)^2}$. 在空间 $Oxyz$ 的哪些点成立等式
\[
|\operatorname{grad} u| = 1\text{？}
\]
\end{prob}

\begin{prob}[M]{4404}
作标量场
\[
u = \sqrt{x^2 + y^2 + (z + 8)^2} + \sqrt{x^2 + y^2 + (z - 8)^2}
\]
的等值面. 求通过点 $M(9, 12, 28)$ 的等值面. 在区域 $x^2 + y^2 + z^2 \leqslant 36$ 内 $\max u$ 等于什么？
\end{prob}

\begin{prob}[M]{4405}
求场
\[
u = \frac{x}{x^2 + y^2 + z^2}
\]
在点 $A(1, 2, 2)$ 及 $B(-3, 1, 0)$ 的梯度之间的夹角 $\varphi$.
\end{prob}

\begin{prob}[M]{4406}
设已知标量场
\[
u = \frac{z}{\sqrt{x^2 + y^2 + z^2}},
\]
作出场的等值面和梯度的等模面.
求区域 $1 < z < 2$ 内的 $\inf u$, $\sup u$, $\inf |\operatorname{grad} u|$, $\sup |\operatorname{grad} u|$.
\end{prob}

\begin{prob}[M]{4407}
精确到高阶无穷小量，求在点 $M_0(x_0, y_0, z_0)$ 处之二无限接近的等值面
\[
u(x, y, z) = c \quad \text{及} \quad u(x, y, z) = c + \Delta c
\]
之间的距离，其中 $u(x_0, y_0, z_0) = c$ ($\operatorname{grad} u(x_0, y_0, z_0) \neq 0$).
\end{prob}

\begin{prob}[L]{4408}
证明公式：
\begin{enumerate}[label=(\alph*), itemsep=2pt, topsep=3pt]
  \item $\operatorname{grad}(u + c) = \operatorname{grad} u$ ($c$ 为常数)；
  \item $\operatorname{grad} cu = c \operatorname{grad} u$ ($c$ 为常数)；
  \item $\operatorname{grad}(u + v) = \operatorname{grad} u + \operatorname{grad} v$；
  \item $\operatorname{grad}(uv) = u \operatorname{grad} v + v \operatorname{grad} u$；
  \item $\operatorname{grad}(u^2) = 2u \operatorname{grad} u$；
  \item $\operatorname{grad} f(u) = f'(u) \operatorname{grad} u$.
\end{enumerate}
\end{prob}

\begin{prob}[L]{4409}
计算：
\begin{enumerate}[label=(\alph*), itemsep=2pt, topsep=3pt]
  \item $\operatorname{grad} r$；
  \item $\operatorname{grad} r^2$；
  \item $\operatorname{grad} \dfrac{1}{r}$，其中 $r = \sqrt{x^2 + y^2 + z^2}$.
\end{enumerate}
\end{prob}

\begin{prob}[S]{4410}
求 $\operatorname{grad} f(r)$，其中 $r = \sqrt{x^2 + y^2 + z^2}$.
\end{prob}

\begin{prob}[S]{4411}
求 $\operatorname{grad}(\boldsymbol{c} \cdot \boldsymbol{r})$，其中 $\boldsymbol{c}$ 为常向量，$\boldsymbol{r}$ 为引自坐标原点的径向量.
\end{prob}

\begin{prob}[S]{4412}
求 $\operatorname{grad}\{|\boldsymbol{c} \times \boldsymbol{r}|^2\}$ ($\boldsymbol{c}$ 为常向量).
\end{prob}

\begin{prob}[M]{4413}
证明公式：
\[
\operatorname{grad} f(u, v) = \frac{\partial f}{\partial u} \operatorname{grad} u + \frac{\partial f}{\partial v} \operatorname{grad} v.
\]
\end{prob}

\begin{prob}[M]{4414.1}
证明公式：
\[
\begin{gathered}
\nabla^2(uv) = u \nabla^2 v + v \nabla^2 u + 2 \nabla u \nabla v, \\
\nabla = \boldsymbol{i}\frac{\partial}{\partial x} + \boldsymbol{j}\frac{\partial}{\partial y} + \boldsymbol{k}\frac{\partial}{\partial z}, \\
\nabla^2 = \nabla \cdot \nabla = \frac{\partial^2}{\partial x^2} + \frac{\partial^2}{\partial y^2} + \frac{\partial^2}{\partial z^2}.
\end{gathered}
\]
\end{prob}

\begin{prob}[M]{4414.2}
证明：若函数 $u = u(x, y, z)$ 在凸区域 $\Omega$ 内可微，且 $|\operatorname{grad} u| \leqslant M$，其中 $M$ 为常数，则对于 $\Omega$ 中任意二点 $A, B$ 有：
\[
|u(A) - u(B)| \leqslant M \rho(A, B),
\]
式中 $\rho(A, B)$ 为 $A$ 与 $B$ 二点之间的距离.
\end{prob}

\begin{prob}[M]{4415}
对于函数 $u = u(x, y, z)$，写出 $\operatorname{grad} u$ 在下列坐标系中的表达式：(a) 圆柱坐标系；\qquad (b) 球坐标系.
\end{prob}

\begin{prob}[M]{4416}
求场 $u = \dfrac{x^2}{a^2} + \dfrac{y^2}{b^2} + \dfrac{z^2}{c^2}$ 在已知点 $M(x, y, z)$ 处沿此点的径向量 $\boldsymbol{r}$ 之方向的导数. 在怎样的情况下，此导数等于梯度的大小？
\end{prob}

\begin{prob}[M]{4417}
求场 $u = \dfrac{1}{r}$（其中 $r = \sqrt{x^2 + y^2 + z^2}$）在方向 $\boldsymbol{l}\{\cos\alpha, \cos\beta, \cos\gamma\}$ 上的导数. 在怎样的情况下，此导数等于零？
\end{prob}

\begin{prob}[M]{4418}
求场 $u = u(x, y, z)$ 在场 $v = v(x, y, z)$ 的梯度方向的导数. 在怎样的情况下，此导数等于零？
\end{prob}

\begin{prob}[M]{4419}
设
\[
u = \arctan \frac{z}{\sqrt{x^2 + y^2}}, \quad \boldsymbol{c} = \boldsymbol{i} + \boldsymbol{j} + \boldsymbol{k},
\]
写出
\[
\boldsymbol{a} = \boldsymbol{c} \times \operatorname{grad} u
\]
通过单位向量 $\boldsymbol{i}, \boldsymbol{j}, \boldsymbol{k}$ 的表达式.
\end{prob}

\begin{prob}[M]{4420}
确定向量场
\[
\boldsymbol{a} = x\boldsymbol{i} + y\boldsymbol{j} + 2z\boldsymbol{k}
\]
的向量线.
\end{prob}

\begin{prob}[M]{4421}
用直接计算的方法证明：向量 $\boldsymbol{a}$ 的散度与直角坐标系的选择无关.
\end{prob}

\begin{prob}[M]{4422}
证明：
\[
\operatorname{div} \boldsymbol{a}(M) = \lim_{d(S) \to 0} \frac{1}{V} \iint_S a_n \,\mathrm{d}S,
\]
其中 $S$ 表示围绕点 $M$ 的封闭曲面，$V$ 是该曲面所围区域的体积，$\boldsymbol{n}$ 为曲面 $S$ 的外法向量，$d(S)$ 为曲面 $S$ 的直径.
\end{prob}

\begin{prob}[M]{4423.1}
求场 $\boldsymbol{a} = \dfrac{-\boldsymbol{i}x + \boldsymbol{j}y + \boldsymbol{k}z}{\sqrt{x^2 + y^2}}$ 在点 $M(3, 4, 5)$ 的散度. 向量 $\boldsymbol{a}$ 通过无穷小球面 $(x - 3)^2 + (y - 4)^2 + (z - 5)^2 = \varepsilon^2$ 的通量 $\Pi$ 近似地等于什么？
\end{prob}

\begin{prob}[S]{4423.2}
求
\[
\operatorname{div} \begin{vmatrix}
\boldsymbol{i} & \boldsymbol{j} & \boldsymbol{k} \\
\dfrac{\partial}{\partial x} & \dfrac{\partial}{\partial y} & \dfrac{\partial}{\partial z} \\
\omega_x & \omega_y & \omega_z
\end{vmatrix}.
\]
\end{prob}

\begin{prob}[L]{4424}
证明：
\begin{enumerate}[label=(\alph*), itemsep=2pt, topsep=3pt]
  \item $\operatorname{div}(\boldsymbol{a} + \boldsymbol{b}) = \operatorname{div} \boldsymbol{a} + \operatorname{div} \boldsymbol{b}$；
  \item $\operatorname{div}(u\boldsymbol{c}) = \boldsymbol{c} \operatorname{grad} u$ ($\boldsymbol{c}$ 为常向量，$u$ 为标量)；
  \item $\operatorname{div}(u\boldsymbol{a}) = u \operatorname{div} \boldsymbol{a} + \boldsymbol{a} \cdot \operatorname{grad} u$.
\end{enumerate}
\end{prob}

\begin{prob}[S]{4425}
求 $\operatorname{div}(\operatorname{grad} u)$.
\end{prob}

\begin{prob}[M]{4426}
求 $\operatorname{div}[\operatorname{grad} f(r)]$，其中 $r = \sqrt{x^2 + y^2 + z^2}$. 在什么情况下 $\operatorname{div}[\operatorname{grad} f(r)] = 0$？
\end{prob}

\begin{prob}[M]{4427}
计算：(a) $\operatorname{div} \boldsymbol{r}$； \qquad (b) $\operatorname{div} \dfrac{\boldsymbol{r}}{r}$.
\end{prob}

\begin{prob}[S]{4428}
计算：$\operatorname{div}[f(r)\boldsymbol{c}]$，式中 $\boldsymbol{c}$ 为常向量.
\end{prob}

\begin{prob}[M]{4429}
求 $\operatorname{div}[f(r)\boldsymbol{r}]$. 在什么情形下该散度等于零？
\end{prob}

\begin{prob}[M]{4430}
求：(a) $\operatorname{div}(u \operatorname{grad} u)$； \qquad (b) $\operatorname{div}(u \operatorname{grad} v)$.
\end{prob}

\begin{prob}[M]{4431}
设流体充满空间并以恒定的角速度 $\omega$ 沿逆时针方向绕 $Oz$ 轴旋转，求速度向量 $\boldsymbol{v}$ 和加速度向量 $\boldsymbol{w}$ 在已知时刻在空间点 $M(x, y, z)$ 处的散度.
\end{prob}

\begin{prob}[M]{4432}
求包含多个引力中心的有限系统所产生的引力场之散度.
\end{prob}

\begin{prob}[M]{4433}
求由极坐标 $r$ 与 $\varphi$ 表示的平面向量 $\boldsymbol{a} = \boldsymbol{a}(r, \varphi)$ 之散度的表达式.
\end{prob}

\begin{prob}[L]{4434}
设
\[
x = f(u, v, w), \quad y = g(u, v, w), \quad z = h(u, v, w),
\]
用正交曲线坐标 $u, v, w$ 表示 $\operatorname{div} \boldsymbol{a}(x, y, z)$.

作为特殊的情形，求用柱坐标和球坐标表示 $\operatorname{div} \boldsymbol{a}$ 的表达式.
\probhint{研究向量 $\boldsymbol{a}$ 通过以曲面 $u = \text{常数}$，$v = \text{常数}$，$w = \text{常数}$ 为界的无穷小平行六面体表面的通量.}
\end{prob}

\begin{prob}[M]{4435}
证明：
(a) $\operatorname{rot}(\boldsymbol{a} + \boldsymbol{b}) = \operatorname{rot} \boldsymbol{a} + \operatorname{rot} \boldsymbol{b}$； \qquad (b) $\operatorname{rot}(u\boldsymbol{a}) = u \operatorname{rot} \boldsymbol{a} + \operatorname{grad} u \times \boldsymbol{a}$.
\end{prob}

\begin{prob}[M]{4436.1}
求：(a) $\operatorname{rot} \boldsymbol{r}$； \qquad (b) $\operatorname{rot}[f(r)\boldsymbol{r}]$.
\end{prob}

\begin{prob}[M]{4436.2}
设
\[
\boldsymbol{a} = \frac{y}{z}\boldsymbol{i} + \frac{z}{x}\boldsymbol{j} + \frac{x}{y}\boldsymbol{k},
\]
求 $\operatorname{rot} \boldsymbol{a}$ 在点 $M(1, 2, -2)$ 的大小和方向.
\end{prob}

\begin{prob}[M]{4437}
求：(a) $\operatorname{rot}[\boldsymbol{c} f(r)]$； \qquad (b) $\operatorname{rot}[\boldsymbol{c} \times f(r)\boldsymbol{r}]$ ($\boldsymbol{c}$ 为常向量).
\end{prob}

\begin{prob}[M]{4438}
证明：$\operatorname{div}(\boldsymbol{a} \times \boldsymbol{b}) = \boldsymbol{b} \cdot \operatorname{rot} \boldsymbol{a} - \boldsymbol{a} \cdot \operatorname{rot} \boldsymbol{b}$.
\end{prob}

\begin{prob}[M]{4439}
求：(a) $\operatorname{rot}(\operatorname{grad} u)$； \qquad (b) $\operatorname{div}(\operatorname{rot} \boldsymbol{a})$.
\end{prob}

\begin{prob}[M]{4440}
设流体充满空间并以恒定的角速度 $\omega$ 围绕轴 $\boldsymbol{l}\{\cos\alpha, \cos\beta, \cos\gamma\}$ 旋转. 求速度向量 $\boldsymbol{v}$ 在已知时刻在空间点 $M(x, y, z)$ 处的旋度.
\end{prob}

\begin{prob}[M]{4441.1}
求由极坐标 $r$ 和 $\varphi$ 表示的平面向量 $\boldsymbol{a} = \boldsymbol{a}(r, \varphi)$ 的旋度的表达式.
\end{prob}

\begin{prob}[M]{4441.2}
写出 $\operatorname{rot} \boldsymbol{a}(x, y, z)$ 在下列坐标系中的表达式：
(a) 圆柱坐标系； \qquad (b) 球坐标系.
\end{prob}

\begin{prob}[M]{4442.1}
求径向量 $\boldsymbol{r}$ 的通量：(a) 通过圆锥体 $x^2 + y^2 \leqslant z^2$ ($0 \leqslant z \leqslant h$) 的侧表面；(b) 通过此圆锥体的底面.
\end{prob}

\begin{prob}[M]{4442.2}
求向量 $\boldsymbol{a} = \boldsymbol{i}yz + \boldsymbol{j}xz + \boldsymbol{k}xy$ 的通量：(a) 通过圆柱体 $x^2 + y^2 \leqslant a^2$ ($0 \leqslant z \leqslant h$) 的侧表面；(b) 通过此圆柱的全表面.
\end{prob}

\begin{prob}[M]{4443}
求径向量 $\boldsymbol{r}$ 通过曲面
\[
z = 1 - \sqrt{x^2 + y^2} \quad (0 \leqslant z \leqslant 1)
\]
的通量.
\end{prob}

\begin{prob}[M]{4444}
求向量 $\boldsymbol{a} = x^2\boldsymbol{i} + y^2\boldsymbol{j} + z^2\boldsymbol{k}$ 通过正八分之一球面 $x^2 + y^2 + z^2 = 1, x \geqslant 0, y \geqslant 0, z \geqslant 0$ 的通量.
\end{prob}

\begin{prob}[M]{4445.1}
求向量 $\boldsymbol{a} = y\boldsymbol{i} + z\boldsymbol{j} + x\boldsymbol{k}$ 通过以平面 $x = 0, y = 0, z = 0, x + y + z = a$ ($a > 0$) 为界的角锥的全表面的通量. 利用奥斯特罗格拉茨基公式验证结果.
\end{prob}

\begin{prob}[M]{4445.2}
求向量 $\boldsymbol{a} = x^3\boldsymbol{i} + y^3\boldsymbol{j} + z^3\boldsymbol{k}$ 通过球面 $x^2 + y^2 + z^2 = x$ 的通量.
\end{prob}

\begin{prob}[M]{4446}
证明：向量 $\boldsymbol{a}$ 通过由方程 $\boldsymbol{r} = \boldsymbol{r}(u, v)$ ($(u, v) \in \Omega$) 给出的曲面 $S$ 的通量等于
\[
\iint_S a_n \,\mathrm{d}S = \iint_\Omega \left( \boldsymbol{a} \frac{\partial \boldsymbol{r}}{\partial u} \frac{\partial \boldsymbol{r}}{\partial v} \right) \mathrm{d}u\,\mathrm{d}v,
\]
式中 $a_n = \boldsymbol{a} \cdot \boldsymbol{n}$，$\boldsymbol{n}$ 为曲面 $S$ 的单位法向量.
\end{prob}

\begin{prob}[M]{4447}
求向量 $\boldsymbol{a} = m \dfrac{\boldsymbol{r}}{r^3}$ ($m$ 为常数) 通过围绕坐标原点的封闭曲面 $S$ 的通量.
\end{prob}

\begin{prob}[M]{4448}
已知向量
\[
\boldsymbol{a}(\boldsymbol{r}) = \sum_{i=1}^n \operatorname{grad}\left( -\frac{e_i}{4\pi r_i} \right),
\]
其中 $e_i$ 为常数，$r_i$ 为点 $M_i$（点源）距动点 $M(\boldsymbol{r})$ 的距离. 求此向量通过围绕点 $M_i$ ($i = 1, 2, \cdots, n$) 的封闭曲面 $S$ 的通量.
\end{prob}

\begin{prob}[M]{4449}
证明：
\[
\iint_S \frac{\partial u}{\partial n} \,\mathrm{d}S = \iiint_V \nabla^2 u \,\mathrm{d}x\,\mathrm{d}y\,\mathrm{d}z,
\]
其中曲面 $S$ 是区域 $V$ 的边界.
\end{prob}

\begin{prob}[M]{4450}
在温度场 $u$ 内，在单位时间内通过面微元 $\mathrm{d}S$ 的热量等于：
\[
\mathrm{d}Q = -k \boldsymbol{n} \cdot \operatorname{grad} u \,\mathrm{d}S,
\]
其中 $k$ 为热导率，$\boldsymbol{n}$ 为曲面 $S$ 的单位法向量，求在单位时间内物体 $V$ 所吸收的热量. 研究温度上升的速度，从而推出物体温度所满足的方程（热传导方程）.
\end{prob}

\begin{prob}[M]{4451}
处于运动状态的不可压缩流体充满区域 $V$. 假定在区域 $V$ 内没有源和汇，试推出连续性方程：
\[
\frac{\partial \rho}{\partial t} + \operatorname{div}(\rho \boldsymbol{v}) = 0,
\]
式中 $\rho = \rho(x, y, z, t)$ 为流体密度，$\boldsymbol{v}$ 为速度向量，$t$ 为时间.
\probhint{研究流体通过包含在 $V$ 中的任意区域 $\omega$ 的流量.}
\end{prob}

\begin{prob}[M]{4452.1}
求向量 $\boldsymbol{a} = \boldsymbol{r}$ 沿着一段螺旋线
\[
\boldsymbol{r} = \boldsymbol{i}a\cos t + \boldsymbol{j}a\sin t + \boldsymbol{k}bt \quad (0 \leqslant t \leqslant 2\pi)
\]
的环量.
\end{prob}

\begin{prob}[M]{4452.2}
求场 $\boldsymbol{a} = \dfrac{1}{y}\boldsymbol{i} + \dfrac{1}{z}\boldsymbol{j} + \dfrac{1}{x}\boldsymbol{k}$ 沿 $M(1, 1, 1)$ 和 $N(2, 4, 8)$ 这两点之间的线段的环量.
\end{prob}

\begin{prob}[M]{4452.3}
求场 $\boldsymbol{a} = \boldsymbol{i}\mathrm{e}^{y-z} + \boldsymbol{j}\mathrm{e}^{z-x} + \boldsymbol{k}\mathrm{e}^{x-y}$ 沿 $O(0, 0, 0)$ 和 $M(1, 3, 5)$ 这两点之间的线段的环量.
\end{prob}

\begin{prob}[M]{4453.1}
求场 $\boldsymbol{a} = (y + z)\boldsymbol{i} + (z + x)\boldsymbol{j} + (x + y)\boldsymbol{k}$ 沿球面 $x^2 + y^2 + z^2 = 25$ 上连接 $M(3, 4, 0)$ 和 $N(0, 0, 5)$ 这两点的最短的一段大圆圆弧的环量.
\end{prob}

\begin{prob}[S]{4453.2}
求向量 $\boldsymbol{a} = f(r)\boldsymbol{r}$（其中 $f$ 为连续函数）沿着弧 $AB$ 的环量.
\end{prob}

\begin{prob}[M]{4454.1}
求向量 $\boldsymbol{a} = -y\boldsymbol{i} + x\boldsymbol{j} + c\boldsymbol{k}$ ($c$ 为常数) 的环量：(a) 沿着圆周 $x^2 + y^2 = 1, z = 0$；(b) 沿着圆周 $(x - 2)^2 + y^2 = 1, z = 0$.
\end{prob}

\begin{prob}[M]{4454.2}
求向量 $\boldsymbol{a} = \operatorname{grad}\left(\arctan\dfrac{y}{x}\right)$ 沿着围线 $C$ 的环量 $\Gamma$：
(a) $C$ 不围绕 $Oz$ 轴； \qquad (b) $C$ 围绕 $Oz$ 轴.
\end{prob}

\begin{prob}[M]{4455}
已知向量场
\[
\boldsymbol{a} = \frac{y}{\sqrt{x}}\boldsymbol{i} - \frac{x}{\sqrt{z}}\boldsymbol{j} + \sqrt{xyz}\boldsymbol{k}.
\]
计算 $\operatorname{rot} \boldsymbol{a}$ 在点 $M(1, 1, 1)$ 的表达式，然后近似地求出向量场沿无穷小圆周
\[
\begin{cases}
(x - 1)^2 + (y - 1)^2 + (z - 1)^2 = \varepsilon^2, \\
(x - 1)\cos\alpha + (y - 1)\cos\beta + (z - 1)\cos\gamma = 0
\end{cases}
\]
（其中 $\cos^2\alpha + \cos^2\beta + \cos^2\gamma = 1$）的环量 $\Gamma$.
\end{prob}

\begin{prob}[L]{4456}
平面定常流由速度向量
\[
\boldsymbol{w} = u(x, y)\boldsymbol{i} + v(x, y)\boldsymbol{j}
\]
描述，试确定：(1) 通过区域 $S$ 的边界（封闭围线）$C$ 流出的流体的量 $Q$（流量）；(2) 速度向量沿着围线 $C$ 的环量 $\Gamma$. 若流体是不可压缩的且流场无旋，则函数 $u$ 和 $v$ 满足怎样的方程？
\end{prob}

\begin{prob}[M]{4457.1}
证明：场
\[
\boldsymbol{a} = yz(2x+y+z)\boldsymbol{i} + xz(x+2y+z)\boldsymbol{j} + xy(x+y+2z)\boldsymbol{k}
\]
是有势场，并求这个场的势.
\end{prob}

\begin{prob}[M]{4457.2}
确认场
\[
\boldsymbol{a} = \frac{2}{(y+z)^{\frac{1}{2}}}\boldsymbol{i} - \frac{x}{(y+z)^{\frac{3}{2}}}\boldsymbol{j} - \frac{x}{(y+z)^{\frac{3}{2}}}\boldsymbol{k}
\]
是有势场，并求该场沿着正八分之一球面上连接点 $M(1,1,3)$ 和 $N(2,4,5)$ 的路径的环量.
\end{prob}

\begin{prob}[S]{4458}
求由位于坐标原点的质量 $m$ 所产生的引力场 $\boldsymbol{a} = -\frac{m}{r^3}\boldsymbol{r}$ 的势.
\end{prob}

\begin{prob}[S]{4459}
设质量 $m_i\ (i=1,2,\cdots,n)$ 的质点位于点 $M_i\ (i=1,2,\cdots,n)$，求此质点系所产生的引力场的势.
\end{prob}

\begin{prob}[M]{4460}
证明：场 $\boldsymbol{a} = f(r)\boldsymbol{r}$（其中 $f(r)$ 是单值连续函数）是有势场. 求这个场的势.
\end{prob}

\begin{prob}[M]{4461}
证明公式：
\[
\operatorname{grad}_P \left\{ \iiint_V \rho(Q) \frac{\mathrm{d}V}{r} \right\} = -\iint_S \rho(Q)\boldsymbol{n}\frac{\mathrm{d}S}{r} + \iiint_V \operatorname{grad}_Q \rho(Q)\frac{\mathrm{d}V}{r},
\]
其中 $S$ 为区域 $V$ 的边界曲面，$\boldsymbol{n}$ 为曲面 $S$ 的外法向量，$r$ 为点 $P(x,y,z)$ 与点 $Q(\xi,\eta,\zeta)$ 之间的距离.
\end{prob}

\begin{prob}[M]{4462}
证明：若 $\boldsymbol{a} = \operatorname{grad} u$，其中
\begin{gather*}
u(x,y,z) = -\frac{1}{4\pi} \iiint\limits_{-\infty}^{+\infty} \frac{\rho(\xi,\eta,\zeta)}{r} \,\mathrm{d}\xi\,\mathrm{d}\eta\,\mathrm{d}\zeta, \\
r = \sqrt{(\xi-x)^2 + (\eta-y)^2 + (\zeta-z)^2},
\end{gather*}
则
\[
\operatorname{div}\boldsymbol{a} = \rho(x,y,z)
\]
（假设相应积分有意义）.
\end{prob}


\end{document}
